 \documentclass[11pt,reqno]{amsbook}
 \usepackage{cmap}
 \usepackage[T1]{fontenc}
 \usepackage[utf8]{inputenc}
 \usepackage{lmodern}
 \usepackage[indonesian]{babel}
 \usepackage{amssymb}
 \usepackage{amscd}
 \usepackage[all]{xy}
 \usepackage{url}
 \usepackage{enumitem}
 \usepackage{tabularx}
 \usepackage{xcolor}
 \definecolor{AccessibleLinkBlue}{HTML}{005A9C}
 \usepackage[pdftex,
            pdfauthor={John M Erdman},
            pdftitle={Analisis Fungsional dan Aljabar Operator: Suatu Pengantar},
            pdfsubject={analisis fungsional dan aljabar operator},
            pdfkeywords={aljabar linear, analisis fungsional, operator, ruang Hilbert, ruang Banach, spektrum, operator kompak, teorema spektral, ruang vektor topologis, ruang Frechet, distribusi, ruang LF, konvolusi, transformasi Fourier, aljabar C bintang, transformasi Gelfand, unitalisasi, barisan eksak, kuasi-invers, elemen positif, identitas aproksimatif, konstruksi GNS, keadaan, representasi, modul Hilbert, ideal esensial, kompaktifikasi, aljabar pengali, teori Fredholm, operator Fredholm, aljabar Calkin, indeks Fredholm, ekstensi, ekstensi semiterbelah, operator Toeplitz, spektrum esensial, pemetaan positif lengkap, aljabar nuklir, teori-K, funktor K nol, proyeksi, ekuivalensi Murray--von Neumann, grup Grothendieck, aljabar-AF, diagram Bratteli, notasi, fungsi, huruf Yunani, huruf Fraktur},
            pdfcreator={OpenAI Codex gpt-5.6-sol, Ultra},
            pdfcreationdate={D:20260824000000+02'00'},
            pdfmoddate={D:20260824000000+02'00'},
            pdflang={id-ID},pdfdisplaydoctitle=true,
            plainpages=false,pdftex,pagebackref=true,colorlinks,
            allcolors=AccessibleLinkBlue,
            bookmarks=true,bookmarksopen=true]{hyperref}
     %plainpages=false corrects page numbers in index
     %pagebackref=true provides back references from bibliography
     %the translated wrapper restores reader navigation bookmarks







\makeindex


 \setlength{\textwidth}{6.5in}
 \setlength{\oddsidemargin}{0in}
 \setlength{\evensidemargin}{0in}
 \setlength{\textheight}{9.25in}
 \setlength{\topmargin}{-0.35in}

 % Michael Barr's DIAGXY.TEX is preserved byte-for-byte under its embedded
 % redistribution notice.  Exact case keeps the build portable.
 \input{DIAGXY.TEX}

%\includeonly{}


\swapnumbers \theoremstyle{plain}
\newtheorem{thm}{Teorema}[section]
\newtheorem{lem}[thm]{Lema}
\newtheorem{prop}[thm]{Proposisi}
\newtheorem{cor}[thm]{Akibat}
\newtheorem{ax}[thm]{Aksioma}

\theoremstyle{definition}
\newtheorem{defn}[thm]{Definisi}
\newtheorem{notn}[thm]{Notasi}
\newtheorem{conv}[thm]{Konvensi}

\theoremstyle{remark}
\newtheorem{rem}[thm]{Catatan}
\newtheorem*{cau}{PERHATIAN}
\newtheorem{fact}[thm]{Fakta}
\newtheorem{exam}[thm]{Contoh}
\newtheorem{exer}[thm]{Latihan}
\newtheorem{prob}[thm]{Soal}
\newtheorem{proj}[thm]{Proyek}

\renewcommand{\proofname}{Bukti}
\renewcommand{\chaptername}{Bab}
\renewcommand{\contentsname}{Daftar Isi}
\renewcommand{\bibname}{Daftar Pustaka}
\renewcommand{\indexname}{Indeks}
\renewcommand{\emailaddrname}{{\itshape Alamat surel}}
\renewcommand{\seeonlyname}{lihat}
\renewcommand{\seename}{lihat juga}
\renewcommand{\alsoname}{lihat juga}



\renewcommand{\thechapter}{\arabic{chapter}}
\renewcommand{\thesection}{\thechapter.\arabic{section}}
\renewcommand{\thesubsection}{\thesection.\arabic{subsection}}


\renewcommand{\labelenumi}{\rm(\alph{enumi})}
%to get enumeration in lower case roman




\numberwithin{equation}{section}





 \newcommand{\field}[1]{\mathbb{#1}}
 \newcommand{\C}{\field{C}}
 \newcommand{\Di}{\field{D}}
 \newcommand{\Dsk}{\field{D}}
 \newcommand{\E}{\field{E}}
 \newcommand{\K}{\field{K}}
 \newcommand{\N}{\field{N}}
 \newcommand{\Po}{\field{P}}
 \newcommand{\Q}{\field{Q}}
 \newcommand{\R}{\field{R}}
 \newcommand{\Sp}{\field{S}}
 \newcommand{\T}{\field{T}}
 \newcommand{\Z}{\field{Z}}



 \newcommand{\sto}{\rightarrow}
      %Diagxy redefined \to . This shortens the arrow.
 \newcommand{\ofml}[1]{\mathfrak{#1}}
      % for families of operators
 \newcommand{\fml}[1]{\mathcal{#1}}
      % for families of functions
 \newcommand{\sfml}[1]{\mathfrak{#1}}
      % for families of sets
 \newcommand{\ftr}[1]{\mathsf{#1}}
      % for functors
 \newcommand{\vc}[1]{\mathbf{#1}}
      % for vectors
 \newcommand{\df}[1]{\textsc{#1}}
      % for terms being defined
 \newcommand{\bs}[1]{\boldsymbol{#1}}
      % to create bold symbols and Greek letters
 \newcommand{\D}{\displaystyle}
 \newcommand{\ns}{\renewcommand{\qed}{}}
 \newcommand{\abs}[1]{\lvert#1\rvert}
      % absolute value
 \newcommand{\bigabs}[1]{\left\lvert#1\right\rvert}
      % large absolute value
  \newcommand{\norm}[1]{\lVert#1\rVert}
      % norm
 \newcommand{\bignorm}[1]{\bigl\lVert#1\bigr\rVert}
      % large norm
 \newcommand{\biggnorm}[1]{\biggl\lVert#1\biggr\rVert}
      % larger norm
 \newcommand{\trinorm}[1]{\vert\mspace{-2mu}\vert
     \mspace{-2mu}\vert#1\vert\mspace{-2mu}\vert
     \mspace{-2mu}\vert}
      % 3-bar norm
 \newcommand{\cstariso}[2]{#1\overset{\,{}_\ast}{\cong}#2}
      % #1 is *-isomorphic to #2
 \newcommand{\cat}[1]{\mathbf{#1}}
      % for categories
 \newcommand{\clo}[1]{\overline{#1}}
      % closure
 \newcommand{\cloco}[1]{\overline{\operatorname{co}}(#1)}
      % closed convex hull
 \newcommand{\co}[1]{\operatorname{co}(#1)}
      % convex hull
 \newcommand{\conj}[1]{\overline{#1}}
      % complex conjugate
 \newcommand{\id}[1]{\operatorname{id}_{#1}}
      % Identity on #1
 \newcommand{\intr}[1]{{#1}^{\circ}}
      % interior
  \newcommand{\lfs}[1]{\mathrm{L_{#1}}}
      % Lebesgue function spaces
 \newcommand{\locint}[1]{\mathrm{L_1^{loc}}({#1})}
      % locally integrable functions on #1
 \newcommand{\M}[2]{\mathbf{M}_{#1}({#2})}
      % for matrix algebras
 \newcommand{\mor}[2]{\operatorname{\mathfrak{Mor}}(#1,#2)}
       % Morphisms from #1 to #2
 \newcommand{\open}[2]{#1\overset{\thickspace{}_\circ}{\subseteq}#2}
      % #1 is an open subset of #2
 \newcommand{\sbsb}[2]{#1_{{}_\sst{#2}}}
      % makes smaller, lower subscript
 \newcommand{\spsp}[2]{#1^{{}^\sst{#2}}}
      % makes smaller higher superscript
 \newcommand{\sst}[1]{{\scriptstyle{#1}}}
      % sub- and superscript size type
 \newcommand{\ssst}[1]{{\scriptscriptstyle{#1}}}
      % second level sub- and superscripts
 \newcommand{\pd}[2]{\dfrac{\partial#1}{\partial#2}}
      % first partial, 1 variable
 \newcommand{\ppd}[2]{\dfrac{\partial^2#1}{\partial#2^2}}
      % second partial, 1 variable
 \newcommand{\ppdd}[3]{\dfrac{\partial^2#1}{\partial#2\partial#3}}
      % second partial, 2 variables
 \newcommand{\wt}[1]{\widetilde{#1}}
 \newcommand{\wh}[1]{\widehat{#1}}
 \makeatletter
 \newcommand{\futurexref}[2]{%
   \@ifundefined{r@#2}{#1}{\ref{#2}}}
 \makeatother
      % Use a live link when the labelled destination has entered this
      % cumulative reader; otherwise retain the printed source number.



 \DeclareMathOperator{\ad}{Ad}
 \DeclareMathOperator{\ba}{ba}
 \DeclareMathOperator{\ca}{ca}
 \DeclareMathOperator{\card}{card}
 \DeclareMathOperator{\cloin}{cl}
 \DeclareMathOperator{\codim}{codim}
 \DeclareMathOperator{\coker}{coker}
 \DeclareMathOperator{\diag}{diag}
 \DeclareMathOperator{\diam}{diam}
 \DeclareMathOperator{\dist}{dist}
 \DeclareMathOperator{\dom}{dom}
 \DeclareMathOperator{\essran}{essran}
 \DeclareMathOperator{\ext}{Ext}
 \DeclareMathOperator{\fin}{Fin}
 \DeclareMathOperator{\Hom}{Hom}
 \DeclareMathOperator{\im}{im}
 \DeclareMathOperator{\ind}{ind}
 \DeclareMathOperator{\intrin}{int}
 \DeclareMathOperator{\inv}{inv}
 \DeclareMathOperator{\mx}{Max}
 \DeclareMathOperator{\ran}{ran}
 \DeclareMathOperator{\rca}{rca}
 \DeclareMathOperator{\spn}{span}
 \DeclareMathOperator{\supp}{supp}
 \DeclareMathOperator{\tr}{tr}







% Separately authored companion-layer presentation environments.
% Repeated source statements remain exact in component files, while duplicate
% labels and index hooks are suppressed locally in the compiled companion copy.
\makeatletter
\newenvironment{o001solution}[3]{%
  \par\bigskip\hrule\medskip
  \phantomsection\hypertarget{#1}{}%
  \noindent{\footnotesize\textbf{ID solusi:} \path{#1}\\
  \textbf{Latihan sumber:} \path{#2}}\par\smallskip
}{\par\medskip\hrule\bigskip}
\newenvironment{o001statement}{%
  \begingroup\let\label\@gobble\let\label@in@display\@gobble\let\index\@gobble
  \begin{quote}\small\noindent\textbf{Pernyataan latihan sumber.}\par\smallskip
}{\end{quote}\endgroup}
\newenvironment{o001answer}{%
  \begin{quote}\noindent\textbf{Jawaban ringkas.}\par\smallskip
}{\end{quote}}
\newenvironment{o001proof}{%
  \begin{quote}\noindent\textbf{Solusi lengkap.}\par\smallskip
}{\hfill\ensuremath{\blacksquare}\end{quote}}
\newenvironment{o001readerwork}[5]{%
  \par\bigskip\hrule\medskip
  \phantomsection\hypertarget{#1}{}%
  \noindent{\footnotesize\textbf{ID dukungan:} \path{#1}\\
  \textbf{Hasil sumber:} \path{#2}; \textbf{petunjuk sumber:} \path{#3}}%
  \par\smallskip
}{\par\medskip\hrule\bigskip}
\newenvironment{o001result}{%
  \begingroup\let\label\@gobble\let\label@in@display\@gobble\let\index\@gobble
  \begin{quote}\small\noindent\textbf{Hasil sumber terpilih.}\par\smallskip
}{\end{quote}\endgroup}
\newenvironment{o001sourcehint}{%
  \begingroup\let\label\@gobble\let\label@in@display\@gobble\let\index\@gobble
  \begin{quote}\small\noindent\textbf{Petunjuk sumber.}\par\smallskip
}{\end{quote}\endgroup}
\makeatother

\begin{document}


 \frontmatter
 \title{ANALISIS FUNGSIONAL \\ DAN \\ ALJABAR OPERATOR: \\ Suatu Pengantar \\ \mbox{} \\ \large Edisi Lengkap Bahasa Indonesia dengan Pendamping Penguasaan}
 \author{John M. Erdman \\
      Portland State University \\
      \mbox{\hphantom{P}} \\
      Versi sumber 4 Oktober 2015 \\
      \mbox{\hphantom{P}} \\
    %  \copyright 2014 John M. Erdman \\
      \mbox{\hphantom{P}} \\
            \mbox{\hphantom{P}} \\
                  \mbox{\hphantom{P}} }


\vskip 2 in

 \email{erdman@pdx.edu}

\maketitle


\begin{center}
\textbf{Edisi Bahasa Indonesia --- teks sumber lengkap}
\end{center}

Karya sumber John M. Erdman ini berlisensi Creative Commons
Attribution--ShareAlike 4.0 International (CC BY-SA 4.0):
\url{https://creativecommons.org/licenses/by-sa/4.0/}.
Terjemahan Bahasa Indonesia dan adaptasi teknis ini juga diterbitkan dengan
lisensi CC BY-SA 4.0. Penerjemahan, penyuntingan teknis, pengindeksan modular,
build, dan QA edisi ini dibantu oleh OpenAI Codex gpt-5.6-sol, Ultra atas arahan
pengguna; kredit John M. Erdman dan kontributor komponen tetap dipertahankan.
Edisi ini tidak
disponsori atau didukung oleh John M. Erdman maupun Portland State University.

\begin{sloppypar}
Sumber resmi:
\url{https://web.pdx.edu/~erdman/FAOA/functional_analysis_operator_algebras_pdf.pdf}
dan
\url{https://web.pdx.edu/~erdman/FAOA/functional_analysis_operator_algebras_web.zip}.
SHA-256 sumber masing-masing adalah
\texttt{f320b16af7448fbb43582c21569840fe\allowbreak657fccf6f31d97f176913fdd0e1eb823}
dan
\texttt{0c667cfa7420b61dda8f8cb4ed9d619d\allowbreak b8abbd1b53d17eafe7d4a2e153342e53}.
\end{sloppypar}

Perubahan pada batas ini: prakata dan semua prosa pembaca dalam Bab 1 sampai
Bab 17 diterjemahkan ke Bahasa Indonesia; struktur, rumus, label, sitasi,
latihan, daftar pustaka, indeks, dan penomoran sumber dipertahankan, dengan
koreksi sumber yang dicatat secara terpisah. Epigraf pihak ketiga, gambar lencana lisensi, dan makro tabel yang
status komponennya tidak cukup jelas tidak digunakan. Berkas
\texttt{DIAGXY.TEX} dipakai tanpa perubahan menurut pemberitahuan distribusinya
sendiri. Halaman dapat
mengalir ulang pada perangkat TeX modern; nomor halaman bukan pengenal tetap.

Batas ini memuat terjemahan teks sumber lengkap, pendamping HTML semantik,
solusi terpisah untuk seluruh 52 latihan, sepuluh pembuktian penguasaan
terpilih, serta jembatan spektral-kompak/SVD. Semua komponen tambahan
mempertahankan provenans dan hak terpisah dan bukan tulisan Erdman.


\vfill\eject

\tableofcontents


 \chapter*{Prakata}

% SOURCE-CORRECTION: PREFACE-C001
Dalam karyanya yang elok, \emph{Hilbert Space Problem
Book}~\cite{Halmos:1982}, Paul Halmos terkenal menekankan bahwa praktik aktif merupakan
syarat mutlak untuk belajar matematika. Halmos tentu bukan satu-satunya yang meyakini hal
itu. Kumpulan catatan ini merupakan pendamping berorientasi aktivitas
untuk mempelajari analisis fungsional linear dan aljabar operator. Catatan ini dimaksudkan
sebagai pendamping pedagogis bagi pemula, pengantar menuju beberapa gagasan utama dalam
bidang analisis ini, kumpulan soal yang menurut saya berguna untuk mempelajari pokok
bahasannya, serta daftar bacaan dan referensi beranotasi.

Sebagian besar hasil dalam pengantar analisis fungsional bersifat lugas dan dapat
diperiksa oleh mahasiswa yang menelaahnya dengan saksama. Bahkan, itulah tujuan utama
catatan ini---meyakinkan pemula bahwa bidang ini dapat dipelajari. Dalam materi berikut
terdapat banyak penanda yang mengajak pembaca untuk berperan aktif: kata-kata seperti
``proposisi'', ``contoh'', ``latihan'', dan ``akibat'', jika tidak diikuti oleh bukti
atau rujukan kepada suatu bukti, merupakan undangan untuk memeriksa pernyataan yang
diajukan. Tentu saja, ada beberapa teorema yang, menurut saya, memerlukan waktu dan usaha
pembuktian yang lebih besar daripada manfaat yang diperoleh. Dalam kasus seperti itu,
apabila saya tidak dapat menyempurnakan bukti bakunya, saya memberikan rujukan alih-alih
mengulang bukti tersebut.

Materi dalam catatan ini mencerminkan isi sebagian dari dua mata kuliah pascasarjana
sepanjang tahun tentang beberapa topik dasar dalam \emph{analisis fungsional (linear)}
dan \emph{aljabar operator}, yang telah berkali-kali saya ajarkan di Portland State
University. Dalam mata kuliah tersebut, selain menggunakan catatan ini, saya meminta
mahasiswa memilih sumber informasi lain---baik buku cetak maupun dokumen daring---yang
sesuai dengan gaya belajar (dan anggaran) masing-masing. Prasyarat untuk mempelajari
catatan ini cukup ringan; bahkan pengenalan sepintas terhadap \emph{aljabar linear},
\emph{analisis real}, dan \emph{topologi} semestinya memadai. Dengan kata lain, istilah
seperti \emph{ruang vektor}, \emph{pemetaan linear}, \emph{limit}, \emph{ukuran Lebesgue}
dan \emph{integral Lebesgue}, \emph{himpunan terbuka}, \emph{himpunan kompak}, serta
\emph{fungsi kontinu} semestinya sudah tidak asing.

Karena \emph{analisis fungsional linear}, setidaknya dalam suatu pengertian, dapat
dipandang sebagai `aljabar linear berdimensi tak hingga', bab pertama catatan ini
dikhususkan bagi tinjauan (yang cukup ringkas) atas beberapa gagasan penting yang,
  dalam dunia angan-angan saya, dibawa pulang mahasiswa dari mata kuliah
aljabar linear pertama mereka. Tiga wawasan utama yang dikembangkan dalam bab pertama
adalah:
\begin{enumerate}
  \item operator yang dapat didiagonalkan pada ruang vektor berdimensi hingga merupakan
  kombinasi linear dari proyeksi;
  \item operator pada ruang hasil kali dalam kompleks berdimensi hingga dapat didiagonalkan secara
  uniter jika dan hanya jika operator itu normal; dan
  \item dua operator normal pada ruang hasil kali dalam kompleks berdimensi hingga ekuivalen secara
  uniter jika dan hanya jika keduanya memiliki nilai-nilai eigen yang sama, masing-masing
  dengan multiplisitas yang sama. % SOURCE-CORRECTION: PREFACE-C002
\end{enumerate}
Saya memandang hasil-hasil tersebut sangat mendasar karena wawasan luar biasa yang lahir
dari berbagai upaya para ahli teori operator selama seabad terakhir untuk membawanya ke
dalam ranah \emph{analisis fungsional}, yakni menggeneralisasikannya ke kasus berdimensi
tak hingga. Salah satu akibat penjelajahan yang subur dan luar biasa rumit ini ialah
ditemukannya banyak hubungan yang mencengangkan, erat, dan kompleks antara \emph{teori
operator} dan bidang-bidang seperti \emph{teori fungsi variabel kompleks}, \emph{topologi
aljabar}, serta \emph{aljabar homologis}. Dalam bagian-bagian selanjutnya dari catatan ini, saya
berusaha menelusuri kembali beberapa langkah tersebut. Kuliah yang menjadi dasar catatan
ini berpuncak pada hasil menakjubkan Brown, Douglas, dan Fillmore yang pertama kali
diterbitkan pada 1973 (lihat~\cite{BrownDF:1973}). Bahkan sekarang, sebagian besar mahasiswa
pascasarjana masih kesulitan menemukan uraian yang cukup mudah dicerna mengenai banyak
materi latar yang diperlukan untuk mengembangkan materi indah ini. Sayangnya, penalaran
transfinit yang serba informal dan masih cukup berhasil di ruang kuliah sangat sukar dialihkan ke kertas
atau media elektronik. Akibatnya, catatan ini masih berada dalam keadaan yang sungguh
belum selesai. Saya terus mengerjakannya; dan mungkin orang lain kelak akan menyelesaikan
apa yang tidak saya selesaikan.

Jika materi aljabar linear dalam bab pertama terasa terlalu ringkas, uraian yang lebih
santai dan lebih menyeluruh dapat ditemukan dalam catatan daring saya~\cite{Erdman:2010}.
% SOURCE-CORRECTION: PREFACE-C003
Demikian pula, materi latar tambahan tentang kalkulus lanjut, ruang metrik dan ruang
topologis, integral Lebesgue, dan sebagainya dapat ditemukan dalam banyak sekali sumber,
termasuk~\cite{Erdman:2005} dan~\cite{Erdman:2007}.

Tentu terdapat sejumlah kelebihan dan kekurangan ketika sebuah dokumen diterbitkan secara
elektronik. Salah satu kelebihannya ialah rujuk silang dapat diikuti dengan cepat melalui
pranala. Menelusuri sebuah buku untuk mencari \emph{lema 3.14.8} dapat memakan waktu
(terutama apabila seorang penulis menerapkan kebiasaan yang sepenuhnya logis tetapi
sangat menjengkelkan, yakni memberi nomor secara terpisah kepada lema, proposisi, teorema,
akibat, dan sebagainya). Kelebihan yang mungkin lebih penting ialah kemampuan untuk
segera memperbaiki kesalahan, menambahkan bagian yang hilang, memperjelas argumen yang
samar, dan membenahi ungkapan yang kurang baik. Kekurangan yang menyertainya ialah bahwa
seorang pembaca yang kembali ke laman web setelah waktu singkat mungkin mendapati segala
sesuatu (halaman, definisi, teorema, bagian) telah memperoleh nomor yang berbeda.
({\LaTeX} adalah alat yang mengagumkan.) Saya akan mengubah tanggal pada halaman judul
untuk memberi tahu pembaca kapan pembaruan tak sepele terakhir dilakukan (yakni pembaruan
yang memengaruhi nomor atau rujuk silang).

Kekurangan paling serius dari kehidupan elektronik ialah ketidaklestariannya. Dalam
kebanyakan kasus, ketika sebuah laman web menghilang, untuk segala keperluan praktis
informasi di dalamnya turut menghilang. Karena alasan ini (dan karena saya ingin materi
ini tersedia secara bebas bagi siapa pun yang menginginkannya), saya menggunakan lisensi
``BerbagiSerupa'' dari \emph{Creative Commons}. % SOURCE-CORRECTION: PREFACE-C004
Saya berharap siapa pun yang menganggap
materi ini berguna akan memperbaiki yang keliru, menambahkan yang hilang, dan menyempurnakan
yang canggung. Saya akan menyertakan berkas sumber {\LaTeX} di situs web saya agar
pemanfaatan dokumen ini, atau bagian-bagiannya, tidak mengharuskan pengetikan ulang tanpa
henti. Mengenai teksnya sendiri, silakan kirim koreksi, saran, keluhan, dan komentar lain
kepada penulis melalui
   \[ \textrm{erdman@pdx.edu} \]


\vfill\eject


\section*{Huruf Yunani}

 \index{huruf!Yunani}%
% SOURCE-CORRECTION: PREFACE-C005
% Markup tabel warisan yang dikecualikan diganti dengan tabularx; data dipertahankan.
\begin{center}
{\renewcommand{\arraystretch}{1.10}%
\begin{tabularx}{\textwidth}{@{}>{\centering\arraybackslash}p{0.23\textwidth}>{\centering\arraybackslash}p{0.27\textwidth}>{\raggedright\arraybackslash}X@{}}
\hline
\textbf{Huruf kapital} & \textbf{Huruf kecil} & \textbf{Nama dalam bahasa Inggris (perkiraan pelafalan)} \\
\hline
A         & $\alpha$                    & Alpha (AL-fuh) \\
B         & $\beta$                     & Beta (BAY-tuh) \\
$\Gamma$  & $\gamma$                    & Gamma (GAM-uh) \\
$\Delta$  & $\delta$                    & Delta (DEL-tuh) \\
$E$       & $\epsilon$ atau $\varepsilon$ & Epsilon (EPP-suh-lon) \\
Z         & $\zeta$                     & Zeta (ZAY-tuh) \\
H         & $\eta$                      & Eta (AY-tuh) \\
$\Theta$  & $\theta$                    & Theta (THAY-tuh) \\
I         & $\iota$                     & Iota (eye-OH-tuh) \\
K         & $\kappa$                    & Kappa (KAP-uh) \\
$\Lambda$ & $\lambda$                   & Lambda (LAM-duh) \\
M         & $\mu$                       & Mu (MYOO) \\
N         & $\nu$                       & Nu (NOO) \\
$\Xi$     & $\xi$                       & Xi (KSEE) \\
O         & o                            & Omicron (OHM-ih-kron) \\
$\Pi$     & $\pi$                       & Pi (PIE) \\
P         & $\rho$                      & Rho (ROH) \\
$\Sigma$  & $\sigma$                    & Sigma (SIG-muh) \\
T         & $\tau$                      & Tau (TAU) \\
Y         & $\upsilon$                  & Upsilon (OOP-suh-lon) \\
$\Phi$    & $\phi$                      & Phi (FEE or FAHY) \\
X         & $\chi$                      & Chi (KHAY) \\
$\Psi$    & $\psi$                      & Psi (PSEE or PSAHY) \\
$\Omega$  & $\omega$                    & Omega (oh-MAY-guh) \\
\hline
\end{tabularx}}
\end{center}








\vfill\eject







\section*{Font Fraktur}

Dalam catatan ini digunakan font Fraktur (paling sering untuk keluarga himpunan dan
keluarga operator). Padanan huruf Roman bagi setiap huruf diberikan di bawah ini. Ketika
menulis dengan tangan atau menyajikan materi di papan tulis, biasanya paling baik
menggantinya dengan huruf Latin kaligrafis.
 \index{font!Fraktur}

% SOURCE-CORRECTION: PREFACE-C006
% Markup tabel warisan yang dikecualikan diganti dengan tabularx; data dipertahankan.
\begin{center}
{\renewcommand{\arraystretch}{1.10}%
\begin{tabularx}{\textwidth}{@{}>{\centering\arraybackslash}X>{\centering\arraybackslash}X>{\centering\arraybackslash}X@{}}
\hline
\textbf{Fraktur} & \textbf{Fraktur} & \textbf{Roman} \\
\textbf{Huruf kapital} & \textbf{Huruf kecil} & \textbf{Huruf kecil} \\
\hline
$\mathfrak A$ & $\mathfrak a$ & a \\
$\mathfrak B$ & $\mathfrak b$ & b \\
$\mathfrak C$ & $\mathfrak c$ & c \\
$\mathfrak D$ & $\mathfrak d$ & d \\
$\mathfrak E$ & $\mathfrak e$ & e \\
$\mathfrak F$ & $\mathfrak f$ & f \\
$\mathfrak G$ & $\mathfrak g$ & g \\
$\mathfrak H$ & $\mathfrak h$ & h \\
$\mathfrak I$ & $\mathfrak i$ & i \\
$\mathfrak J$ & $\mathfrak j$ & j \\
$\mathfrak K$ & $\mathfrak k$ & k \\
$\mathfrak L$ & $\mathfrak l$ & l \\
$\mathfrak M$ & $\mathfrak m$ & m \\
$\mathfrak N$ & $\mathfrak n$ & n \\
$\mathfrak O$ & $\mathfrak o$ & o \\
$\mathfrak P$ & $\mathfrak p$ & p \\
$\mathfrak Q$ & $\mathfrak q$ & q \\
$\mathfrak R$ & $\mathfrak r$ & r \\
$\mathfrak S$ & $\mathfrak s$ & s \\
$\mathfrak T$ & $\mathfrak t$ & t \\
$\mathfrak U$ & $\mathfrak u$ & u \\
$\mathfrak V$ & $\mathfrak v$ & v \\
$\mathfrak W$ & $\mathfrak w$ & w \\
$\mathfrak X$ & $\mathfrak x$ & x \\
$\mathfrak Y$ & $\mathfrak y$ & y \\
$\mathfrak Z$ & $\mathfrak z$ & z \\
\hline
\end{tabularx}}
\end{center}


  \vfill\eject
\phantomsection
\section*{Notasi untuk Himpunan Bilangan}\label{C0009} % SOURCE-CORRECTION: PREFACE-C007
Berikut adalah daftar notasi yang cukup baku untuk beberapa himpunan bilangan yang sering
muncul dalam catatan ini:
 \index{real!garis!subhimpunan khusus dari}%
 \index{bilangan!himpunan khusus dari}%
 \index{C@$\C$ (himpunan bilangan kompleks)}%
 \index{R@$\R$ (himpunan bilangan real)}%
 \index{R@$\R^n$!himpunan tupel-$n$ bilangan real}%
 \index{Q@$\Q$ (himpunan bilangan rasional)}%
 \index{Z@$\Z$ (himpunan bilangan bulat)}%
 \index{N@$\N$ (himpunan bilangan asli)}%
 \index{N@$\N_n$ ($n$ bilangan asli pertama)}%
 \index{D@$\Di$ (cakram satuan terbuka)}%
 \index{T@$\T$ (lingkaran satuan)}%
 \index{S@$\Sp^1$ (lingkaran satuan)}%
\begin{align*}
 &\C \text{ adalah himpunan bilangan kompleks} \\
 &\R \text{ adalah himpunan bilangan real} \\
 &\R^n \text{ adalah himpunan semua tupel-$n$ $(r_1,r_2,\dots,r_n)$ bilangan real} \\
 &\R^+ = \{x \in \R\colon x \ge 0\}, \text{ himpunan bilangan real tak negatif} \\ % SOURCE-CORRECTION: PREFACE-C008
 &\Q \text{ adalah himpunan bilangan rasional} \\
 &\Q^+ = \{x \in \Q\colon x \ge 0\}, \text{ himpunan bilangan rasional tak negatif} \\ % SOURCE-CORRECTION: PREFACE-C009
 &\Z \text{ adalah himpunan bilangan bulat} \\
 &\Z^+ = \{0,1,2,\dots\}, \text{ himpunan bilangan bulat tak negatif} \\ % SOURCE-CORRECTION: PREFACE-C010
 &\N = \{1,2,3,\dots\}\text{, himpunan bilangan asli} \\
 &\N_n = \{1,2,3,\dots,n\}\text{, $n$ bilangan asli pertama} \\ % SOURCE-CORRECTION: PREFACE-C011
 &[a,b] = \{x \in \R\colon a \le x \le b\} \\
 &[a,b) = \{x \in \R\colon a \le x < b\} \\
 &(a,b] = \{x \in \R\colon a < x \le b\} \\
 &(a,b) = \{x \in \R\colon a < x < b\} \\
 &[a,\infty) = \{x \in \R\colon a \le x\} \\
 &(a,\infty) = \{x \in \R\colon a < x\} \\
 &(-\infty,b] = \{x \in \R\colon x \le b\} \\
 &(-\infty,b) = \{x \in \R\colon  x < b\} \\
\index{cakram!satuan terbuka}%
 &\Di = \{(x,y) \in \R^2\colon x^2 + y^2 < 1\},\text{ cakram satuan terbuka} \\
 &\T = \Sp^1 = \{(x,y) \in \R^2\colon x^2 + y^2 = 1\},\text{ lingkaran satuan} \\
\end{align*}



\vfill\eject










\section*{Notasi untuk Fungsi}
\enlargethispage{2\baselineskip}
Fungsi sudah tidak asing sejak kalkulus dasar. Secara informal, suatu \emph{fungsi} terdiri atas sepasang
himpunan dan suatu ``aturan'' yang mengaitkan setiap anggota himpunan pertama (yaitu
\emph{domain}) dengan tepat satu anggota himpunan kedua (yaitu \emph{kodomain}). Walaupun
``definisi'' informal ini tentu memadai untuk sebagian besar keperluan dan jarang menimbulkan
kesalahpahaman, kadang-kadang perumusan yang lebih tepat tetap berguna. Perumusan itu diperoleh
dengan mendefinisikan fungsi melalui suatu jenis khusus relasi antara dua himpunan.

Suatu
 \index{fungsi}%
\df{fungsi} $f$ adalah tripel terurut $(S,T,G)$, dengan $S$ dan $T$ himpunan serta $G$
subhimpunan dari $S \times T$ yang memenuhi:
 \begin{enumerate}
   \item untuk setiap $s \in S$ terdapat $t \in T$ sedemikian sehingga $(s,t) \in
G$; dan
   \item jika $(s,t_1)$ dan $(s,t_2)$ termasuk dalam $G$, maka $t_1 = t_2$. % SOURCE-CORRECTION: PREFACE-C014
 \end{enumerate}
Dalam keadaan ini kita mengatakan bahwa $f$ adalah \emph{fungsi dari} $S$ \emph{ke} $T$ (atau bahwa
$f$ \emph{memetakan} $S$ \emph{ke} $T$) dan menulis $f\colon S \sto T$. Himpunan $S$ adalah
 \index{domain}%
\df{domain} (atau
 \index{masukan!ruang}%
 \index{ruang!masukan}%
\df{ruang masukan}) dari~$f$. Himpunan $T$ adalah
 \index{kodomain}%
\df{kodomain} (atau
 \index{sasaran!ruang}%
 \index{ruang!sasaran}%
\df{ruang sasaran}, atau
 \index{keluaran!ruang}%
 \index{ruang!keluaran}%
\df{ruang keluaran}) dari~$f$. Relasi $G$ adalah
 \index{graf}%
\df{graf} dari~$f$. Untuk menghindari penyebutan graf $G$ secara eksplisit, biasanya
ungkapan ``$(x,y) \in G\,$'' diganti dengan ``$y = f(x)$''; elemen $f(x)$ adalah
 \index{citra!dari suatu titik}%
\df{citra} $x$ di bawah~$f$. Dalam catatan ini (tetapi tidak di semua tempat), kata
 \index{transformasi}%
 \index{konvensi!fungsi = peta = pemetaan = transformasi}%
``transformasi'',
 \index{peta}%
``peta'', dan
 \index{pemetaan}%
``pemetaan'' digunakan sebagai sinonim ``fungsi''. Domain $f$ dinotasikan
 \index{domain@$\dom f$ (domain fungsi $f$)}%
dengan~$\dom f$.

Jika $S$ dan $T$ himpunan, keluarga semua fungsi yang berdomain $S$ dan
berkodomain $T$ dinotasikan
 \index{F@$\fml F(S,T)$ (himpunan fungsi dari $S$ ke~$T$)}%
dengan~$\fml F(S,T)$.

Jika $f \colon S \sto T$ dan $A\subseteq S$, maka
 %\index{<@$\bigl.f\bigr"|_A$ (pembatasan $f$ pada $A$)}%
 \index{<@$f\mid_{{}_A}$ (pembatasan $f$ pada $A$)}%
 \index{pembatasan}%
\df{pembatasan} $f$ pada $A$, yang dinotasikan dengan $\bigl.f\bigr|_A$, adalah
pemetaan dari $A$ ke $T$ yang nilainya pada setiap $x$ di $A$ sama dengan $f(x)$.

Jika $f\colon S \sto T$ dan $A \subseteq S$, maka $f^{\sto} (A)$, yaitu
 \index{<arrowto@$f^{\sto}(A)$ (citra himpunan $A$ di bawah fungsi $f$)}%
 \index{citra!dari suatu himpunan}%
\df{citra} $A$ di bawah $f$ adalah $\{f(x)\colon x \in A\}$. Lazim pula menulis
$f(A)$ sebagai pengganti $f^{\sto}(A)$. Himpunan $f^{\sto}(S)$ adalah
 \index{jangkauan}%
 \index{jangkauan@$\ran f$ (jangkauan fungsi $f$)}%
\df{jangkauan} (atau
 \index{citra!dari suatu fungsi}%
\df{citra}) dari $f$; biasanya kita menulis $\ran f$ sebagai pengganti $f^{\sto}(S)$.

Jika $f\colon S \sto T$ dan $B \subseteq T$, maka
$f^\gets (B)$, yaitu
 \index{<arrowgets@$f^{\gets}(B)$ (pracitra himpunan $B$ di bawah fungsi $f$)}%
 \index{pracitra!dari suatu himpunan}%
 \index{citra!invers!dari suatu himpunan}%
\df{pracitra} $B$ di bawah~$f$ (atau citra inversnya) adalah $\{x \in S\colon f(x) \in B\}$. Dalam banyak buku,
$f^{\gets}(B)$ dinotasikan dengan $f^{-1}(B)$. Notasi ini dapat membingungkan karena menyiratkan
bahwa fungsi $f$ memiliki invers, padahal belum tentu demikian.

Suatu fungsi $f\colon S \sto T$ bersifat
 \index{injektif}%
\df{injektif} (atau
 \index{satu-satu}%
\df{satu-satu}) jika $f(x) = f(y)$ mengakibatkan $x = y$ untuk semua $x$, $y \in S$. Fungsi itu bersifat
 \index{surjektif}%
\df{surjektif} (atau
 \index{onto}%
\df{onto}) jika untuk setiap $t \in T$ terdapat $x \in S$ sedemikian sehingga $f(x) = t$. % SOURCE-CORRECTION: PREFACE-C012
Fungsi itu bersifat
 \index{bijektif}%
\df{bijektif} (atau merupakan
 \index{satu-satu!korespondensi}%
 \index{korespondensi!satu-satu}%
\df{korespondensi satu-satu}) jika fungsi itu injektif sekaligus surjektif. % SOURCE-CORRECTION: PREFACE-C013

Sering kali berguna untuk memandang fungsi sebagai panah dalam diagram. Sebagai contoh, keadaan
$j\colon R \sto U$, $f\colon R \sto S$, $k\colon S \sto T$, $h\colon U \sto T$
dapat direpresentasikan oleh
 \index{diagram}%
diagram berikut.
  \[ \xy
       \square[R`U`S`T;j`f`h`k]
  \endxy \]
Diagram itu dikatakan
 \index{berkomutasi}%
\df{berkomutasi} (atau disebut
 \index{komutatif!diagram}%
 \index{diagram!komutatif}%
\df{diagram komutatif}) jika $h \circ j = k \circ f$.

Diagram tidak harus berbentuk persegi panjang. Sebagai contoh,
 \[ \xy
     \btriangle[R`S`T;f`g`k]
 \endxy \]
adalah diagram komutatif jika $g = k \circ f$.


\vfill\eject













  \endinput


 \cleardoublepage


\mainmatter




 \chapter{ALJABAR LINEAR DAN TEOREMA SPEKTRAL}

Dalam mata kuliah analisis tingkat awal, penekanan diberikan pada perilaku fungsi-fungsi individual (bernilai real atau kompleks). Dalam mata kuliah \emph{analisis fungsional},
fokus beralih ke \emph{ruang-ruang} fungsi tersebut dan kelas-kelas pemetaan tertentu di antara ruang-ruang itu. Ternyata uraian yang ringkas,
mungkin terlampau menyederhanakan, tetapi jelas tidak menyesatkan, tentang materi tersebut ialah \emph{aljabar linear berdimensi tak hingga}. Dari sudut pandang
seorang analis, pencapaian terbesar aljabar linear sebagian besar berupa teorema-teorema tentang operator pada ruang vektor berdimensi hingga
(terutama $\R^n$ dan $\C^n$). Akan tetapi, ruang vektor pemetaan bernilai skalar yang didefinisikan pada berbagai domain biasanya berdimensi tak
hingga. Apakah proses memperumum hasil-hasil mengagumkan dari aljabar linear berdimensi hingga ke ranah berdimensi tak hingga dipandang sebagai
rangkaian komplikasi abstrak yang sangat tidak menyenangkan atau sebagai sumber kaya wawasan menakjubkan, pertanyaan baru yang menarik, dan
petunjuk yang menggugah menuju penerapan di bidang lain sepenuhnya bergantung pada orientasi seseorang terhadap matematika secara umum.

Berdasarkan uraian di atas, akan bermanfaat untuk memulai perjalanan kita ke wilayah baru ini dengan tinjauan singkat atas beberapa keberhasilan
klasik aljabar linear. Tidak jarang mahasiswa yang berhasil menyelesaikan mata kuliah pengantar aljabar linear pada akhirnya hanya memiliki
gagasan yang paling samar tentang pokok bahasan mata kuliah itu. Sebagian kesalahan mungkin terletak pada banyak buku dasar yang mencampuradukkan
secara tidak semestinya dua pokok bahasan yang sangat berbeda, meskipun berkaitan erat: di satu pihak, kajian \emph{ruang vektor} dan pemetaan linear
di antaranya, dan di pihak lain, \emph{ruang hasil kali dalam} beserta pemetaan linearnya. Ruang vektor tidak memiliki gagasan geometri/topologi
tentang \emph{jarak}, \emph{panjang}, \emph{ketegaklurusan}, himpunan \emph{terbuka}, atau \emph{sudut} antara vektor; yang ada hanyalah penjumlahan dan perkalian skalar.
Ruang hasil kali dalam adalah ruang vektor yang dilengkapi struktur-struktur tambahan tersebut. Mari kita tinjau keduanya secara terpisah.


\section{Ruang Vektor dan Dekomposisi Operator yang Dapat Didiagonalkan}
\begin{conv} Dalam catatan ini semua ruang vektor
 \index{konvensi!ruang vektor kompleks atau real}%
 \index{C@$\C$!medan bilangan kompleks}%
 \index{R@$\R$!medan bilangan real}%
akan diasumsikan sebagai ruang vektor atas medan~$\C$ bilangan kompleks (dalam hal ini disebut \emph{ruang vektor kompleks}) atau atas
medan~$\R$ bilangan real (dalam hal ini disebut \emph{ruang vektor real}). Tidak ada medan lain yang akan muncul. Ketika $\K$ muncul, medan itu dapat dipandang
 \index{K@$\K$ (medan bilangan real atau kompleks)}%
sebagai $\C$ ataupun~$\R$. Suatu
 \index{skalar}%
\df{skalar} adalah anggota $\K$; yaitu, suatu \emph{bilangan kompleks} atau suatu \emph{bilangan real}.
\end{conv}

\begin{defn}\label{00001} Tripel $(V,+,M)$ adalah
 \index{vektor!ruang}%
 \index{ruang!vektor}%
\df{ruang vektor} atas $\K$ jika $(V,+)$ adalah grup Abelian dan $M\colon \K \sto \Hom(V)$
adalah homomorfisme gelanggang beridentitas (dengan $\Hom(V)$ gelanggang homomorfisme grup
pada~$V$).
\end{defn}

Untuk memeriksa arti istilah-istilah dalam definisi sebelumnya, lihat tiga bagian pertama dari bab pertama catatan
aljabar linear saya~\cite{Erdman:2010}.

\begin{exer} Definisi \emph{ruang vektor} yang ditemukan dalam banyak buku dasar kurang lebih
sebagai berikut: suatu \emph{ruang vektor atas} $\K$ adalah himpunan $V$ beserta operasi
penjumlahan dan perkalian skalar yang memenuhi aksioma-aksioma berikut:
 \begin{enumerate}
  \item[(1)] jika $x$, $y \in V$, maka $x + y \in V$;
  \item[(2)] $(x + y) + z = x + (y + z)$ untuk setiap
 \index{asosiatif}%
$x$, $y$, $z \in V$ (asosiativitas);
  \item[(3)] terdapat $\vc 0 \in V$ sehingga $x + \vc 0 = x$ untuk setiap
 \index{identitas!aditif}%
$x \in V$ (keberadaan identitas aditif);
  \item[(4)] untuk setiap $x\in V$ terdapat $-x \in V$ sehingga
 \index{aditif!invers}%
 \index{invers!aditif}%
$x + (-x) = \vc 0$ (keberadaan invers aditif);
  \item[(5)] $x + y = y + x$ untuk setiap $x$, $y \in V$
 \index{komutatif}%
(komutativitas);
  \item[(6)] jika $\alpha \in \K$ dan $x \in V$, maka $\alpha x \in V$;
  \item[(7)] $\alpha (x + y) = \alpha x + \alpha y$ untuk setiap $\alpha \in \K$
dan setiap $x$, $y \in V$;
  \item[(8)] $(\alpha + \beta)x = \alpha x + \beta x$ untuk setiap $\alpha$,
$\beta \in \K$ dan setiap $x \in V$;
  \item[(9)] $(\alpha\beta)x = \alpha(\beta x)$ untuk setiap $\alpha$,
$\beta \in \K$ dan setiap $x \in V$; dan
  \item[(10)] $1\,x = x$ untuk setiap $x \in V$.
 \end{enumerate}
Buktikan bahwa definisi ini ekuivalen dengan definisi yang diberikan di atas dalam~\ref{00001}.
\end{exer}

\begin{defn} Subhimpunan $M$ dari ruang vektor $V$ adalah
 \index{subruang!vektor}%
 \index{vektor!subruang}%
\df{subruang vektor} dari $V$ jika merupakan ruang vektor terhadap operasi yang diwarisinya dari~$V$.
\end{defn}

\begin{prop} Subhimpunan tak kosong $M$ dari ruang vektor $V$ adalah subruang vektor dari $V$ jika dan
hanya jika tertutup terhadap penjumlahan dan perkalian skalar. (Artinya: jika $\vc x$ dan
$\vc y$ termasuk dalam $M$, maka $\vc x + \vc y$ juga demikian; dan jika $\vc x$ termasuk dalam $M$ dan $\alpha
\in \K$, maka $\alpha \vc x$ termasuk dalam~$M$.)
\end{prop}

\begin{defn} Vektor $y$ adalah
 \index{linear!kombinasi}%
 \index{kombinasi!linear}%
\df{kombinasi linear} dari vektor-vektor $x_1$, \dots, $x_n$ jika terdapat skalar $\alpha_1$,
\dots, $\alpha_n$ sehingga $y = \sum_{k=1}^n \alpha_k x_k$. Kombinasi linear $\sum_{k=1}^n \alpha_k x_k$ disebut
 \index{linear!kombinasi!trivial}%
 \index{trivial!kombinasi linear}%
\df{trivial} jika semua koefisien $\alpha_1$, \dots, $\alpha_n$ bernilai nol. Jika sekurang-kurangnya
satu $\alpha_k$ tidak sama dengan nol, kombinasi linear itu disebut \df{nontrivial}.
\end{defn}

\begin{defn} Jika $A$ adalah subhimpunan tak kosong dari ruang vektor~$V$, maka $\spn A$, yaitu
 \index{rentang}%
\df{rentang} dari $A$, adalah himpunan semua kombinasi linear unsur-unsur~$A$. Himpunan ini juga sering disebut
 \index{rentang!linear}%
 \index{linear!rentang}%
\df{rentang linear} dari~$A$.
\end{defn}

\begin{notn}\label{bases111} Misalkan $A$ adalah himpunan yang merentang ruang vektor~$V$. Untuk setiap
$x \in V$ terdapat himpunan hingga $S$ yang terdiri atas vektor-vektor dalam $A$, dan untuk setiap unsur $e$ dalam $A$
terdapat skalar $x_e$ sehingga $x = \sum_{e \in S} x_e e$. Jika kita sepakat menetapkan $x_e = 0$
apabila $e \in A \setminus S$, kita dapat pula menulis $x = \sum_{e \in A} x_e e$. Walaupun
notasi ini mungkin memberi kesan bahwa kita menjumlahkan atas himpunan sembarang yang mungkin tak terhitung,
kenyataannya semua kecuali sejumlah hingga suku bernilai nol, sehingga tidak timbul masalah
``kekonvergenan''. Perlakukan $\sum_{e \in A} x_e e$ sebagai jumlah hingga.
Asosiativitas dan komutativitas penjumlahan dalam $V$ membuat ekspresi tersebut tidak ambigu.
\end{notn}

\begin{defn}\label{defn_lin_dep} Subhimpunan $A$ dari suatu ruang vektor disebut
 \index{linear!kebergantungan}%
 \index{bergantung}
\df{bergantung linear} jika vektor nol $\vc 0$ dapat ditulis sebagai kombinasi linear nontrivial
dari unsur-unsur~$A$; yaitu, jika terdapat vektor $x_1, \dots, x_n \in A$
dan skalar $\alpha_1, \dots, \alpha_n$, yang \textbf{tidak semuanya nol,} sehingga $\sum_{k=1}^n
\alpha_kx_k = \vc 0$. Subhimpunan suatu ruang vektor disebut
 \index{linear!kebebasan}%
 \index{bebas}
\df{bebas linear} jika tidak bergantung linear.
\end{defn}

Secara teknis, suatu \emph{himpunan} vektorlah yang bergantung atau bebas linear.
Meskipun demikian, istilah-istilah ini sering digunakan seolah-olah merupakan sifat vektornya
sendiri. Sebagai contoh, jika $S = \{x_1, \dots, x_n\}$ adalah himpunan hingga vektor dalam suatu
ruang vektor, pernyataan ``himpunan $S$ bebas linear'' dan
``vektor-vektor $x_1$, \dots, $x_n$ bebas linear'' dapat digunakan secara bergantian.

\begin{defn}\label{defn_Hamel_basis} Himpunan $B$ yang terdiri atas vektor-vektor dalam ruang vektor $V$ adalah
 \index{basis!Hamel}%
 \index{Hamel!basis}%
\df{basis Hamel} untuk $V$ jika bebas linear dan merentang~$V$.
\end{defn}

\begin{conv}\label{conv_1_1} Ruang vektor yang hanya memuat satu vektor, yaitu vektor nol, adalah ruang vektor
 \index{konvensi!basis untuk ruang vektor nol}%
 \index{vektor!ruang!trivial (atau nol)}%
 \index{trivial!ruang vektor}%
\df{trivial} (atau
 \index{nol!ruang vektor}%
\df{nol}). Terkait ruang ini terdapat sebuah pertanyaan teknis, meskipun tidak terlalu menarik. Apakah ruang ini memiliki basis Hamel?
Dari definisi \emph{kebergantungan linear}~\ref{defn_lin_dep}, jelas bahwa himpunan kosong $\emptyset$ bebas
linear. Himpunan kosong tentu merupakan subhimpunan bebas linear maksimal dari ruang nol (suatu kondisi yang, untuk ruang nontrivial,
ekuivalen dengan menjadi himpunan perentang yang bebas linear). Akan tetapi, sulit untuk berpendapat bahwa himpunan kosong merentang sesuatu.
Jadi, semata-mata sebagai konvensi, kita akan mengatakan bahwa jawaban atas pertanyaan itu adalah \emph{ya}. Ruang vektor nol memiliki basis
Hamel, dan basis itu adalah himpunan kosong!
\end{conv}

Dalam buku pengantar aljabar linear, himpunan perentang yang bebas linear hanya disebut \emph{basis}. Di sini digunakan istilah \emph{basis Hamel}
untuk membedakannya dari \emph{basis ortonormal} dan \emph{basis Schauder}, istilah-istilah yang mengacu pada konsep yang lebih
penting dalam analisis fungsional.

\begin{prop} Misalkan $B$ adalah basis Hamel untuk ruang vektor nontrivial~$V$. Maka setiap unsur dalam $V$ dapat ditulis secara tunggal sebagai kombinasi
linear anggota-anggota~$B$.
\end{prop}

\begin{proof}[\emph{Petunjuk pembuktian}] Eksistensinya jelas. Sederhanakan pembuktian ketunggalan dengan menggunakan konvensi notasi
yang disarankan dalam~\ref{bases111}.   \ns
\end{proof}

\begin{prop}\label{bases025} Misalkan $A$ adalah subhimpunan bebas linear dari ruang
vektor~$V$. Maka terdapat basis Hamel untuk~$V$ yang memuat~$A$.
\end{prop}

\begin{proof}[\emph{Petunjuk pembuktian}] Tunjukkan terlebih dahulu bahwa, untuk ruang nontrivial, subhimpunan bebas linear dari $V$ adalah basis untuk~$V$
jika dan hanya jika merupakan subhimpunan bebas linear maksimal. Kemudian urutkan keluarga subhimpunan bebas linear dari
$V$ yang memuat~$A$ berdasarkan inklusi dan terapkan \emph{lema Zorn}. (Jika Anda belum mengenal \emph{lema Zorn},
lihat Bagian 1.7 dalam catatan aljabar linear saya~\cite{Erdman:2010}.)  \ns
\end{proof}

\begin{cor} Setiap ruang vektor memiliki basis Hamel.
\end{cor}

\begin{defn}\label{000082} Misalkan $M$ dan $N$ adalah subruang dari ruang vektor $V$. Jika
$M \cap N = \{\vc 0\}$ dan $M + N = V$, maka $V$ adalah
 \index{internal!jumlah langsung!dalam ruang vektor}%
 \index{langsung!jumlah!internal}%
 \index{langsung!jumlah!ruang vektor}%
 \index{<binopdirectsum@$\oplus$ (jumlah langsung ruang vektor)}%
\df{jumlah langsung} (\df{internal}) dari $M$ dan $N$. Dalam hal ini kita menulis
   \[ V = M \oplus N\,. \]
Kita mengatakan bahwa $M$ dan $N$ adalah
 \index{komplemen!subruang vektor}%
 \index{vektor!ruang!komplemen}%
\df{subruang vektor komplementer} dan masing-masing merupakan \df{komplemen} (ruang vektor) dari
yang lain. 
 \index{kodimensi}%
\df{Kodimensi} subruang $M$ adalah dimensi komplemennya~$N$.
\end{defn}

\begin{prop}\label{0000823} Misalkan $V$ adalah ruang vektor dan andaikan $V = M \oplus N$. Maka untuk
setiap $v \in V$ terdapat vektor-vektor tunggal $m \in M$ dan $n \in N$ sehingga $v = m + n$.
\end{prop}

\begin{exam} Misalkan $\fml C = \fml C[-1,1]$ adalah ruang vektor semua fungsi kontinu bernilai real
pada interval $[-1,1]$. Fungsi $f$ dalam $\fml C$ disebut \df{genap} jika $f(-x) =
f(x)$ untuk semua $x \in [-1,1]$; fungsi itu disebut \df{ganjil} jika $f(-x) = -f(x)$ untuk semua $x \in [-1,1]$.
Misalkan $\fml C_o = \{f \in \fml C\colon  f \text{ ganjil }\}$ dan $\fml C_e = \{f \in \fml
C\colon  f \text{ genap }\}$. Maka $\fml C = \fml C_o \oplus \fml C_e$.
\end{exam}

\begin{prop}\label{0000824} Jika $M$ adalah subruang dari ruang vektor $V$, maka terdapat subruang $N$ dari
$V$ sehingga $V = M \oplus N$.
\end{prop}

\begin{lem}\label{bases030} Misalkan $V$ adalah ruang vektor dengan basis hingga tak kosong $\{e^1,
\dots,e^n\}$ dan misalkan $v = \sum_{k=1}^n \alpha_k e^k$ adalah vektor dalam~$V$. Jika $p \in \N_n$
dan $\alpha_p \ne 0$, maka
\[
\{e^1, \dots,e^{p-1},v,e^{p+1}, \dots, e^n\}
\]
adalah basis untuk~$V$.
\end{lem}

\begin{prop}\label{bases031} Jika suatu basis untuk ruang vektor $V$ memuat $n$ unsur,
maka setiap subhimpunan bebas linear dari $V$ yang memiliki $n$ unsur juga merupakan basis.
\end{prop}

\begin{proof}[\emph{Petunjuk pembuktian}] Andaikan $\{e^1,\dots,e^n\}$ adalah basis untuk~$V$ dan
$\{v^1, \dots,v^n\}$ bebas linear dalam~$V$. Mulailah dengan menggunakan lema~\ref{bases030}
untuk menunjukkan bahwa (setelah mungkin menomori ulang $e^k$) himpunan $\{v^1, e^2, \dots, e^n\}$
adalah basis untuk~$V$.  \ns
\end{proof}

\begin{cor}\label{bases032} Jika ruang vektor $V$ memiliki basis hingga $B$, maka setiap basis
untuk $V$ berhingga dan memuat jumlah unsur yang sama dengan~$B$.
\end{cor}

\begin{defn} Suatu ruang vektor disebut
 \index{berdimensi hingga}%
\df{berdimensi hingga} jika memiliki basis hingga, dan
 \index{dimensi!ruang vektor}%
 \index{vektor!ruang!dimensi}%
 \index{dim@$\dim V$ (dimensi ruang vektor $V$)}%
\df{dimensi} ruang tersebut adalah banyaknya unsur dalam basis ini (dan dengan demikian dalam basis mana pun) untuk
ruang itu. Dimensi ruang vektor berdimensi hingga $V$ dinotasikan dengan $\dim V$. Berdasarkan
konvensi yang dibuat di atas~\ref{conv_1_1}, ruang vektor nol berdimensi nol. Jika
ruang tersebut tidak memiliki basis hingga, ruang itu disebut
 \index{berdimensi tak hingga}%
\df{berdimensi tak hingga}.
\end{defn}

\begin{defn}\label{defn_op_vsp} Fungsi $T\colon V \sto W$ antara ruang-ruang vektor (atas medan yang sama) disebut
 \index{linear!pemetaan}%
\df{linear} jika $T(u + v) = Tu + Tv$ untuk semua $u$, $v \in V$ dan $T(\alpha v) = \alpha Tv$
untuk semua $\alpha \in \K$ dan $v \in V$. Fungsi linear sering disebut
\emph{transformasi linear} atau \emph{pemetaan linear}. Jika $V = W$, kita mengatakan bahwa pemetaan linear $T$ adalah
 \index{operator!pada ruang vektor}%
\df{operator} pada~$V$. Bergantung pada konteks, operator identitas $x \mapsto x$ pada $V$ dinotasikan
 \index{identitas@$\id V$ atau $I_V$ atau $I$ (operator identitas)}%
dengan $\id V$ atau $I_V$ atau cukup~$I$. Ingat bahwa jika $T\colon V \sto W$ adalah pemetaan linear, maka
 \index{kernel!pemetaan linear}%
 \index{<ker@$\ker T$ (kernel dari~$T$)}%
\df{kernel} dari $T$, yang dinotasikan dengan $\ker T$, adalah $T^{\gets}(\{0\}) := \{x \in V\colon Tx = \vc
0\}$. Selain itu,
 \index{jangkauan!pemetaan linear}%
 \index{<ran@$\ran T$ (jangkauan dari~$T$)}%
\df{jangkauan} dari $T$, yang dinotasikan dengan $\ran T$, adalah $T^{\sto}(V) := \{Tx\colon x \in V\}$.
\end{defn}

\begin{prop} Pemetaan linear $T\colon V \sto W$ bersifat injektif jika dan hanya jika kernelnya nol.
\end{prop}

\begin{prop} Misalkan $T\colon V \sto W$ adalah pemetaan linear injektif antara ruang-ruang vektor. Jika $A$ adalah subhimpunan bebas linear
dari~$V$, maka $T^\sto(A)$ adalah subhimpunan bebas linear dari~$W$.
\end{prop}

\begin{prop} Misalkan $T\colon V \sto W$ adalah pemetaan linear antara ruang-ruang vektor dan $A \subseteq V$.
Maka $T^\sto(\spn A) = \spn T^\sto(A)$.
\end{prop}

\begin{prop} Misalkan $T\colon V \sto W$ adalah pemetaan linear injektif antara ruang-ruang vektor. Jika $B$ adalah basis
untuk subruang $U$ dari~$V$, maka $T^\sto(B)$ adalah basis untuk~$T^\sto(U)$.
\end{prop}

\begin{prop} Dua ruang vektor berdimensi hingga isomorfik jika dan hanya jika keduanya memiliki dimensi yang sama.
\end{prop}

\begin{defn} Misalkan $T$ suatu pemetaan linear antara ruang-ruang vektor. \df{Peringkat} $T$ adalah dimensi jangkauan~$T$, sedangkan
 \index{peringkat}
 \index{nulitas}%
\df{nulitas} $T$ adalah dimensi kernel~$T$.
\end{defn}

\begin{prop} Misalkan $T\colon V \sto W$ adalah pemetaan linear antara ruang-ruang vektor. Jika $V$ berdimensi hingga, maka
  \[ \operatorname{peringkat} T + \operatorname{nulitas} T = \dim V. \]
\end{prop}

\begin{defn} Pemetaan linear $T\colon V \sto W$ antara ruang-ruang vektor disebut
 \index{invertibel!pemetaan linear}%
\df{invertibel} (atau suatu
 \index{isomorfisme!ruang vektor}%
\df{isomorfisme}) jika terdapat pemetaan linear $T^{-1}\colon W \sto V$ sehingga $T^{-1}T
= \id V$ dan $TT^{-1} = \id W$. Dua ruang vektor $V$ dan $W$ disebut
 \index{isomorfik!ruang vektor}%
\df{isomorfik} jika terdapat isomorfisme dari $V$ ke~$W$.
\end{defn}

\begin{prop}\label{inv_lin_maps001} Pemetaan linear $T\colon V \sto W$ antara ruang-ruang vektor
invertibel jika dan hanya jika memiliki invers kiri sekaligus invers kanan.
\end{prop}

\begin{prop}\label{inv_lin_maps004} Pemetaan linear antara ruang-ruang vektor invertibel jika
dan hanya jika bijektif.
\end{prop}

\begin{prop} Misalkan $T$ operator pada ruang vektor berdimensi hingga. Ketiga pernyataan berikut ekuivalen:
 \begin{enumerate}
    \item $T$ adalah isomorfisme;
    \item $T$ injektif; dan
    \item $T$ surjektif.
 \end{enumerate}
\end{prop}

 \begin{defn}\label{defn_similar_vs_ops} Dua operator $R$ dan $T$ pada ruang vektor $V$ disebut
 \index{serupa!operator}%
\df{serupa} jika terdapat operator invertibel $S$ pada $V$ sehingga $R = STS^{-1}$.
\end{defn}

\begin{prop}\label{sim004thm} Jika $V$ adalah ruang vektor, maka keserupaan merupakan relasi ekuivalensi
pada~$\ofml L(V)$.
\end{prop}

Misalkan $V$ dan $W$ adalah ruang vektor berdimensi hingga dengan basis terurut. Andaikan $V$ berdimensi $n$ dengan basis terurut
$\{e^1,e^2, \dots,e^n\}$ dan $W$ berdimensi $m$. Ingat dari pengantar aljabar linear bahwa jika $T\colon V \sto W$ linear,
maka
 \index{matriks!representasi}%
 \index{representasi!matriks}%
\emph{representasi matriks} $[T]$ (diambil terhadap basis terurut dalam $V$ dan~$W$) adalah matriks $m \times n$ $[T]$
yang kolom ke-$k$ ($1 \le k \le n$) adalah vektor kolom $Te^k$ dalam~$W$. Pokoknya ialah bahwa tindakan $T$ pada
vektor $x$ dapat \emph{direpresentasikan} sebagai perkalian $x$ oleh matriks $[T]$; yaitu,
   \[ Tx = [T]x. \]
Persamaan ini mungkin memerlukan sedikit penafsiran. Ruas kiri adalah fungsi $T$ yang dievaluasi di $x$, dengan
hasil evaluasi dipandang sebagai vektor kolom dalam~$W$; ruas kanan adalah matriks $m \times n$ yang dikalikan
dengan matriks $n \times 1$ (yaitu vektor kolom). Jadi, kesamaan yang dinyatakan adalah kesamaan dua matriks $m \times 1$ (vektor kolom).

\begin{defn} Misalkan $V$ adalah ruang vektor berdimensi hingga dan $B = \{e^1, \dots, e^n\}$ adalah basis untuk~$V$. Operator $T$ pada $V$ disebut
 \index{diagonal!operator}%
 \index{operator!diagonal}%
\df{diagonal} jika terdapat skalar $\alpha_1, \dots, \alpha_n$ sehingga $Te^k =
\alpha_ke^k$ untuk setiap $k \in \N_n$. Secara ekuivalen, $T$ diagonal jika representasi matriksnya
$[T] = [T_{ij}]$ memiliki sifat bahwa $T_{ij} = 0$ apabila $i \ne j$.
\end{defn}

Menanyakan apakah operator tertentu pada suatu ruang vektor berdimensi hingga bersifat diagonal,
secara tegas, tidak bermakna. Sebagaimana didefinisikan, sifat operator berupa kediagonalan jelas
\emph{bukan} konsep ruang vektor. Sifat ini hanya bermakna bagi ruang vektor
\emph{yang basisnya telah ditentukan}. Fakta penting ini, meskipun jelas, tampaknya
luput dari perhatian dalam banyak mata kuliah pengantar aljabar linear, mungkin karena fiksasi yang agak berlebihan
pada $\R^n$ dalam mata kuliah tersebut. Berikut sifat \emph{ruang vektor} yang relevan.

\begin{defn} Operator $T$ pada ruang vektor berdimensi hingga $V$ disebut
 \index{dapat didiagonalkan!operator}%
 \index{operator!dapat didiagonalkan}%
\df{dapat didiagonalkan} jika terdapat basis untuk $V$ yang terhadapnya $T$
bersifat diagonal. Secara ekuivalen, operator pada ruang vektor berdimensi hingga \emph{dengan
basis} dapat didiagonalkan jika serupa dengan operator diagonal.
\end{defn}

\begin{defn} Misalkan $V$ adalah ruang vektor dan andaikan $V = M \oplus N$. Kita mengetahui
dari~\ref{0000823} bahwa untuk setiap $\vc v \in V$ terdapat vektor-vektor tunggal $\vc m \in M$
dan $\vc n \in N$ sehingga $\vc v = \vc m + \vc n$. Definisikan fungsi $\sbsb E{MN}\colon V \sto V$ dengan $\sbsb E{MN}v = n$.
Fungsi $\sbsb E{MN}$ adalah
 \index{proyeksi!dalam ruang vektor}%
 \index{E@$\sbsb E{MN}$ (proyeksi sepanjang $M$ ke $N$)}%
\df{proyeksi} $V$ \df{sepanjang} $M$ \df{ke}~$N$. (Kita sering menulis $E$ untuk $\sbsb E{MN}$. Namun ingat bahwa
$E$ bergantung pada $M$ dan~$N$.)
\end{defn}

\begin{prop} Misalkan $V$ adalah ruang vektor dan andaikan $V = M \oplus N$. Jika $E$ adalah
proyeksi $V$ sepanjang $M$ ke $N$, maka
 \begin{enumerate}[label=\rm{(\alph*)}]
  \item $E$ linear;
 \index{idempoten}%
  \item $E^2 = E$ \quad (yaitu, $E$ \df{idempoten});
  \item $\ran E = N$; dan
  \item $\ker E = M$.
 \end{enumerate}
\end{prop}

\begin{prop} Misalkan $V$ adalah ruang vektor dan andaikan $E\colon V \sto V$ adalah fungsi yang
memenuhi
 \begin{enumerate}[label=\rm{(\alph*)}]
  \item $E$ linear, dan
  \item $E^2 = E$.
\end{enumerate}
Maka
 \[V = \ker E \oplus \ran E\]
dan $E$ adalah proyeksi $V$ sepanjang $\ker E$ ke~$\ran E$.
\end{prop}

Penting untuk dicatat bahwa konsekuensi langsung dari dua proposisi terakhir adalah bahwa
fungsi $T\colon V \sto V$ dari ruang vektor berdimensi hingga ke dirinya sendiri merupakan
proyeksi jika dan hanya jika linear dan idempoten.






\begin{prop} Misalkan $T\colon V \sto W$ adalah transformasi linear antara ruang-ruang vektor. Maka
  \begin{enumerate}[label=\rm{(\alph*)}]
    \item $T$ memiliki invers kiri jika dan hanya jika injektif; dan
    \item $T$ memiliki invers kanan jika dan hanya jika surjektif.
  \end{enumerate}
\end{prop}











\begin{prop} Misalkan $V$ adalah ruang vektor dan andaikan $V = M \oplus N$. Jika $E$ adalah
proyeksi $V$ sepanjang $M$ ke $N$, maka $I - E$ adalah proyeksi $V$ sepanjang $N$
ke~$M$.
\end{prop}

Seperti baru saja kita lihat, jika $E$ adalah proyeksi pada ruang vektor $V$, maka operator
identitas pada $V$ dapat ditulis sebagai jumlah dua proyeksi $E$ dan $I - E$ yang
jangkauannya membentuk dekomposisi jumlah langsung ruang tersebut, $V = \ran E \oplus \ran
(I - E)$. Kita dapat memperumum hal ini ke lebih dari dua proyeksi.

\begin{defn} Andaikan pada ruang vektor $V$ terdapat operator-operator proyeksi $E_1$,
\dots, $E_n$ sehingga
\begin{enumerate}
    \item $I_V = E_1 + E_2 + \dots + E_n$ dan
    \item $E_iE_j = 0$ apabila $i \ne j$.
\end{enumerate}
Kita mengatakan bahwa $I_V = E_1 + E_2 + \dots + E_n$ adalah
 \index{resolusi identitas!dalam ruang vektor}%
\df{resolusi identitas}.
\end{defn}

\begin{prop} Jika $I_V = E_1 + E_2 + \dots + E_n$ adalah resolusi identitas pada ruang
vektor $V$, maka $V = \bigoplus_{k=1}^n \ran E_k$.
\end{prop}

\begin{exam} Misalkan $P$ adalah bidang dalam $\R^3$ yang persamaannya $x - z = 0$ dan $L$ adalah garis
yang persamaannya $y = 0$ dan $x = -z$. Misalkan $E$ adalah proyeksi $\R^3$ sepanjang~$L$
ke $P$ dan $F$ adalah proyeksi $\R^3$ sepanjang~$P$ ke~$L$. Maka
 \[[E] = \begin{bmatrix}
              \frac12  & 0 & \frac12 \\
                  0    & 1 &   0     \\
              \frac12  & 0 & \frac12
         \end{bmatrix} \qquad \text{ dan } \qquad
   [F] = \begin{bmatrix}
              \frac12  & 0 & -\frac12 \\
                  0    & 0 &   0     \\
             -\frac12  & 0 &  \frac12
         \end{bmatrix}\,. \]
\end{exam}

\begin{defn} Skalar $\lambda \in \K$ adalah
 \index{nilai eigen}%
\df{nilai eigen} dari operator $T$ pada ruang vektor $V$ jika $\ker(T - \lambda I_V)$
memuat vektor tak nol. Setiap vektor tersebut adalah
 \index{vektor eigen}%
\df{vektor eigen} dari $T$ yang berkaitan dengan $\lambda$, dan $\ker(T - \lambda I_V)$ adalah
 \index{ruang eigen}
\df{ruang eigen} dari $T$ yang berkaitan dengan~$\lambda$. Himpunan semua nilai eigen
operator $T$ disebut
 \index{spektrum titik}%
 \index{spektrum!titik}%
\df{spektrum titik} dan dinotasikan dengan~$\sigma_p(T)$.
\end{defn}

Jika $M$ adalah matriks $n \times n$, maka $\det(M - \lambda I_n)$ (dengan $I_n$ matriks identitas $n
\times n$) adalah polinomial dalam $\lambda$ berderajat~$n$. Inilah
\df{polinomial karakteristik} dari~$M$. Cara baku untuk menghitung nilai-nilai eigen suatu
operator $T$ pada ruang vektor berdimensi hingga adalah mencari akar-akar polinomial
karakteristik dari representasi matriksnya. Sifat multiplikatif fungsi determinan dengan mudah mengakibatkan
bahwa polinomial karakteristik operator $T$ pada ruang vektor $V$ tidak bergantung pada basis yang dipilih untuk $V$ dan,
karena itu, tidak bergantung pada representasi matriks tertentu dari~$T$ yang digunakan.

\vskip 5 pt

\begin{exam} Nilai-nilai eigen operator pada (ruang vektor real) $\R^3$ yang representasi matriksnya
adalah $ \begin{bmatrix} 0  &  0  &  2 \\
                                     0  &  2  &  0 \\
                                     2  &  0  &  0
                      \end{bmatrix}$ ialah $-2$ dan $+2$, dengan nilai yang terakhir memiliki multiplisitas (aljabar maupun
geometris)~$2$. Ruang eigen yang berkaitan dengan nilai eigen negatif dan positif, berturut-turut, adalah
\[
\spn\{(1,0,-1)\}
\quad\text{dan}\quad
\spn\{(1,0,1), (0,1,0)\}.
\]
\end{exam}

\vskip 5 pt

\begin{prop} Setiap operator pada ruang vektor kompleks berdimensi hingga memiliki nilai eigen.
\end{prop}

\begin{exam} Proposisi sebelumnya tidak berlaku untuk ruang real berdimensi hingga.
\end{exam}

\begin{proof}[\emph{Petunjuk pembuktian}] Pertimbangkan rotasi bidang.  \ns
\end{proof}

Fakta utama yang dinyatakan oleh versi ruang vektor berdimensi hingga dari
\emph{teorema spektral} ialah bahwa setiap operator yang dapat didiagonalkan pada ruang demikian dapat
ditulis sebagai kombinasi linear operator-operator proyeksi, dengan koefisien kombinasi
linear tersebut berupa nilai-nilai eigen operator dan jangkauan proyeksinya berupa
ruang eigen yang bersesuaian. Jadi, jika $T$ adalah operator yang dapat didiagonalkan pada ruang vektor
berdimensi hingga $V$, maka $V$ memiliki basis yang terdiri atas vektor-vektor eigen dari~$T$.

Berikut pernyataan formal teorema tersebut.

\begin{thm}[Teorema Spektral: versi ruang vektor berdimensi hingga] Andaikan $T$ adalah
 \index{spektral!teorema!operator yang dapat didiagonalkan pada ruang vektor}%
operator yang dapat didiagonalkan pada ruang vektor berdimensi hingga~$V$. Misalkan $\lambda_1, \dots, \lambda_n$ adalah
nilai-nilai eigen (berbeda) dari~$T$. Maka terdapat resolusi identitas $I_V =
E_1 + \dots + E_n$, dengan untuk setiap $k$ jangkauan proyeksi $E_k$ merupakan
ruang eigen yang berkaitan dengan~$\lambda_k$, dan selanjutnya
     \[ T = \lambda_1E_1 +\dots + \lambda_nE_n\,. \]
\end{thm}

\begin{proof} Pembuktian teorema ini dapat ditemukan dalam \cite{HoffmanK:1971} pada halaman~212. \ns
\end{proof}

\begin{exam} Misalkan $T$ adalah operator pada (ruang vektor real) $\R^2$ yang representasi matriksnya
adalah $\begin{bmatrix} -7 &   8 \\ -16 &  17 \end{bmatrix}$.

\vskip 3 pt

 \begin{enumerate}
  \item Polinomial karakteristik untuk~$T$ adalah
                     $\sbsb cT(\lambda) = \lambda^2 - 10\lambda + 9$.

\vskip 3 pt

  \item Ruang eigen $M_1$ yang berkaitan dengan nilai eigen $1$ adalah $\spn \{(1,1)\}$.

\vskip 3 pt

  \item Ruang eigen $M_2$ yang berkaitan dengan nilai eigen $9$ adalah $\spn \{(1,2)\}$.

\vskip 3 pt

  \item Kita dapat menulis $T$ sebagai kombinasi linear operator-operator proyeksi. Secara khusus,
   \[ T = 1\cdot E_1 + 9\cdot E_2 \text{ dengan } [E_1] = \begin{bmatrix} 2 & -1 \\ 2 & -1 \end{bmatrix}
\text{ dan } [E_2] = \begin{bmatrix} -1 & 1 \\ -2 & 2 \end{bmatrix}\,. \]

\vskip 3 pt

  \item Perhatikan bahwa jumlah $[E_1]$ dan $[E_2]$ adalah matriks identitas dan hasil kali keduanya
adalah matriks nol.

\vskip 3 pt

  \item Matriks $S = \begin{bmatrix} 1 &  1  \\ 1 &  2 \end{bmatrix}$ mendiagonalkan
$[T]$. Artinya, $S^{-1}\,[T]\,S = \begin{bmatrix} 1 &  0  \\ 0 &  9 \end{bmatrix}$.


\vskip 3 pt


  \item Matriks yang merepresentasikan $\sqrt T$ adalah $\begin{bmatrix} -1 &  2  \\ -4 &  5 \end{bmatrix}$.
 \end{enumerate}
\end{exam}

\begin{exer} Misalkan $T$ adalah operator pada $\R^3$ yang representasi matriksnya
    \[
      \begin{bmatrix}
               2  &  0  &  0 \\
              -1  &  3  &  2 \\
               1  & -1  &  0
      \end{bmatrix}.
    \]
  \begin{enumerate}
    \item Tentukan polinomial karakteristik dan polinomial minimal dari~$T$.
    \item Apa yang dapat disimpulkan dari bentuk polinomial minimal tersebut?
    \item Tentukan matriks yang mendiagonalkan~$T$. Apa bentuk diagonal $T$ yang dihasilkan oleh matriks ini?
    \item Tentukan (representasi matriks dari)~$\sqrt T$.
  \end{enumerate}
\end{exer}

\begin{defn} Operator ruang vektor $T$ disebut
 \index{nilpoten}%
\df{nilpoten} jika $T^n = \vc 0$ untuk suatu $n \in \N$.
\end{defn}

Operator pada ruang vektor berdimensi hingga tidak harus dapat didiagonalkan. Jika tidak,
seberapa dekat operator itu dengan operator yang dapat didiagonalkan? Berikut salah satu jawabannya.

\begin{thm} Misalkan $T$ adalah operator pada ruang vektor berdimensi hingga~$V$. Andaikan polinomial minimal
untuk $T$ terfaktor sepenuhnya menjadi faktor-faktor linear
   \[ m_T(x) = (x - \lambda_1)^{r_1} \dots (x - \lambda_k)^{r_k} \]
dengan $\lambda_1, \dots \lambda_k$ merupakan nilai-nilai eigen (berbeda) dari~$T$. Untuk setiap $j$, misalkan
$W_j = \ker(T - \lambda_jI)^{r_j}$ dan $E_j$ adalah proyeksi $V$ ke $W_j$ sepanjang
$W_1 + \dots + W_{j-1} + W_{j+1} + \dots + W_k$. Maka
   \[ V = W_1 \oplus W_2 \oplus \dots \oplus W_k, \]
setiap $W_j$ invarian terhadap $T$, dan $I_V = E_1 + \dots + E_k$. Selanjutnya, operator
   \[ D = \lambda_1E_1 + \dots + \lambda_kE_k \]
dapat didiagonalkan, operator
   \[ N = T - D \]
nilpoten, dan $N$ berkomutasi dengan~$D$.
\end{thm}

\begin{proof} Lihat~\cite{HoffmanK:1971}, halaman 222--223. \ns
\end{proof}

\begin{defn} Karena dalam teorema sebelumnya $T = D + N$, dengan $D$ dapat didiagonalkan dan $N$ nilpoten,
kita menyebut $D$ sebagai
 \index{dapat didiagonalkan!bagian}%
 \index{bagian!dapat didiagonalkan}%
\df{bagian yang dapat didiagonalkan} dari $T$, dan $N$ sebagai
 \index{nilpoten!bagian}%
 \index{bagian!nilpoten}%
\df{bagian nilpoten} dari~$T$.
\end{defn}

\begin{exer} Misalkan $T$ adalah operator pada $\R^3$ yang representasi matriksnya
 \[ \begin{bmatrix}
        3  &  1  & -1 \\
        2  &  2  & -1 \\
        2  &  2  &  0
  \end{bmatrix}. \]

 \begin{enumerate}
  \item Tentukan polinomial karakteristik dan polinomial minimal dari~$T$.
  \item Tentukan ruang-ruang eigen dari~$T$.
  \item Tentukan bagian yang dapat didiagonalkan $D$ dan bagian nilpoten $N$ dari~$T$.
  \item Tentukan matriks yang mendiagonalkan~$D$. Apa bentuk diagonal
dari $D$ yang dihasilkan oleh matriks ini?
  \item Tunjukkan bahwa $D$ berkomutasi dengan~$N$.
 \end{enumerate}
\end{exer}
\section{Operator Normal pada Ruang Hasil Kali Dalam}
\begin{defn} Misalkan $V$ suatu ruang vektor. Sebuah fungsi $s$ yang mengaitkan setiap pasangan vektor
$x$ dan $y$ dalam $V$ dengan suatu skalar
 \index{s@$s(x,y)$ (hasil kali dalam semu)}%
 \index{<s@$s(x,y)$ (hasil kali dalam semu)}%
$s(x,y)$ adalah
 \index{hasil kali dalam semu}%
 \index{hasil kali dalam semu!ruang}%
 \index{hasil kali!dalam semu}%
 \index{ruang!hasil kali dalam semu}%
\df{hasil kali dalam semu} pada $V$ apabila untuk setiap $x$, $y$, $z \in V$ dan $\alpha \in \K$ keempat syarat berikut dipenuhi:
 \begin{enumerate}
   \item $s(x + y , z) = s(x,z) + s(y,z)$;
   \item $s(\alpha x,y) = \alpha s(x,y)$;
   \item $s(x,y) = \conj{s(y,x)}$; dan
   \item $s(x,x) \ge 0$.
 \end{enumerate}
 Jika, selain itu, fungsi $s$ memenuhi
  \begin{enumerate}
    \item Untuk setiap $x$ tak nol dalam $V$, berlaku $s(x,x) > 0$.
  \end{enumerate}
maka $s$ adalah
 \index{hasil kali dalam}%
 \index{hasil kali dalam!ruang}%
 \index{hasil kali!dalam}%
 \index{ruang!hasil kali dalam}%
\df{hasil kali dalam} pada~$V$. Hasil kali dalam dua vektor $x$ dan $y$ biasanya akan kita tuliskan sebagai
 \index{<bracpnt@$\langle x,y \rangle$ (hasil kali dalam)}%
$\langle x,y \rangle$, bukan $s(x,y)$.

\noindent Garis atas pada (c) menyatakan konjugasi kompleks (sehingga berlebihan dalam kasus ruang vektor real).
Syarat (a) dan (b) memperlihatkan bahwa hasil kali dalam semu bersifat linear dalam variabel pertamanya. Syarat (a) dan (b) dalam
proposisi~\ref{00015001} menyatakan bahwa hasil kali dalam kompleks bersifat
 \index{konjugat!linear}%
 \index{linear!konjugat}%
\df{linear konjugat} dalam variabel keduanya. Jika suatu fungsi bernilai skalar dengan dua variabel pada ruang hasil kali dalam semu kompleks
bersifat linear dalam satu variabel dan linear konjugat dalam variabel lainnya, fungsi itu sering disebut
 \index{seskuilinear}%
 \index{linear!seskuilinear}%
\df{seskuilinear}. (Awalan ``sesqui-'' berarti ``satu setengah''.) Istilah \emph{seskuilinear} juga akan kita gunakan untuk
fungsi bilinear pada ruang hasil kali dalam semu real. Secara bersama-sama, syarat (a)--(d) menyatakan bahwa hasil kali dalam semu itu merupakan
\emph{bentuk seskuilinear semidefinit positif dan simetris konjugat}; bersama dengan syarat ketat tambahan di atas, bentuk itu definit positif.
\end{defn}

\begin{notn}\label{norm_notn} Jika $V$ adalah ruang vektor yang telah dilengkapi dengan
hasil kali dalam semu dan $x \in V$, kita memperkenalkan singkatan
  \[ \norm x := \sqrt{s(x,x)} \]
yang dibaca \emph{norma $x$} atau \emph{panjang~$x$}. (Istilah yang agak
optimistis ini dibenarkan, setidaknya ketika $s$ merupakan hasil kali dalam, oleh proposisi~\ref{00015025} di bawah.)
\end{notn}

\begin{prop}\label{00015001} Jika $x$, $y$, dan $z$ merupakan vektor dalam suatu ruang dengan hasil kali dalam semu $s$ dan $\alpha \in \K$, maka
 \begin{enumerate}[label=\rm{(\alph*)}]
   \item $s(x,y + z)  = s(x,y) + s(x,z)$,
   \item $s(x,\alpha y) = \conj\alpha s(x,y)$, dan
   \item jika $s$ merupakan hasil kali dalam, maka $s(x,x) = 0$ jika dan hanya jika $x = \vc 0$.
 \end{enumerate}
\end{prop}

\begin{exam}\label{000150012} Untuk vektor $x = (x_1,x_2,\dots,x_n)$ dan $y = (y_1,y_2,\dots,y_n)$
yang termasuk dalam $\K^n$, definisikan
   \[ \langle x,y \rangle = \sum_{k=1}^n x_k \conj{y_k}\,. \]
Maka
 \index{K@$\K^n$!sebagai ruang hasil kali dalam}%
 \index{hasil kali dalam!ruang!$\K^n$ sebagai}%
$\K^n$ merupakan ruang hasil kali dalam.
\end{exam}

\begin{exam}\label{000150013} Misalkan $l_2 = l_2(\N)$ adalah himpunan semua barisan bilangan
kompleks yang terjumlah secara kuadrat. (Suatu barisan $x = (x_k)_{k=1}^\infty$ disebut
 \index{terjumlah secara kuadrat}%
 \index{dapat dijumlahkan!kuadrat}%
 \index{l@$l_2 = l_2(\N)$!barisan terjumlah secara kuadrat}%
 \index{l@$l_2 = l_2(\N)$!sebagai ruang hasil kali dalam}%
 \index{hasil kali dalam!ruang!$l_2$ sebagai}%
\df{terjumlah secara kuadrat} jika $\sum_{k=1}^\infty \abs{x_k}^2 < \infty$.) (Operasi ruang vektornya didefinisikan
titik demi titik.) Untuk vektor $x = (x_1,x_2,\dots)$ dan $y = (y_1,y_2,\dots)$ dalam $l_2$, definisikan
   \[ \langle x,y \rangle = \sum_{k=1}^\infty x_k \conj{y_k}\,. \]
Maka $l_2$ merupakan ruang hasil kali dalam. (Antara lain, harus ditunjukkan bahwa
deret dalam definisi di atas benar-benar konvergen.) Ruang serupa adalah $l_2(\Z)$, yakni himpunan semua
 \index{l@$l_2(\Z)$!barisan dua sisi terjumlah secara kuadrat}%
 \index{l@$l_2(\Z)$!sebagai ruang hasil kali dalam}%
 \index{hasil kali dalam!ruang!$l_2(\Z)$ sebagai}%
barisan dua sisi yang terjumlah secara kuadrat $x = (\dots,x_{-2},x_{-1},x_0,x_1,x_2,\dots)$ dengan hasil kali dalam yang sudah jelas.
\end{exam}

\begin{exam}\label{000150014} Untuk $a < b$, misalkan $\fml C([a,b])$ adalah keluarga semua
 \index{C@$\fml C([a,b])$!sebagai ruang hasil kali dalam}%
 \index{hasil kali dalam!ruang!$\fml C([a,b])$ sebagai}%
fungsi kontinu bernilai kompleks pada interval $[a,b]$. Untuk setiap $f$, $g \in \fml C([a,b])$, definisikan
   \[ \langle f,g \rangle = \int_a^b f(x) \conj{g(x)}\,dx. \]
Maka $\fml C([a,b])$ merupakan ruang hasil kali dalam.
\end{exam}

Berikut adalah fakta tunggal yang paling berguna mengenai hasil kali dalam semu. Fakta ini merupakan suatu ketaksamaan yang secara beragam
dinisbatkan kepada Cauchy, Schwarz, Bunyakowsky, atau gabungan nama ketiganya.

\begin{thm}[Ketaksamaan Schwarz]\label{00015002} Misalkan $s$ suatu hasil kali dalam semu yang didefinisikan pada ruang vektor~$V$.
 \index{ketaksamaan Schwarz}%
 \index{ketaksamaan!Schwarz}%
Maka ketaksamaan
   \[ \abs{s(x,y)} \le \norm x \,\norm y. \]
berlaku untuk semua vektor $x$ dan~$y$ dalam~$V$.
\end{thm}

\begin{proof}[\emph{Petunjuk pembuktian}] Tetapkan $x$, $y \in V$. Untuk setiap $\alpha \in \K$, kita mengetahui bahwa
  \begin{equation}\label{Schwarz_ineq}
      0 \le s(x - \alpha y, x - \alpha y) .
  \end{equation}
Uraikan ruas kanan~\eqref{Schwarz_ineq} menjadi empat suku. Jika $s(y,x)=0$, kesimpulannya langsung. Jika tidak,
tuliskan $s(y,x)$ dalam bentuk polar: $s(y,x) = r e^{i\theta}$, dengan $r > 0$ dan $\theta \in \R$.
Kemudian, dalam ketaksamaan yang diperoleh, tinjau $\alpha$ yang berbentuk $e^{-i\theta}t$ dengan
$t \in \R$. Perhatikan bahwa kini ruas kanan~\eqref{Schwarz_ineq} merupakan polinom kuadrat dalam~$t$. Apa yang dapat Anda katakan tentang
diskriminannya? \ns
\end{proof}

\begin{prop} Misalkan $V$ suatu ruang vektor dengan hasil kali dalam semu $s$ dan $z \in V$. Maka $s(z,z) = 0$
jika dan hanya jika $s(z,y) = 0$ untuk semua $y \in V$.
\end{prop}

\begin{prop} Misalkan $V$ suatu ruang vektor yang telah dilengkapi dengan hasil kali dalam semu $s$. Maka himpunan
$L := \{z \in V\colon s(z,z) = 0\}$ merupakan subruang vektor dari $V$, dan ruang vektor hasil bagi $V/L$ dapat
dijadikan ruang hasil kali dalam dengan mendefinisikan
    \[ \langle\, [x],[y]\, \rangle := s(x,y) \]
untuk semua $[x]$, $[y] \in V/L$.
\end{prop}

Proposisi berikut menunjukkan bahwa hasil kali dalam bersifat kontinu dalam variabel pertamanya. Simetri konjugat kemudian menjamin
kekontinuan dalam variabel keduanya.

\begin{prop} Jika $(x_n)$ merupakan barisan dalam ruang hasil kali dalam $V$ yang konvergen ke suatu
vektor $a \in V$, maka $\langle x_n,y \rangle \sto \langle a,y \rangle$ untuk setiap $y \in V$.
\end{prop}

\begin{exam} Jika $(a_k)$ merupakan barisan bilangan real yang terjumlah secara kuadrat, maka deret $\sum_{k=1}^\infty k^{-1}a_k$
konvergen mutlak.
\end{exam}

\begin{exer} Misalkan $0 < a < b$.
 \begin{enumerate}
  \item Jika $f$ dan $g$ merupakan fungsi kontinu bernilai real pada
interval $[a,b]$, maka
   \[\left(\int_a^b f(x)g(x)\,dx\right)^2
             \le \int_a^b\bigl(f(x)\bigr)^2\,dx \cdot
                       \int_a^b\bigl(g(x)\bigr)^2\,dx.\]

\vskip 3 pt

  \item Gunakan bagian (a) untuk mencari bilangan $M$ dan $N$ (yang bergantung pada $a$
dan~$b$) sedemikian sehingga
   \[M(b-a) \le \ln\left(\frac ba\right) \le N(b-a).\]

\vskip 3 pt

  \item Gunakan bagian (b) untuk menunjukkan bahwa $2.5 \le e \le 3$.
 \end{enumerate}
\end{exer}

\begin{proof}[\emph{Petunjuk pembuktian}] Untuk (b), cobalah $f(x) = \sqrt x$ dan $g(x) = \dfrac1{\sqrt x}$. Kemudian cobalah
$f(x) = \dfrac1x$ dan $g(x) = 1$. Untuk (c), cobalah $a = 1$ dan $b = 3$. Kemudian cobalah $a = 2$ dan $b = 5$. \ns
\end{proof}

\begin{defn} Misalkan $V$ suatu ruang vektor. Sebuah fungsi
 \index{<unopn@$\norm x$ (norma suatu vektor~$x$)}%
$\norm{\hphantom{x}} \colon V \sto \R \colon x \mapsto \norm{x}$ adalah
 \index{norma}%
\df{norma} pada $V$ jika
 \begin{enumerate}
  \item $\norm{x + y} \le \norm{x} + \norm{y}$ \qquad  untuk semua $x$, $y \in V$;
  \item $\norm{\alpha x} = \abs{\alpha}\,\norm{x}$\qquad untuk semua $x \in V$ dan
$\alpha \in \K$; dan
  \item jika $\norm{x} = 0$, maka $x = \vc 0$.
 \end{enumerate}
Ungkapan $\norm{x}$ dapat dibaca sebagai ``\emph{norma} $x$'' atau ``\emph{panjang}
 \index{panjang}%
$x$''. Jika fungsi $\norm{\hphantom{x}}$ memenuhi (a) dan (b)
di atas (tetapi mungkin tidak memenuhi~(c)), fungsi itu merupakan
 \index{seminorma}%
\df{seminorma} pada~$V$.

Ruang vektor yang telah dilengkapi dengan suatu norma disebut
 \index{bernorma!ruang linear}%
 \index{bernorma!ruang vektor|seeonly{ruang linear bernorma}}%
 \index{ruang!linear bernorma}%
\df{ruang linear bernorma} (atau
 \index{vektor!ruang!bernorma}%
\df{ruang vektor bernorma}). Vektor dalam ruang linear bernorma yang normanya $1$ disebut
 \index{satuan!vektor}%
 \index{vektor!satuan}%
\df{vektor satuan}.
\end{defn}

\begin{exer} Tunjukkan mengapa langsung jelas dari definisi bahwa $\norm{\vc 0} = 0$ dan bahwa norma
(dan seminorma) hanya dapat mengambil nilai nonnegatif.
\end{exer}

Setiap ruang hasil kali dalam merupakan ruang linear bernorma. Dalam bahasa yang lebih tepat, hasil kali dalam pada suatu ruang vektor menginduksi
 \index{norma!diinduksi oleh hasil kali dalam}%
 \index{diinduksi!norma}
suatu norma pada ruang tersebut.

\begin{prop}\label{00015025} Misalkan $V$ suatu ruang hasil kali dalam. Pemetaan $x \mapsto \norm
x$ pada $V$ yang didefinisikan dalam~\ref{norm_notn} merupakan norma pada~$V$.
\end{prop}

Dan setiap ruang linear bernorma merupakan ruang metrik. Ingat kembali definisi berikut.

\begin{defn}\label{C019614} Misalkan $M$ suatu himpunan tak kosong. Sebuah fungsi $d \colon M \times M \sto \R$ adalah
 \index{metrik}%
\df{metrik} (atau
 \index{jarak!fungsi}%
 \index{d@$d(x,y)$ (jarak antara titik $x$ dan $y$)}%
\df{fungsi jarak}) pada $M$ jika untuk semua $x$, $y$, $z \in M$
 \begin{enumerate}
     \item $d(x,y) = d(y,x)$,
     \item $d(x,y) \le d(x,z) + d(z,y)$, \qquad (
 \index{ketaksamaan segitiga}%
 \index{ketaksamaan!segitiga}%
\emph{ketaksamaan segitiga})
     \item $d(x,x) = 0$, dan
     \item $d(x,y) = 0$ hanya jika $x = y$.
 \end{enumerate}
Jika $d$ merupakan metrik pada himpunan $M$, kita mengatakan bahwa pasangan $(M,d)$ merupakan
 \index{metrik!ruang}%
 \index{ruang!metrik}%
\df{ruang metrik}. Notasi ``$d(x,y)$'' dibaca ``jarak antara $x$
dan~$y$''. (Seperti biasa, kita mengganti perumusan yang benar ``Misalkan $(M,d)$ suatu ruang metrik \dots''
dengan ungkapan ``Misalkan $M$ suatu ruang metrik \dots''.)

Jika $d \colon M \times M \sto \R$ memenuhi syarat (1)--(3), tetapi tidak memenuhi (4), maka $d$ merupakan
 \index{pseudometrik}%
\df{pseudometrik} pada~$M$.
\end{defn}

Seperti disebutkan di atas, setiap ruang linear bernorma merupakan ruang metrik. Secara lebih tepat, norma pada suatu ruang vektor
menginduksi suatu
 \index{metrik!diinduksi oleh norma}%
metrik $d$, yang didefinisikan oleh $d(x,y) = \norm{x - y}$. Dengan kata lain, jarak antara
dua vektor adalah panjang selisih keduanya.
  \[ \xy
       \Atriangle/<-`<-`>/[``;x`x-y`y]
  \endxy \]
Jika tidak ada metrik lain yang ditentukan, kita selalu memandang ruang linear bernorma sebagai ruang metrik
 \index{diinduksi!metrik}%
di bawah metrik terinduksi ini. Metrik tersebut disebut \emph{metrik yang diinduksi oleh norma}.

\begin{exam}\label{0001503} Misalkan $V$ suatu ruang linear bernorma. Definisikan
$d\colon V \times V \sto \R$ dengan $d(x,y) = \norm{x - y}$. Maka $d$ merupakan metrik pada~$V$.
Jika $V$ hanya merupakan ruang berseminorma, maka $d$ merupakan pseudometrik.
\end{exam}

\begin{defn} Vektor $x$ dan $y$ dalam ruang hasil kali dalam $V$ disebut
 \index{ortogonal}%
\df{ortogonal} (atau
 \index{tegak lurus}%
\df{tegak lurus}) jika $\langle x,y \rangle = 0$. Dalam hal ini kita menulis
 \index{<binrelperp@$x \perp y$ (vektor ortogonal)}%
$x \perp y$. Subhimpunan $A$ dan $B$ dari $V$ disebut \df{ortogonal} jika $a \perp b$ untuk setiap $a
\in A$ dan $b \in B$. Dalam hal ini kita menulis $A \perp B$.
\end{defn}

\begin{defn} Jika $M$ dan $N$ merupakan subruang dari ruang hasil kali dalam $V$, kita menggunakan notasi
$V = M \oplus N$ untuk menyatakan bukan hanya bahwa $V$ merupakan jumlah langsung (ruang vektor) dari $M$
dan $N$, melainkan juga bahwa $M$ dan $N$ ortogonal. Dengan demikian, kita mengatakan bahwa $V$ merupakan
 \index{internal!jumlah langsung ortogonal}%
 \index{langsung!jumlah!ortogonal}%
 \index{ortogonal!jumlah langsung!ruang hasil kali dalam}%
\df{jumlah langsung ortogonal} (\df{internal}) dari $M$ dan~$N$.
\end{defn}

\begin{prop} Misalkan $a$ suatu vektor dalam ruang hasil kali dalam $V$. Maka $a \perp x$ untuk setiap
$x \in V$ jika dan hanya jika $a = 0$.
\end{prop}

\begin{prop} Dalam ruang hasil kali dalam, $x \perp y$ jika dan hanya jika $\norm{x + \alpha y} = \norm{x - \alpha y}$
untuk semua skalar $\alpha$.
\end{prop}

\begin{prop}[Teorema Pythagoras]\label{00015037} Jika $x \perp y$ dalam suatu ruang hasil
 \index{teorema Pythagoras}%
kali dalam, maka
   \[ \norm{x + y}^2 = \norm x^2 + \norm y^2\,. \]
\end{prop}

\begin{defn}\label{0001504} Misalkan $V$ dan $W$ merupakan ruang hasil kali dalam di atas medan skalar yang sama. Untuk $(v,w)$ dan
$(v',w')$ dalam $V \times W$ dan $\alpha \in \K$, definisikan
 \begin{align*}
      (v,w) + (v',w') &= (v + v', w + w') \\
\intertext{dan}
         \alpha(v,w) &= (\alpha v,\alpha w)\,.
 \end{align*}
Hasilnya adalah suatu ruang vektor, yaitu \emph{jumlah langsung (eksternal)} dari $V$ dan~$W$.
Untuk menjadikannya ruang hasil kali dalam, definisikan
   \[ \langle (v,w),(v',w') \rangle = \langle v,v' \rangle + \langle w,w' \rangle. \]
Definisi ini menjadikan jumlah langsung $V$ dan $W$ suatu ruang hasil kali dalam. Ruang tersebut merupakan
 \index{eksternal!jumlah langsung}%
 \index{langsung!jumlah!eksternal}%
 \index{langsung!jumlah!ruang hasil kali dalam}%
 \index{ortogonal!jumlah langsung!ruang hasil kali dalam}%
 \index{<binopdirectsum@$\oplus$ (jumlah langsung ortogonal)}%
\df{jumlah langsung} dari $V$ dan $W$ yang \df{ortogonal eksternal}, dan dinotasikan dengan $V \oplus W$.
\end{defn}

Perhatikan bahwa notasi $\oplus$ yang sama digunakan baik untuk jumlah langsung internal maupun eksternal,
serta baik untuk jumlah langsung ruang vektor (lihat definisi~\ref{000082}) maupun jumlah langsung ortogonal.
Jadi, ketika melihat simbol $V \oplus W$, penting untuk mengetahui kategori mana yang sedang kita gunakan:
ruang vektor atau ruang hasil kali dalam. Hal ini khususnya penting karena kata ``ortogonal'' lazim dihilangkan
sebagai pewatas ``jumlah langsung'', bahkan ketika sifat ortogonal memang dimaksudkan.

\begin{exam} Dalam $\R^2$, misalkan $M$ adalah sumbu-$x$ dan $L$ adalah garis dengan persamaan $y = x$.
Jika $\R^2$ dipandang sebagai ruang vektor (real), penulisan $\R^2 = M
\oplus L$ benar. Sebaliknya, jika $\R^2$ dipandang sebagai ruang hasil kali dalam (real), maka
$\R^2 \ne M \oplus L$ (karena $M$ dan $L$ tidak tegak lurus).
\end{exam}

\begin{prop} Misalkan $V$ suatu ruang hasil kali dalam. Hasil kali dalam pada $V$, yang dipandang sebagai pemetaan
dari $V \oplus V$ ke $\K$, bersifat kontinu. Demikian pula normanya, yang dipandang sebagai pemetaan dari $V$
ke $\R$.
\end{prop}

Mengenai pembuktian proposisi di atas, perhatikan bahwa pemetaan $(v,v') \mapsto
\norm v + \norm{v'}$, $(v,v') \mapsto \sqrt{\norm v^2 + \norm{v'}^2}$, dan $(v,v')
\mapsto \max\{\norm v, \norm{v'}\}$ semuanya merupakan norma pada $V \oplus V$. Norma manakah yang diinduksi
oleh hasil kali dalam pada $V \oplus V$? Mengapa pilihan norma tidak berpengaruh ketika kita membuktikan
bahwa hasil kali dalam tersebut kontinu?

\begin{prop}[Hukum jajaran genjang]\label{parallelogram_law} Jika $x$ dan $y$ merupakan vektor dalam suatu ruang hasil kali
 \index{hukum jajaran genjang}%
dalam, maka
   \[ \norm{x + y}^2 + \norm{x - y}^2 = 2\norm x^2 + 2\norm y^2\,. \]
\end{prop}

Walaupun setiap hasil kali dalam menginduksi suatu norma, tidak setiap norma berasal dari hasil kali dalam.

\begin{exam}\label{00015061} Tidak ada hasil kali dalam pada ruang $\fml C([0,1])$ yang menginduksi norma seragam.
\end{exam}

\begin{proof}[\emph{Petunjuk pembuktian}] Gunakan proposisi sebelumnya. \ns
\end{proof}

\begin{prop}[Identitas polarisasi]\label{0001507} Jika $x$ dan $y$ merupakan vektor
 \index{identitas polarisasi}%
dalam ruang hasil kali dalam kompleks, maka
 \[ \langle x,y \rangle = \tfrac14 (\norm{x + y}^2 - \norm{x - y}^2
                       + i\,\norm{x + iy}^2 - i\,\norm{x - iy}^2)\,. \]
\end{prop}

Apakah identitas yang benar untuk ruang hasil kali dalam \emph{real}?

\begin{notn}\label{0001511} Misalkan $V$ suatu ruang hasil kali dalam, $x \in V$, dan $A$,
$B \subseteq V$. Jika $x \perp a$ untuk setiap $a \in A$, kita menulis $x \perp A$; dan jika $a
\perp b$ untuk setiap $a \in A$ dan $b \in B$, kita menulis $A \perp B$.
Kita mendefinisikan $A^\perp$, yaitu
 \index{ortogonal!komplemen}%
 \index{komplemen!ortogonal}%
 \index{<unopur@$A^\perp$ (komplemen ortogonal)}%
\df{komplemen ortogonal} dari $A$, sebagai $\{x \in V\colon x \perp A\}$. Karena frasa ``komplemen ortogonal dari $A$''
agak panjang, terutama jika digunakan berulang kali, simbol ``$A^\perp$'' biasanya dilafalkan ``A perp''.
Kita menulis $A^{\perp\perp}$ untuk $\left(A^\perp\right)^\perp$.
\end{notn}

\begin{prop}\label{0001512} Jika $A$ merupakan subhimpunan dari ruang hasil kali dalam $V$, maka
$A^\perp$ merupakan subruang linear tertutup dari~$V$.
\end{prop}

\begin{prop}[Ortonormalisasi Gram--Schmidt] Jika $(v^k)$ merupakan barisan bebas linear dalam
 \index{ortonormalisasi Gram--Schmidt}%
 \index{ortonormalisasi}%
ruang hasil kali dalam~$V$, maka terdapat barisan ortonormal $(e^k)$ dalam~$V$ sedemikian sehingga
$\spn\{v^1, \dots, v^n\} = \spn\{e^1, \dots, e^n\}$ untuk setiap $n \in \N$.
\end{prop}

\begin{cor}\label{0001514} Jika $M$ merupakan subruang dari ruang hasil kali dalam berdimensi
hingga $V$, maka $V = M \oplus M^\perp$.
\end{cor}

Akan berguna untuk mengatakan bahwa suatu barisan \emph{pada akhirnya} memiliki sifat tertentu atau
bahwa barisan itu \emph{sering} memiliki sifat tersebut.

\begin{defn}\label{C013121} Misalkan $(x_n)$ suatu barisan elemen dari himpunan $S$ dan $P$ suatu
sifat yang mungkin dimiliki anggota $S$. Kita mengatakan bahwa barisan $(x_n)$
 \index{pada akhirnya}%
\df{pada akhirnya} memiliki sifat $P$ jika terdapat $n_0 \in \N$ sedemikian sehingga $x_n$ memiliki
sifat $P$ untuk setiap $n \ge n_0$. (Cara lain untuk menyatakan hal yang sama: $x_n$ memiliki
sifat $P$ untuk semua kecuali sejumlah hingga~$n$.)
\end{defn}

\begin{exam}\label{exam13122} Misalkan $l_c$ menyatakan
 \index{l@$l_c$!barisan yang pada akhirnya nol}%
 \index{l@$l_c$!sebagai ruang vektor}%
ruang vektor yang terdiri atas semua barisan bilangan kompleks yang pada akhirnya nol;
yaitu, himpunan semua barisan yang hanya memiliki sejumlah hingga entri tak nol. Barisan tersebut juga disebut barisan dengan
 \index{hingga!tumpuan}%
 \index{tumpuan!hingga}%
\emph{bertumpuan hingga}. Operasi ruang vektornya didefinisikan titik demi titik. Kita menjadikan ruang $l_c$ suatu
 \index{l@$l_c$!sebagai ruang hasil kali dalam}%
 \index{hasil kali dalam!ruang!$l_c$ sebagai}%
ruang hasil kali dalam dengan
mendefinisikan $\langle a,b \rangle = \langle\, (a_n),(b_n)\, \rangle := \sum_{k=1}^\infty a_k \conj{b_k}$.
\end{exam}

\begin{defn}\label{C013131} Misalkan $(x_n)$ suatu barisan elemen dari himpunan $S$ dan $P$ suatu
sifat yang mungkin dimiliki anggota $S$. Kita mengatakan bahwa barisan $(x_n)$
 \index{sering}%
\df{sering} memiliki sifat $P$ jika untuk setiap $k \in \N$ terdapat $n \ge k$ sedemikian sehingga
$x_n$ memiliki sifat~$P$. (Perumusan yang ekuivalen: $x_n$ memiliki sifat $P$ untuk tak hingga banyak
nilai~$n$.)
\end{defn}

\begin{exam}\label{exam13132} Misalkan $V = l_2$ merupakan ruang hasil kali dalam semua barisan bilangan kompleks
yang terjumlah secara kuadrat (lihat contoh~\ref{000150013}) dan $M = l_c$ (lihat contoh~\ref{exam13122}). Maka kesimpulan~\ref{0001514} gagal.
\end{exam}

\begin{defn}\label{def_alg_dual} Suatu
 \index{linear!fungsional}%
 \index{fungsional!linear}%
\df{fungsional linear} pada ruang vektor $V$ adalah pemetaan linear dari $V$ ke medan
skalarnya. Himpunan semua fungsional linear pada $V$ merupakan
 \index{dual!aljabar}%
 \index{aljabar!ruang dual}%
 \index{ruang!dual aljabar}%
\df{ruang dual} (\df{aljabar}) dari~$V$. Kita akan menggunakan
 \index{<unopur@$V^\#$ (ruang dual aljabar)}%
 \index{V@$V^\#$ (ruang dual aljabar)}%
notasi $V^{\#}$ untuk ruang dual aljabar.
\end{defn}

\begin{thm}[Teorema Riesz--Fr\'echet]\label{0001601} Jika $f \in V^{\#}$ dengan $V$ suatu
 \index{teorema Riesz--Fr\'echet!versi berdimensi hingga}%
ruang hasil kali dalam berdimensi hingga, maka terdapat vektor unik $a$ dalam $V$ sedemikian
sehingga
    \[ f(x) = \langle x,a \rangle \]
untuk semua $x$ dalam $V$.
\end{thm}

Sebentar lagi kita akan membuktikan (dalam Teorema~\futurexref{4.5.2}{000342}) bahwa setiap fungsional linear kontinu pada ruang hasil
kali dalam \emph{sebarang} memiliki representasi di atas. Versi berdimensi hingga yang dinyatakan di sini merupakan kasus khusus, karena
setiap pemetaan linear pada ruang hasil kali dalam berdimensi hingga bersifat kontinu (lihat Proposisi~\futurexref{3.3.9}{00015033}).

\begin{defn}\label{0001602} Misalkan $T\colon V \sto W$ suatu transformasi linear antara ruang
hasil kali dalam kompleks. Jika terdapat fungsi $T^*\colon W \sto V$ yang memenuhi
   \[ \langle T\vc v,\vc w \rangle = \langle \vc v,T^*\vc w \rangle \]
untuk semua $\vc v \in V$ dan $\vc w \in W$, maka $T^*$ merupakan
 \index{adjoin!operator pada ruang hasil kali dalam}%
 \index{<unopurstar@$T^*$ (adjoin pada ruang hasil kali dalam)}%
 \index{T@$T^*$ (adjoin pada ruang hasil kali dalam)}%
\df{adjoin} (atau \df{transpos konjugat}, atau \df{konjugat Hermitian}) dari~$T$.
\end{defn}

\begin{prop}\label{00016021} Jika $T\colon V \sto W$ merupakan pemetaan linear antara ruang
hasil kali dalam berdimensi hingga, maka $T^*$ ada.
\end{prop}

\begin{proof}[\emph{Petunjuk pembuktian}] Fungsional $\phi\colon V \times W \sto \C \colon
(v,w) \mapsto \langle Tv,w \rangle$ bersifat seskuilinear. Tetapkan $w \in W$ dan definisikan $\phi_w
\colon V \sto \C \colon v \mapsto \phi(v,w)$. Maka $\phi_w \in V^{\#}$. Gunakan
\emph{teorema Riesz--Fr\'echet}~(\ref{0001601}). \ns
\end{proof}

\begin{prop} Jika $T\colon V \sto W$ merupakan pemetaan linear antara ruang hasil kali dalam
berdimensi hingga, maka fungsi $T^*$ yang didefinisikan di atas bersifat linear dan $T^{**} = T$.
\end{prop}

\begin{thm}[Teorema dasar aljabar linear]\label{00016025} Jika
 \index{dasar!teorema aljabar linear}%
$T\colon V \sto W$ merupakan pemetaan linear antara ruang hasil kali dalam berdimensi hingga, maka
   \[ \ker T^* = (\ran T)^\perp \quad \text{ dan } \quad \ran T^* = (\ker T)^\perp\,. \]
\end{thm}

\begin{defn}\label{000161} Suatu operator $U$ pada ruang hasil kali dalam disebut
 \index{uniter!operator}%
 \index{operator!uniter}%
\df{uniter} jika $UU^* = U^*U = I$, yaitu jika $U^* = U^{-1}$.
\end{defn}

\begin{defn}\label{0001612} Dua operator $R$ dan $T$ pada ruang hasil kali dalam $V$ disebut
 \index{uniter!ekuivalensi}%
 \index{ekuivalen!secara uniter}%
\df{ekuivalen secara uniter} jika terdapat operator uniter $U$ pada $V$ sedemikian sehingga $R =
U^*TU$.
\end{defn}

\begin{prop} Jika $V$ merupakan ruang hasil kali dalam, maka ekuivalensi uniter memang merupakan
relasi ekuivalensi pada~$\ofml L(V)$.
\end{prop}

\begin{defn} Suatu operator $T$ pada ruang hasil kali dalam berdimensi hingga $V$ disebut
 \index{uniter!diagonalisasi}%
 \index{dapat didiagonalkan!secara uniter}%
\df{dapat didiagonalkan secara uniter} jika terdapat basis ortonormal untuk $V$ yang terhadapnya
$T$ bersifat diagonal. Secara ekuivalen, suatu operator pada ruang hasil kali dalam berdimensi hingga
\emph{dengan basis} dapat didiagonalkan secara uniter jika operator itu ekuivalen secara uniter dengan operator
diagonal.
\end{defn}

\begin{defn} Suatu operator $T$ pada ruang hasil kali dalam disebut
 \index{swaadjoin!operator}%
 \index{operator!swaadjoin}%
\df{swaadjoin} (atau
 \index{Hermitian!operator}%
 \index{operator!Hermitian}%
\df{Hermitian}) jika $T^* = T$.
\end{defn}

\begin{defn}\label{000164} Suatu proyeksi $P$ dalam ruang hasil kali dalam disebut
 \index{proyeksi!ortogonal}%
 \index{ortogonal!proyeksi!dalam ruang hasil kali dalam}%
 \index{operator!proyeksi ortogonal}%
\df{proyeksi ortogonal} jika proyeksi itu swaadjoin. Jika $M$ merupakan jangkauan suatu proyeksi
ortogonal, kita akan menggunakan notasi
 \index{PM@$\sbsb PM$ (proyeksi ortogonal ke~$M$)}%
$\sbsb PM$ untuk proyeksi tersebut, bukan notasi yang lebih rumit~$\sbsb E{M^\perp M}$.
\end{defn}

\begin{cau} Proyeksi pada ruang vektor atau ruang linear bernorma bersifat linear dan idempoten,
sedangkan proyeksi ortogonal pada ruang hasil kali dalam bersifat linear, idempoten, dan
swaadjoin. Situasi yang sebenarnya langsung ini agak dirumitkan oleh kecenderungan umum
untuk menyebut proyeksi ortogonal hanya sebagai ``proyeksi''. Bahkan, kelak dalam catatan
ini kita akan menggunakan konvensi tersebut. Dalam ruang hasil kali dalam,
$\oplus$ biasanya menyatakan jumlah langsung ortogonal dan ``proyeksi'' biasanya berarti
``proyeksi ortogonal''. Dalam banyak buku pengantar aljabar linear, seluruh pembahasan
berlangsung dalam $\R^n$, sehingga kadang sulit menentukan apakah pada suatu halaman
tertentu penulis memperlakukan $\R^n$ sebagai ruang vektor atau sebagai ruang hasil kali
dalam.
\end{cau}

\begin{prop} Jika $P$ merupakan proyeksi ortogonal pada ruang hasil kali dalam $V$, maka kita memperoleh
dekomposisi jumlah langsung ortogonal $V = \ker P \oplus \ran P$.
\end{prop}

\begin{defn} Jika $I_V = P_1 + P_2 + \dots + P_n$ merupakan resolusi identitas dalam suatu
ruang hasil kali dalam~$V$ dan setiap $P_k$ merupakan proyeksi ortogonal, kita mengatakan bahwa $I_V =
P_1 + P_2 + \dots + P_n$ merupakan
 \index{resolusi identitas!ortogonal}%
 \index{ortogonal!resolusi identitas}%
\df{resolusi ortogonal dari identitas}.
\end{defn}

\begin{prop} Jika $I_V = P_1 + P_2 + \dots + P_n$ merupakan resolusi ortogonal dari identitas
pada ruang hasil kali dalam $V$, maka $V = \bigoplus_{k=1}^n \ran P_k$.
\end{prop}

\begin{defn} Suatu operator $N$ pada ruang hasil kali dalam disebut
 \index{normal!operator}%
 \index{operator!normal}%
\df{normal} jika $NN^* = N^*N$.
\end{defn}

Dua pencapaian besar aljabar linear adalah \emph{teorema spektral} untuk operator pada
ruang hasil kali dalam (kompleks) berdimensi hingga (lihat \ref{00017}), yang memberikan
syarat perlu dan cukup yang dapat dinyatakan secara sederhana agar suatu operator dapat
didiagonalkan secara uniter, dan teorema~\ref{00018}, yang memberikan klasifikasi lengkap
operator-operator tersebut.

\begin{thm}[Teorema Spektral untuk Ruang Hasil Kali Dalam Kompleks Berdimensi Hingga]\label{00017}
Misalkan $T$ suatu operator pada ruang hasil kali dalam berdimensi hingga~$V$ dengan nilai-nilai eigen (berbeda)
 \index{spektral!teorema!operator normal pada ruang hasil kali dalam}%
$\lambda_1, \dots, \lambda_n$. Maka $T$ dapat didiagonalkan secara uniter jika dan
hanya jika $T$ normal. Jika $T$ normal, maka terdapat resolusi ortogonal dari identitas
$I_V = P_1 + \dots + P_n$, dengan jangkauan proyeksi ortogonal $P_k$ untuk setiap $k$
merupakan ruang eigen yang berkaitan dengan~$\lambda_k$, dan selanjutnya
     \[ T = \lambda_1P_1 +\dots + \lambda_nP_n\,. \]
\end{thm}

\begin{proof} Lihat \cite{Roman:2005}, halaman 227. \ns
\end{proof}

\begin{thm}\label{00018} Dua operator normal pada ruang hasil kali dalam berdimensi hingga
 \index{normal!operator!syarat untuk ekuivalensi uniter}%
ekuivalen secara uniter jika dan hanya jika keduanya memiliki nilai-nilai eigen yang sama,
masing-masing dengan multiplisitas yang sama; yaitu, jika dan hanya jika keduanya memiliki
polinom karakteristik yang sama.
\end{thm}

\begin{proof} Lihat \cite{HoffmanK:1971}, halaman 357. \ns
\end{proof}

Sebagian besar uraian selanjutnya dikhususkan untuk upaya yang sungguh tak sepele dalam mencari
generalisasi yang tepat dari kedua hasil di atas ke ranah berdimensi tak hingga, serta bentang
matematika menakjubkan yang mulai terlihat di sepanjang perjalanan tersebut.

\begin{exer} Misalkan $N = \dfrac13\begin{bmatrix}
                                4+2i & 1-i  & 1-i  \\
                                1-i  & 4+2i & 1-i  \\
                                1-i  & 1-i  & 4+2i \end{bmatrix}$.

\vskip 3 pt
 \begin{enumerate}
  \item Tunjukkan bahwa matriks $N$ bersifat normal.

\vskip 3 pt

  \item Dengan demikian, menurut \emph{teorema spektral}, $N$
dapat dituliskan sebagai kombinasi linear proyeksi-proyeksi ortogonal.
Tunjukkan secara eksplisit cara melakukannya (dan kerjakan perhitungannya).
 \end{enumerate}
\end{exer}

















\endinput

 \chapter{SELINGAN SANGAT SINGKAT TENTANG BAHASA KATEGORI}\label{categories}

Dalam matematika kita mempelajari berbagai hal (objek) dan pemetaan tertentu di antara hal-hal tersebut (morfisme).
Beberapa contohnya ialah himpunan dan fungsi, grup dan homomorfisme, ruang topologis
dan pemetaan kontinu, ruang vektor dan transformasi linear, serta ruang Hilbert dan
pemetaan linear terbatas. Contoh-contoh ini masing-masing berasal dari cabang matematika yang berbeda---teori
himpunan, teori grup, topologi, aljabar linear, dan analisis fungsional.
Namun, bidang-bidang yang berbeda ini memiliki banyak kesamaan: istilah seperti hasil kali,
koproduk, subobjek, hasil bagi, tarik balik, isomorfisme, dan limit proyektif muncul dalam banyak bidang.
Teori kategori merupakan upaya untuk menyatukan dan memformalkan sebagian konsep umum tersebut.
Dalam arti tertentu, teori kategori mempelajari hal-hal yang sama-sama dimiliki oleh berbagai cabang matematika.
Barangkali uraian yang lebih baik ialah: bahasa kategori memberi tahu kita ``bagaimana sesuatu bekerja'',
bukan ``apa sesuatu itu''.

Catatan ini---dan memang setiap buku pada tingkat ini---menggunakan bahasa himpunan di mana-mana tanpa
mensyaratkan kajian terperinci tentang teori himpunan aksiomatik. Demikian pula, kita akan dengan leluasa menggunakan
\emph{bahasa} kategori tanpa terlebih dahulu memulai kajian teori kategori yang memuaskan dari segi
fondasi (teori kategori sendiri merupakan bidang penelitian yang luas dan penting). Kadang-kadang buku teks membuat
bahkan bahasa kategori sulit dipelajari karena sangat mengandalkan latar belakang aljabar pembaca.
Sebagaimana kebanyakan orang merasa nyaman menggunakan bahasa himpunan (terlepas dari apakah mereka pernah mengkaji
\emph{teori} himpunan secara serius), hampir setiap orang seharusnya mendapati bahwa bahasa kategori nyaman digunakan
dan memberi pencerahan tanpa prasyarat yang rumit.

Bagi pembaca yang ingin mendalami pokok ini, saya merekomendasikan pengantar yang sangat ramah ke dunia
kategori, yang terdapat sebagai Bab 3 dalam buku indah karya Semadeni~\cite{Semadeni:1971}. Salah satu buku klasik
yang ditulis oleh seorang pendiri bidang ini ialah~\cite{MacLane:1971}. Buku yang lebih baru ialah~\cite{LawvereS:1997}.

Dengan menunjukkan prinsip-prinsip pemersatu, bahasa kategori kerap memberikan wawasan mencolok tentang
``cara sesuatu bekerja'' dalam matematika. Yang sama pentingnya, kita menjadi lebih efisien karena tidak perlu
mengulangi argumen yang pada hakikatnya sama dalam konteks yang hanya sedikit berbeda.

Untuk sementara, kita hanya akan mendefinisikan ``objek'', ``morfisme'', dan ``funktor'', lalu memberikan beberapa contoh.









\section{Objek dan Morfisme}\label{C0151}
\begin{defn} Misalkan $\sfml A$ suatu kelas yang anggota-anggotanya kita sebut
 \index{objek (dari suatu kategori)}%
\df{objek}. Untuk setiap pasangan objek $(S,T)$, kita kaitkan suatu kelas $\mor ST$ yang
anggota-anggotanya kita sebut
 \index{morfisme!dari suatu kategori}%
\df{morfisme} (atau
 \index{panah}%
\df{panah}) dari $S$ ke $T$. Kita mengasumsikan bahwa $\mor ST$ dan $\mor UV$ saling lepas,
kecuali jika $S = U$ dan $T = V$.

Kita mengandaikan pula adanya operasi $\circ$ (disebut \df{komposisi}) yang
mengaitkan setiap $\alpha \in \mor ST$ dan setiap $\beta \in \mor TU$ dengan suatu morfisme
$\beta \circ \alpha \in \mor SU$ sedemikian rupa sehingga:
 \begin{enumerate}
  \item $\gamma \circ (\beta \circ \alpha) = (\gamma \circ
\beta) \circ \alpha$ apabila $\alpha \in \mor ST$, $\beta \in \mor TU$, dan $\gamma \in
\mor UV$;
  \item untuk setiap objek $S$ terdapat morfisme $I_S \in
\mor SS$ yang memenuhi $\alpha \circ I_S = \alpha$ apabila $\alpha \in \mor ST$, dan $I_S
\circ \beta = \beta$ apabila $\beta \in \mor RS$.
 \end{enumerate}
Dalam keadaan ini, kelas $\sfml A$ bersama keluarga morfisme yang terkait merupakan suatu
 \index{kategori}%
\df{kategori}.

Kita akan mencadangkan notasi
 \index{<arrowto@$S \to^\alpha T$ (morfisme dalam suatu kategori)}%
$S \to^\alpha T$ untuk keadaan ketika $S$ dan $T$ merupakan objek dalam suatu kategori dan
$\alpha$ adalah morfisme anggota $\mor ST$. Seperti halnya pada grup dan ruang vektor,
kita biasanya menghilangkan simbol komposisi $\circ$ dan menulis
  \index{<binopconcat@$\beta\alpha$ (notasi untuk komposisi morfisme)}%
  \index{konvensi!notasi untuk komposisi!morfisme}%
$\beta\alpha$ untuk $\beta \circ \alpha$.

Suatu kategori disebut
 \index{kecil secara lokal}%
 \index{kategori!kecil secara lokal}%
\df{kecil secara lokal} jika $\mor ST$ merupakan himpunan untuk setiap pasangan objek $S$ dan $T$ dalam kategori tersebut. Kategori itu disebut
 \index{kecil}%
 \index{kategori!kecil}%
\df{kecil} jika, selain itu, kelas semua objeknya merupakan suatu himpunan. Dalam catatan ini, kategori-kategori yang kita minati
akan bersifat kecil secara lokal.
\end{defn}

\begin{exam} Kategori
 \index{kategori!$\cat{SET}$ sebagai suatu}%
 \index{SET@$\cat{SET}$!kategori}%
$\cat{SET}$ memiliki himpunan sebagai objek dan fungsi (pemetaan) sebagai morfisme.
\end{exam}

\begin{exam} Kategori
 \index{kategori!$\cat{AbGp}$ sebagai suatu}%
 \index{AbGp@$\cat{AbGp}$!kategori}%
$\cat{AbGp}$ memiliki grup Abelian sebagai objek dan homomorfisme grup sebagai morfisme.
\end{exam}

\begin{exam} Kategori
 \index{kategori!$\cat{VEC}$ sebagai suatu}%
 \index{VEC@$\cat{VEC}$!kategori}%
$\cat{VEC}$ memiliki ruang vektor sebagai objek dan transformasi linear sebagai morfisme.
\end{exam}







\begin{defn}\label{C008111} Misalkan $\le$ suatu relasi pada himpunan tak kosong $S$.
 \begin{enumerate}
  \item Jika relasi
 \index{<le@$\le$ (notasi untuk suatu urutan parsial)}%
 $\le$ bersifat refleksif (yakni, $x \le x$ untuk setiap $x \in S$) dan transitif (yakni, jika $x \le y$ dan $y \le z$, maka $x  \le z$), relasi itu disebut suatu
 \index{refleksif}%
 \index{transitif}%
 \index{praurutan}%
 \index{urutan!pra-}%
\df{praurutan}.
  \item Jika $\le$ merupakan praurutan dan juga bersifat antisimetris (yakni, jika $x \le y$ dan $y \le x$, maka $x = y$), relasi itu disebut suatu
 \index{antisimetris}%
 \index{parsial!urutan}%
 \index{urutan!parsial}%
\df{urutan parsial}.
  \item Elemen $x$ dan $y$ dalam suatu himpunan yang telah diberi praurutan disebut
 \index{dapat dibandingkan}%
\df{dapat dibandingkan} jika $x \le y$ atau $y \le x$.
  \item Jika $\le$ merupakan urutan parsial yang membuat setiap dua elemen dapat dibandingkan, relasi itu disebut suatu
 \index{linear!urutan}%
 \index{urutan!linear}%
\df{urutan linear} (atau suatu
 \index{total!urutan}%
 \index{urutan!total}%
\df{urutan total}).
  \item Jika relasi $\le$ merupakan praurutan (berturut-turut, urutan parsial,
urutan linear) pada $S$, pasangan $(S,\le)$ disebut suatu
 \index{himpunan terpraurut}%
\df{himpunan terpraurut} (berturut-turut,
 \index{himpunan terurut parsial}%
\df{himpunan terurut parsial},
 \index{himpunan terurut linear}%
\df{himpunan terurut linear}).
  \item Subhimpunan terurut linear dari himpunan terurut parsial $(S,\le)$ disebut suatu
 \index{rantai}%
\df{rantai} di~$S$.
 \end{enumerate}
Kita boleh menulis $b \ge a$ sebagai pengganti $a \le b$. Notasi
 \index{<l@$<$ (notasi untuk suatu urutan ketat)}%
$a < b$ (atau, secara ekuivalen, $b > a$) berarti $a \le b$ dan $a \ne b$.
\end{defn}

\begin{defn} Pemetaan $f\colon S \sto T$ antara himpunan-himpunan terpraurut disebut
 \index{urutan!pemetaan pelestari}%
\df{pelestari urutan} jika $x \le y$ di $S$ mengakibatkan $f(x) \le f(y)$ di~$T$.
\end{defn}

\begin{exam} Kategori
 \index{kategori!$\cat{POSET}$ sebagai suatu}%
 \index{POSET@$\cat{POSET}$!kategori}%
$\cat{POSET}$ memiliki himpunan terurut parsial sebagai objek dan pemetaan pelestari urutan sebagai morfisme.
\end{exam}

\begin{notn} Misalkan $f \colon S \sto T$ suatu fungsi antara himpunan. Untuk $A \subseteq S$, kita mendefinisikan
$f^{\sto}(A) = \{f(x)\colon x \in A\}$, dan untuk $B \subseteq T$, kita mendefinisikan
 \index{<superscript@$f^{\sto}(A)$ atau $f(A)$ (citra $A$ di bawah $f$)}%
 \index{<superscript@$f^{\gets}(B)$ (pracitra $B$ di bawah $f$)}%
$f^{\gets}(B) = \{ x \in S\colon f(x) \in B\}$. Kita mengatakan bahwa $f^{\sto}(A)$ adalah
 \index{citra}%
\df{citra $A$ di bawah $f$} dan bahwa $f^{\gets}(B)$ adalah
 \index{pracitra}%
\df{pracitra} (atau
 \index{invers!citra}%
 \index{citra!invers}%
\df{citra invers}) \df{$B$ di bawah~$f$}. Kadang-kadang, untuk menghindari notasi yang berjejal, saya akan menyerah pada kebiasaan meragukan dengan
menulis $f(A)$ sebagai pengganti~$f^{\sto}(A)$.
\end{notn}

\begin{defn} Ingat bahwa keluarga $\sfml T$ yang terdiri atas subhimpunan-subhimpunan suatu himpunan $X$ disebut suatu
 \index{topologi}%
\df{topologi} pada $X$ jika keluarga itu memuat himpunan kosong $\emptyset$ dan $X$ sendiri serta tertutup terhadap pembentukan gabungan
($\bigcup\sfml S \in \sfml T$ apabila $\sfml S \subseteq \sfml T$) dan irisan hingga ($\bigcap\sfml F \in \sfml T$
apabila $\sfml F \subseteq \sfml T$ dan $\sfml F$ berhingga). Himpunan tak kosong $X$ bersama suatu topologi pada $X$ disebut suatu
 \index{topologis!ruang}%
 \index{ruang!topologis}%
\df{ruang topologis}. Jika $(X,\sfml T)$ suatu ruang topologis, setiap anggota $\sfml T$ disebut suatu
 \index{terbuka}%
 \index{<binrel@$\open UX$ ($U$ merupakan subhimpunan terbuka dari~$X$)}%
\df{himpunan terbuka}. Untuk menyatakan bahwa himpunan $U$ merupakan subhimpunan terbuka dari $X$, kita sering menulis $\open UX$. Komplemen dari suatu himpunan terbuka disebut suatu
 \index{tertutup}%
\df{himpunan tertutup}. Fungsi $f\colon X \sto Y$ antara ruang-ruang topologis disebut
 \index{kontinu!fungsi}%
 \index{fungsi!kontinu}%
\df{kontinu} jika pracitra setiap himpunan terbuka bersifat terbuka; yakni, jika $\open{f^{\gets}(V)}X$ apabila $\open VY$.
(Jika Anda belum pernah menjumpai ruang topologis, lihat dua butir pertama dalam bagian pertama
Bab 10 dan bagian pertama Bab 11 pada catatan daring saya~\cite{Erdman:2007}.)
\end{defn}

\begin{exam} Setiap himpunan tak kosong $X$ memiliki topologi terkecil (terkasar); topologi itu terdiri tepat atas dua himpunan,
$\emptyset$ dan~$X$. Topologi ini adalah topologi
 \index{indiskret}%
 \index{topologi!indiskret}%
\df{indiskret} pada~$X$. (Perhatikan ejaan istilah bahasa Inggrisnya: terdapat kata lain dalam bahasa Inggris, ``indiscreet'',
yang bermakna sama sekali berbeda.) Himpunan $X$ juga memiliki topologi terbesar (terhalus), yang terdiri atas
keluarga $\sfml P(X)$ dari semua subhimpunan~$X$. Topologi ini adalah topologi
 \index{diskret}%
 \index{topologi!diskret}%
\df{diskret} pada~$X$.
  \end{exam}

\begin{exam} Kategori
 \index{kategori!$\cat{TOP}$ sebagai suatu}%
 \index{TOP@$\cat{TOP}$!kategori}%
$\cat{TOP}$ memiliki ruang topologis sebagai objek dan fungsi kontinu sebagai morfisme.
\end{exam}

\begin{defn}\label{defn_alg} Misalkan $(A,+,M)$ suatu ruang vektor atas $\K$ yang dilengkapi dengan
operasi biner lain $A \times A \sto A$, dengan $(x,y) \mapsto x \cdot y$, sedemikian
rupa sehingga $(A,+,\cdot\,)$ merupakan gelanggang. (Notasi $x \cdot y$ biasanya disingkat menjadi $xy$.) Jika, selain itu, persamaan
  \begin{equation}\label{eq_def_algebra}
     \alpha (xy) = (\alpha x)y = x(\alpha y)
  \end{equation}
berlaku untuk semua $x$, $y \in A$ dan $\alpha \in \K$, maka $(A,+,M,\,\cdot\,)$ disebut suatu
 \index{aljabar}%
\df{aljabar} atas medan~$\K$ (kadang-kadang disebut \df{aljabar asosiatif
linear}). Seperti biasa, kita menyalahgunakan terminologi dengan menulis hal-hal seperti, ``Misalkan $A$ suatu
aljabar.'' Kita mengatakan bahwa suatu aljabar $A$
 \index{beridentitas!aljabar}%
 \index{aljabar!beridentitas}%
\df{beridentitas} jika gelanggang yang mendasarinya $(A,+,\cdot\,)$ memiliki identitas perkalian; yakni, jika terdapat
 \index{perkalian!identitas}%
 \index{identitas!perkalian}%
suatu elemen $\vc 1_A \ne \vc 0$ di $A$ sehingga $\vc 1_A \cdot x = x \cdot \vc 1_A = x$ untuk semua $x \in A$.
  Aljabar itu disebut
 \index{komutatif!aljabar}%
 \index{aljabar!komutatif}%
\df{komutatif} jika gelanggangnya komutatif; yakni, jika $xy = yx$ untuk semua $x$, $y \in A$.

Subhimpunan $B$ dari suatu aljabar $A$ disebut suatu
 \index{subaljabar}%
\df{subaljabar} dari $A$ jika subhimpunan tersebut merupakan aljabar dengan operasi-operasi yang diwarisinya dari~$A$. Suatu
subaljabar $B$ dari aljabar beridentitas $A$ disebut
 \index{beridentitas!subaljabar}%
 \index{subaljabar!beridentitas}%
\df{subaljabar beridentitas} jika memuat identitas perkalian dari~$A$.

\textbf{PERHATIAN.} Agar menjadi subaljabar beridentitas, 
\emph{tidaklah} cukup jika $B$ memiliki identitas perkaliannya sendiri; subaljabar itu harus memuat identitas dari~$A$. Jadi, suatu
aljabar dapat sekaligus beridentitas dan merupakan subaljabar dari $A$ tanpa menjadi subaljabar beridentitas
dari~$A$.

Pemetaan $f\colon A \sto B$ antara aljabar-aljabar disebut suatu
 \index{homomorfisme!aljabar}%
 \index{aljabar!homomorfisme}%
\df{homomorfisme (aljabar)} jika pemetaan itu merupakan pemetaan linear antara $A$ dan $B$ sebagai ruang vektor yang mempertahankan perkalian
(yakni, $f(xy) = f(x)f(y)$ untuk semua $x$, $y \in A$). Dengan kata lain, homomorfisme aljabar adalah homomorfisme gelanggang linear.
Pemetaan itu disebut
 \index{beridentitas!homomorfisme aljabar}%
 \index{aljabar!homomorfisme!beridentitas}%
 \index{homomorfisme!aljabar!beridentitas}%
\df{homomorfisme (aljabar) beridentitas} jika mempertahankan identitas; yakni, jika $A$ dan $B$ keduanya merupakan aljabar beridentitas dan $f(\vc 1_A) = \vc 1_B$.
 \index{kernel!homomorfisme aljabar}%
\df{Kernel} dari suatu homomorfisme aljabar $f\colon A \sto B$ tentu saja adalah $\{a \in A \colon f(a) = \vc 0\}$.

Jika $f^{-1}$ ada dan juga merupakan homomorfisme aljabar, maka $f$ disebut suatu
 \index{isomorfisme!aljabar}%
\df{isomorfisme} dari $A$ ke $B$. Jika terdapat isomorfisme dari $A$ ke $B$, maka $A$ dan
$B$ disebut
 \index{isomorfik}%
\df{isomorfik}.
\end{defn}

\begin{exam} Kategori
 \index{kategori!$\cat{ALG}$ sebagai suatu}%
 \index{ALG@$\cat{ALG}$!kategori}%
$\cat{ALG}$ memiliki aljabar sebagai objek dan homomorfisme aljabar sebagai morfisme.
\end{exam}

Contoh-contoh sebelumnya merupakan contoh \emph{kategori konkret}---yakni, kategori
yang objek-objeknya berupa himpunan (biasanya bersama struktur tambahan) dan
morfisme-morfismenya berupa fungsi (biasanya dengan mempertahankan struktur tambahan tersebut). Dalam catatan ini,
kategori-kategori yang kita minati bersifat konkret sekaligus kecil secara lokal. Berikut ini (bagi pembaca yang ingin tahu) adalah
definisi yang lebih formal untuk \emph{kategori konkret}.

\begin{defn}\label{C015124} Suatu kategori $\sfml A$ bersama fungsi $\abs{\hphantom{O}}$ yang mengaitkan
setiap objek $A$ di $\sfml A$ dengan suatu himpunan $\abs A$ disebut suatu
 \index{kategori konkret}%
 \index{kategori!konkret}%
\df{kategori konkret} jika syarat-syarat berikut dipenuhi:
 \begin{enumerate}
    \item setiap morfisme $A \to^f B$ merupakan fungsi dari $\abs A$ ke~$\abs B$;
    \item setiap morfisme identitas $I_A$ di $\sfml A$ merupakan fungsi identitas pada~$\abs A$; dan
    \item komposisi morfisme $A \to^f B$ dan $B \to^g C$ sama dengan komposisi
fungsi $f\colon \abs A \sto \abs B$ dan $g\colon \abs B \sto \abs C$.
 \end{enumerate}
Jika $A$ suatu objek dalam kategori konkret $\sfml A$, maka $\abs A$ disebut
 \index{himpunan dasar}%
 \index{himpunan!dasar}%
 \index{<unop@$\abs A$ (himpunan dasar dari objek~$A$)}%
\df{himpunan dasar} dari~$A$.
\end{defn}

Meskipun kategori-kategori yang kita minati dalam catatan ini memang merupakan kategori
konkret, mungkin tetap menarik untuk melihat contoh kategori yang
\emph{tidak disajikan secara konkret}, yakni tanpa terlebih dahulu mewujudkan objek sebagai himpunan dan morfisme sebagai fungsi.

\begin{exam}\label{C015127} Misalkan $G$ suatu monoid (yakni, semigrup beridentitas).
 \index{monoid}%
Tinjau suatu kategori $\cat C$ yang memiliki tepat satu objek, yang kita sebut~$\star$. Karena hanya
ada satu objek, hanya ada satu keluarga morfisme $\mor{\star}{\star}$, yang kita
ambil sebagai~$G$. Komposisi morfisme didefinisikan sebagai perkalian monoid.
Artinya, $a \circ b := ab$ untuk semua $a$, $b \in G$. Jelas bahwa komposisi bersifat asosiatif dan
elemen identitas $G$ merupakan morfisme identitas. Jadi, $\cat C$ adalah suatu kategori.
\end{exam}

\begin{defn} Dalam setiap kategori konkret, kita akan menyebut morfisme injektif sebagai suatu
 \index{monomorfisme}%
\df{monomorfisme} dan morfisme surjektif sebagai suatu
 \index{epimorfisme}%
\df{epimorfisme}.
\end{defn}

\begin{cau} Definisi-definisi di atas mencerminkan penggunaan asli istilah tersebut oleh Bourbaki dan merupakan definisi yang
paling lazim dipakai oleh matematikawan di luar teori kategori itu sendiri; di dalam teori kategori,
``monomorfisme'' berarti ``dapat dibatalkan dari kiri'' dan ``epimorfisme'' berarti ``dapat
dibatalkan dari kanan''. (Perhatikan bahwa istilah \emph{injektif} dan \emph{surjektif} mungkin tidak
bermakna jika diterapkan pada morfisme dalam kategori yang tidak konkret.)

Suatu morfisme $B \to^g C$ disebut
 \index{kiri!dapat dibatalkan dari}%
 \index{dapat dibatalkan!dari kiri}%
\df{dapat dibatalkan dari kiri} jika, setiap kali morfisme $A \to^{f_1} B$ dan $A \to^{f_2} B$ memenuhi
$gf_1 = gf_2$, berlaku $f_1 = f_2$. Mac Lane menyarankan agar morfisme yang dapat dibatalkan dari kiri disebut
 \index{monik}%
morfisme \df{monik}. Perbedaan antara morfisme monik dan monomorfisme ternyata
kecil. Dalam catatan ini, hampir semua morfisme yang kita jumpai bersifat monik jika dan
hanya jika merupakan monomorfisme. Sebagai latihan mudah, buktikan bahwa setiap morfisme injektif
dalam suatu kategori (konkret) bersifat monik. Konversnya kadang-kadang gagal.

Dengan semangat yang sama, Mac Lane menyarankan agar suatu morfisme yang
 \index{kanan!dapat dibatalkan dari}%
 \index{dapat dibatalkan!dari kanan}%
\emph{dapat dibatalkan dari kanan} (yakni, suatu morfisme $A \to^f B$ sedemikian sehingga, apabila
morfisme $B \to^{g_1} C$ dan $B \to^{g_2} C$ memenuhi $g_1f = g_2f$, maka $g_1 = g_2$) disebut suatu
 \index{epik}%
morfisme \df{epik}. Sekali lagi, merupakan latihan mudah untuk menunjukkan bahwa dalam suatu kategori (konkret),
setiap epimorfisme bersifat epik. Namun, konversnya gagal dalam beberapa kategori yang cukup umum.
\end{cau}

\begin{defn} Terminologi untuk invers morfisme dalam kategori pada dasarnya sama
dengan terminologi untuk fungsi. Misalkan $A \to^\alpha B$ dan $B \to^\beta A$ merupakan
morfisme dalam suatu kategori. Jika $\beta \circ \alpha = I_A$, maka $\beta$ adalah suatu
 \index{kiri!invers!morfisme}%
 \index{invers!kiri}%
\df{invers kiri} dari $\alpha$ dan, secara ekuivalen, $\alpha$ adalah suatu
 \index{kanan!invers!morfisme}%
 \index{invers!kanan}%
\df{invers kanan} dari $\beta$. Kita mengatakan bahwa morfisme $\alpha$ merupakan suatu
 \index{isomorfisme!dalam suatu kategori}%
\df{isomorfisme} (atau bersifat
 \index{invertibel!morfisme}%
\df{invertibel}) jika terdapat morfisme $B \to^\beta A$ yang sekaligus merupakan invers kiri dan
invers kanan dari~$\alpha$. Morfisme demikian dinotasikan dengan $\alpha^{-1}$ dan disebut
 \index{<unopur@$\alpha^{-1}$ (invers morfisme $\alpha$)}%
 \index{invers!morfisme}%
\df{invers} dari~$\alpha$.
\end{defn}

\begin{prop} Misalkan $A \to^\alpha B$ suatu morfisme dalam suatu kategori. Jika $\alpha$ memiliki invers kiri $\lambda$ dan
invers kanan $\rho$, maka $\lambda = \rho$ dan, akibatnya, $\alpha$ merupakan suatu isomorfisme.
\end{prop}

Dalam setiap kategori konkret, kita dapat menanyakan apakah setiap morfisme bijektif (yakni, setiap
pemetaan yang sekaligus merupakan monomorfisme dan epimorfisme) adalah suatu isomorfisme. Jawabannya
sering kali merupakan \emph{ya} yang mudah (seperti dalam $\cat{SET}$, $\cat{AbGp}$, dan $\cat{VEC}$) atau
\emph{tidak} yang mudah (misalnya, dalam kategori $\cat{POSET}$
 \index{morfisme bijektif!tidak harus invertibel!dalam $\cat{POSET}$}%
 \index{POSET@$\cat{POSET}$!morfisme bijektif dalam}%
 \index{POSET@$\cat{POSET}$!kategori}%
yang terdiri atas himpunan terurut parsial dan pemetaan pelestari urutan). Namun,
sesekali jawabannya ternyata berupa hasil yang menarik dan mendalam (lihat, misalnya, \emph{teorema pemetaan
terbuka}~\futurexref{6.3.4}{C069414}).

\begin{exam} Dalam kategori $\cat{SET}$,
 \index{morfisme bijektif!invertibel!dalam $\cat{SET}$}%
 \index{SET@$\cat{SET}$!morfisme bijektif dalam}%
setiap morfisme bijektif merupakan suatu isomorfisme.
\end{exam}

\begin{exam}\label{C015216} Kategori
 \index{LAT@$\cat{LAT}$!kategori}%
 \index{kategori!$\cat{LAT}$ sebagai suatu}%
$\cat{LAT}$ memiliki kekisi sebagai objek dan homomorfisme
kekisi sebagai morfisme. (Suatu
 \index{kekisi}%
\df{kekisi} adalah himpunan terurut parsial yang setiap pasangan elemennya memiliki infimum dan supremum, dan suatu
 \index{kekisi!homomorfisme}%
 \index{homomorfisme!kekisi}%
\df{homomorfisme kekisi} adalah pemetaan antara kekisi-kekisi yang mempertahankan infimum dan supremum.)
Dalam kategori ini, morfisme bijektif
 \index{morfisme bijektif!invertibel!dalam $\cat{LAT}$}%
 \index{LAT@$\cat{LAT}$!morfisme bijektif dalam}%
merupakan isomorfisme.
\end{exam}

\begin{exam} Suatu
 \index{homeomorfisme}%
\df{homeomorfisme} adalah bijeksi kontinu $f$ antara ruang-ruang topologis yang inversnya $f^{-1}$ juga kontinu. Homeomorfisme
merupakan isomorfisme dalam kategori $\cat{TOP}$. Dalam kategori ini, tidak setiap bijeksi kontinu merupakan homeomorfisme.
\end{exam}

\begin{exam} Jika, dalam kategori $\cat C$ pada Contoh~\ref{C015127}, monoid $G$ merupakan suatu grup,
maka setiap morfisme dalam $\cat C$ merupakan suatu isomorfisme.
\end{exam}






\section{Funktor}
\begin{defn} Jika $\cat{A}$ dan $\cat{B}$ merupakan kategori, suatu
 \index{kovarian!funktor}%
 \index{funktor!kovarian}%
\df{funktor kovarian} $F$ dari $\cat A$ ke $\cat B$ (ditulis ${\cat A}\to^F{\cat B}$) adalah
sepasang pemetaan: suatu
 \index{objek!pemetaan}%
 \index{pemetaan!objek}%
\df{pemetaan objek} $F$ yang mengaitkan setiap objek $S$ dalam $\cat A$ dengan suatu objek $F(S)$ dalam
$\cat B$, dan suatu
 \index{morfisme!pemetaan}%
 \index{pemetaan!morfisme}%
\df{pemetaan morfisme} (yang juga dinotasikan dengan $F$) yang mengaitkan setiap morfisme $f \in \mor
ST$ dalam $\cat A$ dengan suatu morfisme $F(f) \in \mor {F(S)}{F(T)}$ dalam $\cat B$, sedemikian rupa sehingga
 \begin{enumerate}
  \item $F(g \circ f) = F(g) \circ F(f)$ apabila $g \circ f$ terdefinisi dalam $\cat A$; dan
  \item $F(\id S) = \id{F(S)}$ untuk setiap objek $S$ dalam $\cat A$.
 \end{enumerate}

Definisi suatu
 \index{kontravarian!funktor}%
 \index{funktor!kontravarian}%
\df{funktor kontravarian} ${\cat A}\to^F{\cat B}$ berbeda dari definisi sebelumnya
hanya dalam dua hal. Pertama, pemetaan morfismenya mengaitkan setiap morfisme $f \in \mor ST$ dalam
$\cat A$ dengan suatu morfisme $F(f) \in \mor {F(T)}{F(S)}$ dalam $\cat B$. Kedua, syarat (a)
di atas diganti dengan
 \begin{enumerate}
  \item[(\rm{a'})] $F(g \circ f) = F(f) \circ F(g)$ apabila $g \circ f$
terdefinisi dalam~$\cat A$.
 \end{enumerate}
\end{defn}

\begin{exam}\label{functors005exam} Suatu
 \index{pelupa!funktor}%
 \index{funktor!pelupa}%
\df{funktor pelupa} adalah funktor yang memetakan objek dan morfisme dari suatu kategori $\cat
C$ ke kategori $\cat C'$ yang memiliki struktur atau sifat lebih sedikit. Sebagai contoh, jika $V$
suatu ruang vektor, funktor $F$ yang ``melupakan'' operasi perkalian skalar
pada ruang vektor akan memetakan $V$ ke kategori grup Abelian. (Grup
Abelian $F(V)$ akan memiliki himpunan elemen yang sama dengan ruang vektor $V$ dan
operasi penjumlahan yang sama, tetapi tidak memiliki perkalian skalar.) Suatu pemetaan linear
$T\colon V \sto W$ antara ruang-ruang vektor akan dibawa oleh funktor $F$ menjadi homomorfisme
grup $F(T)$ antara grup-grup Abelian $F(V)$ dan~$F(W)$.

Funktor pelupa juga dapat ``melupakan'' sifat. Jika $G$ suatu objek dalam
kategori grup Abelian, funktor yang ``melupakan'' komutativitas grup Abelian
akan membawa $G$ ke kategori grup.

Pada bagian sebelumnya telah disebutkan bahwa semua kategori yang kita minati dalam
catatan ini merupakan kategori konkret (kategori yang objek-objeknya berupa himpunan dengan struktur
tambahan dan morfisme-morfismenya berupa pemetaan yang, dalam pengertian tertentu, mempertahankan
struktur tambahan tersebut). Kita akan beberapa kali menggunakan jenis khusus funktor
pelupa---funktor yang melupakan seluruh struktur objek kecuali himpunan dasarnya dan yang
melupakan segala sifat morfisme yang mempertahankan struktur. Jika $A$ suatu
objek dalam suatu kategori konkret $\cat C$, kita menotasikan
 \index{<unop@$\abs A$ (funktor pelupa yang bekerja pada~$A$)}%
himpunan dasarnya dengan $\abs{A}$. Jika \mbox{$A~\to^f~B$} suatu morfisme dalam $\cat C$, kita
menotasikan dengan $\abs f$ pemetaan dari $\abs A$ ke $\abs B$ yang dipandang sekadar sebagai fungsi
antara himpunan. Mudah dilihat bahwa $\abs{\hphantom{M}}$\,, yang membawa objek dalam
$\cat C$ ke objek dalam $\cat{SET}$ (kategori himpunan dan pemetaan) serta morfisme dalam $\cat
C$ ke morfisme dalam $\cat{SET}$, merupakan suatu funktor kovarian.

Dalam kategori $\cat{VEC}$ yang terdiri atas ruang vektor dan pemetaan linear, misalnya,
$\abs{\hphantom{M}}$ membuat ruang vektor $V$ ``melupakan'' penjumlahan maupun
perkalian skalarnya ($\abs V$ hanyalah suatu himpunan). Jika $T\colon V \sto W$ suatu transformasi
linear, maka $\abs T\colon \abs V \sto \abs W$ hanyalah pemetaan antara himpunan---pemetaan itu telah
``melupakan'' sifat mempertahankan operasi.
\end{exam}


\begin{defn} Suatu himpunan terurut parsial disebut
 \index{urutan!lengkap}%
 \index{lengkap!urutan}%
\df{lengkap secara urutan} jika setiap subhimpunan tak kosong memiliki supremum (yakni, batas atas
terkecil) dan infimum (batas bawah terbesar).
\end{defn}

\begin{defn} Misalkan $S$ suatu himpunan. Maka
 \index{himpunan kuasa}%
 \index{P@$\sfml P(S)$ (himpunan kuasa dari $S$)}%
\df{himpunan kuasa} dari $S$, yang dinotasikan dengan $\sfml P(S)$, adalah keluarga semua subhimpunan~$S$.
\end{defn}

\begin{exam}[Funktor himpunan kuasa] Misalkan $S$ suatu himpunan tak kosong.
 \begin{enumerate}
  \item Himpunan kuasa $\sfml P(S)$ dari $S$ yang diurutkan secara parsial oleh $\subseteq$
bersifat lengkap secara urutan.
  \item Kelas himpunan terurut parsial yang lengkap secara urutan beserta pemetaan
pelestari urutan merupakan suatu kategori.
 \index{funktor!himpunan kuasa}%
 \index{himpunan kuasa!funktor}%
  \item Untuk setiap fungsi $f$ antara himpunan, misalkan $\sfml P(f) = f^{\sto}$. Maka
$\sfml P$ merupakan funktor kovarian dari kategori himpunan dan fungsi ke kategori
himpunan terurut parsial yang lengkap secara urutan beserta pemetaan pelestari urutan.
  \item Untuk setiap fungsi $f$ antara himpunan, misalkan $\sfml P(f) = f^{\gets}$. Maka
$\sfml P$ merupakan funktor kontravarian dari kategori himpunan dan fungsi ke kategori
himpunan terurut parsial yang lengkap secara urutan beserta pemetaan pelestari urutan.
 \end{enumerate}
\end{exam}

\begin{defn}\label{defn_vsp_adjoint} Misalkan $T\colon V \sto W$ suatu pemetaan linear antara ruang-ruang vektor. Untuk setiap
$g \in W^{\#}$, misalkan $T^{\#}(g) = g\,T$. Perhatikan bahwa $T^{\#}(g) \in V^{\#}$. Pemetaan $T^{\#}$ dari ruang
vektor $W^{\#}$ ke ruang vektor~$V^{\#}$ disebut pemetaan
 \index{adjoin!ruang vektor}%
 \index{vektor!ruang!pemetaan adjoin}%
 \index{<unopur@$T^\#$ (adjoin ruang vektor)}%
\df{adjoin} (ruang vektor) dari~$T$.
\end{defn}

\begin{exam} Pasangan pemetaan $V \mapsto V^{\#}$ (yang membawa suatu ruang vektor ke dual aljabarnya) dan
$T \mapsto T^{\#}$ (yang membawa suatu pemetaan linear ke dualnya) merupakan funktor kontravarian dari kategori $\cat{VEC}$
yang terdiri atas ruang vektor dan pemetaan linear ke kategori itu sendiri.
\end{exam}

\begin{exam}\label{0007606} Jika $\cat C$ suatu kategori, definisikan
 \index{kategori@$\cat{C^2}$ (pasangan objek dan morfisme)}%
$\cat{C^2}$ sebagai kategori yang objek-objeknya merupakan pasangan terurut objek dalam $\cat C$ dan
morfisme-morfismenya merupakan pasangan terurut morfisme dalam~$\cat C$. Jadi, jika $A \to^f C$ dan $B
\to^g D$ merupakan morfisme dalam $\cat C$, maka $(A,B) \to^{(f,g)} (C,D)$ (dengan $(f,g)(a,b) =
\bigl(f(a),g(b)\bigr)$ untuk semua $a \in A$ dan $b \in B$ apabila kategori tersebut konkret) merupakan morfisme dalam~$\cat{C^2}$.
Komposisi morfisme didefinisikan dengan cara yang wajar: $(f,g) \circ (h,j) = (f \circ h, g
\circ j)$. Kita mendefinisikan
 \index{diagonal!funktor}%
 \index{funktor!diagonal}%
\df{funktor diagonal} $\cat C \to^{\ftr D} \cat{C^2}$ dengan $\ftr D(A) := (A,A)$ dan
$\ftr D(f) := (f,f)$. Ini merupakan
funktor kovarian.
\end{exam}

\begin{exam}\label{C029717} Misalkan $X$ dan $Y$ ruang-ruang topologis dan $\phi\colon X \sto Y$
kontinu. Definisikan $\fml C \phi$ pada $\fml C(Y)$ dengan
    \[ \fml C \phi(g) = g \circ \phi \]
untuk semua $g \in \fml C(Y)$. Maka pasangan pemetaan $X \mapsto \fml C(X)$ dan
 \index{C@$\fml C$ (funktor)}%
$\phi \mapsto \fml C \phi$ merupakan funktor kontravarian dari kategori $\cat{TOP}$ yang terdiri atas ruang
topologis dan pemetaan kontinu ke kategori aljabar beridentitas dan homomorfisme
aljabar beridentitas.
\end{exam}











\endinput

 \chapter{RUANG LINEAR BERNORMA}

\section{Norma}
Dalam dunia analisis, penghuni yang paling dominan adalah ruang fungsi, yakni ruang vektor
yang terdiri atas fungsi-fungsi bernilai real atau kompleks. Agar menarik bagi seorang analis,
ruang semacam itu harus dilengkapi dengan suatu topologi. Sering kali topologi tersebut
merupakan topologi metrik, yang pada gilirannya kerap berasal dari suatu norma (lihat
proposisi~\ref{0001503}).

\begin{exam}\label{C023125} Untuk $x = (x_1,\dots, x_n) \in \K^n$, misalkan $\norm{x} =
\left(\sum_{k=1}^n {\abs{x_k}}^2\right)^{1/2}$. Ini adalah
 \index{biasa!norma!pada $\K^n$}%
 \index{norma!biasa!pada $\K^n$}%
\df{norma biasa} (atau
 \index{Euklides!norma!pada $\K^n$}%
 \index{norma!Euklides!pada $\K^n$}%
 \index{K@$\K^n$!sebagai ruang linear bernorma}%
 \index{bernorma!ruang linear!$\K^n$ sebagai}%
\df{norma Euklides}) pada $\K^n$; kecuali jika secara eksplisit dinyatakan sebaliknya,
ketika dipandang sebagai ruang linear bernorma, $\K^n$ selalu dianggap memiliki norma ini.
Jelas bahwa norma ini menginduksi metrik biasa (Euklides) pada~$\K^n$.
\end{exam}

\begin{exam}\label{C023127} Untuk $x = (x_1,\dots,x_n) \in \K^n$, misalkan
 \index{<unopn@$\norm{\hphantom{x}}_1$ (norma-$1$)}%
$\norm{x}_1 = \sum_{k=1}^n \abs{x_k}$. Mudah dilihat bahwa fungsi $x \mapsto \norm{x}_1$
merupakan norma pada~$\K^n$. Norma ini kadang-kadang disebut
 \index{norma!$1$- (pada $\K^n$)}%
 \index{satu@norma-$1$!pada $\K^n$}%
 \index{K@$\K^n$!sebagai ruang linear bernorma}%
 \index{bernorma!ruang linear!$\K^n$ sebagai}%
\df{norma-$1$} pada~$\K^n$. Norma ini menginduksi apa yang disebut \emph{metrik taksi}
pada~$\R^2$.
\end{exam}

Dari enam contoh berikutnya, contoh terakhir adalah yang paling umum. Perhatikan bahwa
semua contoh lainnya merupakan kasus khususnya atau subruang dari kasus-kasus khususnya.

\begin{exam}\label{C023129} Untuk $x = (x_1,\dots,x_n) \in \K^n$, misalkan
 \index{<unopn@$\norm{\hphantom{x}}_u$ (norma seragam)}%
 \index{<unopn@$\norm{\hphantom{x}}_\infty$ (norma seragam)}%
$\norm{x}_\infty = \max\{\,\abs{x_k} \colon 1 \le k \le n\}$. Ini mendefinisikan suatu
norma pada $\K^n$; norma tersebut adalah
 \index{seragam!norma!pada~$\K^n$}%
 \index{norma!seragam!pada $\K^n$}%
 \index{K@$\K^n$!sebagai ruang linear bernorma}%
 \index{bernorma!ruang linear!$\K^n$ sebagai}%
\df{norma seragam} pada~$\K^n$. Norma ini menginduksi metrik seragam pada~$\K^n$.
Notasi alternatif untuk $\norm{x}_\infty$ adalah $\norm{x}_U$.
\end{exam}

\begin{exam}\label{C0231301} Himpunan $l_\infty$ dari semua barisan bilangan kompleks yang
terbatas jelas merupakan ruang vektor di bawah operasi penjumlahan dan perkalian skalar
titik demi titik. Untuk $x = (x_1,x_2,\dots) \in l_\infty$, misalkan
 \index{<unopn@$\norm{\hphantom{x}}_\infty$ (norma seragam)}%
$\norm{x}_\infty = \sup\{\,\abs{x_k} \colon 1 \le k\}$. Ini mendefinisikan suatu norma pada
$l_\infty$; norma tersebut adalah
 \index{seragam!norma!pada~$l_\infty$}%
 \index{norma!seragam!pada $l_\infty$}%
 \index{l@$l_\infty$!sebagai ruang linear bernorma}%
 \index{l@$l_\infty$!barisan terbatas}%
 \index{bernorma!ruang linear!$l_\infty$ sebagai}%
\df{norma seragam} pada~$l_\infty$.
\end{exam}

\begin{exam}\label{C0231302} Salah satu subruang dari~\ref{C0231301} adalah ruang $c$ yang
terdiri atas semua
 \index{c@$c$!sebagai ruang linear bernorma}%
 \index{c@$c$!barisan konvergen}%
 \index{bernorma!ruang linear!$c$ sebagai}%
barisan bilangan kompleks yang konvergen. Di bawah norma seragam, ruang ini merupakan
ruang linear bernorma.
\end{exam}

\begin{exam}\label{C0231303} Di dalam $c$ terdapat ruang linear bernorma lain yang menarik,
yaitu $\sbsb c0$
 \index{c@$c_0$!barisan yang konvergen ke nol}%
 \index{c@$c_0$!sebagai ruang linear bernorma}%
 \index{bernorma!ruang linear!$\sbsb c0$ sebagai}%
yakni ruang semua barisan bilangan kompleks yang konvergen ke nol (sekali lagi dengan
norma seragam).
\end{exam}

\begin{exam}\label{C023131} Misalkan $S$ suatu himpunan tak kosong. Jika $f$ merupakan fungsi
bernilai $\K$ yang terbatas pada $S$, misalkan
  \[ \norm f_u := \sup\{\abs{f(x)} \colon x \in S\}\,. \]
Ini adalah norma pada ruang vektor $\fml B(S)$ dari semua fungsi skalar yang
terbatas pada $S$ dan disebut
 \index{<unopn@$\norm{\hphantom{x}}_u$ (norma seragam)}%
 \index{seragam!norma!pada~$\fml B(S)$}%
 \index{norma!seragam!pada~$\fml B(S)$}%
 \index{B@$\fml B(S)$!fungsi terbatas pada~$S$}%
 \index{B@$\fml B(S)$!sebagai ruang linear bernorma}%
 \index{bernorma!ruang linear!$\fml B(S)$ sebagai}%
\df{norma seragam}. Jelas bahwa norma ini menghasilkan metrik seragam pada~$\fml B(S)$.
Perhatikan bahwa contoh~\ref{C023129} dan~\ref{C0231301} merupakan kasus khusus dari
contoh ini. (Ambil $S = \N_n := \{1,2,\dots,n\}$ atau $S = \N$.)
\end{exam}











\begin{exam}\label{C023133} Mudah untuk membuat perumuman yang cukup luas atas contoh
sebelumnya dengan mengganti medan $\K$ dengan ruang linear bernorma sebarang. Maka,
misalkan $S$ suatu himpunan tak kosong dan $V$ suatu ruang linear bernorma. Untuk $f$
dalam ruang vektor $\fml B(S,V)$ dari semua fungsi bernilai $V$ yang terbatas pada~$S$,
misalkan
  \[ \norm f_u := \sup\{\norm{f(x)} \colon x \in S\}\,. \]
Maka ini adalah suatu norma dan disebut
 \index{<unopn@$\norm{\hphantom{x}}_u$ (norma seragam)}%
 \index{seragam!norma!pada~$\fml B(S,V)$}%
 \index{norma!seragam!pada~$\fml B(S,V)$}%
 \index{B@$\fml B(S,V)$!fungsi bernilai $V$ yang terbatas pada~$S$}%
 \index{B@$\fml B(S,V)$!sebagai ruang linear bernorma}%
 \index{bernorma!ruang linear!$\fml B(S,V)$ sebagai}%
\df{norma seragam} pada $\fml B(S,V)$. (Mengapa kata ``mudah'' dalam kalimat pertama
contoh ini tepat digunakan?)
\end{exam}

\begin{defn} Misalkan $S$ suatu himpunan, $V$ suatu ruang linear bernorma, dan
 \index{F@$\fml F(S,V)$ (fungsi bernilai $V$ pada~$S$)}%
$\fml F(S,V)$ ruang vektor dari semua fungsi bernilai $V$ pada~$S$. Perhatikan suatu
barisan $(f_n)$ dari fungsi-fungsi dalam $\fml F(S,V)$. Jika terdapat fungsi $g$ dalam
$\fml F(S,V)$ sedemikian sehingga
   \[ \sup\{\norm{f_n(x) - g(x)}\colon x \in S\} \sto 0
      \quad \text{ketika $n \sto \infty$,} \]
maka kita mengatakan bahwa barisan $(f_n)$
 \index{<arrowto@$f_n \sto g$ (unif) (konvergensi seragam)}%
 \index{seragam!konvergensi}%
 \index{konvergensi!seragam}%
\df{konvergen seragam} ke~$g$ dan menulis $f_n \sto g \text{\,(unif)}$. Fungsi $g$
adalah
 \index{seragam!limit}%
 \index{limit!seragam}%
\df{limit seragam} dari barisan $(f_n)$. Perhatikan bahwa jika $g$ dan semua $f_n$
termasuk dalam $\fml B(S,V)$, maka konvergensi seragam $(f_n)$ ke $g$ tidak lain adalah
konvergensi $(f_n)$ ke $g$ terhadap metrik seragam.
\end{defn}

Ada banyak cara barisan fungsi dapat konvergen. Boleh dikatakan bahwa dua ragam
konvergensi yang paling umum adalah konvergensi seragam, yang baru saja kita bahas, dan
konvergensi titik demi titik.

\begin{defn} Misalkan $S$ suatu himpunan, $V$ suatu ruang linear bernorma, dan $(f_n)$
suatu barisan dalam $\fml F(S,V)$. Jika terdapat fungsi $g$ sedemikian sehingga
  \[ f_n(x) \sto g(x) \qquad \text{untuk semua $x \in S$}, \]
maka $(f_n)$
 \index{titik demi titik!konvergensi}%
 \index{konvergensi!titik demi titik}%
\df{konvergen titik demi titik} ke~$g$. Dalam hal ini kita menulis
  \[ f_n \sto g \text{\,(ptws)}. \]
Fungsi $g$ adalah
 \index{titik demi titik!limit}%
 \index{limit!titik demi titik}%
 \index{<arrowto@$f_n \sto g$ (ptws) (konvergensi titik demi titik)}%
\df{limit titik demi titik} dari fungsi-fungsi~$f_n$.

Jika $(f_n)$ merupakan barisan naik dari fungsi-fungsi bernilai real (atau real
diperluas) dan $f_n \sto g \text{\,(ptws)}$, kita menulis
 \index{<arrowmonoup@$f_n \uparrow g$ (ptws) (konvergensi monoton titik demi titik)}%
$f_n \uparrow g \text{\,(ptws)}$. Dan jika $(f_n)$ menurun serta memiliki $g$ sebagai
limit titik demi titik, kita menulis
 \index{<arrowmonodown@$f_n \downarrow g$ (ptws) (konvergensi monoton titik demi titik)}%
$f_n \downarrow g \text{\,(ptws)}$.
\end{defn}

Hubungan paling jelas antara kedua jenis konvergensi ini, yang dikenal dari setiap mata
kuliah analisis real, adalah bahwa konvergensi seragam mengakibatkan konvergensi titik demi
titik.

\begin{prop}\label{C023141} Misalkan $S$ suatu himpunan dan $V$ suatu ruang linear bernorma.
Jika barisan $(f_n)$ dalam $\fml F(S,V)$ konvergen seragam ke suatu fungsi $g$ dalam
$\fml F(S,V)$, maka $(f_n)$ konvergen titik demi titik ke~$g$.
\end{prop}

Kebalikannya tidak benar.

\begin{exam} Untuk setiap $n \in \N$, misalkan
$f_n\colon [0,1] \sto \R\colon x \mapsto x^n$. Maka barisan $(f_n)$ konvergen titik demi
titik pada $[0,1]$, tetapi tidak konvergen seragam.
\end{exam}

\begin{exam} Untuk setiap $n \in \N$ dan setiap $x \in \R^+$, misalkan
$f_n(x) = \frac1n\, x$. Maka pada setiap interval berbentuk $[0,a]$ dengan $a > 0$,
barisan $(f_n)$ konvergen seragam ke fungsi konstan~$0$. Pada interval $[0,\infty)$,
barisan itu konvergen titik demi titik ke $0$, tetapi tidak konvergen seragam.
\end{exam}

Ingat pula dari analisis real bahwa limit seragam dari fungsi-fungsi terbatas juga harus
terbatas.

\begin{prop}\label{C023151} Misalkan $S$ suatu himpunan, $V$ suatu ruang linear bernorma,
$(f_n)$ suatu barisan dalam $\fml B(S,V)$, dan $g$ suatu anggota $\fml F(S,V)$. Jika
$f_n \sto g\text{\,(unif)}$, maka $g$ terbatas.
\end{prop}

Dan limit seragam dari fungsi-fungsi kontinu bersifat kontinu.

\begin{prop}\label{C023154} Misalkan $X$ suatu ruang topologis, $V$ suatu ruang linear
bernorma, $(f_n)$ suatu barisan fungsi kontinu bernilai $V$ pada~$X$, dan $g$ suatu anggota
$\fml F(X,V)$. Jika $f_n \sto g\text{\,(unif)}$, maka $g$ kontinu.
\end{prop}

\begin{prop}\label{C023163} Jika $x$ dan $y$ merupakan unsur dari suatu ruang vektor $V$
yang dilengkapi dengan norma $\norm{\hphantom{x}}$, maka
  \[ \bigabs{\strut\,\norm{x} - \norm{y}\,} \le \norm{x - y}\,. \]
\end{prop}

\begin{cor} Norma pada suatu ruang linear bernorma merupakan fungsi kontinu seragam.
 \index{kekontinuan!norma}%
 \index{norma!kekontinuan}%
\end{cor}

\begin{notn} Misalkan $M$ suatu ruang metrik, $a \in M$, dan $r > 0$. Kita menyatakan
$\{x \in M\colon d(x,a) < r\}$, yaitu
 \index{B@$B_r(a)$, $B_r$ (bola terbuka berjari-jari $r$)}%
 \index{terbuka!bola}%
 \index{bola!terbuka}%
bola terbuka berjari-jari $r$ yang berpusat di $a$, dengan $B_r(a)$. Demikian pula, kita
menyatakan $\{x \in M\colon d(x,a) \le r\}$, yaitu
 \index{C@$C_r(a)$, $C_r$ (bola tertutup berjari-jari $r$)}%
 \index{tertutup!bola}%
 \index{bola!tertutup}%
\df{bola tertutup} berjari-jari $r$ yang berpusat di $a$, dengan $C_r(a)$, dan
$\{x \in M\colon d(x,a) = r\}$, yaitu
 \index{S@$S_r(a)$, $S_r$ (sfer berjari-jari $r$)}%
 \index{sfer}%
\df{sfer} berjari-jari $r$ yang berpusat di $a$, dengan $S_r(a)$. Jika ruang metrik tersebut
memiliki titik yang ditandai sebagai nol (khususnya, jika ruang itu merupakan ruang linear bernorma), kita menggunakan
$B_r$, $C_r$, dan $S_r$ masing-masing sebagai singkatan untuk $B_r(0)$, $C_r(0)$, dan $S_r(0)$.
\end{notn}

\begin{prop} Jika $V$ merupakan ruang linear bernorma, $x \in V$, dan $r$, $s > 0$, maka
 \begin{enumerate}[label=\rm{(\alph*)}]
   \item $B_r(\vc 0) = -B_r(\vc 0)$;
   \item $B_{rs}(\vc 0) = rB_s(\vc 0)$;
   \item $x + B_r(\vc 0) = B_r(x)$; dan
   \item $B_r(\vc 0) + B_r(\vc 0) = 2\,B_r(\vc 0)$. (Apakah secara umum
$A + A = 2A$ apabila $A$ merupakan subhimpunan dari suatu ruang vektor?)
 \end{enumerate}
\end{prop}

\begin{defn}\label{C017314} Jika $a$ dan $b$ merupakan vektor dalam ruang vektor $V$, maka
 \index{tertutup!ruas}%
 \index{ruas, tertutup}%
\df{ruas tertutup} antara $a$ dan $b$, yang dinotasikan
 \index{<bracsqr@$[a,b]$ (ruas tertutup dalam ruang vektor)}%
dengan~$[a,b]$, adalah $\{(1 - t)a + tb \colon 0 \le t \le 1\}$.
\end{defn}

\begin{cau}\label{C017317} Perhatikan bahwa terdapat sedikit pertentangan antara notasi
untuk \emph{ruas} tertutup ini, ketika diterapkan pada ruang vektor bilangan real $\R$,
dan notasi biasa untuk \emph{interval} tertutup dalam~$\R$. Dalam $\R$, ruas tertutup
$[a,b]$ sama dengan interval tertutup $[a,b]$ asalkan $a \le b$. Namun, jika $a > b$,
ruas tertutup $[a,b]$ sama dengan ruas $[b,a]$ dan memuat semua bilangan $c$ sedemikian
sehingga $b \le c \le a$, sedangkan interval tertutup $[a,b]$ kosong.
\end{cau}

\begin{defn} Subhimpunan $C$ dari suatu ruang vektor $V$ disebut
 \index{konveks!himpunan}%
\df{konveks} jika ruas tertutup $[a,b]$ termuat dalam $C$ setiap kali $a$, $b \in C$.
\end{defn}

\begin{exam}\label{C023169} Dalam ruang linear bernorma, setiap bola terbuka merupakan
himpunan konveks. Demikian pula setiap bola tertutup.
\end{exam}

\begin{prop}\label{prop_intrsctn_convex} Irisan suatu keluarga subhimpunan konveks dari
ruang vektor bersifat konveks.
\end{prop}

\begin{defn} Misalkan $V$ suatu ruang vektor. Ingat bahwa suatu \emph{kombinasi linear}
dari himpunan hingga $\{x_1, \dots, x_n\}$ yang terdiri atas vektor-vektor dalam $V$ adalah
vektor berbentuk $\sum_{k=1}^n \alpha_k x_k$, dengan
$\alpha_1, \dots, \alpha_n \in \R$. Jika
$\alpha_1 = \alpha_2 = \dots = \alpha_n = 0$, maka kombinasi linear tersebut
\emph{trivial}; jika setidaknya satu $\alpha_k$ tidak sama dengan nol, kombinasi linear
tersebut \emph{nontrivial}. Suatu kombinasi linear
$\sum_{k=1}^n\alpha_k x_k$ dari vektor-vektor $x_1, \dots, x_n$ merupakan
 \index{konveks!kombinasi}%
 \index{kombinasi!konveks}%
\df{kombinasi konveks} jika $\alpha_k \ge 0$ untuk setiap $k$ ($1 \le k \le n$) dan jika
$\sum_{k=1}^n \alpha_k = 1$.
\end{defn}

\begin{defn}\label{smplx005def} Misalkan $A$ suatu subhimpunan tak kosong dari ruang
vektor~$V$. Kita mendefinisikan
 \index{konveks!selubung}%
 \index{selubung!konveks}%
 \index{co@$\co A$ (selubung konveks dari~$A$)}%
\df{selubung konveks} dari $A$, yang dinotasikan dengan $\co A$, sebagai subhimpunan
konveks terkecil dari $V$ yang memuat~$A$.
\end{defn}

Definisi sebelumnya seharusnya langsung menimbulkan kecurigaan. Bagaimana kita tahu bahwa
\emph{ada} subhimpunan konveks ``terkecil'' dari~$V$ yang memuat~$A$? Salah satu calon yang
cukup jelas (mari kita sebut calon ``abstrak'') untuk himpunan tersebut adalah irisan dari
\emph{semua} subhimpunan konveks dari $V$ yang memuat~$A$. (Tentu saja, ini juga mungkin tidak
masuk akal---\emph{kecuali} kita tahu bahwa terdapat setidaknya satu himpunan konveks dalam
$V$ yang memuat~$A$ dan bahwa irisan keluarga semua himpunan konveks yang memuat~$A$ itu
sendiri merupakan himpunan konveks yang memuat~$A$.) Calon masuk akal lainnya (yang akan
kita sebut calon ``konstruktif'') adalah himpunan semua kombinasi konveks dari
unsur-unsur~$A$. (Namun, \emph{apakah} ini suatu himpunan konveks? Dan apakah ini himpunan
terkecil yang memuat~$A$?) Akhirnya, sekalipun kedua pendekatan ini masuk akal, apakah
keduanya merupakan definisi yang ekuivalen untuk \emph{selubung konveks}? Dengan kata
lain, apakah keduanya mendefinisikan himpunan yang sama, dan apakah himpunan itu memenuhi
definisi~\ref{smplx005def}?

\begin{prop}\label{smplx006} Definisi-definisi yang disarankan dalam paragraf sebelumnya,
yang diberi nama ``abstrak'' dan ``konstruktif'', masuk akal dan ekuivalen. Selain itu,
himpunan yang dideskripsikan oleh keduanya memenuhi definisi \emph{selubung konveks} yang diberikan
dalam~\ref{smplx005def}.
\end{prop}

\begin{prop} Jika $T\colon V \sto W$ merupakan pemetaan linear antara ruang-ruang vektor
dan $C$ merupakan subhimpunan konveks dari $V$, maka $T^{\sto}(C)$ merupakan subhimpunan
konveks dari~$W$.
\end{prop}

\begin{prop} Dalam ruang linear bernorma, tutupan setiap himpunan konveks bersifat
konveks.
\end{prop}

\begin{prop} Misalkan $V$ suatu ruang linear bernorma. Untuk setiap $a \in V$, pemetaan
$T_a\colon V \sto V\colon x \mapsto x + a$ (yang disebut
 \index{translasi}%
\df{translasi} oleh~$a$) merupakan homeomorfisme.
\end{prop}

\begin{cor} Jika $U$ merupakan subhimpunan terbuka tak kosong dari ruang linear bernorma
$V$, maka $U - U$ memuat suatu lingkungan dari~$0$.
\end{cor}

\begin{prop} Jika $(x_n)$ merupakan barisan dalam ruang linear bernorma dan $x_n \sto a$,
maka $\D\frac1n\sum_{k=1}^nx_k~\sto~a$.
\end{prop}

Dalam proposisi~\ref{00015025}, kita menunjukkan bagaimana suatu hasil kali dalam pada
ruang vektor menginduksi norma pada ruang tersebut. Wajar untuk bertanya apakah semua
norma dapat dihasilkan dengan cara ini dari suatu hasil kali dalam. Jawabannya adalah
\emph{tidak}. Proposisi berikut memberikan syarat perlu dan cukup yang sangat sederhana
agar suatu norma berasal dari hasil kali dalam.

\begin{prop} Misalkan $V$ suatu ruang linear bernorma. Terdapat hasil kali dalam pada $V$
yang menginduksi norma pada $V$ jika dan hanya jika norma tersebut memenuhi
\emph{hukum jajaran genjang}~\ref{parallelogram_law}.
\end{prop}

\begin{proof}[\emph{Petunjuk pembuktian}] Untuk membuktikan bahwa jika suatu norma memenuhi
\emph{hukum jajaran genjang}, maka norma itu diinduksi oleh hasil kali dalam, gunakan
persamaan yang diberikan dalam \emph{identitas polarisasi}
(proposisi~\ref{0001507}) sebagai definisi. Buktikan terlebih dahulu bahwa
$\langle y,x \rangle = \conj{\langle x,y \rangle}$ untuk semua $x$, $y\in V$ dan bahwa
$\langle x,x \rangle > 0$ apabila $x \ne 0$. Selanjutnya, untuk $z \in V$ sebarang,
definisikan fungsi $f\colon V \sto \C\colon x \mapsto \langle x,z \rangle$. Buktikan bahwa
   \[ f(x + y) + f(x - y) = 2f(x) \tag{$\ast$} \]
untuk semua $x$, $y \in V$. Gunakan ini untuk membuktikan bahwa
$f(\alpha x) = \alpha f(x)$ untuk semua $x \in V$ dan $\alpha \in \R$.
(Mulailah dengan $\alpha$ berupa bilangan asli.) Kemudian tunjukkan bahwa $f$ aditif. (Jika
$u$ dan $v$ merupakan unsur sebarang dari $V$, ambil $x = u + v$ dan $y = u - v$.
Gunakan~$(\ast)$.) Terakhir, buktikan bahwa $f(\alpha x) = \alpha f(x)$ untuk $\alpha$
kompleks dengan menunjukkan bahwa $f(ix) = i\,f(x)$ untuk semua $x \in V$. \ns
\end{proof}

\begin{defn} Misalkan $(X,\sfml T)$ suatu ruang topologis. Suatu keluarga
$\sfml B \subseteq \sfml T$ merupakan
 \index{basis!untuk suatu topologi}%
\df{basis} untuk $\sfml T$ jika setiap anggota $\sfml T$ merupakan gabungan anggota-anggota
$\sfml B$. Dengan kata lain, keluarga $\sfml B$ dari himpunan-himpunan terbuka merupakan
basis untuk $\sfml T$ jika untuk setiap himpunan terbuka $U$ terdapat subkeluarga
$\sfml B'$ dari $\sfml B$ sedemikian sehingga $U = \bigcup \sfml B'$.
\end{defn}

\begin{exam}\label{xmpl_top_metric} Pada bidang $\R^2$ dengan topologi biasa
(Euklides), keluarga semua bola terbuka dalam ruang tersebut merupakan basis untuk
topologinya, karena setiap himpunan terbuka dalam $\R^2$ dapat dinyatakan sebagai gabungan
bola-bola terbuka. Bahkan, dalam ruang metrik mana pun, keluarga bola terbuka merupakan
basis untuk suatu topologi pada ruang tersebut. Ini adalah
 \index{topologi!diinduksi oleh metrik}%
 \index{diinduksi!topologi}%
\df{topologi yang diinduksi oleh metrik}.
\end{exam}

Dalam praktiknya (misalnya dalam ruang metrik), sering kali lebih mudah menentukan suatu
basis untuk topologi daripada menentukan topologi itu sendiri. Namun, penting untuk
menyadari bahwa mungkin terdapat banyak basis yang berbeda untuk topologi yang sama.
Setelah suatu basis tertentu dipilih, kita menyebut anggota-anggotanya
 \index{dasar!himpunan terbuka}%
 \index{terbuka!himpunan!dasar}%
\emph{himpunan terbuka dasar}.

Kata ``menginduksi'' dipandang sebagai relasi transitif. Sebagai contoh, jika diberikan
ruang hasil kali dalam $V$, tidak seorang pun akan menyebut \emph{topologi pada $V$ yang
diinduksi oleh metrik yang diinduksi oleh norma yang diinduksi oleh hasil kali dalam}.
Orang cukup menyebut \emph{topologi pada $V$ yang diinduksi oleh hasil kali dalam}.
Kadang-kadang digunakan ``dibangkitkan oleh'' sebagai pengganti ``diinduksi oleh''.

\begin{prop} Misalkan $V$ suatu ruang vektor dengan dua norma $\norm{\,\cdot\,}_1$ dan
$\norm{\,\cdot\,}_2$. Misalkan $\sfml T_k$ topologi pada $V$ yang diinduksi oleh
$\norm{\,\cdot\,}_k$ (untuk $k = 1,2$). Jika terdapat konstanta $\alpha > 0$ sedemikian
sehingga $\norm{x}_1 \le \alpha \norm{x}_2$ untuk setiap $x \in V$, maka
$\sfml T_1 \subseteq \sfml T_2$.
\end{prop}

\begin{defn} Dua norma pada ruang vektor $V$ disebut
 \index{ekuivalen!norma}%
 \index{norma!ekuivalen}%
\df{ekuivalen} jika terdapat konstanta $\alpha$, $\beta > 0$ sedemikian sehingga untuk
semua $x \in V$
   \[ \alpha\norm{x}_1 \le \norm{x}_2 \le \beta \norm{x}_1\,. \]
\end{defn}

\begin{prop}\label{C023184} Jika $\norm{\hphantom{x}}_1$ dan
$\norm{\hphantom{x}}_2$
 \index{ekuivalen!norma!menginduksi topologi yang identik}%
merupakan norma-norma ekuivalen pada ruang vektor~$V$, maka keduanya menginduksi topologi
yang sama pada~$V$.
\end{prop}

Proposisi~\ref{C023184} memberi kita syarat cukup yang mudah diperiksa agar dua norma
menginduksi topologi yang identik pada ruang vektor. Jadi, misalnya, jika kita hendak
menunjukkan bahwa suatu subhimpunan dari ruang linear bernorma bersifat terbuka, boleh jadi
pembuktiannya dapat disederhanakan dengan mengganti norma yang diberikan dengan norma
ekuivalen.

Demikian pula, andaikan kita hendak memeriksa bahwa suatu fungsi $f$ antara dua ruang
linear bernorma bersifat kontinu. Karena kekontinuan didefinisikan dalam hal himpunan
terbuka dan norma-norma ekuivalen menghasilkan himpunan-himpunan terbuka yang persis sama
(lihat proposisi~\ref{C023184}), kita bebas mengganti norma pada domain $f$ dan kodomain
$f$ dengan norma ekuivalen mana pun yang kita kehendaki. Proses ini kadang-kadang dapat
sangat menyederhanakan argumen. (Kemungkinan penyederhanaan ini, secara kebetulan,
merupakan salah satu keuntungan utama menggunakan definisi topologis untuk kekontinuan
sejak awal.)

\begin{defn}\label{C02319} Misalkan $x = (x_n)$ suatu barisan vektor dalam ruang linear
bernorma $V$. Deret tak hingga $\sum_{k=1}^\infty x_k$ disebut
 \index{konvergen!deret}%
 \index{deret!konvergen}%
\df{konvergen} jika terdapat vektor $b \in V$ sedemikian sehingga
$\norm{b - s_n} \sto 0$ ketika $n \sto \infty$, dengan
$s_n = \sum_{k=1}^n x_k$ sebagai jumlah parsial $\text{ke-}n$ dari barisan~$x$. Deret tersebut
disebut
 \index{mutlak!konvergen!deret}%
 \index{konvergen!mutlak}%
\df{konvergen mutlak} jika $\sum_{k=1}^\infty \norm{x_k} < \infty$.
\end{defn}

\begin{exer}\label{C023191} Misalkan
 \index{l@$l_1$!sebagai ruang linear bernorma}%
 \index{l@$l_1$!barisan yang dapat dijumlahkan mutlak}%
 \index{bernorma!ruang linear!$l_1$ sebagai}%
$l_1$ merupakan himpunan semua barisan $x = (x_n)$ dari bilangan kompleks sedemikian
sehingga deret $\sum x_k$ konvergen mutlak. Jadikan $l_1$ ruang vektor dengan definisi
penjumlahan dan perkalian skalar titik demi titik yang biasa. Untuk setiap $x \in l_1$,
definisikan
 \index{<unopn@$\norm{\hphantom{x}}_1$ (norma-$1$)}%
 \index{<unopn@$\norm{\hphantom{x}}_\infty$ (norma seragam)}%
   \begin{align*}
      \norm{x}_1 &:= \sum_{k=1}^\infty \abs{x_k} \\
   \intertext{dan}
      \norm{x}_\infty &:= \sup\{\abs{x_k}\colon k \in \N\}.
   \end{align*}
(Masing-masing adalah
 \index{satu@norma-$1$!dari suatu barisan}%
 \index{norma!$1$- (dari suatu barisan)}%
norma-$1$ dan
 \index{seragam!norma!dari suatu barisan}%
\emph{norma seragam}.) Tunjukkan bahwa $\norm{\hphantom{x}}_1$ dan
$\norm{\hphantom{x}}_\infty$ keduanya merupakan norma pada~$l_1$. Kemudian buktikan atau
sangkal:
 \begin{enumerate}
  \item[(a)] Jika suatu barisan $(x^i)$ dari vektor-vektor dalam ruang linear bernorma
$\bigl(l_1, \norm{\hphantom{x}}_1\bigr)$ konvergen, maka barisan tersebut juga konvergen
dalam $\bigl(l_1, \norm{\hphantom{x}}_\infty\bigr)$.
  \item[(b)] Jika suatu barisan $(x^i)$ dari vektor-vektor dalam ruang linear bernorma
$\bigl(l_1, \norm{\hphantom{x}}_\infty\bigr)$ konvergen, maka barisan tersebut juga
konvergen dalam $\bigl(l_1, \norm{\hphantom{x}}_1\bigr)$.
 \end{enumerate}
\end{exer}








\section{Pemetaan Linear Terbatas}\label{sec_bdd_lin_maps}
Para analis bekerja dengan objek-objek yang memiliki struktur aljabar sekaligus topologis.
Pada bagian sebelumnya kita mengkaji ruang-ruang vektor yang dilengkapi dengan topologi
yang berasal dari suatu norma. Jika ruang-ruang itu merupakan objek yang sedang
dipertimbangkan, morfisme apakah yang kiranya paling menarik? Jawaban yang masuk akal,
dan memang benar, adalah \emph{pemetaan yang melestarikan struktur topologis maupun
aljabar}; yakni, pemetaan linear kontinu. Kategori yang dihasilkan dinotasikan dengan
 \index{kategori!$\cat{NLS_\infty}$ sebagai suatu}%
 \index{NLSinf@$\cat{NLS_{\boldsymbol\infty}}$!kategori}
$\cat{NLS_{\boldsymbol\infty}}$. Secara topologis, isomorfisme dalam kategori ini adalah
homeomorfisme.

Berikut adalah suatu syarat yang sangat berguna yang, untuk pemetaan linear, ternyata
ekuivalen dengan kekontinuan.

\begin{defn} Suatu pemetaan linear $T\colon V \sto W$ antara ruang-ruang linear bernorma
disebut
 \index{terbatas!pemetaan linear}%
 \index{linear!pemetaan!terbatas}%
\df{terbatas} jika $T(B)$ merupakan subhimpunan terbatas dari $W$ setiap kali $B$
merupakan subhimpunan terbatas dari~$V$. Dengan kata lain, pemetaan linear terbatas
membawa himpunan terbatas ke himpunan terbatas. Kita menyatakan
 \index{B@$\ofml B(V,W)$!pemetaan linear terbatas}%
dengan~$\ofml B(V,W)$ keluarga semua pemetaan linear terbatas dari $V$ ke~$W$. Pemetaan
linear terbatas dari suatu ruang $V$ ke dirinya sendiri sering disebut suatu
 \index{operator}%
 \index{konvensi!semua operator terbatas dan linear}%
\df{operator} (linear terbatas). Keluarga semua operator pada ruang linear bernorma $V$
dinotasikan
 \index{B@$\ofml B(V)$!operator}%
dengan~$\ofml B(V)$. Kelas ruang linear bernorma bersama pemetaan linear terbatas di antara
ruang-ruang tersebut membentuk suatu kategori. Di bawah ini kita memeriksa bahwa kategori
ini tidak lain adalah kategori~$\cat{NLS_{\boldsymbol\infty}}$.
\end{defn}

\begin{cau} Sangat penting untuk menyadari bahwa secara umum suatu pemetaan linear
terbatas \emph{bukan} fungsi terbatas dalam arti biasa dari suatu \emph{fungsi terbatas}
pada himpunan tertentu (lihat contoh~\ref{C023131} dan~\ref{C023133}). Penggunaan kata
``terbatas'' dalam kedua arti yang bertentangan ini mungkin kurang menguntungkan, tetapi
sudah mapan.
\end{cau}

\begin{exam} Pemetaan linear $T\colon \R \sto \R\colon x \mapsto 3x$ merupakan pemetaan
linear terbatas (karena memetakan subhimpunan terbatas dari $\R$ ke subhimpunan terbatas
dari~$\R$), tetapi jika dipandang hanya sebagai fungsi, $T$ tidak terbatas (karena
jangkauannya bukan subhimpunan terbatas dari~$\R$).
\end{exam}

Pengamatan berikut dapat membantu mengurangi kebingungan.

\begin{prop} Suatu pemetaan linear, kecuali jika merupakan pemetaan konstan yang membawa
setiap vektor ke nol, \emph{tidak mungkin} merupakan fungsi terbatas.
\end{prop}

\begin{defn} Misalkan $(M,d)$ dan $(N,\rho)$ ruang-ruang metrik. Suatu fungsi
$f \colon M \sto N$ disebut
 \index{seragam!kekontinuan}%
 \index{kekontinuan!seragam}%
\df{kontinu seragam} jika untuk setiap $\epsilon > 0$ terdapat $\delta > 0$ sedemikian
sehingga $\rho(f(x),f(y)) < \epsilon$ setiap kali $x$, $y \in M$ dan
$d(x,y) < \delta$.
\end{defn}

Salah satu aspek linearitas yang paling mencolok adalah bahwa untuk pemetaan linear,
konsep kekontinuan, kekontinuan di satu titik saja, dan kekontinuan seragam menyatu.
Bahkan, semua konsep tersebut persis sama dengan keterbatasan.

\begin{thm}\label{C023221} Misalkan $T \colon V \sto W$ suatu pemetaan linear antara
ruang-ruang linear bernorma. Maka pernyataan-pernyataan berikut ekuivalen:
  \begin{enumerate}[label=\rm{(\alph*)}]
    \item $T$ terbatas.
    \item Citra bola satuan tertutup di bawah $T$ merupakan himpunan terbatas.
    \item $T$ kontinu seragam pada $V$.
    \item $T$ kontinu pada $V$.
    \item $T$ kontinu di~$\vc 0$.
    \item Terdapat bilangan $M > 0$ sedemikian sehingga
$\norm{Tx} \le M\norm{x}$ untuk semua $x \in V$.
  \end{enumerate}
\end{thm}

\begin{prop} Misalkan $T\colon V \sto W$ suatu pemetaan linear terbatas antara ruang-ruang
linear bernorma dengan domain bukan ruang nol. Maka keempat bilangan berikut ada dan bernilai sama.
 \begin{enumerate}[label=\rm{(\alph*)}]
  \item $\sup\{\norm{Tx}\colon \norm{x} \le 1\}$
  \item $\sup\{\norm{Tx}\colon \norm{x} = 1\}$
  \item $\sup\{\norm{Tx}\,\norm{x}^{-1}\colon x \ne \vc 0\}$
  \item $\inf\{M > 0\colon \norm{Tx} \le M\norm{x}\ \text{ untuk semua $x \in V$}\}$
 \end{enumerate}
\end{prop}

\begin{defn}\label{defn_op_norm} Jika $T$ merupakan pemetaan linear terbatas, maka
$\norm{T}$, yang disebut
 \index{norma!pemetaan linear terbatas}%
 \index{<unopn@$\norm T$ (norma pemetaan linear terbatas~$T$)}%
\df{norma} (atau
 \index{norma!operator}%
 \index{operator!norma}%
\df{norma operator}) dari $T$, didefinisikan sebagai salah satu dari empat ungkapan dalam
proposisi sebelumnya; untuk domain nol, nilainya ditetapkan nol.
\end{defn}

\begin{exam}\label{exam_op_norm_norm} Jika $V$ dan $W$ merupakan ruang-ruang linear
bernorma, maka fungsi
   \[ \norm{\hphantom{T}}\colon \ofml B(V,W) \sto \R
      \colon T \mapsto \norm{T} \]
memang merupakan suatu norma.
\end{exam}

\begin{exam}\label{C023228} Keluarga $\ofml B(V,W)$ dari semua pemetaan linear terbatas
 \index{B@$\ofml B(V,W)$!sebagai ruang linear bernorma}%
 \index{bernorma!ruang linear!$\ofml B(V,W)$ sebagai}%
antara ruang-ruang linear bernorma itu sendiri merupakan ruang linear bernorma di bawah
norma operator yang didefinisikan di atas serta operasi penjumlahan dan perkalian skalar
titik demi titik yang biasa. Yang khususnya penting adalah keluarga
$V^* = \ofml B(V,\K)$ dari anggota-anggota \emph{kontinu} dalam~$V^{\#}$, yaitu
 \index{kontinu!fungsional linear}%
 \index{linear!fungsional!kontinu}%
 \index{fungsional!linear kontinu}%
\emph{fungsional linear kontinu}, yang juga dikenal sebagai
 \index{terbatas!fungsional linear}%
 \index{linear!fungsional!terbatas}%
 \index{fungsional!linear terbatas}%
\emph{fungsional linear terbatas}.
Ingat bahwa dalam~\ref{def_alg_dual} kita mendefinisikan dual aljabar $V^{\#}$ dari ruang
vektor~$V$. Yang lebih penting dalam pembahasan kita di sini adalah $V^*$, yang disebut
 \index{<unopurstar@$V^*$ (ruang dual dari ruang linear bernorma)}%
 \index{V@$V^*$ (ruang dual dari ruang linear bernorma)}%
 \index{dual!ruang}%
 \index{ruang!dual dari ruang linear bernorma}%
\df{ruang dual} dari~$V$. Untuk membedakannya dari dual aljabar dan dual urutan, ruang ini
kadang-kadang disebut
 \index{dual!norma}%
 \index{norma!dual}%
\df{dual norma} atau
 \index{dual!topologis}%
 \index{topologis!dual}%
\df{dual topologis} dari~$V$. Perhatikan bahwa pembedaan ini tidak diperlukan dalam ruang
berdimensi hingga karena di sana $V^* = V^{\#}$ (lihat~\ref{00015033}).
\end{exam}

\begin{prop}\label{C023228a} Keluarga $\ofml B(V,W)$ dari semua pemetaan linear terbatas
 \index{B@$\ofml B(V,W)$!sebagai ruang Banach}%
 \index{Banach!ruang!$\ofml B(V,W)$ sebagai}%
antara ruang-ruang linear bernorma bersifat lengkap (terhadap metrik yang diinduksi oleh
normanya) setiap kali $W$ lengkap. Secara khusus, ruang dual dari setiap ruang linear
bernorma bersifat lengkap.
\end{prop}
















\begin{prop}\label{0003073} Misalkan $U$, $V$, dan $W$ ruang-ruang linear bernorma. Jika
$S \in \ofml B(U,V)$ dan $T \in \ofml B(V,W)$, maka $TS \in \ofml B(U,W)$ dan
$\norm{TS} \le \norm T \norm S$.
\end{prop}

\begin{exam} Pada setiap ruang linear bernorma $V$ yang bukan ruang nol,
 \index{operator!identitas}%
 \index{identitas!operator}%
\emph{operator identitas}
   \[ \id V = I_V =I\colon V \sto V \colon v \mapsto v \]
terbatas dan $\norm{I_V} = 1$.
 \index{operator!nol}%
 \index{nol!operator}%
\emph{Operator nol}
   \[ \vc 0_V = \vc 0\colon V \sto V\colon v \mapsto \vc 0 \]
juga terbatas dan $\norm{0_V} = 0$.
\end{exam}

\begin{exam} Misalkan
$T \colon \R^2 \sto \R^3\colon (x,y) \mapsto (3x, x + 2y, x - 2y)$.
Maka~$\norm{T} = \sqrt{11}$.
\end{exam}

Banyak mahasiswa yang baru saja menyelesaikan mata kuliah kalkulus dasar merasa bahwa
diferensiasi merupakan operasi yang ``lebih baik'' daripada integrasi---mungkin
karena aturan integrasi tampak lebih rumit daripada aturan diferensiasi. Namun,
ketika kita memandang keduanya sebagai pemetaan linear, yang benar justru sebaliknya.

\begin{exam} Sebagai pemetaan linear pada ruang fungsi-fungsi yang dapat didiferensialkan
pada $[0,1]$ (dengan norma seragam), integrasi terbatas; diferensiasi tidak terbatas
(dan karena itu tidak kontinu di titik mana pun!).
\end{exam}


\begin{defn} Misalkan $f\colon V_1 \sto V_2$ suatu fungsi antara ruang-ruang linear
bernorma. Untuk $k = 1,2$, misalkan $\norm{\hphantom{M}}_k$ norma pada~$V_k$ dan $d_k$
metrik pada $V_k$ yang diinduksi oleh~$\norm{\hphantom{M}}_k$ (lihat
contoh~\ref{0001503}). Maka $f$ merupakan
 \index{isometri}%
\df{isometri} (atau
suatu pemetaan \df{isometrik}) jika $d_2(f(x),f(y)) = d_1(x,y)$ untuk semua
 $x$, $y \in V_1$. Pemetaan itu dikatakan
 \index{norma!dipertahankan}%
\df{mempertahankan norma} jika $\norm{f(x)}_2 = \norm{x}_1$ untuk semua $x \in V_1$.
\end{defn}

Sebelumnya dalam bagian ini kita telah membahas kategori
$\cat{NLS_{\boldsymbol\infty}}$ dari ruang linear bernorma dan pemetaan linear kontinu.
Dalam kategori ini, isomorfisme melestarikan struktur topologis (selain struktur ruang
vektor). Apa yang akan terjadi jika kita justru menghendaki isomorfisme melestarikan
struktur metrik dari ruang-ruang linear bernorma; yakni, menjadi isometri? Mudah
dilihat bahwa (dengan adanya linearitas) hal ini ekuivalen dengan meminta agar isomorfisme
mempertahankan norma. Dalam kategori baru ini, yang dinotasikan dengan
 \index{kategori!$\cat{NLS_{\vc 1}}$ sebagai suatu}%
 \index{NLS1@$\cat{NLS_{\vc 1}}$!kategori}
$\cat{NLS_{\vc 1}}$, morfismenya adalah
 \index{kontraktif}%
\df{pemetaan linear kontraktif}, yaitu pemetaan linear $T\colon V \sto W$ antara
ruang-ruang linear bernorma sedemikian sehingga $\norm{Tx} \le \norm x$ untuk semua
$x \in V$. Tampaknya tepat untuk menyebut $\cat{NLS_{\boldsymbol\infty}}$ sebagai
 \index{topologis!kategori}%
 \index{kategori!topologis}%
\emph{kategori topologis} ruang linear bernorma dan $\cat{NLS_{\vc 1}}$ sebagai
 \index{kategori geometris}%
 \index{kategori!geometris}%
\emph{kategori geometris} ruang linear bernorma. Banyak penulis menggunakan istilah
``isomorfisme'' tanpa pengubah untuk isomorfisme dalam kategori topologis
$\cat{NLS_{\boldsymbol\infty}}$ dan
 \index{isometrik!isomorfisme}%
 \index{isomorfisme!isometrik}%
``isomorfisme isometrik'' untuk isomorfisme dalam kategori geometris
$\cat{NLS_{\vc 1}}$. Dalam catatan ini kita akan berfokus pada kategori topologis.
Teori isometrik ruang linear bernorma, meskipun penting, agak lebih khusus.

\begin{exer} Periksalah pernyataan-pernyataan yang belum dibuktikan dalam paragraf
sebelumnya. Secara khusus, buktikan bahwa
 \begin{enumerate}
    \item[(a)] $\cat{NLS_{\boldsymbol\infty}}$ merupakan suatu kategori.
    \item[(b)] Suatu isomorfisme dalam $\cat{NLS_{\boldsymbol\infty}}$ merupakan
isomorfisme ruang vektor sekaligus homeomorfisme.
    \item[(c)] Suatu pemetaan linear antara ruang-ruang linear bernorma merupakan isometri
jika dan hanya jika pemetaan tersebut mempertahankan norma.
    \item[(d)] $\cat{NLS_{\vc 1}}$ merupakan suatu kategori.
    \item[(e)] Suatu isomorfisme dalam $\cat{NLS_{\vc 1}}$ merupakan isometri sekaligus
isomorfisme ruang vektor.
    \item[(f)] Kategori $\cat{NLS_{\vc 1}}$ merupakan subkategori dari
$\cat{NLS_{\boldsymbol\infty}}$ dalam arti bahwa setiap morfisme kategori pertama
termasuk dalam kategori kedua.
 \end{enumerate}
\end{exer}













\section{Ruang Berdimensi Hingga}
Dalam bagian ini dicantumkan beberapa fakta baku dan berguna tentang ruang berdimensi hingga. Bukti hasil-hasil ini
cukup mudah ditemukan. Lihat, misalnya, \cite{BrownP:1970}, Bagian 4.4; \cite{Conway:1990}, Bagian III.3;
\cite{Kreyszig:1978}, Bagian 2.4--5; atau \cite{Rudin:1991}, halaman 16--18.

\begin{defn} Suatu barisan $(x_n)$ dalam ruang metrik disebut
 \index{Cauchy!barisan}%
 \index{barisan!Cauchy}%
\df{barisan Cauchy} jika untuk setiap $\epsilon > 0$ terdapat $n_0 \in \N$ sedemikian sehingga
$d(x_m,x_n) < \epsilon$ apabila $m$, $n \ge n_0$. Syarat ini sering disingkat sebagai
berikut: $d(x_m, x_n) \sto 0$ ketika $m$, $n \sto \infty$ (atau $\lim_{m,n \sto
\infty}d(x_m,x_n) = 0$).
\end{defn}

\begin{defn} Suatu ruang metrik $M$ disebut
 \index{lengkap!ruang metrik}%
 \index{ruang metrik!lengkap}%
 \index{ruang!metrik lengkap}%
\df{lengkap} jika setiap barisan Cauchy dalam $M$ konvergen. (Untuk meninjau kembali fakta-fakta paling dasar tentang ruang
metrik lengkap, lihat bagian pertama Bab 20 dalam catatan daring saya~\cite{Erdman:2007}.)
\end{defn}

\begin{prop}\label{prop_fdnls_Kn} Jika $V$ adalah ruang linear bernorma berdimensi $n$ atas $\K$, maka terdapat suatu pemetaan
 \index{berdimensi hingga!ruang linear bernorma}%
dari $V$ ke $\K^n$ yang sekaligus merupakan homeomorfisme dan isomorfisme ruang vektor.
\end{prop}

\begin{cor} Setiap ruang linear bernorma berdimensi hingga bersifat lengkap.
\end{cor}

\begin{cor}\label{cor_fdvs_closed} Setiap subruang vektor berdimensi hingga dari suatu ruang linear bernorma bersifat tertutup.
\end{cor}

\begin{cor}\label{cor_fdvs_norms_equiv} Sebarang dua norma pada suatu ruang vektor berdimensi hingga bersifat ekuivalen.
\end{cor}

\begin{defn} Subhimpunan $K$ dari ruang topologis $X$ disebut
 \index{kompak}%
\df{kompak} jika setiap keluarga $\sfml U$ yang terdiri atas subhimpunan-subhimpunan terbuka dari $X$ dan gabungannya memuat $K$ memiliki
subkeluarga hingga dari $\sfml U$ yang gabungannya memuat~$K$.
\end{defn}

\begin{prop}\label{00015032} Bola satuan tertutup dalam suatu ruang linear bernorma bersifat kompak jika dan hanya jika
ruang tersebut berdimensi hingga.
\end{prop}

\begin{prop}\label{00015033} Jika $V$ dan $W$ adalah ruang-ruang linear bernorma dan $V$ berdimensi hingga, maka
setiap pemetaan linear $T \colon V \sto W$ bersifat kontinu.
\end{prop}

Misalkan $T$ suatu pemetaan linear terbatas antara ruang-ruang linear bernorma (yang mungkin berdimensi tak hingga). Sama seperti dalam kasus
berdimensi hingga, kernel dan jangkauan $T$ merupakan objek yang sangat penting. Baik dalam kasus berdimensi hingga maupun tak hingga,
keduanya merupakan subruang vektor dari ruang yang memuatnya. Selain itu, dalam kedua kasus tersebut kernel $T$ bersifat tertutup.
(Di bawah suatu fungsi kontinu, pracitra himpunan tertutup bersifat tertutup.) Akan tetapi, terdapat perbedaan besar antara
kedua kasus itu. Dalam kasus berdimensi hingga, jangkauan $T$ harus tertutup (berdasarkan~\ref{cor_fdvs_closed}). Seperti akan kita lihat
dalam Contoh~\futurexref{5.2.14}{exam_ran_nonclosed}, pemetaan linear antara ruang-ruang berdimensi tak hingga tidak harus memiliki jangkauan tertutup.

\section{Hasil Bagi Ruang Linear Bernorma}

\begin{defn}\label{quo001def} Misalkan $A$ adalah objek dalam suatu kategori konkret $\cat C$.
Suatu morfisme surjektif $A \to^\pi B$ dalam $\cat C$ disebut
 \index{hasil bagi!pemetaan!untuk kategori konkret}%
 \index{pi@$\pi$ (pemetaan hasil bagi)}%
\df{pemetaan hasil bagi} untuk $A$ jika suatu fungsi $g\colon B \sto C$ (dalam $\cat{SET}$) merupakan morfisme
(dalam $\cat C$) setiap kali $g \circ \pi$ merupakan morfisme. Suatu objek $B$ dalam $\cat C$ disebut
 \index{hasil bagi!objek}%
\df{objek hasil bagi} untuk $A$ jika merupakan citra suatu pemetaan hasil bagi untuk~$A$.
\end{defn}

\begin{exam} Dalam kategori $\cat{VEC}$ yang terdiri atas ruang-ruang vektor dan pemetaan-pemetaan linear, setiap pemetaan linear
 \index{hasil bagi!pemetaan!dalam $\cat{VEC}$}%
surjektif merupakan pemetaan hasil bagi.
\end{exam}

\begin{exam}\label{C021734} Dalam kategori $\cat{TOP}$, tidak setiap epimorfisme merupakan
 \index{hasil bagi!pemetaan!dalam $\cat{TOP}$}%
pemetaan hasil bagi.
\end{exam}

\begin{proof}[\emph{Petunjuk pembuktian}] Tinjau pemetaan identitas pada bilangan real dengan topologi yang berbeda
pada domain dan kodomainnya. \ns
\end{proof}

Butir berikut, yang semestinya sudah dikenal dari aljabar linear, menunjukkan cara menghasilkan suatu objek hasil bagi tertentu dengan
``mengambil hasil bagi oleh suatu subruang''.

\begin{defn}\label{quo002def} Misalkan $M$ adalah subruang dari ruang vektor $V$. Definisikan suatu
relasi ekuivalensi $\sim$ pada $V$ dengan
   \[ x \sim y \qquad \text{jika dan hanya jika} \qquad y - x \in M. \]
Untuk setiap $x \in V$, misalkan $[x]$ adalah kelas ekuivalensi yang memuat~$x$.
 \index{<bracsqr@$[x]$ (kelas ekuivalensi yang memuat~$x$)}%
 \index{<binop@$V/M$ (hasil bagi $V$ oleh $M$)}%
Misalkan $V/M$ adalah himpunan semua kelas ekuivalensi unsur-unsur~$V$. Untuk $[x]$ dan $[y]$
dalam $V/M$, definisikan
  \[ [x] + [y] := [x+y]\]
dan untuk $\alpha \in \K$ serta $[x] \in V/M$, definisikan
  \[\alpha[x] := [\alpha x].\]
Terhadap operasi-operasi ini, $V/M$ menjadi ruang vektor. Ruang ini adalah
 \index{hasil bagi!ruang vektor}%
 \index{ruang!hasil bagi!vektor}%
\df{ruang hasil bagi} dari $V$ oleh~$M$. Notasi $V/M$ biasanya dibaca ``$V$ modulo~$M$''.
Pemetaan linear
  \[ \pi\colon V \sto V/M\colon x \mapsto [x] \]
disebut
 \index{hasil bagi!pemetaan!untuk ruang vektor}%
 \index{pi@$\pi$ (pemetaan hasil bagi)}%
\df{pemetaan hasil bagi}.
\end{defn}

\begin{exer} Verifikasilah pernyataan-pernyataan dalam Definisi~\ref{quo002def}. Secara khusus,
tunjukkan bahwa $\sim$ merupakan relasi ekuivalensi, bahwa penjumlahan dan perkalian skalar pada
himpunan kelas-kelas ekuivalensi terdefinisi dengan baik, bahwa terhadap operasi-operasi ini $V/M$ merupakan
ruang vektor, bahwa fungsi yang di sini disebut ``pemetaan hasil bagi'' memang merupakan pemetaan hasil bagi
dalam pengertian~\ref{quo001def}, dan bahwa pemetaan hasil bagi ini bersifat linear.
\end{exer}

Hasil berikut disebut \emph{teorema dasar hasil bagi} atau
 \index{hasil bagi!teorema!untuk $\cat{VEC}$}%
 \index{teorema isomorfisme pertama}%
{teorema isomorfisme pertama} untuk ruang-ruang vektor.

\begin{thm}\label{quo005} Misalkan $V$ dan $W$ adalah ruang-ruang vektor dan $M \preceq V$. Jika
$T \in \ofml L(V,W)$ dan $\ker T \supseteq M$, maka terdapat tepat satu $\wt T \in
\ofml L(V/M\,,W)$ yang membuat diagram berikut komutatif.
 \[\xymatrix@+30pt{V \ar[d]_*+{\pi} \ar[dr]^*+{T} & \\
            V/M \ar[r]_*+{\widetilde T} & W  }\]
Selanjutnya, $\wt T$ injektif jika dan hanya jika $\ker T = M$; dan $\wt T$ surjektif
jika dan hanya jika $T$ surjektif.
\end{thm}

\begin{cor} Jika $T\colon V \sto W$ adalah pemetaan linear antara ruang-ruang vektor, maka
$\ran T \cong V/\ker T$.
\end{cor}

\begin{prop}\label{C023711} Misalkan $V$ adalah ruang linear bernorma dan $M$ adalah subruang tertutup
dari~$V$.
 \index{hasil bagi!ruang linear bernorma}%
 \index{bernorma!ruang linear!hasil bagi}%
 \index{ruang!hasil bagi!linear bernorma}%
Maka pemetaan
  \[ \norm{\hphantom{x}} \colon V/M \sto \R \colon [x]
                        \mapsto \inf\{\norm u \colon u \sim x\} \]
merupakan norma pada~$V/M$. Selanjutnya,
 \index{pi@$\pi$ (pemetaan hasil bagi)}%
 \index{hasil bagi!pemetaan!untuk ruang linear bernorma}%
pemetaan hasil bagi
  \[ \pi \colon V \sto V/M \colon x \mapsto [x] \]
merupakan surjeksi linear terbatas dengan $\norm\pi \le 1$.
\end{prop}

\begin{defn} Norma yang didefinisikan pada $V/M$ disebut
 \index{hasil bagi!norma}%
 \index{norma!hasil bagi}%
\df{norma hasil bagi} pada~$V/M$. Terhadap norma ini $V/M$ merupakan ruang linear bernorma dan disebut
 \index{ruang hasil bagi}%
 \index{ruang!hasil bagi}%
\df{ruang hasil bagi} dari $V$ oleh~$M$. Ruang ini sering disebut sebagai ``$V$ modulo $M$''.
\end{defn}

\begin{thm}[Teorema dasar hasil bagi untuk $\cat{NLS_{\boldsymbol\infty}}$]\label{C023717}
 \index{hasil bagi!teorema!untuk $\cat{NLS_{\boldsymbol\infty}}$}%
Misalkan $V$ dan $W$ adalah ruang-ruang linear bernorma dan $M$ adalah subruang tertutup dari~$V$. Jika $T$ merupakan
pemetaan linear terbatas dari $V$ ke $W$ dan $\ker T \supseteq M$, maka terdapat suatu
pemetaan linear terbatas tunggal $\wt T\colon V/M \sto W$ yang membuat diagram berikut komutatif.
  \[ \xymatrix@+30pt{V \ar[d]_*+{\pi} \ar[dr]^*+{T} & \\
             V/M \ar[r]_*+{\wt T} & W  } \]
Selanjutnya: $\norm{\wt T} = \norm{T}$; $\wt T$ injektif jika dan hanya jika $\ker T = M$;
dan $\wt T$ surjektif jika dan hanya jika $T$ surjektif.
\end{thm}

\section{Hasil Kali Ruang Linear Bernorma}\label{C0155}
Hasil kali dan koproduk, seperti hasil bagi, paling baik dideskripsikan dari segi apa yang \emph{dilakukannya}.

\begin{defn}\label{C015511} Misalkan $A_1$ dan $A_2$ adalah objek-objek dalam suatu kategori~$\cat C$. Kita mengatakan bahwa
tripel $(P,\pi_1,\pi_2)$, dengan $P$ suatu objek dan $\pi_k\colon P \sto A_k$ ($k =
1,2$) morfisme-morfisme, merupakan suatu
 \index{hasil kali!dalam suatu kategori}%
\df{hasil kali} dari $A_1$ dan $A_2$ jika untuk setiap objek $B$ dan setiap pasangan morfisme
$f_k\colon B \sto A_k$ ($k = 1,2$) terdapat morfisme tunggal $f\colon B \sto P$ sedemikian
sehingga $f_k = \pi_k \circ f$ untuk $k = 1,2$.

Sudah lazim untuk mengatakan, ``Misalkan $P$ suatu hasil kali dari \dots'' sebagai pengganti, ``Misalkan $(P,\pi_1,\pi_2)$
suatu hasil kali dari \dots''. Hasil kali $A_1$ dan $A_2$ sering ditulis sebagai $A_1 \times
A_2$ atau sebagai $\prod_{k = 1,2} A_k$.
\end{defn}

Dalam suatu kategori tertentu, hasil kali mungkin ada ataupun tidak. Merupakan fakta yang menarik dan
dasar bahwa kapan pun hasil kali itu ada, hasil kali tersebut unik (hingga isomorfisme), sehingga kita
dapat berbicara tanpa ambiguitas tentang \emph{hasil kali} dua objek. Ketika kita mengatakan bahwa suatu
objek kategori yang memenuhi satu atau beberapa syarat bersifat \emph{unik hingga isomorfisme}, tentu saja kita
bermaksud bahwa sebarang dua objek yang memenuhi syarat-syarat tersebut harus isomorfik. Kita
akan sering menggunakan frasa
 \index{esensial!unik}%
 \index{konvensi!unik secara esensial berarti unik hingga isomorfisme}%
 \index{ketunggalan!esensial}%
``unik secara esensial'' untuk ``unik hingga isomorfisme''.

\begin{prop} Dalam sebarang kategori, hasil kali (jika ada) bersifat
 \index{hasil kali!ketunggalan}%
 \index{ketunggalan!hasil kali}%
unik secara esensial.
\end{prop}

\begin{exam}\label{C015517} Dalam kategori $\cat{SET}$,
 \index{hasil kali!dalam $\cat{SET}$}%
 \index{SET@$\cat{SET}$!hasil kali dalam}%
hasil kali dua himpunan $A_1$ dan $A_2$ ada dan sesungguhnya merupakan hasil kali Kartesius biasa
$A_1 \times A_2$ beserta proyeksi koordinat biasa $\pi_k\colon A_1 \times
A_2 \sto A_k \colon (a_1,a_2) \mapsto a_k$.
\end{exam}

\begin{exam}\label{C015521} Jika $V_1$ dan $V_2$ adalah ruang-ruang vektor, kita menjadikan hasil kali Kartesius
$V_1 \times V_2$ suatu ruang vektor sebagai berikut. Definisikan penjumlahan dengan
  \[ (u,v) + (w,x) := (u + w, v + x)  \]
dan perkalian skalar dengan
  \[ \alpha (u,v) := (\alpha u, \alpha v) \,. \]
Hal ini menjadikan $V_1 \times V_2$ suatu ruang vektor, dan ruang ini beserta
proyeksi koordinat biasa $\pi_k\colon V_1 \times V_2 \sto V_k \colon (v_1,v_2)
\mapsto v_k$ ($k = 1,2$) merupakan
 \index{hasil kali!dalam $\cat{VEC}$}%
 \index{VEC@$\cat{VEC}$!hasil kali dalam}%
 \index{vektor!ruang!hasil kali}%
suatu hasil kali dalam kategori $\cat{VEC}$. Ruang ini biasanya disebut
 \index{langsung!jumlah!ruang vektor}%
 \index{jumlah!langsung}%
 \index{vektor!ruang!jumlah langsung}%
 \index{<binopdirectsum@$\oplus$ (jumlah langsung)}%
\df{jumlah langsung} dari $V_1$ dan $V_2$ dan dinotasikan dengan~$V_1 \oplus V_2$.
\end{exam}

Sering kali mungkin dan dikehendaki untuk mengambil hasil kali dari sebarang keluarga objek dalam
suatu kategori. Berikut adalah perumuman Definisi~\ref{C015511}.

\begin{defn}\label{C015524} Misalkan $\bigl(A_\lambda\bigr)_{\lambda \in \Lambda}$ adalah suatu
keluarga terindeks objek-objek dalam kategori~$\cat C$. Kita mengatakan bahwa objek $P$ beserta
suatu keluarga terindeks $\bigl(\pi_\lambda\bigr)_{\lambda \in \Lambda}$ dari morfisme-morfisme
$\pi_\lambda\colon P \sto A_\lambda$ merupakan
 \index{hasil kali!dalam suatu kategori}%
\df{hasil kali} objek-objek $A_\lambda$ jika untuk setiap objek $B$ dan setiap keluarga terindeks
$\bigl(f_\lambda\bigr)_{\lambda \in \Lambda}$ dari morfisme-morfisme $f_\lambda\colon B \sto
A_\lambda$ terdapat pemetaan tunggal $g \colon B \sto P$ sedemikian sehingga $f_\lambda =
\pi_\lambda \circ g$ untuk setiap $\lambda \in \Lambda$.
\end{defn}

Suatu kategori yang di dalamnya sebarang hasil kali ada disebut
 \index{hasil kali!lengkap}%
 \index{lengkap!hasil kali}%
\df{lengkap terhadap hasil kali}. Banyak kategori yang kita jumpai dalam catatan ini lengkap
terhadap hasil kali.

\begin{defn}\label{C015531} Misalkan $\bigl(S_\lambda\bigr)_{\lambda \in \Lambda}$ adalah suatu keluarga
terindeks himpunan.
 \index{hasil kali Kartesius}%
 \index{hasil kali!Kartesius}%
\df{Hasil kali Kartesius} dari keluarga terindeks tersebut, yang dinotasikan dengan
 \index{<<@$\prod S_\lambda$ (hasil kali Kartesius)}%
$\prod_{\lambda \in \Lambda} S_\lambda$ atau cukup $\prod S_\lambda$, adalah himpunan semua
fungsi $f \colon \Lambda \sto \bigcup S_\lambda$ sedemikian sehingga $f(\lambda) \in S_\lambda$
untuk setiap $\lambda \in \Lambda$. Pemetaan-pemetaan $\pi_\lambda \colon \prod S_\lambda \sto
S_\lambda \colon f \mapsto f(\lambda)$ merupakan
 \index{koordinat!proyeksi}%
 \index{pi@$\pi_\lambda$ (proyeksi koordinat)}%
 \index{proyeksi!ke koordinat}%
\df{proyeksi koordinat} kanonik. Dalam banyak kasus, notasi $f_\lambda$ lebih mudah digunakan
daripada~$f(\lambda)$. (Lihat, misalnya, Contoh~\ref{C015537} di bawah.)
\end{defn}

\begin{exam}\label{C015534} Suatu kasus khusus yang sangat penting dari definisi sebelumnya muncul
ketika semua himpunan $S_\lambda$ identik: misalkan $S_\lambda = A$ untuk setiap $\lambda
\in \Lambda$. Dalam hal ini, hasil kali Kartesius tersebut terdiri atas semua fungsi yang memetakan
$\Lambda$ ke~$A$. Artinya, $\prod_{\lambda \in \Lambda} S_\lambda = A^\Lambda$. Perhatikan
pula bahwa dalam kasus ini proyeksi koordinat merupakan
 \index{evaluasi!pemetaan}%
 \index{fungsi!evaluasi pada suatu titik}%
\emph{pemetaan evaluasi}. Untuk setiap $\lambda$, proyeksi koordinat $\pi_\lambda$ membawa
setiap titik $f$ dalam hasil kali tersebut (yaitu, setiap fungsi $f$ dari $\Lambda$ ke~$A$) ke
$f(\lambda)$, nilainya pada~$\lambda$. Singkatnya, setiap proyeksi koordinat merupakan
pemetaan evaluasi pada suatu titik.
\end{exam}

\begin{exam}\label{C015537} Apakah $\R^n$ itu? Ruang ini hanyalah himpunan semua $n$-tupel bilangan real.
Artinya, ruang ini adalah himpunan fungsi dari $\N_n = \{1,2,\dots,n\}$ ke~$\R$. Dengan
kata lain, $\R^n$ hanyalah singkatan untuk~$\R^{\N_n}$. Dalam $\R^n$, biasanya orang menulis $x_j$ untuk
koordinat ke-$j$ dari suatu vektor $x$, bukan~$x(j)$.
\end{exam}

\begin{exam} Dalam kategori $\cat{SET}$,
 \index{hasil kali!dalam $\cat{SET}$}%
 \index{SET@$\cat{SET}$!hasil kali dalam}%
hasil kali suatu keluarga terindeks himpunan ada dan sesungguhnya merupakan hasil kali Kartesius himpunan-himpunan
tersebut beserta proyeksi koordinat kanonik.
\end{exam}

\begin{exam}\label{C015544} Jika $\bigl(V_\lambda\bigr)_{\lambda \in \Lambda}$ adalah suatu keluarga terindeks
ruang-ruang vektor, kita menjadikan hasil kali Kartesius $\prod V_\lambda$ suatu ruang
vektor sebagai berikut. Definisikan penjumlahan dan perkalian skalar secara titik demi titik: untuk $f$, $g \in
\prod V_\lambda$ dan $\alpha \in \K$,
  \[ (f + g)(\lambda) := f(\lambda) + g(\lambda)  \]
dan
  \[ (\alpha f)(\lambda) := \alpha f(\lambda) \,. \]
Hal ini menjadikan $\prod V_\lambda$ suatu ruang vektor, yang kadang-kadang disebut
 \index{langsung!hasil kali}%
 \index{hasil kali!langsung}%
\df{hasil kali langsung} dari ruang-ruang $V_\lambda$, dan ruang ini beserta proyeksi
koordinat kanonik (yang tentu saja merupakan pemetaan linear) adalah suatu
 \index{hasil kali!dalam $\cat{VEC}$}%
 \index{VEC@$\cat{VEC}$!hasil kali dalam}%
 \index{vektor!ruang!hasil kali}%
hasil kali dalam kategori $\cat{VEC}$.
\end{exam}

\begin{exam}\label{C023411} Misalkan $V$ dan $W$ adalah ruang-ruang linear bernorma. Pada hasil kali Kartesius
$V \times W$, definisikan
  \index{<unopn@$\norm{\hphantom{x}}_2$ (norma Euklides atau norma-$2$)}%
  \index{<unopn@$\norm{\hphantom{x}}_1$ (norma-$1$)}%
  \index{<unopn@$\norm{\hphantom{x}}_u$ (norma seragam)}%
 \begin{align*}
    \norm{(x,y)}_2 &= \sqrt{\norm x^2 + \norm y^2},  \\
    \norm{(x,y)}_1 &= \norm x + \norm y,  \\
   \intertext{dan}
    \norm{(x,y)}_u &= \max\{\norm x, \norm y\}\,.
 \end{align*}
Yang pertama adalah
 \index{Euklides!norma!pada hasil kali dua ruang linear bernorma}%
 \index{norma!Euklides!pada hasil kali ruang-ruang linear bernorma}%
\df{norma Euklides} (atau
 \index{dua@norma-$2$}%
 \index{norma!$2$- (pada hasil kali ruang-ruang linear bernorma)}%
\df{norma-$2$}) pada $V \times W$, yang kedua adalah
 \index{satu@norma-$1$!pada hasil kali ruang-ruang linear bernorma}%
 \index{norma!$1$- (pada hasil kali ruang-ruang linear bernorma)}%
\df{norma-$1$}, dan yang terakhir adalah
 \index{seragam!norma!pada hasil kali ruang-ruang linear bernorma}%
 \index{norma!seragam!pada hasil kali ruang-ruang linear bernorma}%
\df{norma seragam}. Memverifikasi bahwa ketiganya merupakan norma pada $V \times W$ sangat mirip dengan
argumen yang diperlukan dalam Contoh~\ref{C023125}, \ref{C023127}, dan~\ref{C023129}. Fakta bahwa
ketiganya merupakan norma-norma ekuivalen merupakan konsekuensi dari pertidaksamaan berikut dan
Proposisi~\ref{C023184}.
   \[ \norm x + \norm y \le \sqrt 2 \sqrt{\norm x^2 + \norm y^2} \le 2 \max\{\norm x,\norm y\} \le 2(\norm x + \norm y) \]
\end{exam}

\begin{conv}\label{C023414} Dalam contoh sebelumnya, kita mendefinisikan tiga norma (ekuivalen) pada
ruang hasil kali $V \times W$. Kita akan mengambil yang kedua, $\norm{\;\;}_1$, sebagai
 \index{konvensi!pemilihan norma hasil kali}%
 \index{hasil kali!norma}%
 \index{norma!hasil kali}%
\df{norma hasil kali} pada $V \times W$. Jadi, kapan pun $V$ dan $W$ merupakan ruang-ruang linear bernorma,
kecuali dinyatakan sebaliknya, kita akan memandang hasil kali $V \times W$ sebagai ruang
linear bernorma terhadap norma ini. Biasanya kita cukup menulis $\norm{(x,y)}$, bukan
$\norm{(x,y)}_1$. Hasil kali ruang-ruang linear bernorma $V$ dan $W$ biasanya dinotasikan
dengan $V \oplus W$ dan disebut
 \index{hasil kali!ruang linear bernorma}%
 \index{bernorma!ruang linear!hasil kali}%
 \index{<binopdirectsum@$\oplus$ (jumlah langsung)}%
 \index{langsung!jumlah!ruang linear bernorma}%
 \index{bernorma!ruang linear!jumlah langsung}%
\df{jumlah langsung} dari $V$ dan~$W$. (Dalam kasus khusus $V = W = \R$, metrik apakah yang
diinduksi oleh norma hasil kali pada~$\R^2$?)
\end{conv}

\begin{exam}\label{exam_prod_nls} Tunjukkan bahwa hasil kali $V \oplus W$ dari ruang-ruang linear bernorma memang merupakan suatu hasil kali
 \index{hasil kali!dalam $\cat{NLS_{\boldsymbol\infty}}$}%
 \index{NLSinf@$\cat{NLS_{\boldsymbol\infty}}$!hasil kali dalam}%
dalam kategori $\cat{NLS_{\boldsymbol\infty}}$ yang terdiri atas ruang-ruang linear bernorma dan pemetaan-pemetaan
linear terbatas.
\end{exam}

\begin{prop} Misalkan $\bigl((x_n,y_n)\bigr)$ adalah barisan dalam jumlah langsung $V \oplus W$ dari dua ruang linear
bernorma. Buktikan bahwa $\bigl((x_n,y_n)\bigr)$ konvergen ke titik $(a,b)$ dalam $V \oplus W$ jika dan hanya
jika $x_n \sto a$ dalam $V$ dan $y_n \sto b$ dalam~$W$.
\end{prop}

\begin{prop}\label{C023427} Penjumlahan merupakan operasi kontinu pada suatu ruang linear bernorma~$V$.
 \index{kekontinuan!penjumlahan}%
 \index{penjumlahan!kekontinuan}%
Artinya, pemetaan
   \[ A \colon V \oplus V \sto V \colon (x,y) \mapsto x + y \]
bersifat kontinu.
\end{prop}

\begin{exer} Berikan bukti yang \emph{sangat} singkat (tanpa $\epsilon$, $\delta$, ataupun himpunan terbuka) bahwa jika
$(x_n)$ dan $(y_n)$ merupakan barisan-barisan dalam suatu ruang linear bernorma yang masing-masing konvergen ke $a$ dan $b$,
maka $x_n + y_n \sto a + b$.
\end{exer}

\begin{prop}\label{C023434} Perkalian skalar merupakan operasi kontinu pada suatu ruang
 \index{kekontinuan!perkalian skalar}%
 \index{skalar!perkalian!kekontinuan}%
linear bernorma~$V$, dalam arti bahwa pemetaan
   \[ S \colon \K \times V \sto V \colon  (\alpha,x) \mapsto \alpha x \]
bersifat kontinu.
\end{prop}

\begin{prop} Jika $B$ dan $C$ adalah subhimpunan-subhimpunan dari suatu ruang linear bernorma dan $\alpha$ adalah skalar, maka
 \begin{enumerate}
   \item $\clo{\alpha B} = \alpha \clo B$; dan
   \item $\clo B + \clo C  \subseteq  \clo{B + C}$.
 \end{enumerate}
\end{prop}

\begin{exam} Jika $B$ dan $C$ adalah subhimpunan-subhimpunan tertutup dari suatu ruang linear bernorma, tidak
selalu berlaku bahwa $B + C$ tertutup (dan karena itu $\clo B + \clo C$ dan $\clo{B +
C}$ tidak harus sama).
\end{exam}

\begin{prop} Misalkan $C_1$ dan $C_2$ adalah subhimpunan-subhimpunan kompak dari suatu ruang linear bernorma~$V$. Maka
   \begin{enumerate}
     \item $C_1 \times C_2$ adalah subhimpunan kompak dari $V \times V$, dan
     \item $C_1 + C_2$ adalah subhimpunan kompak dari~$V$.
   \end{enumerate}
\end{prop}

\begin{exam}\label{C023441} Misalkan $X$ suatu ruang topologis. Maka keluarga
$\fml C(X)$ dari semua fungsi kontinu bernilai skalar pada~$X$ merupakan suatu
 \index{C@$\fml C(X)$!sebagai ruang vektor}%
ruang vektor terhadap operasi penjumlahan dan perkalian skalar biasa yang didefinisikan titik demi titik.
\end{exam}

\begin{exam}\label{C023443} Misalkan $X$ suatu ruang topologis. Kita menotasikan
 \index{C@$\fml C_b(X)$!fungsi kontinu terbatas pada $X$}%
dengan $\fml C_b(X)$ keluarga semua fungsi kontinu terbatas bernilai skalar pada~$X$. Keluarga ini merupakan
 \index{C@$\fml C_b(X)$!sebagai ruang linear bernorma}%
 \index{bernorma!ruang linear!$\fml C_b(X)$ sebagai suatu}%
ruang linear bernorma terhadap norma seragam. Bahkan, ruang ini merupakan subruang dari~$\fml B(X)$
(lihat Contoh~\ref{C023131}).
\end{exam}

\begin{exam}\label{C023464} Jika $S$ adalah suatu himpunan dan $a \in S$, maka fungsional evaluasi
     \[ E_a\colon \fml B(S) \sto \K\colon f \mapsto f(a) \]
merupakan fungsional linear terbatas pada ruang $\fml B(S)$ dari semua fungsi terbatas bernilai
skalar pada~$S$. Jika $S$ kebetulan merupakan ruang topologis, kita juga dapat memandang $E_a$
sebagai anggota ruang dual dari $\fml C_b(S)$. (Dalam kedua kasus tersebut, berapakah~$\norm{E_a}$?)
\end{exam}

\begin{defn}\label{def_norm_alg} Misalkan $A$ adalah suatu aljabar yang padanya telah didefinisikan sebuah norma. Andaikan bahwa sebagai tambahan,
 \index{submultiplikatif}%
sifat \df{submultiplikatif}
  \[ \norm{xy} \le \norm x\,\norm y \]
dipenuhi untuk semua $x$, $y \in A$. Maka $A$ merupakan suatu
 \index{bernorma!aljabar}%
 \index{aljabar!bernorma}%
\df{aljabar bernorma}. Kita menetapkan satu syarat tambahan: jika aljabar $A$ beridentitas, maka
   \[ \norm{\vc 1} = 1\,. \]
Dalam kasus ini, $A$ merupakan
 \index{beridentitas!aljabar bernorma}%
 \index{bernorma!aljabar!beridentitas}%
 \index{aljabar!bernorma beridentitas}%
\df{aljabar bernorma beridentitas}.
\end{defn}

\begin{exam}\label{C023448} Dalam Contoh~\ref{C023131}, telah kita tunjukkan bahwa keluarga
 \index{B@$\fml B(S)$!sebagai aljabar bernorma}%
 \index{bernorma!aljabar!$\fml B(S)$ sebagai suatu}%
$\fml B(S)$ dari fungsi-fungsi terbatas bernilai skalar pada suatu himpunan $S$ merupakan ruang linear bernorma terhadap
norma seragam. Keluarga ini juga merupakan aljabar bernorma komutatif beridentitas.
\end{exam}

\begin{prop} Perkalian merupakan operasi kontinu pada suatu aljabar bernorma~$V$, dalam arti
bahwa pemetaan
   \[ M \colon V \times V \sto V \colon (x,y) \mapsto  xy \]
bersifat kontinu.
\end{prop}

\begin{proof}[\emph{Petunjuk pembuktian}] Jika Anda dapat membuktikan bahwa perkalian pada bilangan real
bersifat kontinu, hampir pasti Anda juga dapat membuktikan bahwa perkalian bersifat kontinu pada sebarang aljabar
bernorma. \ns
\end{proof}

\begin{exam}\label{C023454} Keluarga $\fml C_b(X)$ dari fungsi-fungsi kontinu terbatas bernilai skalar
 \index{C@$\fml C_b(X)$!sebagai aljabar bernorma}%
 \index{bernorma!aljabar!$\fml C_b(X)$ sebagai suatu}%
pada suatu ruang topologis~$X$, terhadap operasi-operasi biasa yang didefinisikan titik demi titik dan
norma seragam, merupakan aljabar bernorma komutatif beridentitas. (Lihat Contoh~\ref{C023443}.)
\end{exam}

\begin{exam}\label{C023456} Keluarga $\ofml B(V)$ dari semua operator (linear terbatas)
 \index{B@$\ofml B(V)$!sebagai aljabar bernorma beridentitas}%
 \index{bernorma!aljabar!$\ofml B(V)$ sebagai suatu}%
pada suatu ruang linear bernorma bukan nol $V$ merupakan aljabar bernorma beridentitas. (Lihat~\ref{C023228} dan~\ref{C023228a}.)
\end{exam}

\section{Koproduk Ruang Linear Bernorma}\label{C0156}
Koproduk serupa dengan hasil kali---kecuali bahwa semua anak panah dibalik.

\begin{defn}\label{C015611} Misalkan $A_1$ dan $A_2$ adalah objek-objek dalam suatu kategori~$\cat C$. Tripel
$(Q,\iota_1,\iota_2)$, dengan $Q$ suatu objek dan $\iota_k\colon A_k \sto Q$ morfisme untuk $k=1,2$, disebut
 \index{koproduk}%
\df{koproduk} dari $A_1$ dan $A_2$ jika untuk setiap objek $B$ dan setiap pasangan morfisme
$f_k\colon A_k \sto B$ (k = 1,2) terdapat morfisme tunggal $f\colon Q \sto B$ sedemikian
sehingga $f_k = f \circ \iota_k$ untuk $k = 1,2$.
\end{defn}

\begin{prop} Dalam sebarang kategori, koproduk (jika ada) bersifat
 \index{koproduk!ketunggalan}%
unik secara esensial.
\end{prop}

\begin{defn}\label{C0006127} Misalkan $A$ dan $B$ adalah himpunan-himpunan. Ketika $A$ dan $B$ saling lepas, kita
sering menggunakan notasi
 \index{<binopuniond@$\uplus$, $\biguplus$ (gabungan saling lepas)}%
$A \uplus B$ sebagai pengganti $A \cup B$ (untuk menekankan bahwa $A$ dan~$B$ saling lepas). Ketika
$C = A \uplus B$, kita mengatakan bahwa $C$ merupakan
 \index{saling lepas!gabungan}%
 \index{gabungan!saling lepas}%
\df{gabungan saling lepas} dari $A$ dan~$B$. Demikian pula, kita sering memilih untuk menulis gabungan dari
suatu keluarga himpunan $\sfml A$ yang \emph{saling lepas berpasangan} sebagai $\biguplus\sfml A$. Notasi
$\biguplus_{\lambda \in \Lambda} A_\lambda$ juga dapat digunakan untuk menyatakan gabungan dari suatu
keluarga terindeks \emph{saling lepas berpasangan} $\{A_\lambda\colon  \lambda \in \Lambda\}$ dari
himpunan-himpunan. Ketika $C = \biguplus\sfml A$ atau $C = \biguplus_{\lambda \in \Lambda} A_\lambda$, kita
mengatakan bahwa $C$ merupakan \df{gabungan saling lepas} dari keluarga yang bersesuaian.
\end{defn}

\begin{defn} Dalam definisi sebelumnya, kita memperkenalkan notasi $A \uplus B$ untuk
gabungan dua himpunan saling lepas $A$ dan~$B$. Sekarang kita memperluas notasi ini dan memperbolehkan
diri kita mengambil \emph{gabungan saling lepas} dari himpunan $A$ dan $B$ sekalipun keduanya tidak
saling lepas. Kita ``membuat keduanya saling lepas'' dengan mengidentifikasi (dengan cara yang jelas) himpunan $A$ dengan
himpunan $A' = \{(a,1)\colon a \in A\}$ dan $B$ dengan himpunan $B' = \{(b,2)\colon b \in B\}$. Kemudian kita
menotasikan gabungan $A'$ dan $B'$, yang memang saling lepas, dengan $A \uplus B$ dan menyebutnya
 \index{saling lepas!gabungan}%
 \index{gabungan!saling lepas}%
 \index{<binopuniond@$\uplus$, $\biguplus$ (gabungan saling lepas)}%
\df{gabungan saling lepas} dari $A$ dan~$B$.

Secara umum, jika $\bigl(A_\lambda\bigr)_{\lambda \in \Lambda}$ adalah suatu keluarga terindeks
himpunan,
 \index{saling lepas!gabungan}%
 \index{gabungan!saling lepas}%
\df{gabungan saling lepas}nya, \smash[b]{$\biguplus\limits_{\lambda \in \Lambda}A_\lambda$}, didefinisikan sebagai
$\{(a,\lambda)\colon \lambda \in \Lambda, a \in A_\lambda\}$.
\end{defn}

\begin{exam}\label{C015617} Dalam kategori $\cat{SET}$,
 \index{koproduk!dalam $\cat{SET}$}%
 \index{SET@$\cat{SET}$!koproduk dalam}%
koproduk dua himpunan $A_1$ dan $A_2$ ada dan merupakan gabungan saling lepas $A_1 \uplus A_2$
beserta pemetaan inklusi yang jelas $\iota_k\colon A_k  \sto A_1 \uplus A_2$ ($k = 1,2$).
\end{exam}

Menarik untuk mengamati bahwa meskipun hasil kali dan koproduk suatu koleksi hingga
objek dalam kategori $\cat{SET}$ sangat berbeda, dalam kategori $\cat{VEC}$ yang lebih kompleks
keduanya ternyata merupakan hal yang persis sama.

\begin{exam} Jika $V_1$ dan $V_2$ adalah ruang-ruang vektor, jadikan hasil kali Kartesius $V_1 \times V_2$
suatu ruang vektor seperti dalam Contoh~\ref{C015521}. Ruang ini beserta
injeksi-injeksi yang jelas merupakan suatu
 \index{koproduk!dalam $\cat{VEC}$}%
 \index{VEC@$\cat{VEC}$!koproduk dalam}%
 \index{vektor!ruang!koproduk}%
koproduk dalam kategori $\cat{VEC}$.
\end{exam}

Sering kali mungkin dan dikehendaki untuk mengambil koproduk dari sebarang keluarga objek
dalam suatu kategori. Berikut adalah perumuman Definisi~\ref{C015611}.

\begin{defn}\label{C015624} Misalkan $\bigl(A_\lambda\bigr)_{\lambda \in \Lambda}$ adalah suatu keluarga terindeks
objek-objek dalam kategori~$\cat C$. Kita mengatakan bahwa objek $C$ beserta suatu
keluarga terindeks $\bigl(\iota_\lambda\bigr)_{\lambda \in \Lambda}$ dari morfisme-morfisme
$\iota_\lambda\colon A_\lambda \sto C$ merupakan
 \index{koproduk!dalam suatu kategori}%
\df{koproduk} dari objek-objek $A_\lambda$ jika untuk setiap objek $B$ dan setiap keluarga
terindeks $\bigl(f_\lambda\bigr)_{\lambda \in \Lambda}$ dari morfisme-morfisme $f_\lambda\colon
A_\lambda \sto B$ terdapat pemetaan tunggal $g \colon C \sto B$ sedemikian sehingga $f_\lambda = g
\circ \iota_\lambda$ untuk setiap $\lambda \in \Lambda$. Notasi biasa untuk
koproduk objek-objek $A_\lambda$ adalah
 \index{<<@$\coprod A_\lambda$ (koproduk)}%
$\coprod_{\lambda \in \Lambda} A_\lambda$.
\end{defn}

\begin{exam} Dalam kategori $\cat{SET}$,
 \index{koproduk!dalam $\cat{SET}$}%
 \index{SET@$\cat{SET}$!koproduk dalam}%
koproduk suatu keluarga terindeks himpunan ada dan merupakan gabungan saling lepas himpunan-himpunan tersebut.
\end{exam}

\begin{defn} Misalkan $S$ suatu himpunan dan $V$ suatu ruang vektor.
 \index{tumpuan!fungsi}%
\df{Tumpuan} suatu fungsi $f \colon S \sto V$ adalah $\{s \in S \colon f(s) \ne \vc 0\}$.
\end{defn}

\begin{exam}\label{X_dir_sum_VEC} Jika $\bigl(V_\lambda\bigr)_{\lambda \in \Lambda}$ adalah suatu keluarga terindeks
ruang-ruang vektor, kita menjadikan hasil kali Kartesius tersebut suatu ruang vektor seperti dalam Contoh~\ref{C015544}. Himpunan fungsi $f$
dalam $\prod V_\lambda$ yang memiliki tumpuan hingga (yaitu, yang bernilai tak nol hanya pada berhingga banyak unsur domainnya) jelas
merupakan subruang dari $\prod V_\lambda$. Subruang ini adalah
 \index{langsung!jumlah}%
 \index{jumlah!langsung}%
\df{jumlah langsung} dari ruang-ruang $V_\lambda$. Ruang ini dinotasikan dengan $\bigoplus_{\lambda \in
\Lambda} V_\lambda$. Ruang ini beserta pemetaan-pemetaan injektif yang jelas (yang
bersifat linear) merupakan suatu
 \index{koproduk!dalam $\cat{VEC}$}%
 \index{VEC@$\cat{VEC}$!koproduk dalam}%
 \index{vektor!ruang!koproduk}%
koproduk dalam kategori $\cat{VEC}$.
\end{exam}

\begin{exer} Apa saja koproduk dalam kategori $\cat{NLS_{\boldsymbol\infty}}$?
\end{exer}

\begin{exer} Tentukan hasil kali dan koproduk dua ruang $V$ dan $W$ dalam kategori
$\cat{NLS_{\vc 1}}$ yang terdiri atas ruang-ruang linear bernorma dan kontraksi-kontraksi linear.
\end{exer}




















\section{Jaring}\label{C0314}
Jaring (kadang-kadang disebut \emph{barisan tergeneralisasi}) berguna untuk mencirikan banyak
sifat ruang topologis umum. Dalam ruang semacam itu, jaring pada dasarnya digunakan
dengan cara yang sama seperti barisan digunakan dalam ruang metrik.

\begin{defn} Himpunan terpraurut yang setiap pasangan elemennya memiliki batas atas disebut
 \index{terarah!himpunan}%
 \index{himpunan!terarah}%
\df{himpunan terarah}.
\end{defn}

\begin{exam}\label{C031414} Jika $S$ suatu himpunan, maka $\fin S$ merupakan himpunan terarah terhadap
inklusi $\subseteq$.
\end{exam}

\begin{defn} Misalkan $S$ suatu himpunan dan $\Lambda$ suatu himpunan terarah. Pemetaan
$x\colon \Lambda \sto S$ disebut suatu
 \index{jaring}%
\df{jaring} di $S$ (atau suatu \emph{jaring anggota-anggota $S$}). Nilai $x$ di $\lambda \in
\Lambda$ biasanya ditulis $\sbsb x\lambda$ (alih-alih $x(\lambda)$), dan jaring $x$
sendiri sering dinotasikan dengan $\bigl(\sbsb x\lambda\bigr)_{\lambda \in \Lambda}$ atau cukup
$\bigl(\sbsb x\lambda\bigr)$.
\end{defn}

\begin{exam} Contoh jaring yang paling jelas ialah barisan. Barisan merupakan fungsi yang domainnya adalah himpunan (terpraurut) $\N$ dari bilangan asli.
\end{exam}

\begin{defn} Suatu jaring $x = \bigl(\sbsb x\lambda\bigr)$ dikatakan
 \index{pada akhirnya}%
\df{pada akhirnya} memiliki suatu sifat tertentu jika terdapat $\lambda_0 \in \Lambda$ sedemikian sehingga $\sbsb x\lambda$ memiliki sifat itu
apabila $\lambda \ge \lambda_0$; dan jaring itu
 \index{sering}%
\df{sering} memiliki sifat tersebut jika untuk setiap $\lambda_0 \in \Lambda$ terdapat $\lambda \in \Lambda$ sedemikian sehingga $\lambda \ge \lambda_0$
dan $\sbsb x\lambda$ memiliki sifat itu.
\end{defn}

\begin{defn} Suatu
 \index{lingkungan}%
\df{lingkungan} dari sebuah titik dalam ruang topologis adalah sebarang himpunan yang memuat suatu himpunan terbuka yang memuat titik tersebut.
\end{defn}

\begin{defn} Suatu jaring $x$ dalam ruang topologis $X$
 \index{kekonvergenan!jaring}%
 \index{jaring!kekonvergenan}%
\df{konvergen} ke titik $a \in X$ jika jaring itu pada akhirnya berada dalam setiap lingkungan~$a$. Dalam
hal ini $a$ adalah
 \index{limit!jaring}%
 \index{jaring!limit}%
\df{limit} dari $x$, dan kita menulis
   \[ x_{{}_{\sst \lambda}} \to_{\lambda \in \Lambda}a \]
atau cukup
 \index{<arrowto@$x_\lambda \sto a$ (kekonvergenan jaring)}%
$x_{{}_{\sst \lambda}} \sto a$. Jika limit bersifat unik, kita juga dapat menggunakan notasi
$\lim_{\lambda \in \Lambda} \sbsb x\lambda = a$ atau, secara lebih sederhana, $\lim \sbsb x\lambda =
a$.

Titik $a$ di $X$ merupakan
 \index{titik gugus!jaring}%
 \index{jaring!titik gugus}%
\df{titik gugus} dari jaring $x$ jika $x$ sering berada dalam setiap lingkungan~$a$.
\end{defn}

\begin{defn} Misalkan $(X,\sfml T)$ suatu ruang topologis. Subkeluarga $\sfml S \subseteq
\sfml T$ merupakan suatu
 \index{subbasis}%
\df{subbasis} untuk topologi $\sfml T$ jika keluarga semua irisan hingga dari
anggota-anggota $\sfml S$ merupakan basis untuk~$\sfml T$.
\end{defn}

\begin{exam}\label{C031429} Misalkan $X$ suatu ruang topologis dan $a \in X$. Suatu
 \index{lingkungan!basis untuk}%
 \index{basis!lingkungan suatu titik}%
\df{basis} bagi keluarga lingkungan titik $a$ adalah keluarga $\sfml B_a$ yang terdiri atas
lingkungan-lingkungan $a$ dengan sifat bahwa setiap lingkungan $a$ memuat sekurang-kurangnya
satu anggota~$\sfml B_a$. Secara umum, terdapat banyak pilihan basis lingkungan di
suatu titik. Setelah basis semacam itu dipilih, kita menyebut anggota-anggotanya sebagai
 \index{dasar!lingkungan}%
 \index{lingkungan!dasar}%
\emph{lingkungan dasar} dari~$a$.

Misalkan $\sfml B_a$ suatu basis bagi keluarga lingkungan di $a$. Urutkan $\sfml B_a$ berdasarkan
relasi memuat (kebalikan dari inklusi); yakni, untuk $U$, $V \in \sfml B_a$, tetapkan
  \[ U \preceq V \text{ jika dan hanya jika } U \supseteq V\,. \]
Dengan demikian, $\sfml B_a$ menjadi himpunan terarah. Sekarang (dengan menggunakan aksioma pilihan) pilih satu
elemen $x_{{}_U}$ dari setiap himpunan $U \in \sfml B_a$. Maka $(\sbsb xU)$ merupakan jaring di $X$
dan $\sbsb xU \sto a$.
\end{exam}

Dua proposisi berikut menjamin bahwa agar suatu jaring konvergen ke titik $a$ dalam
ruang topologis, cukup jika jaring itu pada akhirnya berada dalam setiap lingkungan dasar (atau bahkan
lingkungan subdasar) dari~$a$.

\begin{prop}\label{C031431} Misalkan $\sfml B$ suatu basis bagi topologi pada ruang topologis
$X$ dan $a$ suatu titik di~$X$. Suatu jaring $(x_\lambda)$ konvergen ke $a$ jika dan hanya
jika jaring itu pada akhirnya berada dalam setiap lingkungan~$a$ yang termasuk dalam~$\sfml B$.
\end{prop}

\begin{prop}\label{C031434} Misalkan $\sfml S$ suatu subbasis bagi topologi pada ruang topologis
$X$ dan $a$ suatu titik di~$X$. Suatu jaring $(x_\lambda)$ konvergen ke $a$ jika dan hanya jika
jaring itu pada akhirnya berada dalam setiap lingkungan~$a$ yang termasuk dalam~$\sfml S$.
\end{prop}

\begin{exam}\label{C031437} Misalkan $J = [a,b]$ suatu interval tetap pada garis real,
$\vc x = (x_0,x_1,\dots,x_n)$ suatu $(n+1)$-tupel titik-titik di $J$, dan $\vc t =
(t_1,\dots,t_n)$ suatu $n$-tupel. Pasangan $(\vc x;\vc t)$ merupakan suatu
 \index{partisi!dengan pilihan}%
 \index{pilihan}%
\df{partisi dengan pilihan} dari interval $J$ jika:
 \begin{enumerate}
  \item $x_{k-1} < x_k$ untuk $1 \le k \le n$;
  \item $t_k \in [x_{k-1},x_k]$ untuk $1 \le k \le n$;
  \item $x_0 = a$; dan
  \item $x_n = b$.
 \end{enumerate}
Gagasannya ialah bahwa $\vc x$ \emph{mempartisi} interval tersebut menjadi subinterval-subinterval dan $t_k$ adalah
titik yang \emph{dipilih} dari subinterval $\text{ke-}k$. Jika $\vc x =
(x_0,x_1,\dots,x_n)$, maka $\{\vc x\}$ menyatakan \emph{himpunan} $\{x_0,x_1,\dots,x_n\}$.
Misalkan $P = (\vc x;\vc s)$ dan $Q = (\vc y;\vc t)$ partisi-partisi (dengan pilihan) dari
interval~$J$. Kita menulis $P \preccurlyeq Q$ dan mengatakan bahwa $Q$ merupakan suatu
 \index{penghalusan}%
\df{penghalusan} dari $P$ jika $\{\vc x\} \subseteq \{\vc y\}$. Terhadap relasi
$\preccurlyeq$, keluarga $\sfml P$ dari partisi-partisi dengan pilihan pada $J$ merupakan himpunan terarah
(tetapi \emph{bukan} terurut parsial).

Sekarang andaikan $f$ suatu fungsi terbatas pada interval $J$ dan $P = (\vc x;\vc t)$ suatu
partisi $J$ menjadi $n$ subinterval. Misalkan $\Delta x_k := x_k - x_{k-1}$ untuk $1 \le k
\le n$. Lalu definisikan
   \[ S_f(P) := \sum_{k=1}^n f(t_k)\,\Delta x_k\,. \]
Setiap jumlah $S_f(P)$ semacam itu merupakan suatu
 \index{Riemann!jumlah}%
 \index{jumlah!Riemann}%
\df{jumlah Riemann} dari $f$ pada~$J$. Perhatikan bahwa karena $\sfml P$ merupakan himpunan terarah, $S_f$ adalah suatu
jaring bilangan real. Jika jaring $S_f$ dari jumlah-jumlah Riemann konvergen, kita mengatakan bahwa fungsi
$f$
 \index{Riemann!dapat diintegralkan}%
 \index{dapat diintegralkan!secara Riemann}%
\df{dapat diintegralkan secara Riemann}. Limit jaring $S_f$ (jika ada) merupakan
 \index{Riemann!integral}%
 \index{integral!Riemann}%
\df{integral Riemann} dari $f$ pada interval~$J$ dan dinotasikan dengan $\int_a^b f(x)\,dx$.
\end{exam}

\begin{exam}\label{exam_net_indiscrete} Dalam ruang topologis $X$ dengan topologi
indiskret (yang himpunan terbukanya hanya $X$ dan himpunan kosong), setiap jaring di ruang tersebut
konvergen ke setiap titik di ruang itu.
\end{exam}

\begin{exam} Dalam ruang topologis $X$ dengan topologi diskret (yang setiap subhimpunan $X$ bersifat terbuka), suatu jaring
di ruang tersebut konvergen jika dan hanya jika jaring itu pada akhirnya konstan.
 \end{exam}

\begin{defn} Suatu ruang topologis
 \index{Hausdorff}%
 \index{pemisahan!aksioma!Hausdorff}%
 \index{topologis!ruang!Hausdorff}%
 \index{topologi!Hausdorff}%
\df{Hausdorff} adalah ruang yang setiap pasangan titik berbedanya dapat
dipisahkan oleh himpunan-himpunan terbuka; yakni, untuk setiap pasangan titik berbeda $x$ dan $y$ di ruang tersebut, terdapat
lingkungan terbuka dari $x$ dan $y$ yang saling lepas.
\end{defn}

Seperti ditunjukkan oleh Contoh~\ref{exam_net_indiscrete}, jaring dalam ruang topologis umum tidak harus memiliki
limit yang unik. Namun, suatu ruang tidak memerlukan asumsi yang terlalu membatasi untuk menyingkirkan
patologi ini. Agar limit bersifat unik, cukup disyaratkan bahwa ruang itu Hausdorff.
Menariknya, ternyata syarat ini juga perlu.

\begin{prop}\label{C031441} Jika $X$ suatu ruang topologis, maka pernyataan-pernyataan berikut ekuivalen:
 \begin{enumerate}
  \item $X$ bersifat Hausdorff.
  \item Tidak ada jaring di $X$ yang dapat konvergen ke lebih dari satu titik.
  \item Diagonal di $X \times X$ bersifat tertutup.
 \end{enumerate}
\end{prop}
(Adapun
 \index{diagonal}%
\df{diagonal} dari suatu himpunan $S$ adalah $\{(s,s) \colon s \in S\} \subseteq S \times S$.)

\begin{exam} Ingat kembali dari kalkulus dasar bahwa
 \index{parsial!jumlah}%
 \index{jumlah!parsial}%
 \index{deret tak hingga!jumlah parsial}%
 \index{deret!tak hingga}%
\df{jumlah parsial} $\text{ke-}n$ dari suatu deret tak hingga $\sum_{k=1}^\infty x_k$ didefinisikan sebagai $s_n =
\sum_{k=1}^n x_k$. Deret tersebut dikatakan
 \index{kekonvergenan!deret tak hingga}%
 \index{deret tak hingga!kekonvergenan}%
\df{konvergen} jika barisan $(s_n)$ dari jumlah-jumlah parsialnya konvergen dan, jika deret tersebut
konvergen,
 \index{jumlah!deret tak hingga}%
 \index{deret tak hingga!jumlah}%
 \index{deret!jumlah}%
\emph{jumlah}nya adalah limit barisan $(s_n)$ dari jumlah-jumlah parsial. Jadi, sebagai contoh,
deret tak hingga $\sum_{k=1}^\infty 2^{-k}$ konvergen dan jumlahnya adalah~$1$.
\end{exam}

Gagasan \emph{menjumlahkan} suatu deret tak hingga sangat bergantung pada fakta bahwa
jumlah-jumlah parsial deret tersebut terurut linear oleh inklusi. ``Jumlah-jumlah parsial'' dari suatu
jaring bilangan tidak harus terurut parsial, sehingga kita memerlukan gagasan baru tentang
``penjumlahan'', yang kadang-kadang disebut \emph{penjumlahan tak berurutan}.

\begin{notn} Jika $A$ suatu himpunan, kita notasikan
 \index{fin@$\fin A$ (keluarga subhimpunan hingga dari suatu himpunan $A$)}%
dengan $\fin A$ keluarga semua subhimpunan hingga dari~$A$. (Perhatikan bahwa keluarga ini merupakan himpunan terarah terhadap
relasi~$\subseteq$.)
\end{notn}

\begin{defn}\label{def_summable_R} Misalkan $A \subseteq \R$. Untuk setiap $F \in \fin A$, definisikan $S_F$ sebagai jumlah
bilangan-bilangan dalam~$F$, yakni, $S_F = \sum F$. Maka $S = \bigl(S_F\bigr)_{F \in \fin A}$
merupakan jaring di~$\R$. Jika jaring ini konvergen, himpunan bilangan $A$ dikatakan
 \index{dapat dijumlahkan!himpunan bilangan real}%
\df{dapat dijumlahkan}; limit jaring tersebut merupakan
 \index{jumlah!himpunan yang dapat dijumlahkan}%
\df{jumlah} dari $A$ dan dinotasikan dengan $\sum A$.
\end{defn}

\begin{defn} Suatu jaring $(x_\lambda)$ dari bilangan real disebut
 \index{menaik!jaring dalam $\R$}%
 \index{jaring!menaik}%
\df{menaik} jika $\lambda \le \mu$ mengakibatkan $x_\lambda \le x_\mu$. Suatu jaring bilangan real (atau kompleks)
disebut
 \index{terbatas!jaring dalam $\K$}%
 \index{jaring!terbatas}%
\df{terbatas} jika citranya terbatas.
\end{defn}

\begin{exam}\label{C031454} Misalkan $f(n) = 2^{-n}$ untuk setiap $n \in \N$, dan definisikan $x \colon
\fin(\N) \sto \R$ dengan $x(A) = \sum_{n \in A}f(n)$.
 \begin{enumerate}
  \item Jaring $x$ menaik.
  \item Jaring $x$ terbatas.
  \item Jaring $x$ konvergen ke 1.
  \item Himpunan bilangan $\{\frac12,\frac14,\frac18, \frac1{16},\dots\}$ dapat dijumlahkan dan
jumlahnya adalah~$1$.
  \end{enumerate}
\end{exam}

\begin{exam} Dalam ruang metrik, setiap barisan konvergen bersifat terbatas. Berikan contoh yang menunjukkan
bahwa suatu jaring konvergen dalam ruang metrik (bahkan di~$\R$) tidak harus terbatas.
\end{exam}

Dari kalkulus dasar kita telah mengenal bahwa barisan menaik yang terbatas di $\R$
bersifat konvergen. Contoh di atas merupakan kasus khusus dari pengamatan yang lebih umum: jaring
menaik yang terbatas di $\R$ bersifat konvergen.

\begin{prop}\label{C031461} Setiap jaring menaik yang terbatas
$\bigl(x_{{}_{\sst \lambda}}\bigr)_{\lambda \in \Lambda}$ dari bilangan real bersifat konvergen dan
  \[ \lim x_{{}_{\sst \lambda}} = \sup\{x_{{}_{\sst \lambda}}\colon \lambda \in \Lambda\}\,. \]
\end{prop}

\begin{exam} Misalkan $(x_k)$ suatu barisan bilangan real positif yang berbeda-beda dan $A = \{x_k \colon k \in \N\}$.
Maka $\sum A$ (sebagaimana didefinisikan di atas) sama dengan jumlah deret $\sum_{k=1}^\infty
x_k$.
\end{exam}

\begin{exam} Himpunan $\bigl\{\frac1k \colon k \in \N \bigr\}$ tidak dapat dijumlahkan dan deret tak hingga
$\sum_{k=1}^\infty \frac1k$ tidak konvergen.
\end{exam}

\begin{exam} Himpunan $\bigl\{(-1)^k \frac1k\colon k \in \N \bigr\}$ tidak dapat dijumlahkan, tetapi
deret tak hingga $\sum_{k=1}^\infty (-1)^k \frac1k$ memang konvergen.
\end{exam}

\begin{prop} Setiap subhimpunan bilangan real yang dapat dijumlahkan bersifat terhitung.
\end{prop}

\begin{proof}[\emph{Petunjuk pembuktian.}] Jika $A$ suatu himpunan bilangan real yang dapat dijumlahkan dan $n \in \N$,
berapa banyak anggota $A$ yang dapat lebih besar daripada~$1/n$? \ns
\end{proof}

Dalam Proposisi~\ref{C031441}, kita telah melihat satu contoh cara jaring
mencirikan sifat topologis: suatu ruang bersifat Hausdorff jika dan hanya jika setiap jaring
di ruang tersebut memiliki paling banyak satu limit. Berikut ini sifat topologis lain---kekontinuan di suatu titik---yang dapat dengan mudah dicirikan menggunakan
jaring.

\begin{prop}\label{C031471} Misalkan $X$ dan $Y$ ruang-ruang topologis. Suatu fungsi
 \index{kekontinuan!karakterisasi oleh jaring}%
$f\colon X \sto Y$ kontinu di titik $a$ dalam $X$ jika dan hanya jika $f\bigl(\sbsb x\lambda\bigr) \sto f(a)$ di $Y$
setiap kali $\bigl(\sbsb x\lambda\bigr)$ merupakan jaring yang konvergen ke $a$ di~$X$.
\end{prop}

Dalam ruang metrik, kita mencirikan tutupan dan interior himpunan menggunakan barisan. Untuk
ruang topologis umum, kita harus mengganti barisan dengan jaring.

Dalam mendefinisikan istilah ``interior'' dan ``tutupan'', kita memiliki pilihan serupa dengan pilihan ketika kita mendefinisikan ``selubung konveks'' dalam
\ref{smplx005def}. Terdapat definisi `abstrak' dan definisi `konstruktif'. Mungkin akan membantu jika Anda mengingat kembali
Proposisi~\ref{smplx006} dan paragraf yang mendahuluinya.

\begin{defn}\label{def_interior} Untuk suatu himpunan $A$ dalam ruang topologis $X$,
 \index{interior}%
\df{interior} himpunan tersebut adalah subhimpunan terbuka terbesar dari $X$ yang termuat dalam~$A$; yakni,
gabungan semua subhimpunan terbuka dari $X$ yang termuat dalam~$A$. Interior $A$
 \index{<unopur@$\intr A$ (interior dari~$A$)}%
dinotasikan dengan $\intr A$.
\end{defn}

\begin{prop} Definisi di atas bermakna. Lebih lanjut, suatu titik $a$ dalam ruang topologis $X$ berada dalam interior suatu subhimpunan
 \index{interior!karakterisasi oleh jaring}%
 \index{karakterisasi!oleh jaring!interior}%
 \index{jaring!karakterisasi!interior}%
$A$ dari $X$ jika dan hanya jika setiap jaring di $X$ yang konvergen ke $a$ pada akhirnya berada dalam~$A$.
\end{prop}

\begin{proof}[\emph{Petunjuk pembuktian}] Satu arah mudah. Untuk arah lainnya, gunakan Proposisi~\ref{C031431} dan jenis
jaring yang dikonstruksi dalam Contoh~\ref{C031429}. \ns
\end{proof}

\begin{defn}\label{def_closure} Untuk suatu himpunan $A$ dalam ruang topologis $X$,
 \index{tutupan}%
\df{tutupan} himpunan tersebut adalah subhimpunan tertutup terkecil dari $X$ yang memuat~$A$; yakni,
irisan semua subhimpunan tertutup dari $X$ yang memuat~$A$. Tutupan $A$
 \index{<unopu@$\clo A$ (tutupan~$A$)}%
dinotasikan dengan~$\clo A$.
\end{defn}

\begin{prop}\label{C031481} Definisi di atas bermakna. Lebih lanjut, suatu titik $b$ dalam ruang topologis $X$ termasuk dalam
 \index{tutupan!karakterisasi oleh jaring}%
 \index{karakterisasi!oleh jaring!tutupan}%
 \index{jaring!karakterisasi!tutupan}%
tutupan suatu subhimpunan $A$ dari~$X$ jika dan hanya jika terdapat suatu jaring di $A$ yang konvergen ke~$b$.
\end{prop}

\begin{cor} Suatu subhimpunan $A$ dari ruang topologis bersifat tertutup jika dan hanya jika limit setiap jaring konvergen di $A$ termasuk dalam~$A$.
\end{cor}



















\section{Teorema-Teorema Hahn--Banach}

Perlu dikatakan sejak awal bahwa merujuk pada \emph{teorema Hahn--Banach} seolah-olah hanya ada satu teorema sebenarnya agak menyesatkan.
Seorang penulis yang menyatakan sedang menggunakan \emph{teorema Hahn--Banach} mungkin menggunakan salah satu dari sekitar selusin versi
teorema tersebut atau salah satu korolarinya. Apa kesamaan berbagai versi itu? Pada dasarnya, versi-versi tersebut menjamin bahwa ruang
linear bernorma pada umumnya memiliki cukup banyak fungsional linear terbatas sehingga ruang dualnya menjadi alat yang berguna
untuk mempelajari ruang itu sendiri. Artinya, versi-versi tersebut menjamin adanya teori dualitas yang menarik dan produktif.
Di bawah ini kita membahas empat versi `teorema' tersebut dan tiga korolarinya. Kita akan mengikuti praktik yang lazim:
ketika dalam pembahasan selanjutnya kita menggunakan salah satu hasil ini, kita akan mengatakan bahwa kita menggunakan `teorema' \emph{Hahn--Banach}.

Seperti yang sering terjadi dalam matematika, pengaitan hasil-hasil ini dengan Hahn dan Banach agak ganjil. Tampaknya,
gagasan-gagasan dasarnya mula-mula berasal dari Eduard Helly. Untuk sejarah menarik tentang hasil-hasil `Hahn--Banach',
lihat~\cite{Hochstadt:1980}.

Sebagai catatan, dalam nama ``Banach'', kedua huruf ``a'' dilafalkan seperti huruf ``a'' dalam kata Indonesia \emph{bapak}, sedangkan ``ch''
kira-kira dilafalkan seperti \emph{kh}; tekanan jatuh pada suku kata pertama (BAH-nakh).

\begin{defn} Suatu fungsi bernilai real $f$ pada ruang vektor real $V$ merupakan suatu
 \index{sublinear}%
\df{fungsional sublinear} jika $f(x + y) \le f(x) + f(y)$ dan $f(\alpha x) = \alpha f(x)$
untuk semua $x$, $y \in V$ dan $\alpha \ge 0$.
\end{defn}

\begin{exam} Suatu norma pada ruang vektor real merupakan fungsional sublinear.
\end{exam}

Versi pertama \emph{teorema Hahn--Banach} kita murni merupakan aljabar linear dan hanya berlaku untuk ruang vektor \emph{real}.
Versi ini memungkinkan kita memperpanjang fungsional linear tertentu pada ruang semacam itu dari subruang ke seluruh ruang.

\begin{thm}[Teorema Hahn--Banach I]\label{HBTI} Misalkan $M$ suatu subruang dari ruang vektor real $V$ dan $p$
 \index{teorema Hahn--Banach}%
suatu fungsional sublinear pada~$V$. Jika $f$ suatu fungsional linear pada $M$ sedemikian sehingga $f \le p$ pada $M$, maka
$f$ memiliki perpanjangan $g$ ke seluruh $V$ sedemikian sehingga $g \le p$ pada~$V$.
\end{thm}

\begin{proof}[\emph{Petunjuk pembuktian}] Dalam teorema ini, notasi $h \le p$ berarti bahwa $\dom h \subseteq \dom p = V$
dan bahwa $h(x) \le p(x)$ untuk semua $x \in \dom h$. Dalam hal ini kita mengatakan bahwa $h$
 \index{didominasi}%
\df{didominasi} oleh~$p$. Kita menulis
 \index{<binrelle@$f \preccurlyeq g$ ($g$ merupakan perpanjangan~$f$)}%
$f \preccurlyeq g$ jika $g$ merupakan perpanjangan~$f$ (yakni,
jika $f(x) = g(x)$ untuk semua $x \in \dom f$).

Misalkan $\fml S$ keluarga semua fungsional linear bernilai real pada subruang-subruang $V$ yang merupakan perpanjangan~$f$ dan
didominasi oleh~$p$. Keluarga $\fml S$ terurut parsial oleh $\preccurlyeq$. Gunakan \emph{lema Zorn} untuk menghasilkan suatu
elemen maksimal $g$ dari~$\fml S$. Pembuktian selesai jika Anda dapat menunjukkan bahwa $\dom{g} = V$.

Untuk itu, andaikan demi kontradiksi bahwa terdapat vektor $c \in V$ yang tidak termasuk dalam~$D = \dom{g}$. Buktikan terlebih dahulu bahwa
   \begin{equation}\label{eq_HBTI}
      g(y) - p(y - c) \le - g(x) + p(x + c)
   \end{equation}
untuk semua $x$, $y \in D$. Simpulkan bahwa terdapat bilangan $\gamma$ yang terletak di antara supremum semua bilangan
di ruas kiri~\eqref{eq_HBTI} dan infimum semua bilangan di ruas kanan.

Misalkan $E$ rentang dari $D \cup \{c\}$. Definisikan $h$ pada $E$ dengan $h(x + \alpha c) = g(x) + \alpha\gamma$ untuk $x \in D$ dan
$\alpha \in \R$. Sekarang tunjukkan bahwa $h \in \fml S$. (Ambil $\alpha > 0$ dan tinjau secara terpisah $h(x + \alpha c)$ dan $h(x - \alpha c)$.) \ns
\end{proof}

Teorema berikutnya adalah versi Teorema~\ref{HBTI} untuk ruang dengan skalar kompleks. Teorema ini kadang-kadang disebut
 \index{teorema Bohnenblust--Sobczyk--Suhomlinov}%
\emph{teorema Bohnenblust--Sobczyk--Suhomlinov}.

\begin{thm}[Teorema Hahn--Banach II]\label{HBTII} Misalkan $V$ suatu ruang vektor kompleks, $p$ suatu seminorma pada~$V$, dan $M$ suatu subruang dari~$V$.
Jika $f\colon M \sto \C$ suatu fungsional linear sedemikian sehingga $\abs f \le p$, maka $f$ memiliki perpanjangan $g$ yang merupakan fungsional linear
pada seluruh~$V$ dan memenuhi $\abs{g} \le p$.
\end{thm}

\begin{proof}[\emph{Petunjuk pembuktian}] Hasil ini pada dasarnya merupakan korolar langsung dari versi real
\emph{teorema Hahn--Banach} yang diberikan di atas dalam~\ref{HBTI}. Misalkan $V_{\R}$ ruang vektor real yang terkait dengan~$V$ (vektornya sama,
tetapi hanya menggunakan $\R$ sebagai medan skalar). Perhatikan bahwa jika $u$ adalah bagian real dari $f$ (yakni,
 \index{R@$\Re(z)$ (bagian real dari $z \in \C$)}%
$u(x) = \Re(f(x))$ untuk setiap $x \in M$), maka $u$ merupakan fungsional $\R$-linear pada $M$ (yakni, $u(x + y) = u(x) + u(y)$ dan $u(\alpha x) = \alpha u(x)$
untuk semua $x$, $y \in M$ dan semua $\alpha \in \R$), serta $f(x) = u(x) - i\,u(ix)$. Sebaliknya, jika $u$ suatu fungsional $\R$-linear pada suatu ruang vektor,
maka $f(x) := u(x) - i\,u(ix)$ mendefinisikan suatu fungsional linear (bernilai kompleks) pada ruang tersebut. Perhatikan juga bahwa jika $u$ merupakan bagian real dari~$f$,
maka $\abs{u(x)} \le p$ untuk semua $x$ jika dan hanya jika $\abs{f(x)} \le p$ untuk semua~$x$. Dengan mengingat pengamatan-pengamatan ini, cukup terapkan~\ref{HBTI}
pada bagian real dari $f$ di~$V_\R$. \ns
\end{proof}

Berikutnya mungkin merupakan versi \emph{teorema Hahn--Banach} yang paling sering digunakan.

\begin{thm}[Teorema Hahn--Banach III]\label{HBTIII} Misalkan $M$ suatu subruang vektor dari ruang linear bernorma (real atau
 \index{teorema Hahn--Banach}%
kompleks) $V$. Maka setiap fungsional linear kontinu $f$ pada $M$ memiliki
perpanjangan $g$ ke seluruh $V$ sedemikian sehingga $\norm{g} = \norm f$.
\end{thm}

\begin{proof}[\emph{Petunjuk pembuktian}] Tinjau fungsi $p\colon V \sto \R \colon x \mapsto \norm f \norm x$. \ns
\end{proof}

Ketika membuktikan suatu teorema, kadang-kadang berguna untuk mempertimbangkan kemungkinan adanya pembuktian yang lebih ``alami'' atau bahkan lebih ``jelas''.
Sesekali kita memperoleh imbalan berupa penyederhanaan yang mencerahkan dan penjelasan yang lebih jernih atas materi yang rumit. Namun, kadang-kadang kita justru terkecoh.

\begin{exer} Kawan Anda, Fred R. Dimm, mengira bahwa ia telah menemukan pembuktian langsung yang sangat sederhana untuk
versi teorema Hahn--Banach yang diberikan dalam Teorema~\ref{HBTIII}: Suatu fungsional
linear kontinu $f$ yang didefinisikan pada subruang (linear) $M$ dari ruang linear bernorma~$V$ dapat
diperpanjang menjadi fungsional linear kontinu $g$ pada seluruh $V$ tanpa memperbesar
normanya. Menurutnya, kita hanya perlu mengambil $N$ sebagai sebarang komplemen ruang vektor
dari~$M$ dan mengambil $B$ sebagai suatu basis (Hamel) untuk $N$. Definisikan $g(e) = 0$ untuk setiap vektor $e
\in B$ dan $g = f$ pada~$M$. Kemudian perpanjang $g$ secara linear ke seluruh~$V$. Tentunya, kata Fred, $g$ linear berdasarkan konstruksinya dan normanya tidak bertambah karena kita
telah menjadikannya nol di tempat-tempat yang sebelumnya belum didefinisikan. Tunjukkan kepada Fred letak kekeliruannya.
\end{exer}

\begin{thm}[Teorema Hahn--Banach IV]\label{HBTIV} Misalkan $V$ suatu ruang linear bernorma (real atau kompleks),
 \index{teorema Hahn--Banach}%
$M$ suatu subruang vektor dari $V$, dan andaikan $y \in V$ sedemikian sehingga
 \index{dist@$\dist(y,M)$ (jarak antara $y$ dan $M$)}%
$d := \dist(y,M) > 0$. Maka terdapat fungsional $g \in V^*$ sedemikian sehingga $g^{\sto}(M) = \{0\}$,
$g(y) = d$, dan $\norm g = 1$.
\end{thm}

\begin{proof}[\emph{Petunjuk pembuktian}] Misalkan $N := \spn{(M \cup \{y\})}$ dan definisikan
$f\colon N \sto \K \colon x + \alpha y \mapsto \alpha d$ ($x \in M$, $\alpha \in \K$). \ns
\end{proof}

Kasus khusus yang paling berguna dari teorema di atas adalah ketika $M = \{\vc 0\}$ dan $y \ne \vc 0$.
Hal ini memberi tahu kita bahwa satu-satunya cara agar $f(x)$ lenyap untuk setiap $f \in V^*$ adalah jika $x$
merupakan vektor nol.

\begin{cor}\label{HBCorI} Misalkan $V$ suatu ruang linear bernorma dan $x \in V$. Jika $f(x) = 0$ untuk
setiap $f \in V^*$, maka $x = \vc 0$.
\end{cor}

\begin{defn} Suatu keluarga $\fml G(S)$ dari fungsi-fungsi bernilai skalar pada himpunan~$S$
 \index{pemisahan!pasangan titik oleh keluarga fungsi}%
\df{memisahkan} titik-titik $S$ jika untuk setiap pasangan titik berbeda $x$ dan $y$ di $S$, terdapat
fungsi $f \in \fml G(S)$ sedemikian sehingga $f(x) \ne f(y)$.
\end{defn}

\begin{cor}\label{HBCorII} Jika $V$ suatu ruang linear bernorma, maka $V^*$ memisahkan titik-titik~$V$.
\end{cor}

\begin{cor}\label{HBCorIII} Misalkan $V$ suatu ruang linear bernorma dan $\vc 0 \ne x \in V$. Maka terdapat
$f \in V^*$ sedemikian sehingga $\norm f = 1$ dan $f(x) = \norm x$.
\end{cor}

\begin{prop} Misalkan $V$ suatu ruang linear bernorma, $\{x_1,\dots,x_n\}$ suatu subhimpunan yang bebas linear
dari $V$, dan $\alpha_1, \dots, \alpha_n$ skalar-skalar. Maka terdapat $f$ dalam
$V^*$ sedemikian sehingga $f(x_k) = \alpha_k$ untuk $1\le k \le n$.
\end{prop}

\begin{proof}[\emph{Petunjuk pembuktian}] Untuk $x = \sum_{k=1}^n \beta_k x_k$, tetapkan $g(x) =
\sum_{k=1}^n \alpha_k \beta_k$. \ns
\end{proof}

\begin{prop}\label{C067131} Misalkan $V$ suatu ruang linear bernorma dan $\vc 0 \ne x \in V$. Maka
   \[ \norm x = \max\{\abs{f(x)}\colon f \in V^* \text{ dan } \norm f \le 1\}\,. \]
\end{prop}

Penggunaan $\max$ alih-alih $\sup$ dalam proposisi di atas disengaja. Proposisi tersebut menyatakan bahwa
terdapat $f$ dalam bola satuan tertutup dari $V^*$ sedemikian sehingga $f(x) = \norm x$.

\begin{defn} Suatu subhimpunan $A$ dari ruang topologis $X$ disebut
 \index{padat}%
\df{padat} di $X$ jika $\clo A = X$. Ruang $X$ disebut
 \index{separabel}%
\df{separabel} jika ruang itu memuat suatu subhimpunan padat yang terhitung.
\end{defn}

\begin{exam} Ruang $\K^n$ bersifat separabel terhadap topologi biasa; yakni,
 \index{biasa!topologi!untuk $\K^n$}%
 \index{topologi!pada $\K^n$}%
 \index{K@$\K^n$!topologi biasa pada}%
topologi yang diinduksi oleh norma biasa tersebut. (Lihat \ref{C023125}, \ref{0001503}, dan \ref{xmpl_top_metric}.)
\end{exam}

\begin{exam} Dengan topologi diskret, garis real $\R$ tidak separabel.
\end{exam}

\begin{prop} Misalkan $V$ suatu ruang linear bernorma. Jika $V^*$ separabel, maka $V$ juga separabel.
\end{prop}




















\endinput

 \chapter{RUANG HILBERT}

Untuk fakta-fakta paling mendasar tentang ruang Hilbert, saya tidak dapat memikirkan pengantar yang lebih baik daripada dua bab pertama
dari buku klasik Halmos~\cite{Halmos:1951}.  Pengantar lain yang baik adalah~\cite{Young:1988}.  Bab 3 dan 9 dari~\cite{Kreyszig:1978}
serta bab 4 dari~\cite{Ha:2006} mungkin juga berguna bagi Anda.  Sebagai langkah awal menuju penjelajahan yang lebih serius
atas seluk-beluk ruang Hilbert, lihat~\cite{Conway:1990}, \cite{Conway:2000}, dan~\cite{Helmberg:1969}.


\section{Definisi dan Contoh}

Ruang Hilbert adalah \emph{ruang hasil kali dalam lengkap}.  Secara lebih terperinci:
\begin{defn} Dalam proposisi~\ref{00015025} kita menunjukkan bagaimana hasil kali dalam pada suatu ruang vektor menginduksi norma pada ruang itu dan dalam
proposisi~\ref{0001503} bagaimana norma tersebut selanjutnya menginduksi metrik.  Jika suatu ruang hasil kali dalam lengkap terhadap metrik ini,
ruang itu disebut
  \index{Hilbert!ruang}%
  \index{ruang!Hilbert}%
\df{ruang Hilbert}.  Demikian pula,
  \index{Banach!ruang}%
  \index{ruang!Banach}%
\df{ruang Banach} adalah ruang linear bernorma lengkap dan
  \index{Banach!aljabar}%
  \index{aljabar!Banach}%
\df{aljabar Banach} adalah aljabar bernorma lengkap.
\end{defn}

\begin{exam} Dengan hasil kali dalam yang didefinisikan dalam
 \index{K@$\K^n$!sebagai ruang Hilbert}%
 \index{Hilbert!ruang!$\K^n$ sebagai}%
contoh~\ref{000150012}, $\K^n$ merupakan ruang Hilbert.
\end{exam}

\begin{exam}\label{exam150013} Dengan hasil kali dalam yang diberikan dalam
 \index{l@$l_2 = l_2(\N)$!sebagai ruang Hilbert}%
 \index{l@$l_2(\Z)$!sebagai ruang Hilbert}%
 \index{Hilbert!ruang!$l_2 = l_2(\N)$, $l_2(\Z)$ sebagai}%
 \index{l@$l_2 = l_2(\N)$!barisan yang kuadratnya dapat dijumlahkan}%
 \index{l@$l_2 = l_2(\N)$!sebagai ruang hasil kali dalam}%
 \index{hasil kali dalam!ruang!$l_2$ sebagai}%
contoh~\ref{000150013}, $l_2 = l_2(\N)$ dan $l_2(\Z)$ merupakan ruang Hilbert.
\end{exam}

\begin{exam} Dengan hasil kali dalam yang didefinisikan dalam contoh~\ref{000150014}, $\fml C([a,b])$ \emph{bukan} ruang Hilbert.
\end{exam}

\begin{exam}\label{exam_CX} Misalkan $\fml C(X) = \fml C(X,\C)$ adalah keluarga semua fungsi kontinu bernilai kompleks pada ruang
Hausdorff kompak~$X$.  Dengan
 \index{seragam!norma!pada $\fml C(X)$}%
 \index{norma!seragam}%
 \index{C@$\fml C(X)$!sebagai ruang Banach}%
 \index{C@$\fml C(X)$!fungsi kontinu bernilai kompleks pada~$X$}%
\df{norma seragam}
   \[ \norm f_u := \sup\{\abs{f(x)}\colon x \in X\} \]
$\fml C(X)$ merupakan ruang Banach.  Jika ruang $X$ tidak kompak, dengan cara yang sama kita dapat menjadikan $\fml C_b(X)$, yakni himpunan
 \index{C@$\fml C_b(X)$!sebagai ruang Banach}%
semua fungsi kontinu bernilai kompleks yang \emph{terbatas} pada $X$ sebagai ruang Banach.
Jika medan skalarnya adalah $\R$, kita akan
 \index{C@$\fml C(X,\R)$!fungsi kontinu bernilai real pada~$X$}%
 \index{C@$\fml C(X,\R)$!sebagai ruang Banach}%
menuliskan $\fml C(X,\R)$ untuk keluarga fungsi kontinu bernilai real pada~$X$.  Keluarga ini juga merupakan ruang Banach.
\end{exam}

\begin{exam}\label{000213} Ruang hasil kali dalam $l_c$ yang terdiri atas semua barisan $(a_n)$ bilangan kompleks yang pada akhirnya bernilai nol
(lihat contoh~\ref{exam13122}) bukan ruang Hilbert.
\end{exam}

\begin{exam}\label{00022} Misalkan $\mu$ adalah ukuran positif pada suatu $\sigma$-aljabar $\sfml A$ yang terdiri atas
subhimpunan-subhimpunan dari himpunan~$S$.  Jika Anda belum mengenal teori ukuran umum, pengantar singkat tentang ukuran Lebesgue
pada garis real seharusnya memadai untuk contoh-contoh yang muncul dalam catatan ini.  Suatu fungsi bernilai kompleks $f$ pada $S$ disebut
 \index{terukur}%
\df{terukur} jika prapeta oleh $f$ dari setiap himpunan Borel (secara ekuivalen, setiap himpunan
terbuka) di $\C$ termasuk dalam~$\sfml A$. Kita mendefinisikan relasi ekuivalensi $\sim$ pada keluarga
fungsi terukur bernilai kompleks dengan menetapkan $f \sim g$ apabila $f$ dan $g$ berbeda pada suatu himpunan
berukuran nol, yaitu apabila $\mu(\{x \in S\colon f(x) \ne g(x)\}) = 0$.  Kita menggunakan
notasi konvensional dan menyatakan kelas ekuivalensi yang memuat $f$ dengan $f$ itu sendiri (alih-alih
dengan sesuatu yang lebih logis seperti~$[f]$).  Keluarga (kelas-kelas ekuivalensi dari) fungsi
terukur bernilai kompleks pada $S$ dinyatakan dengan $\fml M(S,\mu)$. Suatu fungsi $f \in \fml M(S,\mu)$ disebut
 \index{terintegralkan kuadrat}%
 \index{M@$\fml M(S,\mu)$ (kelas ekuivalensi fungsi terukur)}%
 \index{L@$\fml L_2(S,\mu)$!fungsi terintegralkan kuadrat}%
\df{terintegralkan kuadrat} jika $\int_S \abs{f(x)}^2\,d\mu(x) < \infty$.  Keluarga
(kelas-kelas ekuivalensi dari) fungsi terintegralkan kuadrat pada $S$ dinyatakan dengan~$\lfs 2(S,\mu)$.
 \index{L@$\lfs 2(S,\mu)$!kelas fungsi terintegralkan kuadrat}%
 \index{L@$\lfs 2(S,\mu)$!sebagai ruang Hilbert}%
 \index{Hilbert!ruang!$\lfs 2(S,\mu)$ sebagai}%
Untuk setiap $f$, $g \in \lfs 2(S,\mu)$, definisikan
   \[ \langle f,g \rangle := \int_S f(x) \conj{g(x)}\,d\mu(x)\,. \]
Dengan hasil kali dalam ini (dan operasi ruang vektor titik demi titik yang sudah jelas), $\lfs 2(S,\mu)$ merupakan ruang Hilbert.
\end{exam}

\begin{exam}\label{exam_L1S} Gunakan notasi seperti dalam contoh sebelumnya.  Suatu fungsi $f \in \fml M(S,\mu)$ disebut
 \index{terintegralkan}%
 \index{L@$\fml L_1(S,\mu)$!fungsi terintegralkan}%
\df{terintegralkan} jika $\int_S \abs{f(x)}\,d\mu(x) < \infty$.  Keluarga
(kelas-kelas ekuivalensi dari) fungsi terintegralkan pada $S$ dinyatakan dengan~$\lfs 1(S,\mu)$.
 \index{L@$\lfs 1(S,\mu)$!kelas fungsi terintegralkan}%
 \index{L@$\lfs 1(S,\mu)$!sebagai ruang Banach}%
 \index{Banach!ruang!$\lfs 1(S,\mu)$ sebagai}%
Untuk setiap $f \in \lfs 1(S,\mu)$, definisikan
    \[ {\norm f}_1 := \int_S \abs{f(x)}\,d\mu(x)\,. \]
Dengan norma ini (dan operasi ruang vektor titik demi titik yang sudah jelas), $\lfs 1(S,\mu)$ merupakan ruang Banach.
\end{exam}

\begin{defn} Jika $J$ adalah sembarang interval di $\R$, suatu fungsi bernilai real (atau kompleks) $f$ pada $J$ disebut
 \index{mutlak!kekontinuan!fungsi}%
 \index{kekontinuan!mutlak!fungsi}%
 \index{kontinu!mutlak}%
\df{kontinu mutlak} jika memenuhi syarat berikut:
  \quote untuk setiap $\epsilon > 0$ terdapat $\delta > 0$ sedemikian
 sehingga jika $(a_1,b_1), \dots,(a_n,b_n)$ \\
 merupakan subinterval-subinterval tak kosong yang saling lepas dari $J$ dengan $\sum_{k=1}^n(b_k - a_k) < \delta$, \\
 maka $\sum_{k=1}^n\abs{f(b_k) - f(a_k)}~<~\epsilon$.
  \endquote
\end{defn}

\begin{exam}\label{exam_Hsp_abs_cont} Misalkan $\lambda$ adalah ukuran Lebesgue pada interval $[0,1]$ dan $H$ adalah himpunan semua
fungsi kontinu mutlak pada $[0,1]$ sedemikian sehingga $f\,'$ termasuk dalam $\lfs 2([0,1],\lambda)$ dan $f(0) = 0$.
 \index{Hilbert!ruang!fungsi kontinu mutlak}%
 \index{mutlak!fungsi kontinu!ruang Hilbert}%
Untuk $f$ dan $g$ dalam $H$, definisikan
   \[ \langle f,g \rangle = \int_0^1 f\,'(t)\clo{g\,'(t)}\,dt. \]
Ini merupakan hasil kali dalam pada $H$ yang menjadikan $H$ ruang Hilbert.
\end{exam}

\begin{exam} Jika $H$ dan $K$ merupakan ruang Hilbert, jumlah langsung (ortogonal eksternal) keduanya $H \oplus K$
 \index{Hilbert!ruang!jumlah langsung ruang Hilbert sebagai}%
 \index{langsung!jumlah!ruang Hilbert}%
(sebagaimana didefinisikan dalam~\ref{0001504}) juga merupakan ruang Hilbert.
\end{exam}


















\section{Penjumlahan Tak Berurutan}

Mendefinisikan \emph{penjumlahan tak berurutan} dalam ruang linear bernorma tidak lebih sulit daripada dalam~$\R$.  Kita memperumum
definisi~\ref{def_summable_R}.

\begin{defn} Misalkan $V$ adalah ruang linear bernorma dan $A \subseteq V$.
Untuk setiap $F \in \fin A$, definisikan
  \[ \sbsb sF = \sum F\]
di mana $\sum F$ adalah jumlah dari vektor-vektor (yang berhingga banyaknya) dalam~$F$.
Dengan demikian, $s = \bigl(\sbsb sF\bigr)_{F \in \fin A}$ adalah jaring dalam~$V$. Jika jaring ini konvergen, himpunan
$A$ disebut
 \index{dapat dijumlahkan!himpunan vektor}%
\df{dapat dijumlahkan}; limit jaring tersebut adalah
 \index{jumlah!himpunan yang dapat dijumlahkan}%
 \index{<sum@$\sum A$ (jumlah suatu himpunan vektor)}%
\df{jumlah} dari $A$ dan dinyatakan dengan~$\sum A$. Jika $\{\norm a\colon a \in A\}$ dapat dijumlahkan, kita mengatakan bahwa himpunan $A$
 \index{mutlak!dapat dijumlahkan}%
 \index{dapat dijumlahkan!secara mutlak}%
\df{dapat dijumlahkan secara mutlak}.

Untuk mengakomodasi jumlah yang memuat vektor yang sama lebih dari satu kali, kita perlu menggunakan keluarga vektor terindeks.  Keluarga demikian
memerlukan pendekatan yang sedikit berbeda. Misalkan, sebagai contoh, $A = \bigl(x_i\bigr)_{i \in I}$ dengan $I$
sebagai sembarang himpunan indeks. Maka untuk setiap $F \in \fin I$,
  \[ \sbsb sF = \sum_{i \in F}x_i . \]
Seperti di atas, $s$ adalah jaring dalam~$V$. Jika jaring ini konvergen, $\bigl(x_i\bigr)_{i \in I}$ merupakan
 \index{dapat dijumlahkan!keluarga vektor terindeks}%
keluarga yang dapat dijumlahkan dan limitnya adalah
 \index{jumlah!keluarga terindeks}%
 \index{<sum@$\sum_{i \in I} x_i$ (jumlah suatu keluarga terindeks)}%
jumlah keluarga terindeks tersebut, yang dinyatakan dengan $\sum_{i \in I} x_i$.  Cara lain untuk mengatakan bahwa $\bigl(x_i\bigr)_{i \in I}$
dapat dijumlahkan adalah dengan mengatakan bahwa ``deret'' $\sum_{i \in I}x_i$
 \index{konvergensi!deret tak hingga}%
 \index{jumlah!deret tak hingga}%
\df{konvergen} atau bahwa ``jumlah'' $\sum_{i \in I}x_i$ \df{ada}.  Suatu keluarga terindeks $\bigl(x_i\bigr)_{i \in I}$
 \index{mutlak!dapat dijumlahkan}%
 \index{dapat dijumlahkan!secara mutlak}%
\df{dapat dijumlahkan secara mutlak} jika keluarga terindeks $\bigl(\norm{x_i}\bigr)_{i \in I}$ dapat dijumlahkan.  Cara lain untuk mengatakan
bahwa $\bigl(x_i\bigr)_{i \in I}$ dapat dijumlahkan secara mutlak adalah dengan mengatakan bahwa ``deret'' $\sum_{i \in I}x_i$
 \index{konvergensi!mutlak}%
 \index{mutlak!konvergensi}%
\df{konvergen mutlak}.

Seperti biasa, barisan diperlakukan sebagai kasus tersendiri.  Misalkan $(a_k)$ adalah barisan dalam~$V$.  Jika deret tak hingga
$\sum_{k=1}^\infty a_k$ (yaitu, barisan
 \index{parsial!jumlah}%
 \index{jumlah!parsial}%
 \index{deret tak hingga!jumlah parsial}%
 \index{deret!tak hingga}%
jumlah parsial $s_n = \sum_{k=1}^n a_k$) konvergen ke suatu
vektor $b$ dalam $V$, kita mengatakan bahwa barisan $(a_k)$
 \index{deret tak hingga!dalam ruang linear bernorma}%
 \index{dapat dijumlahkan!barisan}%
 \index{barisan!dapat dijumlahkan}%
\df{dapat dijumlahkan} atau, secara ekuivalen, bahwa deret $\sum_{k=1}^\infty a_k$ merupakan
 \index{konvergensi!deret}%
 \index{deret!konvergen}%
\df{deret konvergen}.  Vektor $b$ disebut
 \index{jumlah!deret tak hingga}%
 \index{deret!jumlah}%
 \index{deret tak hingga!jumlah}%
\df{jumlah} deret $\sum_{k=1}^\infty a_k$ dan kita menulis
  \[ \sum_{k=1}^\infty a_k = b\,. \]
Jelas bahwa syarat perlu dan cukup agar suatu deret $\sum_{k=1}^\infty a_k$
konvergen atau, secara ekuivalen, agar barisan $(a_k)$ dapat dijumlahkan, adalah adanya suatu
vektor $b$ dalam $V$ sedemikian sehingga
  \begin{equation}\label{cond_ser_conv}
     \biggnorm{b - \sum_{k=1}^n a_k} \sto 0 \qquad \text{ketika} \qquad n \sto \infty.
  \end{equation}
\end{defn}

\begin{prop}\label{prop_abs_sum_implies_sum} Suatu ruang linear bernorma $V$ lengkap jika dan hanya jika setiap barisan yang
dapat dijumlahkan secara mutlak dalam $V$ dapat dijumlahkan.
\end{prop}

\begin{proof}[\emph{Petunjuk pembuktian}] Misalkan $(x_n)$ adalah barisan Cauchy dalam suatu ruang linear bernorma tempat setiap barisan yang
dapat dijumlahkan secara mutlak juga dapat dijumlahkan.  Carilah suatu subbarisan $\bigl(x_{n_k}\bigr)$ dari $(x_n)$ sedemikian sehingga $\norm{x_n - x_m} < 2^{-k}$
apabila $m$, $n \ge n_k$.  Jika kita menetapkan $y_1 = x_{n_1}$ dan $y_j = x_{n_j} - x_{n_{j-1}}$ untuk $j > 1$, maka $(y_j)$
dapat dijumlahkan secara mutlak.   \ns
\end{proof}

\begin{prop} Jika $\bigl(x_i\bigr)_{i \in I}$ dan $\bigl(y_i\bigr)_{i \in I}$ merupakan
keluarga terindeks yang dapat dijumlahkan dalam suatu ruang linear bernorma dan $\alpha$ adalah skalar, maka
$\bigl(\alpha x_i\bigr)_{i \in I}$ dan $\bigl(x_i + y_i\bigr)_{i \in I}$
dapat dijumlahkan; selain itu,
  \[ \sum_{i \in I} \alpha x_i = \alpha \, \sum_{i \in I} x_i \]
dan
  \[ \sum_{i \in I} x_i + \sum_{i \in I} y_i =
                         \sum_{i \in I} (x_i + y_i). \]
\end{prop}

\begin{rem} Dalam buku teks sering dijumpai suatu versi dari
``proposisi'' berikut.
\vskip 10 pt
\begin{quote}
Jika $\sum_{i \in I}x_i = u$ dan $\sum_{i \in I}x_i = v$ dalam suatu
ruang Hilbert (atau Banach), maka $u = v$.
\end{quote}

 \vskip 10 pt
\noindent (Saya pernah melihat bukti sepanjang 6 baris untuk pernyataan ini.) Pernyataan ini merupakan
konsekuensi sepele dari hasil yang mana?
\end{rem}

\begin{prop} Jika $\bigl(x_i\bigr)_{i \in I}$ adalah keluarga terindeks vektor yang dapat dijumlahkan dalam ruang Hilbert $H$ dan
$y$ adalah vektor dalam $H$, maka $\bigl(\langle x_i,y \rangle\bigr)_{i \in I}$ adalah keluarga terindeks skalar yang dapat dijumlahkan dan
   \[ \sum_{i \in I}\langle x_i,y \rangle  = \biggl < \sum_{i \in I} x_i, y \biggr >. \]
\end{prop}

Berikut ini merupakan perumuman dari contoh~\ref{exam150013}.

\begin{exam} Misalkan $I$ adalah sembarang himpunan tak kosong dan $l^2(I)$ adalah himpunan semua fungsi bernilai kompleks $x$ pada $I$ sedemikian sehingga
$\sum_{i \in I}\abs{x_i}^2 < \infty$. Untuk semua $x$, $y \in l^2(I)$, definisikan
    \[ \langle x,y \rangle = \sum_{i \in I}x(i)\conj{y(i)}.\]
  \begin{enumerate}
    \item Jika $x \in l^2(I)$, maka $x(i) = 0$ untuk semua indeks $i$, kecuali mungkin indeks-indeks dalam suatu himpunan terhitung.
    \item Pemetaan $(x,y) \mapsto \langle x,y \rangle$ merupakan hasil kali dalam pada ruang vektor $l^2(I)$.
    \item Ruang $l^2(I)$ merupakan ruang Hilbert.
  \end{enumerate}
\end{exam}

\begin{prop}\label{prop_sum_conv} Misalkan $(x_n)$ adalah barisan dalam suatu ruang Hilbert.  Jika barisan tersebut, ketika dipandang sebagai keluarga terindeks,
dapat dijumlahkan, maka deret tak hingga $\sum_{k=1}^\infty x_k$ konvergen.
\end{prop}

\begin{exam} Kebalikan dari proposisi~\ref{prop_sum_conv} tidak berlaku.
\end{exam}

Jumlah langsung dari dua ruang Hilbert merupakan ruang Hilbert.

\begin{prop} Jika $H$ dan $K$ merupakan ruang Hilbert, maka jumlah langsung (ortogonal) keduanya (sebagaimana didefinisikan dalam~\ref{0001504})
merupakan ruang Hilbert.
\end{prop}

Kita juga dapat mengambil hasil kali lebih dari dua ruang Hilbert---bahkan hasil kali suatu keluarga tak hingga ruang Hilbert.

\begin{prop} Misalkan $\bigl(H_i\bigr)_{i \in I}$ adalah keluarga terindeks ruang Hilbert.  Suatu fungsi $x\colon I \sto \bigcup H_i$
termasuk dalam $H = \D\bigoplus_{i \in I}H_i$ jika $x(i) \in H_i$ untuk setiap $i$ dan jika $\sum_{i \in I}\norm{x(i)}^2 < \infty$.
Mendefinisikan $\langle \hphantom x, \hphantom y \rangle$ pada $H$ dengan $\langle x, y \rangle := \sum_{i \in I} \langle
x(i), y(i) \rangle$ menjadikan $H$ ruang Hilbert.
\end{prop}

\begin{defn} Ruang Hilbert yang didefinisikan dalam proposisi sebelumnya adalah
 \index{eksternal!jumlah langsung ortogonal}%
 \index{ortogonal!jumlah langsung!ruang Hilbert}%
 \index{langsung!jumlah!ruang Hilbert}%
 \index{Hilbert!ruang!jumlah langsung}%
 \index{jumlah!langsung}%
\df{jumlah langsung ortogonal (eksternal)} dari ruang-ruang Hilbert~$H_i$.
\end{defn}

\begin{exer}  Apakah jumlah langsung yang didefinisikan di atas (dengan morfisme yang sesuai) merupakan hasil kali dalam kategori $\cat{HSp}$
yang objeknya ruang Hilbert dan morfismenya pemetaan linear terbatas?  Apakah jumlah langsung itu merupakan koproduk?
\end{exer}


























\section{Geometri Ruang Hilbert}

\begin{conv}  Dalam konteks ruang Hilbert, kata ``subruang'' akan selalu berarti
 \index{konvensi!subruang dari ruang Hilbert bersifat tertutup}%
 \index{subruang!ruang Hilbert}%
\emph{subruang vektor tertutup}.  Alasannya adalah bahwa kita ingin agar subruang dari suatu ruang Hilbert merupakan ruang Hilbert tersendiri; dengan kata lain,
subruang dari suatu ruang Hilbert seharusnya merupakan subobjek dalam kategori ruang Hilbert (dan morfisme yang sesuai).  Untuk menunjukkan bahwa $M$ adalah subruang dari $H$, kita menulis
 \index{<binrelle@$\preccurlyeq$ (subruang dari ruang Banach atau Hilbert)}%
$M \preccurlyeq H$.  Subruang vektor (yang tidak harus tertutup) dari suatu ruang Hilbert juga
 \index{linear!subruang}%
 \index{vektor!subruang}%
 \index{linear!manifold}%
 \index{manifold!linear}%
disebut \emph{subruang linear} atau \emph{manifold linear}.
\end{conv}

\begin{defn} Misalkan $A$ adalah subhimpunan tak kosong dari ruang Banach (atau Hilbert)~$B$.  Kita mendefinisikan
 \index{tertutup!rentang linear}%
 \index{linear!rentang!tertutup}%
 \index{<unops@$\bigvee A$ (rentang linear tertutup dari~$A$)}%
\df{rentang linear tertutup} dari $A$ (dinyatakan dengan $\bigvee A$) sebagai subruang terkecil dari $B$ yang memuat~$A$; yaitu,
irisan semua subruang dari $B$ yang memuat~$A$.
\end{defn}

Seperti pada definisi \ref{smplx005def}, \ref{def_interior}, dan \ref{def_closure}, kita perlu menunjukkan bahwa definisi sebelumnya bermakna.

\begin{prop} Definisi sebelumnya bermakna.  Lebih lanjut, suatu vektor berada dalam rentang linear tertutup dari $A$ jika dan hanya jika vektor itu termasuk
dalam penutupan rentang linear dari~$A$.
\end{prop}

\begin{thm}[Teorema Vektor Peminimum] Jika $C$ adalah subhimpunan konveks tertutup tak kosong dari suatu
 \index{teorema vektor peminimum}%
ruang Hilbert $H$ dan $a \in C^c$, maka terdapat tepat satu $b \in C$ sedemikian sehingga
$\norm{b - a} \le \norm{x - a}$ untuk setiap $x \in C$.
 \[\xymatrix{
     &&&&&{}  \\
     &&&& x   \\
     a\ar[rrrru]\ar[rrr] &&& b  \\
     &&&&&C   \\
     &&&&{}\ar@{-}@/^3.5pc/[uuuu]       }\]
\end{thm}

\begin{proof}[\emph{Petunjuk pembuktian}]  Untuk bagian eksistensi, misalkan $d = \inf\{\norm{x - a}\colon x \in C\}$.  Pilih suatu barisan
$(y_n)$ vektor dalam $C$ sedemikian sehingga $\norm{y_n - a} \sto d$.  Untuk membuktikan bahwa barisan $(y_n)$ adalah barisan Cauchy, terapkan
\emph{hukum jajaran genjang}~\ref{parallelogram_law} pada vektor $y_m - a$ dan $y_n - a$ untuk $m$, $n \in \N$.  Bagian keunikan
serupa: terapkan \emph{hukum jajaran genjang} pada vektor $b - a$ dan $c - a$, dengan $b$ dan $c$ sebagai dua vektor yang
memenuhi kesimpulan teorema.  \ns
\end{proof}

\begin{exam} Ruang vektor $\R^2$ dengan metrik seragam merupakan ruang Banach. Untuk melihat
bahwa \emph{Teorema Vektor Peminimum} tidak berlaku dalam ruang ini, ambillah $C$ sebagai
bola satuan tertutup berpusat di titik asal dan $a$ sebagai titik~$(2,0)$.
\end{exam}

\begin{exam} Himpunan-himpunan
   \[ C_1 = \biggl\{f \in \fml C([0,1])\colon \int_0^{1/2}f - \int_{1/2}^1f = 1\biggr\} \qquad \text{dan} \qquad
                               C_2 = \biggl\{f \in \lfs 1([0,1],\lambda)\colon \int_0^1f = 1\biggr\} \]
merupakan contoh yang menunjukkan bahwa baik pernyataan eksistensi maupun pernyataan keunikan dalam \emph{Teorema Vektor Peminimum}
tidak harus berlaku dalam ruang Banach. (Simbol $\lambda$ menyatakan ukuran Lebesgue.)
\end{exam}

\begin{thm}[Teorema Dekomposisi Vektor]\label{thm_vec_decomp} Misalkan $H$ adalah ruang Hilbert dan
 \index{vektor!teorema dekomposisi}%
$M$ adalah subruang (linear tertutup) dari~$H$.
Maka untuk setiap $x \in H$ terdapat tepat satu pasangan vektor $y \in M$ dan $z \in M^\perp$ sedemikian sehingga $x = y + z$.
\end{thm}

\begin{proof}[\emph{Petunjuk pembuktian}] Asumsikan $x \notin M$ (jika tidak, hasilnya sepele). Jelaskan bagaimana kita
mengetahui bahwa terdapat tepat satu vektor $y \in M$ sedemikian sehingga
   \begin{equation} \norm{x - y} \le \norm{x - m} \end{equation}
untuk semua $m \in M$.  Jelaskan mengapa, untuk bagian eksistensi dari bukti, cukup dibuktikan
bahwa $x - y$ tegak lurus terhadap setiap vektor \emph{satuan} $m \in M$.  Jelaskan mengapa benar
bahwa untuk setiap vektor satuan $m \in M$,
   \begin{equation} \norm{x - y}^2 \le \norm{x - (y + \lambda m)}^2 \end{equation}
di mana $\lambda := \langle x - y, m \rangle$.  \ns
\end{proof}

\begin{cor} Jika $M$ merupakan subruang dari suatu ruang Hilbert~$H$, maka $H = M \oplus M^\perp$.
\end{cor}

\begin{exam} Hasil sebelumnya menyatakan bahwa setiap subruang dari suatu ruang Hilbert merupakan suku
langsung. Hasil ini tidak harus benar jika $M$ hanya diasumsikan sebagai subruang linear dari ruang tersebut.
Sebagai contoh, perhatikan bahwa $M = l_c$ (lihat contoh~\ref{000213}) adalah subruang linear dari
ruang Hilbert $l_2$ (lihat contoh~\ref{exam150013}), tetapi $M^\perp = \{\vc 0\}$.
\end{exam}

\begin{cor} Suatu subruang vektor $A$ dari ruang Hilbert $H$ padat dalam $H$ jika dan hanya jika $A^\perp = \{0\}$.
\end{cor}

\begin{prop} Setiap subruang sejati dari suatu ruang Hilbert memiliki interior kosong.
\end{prop}

\begin{prop} Jika suatu basis Hamel bagi ruang Hilbert tidak hingga, maka basis tersebut tak terhitung.
\end{prop}

\begin{proof}[\emph{Petunjuk pembuktian}] Andaikan basis tak hingga terhitungnya ialah $\{e^1, e^2, e^3, \dots \}$.
Untuk setiap $n \in \N$, misalkan $M_n = \spn\{e^1, e^2, \dots, e^n\}$.  Terapkan \emph{teorema kategori Baire} (lihat~\cite{Erdman:2007},
teorema 29.3.20) pada $\bigcup_{n=1}^\infty M_n$ dengan mengingat proposisi sebelumnya.   \ns
\end{proof}

\begin{prop}\label{00023} Misalkan $H$ adalah ruang Hilbert. Maka pernyataan-pernyataan berikut berlaku:
 \begin{enumerate}
  \item jika $A \subseteq H$, maka $A \subseteq A^{\perp\perp}$;
  \item jika $A \subseteq B \subseteq H$, maka $B^\perp \subseteq A^\perp$;
  \item jika $M$ adalah subruang dari $H$, maka $M = M^{\perp\perp}$; dan
  \item jika $A \subseteq H$, maka $\bigvee A = A^{\perp\perp}$.
 \end{enumerate}
\end{prop}

\begin{prop} Misalkan $M$ dan $N$ adalah subruang dari suatu ruang Hilbert.  Maka
 \begin{enumerate}
  \item  $(M + N)^\perp = (M \cup N)^\perp =  M^\perp \cap N^\perp$, dan
  \item  $(M \cap N)^\perp = \clo{M^\perp + N^\perp}$.
 \end{enumerate}
\end{prop}

\begin{prop} Misalkan $V$ adalah ruang linear bernorma.  Terdapat hasil kali dalam pada $V$ yang menginduksi norma pada $V$
jika dan hanya jika norma tersebut memenuhi \emph{hukum jajaran genjang}~\ref{parallelogram_law}.
\end{prop}

\begin{proof}[\emph{Petunjuk pembuktian}]  Kita mengatakan bahwa suatu hasil kali dalam $\langle \,\cdot\,\,,\,\cdot\, \rangle$
pada $V$ menginduksi norma $\norm{\,\cdot\,}$ jika $\norm x = \sqrt{\langle x,x \rangle}$ berlaku untuk setiap $x$ dalam $V$.
Untuk membuktikan bahwa jika suatu norma memenuhi \emph{hukum jajaran genjang}, maka norma tersebut diinduksi oleh hasil kali dalam,
gunakan persamaan dalam \emph{identitas polarisasi}~\ref{0001507} sebagai definisi.
Buktikan terlebih dahulu bahwa $\langle y,x \rangle = \conj{\langle x,y \rangle}$ untuk semua $x$, $y\in V$ dan bahwa
$\langle x,x \rangle > 0$ apabila $x \ne 0$.  Selanjutnya, untuk sembarang $z \in V$, definisikan fungsi
$f\colon V \sto \C\colon x \mapsto \langle x,z \rangle$.  Buktikan bahwa
   \[ f(x + y) + f(x - y) = 2f(x) \tag{$\ast$} \]
untuk semua $x$, $y \in V$.  Gunakan ini untuk membuktikan bahwa $f(\alpha x) = \alpha f(x)$ untuk semua
$x \in V$ dan $\alpha \in \R$.  (Mulailah dengan $\alpha$ berupa bilangan asli.)  Kemudian
tunjukkan bahwa $f$ aditif. (Jika $u$ dan $v$ adalah sembarang anggota $V$, misalkan $x = u + v$
dan $y = u - v$.  Gunakan~$(\ast)$.)  Terakhir, buktikan bahwa $f(\alpha x) = \alpha f(x)$ untuk
$\alpha$ kompleks dengan menunjukkan bahwa $f(ix) = i\,f(x)$ untuk semua $x \in V$.  \ns
\end{proof}
















\section{Himpunan Ortonormal dan Basis}\label{sec_onbases}
\begin{defn} Suatu subhimpunan tak kosong $E$ dari ruang Hilbert bersifat
 \index{ortonormal}%
\df{ortonormal} jika $e \perp f$ untuk setiap pasangan vektor berbeda $e$ dan $f$ di $E$, dan
$\norm e = 1$ untuk setiap $e \in E$.
\end{defn}

\begin{prop} Dalam ruang hasil kali dalam, setiap himpunan ortonormal bebas linear.
\end{prop}

\begin{prop}[Pertidaksamaan Bessel]\label{C063141} Misalkan $H$ ruang Hilbert dan $E$ suatu
 \index{pertidaksamaan Bessel}%
subhimpunan ortonormal dari~$H$. Untuk setiap $x \in H$,
  \[ \sum_{e \in E} \norm{\langle x,e \rangle}^2 \le \norm x^2\,. \]
\end{prop}

\begin{cor} Jika $E$ subhimpunan ortonormal dari ruang Hilbert $H$ dan $x \in H$, maka
$\{e \in E\colon \langle x,e \rangle~\ne~0\}$ terhitung.
\end{cor}

\begin{prop}\label{C063144} Misalkan $E$ himpunan ortonormal dalam ruang Hilbert~$H$. Pernyataan-pernyataan
berikut ekuivalen.
 \begin{enumerate}
  \item $E$ adalah himpunan ortonormal maksimal.
  \item Jika $x \perp E$, maka $x = 0$ ($E$ bersifat \emph{total}).
 \index{total!himpunan ortonormal}%
  \item $\bigvee E = H$ ($E$ adalah \emph{himpunan ortonormal lengkap}).
 \index{lengkap!himpunan ortonormal}%
  \item $\D x = \sum_{e \in E} \langle x,e \rangle e$ untuk semua $x \in H$. \emph{(Ekspansi
Fourier)}
 \index{Fourier!ekspansi}%
  \item $\D \langle x,y \rangle = \sum_{e \in E} \langle x,e \rangle \langle e,y \rangle$
untuk semua $x,y \in H$. \emph{(Identitas Parseval.)}
 \index{identitas Parseval}%
  \item $\D \norm{x}^2 = \sum_{e \in E} \abs{\langle x,e \rangle}^2$ untuk semua $x \in H$.
\emph{(Identitas Parseval.)}
 \end{enumerate}
\end{prop}

\begin{defn}\label{C063147} Jika $H$ ruang Hilbert, maka himpunan ortonormal $E$ yang memenuhi
salah satu (dan karenanya semua) kondisi dalam proposisi~\ref{C063144} disebut
 \index{ortonormal!basis}%
 \index{basis!ortonormal}%
\df{basis ortonormal} untuk~$H$ (atau
 \index{lengkap!himpunan ortonormal}%
 \index{ortonormal!himpunan!lengkap}%
\df{himpunan ortonormal lengkap} dalam~$H$, atau
 \index{Hilbert!ruang!basis untuk}%
 \index{basis!untuk ruang Hilbert}%
\df{basis ruang Hilbert} untuk~$H$). Perhatikan bahwa konsep ini sangat berbeda dari
basis biasa (Hamel) untuk ruang vektor (lihat~\ref{defn_Hamel_basis}). Untuk ruang Hilbert,
vektor-vektor basis harus memiliki panjang satu dan saling tegak lurus
(kondisi yang tidak bermakna dalam ruang vektor); tetapi \emph{tidak} disyaratkan bahwa
setiap vektor dalam ruang tersebut merupakan kombinasi linear dari vektor-vektor basis.
\end{defn}

\begin{prop} Jika $E$ subhimpunan ortonormal dari ruang Hilbert $H$, maka terdapat
basis ortonormal untuk $H$ yang memuat~$E$.
\end{prop}

\begin{cor} Setiap ruang Hilbert tak nol memiliki basis.
\end{cor}

\begin{prop} Jika $E$ subhimpunan ortonormal dari ruang Hilbert $H$ dan $x \in H$, maka
$\{\langle x,e \rangle e: e \in E\}$ dapat dijumlahkan.
\end{prop}

\begin{exam}\label{C063149} Untuk setiap $n \in \N$, misalkan $\vc e^n$ adalah barisan dalam $l_2$ yang
koordinat ke-$n$-nya bernilai $1$ dan semua koordinat lainnya bernilai~$0$. Maka $\{\vc e^n\colon n \in \N\}$
adalah basis ortonormal untuk~$l_2$. Basis ini disebut
 \index{biasa!basis ortonormal untuk~$l_2$}%
 \index{ortonormal!basis!biasa (atau standar) untuk~$l_2$}%
 \index{basis!standar (atau biasa)}%
\df{biasa} (atau
 \index{standar!basis ortonormal}%
\df{standar}); basis tersebut adalah \df{basis ortonormal} untuk~$l_2$.
\end{exam}

\begin{exer} Suatu fungsi $f\colon \R \sto \C$ berbentuk
   \[ f(t) = a_0 + \sum_{k=1}^n (a_k \cos kt + b_k \sin kt) \]
dengan $a_0, \dots, a_n, b_1, \dots, b_n$ bilangan kompleks disebut
 \index{polinom trigonometri}%
 \index{polinom!trigonometri}%
\df{polinom trigonometri}.
 \begin{enumerate}
  \item Jelaskan mengapa fungsi bernilai kompleks yang $2\pi$-periodik pada garis real $\R$
sering diidentifikasi dengan fungsi bernilai kompleks pada lingkaran satuan~$\T$.
  \item Beberapa penulis menyebut \emph{polinom trigonometri} sebagai fungsi berbentuk
    \[ f(t) = \sum_{k=-n}^n c_k e^{ikt} \]
 dengan $c_{-n}, \dots, c_{-1}, c_0, c_1, \dots, c_n$ bilangan kompleks. Jelaskan dengan cermat mengapa bentuk ini persis
sama dengan definisi sebelumnya.
  \item Jelaskan alasan penggunaan istilah \emph{polinom trigonometri}.
  \item Buktikan bahwa setiap fungsi kontinu yang $2\pi$-periodik pada $\R$ merupakan limit seragam
dari suatu barisan polinom trigonometri.
 \end{enumerate}
\end{exer}

\begin{prop}\label{prop_unique_onbasis} Misalkan $E$ basis untuk ruang Hilbert $H$ dan $x \in H$.
Jika $x = \sum_{e \in E} \alpha_e e$, maka $\alpha_e = \langle x,e \rangle$ untuk setiap $e \in E$.
\end{prop}

\begin{prop} Tinjau ruang Hilbert $\lfs 2([0,2\pi],\C)$ yang terdiri atas (kelas-kelas ekuivalensi dari) fungsi
 \index{L@$\lfs 2([0,2\pi],\C)$}%
bernilai kompleks yang terintegralkan kuadrat pada $[0,2\pi]$, dengan hasil kali dalam
  \[ \langle f,g \rangle = \frac1{2\pi}\int_0^{2\pi} f(x) \conj{g(x)}\,dx. \]
Untuk setiap bilangan bulat $n$ (positif, negatif, atau nol), misalkan $\vc e^n(x) = e^{inx}$ untuk semua $x
\in [0,2\pi]$. Maka himpunan fungsi $\vc e^n$ untuk $n \in \Z$ merupakan
basis ortonormal untuk $\lfs 2([0,2\pi],\C)$.
\end{prop}

\begin{exam} Jumlah deret tak hingga $\D\sum_{k=1}^\infty \frac1{k^2}$ adalah $\dfrac{\pi^2}6$\,.
\end{exam}

\begin{proof}[\emph{Petunjuk pembuktian}] Misalkan $f(x) = x$ untuk $0 \le x \le 2\pi$. Pandang $f$ sebagai anggota
$\lfs 2([0,2\pi])$. Tentukan $\norm{f}$ dengan dua cara. \ns
\end{proof}

Proposisi berikut memungkinkan kita mendefinisikan gagasan \emph{dimensi} ruang Hilbert. (Pada ruang berdimensi
hingga, gagasan ini sama dengan konsep dalam ruang vektor.)

\begin{prop} Sebarang dua basis ortonormal dari suatu ruang Hilbert memiliki kardinalitas yang sama (artinya, keduanya
berkorespondensi satu-satu).
\end{prop}

\begin{proof} Lihat \cite{Halmos:1951}, halaman 29, Teorema 1. \ns
\end{proof}

\begin{defn}
 \index{dimensi}%
Yang dimaksud dengan \df{dimensi} ruang Hilbert adalah kardinalitas dari sebarang basis ortonormal untuk ruang tersebut.
\end{defn}

\begin{prop}\label{prop_Schauder_basis_Hsp} Ruang Hilbert separabel jika dan hanya jika ruang tersebut memiliki
 \index{separabel!ruang Hilbert}%
 \index{Hilbert!ruang!separabel}%
basis ortonormal terhitung.
\end{prop}

\begin{proof} Lihat \cite{Kreyszig:1978}, halaman 171, Teorema 3.6-4. \ns
\end{proof}





















\section{Teorema Riesz--Fr\'echet}
\begin{exam}\label{000341} Misalkan $H$ ruang Hilbert dan $a \in H$. Definisikan $\psi_a\colon
H \sto \C\colon x \mapsto \langle x,a \rangle$. Maka $\psi_a \in H^*$.
\end{exam}

Sekarang kita menggeneralisasi teorema~\ref{0001601} ke ranah berdimensi tak hingga. Hasil ini
kadang-kadang disebut \emph{Teorema Representasi Riesz} (yang dapat menimbulkan kebingungan dengan
 \index{representasi!fungsional linear kontinu}%
 \index{representasi!teorema kecil Riesz}%
hasil yang lebih substansial mengenai representasi fungsional linear tertentu sebagai ukuran)
atau \emph{Teorema Representasi Riesz Kecil}. Teorema ini menyatakan bahwa \emph{satu-satunya}
fungsional linear kontinu pada ruang Hilbert adalah fungsi $\psi_a$ yang didefinisikan dalam~\ref{000341}.
Dengan kata lain, jika diberikan fungsional linear kontinu $f$ pada ruang Hilbert, terdapat vektor unik
$a$ dalam ruang tersebut sehingga nilai $f$ pada vektor $x$ dapat dinyatakan hanya dengan mengambil
hasil kali dalam $x$ dengan~$a$.

\begin{thm}[Teorema Riesz--Fr\'echet]\label{000342} Jika $f \in H^*$ dengan $H$ ruang Hilbert,
 \index{teorema Riesz--Fr\'echet}%
maka terdapat vektor unik $a$ dalam $H$ sehingga $f = \psi_a$. Selain itu,
$\norm a = \norm{\psi_a}$.
\end{thm}

\begin{proof}[\emph{Petunjuk pembuktian}] Untuk kasus $f \ne 0$, pilih vektor satuan $z$ dalam $(\ker f)^\perp$.
Perhatikan bahwa untuk setiap $x \in H$, vektor $x - \frac{f(x)}{f(z)}\,z$ termasuk dalam $\ker f$. \ns
\end{proof}

\begin{cor} Kernel setiap fungsional linear terbatas tak nol pada ruang Hilbert memiliki kodimensi~$1$.
\end{cor}

\begin{proof}[\emph{Petunjuk pembuktian}]
 \index{kodimensi}%
Yang dimaksud dengan \df{kodimensi} suatu subruang dari ruang Hilbert adalah dimensi komplemen ortogonalnya. Pernyataan ini sebenarnya
merupakan akibat dari \emph{pembuktian} \emph{Teorema Riesz--Fr\'echet}. Dalam pembuktian tersebut kita mendapati bahwa vektor
yang merepresentasikan fungsional linear terbatas $f$ merupakan kelipatan (tak nol) dari vektor satuan dalam $(\ker f)^\perp$.
Apa yang akan terjadi jika terdapat dua vektor bebas linear dalam $(\ker f)^\perp$? \ns
\end{proof}

\begin{defn} Ingat bahwa pemetaan $T\colon V \sto W$ antara ruang-ruang vektor kompleks bersifat
 \index{linear konjugat}%
 \index{linear!konjugat}%
\df{linear konjugat} jika $T(x + y) = Tx + Ty$ dan $T(\alpha x) = \conj{\alpha}\,Tx$ untuk semua
$x$, $y \in V$ dan $\alpha \in \C$. Pemetaan linear konjugat yang bijektif disebut
 \index{antiisomorfisme}%
 \index{isomorfisme!anti-}%
\df{antiisomorfisme}.
\end{defn}

\begin{exam}\label{0022057} Untuk ruang Hilbert $H$, pemetaan
$\psi\colon H \sto H^*\colon a \mapsto \psi_a$ yang didefinisikan dalam contoh~\ref{000341} merupakan
antiisomorfisme isometrik antara $H$ dan ruang dualnya.
\end{exam}

\begin{defn} Pemetaan linear konjugat $C\colon V \sto V$ dari ruang vektor ke dirinya sendiri
yang memenuhi $C^2 = \id V$ disebut
 \index{konjugasi}
\df{konjugasi} pada~$V$.
\end{defn}

\begin{exam} Konjugasi kompleks $z \mapsto \conj z$ merupakan contoh konjugasi dalam
ruang vektor~$\C$.
\end{exam}

\begin{exam}\label{0022067} Misalkan $H$ ruang Hilbert dan $E$ basis untuk~$H$. Maka
pemetaan $x \mapsto \sum_{e \in E}\langle e,x \rangle e$ merupakan konjugasi dan isometri pada~$H$.
\end{exam}

Istilah ``antiisomorfisme'' agak menyesatkan. Istilah tersebut mengesankan sesuatu yang sama sekali
berbeda dari isomorfisme. Padahal, antiisomorfisme hampir sebaik isomorfisme.
Pada intinya, proposisi berikut menyatakan bahwa ruang Hilbert dan ruang dualnya isomorfik
\emph{karena} keduanya antiisomorfik.

\begin{prop}\label{0022071} Misalkan $H$ ruang Hilbert, $\psi\colon H \sto H^*$ adalah
antiisomorfisme yang didefinisikan dalam~\ref{000341} (lihat~\ref{0022057}), dan $C$ adalah konjugasi
yang didefinisikan dalam~\ref{0022067}. Maka komposisi $C \psi^{-1}$ merupakan isomorfisme isometrik
dari $H^*$ ke~$H$.
\end{prop}

\begin{cor}\label{cor_Hsp_iso_dual} Setiap ruang Hilbert isomorfik secara isometrik dengan ruang dualnya.
\end{cor}

\begin{prop} Misalkan $H$ ruang Hilbert dan $\psi\colon H \sto H^*$ antiisomorfisme
yang didefinisikan dalam~\ref{000341} (lihat~\ref{0022057}). Jika kita mendefinisikan $\langle f,g \rangle :=
\langle \psi^{-1}g, \psi^{-1}f \rangle$ untuk semua $f$, $g \in H^*$, maka $H^*$ menjadi
ruang Hilbert yang isomorfik secara isometrik dengan~$H$. Norma yang dihasilkan pada $H^*$ adalah norma
yang biasa digunakan.
\end{prop}

\begin{cor} Setiap ruang Hilbert bersifat refleksif.
\end{cor}

\begin{exam} Misalkan $H$ himpunan semua fungsi bernilai kompleks $f$ yang kontinu mutlak
pada $[0,1]$ sedemikian sehingga $f(0) = 0$ dan $f' \in
\lfs 2([0,1],\lambda)$. Definisikan hasil kali dalam pada $H$ dengan
 \[\langle f,g \rangle := \int_0^1 f'(t)\,\clo{g'(t)}\,dt.\]
Kita telah mengetahui bahwa $H$ ruang Hilbert (lihat contoh~\ref{exam_Hsp_abs_cont}).
Tetapkan $t \in (0,1]$. Definisikan
   \[ E_t\colon H \sto \C\colon f \mapsto f(t). \]
Maka $E_t$ merupakan fungsional linear terbatas pada~$H$. Berapakah $\norm{E_t}$?
Vektor $g$ manakah dalam $H$ yang merepresentasikan fungsional~$E_t$?
\end{exam}

\begin{exer} Pada suatu sore, kawan Anda, Fred R. Dimm, datang membawa masalah. ``Begini,'' katanya,
``Pada $\lfs 2 = \lfs 2(\R,\R)$, fungsional evaluasi $E_0$, yang mengevaluasi setiap anggota $\lfs 2$
di~$0$, jelas merupakan fungsional linear terbatas. Jadi, menurut Teorema Representasi Riesz,
seharusnya ada fungsi $g$ dalam $\lfs 2$ sedemikian sehingga $f(0) = \langle f,g \rangle = \int_{-\infty}^\infty fg$
untuk semua $f$ dalam~$\lfs 2$. Namun, itu adalah sifat `fungsi~$\delta$', yang katanya
tidak ada. Di mana letak kesalahan saya?'' Berilah Freddy nasihat (yang membantu).
\end{exer}

\begin{exer} Masih pada sore yang sama, Freddy kembali. Setelah mempertimbangkan
nasihat yang Anda berikan (dalam soal sebelumnya), ia merevisi
pertanyaannya dengan menjadikan domain $E_0$ sebagai himpunan fungsi bernilai real yang terbatas dan
kontinu pada~$\R$. Dengan sabar Anda
menjelaskan kepada Freddy bahwa sekarang setidaknya ada dua kekeliruan dalam caranya
menggunakan `fungsi~$\delta$'. Apakah kedua kekeliruan tersebut?
\end{exer}

\begin{exer} Mimpi buruk itu berlanjut. Sekarang pukul 23.00 dan Freddy kembali lagi.
Kali ini (setelah, untungnya, menyerah soal fungsi~$\delta$) ia ingin
menerapkan teorema representasi pada fungsional ``evaluasi
pada nol'' pada ruang~$l_2(\Z)$. Dapatkah ia melakukannya? Jelaskan.
\end{exer}



























\section{Topologi Kuat dan Lemah pada Ruang Hilbert}
Kita mengingat kembali secara singkat definisi dan fakta-fakta penting mengenai \emph{topologi lemah}. Untuk pengantar yang lebih terperinci
mengenai kelas topologi ini, lihat buku teks topologi mana pun yang baik atau bagian 11.4 dari catatan saya~\cite{Erdman:2007}.

\begin{defn}\label{C021414} Misalkan $S$ suatu himpunan, untuk setiap $\alpha \in A$ (dengan
$A$ sebarang himpunan indeks) $X_\alpha$ merupakan ruang topologis, dan $f_\alpha
\colon S \sto X_\alpha$ untuk setiap $\alpha \in A$. Misalkan
   \[ \sfml S = \{{f_\alpha}^\gets(U_\alpha) \colon \open{U_\alpha}{X_\alpha}
                                             \text{ dan } \alpha \in A\}\,. \]
Gunakan keluarga $\sfml S$ sebagai subbasis bagi suatu topologi pada~$S$. Topologi ini disebut
 \index{lemah!topologi!diinduksi oleh keluarga fungsi}%
 \index{topologi!lemah!diinduksi oleh keluarga fungsi}%
\df{topologi lemah} yang diinduksi oleh (atau ditentukan oleh) fungsi-fungsi~$f_\alpha$.
\end{defn}

\begin{prop}\label{C021421} Dalam topologi lemah (yang didefinisikan di atas) pada himpunan $S$, setiap fungsi $f_\alpha$ kontinu;
bahkan, topologi lemah adalah topologi paling lemah pada
 \index{lemah!topologi!karakterisasi}%
 \index{topologi!lemah!karakterisasi}%
$S$ yang menjadikan fungsi-fungsi tersebut kontinu.
\end{prop}

\begin{prop}\label{C021424} Misalkan $X$ ruang topologis dengan topologi lemah yang ditentukan
oleh keluarga fungsi $\fml F$. Buktikan bahwa fungsi $g\colon W \sto X$, dengan $W$
ruang topologis, kontinu jika dan hanya jika $f \circ g$ kontinu untuk setiap $f
\in \fml F$.
\end{prop}

\begin{defn}\label{C021437} Misalkan $\bigl(A_\lambda\bigr)_{\lambda \in \Lambda}$ keluarga terindeks
ruang topologis. Topologi lemah pada hasil kali Kartesius $\prod
A_\lambda$ yang diinduksi oleh keluarga proyeksi koordinat $\pi_\lambda$ (lihat
definisi~\ref{C015531}) disebut
 \index{hasil kali!topologi}%
 \index{topologi!hasil kali}%
\df{topologi hasil kali}.
\end{defn}

\begin{prop} Hasil kali dari suatu keluarga ruang Hausdorff adalah Hausdorff.
\end{prop}

\begin{exam}\label{C031473} Misalkan $Y$ ruang topologis dengan topologi lemah (lihat
definisi~\ref{C021414}) yang ditentukan oleh keluarga fungsi terindeks $(f_\alpha)_{\alpha \in A}$,
dengan sifat bahwa, untuk setiap $\alpha \in A$, kodomain $f_\alpha$ berupa ruang topologis~$X_\alpha$.
Maka jaring $\bigl(y_{{}_\lambda}\bigr)_{\lambda \in \Lambda}$ dalam $Y$
konvergen ke titik $a \in Y$ jika dan hanya jika
\mbox{$f_\alpha(y_{{}_\lambda})~\to_{\lambda \in \Lambda}~f_\alpha(a)$} untuk setiap $\alpha
\in A$.
\end{exam}

\begin{exam}\label{C031475} Jika $Y = \prod_{\alpha \in A} X_\alpha$ adalah hasil kali dari ruang-ruang topologis tak kosong
dengan topologi hasil kali (lihat definisi~\ref{C021437}), maka suatu jaring
$\bigl(\sbsb y\lambda\bigr)$ dalam $Y$ konvergen ke $a \in Y$ jika dan hanya jika $\bigl(\sbsb
y\lambda\bigr)_\alpha \sto a_\alpha$ untuk setiap $\alpha \in A$.
\end{exam}

\begin{exam}\label{exam_prod_top_norm} Jika $V$ dan $W$ ruang linear bernorma, maka topologi hasil kali pada
$V \times W$ sama dengan topologi yang diinduksi oleh norma hasil kali pada $V \oplus W$.
\end{exam}

Contoh berikut menggambarkan ketidakcukupan barisan ketika kita menangani ruang topologis umum.

\begin{exam} Misalkan $\fml F = \fml F([0,1])$ adalah himpunan semua fungsi $f\colon [0,1] \sto \R$, dan beri
$\fml F$ topologi hasil kali, yaitu topologi lemah yang ditentukan oleh fungsional-fungsional
evaluasi
    \[ E_x\colon \fml F \sto \R\colon f \mapsto f(x) \]
dengan $x \in [0,1]$. Jadi, lingkungan terbuka dasar dari suatu titik $g \in \fml F$ ditentukan
oleh himpunan titik hingga $A \in \fin[0,1]$ dan bilangan $\epsilon > 0$:
  \[ U(g\,;\,A\,;\, \epsilon) := \{f \in \fml F\colon \abs{f(t) - g(t)} < \epsilon
                           \text{ untuk semua $t \in A$}\}. \]
(Apa nama yang biasa digunakan untuk topologi ini?) Misalkan $\fml G$ himpunan semua fungsi dalam $\fml F$
yang bertumpuan hingga. Maka
  \begin{enumerate}
    \item fungsi konstan $\vc 1$ termasuk dalam $\clo{\fml G}$,
    \item terdapat jaring fungsi dalam $\fml G$ yang konvergen ke~$\vc 1$, tetapi
    \item tidak ada barisan dalam $\fml G$ yang konvergen ke~$\vc 1$.
  \end{enumerate}
\end{exam}

\begin{defn} Jaring $(x_{{}_{\sst\lambda}})$ dalam ruang Hilbert $H$ dikatakan
 \index{konvergensi!dalam ruang Hilbert!lemah}%
 \index{lemah!konvergensi!dalam ruang Hilbert}%
\df{konvergen secara lemah} ke titik $a$ dalam $H$ jika $\langle x{{}_{\sst\lambda}}, y \rangle \to \langle a,y \rangle$
untuk setiap $y \in H$. Dalam hal ini kita menulis $x_{{}_{\sst\lambda}} \to^w a$. Dalam ruang Hilbert, suatu jaring dikatakan
 \index{konvergensi!dalam ruang Hilbert!kuat}%
 \index{kuat!konvergensi!dalam ruang Hilbert}%
\df{konvergen secara kuat} (atau
 \index{konvergensi!dalam norma}%
 \index{norma!konvergensi}%
\df{konvergen dalam norma}) jika jaring tersebut konvergen terhadap topologi yang diinduksi oleh norma. Inilah jenis
konvergensi yang biasa dalam ruang Hilbert. Jika kita ingin menekankan perbedaan di antara cara-cara konvergensi, kita dapat menulis
$x_{{}_{\sst\lambda}} \to^s a$ untuk konvergensi kuat.
\end{defn}

\begin{exer} Jelaskan mengapa topologi pada ruang Hilbert yang dibangkitkan oleh konvergensi lemah jaring sesungguhnya merupakan
topologi lemah dalam pengertian topologis yang lazim.
\end{exer}

Ketika kita mengatakan bahwa suatu subhimpunan dari ruang Hilbert $H$ bersifat
 \index{lemah!ketertutupan}%
 \index{tertutup!secara lemah}%
\df{tertutup secara lemah}, maksudnya subhimpunan tersebut tertutup terhadap topologi lemah pada $H$; suatu himpunan bersifat
 \index{lemah!kekompakan}%
 \index{kompak!secara lemah}%
\df{kompak secara lemah} jika kompak dalam topologi lemah pada~$H$; dan seterusnya. Suatu fungsi
$f\colon H \sto H$ bersifat
 \index{lemah!kontinuitas}%
 \index{kontinu!secara lemah}%
\df{kontinu secara lemah} jika kontinu sebagai pemetaan dari $H$ yang dilengkapi topologi lemah ke $H$
yang dilengkapi topologi lemah.

\begin{exer} Misalkan $H$ ruang Hilbert dan $(x_{{}_{\sst\lambda}})$ jaring dalam $H$.
 \begin{enumerate}
  \item Jika $x_{{}_{\sst\lambda}} \to^s  a$ dalam $H$, maka $x_{{}_{\sst\lambda}} \to^w a$ dalam $H$.
  \item Topologi kuat (norma) pada $H$ lebih kuat (lebih besar) daripada topologi lemah.
  \item Tunjukkan dengan contoh bahwa kebalikan dari (a) tidak berlaku.
  \item Apakah norma pada $H$ kontinu secara lemah?
  \item Jika $x_{{}_{\sst\lambda}} \to^w  a$ dalam $H$ dan jika $\norm{x_{{}_{\sst\lambda}}} \to \norm{a}$, maka
$x_{{}_{\sst\lambda}}\to^s a$ dalam~$H$.
  \item Jika pemetaan linear $T\colon H \sto H$ kontinu (secara kuat), maka pemetaan tersebut kontinu secara lemah.
 \end{enumerate}
\end{exer}

\begin{exer} Misalkan $H$ ruang Hilbert berdimensi tak hingga. Tinjau subhimpunan-subhimpunan berikut dari~$H$:
 \begin{enumerate}[label=\rm{(\alph*)}]
  \item bola satuan terbuka;
  \item bola satuan tertutup;
  \item sfer satuan;
  \item subruang linear tertutup.
 \end{enumerate}
Tentukan untuk setiap himpunan tersebut apakah masing-masing bersifat: tertutup secara kuat, tertutup secara lemah,
terbuka secara kuat, terbuka secara lemah, kompak secara kuat, kompak secara lemah. (Salah satu bagian
terlalu sulit untuk saat ini: \emph{Teorema Alaoglu}~\futurexref{6.2.9}{C067441} akan menunjukkan bahwa
bola satuan tertutup kompak secara lemah.)
\end{exer}

\begin{exer} Misalkan $\{e^n\colon n \in \N\}$ basis biasa untuk $l_2$. Untuk
setiap $n \in \N$, misalkan $a_n = \sqrt n\, e^n$. Manakah di antara pernyataan berikut
yang benar?
 \begin{enumerate}
  \item $0$ adalah titik akumulasi lemah dari himpunan
$\{a_n\colon n \in \N\}$.
  \item $0$ adalah titik gugus lemah dari barisan~$(a_n)$.
  \item $0$ adalah limit lemah dari suatu subbarisan dari~$(a_n)$.
 \end{enumerate}
\end{exer}










\section{Morfisme Universal}
Banyak bagian matematika melibatkan pembentukan objek baru dari objek yang sudah ada---misalnya
hasil kali, koproduk, hasil bagi, pelengkapan, kompaktifikasi, dan unitalisasi.
Sering kali suatu konstruksi dapat---dan sangat baik jika dapat---dicirikan melalui
sebuah diagram yang menjelaskan apa yang “dilakukan” oleh objek hasil konstruksi, alih-alih
menjelaskan apa objek itu atau bagaimana objek tersebut dibangun. Diagram semacam ini disebut
 \index{universal!pemetaan!diagram}%
 \index{pemetaan!universal}%
\df{diagram pemetaan universal}, dan diagram itu menjelaskan
 \index{universal!pemetaan!sifat}%
 \index{pemetaan!diagram!universal}%
\df{sifat universal} objek yang dibangun.
Secara khusus, suatu konstruksi biasanya dapat dicirikan oleh adanya
sebuah morfisme tunggal yang memiliki sifat tertentu. Karena morfisme ini beserta
sifat yang bersesuaian dengannya mencirikan secara tunggal konstruksi tersebut, keduanya masing-masing
disebut \emph{morfisme universal} dan \emph{sifat universal}. Definisi berikut
merupakan salah satu cara untuk menjelaskan peran morfisme semacam itu. Jika ini pertama kalinya Anda
menjumpai konsep \emph{universal}, jangan khawatir karena sifatnya memang cukup abstrak. Menurut saya,
tidak banyak orang yang langsung merasa nyaman dengan definisi ini sebelum menjumpai setidaknya
setengah lusin contoh dalam berbagai konteks. (Ada definisi yang bahkan lebih teknis untuk konsep
penting ini, misalnya \emph{objek awal atau objek terminal dalam kategori koma}. Jika tertarik
menelusurinya, carilah “universal property” di Wikipedia~\cite{wiki:xxx}.)

\begin{defn}\label{00078} Misalkan $\cat A \to^{\ftr{F}} \cat B$ suatu funktor antara kategori
$\cat A$ dan $\cat B$, dan misalkan $B$ suatu objek di~$\cat B$. Pasangan $(A,u)$, dengan $A$
suatu objek di $\cat A$ dan $u$ suatu morfisme di $\cat B$ dari $B$ ke~$\ftr F(A)$, disebut
 \index{universal!morfisme}%
 \index{morfisme!universal}%
\df{morfisme universal} untuk $B$ (terhadap funktor~$\ftr F$) jika untuk setiap
objek $A'$ di $\cat A$ dan setiap morfisme $B \to^f \ftr F(A')$ di $\cat B$ terdapat
sebuah morfisme tunggal $A \to^{\tilde f} A'$ di $\cat A$ sehingga diagram berikut komutatif.
 \begin{equation}\label{00078i}
   \xy
     \qtriangle/{>}`{>}`{-->}/[B`\ftr F(A)`\ftr F(A');u`f`\ftr F(\tilde f)]
     \morphism(1000,500)|r|/{-->}/<0,-500>[A`A';\tilde f]
   \endxy
 \end{equation}
Dalam konteks ini, objek $A$ sering disebut
 \index{universal!objek}%
 \index{objek!universal}%
\df{objek universal} di~$\cat A$.
\end{defn}

Berikut adalah contoh pertama sifat semacam itu.

\begin{defn} Misalkan $F$ suatu objek dalam kategori konkret $\cat C$, dan
$\iota\colon S \sto \abs F$ suatu pemetaan inklusi yang domainnya merupakan himpunan tak kosong~$S$.
Kita mengatakan bahwa objek $F$
 \index{bebas!objek}%
 \index{objek!bebas}%
\df{bebas pada} himpunan $S$ (atau bahwa $F$ adalah \df{objek bebas yang dibangkitkan oleh}~$S$)
jika untuk setiap objek $A$ di $\cat C$ dan setiap pemetaan $f\colon S \sto \abs A$ terdapat
sebuah morfisme tunggal $\widetilde f_\iota\colon F \sto A$ di $\cat C$ sehingga
$\abs{\wt f_\iota} \circ \iota = f$.
 \[\xy
     \qtriangle/{>}`{>}`{-->}/[S`\abs F`\abs A;\iota`f`\abs{\wt f_\iota}]
     \morphism(1000,500)|r|/{-->}/<0,-500>[F`A;\widetilde f_\iota]
   \endxy\]
Kita akan menaruh perhatian pada
 \index{bebas!ruang vektor}%
 \index{ruang!vektor!bebas}%
\df{ruang vektor bebas}, yaitu objek bebas dalam kategori $\cat{VEC}$ yang objeknya adalah
ruang vektor dan morfismenya adalah pemetaan linear. Tentu saja, sekadar
\emph{mendefinisikan} suatu konsep tidak menjamin keberadaannya. Ternyata ruang vektor bebas
memang ada pada setiap himpunan. (Lihat contoh~\ref{ump001z}.)
\end{defn}

\vskip 5 pt

\begin{exer} Dalam definisi sebelumnya, rujukan kepada funktor pelupa sering
dihilangkan dan diagram yang menyertainya sering digambar sebagai berikut:
 \[\xy
     \qtriangle/{>}`{>}`{-->}/[S`F`A;\iota`f`\wt f_\iota]
   \endxy\]
Diagram ini tentu terlihat jauh lebih sederhana. Menurut Anda, mengapa saya memilih versi yang
lebih rumit?
\end{exer}

\vskip 5 pt

\begin{prop} Jika dua objek dalam suatu kategori konkret bebas pada himpunan yang sama,
maka kedua objek tersebut isomorfik.
\end{prop}

\vskip 5 pt

\begin{defn} Misalkan $A$ subhimpunan dari himpunan tak kosong~$S$, dan misalkan $\K$
adalah $\R$ atau~$\C$. Definisikan $\sbsb\chi A \colon S \sto \K$, yaitu
 \index{karakteristik!fungsi}%
 \index{<@$\sbsb{\chi}{A}$ (fungsi karakteristik dari $A$)}%
 \index{fungsi!karakteristik}%
\df{fungsi karakteristik} dari~$A$, dengan
  \[ \sbsb{\chi}{A}(x) = \begin{cases}
                                  1, &\text{jika $x \in A$} \\
                                  0, &\text{selain itu}
                              \end{cases} \]
\end{defn}

\begin{exam}\label{ump001z} Jika $S$ adalah sebarang himpunan tak kosong,
maka terdapat ruang vektor $V$ atas~$\K$ yang bebas pada~$S$. Ruang vektor ini
unik (hingga isomorfisme).
\end{exam}

\begin{proof}[\emph{Petunjuk pembuktian}] Untuk himpunan $S$, ambil $V$ sebagai himpunan
semua fungsi bernilai $\K$ pada $S$ yang bertumpuan hingga. Definisikan penjumlahan dan
perkalian skalar secara titik demi titik. Pemetaan $\iota\colon s \mapsto \sbsb \chi {\{s\}}$
yang membawa setiap unsur $s \in S$ ke fungsi karakteristik dari~$\{s\}$ adalah injeksi
yang diinginkan. Untuk membuktikan bahwa $V$ bebas pada $S$, tunjukkan bahwa untuk setiap
ruang vektor $W$ dan setiap fungsi \hbox{$S~\to^{f}~\abs W$} terdapat pemetaan linear
tunggal $V \to^{\wt f} W$ yang membuat diagram berikut komutatif.
 \[\xy
     \qtriangle/{>}`{>}`{-->}/[S`\abs V`\abs W;\iota`f`\abs{\wt f}]
     \morphism(1100,500)|r|/{-->}/<0,-500>[V`W;\wt f]
   \endxy\]
\ns \end{proof}

\begin{prop} Setiap ruang vektor adalah bebas.
\end{prop}

\begin{proof}[\emph{Petunjuk pembuktian}] Tentu saja, sebagian dari persoalannya adalah
menentukan himpunan $S$ yang membuat ruang vektor tersebut bebas. \ns
\end{proof}


\begin{exam} Misalkan $S$ suatu himpunan tak kosong dan $S' = S \cup \{\vc 1\}$, dengan
$\vc 1$ suatu unsur yang tidak termasuk dalam~$S$. Sebuah
 \index{kata}%
\df{kata} dalam bahasa~$S$ adalah suatu barisan $s$ berunsur di $S'$ yang akhirnya selalu
bernilai $\vc 1$ dan memenuhi syarat: jika $s_k = \vc 1$, maka $s_{k+1} = \vc 1$.
Barisan konstan $(\vc 1, \vc 1, \vc 1, \dots)$ disebut
 \index{kata kosong}%
 \index{kata!kosong}%
\df{kata kosong}. Misalkan $F$ adalah himpunan semua kata dalam bahasa~$S$. Andaikan
$s = (s_1, \dots, s_m, \vc 1, \vc 1, \dots)$ dan
$t = (t_1, \dots, t_n, \vc 1, \vc 1, \dots)$ adalah kata (dengan
$s_1$, \dots, $s_m$, $t_1$, \dots, $t_n \in S$). Definisikan
   \[ s \ast t := (s_1, \dots, s_m, t_1, \dots, t_n, \vc 1, \vc 1, \dots). \]
Operasi ini disebut
 \index{konkatenasi}%
\df{konkatenasi}. Mudah dilihat bahwa himpunan $F$ dengan operasi konkatenasi merupakan
 \index{monoid}%
monoid (semigrup dengan unsur identitas), dengan kata kosong sebagai unsur identitas. Inilah
 \index{bebas!monoid}%
 \index{monoid!bebas}%
\df{monoid bebas} yang dibangkitkan oleh~$S$. Jika kata kosong dikeluarkan, kita memperoleh
 \index{bebas!semigrup}%
 \index{semigrup!bebas}%
\df{semigrup bebas} yang dibangkitkan oleh~$S$. Diagram terkait adalah
 \[\xy
     \qtriangle/{>}`{>}`{-->}/[S`\abs F`\abs G;\iota`f`\abs{\tilde f}]
     \morphism(1100,500)|r|/{-->}/<0,-500>[F`G;\tilde f]
   \endxy\]
di mana $\iota$ adalah injeksi yang jelas $s \mapsto (s,\vc1, \vc 1, \dots)$
(biasanya diperlakukan sebagai pemetaan inklusi), $G$ adalah sebarang semigrup,
$f\colon S \sto G$ adalah sebarang fungsi, dan $\abs{\hphantom{O}}$ adalah funktor pelupa
dari kategori monoid dan homomorfisme (atau kategori semigrup dan homomorfisme) ke kategori
$\cat{SET}$.
\end{exam}

\begin{exam} Ingat kembali penyajian biasa untuk koproduk dua objek dalam suatu kategori
dari definisi~\ref{C015611}. Jika $A_1$ dan $A_2$ adalah dua objek dalam kategori $\cat C$,
maka sebuah
 \index{koproduk}%
\df{koproduk} dari $A_1$ dan $A_2$ adalah tripel $(Q,\iota_1,\iota_2)$, dengan $Q$ suatu
objek di $\cat C$ dan $A_k \to^{\iota_k} Q$ ($k = 1,2$) morfisme di $\cat C$, yang memenuhi
syarat berikut: jika $B$ adalah sebarang objek di $\cat C$ dan
$A_k \to^{f_k} B$ ($k = 1$,$2$) adalah sebarang morfisme di~$\cat C$, maka terdapat sebuah
morfisme tunggal $Q \to^f B$ di $\cat C$ yang membuat diagram berikut komutatif.
  \begin{equation}
    \xy
      \Atrianglepair/<-`<-`<-`>`<-/[B`A_1`Q`A_2;f_1`f`f_2`\iota_1`\iota_2]
    \endxy
   \end{equation}

Sekilas mungkin tidak jelas bahwa konstruksi ini bersifat \emph{universal} dalam arti
definisi~\ref{00078}. Untuk melihatnya, misalkan $\ftr D$ adalah funktor diagonal dari
kategori $\cat C$ ke kategori pasangan $\cat{C^2}$ (lihat contoh~\ref{0007606}).
Andaikan $(Q,\iota_1,\iota_2)$ adalah koproduk objek $A_1$ dan $A_2$ di $\cat C$ dalam
pengertian di atas. Maka $A = (A_1,A_2)$ adalah objek di $\cat{C^2}$,
$A \to^{\iota} \ftr D(Q)$ adalah morfisme di $\cat{C^2}$, dan pasangan
$(Q,\iota)$ bersifat universal dalam arti~\ref{00078}. Diagram yang bersesuaian dengan
diagram~\eqref{00078i} adalah
 \begin{equation}
   \xy
     \qtriangle/{>}`{>}`{-->}/[A`\ftr D(Q)`\ftr D(B);\iota`f`\ftr D(\tilde f)]
     \morphism(1000,500)|r|/{-->}/<0,-500>[Q`B;\tilde f]
   \endxy
 \end{equation}
di mana $B$ adalah sebarang objek di $\cat C$ dan, untuk $k = 1,2$,
$A_k \to^{f_k} B$ adalah sebarang morfisme di $\cat C$, sehingga
$f = (f_1,f_2)$ merupakan morfisme di $\cat{C^2}$.
\end{exam}

\begin{exam} Koproduk dua objek $H$ dan $K$ dalam kategori $\cat{HIL}$, yaitu kategori
 \index{koproduk!dalam $\cat{HIL}$}%
 \index{hil@$\cat{HIL}$!koproduk dalam}%
ruang Hilbert dan pemetaan linear terbatas (dan, secara lebih umum, dalam kategori ruang
hasil kali dalam dan pemetaan linear), adalah jumlah langsung ortogonal eksternalnya
$H \oplus K$ (lihat~\ref{0001504}).
\end{exam}


\begin{exam}\label{00078036} Misalkan $A$ dan $B$ ruang Banach. Pada hasil kali Kartesius
 \index{koproduk!dalam $\cat{BAN_\infty}$}%
 \index{ban@$\cat{BAN_\infty}$!koproduk dalam}%
$A \times B$, definisikan penjumlahan dan perkalian skalar secara titik demi titik. Untuk
setiap $(a,b) \in A \times B$, tetapkan $\norm{(a,b)} = \max\{\norm a,\norm b\}$. Dengan
demikian, $A \times B$ menjadi ruang Banach yang dinotasikan dengan $A \oplus B$ dan
disebut
 \index{jumlah!langsung!ruang Banach}%
\df{jumlah langsung} dari $A$ dan~$B$. Jumlah langsung ini merupakan koproduk dalam
kategori topologis $\cat{BAN_\infty}$ dari ruang Banach, tetapi bukan dalam kategori
geometris yang bersesuaian, yaitu~$\cat{BAN_1}$.
\end{exam}

\begin{exam}\label{00078037} Untuk membangun koproduk dalam kategori geometris
 \index{koproduk!dalam $\cat{BAN_1}$}%
 \index{ban@$\cat{BAN_1}$!koproduk dalam}%
$\cat{BAN_1}$ dari ruang Banach, definisikan operasi ruang vektor secara titik demi titik
pada $A \times B$, tetapi gunakan norma
$\norm{(a,b)}_1 = \norm a + \norm b$ untuk setiap $(a,b) \in A \times B$.
\end{exam}

Hampir setiap konsep dalam teori kategori memiliki konsep dual---konsep yang diperoleh
dengan membalik arah semua panah. Sebagai contoh, kita dapat membalik semua panah dalam
diagram~\eqref{00078i}.


\begin{defn}\label{000781} Misalkan $\cat A \to^{\ftr{F}} \cat B$ suatu funktor antara
kategori $\cat A$ dan $\cat B$, dan misalkan $B$ suatu objek di~$\cat B$. Pasangan $(A,u)$,
dengan $A$ suatu objek di $\cat A$ dan $u$ suatu morfisme di $\cat B$ dari $\ftr F(A)$
ke~$B$, disebut
 \index{kouniversal!morfisme}%
 \index{morfisme!kouniversal}%
\df{morfisme kouniversal untuk $B$ (terhadap~$\ftr F$)} jika untuk setiap objek $A'$ di
$\cat A$ dan setiap morfisme $\ftr F(A') \to^f B$ di $\cat B$ terdapat sebuah morfisme
tunggal $A' \to^{\tilde f} A$ di $\cat A$ sehingga diagram berikut komutatif.
 \begin{equation}\label{000781i}
   \xy
     \qtriangle/{<-}`{<-}`{<--}/[B`\ftr F(A)`\ftr F(A');u`f`\ftr F(\tilde f)]
     \morphism(1000,500)|r|/{<--}/<0,-500>[A`A';\tilde f]
   \endxy
 \end{equation}
Sebagian penulis membalik konvensi ini dan menyebut morfisme dalam~\ref{00078}
\emph{kouniversal}, sedangkan morfisme yang didefinisikan di sini disebut
\emph{universal}. Penulis lain, termasuk penulis buku ini, menyebut keduanya
\emph{morfisme universal}.
\end{defn}

\begin{conv} Morfisme dari kedua jenis yang didefinisikan dalam~\ref{00078}
dan~\ref{000781} akan disebut
 \index{konvensi!tentang universalitas}%
\emph{morfisme universal}.
\end{conv}

\begin{exam}\label{0007812} Ingat kembali pendekatan kategoris terhadap hasil kali dari
definisi~\ref{C015511}. Jika $A_1$ dan $A_2$ adalah dua objek dalam kategori $\cat C$,
maka sebuah
 \index{hasil kali}%
\df{hasil kali} dari $A_1$ dan $A_2$ adalah tripel $(P,\pi_1,\pi_2)$, dengan $P$ suatu
objek di $\cat C$ dan $P \to^{\pi_k} A_k$ ($k = 1,2$) morfisme di $\cat C$, yang memenuhi
syarat berikut: jika $B$ adalah sebarang objek di $\cat C$ dan
$B \to^{f_k} A_k$ ($k = 1,2$) adalah sebarang morfisme di~$\cat C$, maka terdapat sebuah
morfisme tunggal $B \to^f P$ di $\cat C$ yang membuat diagram berikut komutatif.
  \begin{equation}
    \xy
      \Atrianglepair/>`>`>`<-`>/[B`A_1`P`A_2;f_1`f`f_2`\pi_1`\pi_2]
    \endxy
   \end{equation}
Konstruksi ini bersifat (ko)universal dalam arti definisi~\ref{000781}.
\end{exam}

\begin{exam} Hasil kali dua objek $H$ dan $K$ dalam kategori $\cat{HIL}$, yaitu kategori
 \index{hasil kali!dalam $\cat{HIL}$}%
 \index{hil@$\cat{HIL}$!hasil kali dalam}%
ruang Hilbert dan pemetaan linear terbatas (dan, secara lebih umum, dalam kategori ruang
hasil kali dalam dan pemetaan linear), adalah jumlah langsung ortogonal eksternalnya
$H \oplus K$ (lihat~\ref{0001504}).
\end{exam}


\begin{exam} Jika $A$ dan $B$ adalah aljabar Banach, jumlah langsungnya $A \oplus B$
merupakan
 \index{hasil kali!dalam $\cat{BA_\infty}$}%
 \index{ba@$\cat{BA_\infty}$!hasil kali dalam}%
 \index{hasil kali!dalam $\cat{BA_1}$}%
 \index{ba@$\cat{BA_1}$!hasil kali dalam}%
hasil kali dalam kategori topologis maupun geometris dari aljabar Banach, yaitu
$\cat{BA_\infty}$ dan $\cat{BA_1}$. Bandingkan dengan keadaan yang dibahas dalam
contoh~\ref{00078036} dan~\ref{00078037}.
\end{exam}

Konstruksi universal dalam kategori apa pun menghasilkan objek yang unik hingga isomorfisme.

\begin{prop}\label{C035562} Objek universal dalam suatu kategori
 \index{universal!objek!keunikan}%
pada hakikatnya unik.
\end{prop}













\section{Pelengkapan Ruang Hasil Kali Dalam}
Sedikit tinjauan ulang tentang pelengkapan ruang metrik akan berguna di sini. Kadang-kadang
kita mungkin tergoda untuk menganggap bahwa pelengkapan suatu ruang metrik $M$ adalah ruang
metrik lengkap yang memiliki subhimpunan padat yang homeomorfik dengan $M$. Masalah dengan
pencirian ini adalah bahwa pencirian tersebut dapat menghasilkan dua “pelengkapan” ruang
metrik yang bahkan tidak homeomorfik.

\begin{exam} Misalkan $M$ adalah interval $(0,\infty)$ dengan metrik biasa. Terdapat dua
ruang metrik lengkap $N_1$ dan $N_2$ sehingga
  \begin{enumerate}
    \item $N_1$ dan $N_2$ \emph{tidak} homeomorfik,
    \item $M$ homeomorfik dengan suatu subhimpunan padat dari $N_1$, dan
    \item $M$ homeomorfik dengan suatu subhimpunan padat dari~$N_2$.
  \end{enumerate}
\end{exam}

Berikut adalah definisi yang tepat.

\begin{defn}\label{C035511} Misalkan $M$ dan $N$ ruang metrik. Kita mengatakan bahwa $N$
adalah
 \index{pelengkapan!ruang metrik}%
 \index{metrik!ruang!pelengkapan}%
\df{pelengkapan} dari $M$ jika $N$ lengkap dan $M$ isometrik dengan suatu subhimpunan
padat dari~$N$.
\end{defn}

Dalam konstruksi bilangan real dari bilangan rasional, salah satu pendekatan adalah
memperlakukan himpunan bilangan rasional sebagai ruang metrik dan mendefinisikan himpunan
bilangan real sebagai pelengkapannya. Salah satu cara baku untuk menghasilkan pelengkapan
memakai kelas-kelas ekuivalensi barisan Cauchy bilangan rasional. Teknik ini berlaku pula
untuk ruang metrik: cara elementer untuk melengkapi sebarang ruang metrik dimulai dengan
mendefinisikan relasi ekuivalensi pada himpunan barisan Cauchy yang berunsur di ruang
tersebut. Berikut adalah pendekatan yang sedikit kurang elementer, tetapi jauh lebih sederhana.

\begin{prop} Setiap ruang metrik $M$ isometrik dengan suatu subruang dari~$\fml B(M)$.
\end{prop}

\begin{proof}[\emph{Petunjuk pembuktian}] Jika $(M,d)$ adalah ruang metrik, tetapkan
$a \in M$. Untuk setiap $x \in M$, definisikan $\phi_x\colon M \sto \R$ dengan
$\phi_x(u) = d(x,u) - d(u,a)$. Tunjukkan terlebih dahulu bahwa
$\phi_x \in \fml B(M)$ untuk setiap $x \in M$. Kemudian tunjukkan bahwa
$\phi\colon M \sto \fml B(M)$ adalah isometri. \ns
\end{proof}

\begin{cor}\label{C035521} Setiap ruang metrik mempunyai pelengkapan.
\end{cor}

Proposisi berikut menunjukkan bahwa pelengkapan ruang metrik bersifat universal dalam
arti definisi~\ref{00078}. Kategori yang dimaksud dalam hasil ini adalah kategori semua
ruang metrik dan pemetaan kontinu seragam, beserta subkategorinya yang terdiri atas semua
ruang metrik lengkap dan pemetaan kontinu seragam. Di sini, funktor pelupa
$\abs{\hphantom{O}}$ hanya melupakan kelengkapan objek (dan tidak mengubah morfisme).

\begin{prop}\label{C035557} Misalkan $M$ ruang metrik dan $\clo M$ pelengkapannya. Jika
$N$ ruang metrik lengkap dan $f\colon M \sto N$ kontinu seragam, maka terdapat sebuah
fungsi kontinu seragam tunggal $\clo f\colon \clo M \sto N$ yang membuat diagram berikut
komutatif.
 \[\xy
     \qtriangle/{>}`{>}`{-->}/[M`\abs{\clo M}`\abs N;\iota`f`\abs{\clo f}]
     \morphism(1100,500)|r|/{-->}/<0,-500>[\clo M`N;\clo f]
   \endxy\]
\end{prop}

Akibat berikut dari proposisi~\ref{C035557} dan~\ref{C035562} memungkinkan kita berbicara
tentang \emph{pelengkapan} suatu ruang metrik.

\begin{cor}\label{C035564} Pelengkapan ruang metrik unik (hingga isometri).
\end{cor}

Sekarang, bagaimana dengan pelengkapan ruang hasil kali dalam? Seperti telah kita lihat,
setiap ruang hasil kali dalam adalah ruang metrik, sehingga ruang itu mempunyai pelengkapan
(sebagai ruang metrik). Pertanyaannya: apakah pelengkapan ini merupakan ruang Hilbert?
Jawabannya, dengan senang hati, adalah \emph{ya}.

\begin{thm} Misalkan $V$ ruang hasil kali dalam dan $H$ pelengkapannya sebagai ruang
metrik. Maka terdapat hasil kali dalam pada $H$ yang
  \begin{enumerate}
    \item merupakan perpanjangan dari hasil kali dalam pada~$V$, dan
    \item menginduksi metrik pada~$H$.
  \end{enumerate}
\end{thm}


















\endinput

 \chapter{OPERATOR PADA RUANG HILBERT}

\section{Pemetaan Linear Invertibel dan Isometri}

\begin{defn} Suatu pemetaan linear $T\colon H \sto K$ antara dua ruang hasil kali dalam dikatakan
 \index{hasil kali dalam!mempertahankan}%
\df{mempertahankan hasil kali dalam} jika $\langle Tx, Ty \rangle = \langle x, y \rangle$ untuk semua $x$, $y \in H$.
\end{defn}

\begin{prop}\label{prop_isometry_equiv} Misalkan $T\colon H \sto K$ suatu pemetaan linear antara ruang-ruang hasil kali dalam. Maka
pernyataan-pernyataan berikut ekuivalen:
 \begin{enumerate}[label=\rm{(\alph*)}]
   \item $T$ mempertahankan hasil kali dalam;
   \item $T$ mempertahankan norma; dan
   \item $T$ merupakan isometri.
 \end{enumerate}
\end{prop}

Sama seperti pada ruang linear bernorma (lihat bagian paling akhir dari seksi~\ref{sec_bdd_lin_maps}),
terdapat (sekurang-kurangnya) dua kategori penting yang objek-objeknya adalah ruang Hilbert:
 \begin{enumerate}
  \item[(a)] $\cat{HSp} =$ ruang Hilbert dan pemetaan linear terbatas, dan
   \index{HSp@$\cat{HSp}$!kategori}%
   \index{kategori!$\cat{HSp}$ sebagai suatu}%
  \item[(b)] $\cat{HSp_1} =$ ruang Hilbert dan kontraksi linear.
   \index{HSp1@$\cat{HSp_1}$!kategori}%
   \index{kategori!$\cat{HSp_1}$ sebagai suatu}%
 \end{enumerate}
Dalam kategori apa pun, lazim untuk mendefinisikan \emph{isomorfisme} sebagai \emph{morfisme invertibel}. Karena hampir selalu
 \index{kategori!topologis}%
kategori (topologis) ruang Hilbert dan pemetaan linear terbataslah yang kita jumpai
dalam catatan ini, tampaknya masuk akal untuk menerapkan kata ``isomorfisme'' pada
isomorfisme dalam kategori ini---dengan kata lain, pada pemetaan invertibel; sedangkan isomorfisme
dalam kategori yang lebih ketat, yaitu
 \index{kategori!geometris}%
kategori (geometris) ruang Hilbert dan kontraksi linear, dapat secara wajar
disebut \emph{isomorfisme isometrik}. Namun, ini \emph{bukan} konvensi yang dipakai. Ketika kebanyakan matematikawan memikirkan
``isomorfisme ruang Hilbert'', yang mereka maksud adalah pemetaan yang mempertahankan seluruh struktur ruang Hilbert---termasuk
hasil kali dalam. Dengan demikian (menggunakan~\ref{prop_isometry_equiv}), kata
 \index{isomorfisme!antara ruang Hilbert}%
 \index{isomorfisme!dalam kategori $\cat{HSp}$}%
 \index{HSp@$\cat{HSp}$!isomorfisme dalam}%
 \index{isomorfisme!dalam kategori $\cat{HSp_1}$}%
 \index{HSp1@$\cat{HSp_1}$!isomorfisme dalam}%
``isomorfisme'', apabila diterapkan pada pemetaan antara ruang-ruang Hilbert, di hampir semua tempat berarti \emph{isomorfisme isometrik}.
Akibatnya, isomorfisme dalam kategori yang lebih umum $\cat{HSp}$ disebut \emph{pemetaan (linear terbatas) invertibel}.
Ingat pula: Kita mencadangkan kata
 \index{operator}%
 \index{konvensi!semua operator terbatas dan linear}%
``operator'' bagi pemetaan linear terbatas dari suatu ruang \emph{ke ruang itu sendiri}.

\begin{exam} Definisikan suatu operator $T$ pada ruang Hilbert $H = L_2([0,\infty))$ dengan
   \[ Tf(t) = \begin{cases}   f(t-1),  &\text{ jika $t \ge 1$} \\
                                   0,  &\text{ jika $0 \le t < 1$.}
           \end{cases} \]
Maka $T$ merupakan isometri, tetapi bukan isomorfisme isometrik. Sebaliknya, jika
$H = L_2(\R)$ dan $T$ memetakan setiap $f \in H$ ke fungsi $t \mapsto f(t-1)$, maka $T$ merupakan isomorfisme isometrik.
\end{exam}

\begin{exam} Misalkan $H$ ruang Hilbert yang didefinisikan dalam contoh~\ref{exam_Hsp_abs_cont}. Maka pemetaan diferensiasi
$D\colon f \mapsto f\,'$ dari $H$ ke $L_2([0,1])$ merupakan isomorfisme isometrik. Apa inversnya?
\end{exam}

\begin{exam} Misalkan $\{(X_i,\mu_i) \colon i \in I\}$ suatu keluarga ruang ukur $\sigma$-hingga. Jadikan gabungan saling lepas
$\D X = \biguplus_{i \in I}X_i$ sebagai ruang ukur $(X,\mu)$ dengan cara biasa. Maka $L_2(X,\mu)$ isomorfik secara isometrik
dengan jumlah langsung ruang-ruang Hilbert $L_2(X_i,\mu_i)$.
\end{exam}

\begin{exam} Misalkan $\mu$ ukuran Lebesgue pada $[-\frac\pi2, \frac\pi2]$ dan $\nu$ ukuran Lebesgue pada $\R$.
Misalkan $(Uf)(x) = \dfrac{f(\arctan x)}{\sqrt{1+x^2}}$ untuk semua $f \in L_2(\mu)$ dan semua $x \in \R$. Maka $U$ merupakan
isomorfisme isometrik antara $L_2([-\frac\pi2, \frac\pi2], \mu)$ dan $L_2(\R, \nu)$.
\end{exam}

\begin{prop} Setiap surjeksi yang mempertahankan hasil kali dalam dari suatu ruang Hilbert ke ruang Hilbert lainnya secara otomatis linear.
\end{prop}

\begin{prop}\label{prop_norm_lin_map} Jika $T\colon H \sto K$ merupakan pemetaan linear terbatas antara ruang-ruang Hilbert tak nol, maka
   \[ \norm T = \sup\{\abs{\langle Tx,y \rangle}\colon \norm x = \norm y = 1\}\,. \]
\end{prop}

\begin{proof}[\emph{Petunjuk pembuktian}] Tunjukkan terlebih dahulu bahwa $\norm x = \sup\{\abs{\langle x,y \rangle}\colon \norm y = 1\}$. \ns
\end{proof}

\begin{defn} Suatu
 \index{kurva}%
\df{kurva} dalam ruang Hilbert $H$ adalah pemetaan kontinu dari $[0,1]$ ke~$H$. Suatu kurva dikatakan
 \index{sederhana!kurva}%
\df{sederhana} jika kurva itu injektif. Jika $c$ suatu kurva dan $0 \le a < b \le 1$, maka
 \index{tali busur}%
\df{tali busur} dari $c$ yang bersesuaian dengan interval $[a,b]$ adalah vektor $c(b) - c(a)$. Dua tali busur dikatakan
\df{tak bertumpang tindih} jika interval parameter yang terkait dengannya paling banyak memiliki satu titik ujung bersama.
\end{defn}

Dalam \cite{Halmos:1982}, Halmos menggambarkan betapa luasnya ruang Hilbert berdimensi tak hingga melalui
contoh elegan berikut.

\begin{exam} Dalam setiap ruang Hilbert berdimensi tak hingga terdapat kurva sederhana yang berbelok membentuk sudut siku-siku
di setiap titik, dalam arti bahwa setiap pasangan tali busur yang tak bertumpang tindih saling tegak lurus.
\end{exam}

\begin{proof}[\emph{Petunjuk pembuktian}] Tinjau fungsi-fungsi karakteristik dalam $L_2([0,1])$. \ns
\end{proof}
























\section{Operator dan Adjoinnya}

\begin{prop}\label{C0635103} Misalkan $S$, $T \in \ofml B(H,K)$, dengan $H$ dan $K$ ruang
Hilbert. Jika $\langle Sx,y \rangle = \langle Tx,y \rangle$ untuk setiap $x \in H$ dan $y \in
K$, maka $S = T$.
\end{prop}

\begin{defn} Misalkan $T$ suatu operator pada ruang Hilbert~$H$. Maka $Q_T$, yaitu
 \index{bentuk kuadratik}%
 \index{bentuk!kuadratik}%
 \index{Q@$Q_T$ (bentuk kuadratik yang terkait dengan~$T$)}%
\df{bentuk kuadratik yang terkait dengan~$T$}, adalah fungsi bernilai skalar yang didefinisikan oleh
   \[ Q_T(x) := \langle Tx,x \rangle \]
untuk setiap $x \in H$.
\end{defn}

Suatu operator pada ruang Hilbert kompleks adalah operator nol jika dan hanya jika bentuk kuadratik yang terkait dengannya adalah fungsi nol.

\begin{prop}\label{C0635105} Jika $H$ adalah ruang Hilbert \emph{kompleks} dan $T \in \ofml B(H)$, maka
$T = \vc 0$ jika dan hanya jika $Q_T = 0$.
\end{prop}

\begin{proof}[\emph{Petunjuk pembuktian}] Dalam hipotesis $Q_T(z) = 0$ untuk setiap $z \in H$, gantilah $z$ mula-mula dengan $y + x$
dan kemudian dengan $y + ix$.  \ns
\end{proof}

Proposisi sebelumnya merupakan salah satu dari sedikit proposisi yang akan kita jumpai yang tidak berlaku sekaligus untuk ruang Hilbert real dan kompleks.

\begin{exam} Proposisi~\ref{C0635105} tidak benar jika dalam hipotesis kata ``kompleks'' diganti dengan ``real''.
\end{exam}

\begin{defn} Suatu fungsi $T\colon V \sto W$ antara ruang vektor kompleks disebut
 \index{konjugat!linear}%
 \index{linear!konjugat}%
\df{linear konjugat} jika fungsi itu aditif ($T(x+y) = Tx + Ty$) dan memenuhi $T(\alpha x) =
\conj\alpha\, Tx$ untuk semua $x \in V$ dan $\alpha \in \C$. Suatu fungsi bernilai kompleks dengan dua
variabel $\phi\colon V \times W \sto \C$ disebut
 \index{fungsional seskuilinear}%
 \index{fungsional!seskuilinear}%
\df{fungsional seskuilinear} jika fungsi itu linear dalam variabel pertamanya dan linear konjugat
dalam variabel keduanya.

Jika $H$ dan $K$ adalah ruang linear bernorma, suatu fungsional seskuilinear $\phi$ pada $H \times K$
disebut
 \index{terbatas!fungsional seskuilinear}%
 \index{fungsional seskuilinear!terbatas}%
 \index{fungsional!seskuilinear terbatas}%
\df{terbatas} jika terdapat suatu konstanta $M > 0$ sedemikian sehingga
    \[ \abs{\phi(x,y)} \le M\norm x\norm y \]
untuk semua $x \in H$ dan $y \in K$.
\end{defn}

\begin{prop} Jika $\phi \colon H \oplus K \sto \C$ adalah fungsional seskuilinear terbatas pada
jumlah langsung dua ruang Hilbert tak nol, maka bilangan-bilangan berikut semuanya ada dan sama:
 \begin{enumerate}[label=\rm{(\alph*)}]
    \item $\sup\{\abs{\phi(x,y)}\colon \norm x \le 1, \norm y \le 1\}$
 \vskip 3 pt
    \item $\sup\{\abs{\phi(x,y)}\colon \norm x = \norm y = 1\}$
 \vskip 3 pt
    \item $\sup\left\{\dfrac{\abs{\phi(x,y)}}{\norm x\,\norm y}\colon x,y \ne 0\right\}$
 \vskip 3 pt
    \item $\inf\{M > 0\colon \abs{\phi(x,y)} \le M\norm x\,\norm y
                    \text{ untuk semua $x \in H$, $y \in K$}\}$.
 \end{enumerate}
\end{prop}

\begin{proof}[\emph{Petunjuk pembuktian}] Pembuktiannya hampir sama persis dengan hasil yang bersesuaian
untuk pemetaan linear (lihat~\ref{C023221}). \ns
\end{proof}

\begin{defn} Misalkan $\phi \colon H \oplus K \sto \C$ fungsional seskuilinear terbatas pada
jumlah langsung dua ruang Hilbert tak nol. Kita mendefinisikan $\norm \phi$, yaitu
 \index{norma!fungsional seskuilinear terbatas}%
\emph{norma} dari $\phi$, sebagai salah satu di antara ekspresi-ekspresi (yang sama) dalam hasil sebelumnya.
\end{defn}

\begin{prop} Misalkan $T\colon H \sto K$ pemetaan linear terbatas antara ruang Hilbert tak nol. Maka
   \[ \phi \colon H \oplus K \sto \C\colon (x,y) \mapsto \langle Tx,y \rangle \]
adalah fungsional seskuilinear terbatas pada $H \oplus K$ dan $\norm\phi = \norm T$.
\end{prop}

\begin{prop} Misalkan $\phi \colon H \times K \sto \C$ fungsional seskuilinear terbatas pada
hasil kali dua ruang Hilbert tak nol. Maka terdapat pemetaan linear terbatas yang unik $T \in \ofml B(H,K)$ dan
$S \in \ofml B(K,H)$ sedemikian sehingga
    \[ \phi(x,y) = \langle Tx,y \rangle = \langle x,Sy \rangle \]
untuk semua $x \in H$ dan $y \in K$. Selain itu, $\norm T = \norm S = \norm\phi$.
\end{prop}

\begin{proof}[\emph{Petunjuk pembuktian}] Tunjukkan bahwa untuk setiap $x \in H$, pemetaan
$y \mapsto \conj{\phi(x,y)}$ adalah fungsional linear terbatas pada~$K$. \ns
\end{proof}

\begin{defn}\label{def000926} Misalkan $T\colon H \sto K$ pemetaan linear terbatas antara ruang-ruang Hilbert. Pemetaan
$(x,y) \mapsto \langle Tx,y \rangle$ dari $H \oplus K$ ke dalam $\C$ adalah fungsional
seskuilinear terbatas. Berdasarkan proposisi sebelumnya, terdapat pemetaan linear terbatas yang
unik
 \index{adjoin!pemetaan linear!antara ruang Hilbert}%
 \index{linear!pemetaan!adjoin}%
 \index{<unopurstar@$T^*$ (adjoin ruang Hilbert)}%
$T^*\colon K \sto H$ yang disebut \df{adjoin} dari $T$ sedemikian sehingga
   \[ \langle Tx,y \rangle = \langle x,T^*y \rangle \]
untuk semua $x \in H$ dan $y \in K$.
\end{defn}

\begin{exam}\label{C063526} Ingat dari Contoh~\ref{exam150013} bahwa keluarga $l_2$ dari semua
barisan bilangan kompleks yang kuadratnya dapat dijumlahkan merupakan ruang Hilbert. Misalkan
    \[ S\colon l_2 \sto l_2\colon (x_1,x_2,x_3,\dots) \mapsto (0,x_1,x_2,\dots)\,. \]
Maka $S$ adalah operator pada $l_2$, yang disebut
 \index{geser unilateral}%
 \index{S@$S$ (operator geser unilateral)}%
 \index{operator!geser unilateral}%
\df{operator geser unilateral}, dan $\norm S = 1$. Adjoin $S^*$ dari operator geser
unilateral
 \index{geser unilateral!adjoin}%
 \index{adjoin!operator geser unilateral}%
 \index{operator!geser unilateral!adjoin}%
diberikan oleh
  \[ S^*\colon l_2 \sto l_2\colon (x_1,x_2,x_3,\dots) \mapsto (x_2,x_3,x_4,\dots). \]
\end{exam}

\begin{exam} Ambil $\bigl(e^n\bigr)_{n=1}^\infty$ sebagai basis ortonormal bagi ruang Hilbert separabel~$H$ dan
$\bigl(\alpha_n\bigr)_{n=1}^\infty$ barisan terbatas bilangan kompleks. Misalkan $Te^n = \alpha_n e^n$ untuk setiap~$n$.
Maka $T$ dapat diperluas secara unik menjadi operator pada~$H$. Berapakah norma operator ini? Apa adjoinnya?
\end{exam}

\begin{defn} Operator yang dibahas dalam contoh sebelumnya dikenal sebagai
 \index{operator!diagonal}%
 \index{diagonal!operator}%
\df{operator diagonal}. Notasi yang lazim untuk operator ini adalah
 \index{diag@$\diag(\alpha_1,\alpha_2,\dots)$ (operator diagonal)}%
$\diag(\alpha_1,\alpha_2,\dots)$.
\end{defn}

Dalam kategori ruang hasil kali dalam dan pemetaan linear terbatas, baik kernel maupun jangkauan suatu morfisme merupakan objek dalam kategori tersebut.
Namun, seperti ditunjukkan oleh contoh berikut, hal ini tidak harus berlaku untuk ruang Hilbert. Walaupun operator pada ruang Hilbert selalu mempunyai kernel
yang sendiri merupakan ruang Hilbert, jangkauannya mungkin tidak tertutup dan akibatnya tidak lengkap.

\begin{exam}\label{exam_ran_nonclosed} Jangkauan operator diagonal $\diag\bigl(1,\frac12,\frac13, \dots\bigr)$ pada
 \index{tertutup!jangkauan!operator tanpa}%
 \index{operator!tanpa jangkauan tertutup}%
ruang Hilbert $l_2$ tidak tertutup.
\end{exam}

\begin{exam}\label{C063527} Misalkan $(S,\fml A, \mu)$ ruang ukur positif yang sigma-hingga dan $L_2(S,\mu)$
ruang Hilbert dari semua (kelas ekuivalensi) fungsi bernilai kompleks pada $S$
yang terintegralkan kuadrat terhadap~$\mu$. Misalkan $\phi$ fungsi bernilai kompleks yang
terbatas secara esensial dan terukur terhadap $\mu$ pada~$S$. Definisikan
 \index{M@$M_\phi$ (operator perkalian)}%
$M_\phi$ pada $L_2(S,\mu)$ dengan $M_\phi(f) := \phi f$. Maka $M_\phi$ adalah operator pada $L_2(S,\mu)$; operator ini disebut
 \index{operator perkalian!pada $L_2(S,\mu)$}%
 \index{operator!perkalian!pada $L_2(S,\mu)$}%
\df{operator perkalian}. Normanya diberikan oleh $\norm{M_\phi} = \norm\phi_\infty$ dan
 \index{operator perkalian!adjoin}%
 \index{operator!perkalian!adjoin}%
 \index{adjoin!operator perkalian}%
adjoinnya oleh ${M_\phi}^*~=~M_{\conj\phi}$.
\end{exam}

Akan berguna untuk sedikit memperluas Definisi~\ref{0001612}.

\begin{defn} Operator $S$ dan $T$ masing-masing pada ruang Hilbert $H$ dan $K$ disebut
 \index{uniter!ekuivalensi}%
 \index{ekuivalen!secara uniter}%
\df{ekuivalen secara uniter} jika terdapat isomorfisme isometrik $U\colon H \sto K$ sedemikian sehingga $T = USU^{-1}$.

\end{defn}

\begin{exam} Misalkan $\mu$ ukuran Lebesgue pada $[-\frac\pi2, \frac\pi2]$ dan $\nu$ ukuran Lebesgue pada $\R$.
Misalkan $\phi(x) = x$ untuk $x \in [-\frac\pi2, \frac\pi2]$ dan $\psi(x) = \arctan x$ untuk $x \in \R$. Maka operator
perkalian $M_\phi$ dan $M_\psi$ (masing-masing pada $L_2([-\frac\pi2, \frac\pi2], \mu)$ dan $L_2(\R, \nu)$)
ekuivalen secara uniter.
\end{exam}

\begin{exam} Misalkan $M_\phi$ operator perkalian pada $L_2(S,\mu)$, dengan $(S,\fml A,\mu)$
ruang ukuran $\sigma$-hingga seperti dalam Contoh~\ref{C063527}. Maka $M_\phi$ idempoten jika dan hanya jika
$\phi$ adalah fungsi karakteristik dari suatu himpunan dalam~$\fml A$.
\end{exam}

Kita dapat memandang
 \index{spektral!teorema}%
\df{teorema spektral} dalam aljabar linear dasar sebagai pernyataan berikut: \textbf{Setiap operator normal pada ruang hasil kali dalam kompleks
berdimensi hingga ekuivalen secara uniter dengan suatu operator perkalian.} Yang benar-benar
menakjubkan adalah bagaimana pernyataan ini digeneralisasi ke ruang Hilbert berdimensi tak hingga---pokok bahasan
yang akan dibicarakan jauh kemudian.

\begin{exer} Misalkan $\N_3 = \{1,2,3\}$ dan $\mu$ ukuran pencacahan pada $\N_3$.
 \begin{enumerate}
  \item[(a)] Identifikasilah ruang Hilbert~$L_2(\N_3,\mu)$.
  \item[(b)] Identifikasilah operator perkalian pada~$L_2(\N_3,\mu)$.
  \item[(c)] Misalkan $T$ operator linear pada $\C^3$ yang representasi
matriksnya adalah
   \[ \begin{bmatrix}
          2  &  0  &  1 \\
          0  &  2  & -1 \\
          1  & -1  &  1
      \end{bmatrix}.\]
Gunakan operator $T$ untuk mengilustrasikan rumusan teorema spektral sebelumnya.
 \end{enumerate}
\end{exer}

\begin{exam}\label{exam_int_op} Misalkan $L_2 = L_2([0,1],\lambda)$ ruang Hilbert dari semua (kelas ekuivalensi)
fungsi bernilai kompleks pada $[0,1]$ yang terintegralkan kuadrat terhadap ukuran Lebesgue~$\lambda$.
Jika $k\colon [0,1] \times [0,1] \sto \C$ fungsi terukur Borel yang terbatas, definisikan $K$ pada $L_2$ dengan
   \[ (Kf)(x) = \int_0^1 k(x,y)f(y)\,dy\,. \]
Maka $K$ adalah
 \index{integral!operator}%
 \index{operator!integral}%
\df{operator integral} dan
 \index{kernel!operator integral}%
 \index{integral!operator!kernel}%
\df{kernel} dari operator tersebut adalah fungsi~$k$. (Ini merupakan penggunaan lain dari kata ``kernel''; penggunaan ini sama sekali tidak berkaitan
dengan penggunaan kata yang lebih umum: $\ker K = K^{\gets}(\{\vc 0\})$---lihat Definisi~\ref{defn_op_vsp}\,).
Adjoin $K^*$ dari operator $K$ juga merupakan operator integral pada $L_2$ dan
 \index{adjoin!operator integral}%
 \index{integral!operator!adjoin}%
 \index{operator!integral!adjoin}%
kernelnya adalah fungsi $k^*$ yang didefinisikan pada $[0,1] \times [0,1]$ oleh $k^*(x,y) = \conj{k(y,x)}$.
\end{exam}

\begin{proof}[\emph{Petunjuk pembuktian}] Verifikasilah bahwa untuk $f \in L_2$
   \[ \abs{(Kf)(x)} \le \norm{k}_\infty^{1/2}\left[ \int_0^1 \abs{k(x,y)}\abs{f(y)}^2\,dy\right]^{1/2}.\]
\ns
\end{proof}

\begin{exam}\label{000319} Misalkan $H = L_2([0,1])$ ruang Hilbert real dari semua (kelas ekuivalensi) fungsi
bernilai real pada $[0,1]$ yang terintegralkan kuadrat terhadap ukuran Lebesgue. Definisikan $V$ pada $H$ dengan
   \begin{equation}\label{000319i}
       Vf(x) = \int_0^x f(t)\,dt\,.
   \end{equation}
Maka $V$ merupakan operator pada~$H$. Operator ini adalah
 \index{operator Volterra}%
 \index{operator!Volterra}%
\df{operator Volterra} dan merupakan contoh operator integral. (Apa kernel~$k$\, untuk operator ini?)
\end{exam}

\begin{exer} Misalkan $V$ adalah \emph{operator Volterra} yang didefinisikan di atas.
 \begin{enumerate}
  \item[(a)] Tunjukkan bahwa $V$ injektif.
  \item[(b)] Hitunglah $V^*$.
  \item[(c)] Tentukan $V + V^*$ beserta jangkauannya.
 \end{enumerate}
\end{exer}

\begin{prop}\label{prop_adj_invol} Jika $S$ dan $T$ adalah operator pada ruang Hilbert kompleks~$H$ dan $\alpha \in \C$, maka
 \begin{enumerate}[label=\rm{(\alph*)}]
  \item $(S + T)^* = S^* + T^*$;
  \item $(\alpha T)^* = \conj\alpha T^*$;
  \item $T^{**} = T$; dan
  \item $(TS)^* = S^*T^*$.
 \end{enumerate}
\end{prop}

\begin{prop}\label{C063544} Misalkan $T$ operator pada ruang Hilbert~$H$. Maka
   \begin{enumerate}[label=\rm{(\alph*)}]
     \item $\norm{T^*} = \norm T$ dan
     \item $\norm{T^*T} = {\norm T}^2$.
   \end{enumerate}
\end{prop}

\begin{proof}[\emph{Petunjuk pembuktian}] Perhatikan bahwa karena $T = T^{**}$, untuk membuktikan (a) cukup ditunjukkan bahwa $\norm T \le \norm{T^*}$.
Untuk itu, amati bahwa ${\norm{Tx}}^2 = \langle T^*Tx,x \rangle \le \norm{T^*}\norm T$ apabila $x \in H$ dan $\norm x \le 1$. Dari sini
simpulkan bahwa ${\norm T}^2 \le \norm{T^*T} \le \norm{T^*}\norm T$.  \ns
\end{proof}

Syarat (b) yang tampak cukup biasa dalam proposisi sebelumnya kelak ternyata menjadi \emph{sifat} yang sangat penting
ketika aljabar-$C^*$ pertama kali diperkenalkan dalam Bab~\futurexref{7}{chap_cpt_ops}.

\begin{prop}\label{prop_ker_adj_perp_ran_op} Jika $T$ adalah operator pada ruang Hilbert, maka
 \begin{enumerate}[label=\rm{(\alph*)}]
   \item $\ker T^* = (\ran T)^\perp$,
   \item $\clo{\ran T^*} = (\ker T)^\perp$,
   \item $\ker T = (\ran T^*)^\perp$, dan
   \item $\clo{\ran T} = (\ker T^*)^\perp$.
 \end{enumerate}
\end{prop}

Bandingkan hasil sebelumnya dengan Teorema~\ref{00016025}.

\begin{defn}\label{def_bdd_away_zero} Suatu pemetaan linear terbatas $T \colon V \sto W$ antara ruang linear bernorma disebut
 \index{terbatas!jauh dari nol}%
 \index{nol!terbatas jauh dari}%
\df{terbatas jauh dari nol} (atau
 \index{terbatas!dari bawah}%
\df{terbatas dari bawah}) jika terdapat suatu bilangan $\delta > 0$ sedemikian sehingga $\norm{Tx} \ge \delta\norm{x}$ untuk setiap $x \in V$.
\end{defn}

\begin{exam} Jelas bahwa jika suatu pemetaan linear terbatas memiliki sifat terbatas jauh dari nol, maka pemetaan itu injektif. Operator
$T\colon x \mapsto (x_1, \frac12 x_2, \frac13 x_3, \dots)$ yang didefinisikan pada ruang Hilbert $l_2$ menunjukkan bahwa sifat terbatas
jauh dari nol merupakan syarat yang sungguh lebih kuat daripada sifat injektif.
\end{exam}

\begin{prop}\label{prop_bafz_closed_rng} Setiap pemetaan linear terbatas antara ruang Hilbert yang terbatas jauh dari
nol memiliki jangkauan tertutup.
\end{prop}

\begin{prop}\label{prop_nasc_op_invert} Suatu operator pada ruang Hilbert invertibel jika dan hanya jika
operator itu terbatas jauh dari nol dan memiliki jangkauan padat.
\end{prop}

\begin{prop} Pasangan pemetaan $H \mapsto H$ dan $T \mapsto T^*$ yang memetakan setiap ruang Hilbert ke
 \index{adjoin!sebagai funktor pada $\cat{HSp}$}%
 \index{funktor!adjoin}%
dirinya sendiri dan setiap pemetaan linear terbatas antara ruang Hilbert ke adjoinnya merupakan funktor kontravarian
dari kategori $\cat{HSp}$ yang terdiri atas ruang Hilbert dan pemetaan linear terbatas ke kategori itu sendiri.
 \index{HSp@$\cat{HSp}$!funktor adjoin pada}%
\end{prop}

\begin{prop} Misalkan $H$ ruang Hilbert, $T \in \ofml B(H)$, $M = H \oplus \{0\}$, dan $N$ adalah graf dari~$T$.
 \begin{enumerate}[label=\rm{(\alph*)}]
  \item Maka $N$ adalah subruang linear tertutup dari $H \oplus H$.
  \item $T$ injektif jika dan hanya jika $M \cap N = \{(0,0)\}$.
  \item $\ran T$ padat dalam $H$ jika dan hanya jika $M + N$ padat dalam $H \oplus H$.
  \item $T$ surjektif jika dan hanya jika $M + N = H \oplus H$.
 \end{enumerate}
\end{prop}

\begin{exam} Terdapat ruang Hilbert $H$ dengan subruang $M$ dan $N$ sedemikian sehingga $M + N$ tidak tertutup.
\end{exam}

\begin{prop} Misalkan $H$ dan $K$ ruang Hilbert dan $V$ adalah subruang vektor dari~$H$. Setiap pemetaan linear terbatas
$T\colon V \sto K$ dapat diperluas menjadi pemetaan linear terbatas dari $\clo V$, yaitu tutupan $V$, tanpa memperbesar normanya.
\end{prop}

\begin{exer} Misalkan $\fml E$ adalah basis ortonormal bagi ruang Hilbert~$H$. Bahas transformasi linear $T$ pada $H$
sedemikian sehingga $T(e) = 0$ untuk setiap $e \in \fml E$. Berikan contoh tak trivial.
\end{exer}






















\section{Aljabar dengan Involusi}
\begin{defn}\label{00109}  Suatu
 \index{involusi}%
\df{involusi} pada suatu aljabar $A$ adalah suatu pemetaan $x \mapsto x^\ast$ dari $A$ ke $A$ yang memenuhi
  \begin{enumerate}
    \item[(a)] $(x + y)^\ast = x^\ast + y^\ast$,
    \item[(b)] $(\alpha x)^* = \conj\alpha\, x^*$,
    \item[(c)] $x^{\ast\ast} = x$, dan
    \item[(d)] $(xy)^\ast = y^\ast x^\ast$
  \end{enumerate}
untuk semua $x$, $y \in A$, dan $\alpha \in \C$.  Suatu aljabar yang padanya telah didefinisikan suatu involusi adalah
 \index{star@$*\,$-aljabar}%
 \index{<star@$*\,$-aljabar}%
 \index{aljabar!dengan involusi}%
 \index{aljabar!$*\,$-}%
\df{aljabar-$*\,$} (dibaca ``aljabar bintang''). Jika $a$ adalah elemen dari suatu aljabar-$*\,$, maka $a^*$ disebut
 \index{adjoin}%
\df{adjoin} dari~$a$.
\end{defn}

\begin{exam} Dalam aljabar $\C$ dari bilangan kompleks, pemetaan $z \mapsto \conj z$ yang memetakan suatu
bilangan ke konjugat kompleksnya merupakan suatu involusi.
\end{exam}

\begin{exam}\label{0010902} Pemetaan yang memetakan suatu matriks $n \times n$ ke transpos konjugatnya
merupakan suatu involusi pada aljabar beridentitas~$M_n$.
\end{exam}

\begin{exam} Misalkan $X$ suatu ruang Hausdorff kompak.  Pemetaan $f \mapsto \conj f$ yang memetakan suatu
fungsi ke konjugat kompleksnya merupakan suatu involusi dalam aljabar~$\fml C(X)$.
\end{exam}

\begin{exam} Pemetaan $T \mapsto T^*$ yang memetakan suatu operator ruang Hilbert ke adjoinnya merupakan suatu
involusi dalam aljabar $\ofml B(H)$ (lihat proposisi~\ref{prop_adj_invol}).
\end{exam}

\begin{prop} Misalkan $a$ dan $b$ elemen dari suatu aljabar-$*\,$.  Maka $a$ berkomutasi dengan $b$
jika dan hanya jika $a^*$ berkomutasi dengan $b^*$.
\end{prop}

\begin{prop} Dalam suatu aljabar-$*\,$ beridentitas, $\vc 1^* = \vc 1$.
\end{prop}

\begin{prop}\label{00109125} Jika suatu aljabar-$*\,$ $A$ memiliki identitas perkalian kiri $e$,
maka $A$ beridentitas dan $e = \vc 1_A$.
\end{prop}

\begin{prop}\label{00109128fa} Misalkan $a$ elemen dari suatu aljabar-$*\,$ beridentitas.  Maka $a^*$ invertibel jika dan hanya jika $a$ invertibel.
 \index{invers!dari suatu adjoin}%
 \index{invers!adjoin dari suatu}%
 \index{adjoin!invers dari suatu}%
 \index{adjoin!dari suatu invers}%
Dan ketika $a$ invertibel, berlaku
   \[ (a^*)^{-1} = \bigl(a^{-1}\bigr)^*\,. \]
\end{prop}

\begin{defn} Suatu homomorfisme aljabar $\phi$ di antara aljabar-$*\,$ yang mempertahankan
involusi (yaitu, sedemikian sehingga $\phi(a^*) = (\phi(a))^*$) adalah suatu
 \index{<star@homomorfisme-$*\,$}%
 \index{star@homomorfisme-$*\,$}%
 \index{homomorfisme!$*\,$-}%
\df{homomorfisme-$*\,$} (dibaca ``homomorfisme bintang''). Suatu homomorfisme-$*\,$ $\phi\colon A \sto B$ di antara aljabar beridentitas disebut
 \index{beridentitas!homomorfisme-$*\,$}%
 \index{morfisme bijektif!dapat dibalik!dalam aljabar-$*\,$}%
 \index{<star@homomorfisme-$*\,$!beridentitas}%
\df{beridentitas} jika $\phi(\vc 1_A) = \vc 1_B$.  Dalam kategori aljabar-$*\,$ dan homomorfisme-$*\,$, isomorfisme (yang untuk penegasan disebut
 \index{<star@isomorfisme-$*\,$}%
 \index{star@isomorfisme-$*\,$}%
 \index{isomorfisme!$*\,$-}
\df{isomorfisme-$*\,$}) adalah homomorfisme-$*\,$ bijektif.
\end{defn}

\begin{prop}\label{0010952} Misalkan $\phi\colon A \sto B$ suatu homomorfisme-$*\,$ di antara aljabar-$*\,$. Maka
kernel $\phi$ adalah suatu ideal-$*\,$ dalam $A$, dan citra $\phi$ adalah suatu
subaljabar-$*\,$ dari~$B$.
\end{prop}
















\section{Operator Swaadjoin}
\begin{defn} Suatu elemen $a$ dari aljabar-$*\,$ $A$ bersifat
 \index{swaadjoin!elemen}%
\df{swaadjoin} (atau
 \index{Hermitian!elemen}%
\df{Hermitian}) jika $a^* = a$.  Elemen ini bersifat
 \index{normal!elemen}%
\df{normal} jika $a^*a = aa^*$. Jika aljabarnya beridentitas, elemen ini bersifat
 \index{uniter!elemen}%
\df{uniter} jika $a^*a = aa^* = \vc 1$. Himpunan semua elemen swaadjoin dari $A$
dinotasikan dengan
 \index{H@$\fml H(A)$ (elemen swaadjoin dari suatu aljabar-$*\,$)}%
$\fml H(A)$, himpunan semua elemen normal
 \index{N@$\fml N(A)$ (elemen normal dari suatu aljabar-$*\,$)}%
dengan~$\fml N(A)$, dan, ketika aljabarnya beridentitas, himpunan semua elemen uniter
 \index{U@$\fml U(A)$ (elemen uniter dari suatu aljabar-$*\,$)}%
dengan~$\fml U(A)$.
\end{defn}

\begin{prop}\label{001093} Misalkan $a$ elemen dari suatu aljabar-$*\,$ kompleks.  Maka terdapat
elemen swaadjoin unik $u$ dan $v$ sedemikian sehingga $a = u + iv$.
\end{prop}

\begin{proof}[\emph{Petunjuk pembuktian}] Pikirkan kasus khusus penulisan suatu bilangan kompleks
dalam bentuk bagian real dan imajinernya.  \ns
\end{proof}

\begin{defn} Ambil $S$ suatu himpunan bagian dari aljabar-$*\,$~$A$.  Maka
$S^*~=~\{s^*\colon~s~\in~S\}$. Himpunan $S$ disebut
 \index{swaadjoin!himpunan bagian dari suatu aljabar-$*\,$}%
 \index{<unopurstar@$S^*$ (himpunan adjoin dari elemen-elemen~$S$)}%
\df{swaadjoin} jika $S^* = S$.

Suatu subaljabar swaadjoin tak kosong dari $A$ disebut
 \index{<star@subaljabar-$*\,$ (subaljabar swaadjoin)}%
 \index{star@subaljabar-$*\,$ (subaljabar swaadjoin)}%
 \index{subaljabar!$*\,$-}%
 \index{swaadjoin!subaljabar}%
\df{subaljabar-$*\,$}. Istilah sumber yang bersinonim juga diterjemahkan sebagai \df{subaljabar-$*\,$}.
\end{defn}

\begin{cau} Definisi sebelumnya \emph{tidak} menyatakan bahwa elemen-elemen dari suatu
himpunan bagian swaadjoin dari suatu aljabar-$*\,$ dengan sendirinya bersifat swaadjoin.
\end{cau}

 \begin{prop} Suatu operator $T$ pada ruang Hilbert kompleks $H$ bersifat swaadjoin jika dan hanya jika
 bentuk kuadratik terkaitnya $Q_T$ bernilai real.
 \end{prop}

\begin{proof}[\emph{Petunjuk pembuktian}] Gunakan trik yang sama seperti dalam~\ref{C0635105}.  Dalam hipotesis bahwa
$Q_T(z)$ selalu real, gantilah $z$ mula-mula dengan $y + x$, lalu dengan $y + ix$.   \ns
\end{proof}

\begin{defn} Misalkan $T$ suatu operator pada ruang Hilbert tak nol~$H$.  Adapun
 \index{numerik!jangkauan}%
 \index{jangkauan!numerik}%
 \index{W@$W(T)$ (jangkauan numerik)}%
\df{jangkauan numerik} dari $T$, yang dinotasikan dengan $W(T)$, adalah jangkauan yang diperoleh pada sfer satuan dari $H$
dengan membatasi bentuk kuadratik yang terkait dengan~$T$. Yaitu,
   \[ W(T) := \{Q_T(x) \colon x \in H \text{ dan } \norm x = 1\}. \]
Adapun
 \index{numerik!radius}%
 \index{radius!numerik}%
 \index{w@$w(T)$ (radius numerik)}%
\df{radius numerik} dari $T$, yang dinotasikan dengan $w(T)$, adalah $\sup\{\abs \lambda\colon \lambda \in W(T)\}$.
\end{defn}

\begin{prop} Jika $T$ suatu operator swaadjoin pada ruang Hilbert tak nol, maka $\norm T = w(T)$.
\end{prop}

\begin{proof}[\emph{Petunjuk pembuktian}] Bagian yang sulit adalah menunjukkan bahwa $\norm T \le w(T)$.  Untuk itu, tunjukkan terlebih dahulu bahwa
   \begin{equation}
             Q_T(x + y) - Q_T(x - y) =   4 \Re\langle Tx,y \rangle
   \end{equation}
untuk semua vektor $x$ dan $y$ dalam ruang Hilbert, di mana $Q_T$ merupakan bentuk kuadratik yang terkait dengan $T$ dan $\Re z$ adalah
bagian real dari bilangan kompleks~$z$.  Gunakan ini untuk memverifikasi bahwa jika $\norm x = \norm y = 1$, maka
   \begin{equation}\label{num_ran_saop_eqn2}
           \Re\langle Tx,y \rangle \le w(T).
   \end{equation}
Tuliskan bilangan kompleks $\langle Tx,y \rangle$ dalam bentuk polar $re^{i\theta}$.  Jelas bahwa jika kita mengganti
$x$ dengan $e^{-i\theta}x$ dalam pertidaksamaan~\eqref{num_ran_saop_eqn2}, pertidaksamaan yang dihasilkan tetap berlaku kapan pun $\norm x = \norm y = 1$.

Terakhir, perhatikan bahwa jika $\norm x = 1$ dan $y$ adalah sebarang vektor dalam $H$ sedemikian sehingga $Tx = \norm{Tx}\,y$, maka
   \begin{equation}
             \norm{Tx} = \langle Tx,y \rangle\,.
   \end{equation}
Dengan demikian, hasil yang diinginkan segera diperoleh.         \ns
\end{proof}

Dalam proposisi~\ref{C0635105}, kita melihat bahwa suatu operator $T$ pada ruang Hilbert \emph{kompleks} merupakan operator nol jika dan hanya jika
bentuk kuadratik terkaitnya bernilai nol.  Dari proposisi sebelumnya mudah diperoleh akibat bahwa ekuivalensi ini juga berlaku
untuk operator swaadjoin pada ruang Hilbert real.

\begin{cor}\label{real_HSp_zero_Qform} Misalkan $T$ suatu operator swaadjoin pada ruang Hilbert.  Maka $T = \vc 0$ jika dan
hanya jika bentuk kuadratik terkaitnya $Q_T$ bernilai nol.
\end{cor}

\begin{defn} Suatu operator $T$ pada ruang Hilbert $H$ bersifat
 \index{positif!operator}%
 \index{operator!positif}%
\df{positif} jika operator tersebut swaadjoin dan bentuk kuadratik terkaitnya positif ($Q_T(x) \ge 0$ untuk setiap $x \in H$).
\end{defn}

\begin{exer} Misalkan $H$ ruang Hilbert kompleks dan $T$ operator positif pada~$H$.  Definisikan
   \[ \langle x,y \rangle_1 := \langle Tx,y \rangle \qquad \text{ dan } \qquad \norm{x}_1 \equiv \sqrt{\langle x,x \rangle_1} \]
untuk semua $x \in H$.
 \begin{enumerate}
  \item[(a)] Buktikan bahwa $\norm{\hphantom O}_1$ merupakan norma jika dan hanya jika $\ker T = \{0\}$.
  \item[(b)] Misalkan $\ker T = \{0\}$. Tunjukkan bahwa norma $\norm{\hphantom O}$ dan
$\norm{\hphantom O}_1$ menginduksi topologi yang sama pada~$H$ jika dan hanya jika $T$ invertibel dalam
$\ofml B(H)$. \emph{Petunjuk.}  Ketika terdapat dua topologi pada~$H$, sering kali berguna untuk
mempertimbangkan operator identitas pada~$H$.
  \item[(c)] Misalkan $T$ invertibel dalam $\ofml B(H)$.  Untuk $S \in \ofml B(H)$, carilah suatu rumus bagi
adjoin dari $S$ terhadap hasil kali dalam $\langle \hphantom x , \hphantom y \rangle_1$.
 \end{enumerate}
\end{exer}


















\section{Proyeksi}
\begin{conv} Saat ini perhatian kita terutama difokuskan pada
 \index{konvensi!proyeksi di ruang Hilbert bersifat ortogonal}%
operator pada ruang Hilbert.  Dalam konteks ini, istilah \emph{proyeksi} selalu dimaknai
 \index{proyeksi}%
sebagai \emph{proyeksi ortogonal} (lihat definisi~\ref{000164}). Jadi, suatu operator pada ruang Hilbert
disebut proyeksi jika $P^2 = P$ dan $P^* = P$.
\end{conv}

Kita memperumum definisi ``proyeksi (ortogonal)'' dari ruang Hilbert ke aljabar-$*\,$.

\begin{defn} Suatu
 \index{proyeksi!dalam suatu aljabar-$*\,$}%
\df{proyeksi} dalam suatu aljabar-$*\,$ $A$ adalah suatu elemen $p$ dari aljabar tersebut yang bersifat
idempoten ($p^2 = p$) dan swaadjoin ($p^* = p$).  Himpunan semua proyeksi dalam $A$
dinotasikan
 \index{P@$\fml P(A)$!proyeksi dalam suatu aljabar}%
dengan~$\fml P(A)$.  Dalam kasus $\ofml B(H)$, yaitu operator terbatas pada suatu ruang Hilbert,
kita menulis $\fml P(H)$,
 \index{p@$\fml P(H)$ atau $\fml P(\ofml B(H))$!proyeksi ortogonal pada ruang Hilbert}%
atau, jika $H$ telah dipahami, cukup $\fml P$, untuk notasi yang lebih informatif $\fml P(\ofml B(H))$.
\end{defn}

\begin{prop} Setiap operator pada ruang Hilbert yang merupakan isometri pada komplemen ortogonal dari kernelnya memiliki jangkauan tertutup.
\end{prop}

\begin{proof}[\emph{Petunjuk pembuktian}] Perhatikan terlebih dahulu bahwa jika $M = (\ker T)^\perp$, dengan $T$ adalah operator yang dimaksud, maka
$\ran T = T^{\sto}(M)$.    \ns
\end{proof}

\begin{prop} Misalkan $P$ suatu proyeksi pada ruang Hilbert $H$. Maka
 \begin{enumerate}[label=\rm{(\alph*)}]
  \item $Px = x$ jika dan hanya jika $x \in \ran P$\,;
  \item $\ker P = (\ran P)^\perp$; dan
  \item $H = \ker P \oplus \ran P$.
 \end{enumerate}
\end{prop}

Pada bagian (c) dari proposisi sebelumnya, lambang $\oplus$ tentu saja menyatakan
jumlah langsung \emph{ortogonal} (lihat paragraf setelah~\ref{000164}).

\begin{prop} Misalkan $M$ suatu subruang dari ruang Hilbert $H$.  Jika $P \in \ofml B(H)$, jika $Px = x$
untuk setiap $x \in M$, dan jika $Px = \vc 0$ untuk setiap $x \in M^\perp$, maka $P$ adalah
proyeksi dari $H$ pada~$M$.
\end{prop}

Pada bagian selebihnya dari seksi ini, penting untuk selalu mengingat bahwa ruang $\ofml B(H)$ dari operator pada ruang
Hilbert merupakan \emph{contoh} aljabar-$*\,$, dan bahwa proyeksi (ortogonal) pada suatu ruang Hilbert merupakan
\emph{contoh} proyeksi dalam aljabar-$*\,$.  Jadi, di satu sisi, semua hal yang dinyatakan untuk
proyeksi semacam itu dalam aljabar-$*\,$ berlaku bagi proyeksi ruang Hilbert. Karena tidak ada nama yang lebih baik, kita akan mengatakan bahwa
hasil-hasil ini mendokumentasikan struktur
 \index{pandangan abstrak tentang proyeksi}%
\df{abstrak} dari proyeksi.  Di sisi lain, proyeksi ruang Hilbert memiliki struktur tambahan yang
tidak bermakna dalam aljabar-$*\,$ umum: proyeksi tersebut memiliki \emph{domain} dan \emph{jangkauan}, serta \emph{bekerja pada vektor}.
Ketika menyelidiki hal-hal semacam itu, kita akan mengatakan bahwa kita membahas struktur
 \index{pandangan spasial tentang proyeksi}%
\df{spasial} (atau \df{konkret}) dari proyeksi.  Perhatikan bagaimana kedua pandangan ini muncul berselang-seling dalam bagian selebihnya dari seksi ini.

Hasil pertama kita memberikan syarat perlu dan cukup agar jumlah proyeksi merupakan suatu proyeksi.

\begin{prop}\label{prop_sum_projs} Misalkan $p$ dan $q$ proyeksi dalam suatu aljabar-$*\,$. Maka pernyataan berikut ekuivalen:
 \begin{enumerate}[label=\rm{(\alph*)}]
  \item $pq = \vc 0$;
  \item $qp = \vc 0$;
  \item $qp = -pq$;
 \index{proyeksi!jumlah dari}%
  \item $p + q$ merupakan suatu proyeksi.
 \end{enumerate}
\end{prop}

\begin{defn}\label{def_perp_projs} Misalkan $p$ dan $q$ proyeksi dalam suatu aljabar-$*\,$.  Jika salah satu syarat
dalam hasil sebelumnya berlaku, maka kita mengatakan bahwa $p$ dan $q$ bersifat
 \index{proyeksi!ortogonal}%
 \index{proyeksi!keortogonalan dari dua proyeksi}%
 \index{ortogonal!proyeksi!dalam suatu aljabar-$*\,$}%
\df{ortogonal} dan menulis
 \index{<binrel@$p \perp q$ (keortogonalan proyeksi)}%
$p \perp q$.  (Jadi, untuk operator pada suatu ruang Hilbert, kita dapat dengan tepat, meskipun tidak nyaman, menyebut
\emph{proyeksi ortogonal yang ortogonal}!)
\end{defn}

Berikutnya adalah hasil spasial yang bersesuaian untuk jumlah dua proyeksi.

\begin{prop} Misalkan $P$ dan $Q$ proyeksi pada ruang Hilbert~$H$. Maka $P \perp Q$
 \index{ortogonal!proyeksi ruang Hilbert}%
jika dan hanya jika $\ran P \perp \ran Q$. Dalam kasus ini, $P + Q$ merupakan suatu proyeksi
yang kernelnya adalah $\ker P \cap \ker Q$ dan jangkauannya adalah $\ran P + \ran Q$.
\end{prop}

\begin{exam} Pada suatu ruang Hilbert, proyeksi (ortogonal) tidak harus berkomutasi. Sebagai contoh,
misalkan $P$ proyeksi dari ruang Hilbert (real) $\R^2$ pada garis $y = x$ dan $Q$
adalah proyeksi dari $\R^2$ pada sumbu-$x$. Maka $PQ \ne QP$.
\end{exam}

Hasil kali dua proyeksi merupakan suatu proyeksi jika dan hanya jika keduanya berkomutasi.

\begin{prop} Misalkan $p$ dan $q$ proyeksi dalam suatu aljabar-$*\,$. Maka $pq$ merupakan suatu
 \index{proyeksi!hasil kali dari}%
proyeksi jika dan hanya jika $pq = qp$.
\end{prop}

\begin{prop} Misalkan $P$ dan $Q$ proyeksi pada ruang Hilbert~$H$. Jika $PQ = QP$, maka
$PQ$ merupakan suatu proyeksi yang kernelnya adalah $\ker P + \ker Q$ dan jangkauannya adalah $\ran P \cap
\ran Q$.
\end{prop}

\begin{prop} Misalkan $p$ dan $q$ proyeksi dalam suatu aljabar-$*\,$. Maka pernyataan berikut ekuivalen:
 \begin{enumerate}[label=\rm{(\alph*)}]
  \item $pq = p$;
  \item $qp = p$;
 \index{proyeksi!selisih dari}%
  \item $q - p$ merupakan suatu proyeksi.
 \end{enumerate}
\end{prop}

\begin{defn}\label{0022251} Misalkan $p$ dan $q$ proyeksi dalam suatu aljabar-$*\,$.  Jika salah satu
dari syarat dalam hasil sebelumnya berlaku, maka kita
 \index{<binrelle@$p \preceq q$ (urutan proyeksi dalam suatu aljabar-$*\,$)}%
 \index{proyeksi!urutan dari}%
 \index{urutan!parsial!dari proyeksi}%
 \index{parsial!urutan!dari proyeksi}%
menulis $p \preceq q$.
\end{defn}

\begin{prop} Misalkan $P$ dan $Q$ proyeksi pada ruang Hilbert~$H$. Maka pernyataan berikut ekuivalen:
 \begin{enumerate}[label=\rm{(\alph*)}]
  \item $P \preceq Q$;
  \item $\norm{Px} \le \norm{Qx}$ untuk semua $x \in H$; dan
  \item $\ran P \subseteq \ran Q$.
 \end{enumerate}
Dalam kasus ini, $Q - P$ merupakan suatu proyeksi yang kernelnya adalah $\ran P + \ker Q$ dan jangkauannya adalah $\ran Q \ominus \ran P$.
\end{prop}

\noindent\emph{Notasi:} Misalkan $H$, $M$, dan $N$ subruang dari suatu ruang Hilbert. Pernyataan
 \index{<binopdirectdiff@$\ominus$ (selisih langsung)}%
$H = M \oplus N$ dapat ditulis ulang sebagai $M = H \ominus N$ (atau $N = H \ominus M$).

\begin{prop} Operasi $\preceq$ yang didefinisikan dalam~\ref{0022251} untuk proyeksi pada suatu
aljabar-$*\,$ beridentitas $A$ merupakan urutan parsial pada $\fml P(A)$. Jika $p$ adalah suatu proyeksi dalam $A$,
maka $\vc 0 \preceq p \preceq \vc 1$.
\end{prop}

\begin{prop} Misalkan $p$ dan $q$ proyeksi dalam suatu aljabar-$*\,$~$A$. Jika $pq = qp$, maka
infimum dari $p$ dan $q$, yang kita notasikan
 \index{<binrel@$p \curlywedge q$ (infimum proyeksi)}%
 \index{infimum!dari proyeksi dalam suatu aljabar-$*\,$}%
 \index{proyeksi!infimum dari}%
dengan $p \curlywedge q$, ada terhadap urutan parsial $\preceq$ dan $p
\curlywedge q = pq$.  Infimum $p \curlywedge q$ dapat ada bahkan ketika $p$ dan $q$ tidak
berkomutasi.  Syarat perlu dan cukup agar $p \perp q$ berlaku adalah bahwa $p
\curlywedge q = \vc 0$ dan $pq = qp$ keduanya berlaku.
\end{prop}

\begin{prop} Misalkan $p$ dan $q$ proyeksi dalam suatu aljabar-$*\,$~$A$. Jika $p \perp q$, maka
supremum dari $p$ dan $q$, yang kita notasikan
 \index{<binrel@$p \curlyvee q$ (supremum proyeksi)}%
 \index{supremum!dari proyeksi dalam suatu aljabar-$*\,$}%
 \index{proyeksi!supremum dari}%
dengan $p \curlyvee q$, ada terhadap urutan parsial $\preceq$ dan $p
\curlyvee q = p + q$.  Supremum $p \curlyvee q$ dapat ada bahkan ketika $p$ dan $q$ tidak
ortogonal.
\end{prop}

\begin{prop} Misalkan $A$ suatu aljabar-$C^*$, $a \in A$ suatu elemen positif, dan $0 < \epsilon \le \frac12$. Jika
$\norm{a^2 - a} < \epsilon/2$, maka terdapat suatu proyeksi $p$ dalam $A$ sedemikian sehingga
$\norm{p - a} < \epsilon$.
\end{prop}
 
 
 
 
 



 
 

\section{Operator Normal}

\begin{prop}\label{prop_nasc_normal} Suatu operator $T$ pada ruang Hilbert $H$ bersifat normal jika dan hanya jika
$\norm{Tx} = \norm{T^*x}$ untuk semua $x \in H$.
\end{prop}

\begin{proof}[\emph{Petunjuk pembuktian}] Gunakan Korolar~\ref{real_HSp_zero_Qform}.  \ns
\end{proof}

\begin{cor} Jika $T$ merupakan operator normal pada suatu ruang Hilbert, maka $\ker T = \ker T^*$.
\end{cor}

\begin{prop} Jika $T$ merupakan operator normal pada suatu ruang Hilbert, maka $\norm{T^2} = {\norm T}^2$.
\end{prop}

\begin{proof}[\emph{Petunjuk pembuktian}] Substitusikan $Tx$ sebagai pengganti $x$ dalam Proposisi~\ref{prop_nasc_normal}. Gunakan
Proposisi~\ref{C063544}.  \ns
\end{proof}

\begin{prop} Jika $T$ merupakan operator normal pada ruang Hilbert $H$, maka
   \[ H = \clo{\ran T} \oplus \ker T\,.  \]
\end{prop}

\begin{cor} Suatu operator normal pada ruang Hilbert memiliki jangkauan padat jika dan hanya jika operator itu injektif.
\end{cor}

\begin{exam} Operator \emph{geser unilateral} (\ref{C063526}) merupakan contoh operator pada ruang Hilbert yang injektif
tetapi tidak memiliki jangkauan padat. Adjoinnya memiliki jangkauan padat (bahkan, bersifat surjektif), tetapi tidak injektif.
\end{exam}

\begin{cor} Suatu operator normal pada ruang Hilbert bersifat invertibel jika dan hanya jika operator tersebut terbatas jauh dari nol.
\end{cor}
















\section{Operator Berperingkat Hingga}

\begin{defn}
 \index{peringkat}%
\df{peringkat} suatu pemetaan linear adalah dimensi (ruang vektor) dari jangkauannya. Jadi, suatu operator
 \index{hingga!peringkat}%
 \index{operator!berperingkat hingga}%
\df{berperingkat hingga} adalah operator yang memiliki jangkauan berdimensi hingga. Kita menyatakan dengan $\ofml{FR}(V)$
 \index{fr@$\ofml{FR}(V)$ (operator berperingkat hingga pada~$V$)}%
keluarga operator berperingkat hingga pada suatu ruang vektor $V$.
\end{defn}

\begin{exam}\label{0022431fa} Untuk vektor $x$ dan $y$ dalam ruang Hilbert $H$, definisikan
   \[ x \otimes y\colon H \sto H \colon z \mapsto \langle z,y \rangle x\,. \]
Jika $x$ dan $y$ merupakan vektor tak nol dalam $H$, maka $x \otimes y$ merupakan operator berperingkat satu pada~$H$.
\end{exam}

Dua proposisi berikutnya mengidentifikasi adjoin dan komposisi operator-operator berperingkat satu ini.

\begin{prop} Jika $u$ dan $v$ merupakan vektor dalam suatu ruang Hilbert, maka
   \[ (u \otimes v)^* = v \otimes u\,. \]
\end{prop}

\begin{prop} Jika $u$, $v$, $x$, dan $y$ merupakan vektor dalam suatu ruang Hilbert, maka
    \[ (u \otimes v)(x \otimes y) = \langle x,v \rangle (u \otimes y)\,. \]
\end{prop}

\begin{prop} Jika $x$ merupakan vektor dalam ruang Hilbert $H$, maka $x \otimes x$ merupakan proyeksi berperingkat satu jika dan
hanya jika $x$ merupakan vektor satuan.
\end{prop}

\begin{prop} Andaikan bahwa $\{e^1, e^2, e^3, \dots, e^n\}$ merupakan himpunan ortonormal dalam ruang Hilbert~$H$. Misalkan
$M = \spn\{e^1, \dots, e^n\}$. Maka $P = \sum_{k=1}^n e^k \otimes e^k$ merupakan proyeksi ortogonal dari $H$ ke~$M$.
\end{prop}

\begin{prop}\label{prop_op_act_otimes1} Jika $u$ dan $v$ merupakan vektor dalam suatu ruang Hilbert dan $T$ merupakan operator pada ruang Hilbert yang sama, yaitu~$H$, maka
   \[ T\,(u \otimes v) = (Tu) \otimes v\,. \]
\end{prop}

\begin{prop}\label{prop_op_act_otimes2} Jika $u$ dan $v$ merupakan vektor dalam suatu ruang Hilbert dan $T$ merupakan operator pada ruang Hilbert yang sama, yaitu~$H$, maka
   \[ (u \otimes v)\,T = u \otimes T^*v\,. \]
\end{prop}

Kita akan segera melihat bahwa keluarga semua operator berperingkat hingga pada suatu ruang Hilbert tak nol merupakan ideal-$*\,$ minimal dalam
aljabar-$*\,$ $\ofml B(H)$. Sebagai persiapan, kita meninjau beberapa fakta dasar tentang ideal dalam aljabar.

\begin{defn}\label{000551} Suatu
 \index{ideal!kiri}%
 \index{kiri!ideal}%
\df{ideal kiri} dalam suatu aljabar $A$ adalah subruang vektor $J$ dari $A$ sedemikian sehingga $AJ \subseteq J$. (Untuk
 \index{ideal!kanan}%
 \index{kanan!ideal}%
\df{ideal kanan}, tentu saja, kita mensyaratkan $JA \subseteq J$.) Kita mengatakan bahwa $J$ merupakan suatu
 \index{ideal!dalam suatu aljabar}%
\df{ideal} jika ideal tersebut merupakan ideal dua sisi, yaitu ideal kiri sekaligus ideal kanan. Suatu ideal disebut
 \index{sejati!ideal}%
 \index{ideal!sejati}%
\df{sejati} jika ideal itu merupakan subhimpunan sejati dari~$A$.

Ideal $\{0\}$ dan $A$ sering disebut sebagai
 \index{trivial!ideal}%
 \index{ideal!trivial}%
\df{ideal trivial} dari~$A$. Aljabar $A$ disebut
 \index{sederhana!aljabar}%
 \index{aljabar!sederhana}%
\df{sederhana} jika tidak memiliki ideal tak trivial.

Suatu
 \index{ideal maksimal}%
 \index{ideal!maksimal}%
ideal \df{maksimal} adalah ideal sejati yang tidak termuat secara sejati dalam ideal sejati lain.
Kita menyatakan keluarga semua ideal maksimal dalam suatu aljabar $A$
 \index{maximal@$\mx A$!himpunan ideal maksimal dalam suatu aljabar}%
dengan~$\mx A$. Suatu
 \index{ideal minimal}%
 \index{ideal!minimal}%
ideal \df{minimal} adalah ideal tak nol yang tidak memuat secara sejati ideal tak nol lain.
\end{defn}

 \index{konvensi!ideal bersifat dua sisi}%
\begin{conv} Setiap kali kita merujuk pada suatu \emph{ideal} dalam aljabar, ideal tersebut dipahami sebagai
ideal dua sisi (kecuali dinyatakan sebaliknya).
\end{conv}

\begin{prop} Tidak satu pun elemen invertibel dalam aljabar beridentitas dapat termasuk dalam suatu ideal sejati.
\end{prop}

 \begin{prop}\label{000552} Setiap ideal sejati dalam aljabar beridentitas $A$ termuat dalam suatu
 ideal maksimal. Jadi, khususnya, $\mx A$ tak kosong apabila $A$ merupakan aljabar beridentitas.
\end{prop}

\begin{proof}[\emph{Petunjuk pembuktian}] Lema Zorn. \ns  \end{proof}

\begin{prop}\label{000553} Misalkan $a$ suatu elemen dari aljabar komutatif~$A$. Jika $A$
beridentitas dan $a$ tidak invertibel, maka $aA$ merupakan ideal sejati dalam~$A$.
\end{prop}

\begin{defn}\label{000554} Ideal $aA$ dalam proposisi sebelumnya adalah
 \index{ideal utama}%
 \index{ideal!utama}%
\df{ideal utama} yang dibangkitkan oleh~$a$.
\end{defn}

\begin{defn}\label{000555} Misalkan $J$ suatu ideal dalam aljabar~$A$. Definisikan suatu
relasi ekuivalensi $\sim$ pada $A$ dengan
  \[ a \sim b \qquad \text{jika dan hanya jika } \qquad b-a \in J. \]
Untuk setiap $a \in A$, misalkan $[a]$ adalah kelas ekuivalensi yang memuat~$a$. Misalkan $A/J$ adalah
himpunan semua kelas ekuivalensi elemen-elemen~$A$. Untuk $[a]$ dan $[b]$ dalam $A/J$, definisikan
  \[ [a] + [b] := [a+b] \qquad \text{ dan } \qquad [a][b] := [ab] \]
dan untuk $\alpha \in \C$ dan $[a] \in A/J$, definisikan
  \[\alpha[a] := [\alpha a] \,. \]
Terhadap operasi-operasi ini, $A/J$ menjadi suatu aljabar. Aljabar tersebut adalah
 \index{hasil bagi!aljabar}%
 \index{aljabar!hasil bagi}%
\df{aljabar hasil bagi} dari $A$ oleh~$J$. Notasi $A/J$ biasanya dibaca ``$A$ modulo~$J$''.
Homomorfisme aljabar surjektif
  \[ \pi\colon A \sto A/J\colon a \mapsto [a]\]
disebut
 \index{hasil bagi!pemetaan!untuk aljabar}%
\df{pemetaan hasil bagi}.
\end{defn}

\begin{prop} Definisi sebelumnya bermakna dan pernyataan-pernyataan di dalamnya benar.
Lebih lanjut, jika ideal $J$ bersifat sejati, maka aljabar hasil bagi $A/J$ beridentitas jika $A$ beridentitas.
\end{prop}

\begin{proof}[\emph{Petunjuk pembuktian}] Anda perlu menunjukkan bahwa:
 \begin{enumerate}
  \item[(a)] $\sim$ merupakan relasi ekuivalensi.
  \item[(b)] Penjumlahan dan perkalian kelas-kelas ekuivalensi terdefinisi dengan baik.
  \item[(c)] Perkalian suatu kelas ekuivalensi dengan skalar terdefinisi dengan baik.
  \item[(d)] $A/J$ merupakan suatu aljabar.
  \item[(e)] ``Pemetaan hasil bagi'' $\pi$ memang merupakan homomorfisme aljabar surjektif.
 \end{enumerate}
(Pada tahap mana diperlukan bahwa kita mengambil hasil bagi oleh suatu ideal, bukan sekadar subaljabar?) \ns
\end{proof}

\begin{defn} Dalam suatu aljabar $A$ dengan involusi,
 \index{<star@ideal-$*\,$}%
 \index{star@ideal-$*\,$}%
 \index{ideal!$*\,$-}%
suatu \df{ideal-$*\,$} adalah ideal swaadjoin dalam~$A$.
\end{defn}

\begin{prop}\label{0010953} Andaikan $J$ suatu ideal-$*\,$ dalam aljabar-$*\,$ $A$. Dalam aljabar
hasil bagi $A/J$, sebagaimana dikembangkan dalam~\ref{000555}, definisikan $[a]^* := [a^*]$ untuk setiap $[a] \in A/J$.
 \index{hasil bagi!aljabar-$*\,$}%
 \index{<star@aljabar-$*\,$!hasil bagi}%
Ini merupakan operasi uner yang terdefinisi dengan baik pada aljabar hasil bagi $A/J$ yang terbentuk; aljabar tersebut menjadi aljabar-$*\,$, dan
pemetaan hasil bagi $a \mapsto [a]$ merupakan homomorfisme-$*\,$. Aljabar-$*\,$ $A/J$ adalah, tentu
saja,
 \index{hasil bagi!aljabar-$*\,$}%
 \index{<star@aljabar-$*\,$!hasil bagi}%
 \index{star@aljabar-$*\,$!hasil bagi}%
\df{hasil bagi} dari $A$ oleh~$J$.
\end{prop}

\begin{prop} Misalkan $H$ suatu ruang Hilbert. Maka keluarga $\ofml{FR}(H)$ dari semua operator berperingkat hingga pada $H$
 \index{star@ideal-$*\,$!$\ofml{FR}(H)$ sebagai suatu}%
 \index{FR@$\ofml{FR}(H)$!sebagai ideal-$*\,$}%
merupakan ideal-$*\,$ dalam aljabar-$*\,$~$\ofml B(H)$.
\end{prop}

\begin{prop}\label{prop_frop_sum_rnk1} Misalkan $T$ suatu operator berperingkat $n$ pada ruang Hilbert~$H$. Maka terdapat
vektor $u^1$, \dots, $u^n$, $v^1$, \dots, $v^n$ dalam $H$ sedemikian sehingga
   \[ T = \sum_{k=1}^n u^k \otimes v^k\,. \]
\end{prop}

\begin{proof}[\emph{Petunjuk pembuktian}] Misalkan $\{e^1, \dots, e^n\}$ suatu basis ortonormal bagi jangkauan~$T$. Untuk
sembarang vektor $x \in H$, tuliskan \emph{ekspansi Fourier} dari $Tx$ (lihat Proposisi~\ref{C063144}(d)\,).  \ns
\end{proof}

\begin{cor} Setiap operator berperingkat satu pada suatu ruang Hilbert berbentuk $u \otimes v$ untuk suatu pasangan vektor tak nol
$u$ dan $v$ dalam~$H$.
\end{cor}

\begin{cor} Setiap operator berperingkat hingga pada suatu ruang Hilbert merupakan jumlah dari (sejumlah berhingga) operator berperingkat satu.
\end{cor}

\begin{prop} Jika $H \ne \{0\}$ merupakan ruang Hilbert, maka keluarga $\ofml{FR}(H)$ dari operator berperingkat hingga pada $H$
merupakan ideal minimal dalam~$\ofml B(H)$.
\end{prop}

\begin{proof}[\emph{Petunjuk pembuktian}] Misalkan $\ofml J$ suatu ideal tak nol dalam $\ofml B(H)$. Kita perlu menunjukkan bahwa $\ofml J$
memuat~$\ofml{FR}(H)$. Berdasarkan Proposisi~\ref{prop_frop_sum_rnk1}, cukup dibuktikan bahwa $\ofml J$ memuat setiap
operator berbentuk $u \otimes v$. Untuk itu, misalkan $T$ sebarang anggota tak nol dari $\ofml J$, misalkan $y$ suatu vektor satuan dalam
jangkauan $T$, dan misalkan $x$ sebarang vektor dalam $H$ sedemikian sehingga $y = Tx$. Tinjau operator $(u \otimes y)T(x \otimes v)$. \ns
\end{proof}

Meskipun himpunan operator berperingkat hingga merupakan ideal dalam aljabar Banach yang terdiri atas operator pada suatu ruang Hilbert berdimensi tak hingga, himpunan itu bukan ideal
tertutup dan, sebagai akibatnya, ternyata bukan himpunan yang tepat untuk membentuk hasil bagi dengan harapan memperoleh suatu objek hasil bagi dalam
kategori aljabar Banach. Tutupan ideal operator berperingkat hingga pada suatu ruang Hilbert adalah ideal \emph{operator
kompak}. Dalam arti tertentu, langkah berikutnya yang wajar ialah mengkaji sifat-sifat kelas operator ini; di dalamnya kita menjumpai
beberapa operator terpenting dalam penerapan, misalnya operator integral. Namun, pemaparan yang lebih lancar dan lebih gamblang
dimungkinkan jika kita memiliki sedikit lebih banyak perangkat teknis. Oleh karena itu, dalam bab berikutnya kita akan berfokus pada
pengembangan beberapa alat baku yang digunakan dalam analisis fungsional.
























\endinput

 \chapter{RUANG BANACH}

\section{Transformasi Alami}
 Ingat bahwa jika $V$ suatu ruang linear bernorma, maka
 \index{dual!ruang}%
 \index{ruang!dual}%
 \index{<unopurstar@$V^*$ (ruang dual dari suatu ruang linear bernorma)}%
\emph{ruang dual} $V^*$ adalah himpunan semua fungsional linear kontinu pada~$V$.  Pada $V^*$, penjumlahan dan
perkalian skalar didefinisikan secara titik demi titik. Normanya adalah norma operator biasa: jika $f \in V^*$, maka
 \[ \norm{f} := \inf\{M > 0\colon \abs{f(x)} \le M\norm{x}\quad \text{untuk setiap $x \in V$}\} = \sup\{\abs{f(x)}\colon \norm{x} \le 1\}. \]
Di bawah metrik yang diinduksi oleh norma ini, ruang dual $V^*$ lengkap dan karena itu merupakan ruang Banach.

\begin{defn}\label{defn_nls_adjoint} Jika $T\colon V \sto W$ suatu pemetaan linear kontinu antara ruang-ruang linear bernorma,
 \index{<unopurstar@$T^*$ (adjoin ruang linear bernorma)}%
$T^*\colon W^* \sto V^*$ didefinisikan seperti pada ruang vektor: $T^*(g) = g\,T$ untuk setiap
$g \in W^*$.  Pemetaan $T^*$ disebut
 \index{adjoin!pemetaan linear!antara ruang linear bernorma}%
 \index{linear!pemetaan!adjoin dari}%
\df{adjoin} dari~$T$.
\end{defn}

\begin{exam}\label{exam_dual_functor} Adjoin $T^*$ dari suatu pemetaan linear terbatas antara ruang-ruang linear bernorma $V$ dan
$W$ sendiri merupakan pemetaan linear terbatas antara $W^*$ dan~$V^*$. Dalam kategori ruang linear bernorma dan pemetaan linear
 \index{funktor!dualitas}%
 \index{funktor dualitas}%
kontinu, pasangan pemetaan $V \mapsto V^*$ (pemetaan objek) dan $T \mapsto T^*$ (pemetaan morfisme) merupakan suatu funktor
kontravarian dari $\cat{NLS_\infty}$ ke dirinya sendiri.
\end{exam}

Perhatikan bahwa Definisi~\ref{defn_nls_adjoint} berbeda dari definisi adjoin suatu pemetaan
antara ruang-ruang Hilbert (lihat~\ref{def000926}).  Karena setiap ruang Hilbert merupakan ruang linear bernorma, hal ini dapat saja
menimbulkan kebingungan.  Ada baiknya kita meyakinkan diri tentang fakta sederhana berikut: karena anti-isomorfisme (dan akibatnya
isomorfisme) antara ruang Hilbert dan dualnya yang diberikan oleh \emph{teorema Riesz--Fr\'echet} (lihat~\ref{0022057}
dan~\ref{cor_Hsp_iso_dual}), adjoin ruang linear bernorma dari suatu pemetaan linear terbatas \emph{sangat} mirip (meskipun
tentu tidak identik) dengan adjoin ruang Hilbertnya.

\begin{exam}\label{C057717}
Dual $V^{**}$ dari dual $V^*$ suatu ruang linear bernorma $V$ adalah \df{dual kedua} dari~$V$
 \index{dual!kedua}%
dan $T^{**} := \bigl(T^*\bigr)^*$ merupakan pemetaan linear terbatas dari $V^{**}$ ke~$W^{**}$.  Pada kategori
$\cat{BAN_{\boldsymbol\infty}}$ yang terdiri atas ruang Banach dan transformasi linear terbatas,
 \index{kedua!dual!funktor untuk ruang linear bernorma}%
 \index{funktor!dual kedua!untuk ruang linear bernorma}%
\df{funktor dual kedua} mencakup pemetaan objek $B \mapsto B^{**}$ dan pemetaan morfisme $T \mapsto T^{**}$.
Funktor ini kovarian.
\end{exam}

\begin{defn}\label{defn_nat_transf} Misalkan $\cat A$ dan $\cat B$ kategori, serta $F$,$G \colon \cat A \sto \cat B$
funktor kovarian.  Suatu
 \index{alami!transformasi}%
 \index{transformasi!alami}%
\df{transformasi alami} dari $F$ ke $G$ adalah pemetaan $\tau$ yang memasangkan kepada setiap objek
$A$ dalam $\cat A$ suatu morfisme $\tau_A \colon F(A) \sto G(A)$ dalam $\cat B$ sedemikian sehingga
untuk setiap morfisme $\alpha \colon A \sto A'$ dalam $\cat A$, diagram berikut komutatif.
   \[ \CD
      F(A) @>  F(\alpha)>>  F(A')  \\
         @V\tau_AVV                @VV\tau_{A'}V  \\
      G(A) @>> G(\alpha)>  G(A')
   \endCD \]
Kita menyatakan transformasi semacam itu dengan $\tau \colon F \sto G$. (Definisi transformasi alami
antara dua funktor kontravarian seharusnya jelas: cukup balikkan
panah-panah horizontal dalam diagram sebelumnya.)

Suatu transformasi alami $\tau \colon  F \sto G$ disebut
 \index{alami!ekuivalensi}%
 \index{ekuivalensi!alami}%
\df{ekuivalensi alami} jika setiap morfisme $\tau_A$ merupakan isomorfisme dalam~$\cat B$.
\end{defn}

\begin{defn}\label{defn_nat_embed} Jika $V$ suatu ruang linear bernorma, maka pemetaan $\sbsb{\tau}V \colon V \sto V^{**} \colon x
\mapsto x^{**}$, dengan $x^{**}(f) = f(x)$ untuk setiap $f \in V^*$, disebut
 \index{tauv@$\sbsb{\tau}V$ (pembenaman alami ruang linear bernorma)}%
 \index{alami!pembenaman!untuk ruang linear bernorma}%
 \index{pembenaman!alami}%
\df{pembenaman alami} dari $V$ ke~$V^{**}$. (Tidak sepenuhnya jelas bahwa pemetaan ini
injektif. Kita akan menunjukkannya dalam Proposisi~\ref{C067137}.)
\end{defn}

\begin{exam}\label{exam_2dual_nat} Dalam kategori $\cat{BAN_{\boldsymbol\infty}}$ yang terdiri atas ruang Banach
dan transformasi linear terbatas, misalkan $I$ adalah funktor identitas dan $(\,\cdot\,)^{**}$ adalah
funktor dual kedua.  Maka
 \index{alami!pembenaman!sebagai funktor}%
 \index{funktor!pembenaman alami sebagai}%
 \index{taub@$\tau_B$ (pembenaman alami $B$ ke $B^{**}$)}%
\emph{pembenaman alami} $\tau$ merupakan transformasi alami dari $I$ ke~$(\,\cdot\,)^{**}$.
\end{exam}

\begin{exam} Jelaskan secara tepat apa yang dimaksud ketika orang mengatakan bahwa isomorfisme
    \[ \fml C(X \uplus Y) \cong \fml C(X) \times \fml C(Y) \]
merupakan ekuivalensi alami.  Kategori-kategori yang terlibat adalah $\cat{CpH}$ (ruang Hausdorff kompak
dan pemetaan kontinu) serta $\cat{BAN_\infty}$ (ruang Banach dan transformasi linear
terbatas).
\end{exam}

\begin{exam}\label{C057724} Pada kategori $\cat{CpH}$ yang terdiri atas ruang Hausdorff kompak dan
pemetaan kontinu, misalkan $I$ funktor identitas dan $\fml C^*$ funktor yang
memetakan ruang Hausdorff kompak $X$ ke ruang Banach $\fml (C(X))^*$, serta memetakan suatu
fungsi kontinu $\phi\colon X \sto Y$ antara ruang-ruang Hausdorff kompak ke
transformasi linear kontraktif $\fml C^*\phi = (\fml C\phi)^*$. Definisikan
 \index{evaluasi!pemetaan}%
\emph{pemetaan evaluasi} $E$ dengan
  \[ E_X\colon X \sto \fml C^*(X)\colon x \mapsto E_X(x) \]
dengan $(E_X(x))(f) = f(x)$ untuk setiap $f \in \fml C(X)$.  (Notasi ini dapat merepotkan.  Dalam praktik, biasanya orang menulis
$E_x$ saja untuk $E_X(x)$.  Jadi kita memperoleh $E_x(f) = f(x)$, yang tentu tampak lebih baik daripada $(E_X(x))(f) = f(x)$.  Dalam kebanyakan
konteks, penyederhanaan notasi ini tampaknya tidak mungkin menimbulkan kebingungan.)
 \begin{enumerate}
  \item[(a)]  Perhatikan bahwa definisi ini memang bermakna. Jika $x \in X$, maka $E_X(x)$ benar-benar
merupakan anggota~$\fml C^*(X)$.
  \item[(b)] Tunjukkan bahwa pemetaan $E$ membuat diagram berikut komutatif:
     \[ \xymatrix@+30pt{{X}\ar[r]^*+{\phi} \ar[d]_*+{E_X}
            & {Y} \ar[d]^*+{E_Y} \\
          {\fml C^*(X)} \ar[r]_*+{\fml C^*(\phi)} & {\fml C^*(Y)}   }
     \]
 \end{enumerate}
\end{exam}

Contoh sebelumnya mungkin terasa agak mengecewakan.  Tampaknya kita sedang
berusaha membentuk suatu transformasi alami antara funktor identitas (pada
$\cat{CpH}$) dan funktor $\fml C^*$.  Masalahnya, tentu saja, ialah kedua
funktor tersebut memiliki kategori sasaran yang berbeda.  Pertanyaan ini akan kita bahas kelak.
Pemetaan evaluasi $E_X\colon X \sto \fml C^*(X)$ tentu tidak surjektif.
Perhatikan, misalnya, bahwa setiap fungsional berbentuk $E_X(x)$
mempertahankan perkalian (selain penjumlahan dan perkalian skalar). Salah satu akibatnya
ialah bahwa semua fungsional semacam itu terletak pada sfer satuan $\fml C^*(X)$.
Dengan suatu topologi baru (yang disebut \emph{topologi lemah-bintang}---lihat~\ref{C067434}) pada
$\fml C^*(X)$, jangkauan $E_X$ ternyata merupakan ruang Hausdorff kompak yang secara alami
ekuivalen dengan $X$ sendiri.  Yang lebih luar biasa lagi, pada akhirnya kita akan menunjukkan bahwa jangkauan $E_X$
dapat dipandang sebagai ruang ideal maksimal dari aljabar~$\fml C(X)$. (Lihat~\futurexref{11.2.20}{000731}.)

\begin{prop}\label{C067137} Jika $V$ suatu ruang linear bernorma, maka pembenaman alami
$\sbsb{\tau}V \colon V \sto V^{**}\colon x \mapsto x^{**}$ (yang didefinisikan dalam~\ref{defn_nat_embed})
merupakan pemetaan linear isometrik.
\end{prop}

\begin{defn} Suatu ruang Banach $B$ disebut
 \index{refleksif!ruang Banach}%
\df{refleksif} jika pembenaman alami $\sbsb{\tau}B\colon B \sto B^{**}$ (yang didefinisikan dalam proposisi sebelumnya) surjektif.
\end{defn}

Contoh~\ref{exam_2dual_nat} menunjukkan bahwa dalam kategori ruang Banach dan pemetaan linear terbatas,
pemetaan $\tau\colon B \mapsto \sbsb{\tau}B$ merupakan transformasi alami dari
funktor identitas ke funktor dual kedua.  Dengan menggunakan \emph{teorema Hahn--Banach}, kita dapat
mengatakan lebih banyak.

\begin{prop} Dalam kategori ruang Banach refleksif dan pemetaan linear terbatas, pemetaan
$\tau\colon B \mapsto \sbsb{\tau}B$ merupakan ekuivalensi alami antara funktor identitas dan
funktor dual kedua.
\end{prop}

\begin{exam} Setiap ruang Hilbert bersifat refleksif.
\end{exam}

\begin{exer} Misalkan $V$ dan $W$ ruang-ruang linear bernorma, $U$ subhimpunan konveks takkosong dari $V$,
dan $f \colon U \sto W$.  Tanpa menggunakan bentuk apa pun dari \emph{teorema nilai rata-rata}, tunjukkan
bahwa jika $df = \vc 0$ pada $U$, maka $f$ konstan pada~$U$.  (Untuk definisi $df$, yaitu
\emph{diferensial} dari~$f$, lihat Bagian 13.1--2 dalam catatan saya~\cite{Erdman:2007}.)
\end{exer}

\begin{exam}\label{C067181} Misalkan $l_\infty$ ruang Banach dari semua barisan bilangan
real yang terbatas, dan $c$ subruang $l_\infty$ yang terdiri atas semua barisan $x = (x_n)$
sedemikian sehingga $\lim_{n \sto \infty} x_n$ ada. Definisikan
  \[ f\colon c \sto \R\colon x \mapsto \lim_{n\sto \infty} x_n. \]
Maka $f$ merupakan fungsional linear kontinu pada $c$ dan $\norm f = 1$.  Berdasarkan
\emph{teorema Hahn--Banach}, $f$ memiliki perluasan ke $\wh f \in {l_\infty}^*$ dengan
$\norm{\wh f} = 1$. Untuk setiap
subhimpunan $A$ dari $\N$, definisikan barisan $\bigl(x^A_n\bigr)$ dengan $x^A_n = \left\{%
\begin{array}{ll}
    1, & \hbox{jika $n \in A$;} \\
    0, & \hbox{jika $n \notin A$.} \\
\end{array}%
\right.$ Selanjutnya, definisikan $ \mu\colon \sfml P(\N) \sto \R \colon A \mapsto \wh
f\bigl(x^A_n\bigr)$. Maka $\mu$ merupakan fungsi himpunan aditif hingga, tetapi tidak aditif
terhitung.
\end{exam}













\section{Teorema Alaoglu}
\begin{notn} Misalkan $V$ suatu ruang linear bernorma, $M \subseteq V$, dan $F \subseteq V^\ast$.
Kita definisikan
 \begin{align*}
         M^\perp &:= \{f \in V^*\colon f(x)=0 \text{ untuk setiap $x \in M$}\} \\
         F_\perp &:= \{x \in V\colon f(x)=0 \text{ untuk setiap $f \in F$}\}
  \end{align*}
$M^\perp$ dapat dibaca ``M perp'' atau ``M perp atas'', sedangkan $F_\perp$ dibaca ``F perp'' atau ``F perp bawah''.
Himpunan $M^\perp$ adalah
 \index{anihilator}%
 \index{<unopur@$M^\perp$ (anihilator suatu himpunan)}%
\df{anihilator} dari~$M$, dan $F_\perp$ adalah
 \index{praanihilator}%
 \index{<unopvr@$F_\perp$ (praanihilator suatu himpunan)}%
\df{praanihilator} dari~$F$.
\end{notn}

\begin{prop}[Anihilator]  Misalkan $B$ suatu ruang Banach, $M \subseteq B$, dan $F \subseteq B^*$. Maka
 \begin{enumerate}[label=\rm{(\alph*)}]
  \item Jika $M \subseteq N \subseteq B$, maka $N^\perp \subseteq M^\perp$.
  \item Jika $F \subseteq G \subseteq B^*$, maka $G_\perp \subseteq F_\perp$.
  \item $M^\perp$ merupakan subruang linear tertutup dari $B^*$.
  \item $F_\perp$ merupakan subruang linear tertutup dari $B$.
  \item $M \subseteq M^\perp{}_\perp$.
  \item $F \subseteq F_\perp{}^\perp$.
  \item $M^\perp{}_\perp = \bigvee M$.
  \item $M^\perp = B^*$ jika dan hanya jika $M \subseteq \{0\}$.
  \item $F_\perp = B$ jika dan hanya jika $F \subseteq \{0\}$.
  \item Jika $M$ suatu subruang linear dari $B$, maka $M^\perp = \{0\}$ jika dan hanya jika $M$
padat dalam $B$.
  \item\label{Fpp_reflexive_case} Jika $B$ refleksif, maka $F_\perp{}^\perp = \bigvee F$.
 \end{enumerate}
\end{prop}

\begin{proof}[\emph{Petunjuk pembuktian}] Untuk (k), tunjukkan bahwa $\tau(F_\perp) = F^\perp$ (dengan $\tau$
pembenaman alami) dan karena itu $F_\perp{}^\perp = F^\perp{}_\perp$.  \ns
\end{proof}

\begin{exam} Misalkan $B$ suatu ruang Banach yang \emph{tidak} refleksif.  Misalkan $F = \ran \tau$ (dengan $\tau$
pembenaman alami).  Maka $F_\perp{}^\perp = B^{**} \ne \bigvee F$, sehingga jelas bahwa~\ref{Fpp_reflexive_case}
dalam proposisi sebelumnya tidak harus berlaku dalam kasus nonrefleksif.
\end{exam}

Hasil berikut adalah analog ruang Banach dari \emph{teorema dasar aljabar linear}~\ref{00016025} dan
Proposisi~\ref{prop_ker_adj_perp_ran_op}.

\begin{prop}\label{prop_FTLA_Bsp} Misalkan $T \in \ofml B(B,C)$, dengan $B$ dan $C$ ruang-ruang Banach. Maka
 \begin{enumerate}[label=\rm{(\alph*)}]
  \item $\ker T^\ast = (\ran T)^\perp$.
  \item $\ker T = (\ran T^*)_\perp$.
\vskip 3 pt
  \item $\clo{\ran T} = (\ker T^*)_\perp$.
\vskip 3 pt
  \item $\clo{\ran T^*} \subseteq (\ker T)^\perp$.
 \end{enumerate}
\end{prop}

\begin{exam} Suatu contoh operator ruang Banach $T$ yang membuat $\clo{\ran T^*}$ dan $(\ker T)^\perp$ berbeda
diberikan dalam Latihan 7, halaman 243 dari~\cite{BrownP:1970}.
\end{exam}

Contoh berikut tentu merupakan akibat langsung yang sepele dari Proposisi~\ref{00015032}, tetapi cobalah memberikan
pembuktian langsung yang sederhana.

\begin{exam}\label{X_l2ball_not_cpt} Bola satuan tertutup dalam $l_2$ tidak kompak.
\end{exam}

Topologi norma biasa pada ruang linear bernorma ternyata terlalu besar untuk beberapa penerapan.  Seperti telah kita lihat,
topologi ini memiliki begitu banyak himpunan terbuka sehingga bola satuan tertutup tidak kompak (kecuali dalam kasus berdimensi hingga). Di
sisi lain, terlalu sedikit himpunan terbuka (sebagai contoh ekstrem, topologi indiskret) akan merusak teori dualitas karena tidak
menyediakan cukup banyak fungsional linear kontinu.  \emph{Topologi lemah-bintang} (yang didefinisikan
di bawah) pada dual suatu ruang linear bernorma ternyata tepat ukurannya.  Topologi ini memiliki cukup banyak himpunan terbuka untuk menjamin
teori dualitas yang kuat dan cukup sedikit himpunan terbuka untuk membuat bola satuan tertutup kompak
(lihat \emph{teorema Alaoglu}~\ref{C067441}).

\begin{defn}\label{C067434} Misalkan $V$ suatu ruang linear bernorma. 
 \index{lemah!topologi!pada ruang bernorma}%
 \index{topologi!lemah!pada ruang linear bernorma}%
\df{topologi lemah} pada $V$ adalah topologi lemah yang diinduksi oleh unsur-unsur~$V^*$. (Jadi,
topologi lemah pada ruang dual $V^*$ adalah topologi lemah yang diinduksi oleh
unsur-unsur~$V^{**}$.) Ketika suatu jaring $(x_\lambda)$ konvergen secara lemah ke vektor $a$ dalam $V$,
kita menulis
 \index{<arrowright@$x_\lambda \to^w a$ (konvergensi lemah)}%
$x_\lambda \to^w a$. 
 \index{w*@$w^*$-topologi (topologi lemah-bintang)}%
 \index{topologi!lemah-bintang}%
\df{topologi-$w^*$} (dibaca \emph{topologi lemah-bintang}) pada ruang dual $V^*$ adalah
topologi lemah yang diinduksi oleh unsur-unsur $\ran \tau$, dengan $\tau$ pembenaman alami
dari $V$ ke~$V^{**}$.  Ketika suatu jaring $(f_\lambda)$ konvergen secara lemah-bintang ke vektor $g$
dalam $V^*$, kita menulis
 \index{<arrowright@$f_\lambda \to^{w^*} g$ (konvergensi lemah-bintang)}%
$f_\lambda \to^{w^*} g$. Perhatikan bahwa ketika suatu ruang Banach $B$ refleksif (khususnya, untuk ruang
Hilbert), topologi lemah dan lemah-bintang pada $B^*$ identik.
\end{defn}

\begin{prop} Misalkan $f \in V^*$, dengan $V$ suatu ruang Banach.  Untuk setiap $\epsilon > 0$ dan setiap
subhimpunan hingga $F \subseteq V$, definisikan
  \[ U(f;F;\epsilon)
           = \{g \in V^*\colon \abs{f(x) - g(x)} < \epsilon \text{ untuk setiap $x \in F$}\}. \]
Keluarga semua himpunan semacam itu merupakan basis bagi topologi-$w^*$ pada~$V^*$.
\end{prop}

Hasil berikut, \emph{teorema Alaoglu}, menyatakan kekompakan-$w^*$ dari bola satuan tertutup dalam
dual suatu ruang linear bernorma $V$ dan merupakan salah satu teorema yang paling berguna dalam analisis fungsional.  Karena itu,
pada awalnya mungkin agak mengejutkan bahwa pembuktian lengkap yang dijelaskan secara cermat
begitu sulit ditemukan.  Pembuktian standar pada dasarnya sangat sederhana; intinya hanyalah mengenali
bahwa anggota-anggota suatu keluarga $\fml F$ dari fungsi bernilai skalar dapat dipandang
  \begin{enumerate}
     \item[(i)] sebagai fungsi yang didefinisikan pada hasil kali takterhitung dari subhimpunan kompak medan skalar; atau
     \item[(ii)] sebagai anggota bola satuan tertutup dari~$V^*$.
  \end{enumerate}
Inti pembuktiannya adalah meyakinkan diri bahwa konvergensi suatu jaring fungsi semacam itu dalam ruang hasil kali
tepat ``sama'' dengan konvergensi-$w^*$ jaring tersebut dalam bola satuan tertutup dari~$V^*$.  Upaya untuk sekaligus
mengidentifikasi dan membedakan dua himpunan fungsi ini menantang secara notasi.  Jadi, saran terbaik saya ialah mengerjakannya
sendiri sampai Anda melihat betapa cerdik, tetapi pada dasarnya sederhana, argumen tersebut.

\begin{thm}[Teorema Alaoglu]\label{C067441} Bola satuan tertutup dalam dual suatu ruang linear bernorma
 \index{teorema Alaoglu}%
kompak dalam topologi-$w^*$.
\end{thm}

\begin{proof}[\emph{Petunjuk pembuktian}] Misalkan $V$ suatu ruang linear bernorma, dan nyatakan dengan $[V^*]_1$ bola satuan tertutup dari
ruang dualnya.  Untuk setiap $x \in V$, definisikan $D_x = \{\lambda \in \K \colon \abs\lambda \le \norm x\}$.  Hasil kali
\smash[b]{$\mathbf P = \prod\limits_{x \in V} D_x$} kompak berdasarkan \emph{teorema Tychonoff}.  Tinjau fungsi
   \[ \Phi\colon [V^*]_1 \sto \mathbf P \colon f \mapsto \bigl(f(x)\bigr)_{x \in V} \]
dengan $[V^*]_1$ dilengkapi topologi-$w^*$ dan $\mathbf P$ dilengkapi topologi hasil kali (lihat~\ref{C021437}).  Di sini kita menulis
$\bigl(f(x)\bigr)_{x \in V}$ untuk keluarga bilangan terindeks (atau barisan tergeneralisasi), bukan notasi yang lebih lazim
$\bigl(f_x\bigr)_{x \in V}$.  (Untuk pembahasan lebih lanjut tentang keluarga terindeks, lihat~\cite{Erdman:2007}, Bagian~7.2.)  Mudah dilihat bahwa
$\Phi$ injektif.  Untuk menunjukkan bahwa pemetaan itu kontinu, misalkan $\bigl(f_\lambda\bigr)_{\lambda \in \Lambda}$ suatu jaring dalam $[V^*]_1$
dan $g$ suatu fungsi dalam $[V^*]_1$ sedemikian sehingga $f_\lambda \to^{w^*} g$.  Perhatikan bahwa hal ini terjadi jika dan hanya jika
$\pi_x(f_\lambda) \sto \pi_x(g)$ untuk setiap $x \in V$ (dengan $\pi_x$ proyeksi koordinat pada~$\mathbf P$---lihat~\ref{C015531}).
Terakhir, misalkan $\bigl(f_\lambda\bigr)_{\lambda \in \Lambda}$ suatu jaring dalam $[V^*]_1$ sedemikian sehingga $\bigl(\Phi(f_\lambda)\bigr)_{\lambda \in \Lambda}$
konvergen dalam~$\mathbf P$.  Tunjukkan bahwa jaring $\bigl(f_\lambda(x)\bigr)_{\lambda \in \Lambda}$ konvergen dalam $\K$ untuk setiap $x \in V$.
Misalkan $g(x) = \lim\limits_\lambda f_\lambda(x)$ untuk setiap $x \in V$, lalu buktikan bahwa $g$ linear, $\norm g \le 1$,
$f_\lambda \to^{w^*} g$ dalam $[V^*]_1$, dan $\Phi(f_\lambda) \sto \Phi(g)$ dalam~$\mathbf P$.  \ns
\end{proof}

Sepengetahuan saya, versi pembuktian standar ini yang paling ramah pembaca terdapat dalam~\cite{BrownP:1977}, Teorema 15.11.
Pembuktian yang jauh lebih singkat, jauh lebih mudah secara notasi, jauh lebih elegan, dan jauh kurang mencerahkan tersedia dalam~\cite{Pedersen:1995}, Teorema 2.5.2.
Menurut saya, pembuktian itu kurang mencerahkan karena bergantung pada jaring universal.  Suatu jaring dalam himpunan $S$ disebut
 \index{universal!jaring}%
 \index{jaring!universal}%
\df{universal} jika untuk setiap subhimpunan $T$ dari $S$, jaring tersebut pada akhirnya berada dalam $T$ atau pada akhirnya berada dalam komplemennya. Meskipun
setiap jaring dapat dibuktikan memiliki subjaring universal, bahkan mencari satu contoh nontrivial dari jaring universal (yakni, jaring yang tidak
pada akhirnya konstan) memerlukan \emph{aksioma pilihan}.






\section{Teorema Pemetaan Terbuka}\label{C0694}
Salah satu pertanyaan yang dapat (dan patut) diajukan tentang setiap kategori konkret ialah apakah setiap morfisme bijektif selalu merupakan
isomorfisme.  Dalam beberapa kategori jawabannya jelas \,\emph{ya} (misalnya, dalam $\cat{VEC}$).  Dalam kategori lain
jawabannya jelas \,\emph{tidak} (dalam $\cat{TOP}$, misalnya, tinjau pemetaan identitas dari bilangan
real bertopologi diskret ke bilangan real bertopologi biasa).  Dalam kategori-kategori yang muncul dalam
kajian analisis, aljabar sering menarik dengan penuh semangat ke satu arah, sementara topologi melawan dengan kuat dan menarik
ke arah yang lain.  Di sini pertanyaan tersebut dapat menjadi sangat menarik---dan jawabannya sama sekali tidak sepele.  Dalam kategori
$\cat{BAN_\infty}$ yang terdiri atas ruang Banach dan pemetaan linear kontinu, pertanyaan ini ternyata memiliki kedalaman yang menarik.  Dalam
kasus ini aljabar menang; jawabannya ya.  Hal ini ternyata merupakan akibat langsung dari \emph{teorema
pemetaan terbuka} yang termasyhur.  Ketika Anda menelaah pembuktian hasil berikut, perhatikan bahwa untuk pemetaan yang dibahas,
kelengkapan domain dan kelengkapan kodomain sama-sama esensial---dan dengan alasan yang sangat berbeda.

\begin{defn} Suatu pemetaan $f \colon X \sto Y$ antara ruang-ruang topologis disebut
 \index{terbuka!pemetaan}%
 \index{pemetaan!terbuka}%
\df{terbuka} jika memetakan himpunan terbuka ke himpunan terbuka; yaitu, $f$ merupakan pemetaan terbuka apabila $f^{\sto}(U)$ terbuka dalam $Y$
setiap kali $U$ terbuka dalam~$X$.
\end{defn}

\emph{Teorema pemetaan terbuka} menyatakan bahwa surjeksi linear kontinu antara ruang-ruang Banach selalu merupakan pemetaan terbuka.
Ada satu rincian teknis yang membuat pembuktian hasil ini agak rumit.  Misalkan $T$ suatu pemetaan linear kontinu
dari ruang Banach $B$ ke ruang linear bernorma~$V$, dan $R$ adalah citra oleh $T$ dari suatu bola terbuka di sekitar
titik asal dalam~$B$.  Kita perlu mengetahui bahwa jika tutupan $R$ memuat suatu bola terbuka di sekitar titik asal dalam $V$, maka
$R$ sendiri memuat bola yang sama. Akan memudahkan jika bagian informasi yang krusial ini kita pisahkan sebagai proposisi berikut.

\begin{notn} Dalam ruang linear bernorma $V$, kita menyatakan dengan
 \index{<unop@$(V)_r$ (bola terbuka berjari-jari $r$ di sekitar titik asal)}%
$(V)_r$ bola terbuka di sekitar titik asal dalam $V$ yang berjari-jari~$r$.
\end{notn}

\begin{prop}\label{prop_for_OMT} Misalkan $T\colon B \sto V$ suatu pemetaan linear kontinu dari ruang Banach ke ruang linear bernorma.
Jika $r$, $s > 0$ dan $(V)_s \subseteq \clo{T\bigl(\,(B)_r\,\bigr)}$, maka $(V)_s \subseteq T\bigl(\,(B)_r\,\bigr)$.
\end{prop}

\begin{proof}[\emph{Petunjuk pembuktian}] Pertama, buktikan bahwa tanpa mengurangi keumuman kita dapat mengasumsikan $r = s = 1$.  Dengan demikian,
hipotesis kita adalah $(V)_1 \subseteq \clo{T\bigl(\,(B)_1\,\bigr)}$.

Misalkan $z \in (V)_1$.  Pembuktian selesai jika kita dapat menunjukkan bahwa $z \in T\bigl(\,(B)_1\,\bigr)$.  Pilih $\delta > 0$ sedemikian
sehingga $\norm z < 1 - \delta$, dan misalkan $y = (1 - \delta)^{-1}z$.  Bangun suatu barisan $(v_k)_{k=0}^\infty$ dari vektor-vektor dalam~$V$
yang memenuhi
   \begin{enumerate}
     \item[(i)] $\norm{y - v_k} < \delta^k$ untuk $k = 0,1,2,\dots$; dan
     \item[(ii)] $v_k - v_{k-1} \in T\bigl(\,\delta^{k-1}(B)_1\,\bigr)$ untuk $k = 1,2,3,\dots$.
   \end{enumerate}
Untuk melakukannya, mulailah dengan $v_0 = \vc 0$. Gunakan hipotesis untuk menyimpulkan bahwa $y - v_0 \in \clo{T\bigl(\,(B)_1\,\bigr)}$ dan
karena itu terdapat vektor $w_0 \in T\bigl(\,(B)_1\,\bigr)$ sedemikian sehingga $\norm{(y - v_0) - w_0} < \delta$.  Misalkan
$v_1 = v_0 + w_0$, lalu periksa bahwa (i) dan (ii) berlaku untuk $k = 1$.

Kemudian cari $w_1 \in T\bigl(\,\delta (B)_1\,\bigr)$ sedemikian sehingga $\norm{(y - v_1) - w_1} < \delta^2$, dan misalkan $v_2 = v_1 + w_1$.
Periksa bahwa (i) dan (ii) berlaku untuk $k= 2$.  Lanjutkan secara induktif.

Untuk setiap $n \in \N$, pilih vektor $u_n \in \delta^{n-1}(B)_1$ sedemikian sehingga $T(u_n) = v_n - v_{n-1}$. Tunjukkan bahwa barisan
$(u_n)$ dapat dijumlahkan (gunakan~\ref{prop_abs_sum_implies_sum}), dan misalkan $x = \sum_{n=1}^\infty u_n$.  Selesaikan pembuktian dengan menunjukkan
bahwa $Tx = y$ dan bahwa $z \in T\bigl(\,(B)_1\,\bigr)$.   \ns
\end{proof}

\begin{thm}[Teorema pemetaan terbuka]\label{C069414} Setiap surjeksi linear terbatas antara
 \index{terbuka!pemetaan!teorema}%
 \index{morfisme bijektif!dapat diinverskan!dalam $\cat{BAN_\infty}$}%
ruang-ruang Banach merupakan pemetaan terbuka.
\end{thm}

\begin{proof}[\emph{Petunjuk pembuktian}] Misalkan $T\colon B \sto C$ suatu surjeksi linear terbatas antara ruang-ruang Banach.
Mulailah dengan membuktikan klaim berikut: yang perlu kita tunjukkan hanyalah bahwa $T\bigl((B)_1\bigr)$ memuat suatu bola terbuka di sekitar
titik asal dalam~$C$.  Untuk sementara, misalkan kita mengetahui bahwa bola semacam itu ada.  Namai bola itu~$W$.  Diberikan suatu subhimpunan terbuka $U$ dari $B$, kita
ingin menunjukkan bahwa citranya oleh $T$ terbuka dalam~$C$.  Untuk itu, ambil sebarang titik $y$ dalam $T(U)$ dan pilih $x$
dalam $U$ sedemikian sehingga $y = Tx$.  Karena $U - x$ merupakan lingkungan titik asal dalam $B$, kita dapat mencari bilangan $r > 0$ sedemikian sehingga
$r\,(B)_1 \subseteq U - x$.  Buktikan bahwa $T(U)$ terbuka dengan menunjukkan bahwa $y + rW \subseteq T(U)$.

Untuk membuktikan klaim di atas, misalkan $V = T\bigl((B)_1\bigr)$ dan perhatikan bahwa karena $T$ surjektif, $C = T(B) =
\bigcup_{k=1}^\infty k\clo V$.  Gunakan versi \emph{teorema kategori Baire} yang menyatakan bahwa jika suatu ruang metrik lengkap takkosong
merupakan gabungan suatu barisan himpunan tertutup, maka sekurang-kurangnya salah satu himpunan tersebut memiliki interior takkosong.
(Lihat, misalnya, Teorema 29.3.21 dalam~\cite{Erdman:2007}.)  Pilih $k \in \N$ sedemikian sehingga $\bigl(k\clo V\bigr)^o \neq \emptyset$.
Dengan demikian, terdapat $y \in \clo V$ dan $r > 0$ sedemikian sehingga $y + (C)_r \subseteq k\clo V$.  Dari kekonveksan
$k\clo V$, tidak sulit melihat bahwa $(C)_r \subseteq k\clo V$.  (Tuliskan sebarang titik $z$ dalam $(C)_r$ sebagai
$\frac12(y + z) + \frac12(-y + z)$.) Sekarang diperoleh $(C)_{r/k} \subseteq \clo V$.  Gunakan Proposisi~\ref{prop_for_OMT}.
\ns \end{proof}

\begin{cor}\label{C069417} Setiap bijeksi linear terbatas antara ruang-ruang Banach merupakan suatu
isomorfisme.
\end{cor}

\begin{cor} Setiap surjeksi linear terbatas antara ruang-ruang Banach
 \index{BAN@$\cat{BAN_\infty}$!hasil bagi dalam}%
 \index{hasil bagi!pemetaan!dalam $\cat{BAN_\infty}$}%
merupakan pemetaan hasil bagi dalam~$\cat{BAN_\infty}$.
\end{cor}

\begin{exam} Misalkan $\fml C_u$ himpunan fungsi kontinu bernilai real pada $[0,1]$ dengan
norma seragam, dan $\fml C_1$ himpunan fungsi yang sama dengan norma $\mathrm{L_1}$
(\,$\norm{f}_1 = \int_0^1\abs{f}\,d\lambda$\,). Maka pemetaan identitas $I\colon \fml C_u
\sto \fml C_1$ merupakan bijeksi linear kontinu, tetapi tidak terbuka. (Mengapa hal ini tidak
bertentangan dengan \emph{teorema pemetaan terbuka}?)
\end{exam}

\begin{exam} Misalkan $l$ himpunan barisan bilangan real yang pada akhirnya bernilai nol,
$V$ adalah $l$ yang dilengkapi norma $l_1$, dan $W$ adalah $l$ dengan norma seragam.  Pemetaan
identitas $I\colon V \sto W$ merupakan bijeksi linear terbatas, tetapi bukan suatu isomorfisme.
\end{exam}

\begin{exam} Pemetaan $T\colon \R^2 \sto \R\colon (x,y) \mapsto x$ menunjukkan bahwa meskipun
surjeksi linear terbatas antara ruang-ruang Banach memetakan himpunan terbuka ke himpunan terbuka, pemetaan tersebut tidak harus
memetakan himpunan tertutup ke himpunan tertutup.
\end{exam}

\begin{prop} Misalkan $V$ suatu ruang vektor, dan misalkan $\norm{\cdot}_1$ dan
$\norm{\cdot}_2$ adalah norma-norma pada~$V$ dengan topologi masing-masing $\ofml T_1$
dan~$\ofml T_2$.  Jika $V$ lengkap terhadap kedua norma dan jika $\ofml T_1
\supseteq \ofml T_2$, maka $\ofml T_1 = \ofml T_2$.
\end{prop}

\begin{exer} Apakah terdapat suatu barisan $(a_n)$ dari bilangan kompleks yang memenuhi
syarat berikut: suatu barisan $(x_n)$ dalam $\C$ dapat dijumlahkan secara mutlak jika dan hanya jika
$(a_nx_n)$ terbatas? \emph{Petunjuk.} Tinjau $ T\colon l_\infty \sto l_1\colon (x_n)
\mapsto (x_n/a_n)$.
\end{exer}

\begin{exam}\label{C069431} Misalkan $\fml C^2([a,b])$ keluarga semua fungsi bernilai real yang dua kali
dapat didiferensialkan secara kontinu pada interval $[a,b]$. Pandang keluarga ini sebagai ruang vektor
dengan cara biasa. Untuk setiap $f \in \fml C^2([a,b])$, definisikan
    \[ \norm{f} = \norm{f}_u + \norm{f'}_u + \norm{f''}_u\,. \]
Ini memang suatu norma, dan dengan norma ini $\fml C^2([a,b])$ merupakan ruang Banach.
\end{exam}

\begin{exam}\label{C069432} Misalkan $\fml C^2([a,b])$ ruang Banach yang diberikan dalam
Contoh~\ref{C069431}, dan misalkan $p_0$ dan $p_1$ merupakan anggota $\fml C([a,b])$.   Definisikan
  \[T\colon C^2([a,b]) \sto C([a,b])\colon
                      f \mapsto f'' + p_1f' + p_0f.\]
Maka $T$ merupakan pemetaan linear terbatas.
\end{exam}

\begin{exer} Tinjau suatu persamaan diferensial berbentuk
   \[ y'' + p_1\,y' + p_0\,y = q  \tag{$\ast$} \]
dengan $p_0$, $p_1$, dan $q$ fungsi-fungsi bernilai real kontinu (yang tetap) pada interval
$[a,b]$.  Nyatakan secara tepat, lalu buktikan, pernyataan bahwa \emph{solusi-solusi $(\ast)$
bergantung secara kontinu pada nilai awal}.  Anda boleh menggunakan tanpa bukti teorema standar berikut:
untuk setiap titik $c$ dalam $[a,b]$ dan setiap pasangan bilangan real $a_0$ dan $a_1$,
terdapat solusi tunggal untuk $(\ast)$ sedemikian sehingga $y(c) = a_0$ dan $y'(c) = a_1$.
\emph{Petunjuk.} Untuk $c \in [a,b]$ yang tetap, tinjau pemetaan
  \[ S\colon \fml C^2([a,b]) \sto \fml C([a,b]) \times \R^2 \colon
                                        f \mapsto \bigl(Tf,f(c),f'(c)\bigr) \]
dengan $\fml C^2([a,b])$ ruang Banach dari Contoh~\ref{C069431} dan $T$ pemetaan
linear terbatas dari Contoh~\ref{C069432}.
\end{exer}

\begin{defn} Suatu pemetaan linear terbatas $T\colon V \sto W$ antara ruang-ruang linear bernorma disebut
 \index{terbatas!jauh dari nol}%
\df{terbatas jauh dari nol} (atau
 \index{terbatas!dari bawah}%
\df{terbatas dari bawah}) jika terdapat bilangan $\delta > 0$ sedemikian sehingga $\norm{Tx} \ge
\delta\norm{x}$ untuk setiap $x \in V$. (Ingat bahwa kata ``terbatas'' memiliki satu arti ketika
menerangkan suatu pemetaan linear dan arti yang sangat berbeda ketika diterapkan pada fungsi umum;
demikian pula ungkapan ``terbatas jauh dari nol''. Lihat Definisi~\ref{def_bdd_away_zero}.)
\end{defn}

\begin{prop} Suatu pemetaan linear terbatas antara ruang-ruang Banach bersifat terbatas jauh dari nol jika dan
hanya jika pemetaan itu memiliki jangkauan tertutup dan kernel nol.
\end{prop}


















\section{Teorema Graf Tertutup}
Syarat perlu dan cukup yang sederhana agar suatu pemetaan linear $T\colon B \sto C$ antara ruang-ruang Banach
bersifat kontinu ialah grafnya tertutup.  Dalam satu arah, hal ini hampir jelas.  Kebalikannya
disebut \emph{teorema graf tertutup}.  Ketika kita mensyaratkan graf $T$
 \index{graf!tertutup}%
``tertutup'', maksud kita ialah tertutup sebagai subhimpunan $B \times C$ dengan topologi hasil kali.  Ingat (dari~\ref{exam_prod_top_norm})
bahwa topologi pada jumlah langsung $B \oplus C$ yang diinduksi oleh norma hasil kali sama dengan topologi hasil kali pada $B \times C$.

\begin{exer} Jika $T\colon B \sto C$ suatu pemetaan linear kontinu antara ruang-ruang Banach, maka grafnya tertutup.
\end{exer}

\begin{thm}[Teorema graf tertutup]\label{thm_CGT} Misalkan $T\colon B \sto C$ suatu pemetaan linear antara ruang-ruang Banach.
 \index{tertutup!teorema graf}%
Jika graf $T$ tertutup, maka $T$ kontinu.
\end{thm}

\begin{proof}[\emph{Petunjuk pembuktian}] Misalkan $G = \{(x,Tx)\colon x \in B\}$ graf~$T$. Terapkan
(Korolar~\ref{C069417} dari) \emph{teorema pemetaan terbuka} pada pemetaan
  \[ \pi\colon G \sto B\colon (x,Tx) \mapsto x\,. \]
\ns \end{proof}

Proposisi berikut memberikan syarat perlu dan cukup yang sangat sederhana agar graf suatu pemetaan linear antara ruang-ruang Banach
tertutup.

\begin{prop}\label{prop_nasc_closed_graph} Misalkan $T\colon B \sto C$ suatu pemetaan linear antara ruang-ruang Banach. Maka $T$
kontinu jika dan hanya jika syarat berikut dipenuhi:
   \[ \text{jika $x_n \sto 0$ dalam $B$ dan $Tx_n \sto c$ dalam $C$, maka $c = 0$.} \]
\end{prop}

\begin{prop} Misalkan $T \colon B \sto C$ suatu pemetaan linear antara ruang-ruang Banach sedemikian sehingga jika
$x_n \sto 0$ dalam $B$, maka $(g \circ T)(x_n) \sto 0$ untuk setiap $g \in C^*$.  Maka $T$
terbatas.
\end{prop}

\begin{prop} Misalkan $T \colon B \sto C$ suatu fungsi linear antara ruang-ruang Banach. Definisikan
$T^*f = f \circ T$ untuk setiap $f \in C^*$.  Jika $T^*$ memetakan $C^*$ ke dalam $B^*$, maka $T$
kontinu.
\end{prop}

\begin{defn}\label{C069624} Suatu pemetaan linear $T$ dari ruang vektor ke ruang itu sendiri disebut
 \index{idempoten}%
\df{idempoten} jika $T^2 = T$.
\end{defn}

\begin{prop}\label{prop_cont_idem_op} Suatu pemetaan linear idempoten dari ruang Banach ke ruang itu sendiri
kontinu jika dan hanya jika kernel dan jangkauannya tertutup.
\end{prop}

\begin{proof}[\emph{Petunjuk pembuktian}] Untuk bagian ``jika'', gunakan Proposisi~\ref{prop_nasc_closed_graph}.  \ns
\end{proof}

Proposisi berikut merupakan penerapan \emph{teorema graf tertutup} yang berguna pada operator ruang Hilbert.

\begin{prop} Misalkan $H$ suatu ruang Hilbert, serta $S$ dan $T$ fungsi dari $H$ ke dirinya sendiri sedemikian sehingga
    \[ \langle Sx,y \rangle = \langle x,Ty \rangle \]
untuk setiap $x$, $y \in H$.  Maka $S$ dan $T$ merupakan operator linear terbatas.
\end{prop}

\begin{proof}[\emph{Petunjuk pembuktian}] Untuk membuktikan linearitas pemetaan $S$, tinjau hasil kali dalam dari vektor-vektor
$S(\alpha x + \beta y) - \alpha Sx - \beta Sy$ dan $z$, dengan $x$, $y$, $z \in H$ dan $\alpha$, $\beta \in \K$.
Untuk membuktikan kekontinuan $S$, tunjukkan bahwa grafnya tertutup.  \ns
\end{proof}

Ingat bahwa dalam proposisi sebelumnya, operator $T$ adalah \emph{adjoin (ruang Hilbert)} dari $S$;
yaitu, $T = S^*$.

\begin{defn} Misalkan $a < b$ dalam~$\R$. Suatu fungsi $f\colon [a,b] \sto \R$ disebut
 \index{dapat didiferensialkan secara kontinu}%
 \index{diferensiabel!secara kontinu}%
dapat didiferensialkan secara kontinu jika fungsi itu dapat didiferensialkan pada (suatu himpunan terbuka yang memuat) $[a,b]$
dan turunannya $f\,'$ kontinu pada~$[a,b]$.  Himpunan semua fungsi bernilai real yang dapat
didiferensialkan secara kontinu pada $[a,b]$ dinotasikan dengan $\fml C^1(\,[a,b]\,)$.
\end{defn}

\begin{exam} Misalkan $D\colon \fml C^1(\,[0,1]\,) \sto C(\,[0,1]\,)\colon f \mapsto f\,'$, dengan
baik $C^1(\,[0,1]\,)$ maupun $C(\,[0,1]\,)$ dilengkapi norma seragam. Pemetaan
$D$ linear dan memiliki graf tertutup, tetapi tidak kontinu.  (Mengapa hal ini tidak bertentangan dengan
\emph{teorema graf tertutup}?  Apa yang dapat kita simpulkan tentang~$\fml C^1(\,[0,1]\,)$?)
\end{exam}

\begin{prop} Misalkan $S\colon A \sto B$ dan $T\colon B \sto C$ pemetaan linear antara
ruang-ruang Banach, serta $T$ terbatas dan injektif.  Maka $S$ terbatas jika dan hanya jika $TS$ terbatas.
\end{prop}











\section{Dualitas Ruang Banach}

\begin{conv}\label{conv_subsp_Bsp} Dalam konteks ruang Banach, kita menggunakan konvensi yang sama untuk pemakaian kata
``subruang'' seperti yang kita gunakan untuk ruang Hilbert. Kata ``subruang'' akan selalu berarti
 \index{subruang!dari ruang Banach}%
 \index{konvensi!subruang ruang Banach bersifat tertutup}%
\emph{subruang vektor tertutup}. Untuk menunjukkan bahwa $M$ adalah subruang dari suatu ruang Banach $B$, kita menulis
 \index{<binrelle@$\preccurlyeq$ (subruang dari ruang Banach atau Hilbert)}%
$M \preccurlyeq B$.
\end{conv}

\begin{prop} Jika $M$ adalah subruang dari suatu ruang Banach $B$,
 \index{hasil bagi!ruang Banach}%
 \index{Banach!ruang!hasil bagi}%
 \index{ruang!hasil bagi!Banach}%
maka ruang linear bernorma hasil bagi $B/M$ juga merupakan ruang Banach.
\end{prop}

\begin{proof}[\emph{Petunjuk pembuktian}] Untuk menunjukkan bahwa $B/M$ lengkap, cukup ditunjukkan bahwa dari
sebarang barisan Cauchy $(u_n)$ dalam $B/M$ kita dapat mengambil suatu subbarisan yang konvergen. Untuk itu,
perhatikan bahwa jika $(u_n)$ adalah barisan Cauchy, maka untuk setiap $k \in \N$ terdapat $N_k \in \N$ sedemikian sehingga
$\norm{u_n - u_m} < 2^{-k}$ setiap kali $m$, $n \ge N_k$. Sekarang gunakan induksi. Pilih $n_1 = N_1$.
Setelah $n_1$, $n_2$, \dots, $n_k$ dipilih sedemikian sehingga $n_1 < n_2 < \dots < n_k$ dan $N_j \le n_j$ untuk
$1 \le j \le k$, pilih $n_{k+1} = \max\{N_{k+1}, n_k + 1\}$. Dengan demikian diperoleh subbarisan
$\bigl(u_{n_k}\bigr)$ yang memenuhi $\norm{u_{n_k} - u_{n_{k+1}}} < 2^{-k}$ untuk setiap $k \in \N$.
Dari setiap kelas ekuivalensi $u_{n_k}$, pilih suatu vektor $x_k$ dalam $B$ sedemikian sehingga $\norm{x_k - x_{k+1}} < 2^{-k+1}$.
Kemudian Proposisi~\ref{prop_abs_sum_implies_sum} menunjukkan bahwa barisan $(x_k - x_{k+1})$ dapat
dijumlahkan. Dari sini, simpulkan bahwa barisan $(x_k)$ dan $\bigl(u_{n_k}\bigr)$ konvergen.
(Rincian argumen ini dituliskan dengan baik dalam pembuktian Teorema 3.4.13 pada~\cite{BrownP:1970}.) \ns
\end{proof}

\begin{prop} Jika $J$ adalah ideal tertutup dalam suatu aljabar Banach $A$,
 \index{hasil bagi!aljabar Banach}%
 \index{Banach!aljabar!hasil bagi}%
 \index{aljabar!hasil bagi!Banach}%
maka $A/J$ merupakan aljabar Banach.
\end{prop}

\begin{thm}[Teorema hasil bagi fundamental untuk $\cat{BALG}$] Misalkan $A$ dan $B$ aljabar Banach
 \index{BALG@$\cat{BALG}$!hasil bagi dalam}%
 \index{hasil bagi!teorema!untuk $\cat{BALG}$}%
dan $J$ ideal tertutup sejati dalam~$A$. Jika $\phi$ suatu homomorfisme dari $A$ ke $B$ dan
$\ker\phi \supseteq J$, maka terdapat homomorfisme tunggal $\wt\phi\colon A/J \sto
B$ yang membuat diagram berikut komutatif.
   \[ \xymatrix@+30pt{A \ar[d]_*+{\pi} \ar[dr]^*+{\phi} & \\
                           A/J \ar[r]_*+{\widetilde\phi} & B  }\]
Selain itu, $\wt\phi$ injektif jika dan hanya jika $\ker\phi = J$; dan $\wt\phi$
surjektif jika dan hanya jika $\phi$ surjektif.
\end{thm}

\begin{prop}\label{C037821} Misalkan $J$ ideal tertutup sejati dalam suatu aljabar Banach komutatif
beridentitas~$A$. Maka $J$ maksimal jika dan hanya jika $A/J$ adalah medan.
\end{prop}

Proposisi berikut menyatakan syarat-syarat yang cukup sederhana agar suatu epimorfisme antara ruang Banach
ternyata menjadi faktor dalam morfisme-morfisme lain. Sepintas lalu, kegunaan hasil semacam ini
mungkin tampak agak meragukan. Namun, ketika membuktikan Proposisi~\ref{00151231}, perhatikan bahwa
inilah \emph{tepatnya} yang Anda perlukan.

\begin{prop}\label{prop_surj_as_factor} Tinjau kategori ruang Banach
dan pemetaan linear kontinu. Jika $T\colon A \sto C$, $R\colon A \sto B$, $R$
surjektif, dan $\ker R \subseteq \ker T$, maka terdapat
$S\colon B \sto  C$ tunggal sedemikian sehingga $SR = T$. Selain itu, $S$ injektif jika dan
hanya jika $\ker R = \ker T$.
\end{prop}

\begin{proof}[\emph{Petunjuk pembuktian}] Tinjau diagram komutatif berikut.
Gunakan \emph{teorema pemetaan terbuka}.
 \[\xy
    \scalefactor{1.4}%
    \Atrianglepair/>`>`>`<-`>/[A`B`A/\ker R`C;R`\pi`T`\wt R`\wt T]
 \endxy\]   \ns
\end{proof}

Menarik untuk diperhatikan bahwa hasil sebelumnya dapat dinyatakan dengan cara yang
sedikit memperkuat klaim keunikannya. Sebagaimana dinyatakan, kesimpulannya ialah bahwa terdapat
pemetaan linear terbatas $S$ tunggal sedemikian sehingga $SR = T$. Tanpa kerja tambahan, kita dapat merumuskan ulang
proposisi tersebut sebagai berikut.

\begin{prop}\label{prop_surj_as_factor2} Tinjau kategori ruang Banach
dan pemetaan linear kontinu. Jika $T\colon A \sto C$, $R\colon A \sto B$, $R$
surjektif, dan $\ker R \subseteq \ker T$, maka terdapat fungsi tunggal
$S$ yang memetakan $B$ ke $C$ sedemikian sehingga $S\circ R = T$. Fungsi $S$ harus
terbatas dan linear. Selain itu, fungsi tersebut injektif jika dan hanya jika $\ker R = \ker T$.
\end{prop}

Apakah versi kedua hasil ini, yang lebih terperinci, layak dinyatakan? Menurut saya, jawabannya bergantung pada
arah yang hendak dituju. Memilih versi yang lebih panjang hanya dengan alasan bahwa versi itu lebih umum
tampaknya tidak menarik. Di pihak lain, andaikan kemudian kita memerlukan hasil berikut.

\begin{prop}\label{prop_factor_blm} Jika suatu pemetaan linear terbatas $T\colon A \sto C$ antara ruang-ruang Banach dapat difaktorkan
menjadi komposisi dua fungsi $T = f \circ R$, dengan $R \in \ofml B(A,B)$ surjektif,
maka $f \colon B \sto C$ terbatas dan linear.
\end{prop}

\begin{proof}[\emph{Petunjuk pembuktian}] Untuk menerapkan~\ref{prop_surj_as_factor2}, yang diperlukan di sini hanyalah
memperhatikan bahwa $f(0) = 0$ sehingga $\ker R \subseteq \ker T$.   \ns
\end{proof}

Karena hasil sebelumnya merupakan akibat yang hampir langsung dari~\ref{prop_surj_as_factor2}, ada baiknya
versi yang lebih terperinci ini tersedia. Orang mungkin berpendapat bahwa meskipun~\ref{prop_factor_blm} barangkali tidak
sepenuhnya jelas dari pernyataan~\ref{prop_surj_as_factor}, hasil itu tentu jelas dari
pembuktiannya, jadi untuk apa memakai versi yang lebih rumit? Hal-hal semacam ini tentu merupakan soal selera.

\begin{defn}\label{defn_Bsp_exact_seq} Suatu barisan ruang Banach dan pemetaan linear kontinu
   \[\xymatrix{{\cdots}\ar[r] & B_{n-1}\ar[r]^{T_n}
          & B_n\ar[r]^{T_{n+1}} & B_{n+1}\ar[r] & {\cdots}
   }\]
dikatakan
 \index{barisan!eksak}%
 \index{eksak!barisan}%
\df{eksak di} $B_n$ jika $\ran T_n = \ker T_{n+1}$. Suatu barisan disebut \df{eksak} jika barisan tersebut
eksak di setiap ruang Banach penyusunnya. Suatu barisan ruang Banach dan
pemetaan linear kontinu berbentuk
   \begin{equation}\label{001500202i}\xymatrix{
      0 \ar[r] & A \ar[r]^S & B\ar[r]^T & C\ar[r] & 0
   }\end{equation}
 \index{barisan!eksak pendek}%
 \index{barisan eksak pendek}%
yang eksak di $A$, $B$, dan $C$ disebut \df{barisan eksak pendek}. (Di sini $\vc 0$ menyatakan ruang Banach trivial
berdimensi nol, dan panah-panah tanpa label adalah pemetaan yang jelas.)
\end{defn}

Definisi-definisi sebelumnya diberikan untuk
 \index{BAN@$\cat{BAN_\infty}$!kategori}%
 \index{kategori!$\cat{BAN_\infty}$ sebagai suatu}%
kategori $\cat{BAN_\infty}$ dari ruang Banach dan pemetaan linear kontinu. Tentu saja,
gagasan tentang suatu \emph{barisan eksak} bermakna dalam banyak situasi---misalnya, dalam
kategori aljabar Banach, aljabar-$C^*$, ruang Hilbert, ruang vektor, grup Abel,
dan modul, di antara yang lainnya.

\begin{prop}\label{00151231} Tinjau diagram berikut dalam kategori ruang Banach
dan pemetaan linear kontinu
 \[\xymatrix{
      0\ar[r] & A\ar[d]^f\ar[r]^j & B\ar[d]^g\ar[r]^k
              & C\ar@{-->}[d]^h \ar[r] & 0 \\
      0\ar[r] & A'\ar[r]_{j'} & B'\ar[r]_{k'}
      & C'\ar[r] & 0
 }\]
Jika baris-barisnya eksak dan persegi kirinya komutatif, maka terdapat pemetaan linear kontinu
$h\colon C \sto C'$ tunggal yang membuat persegi kanannya komutatif.
\end{prop}

\begin{prop}[Lema Lima]\label{00151232}
 \index{lema lima}%
Andaikan dalam diagram ruang Banach dan pemetaan linear kontinu berikut
 \[\xymatrix{
      0\ar[r] & A\ar[d]^f\ar[r]^j & B\ar[d]^g\ar[r]^k
              & C\ar[d]^h \ar[r] & 0 \\
      0\ar[r] & A'\ar[r]_{j'} & B'\ar[r]_{k'}
      & C'\ar[r] & 0
 }\]
baris-barisnya eksak dan persegi-perseginya komutatif. Maka pernyataan-pernyataan berikut berlaku.
 \begin{enumerate}[label=\rm{(\alph*)}]
  \item Jika $g$ surjektif, maka $h$ juga surjektif.
  \item Jika $f$ surjektif dan $\ran j \supseteq \ker g$, maka $h$ injektif.
  \item Jika $f$ dan $h$ surjektif, maka $g$ juga surjektif.
  \item Jika $f$ dan $h$ injektif, maka $g$ juga injektif.
 \end{enumerate}
\end{prop}

Ada banyak keuntungan bekerja dengan operator pada ruang Hilbert atau Banach yang memiliki jangkauan tertutup.
Salah satu keuntungan langsung ialah bahwa kita memperoleh kembali kesederhanaan \emph{teorema dasar aljabar linear}
(\ref{00016025}) dari versi-versi yang agak lebih rumit dalam Proposisi~\ref{prop_ker_adj_perp_ran_op}
dan Teorema~\ref{prop_FTLA_Bsp}. Di pihak lain, ada pula kerugiannya. Operator-operator itu tidak membentuk suatu kategori; komposisi
dua operator berjangkauan tertutup dapat gagal memiliki jangkauan tertutup.

\begin{prop} Jika suatu pemetaan linear terbatas $T\colon B \sto C$ antara dua ruang Banach memiliki jangkauan tertutup, maka
adjoinnya~$T^*$ juga demikian dan $\ran T^* = (\ker T)^\perp$.
\end{prop}

\begin{proof}[\emph{Petunjuk pembuktian}] Simpulkan dari Proposisi~\ref{prop_FTLA_Bsp}(d) bahwa kita hanya perlu menunjukkan bahwa
$(\ker T)^\perp \subseteq \ran T^*$. Misalkan $f$ suatu fungsional sebarang dalam $(\ker T)^\perp$. Yang harus dicari adalah suatu
fungsional $g$ dalam $C^*$ sedemikian sehingga $f = T^*g$. Untuk itu, gunakan \emph{teorema pemetaan terbuka} untuk menunjukkan bahwa pemetaan $\wt T_\ast$
dalam diagram berikut invertibel. (Pemetaan $T_\ast\colon B \sto \ran T$ dan $T$ sama persis, kecuali bahwa
kodomainnya berbeda.)

 \[\xy
    \scalefactor{1.4}%
    \Atrianglepair/>`>`>`<-`>/[B`\ran T`B/\ker T`\K;T_\ast`\pi`f`\wt{T_\ast}`\wt f]
 \endxy\]

Misalkan $g_0 := \wt f\,{\wt T_\ast}^{-1}$. Gunakan \emph{teorema Hahn--Banach} untuk menghasilkan~$g$ yang diinginkan.   \ns
\end{proof}

\begin{prop} Dalam kategori $\cat{BAN_\infty}$, jika barisan $\xymatrix{A \ar[r]^S & B \ar[r]^T & C}$ eksak di $B$
dan $T$ memiliki jangkauan tertutup, maka $\xymatrix{C^* \ar[r]^{T^*} & B^* \ar[r]^{S^*} & A^*}$ eksak di~$B^*$.
\end{prop}

\begin{defn} Suatu funktor dari kategori ruang Banach dan pemetaan linear kontinu ke dirinya sendiri disebut
 \index{eksak!funktor}%
 \index{funktor!eksak}%
\df{eksak} jika funktor tersebut membawa barisan eksak pendek ke barisan eksak pendek.
\end{defn}

\begin{exam}\label{exam_dualfunctor_exact} Funktor $B \mapsto B^*$, $T \mapsto T^*$ (lihat~\ref{exam_dual_functor})
dari kategori ruang Banach dan pemetaan linear kontinu ke dirinya sendiri bersifat eksak.
\end{exam}

\begin{exam} Misalkan $B$ suatu ruang Banach dan $\sbsb{\tau}B\colon x \mapsto x^{**}$ pembenaman alami dari $B$ ke~$B^{**}$
(lihat Definisi~\ref{defn_nat_embed}). Maka ruang hasil bagi $B^{**}/\ran\sbsb{\tau}B$ merupakan ruang Banach.
\end{exam}

\begin{notn} Kita akan menyatakan dengan $\clo B$ ruang hasil bagi
$B^{**}/\ran \sbsb{\tau}B$ dalam contoh sebelumnya.
\end{notn}

\begin{prop} Jika $T\colon B \sto C$ suatu pemetaan linear kontinu antara ruang-ruang Banach, maka terdapat pemetaan linear kontinu tunggal
$\clo T\colon \clo B \sto \clo C$ sedemikian sehingga $\clo T(\phi + \ran\sbsb{\tau}B) = (T^{**}\phi) + \ran\sbsb{\tau}C$ untuk setiap $\phi \in B^{**}$.
\end{prop}

\begin{proof}[\emph{Petunjuk pembuktian}] Terapkan~\ref{00151231} pada diagram berikut:
 \begin{equation}\label{eqn_exactCD_Bbar}\xymatrix{
      0\ar[r] & B\ar[d]^T\ar[r]^{\sbsb{\tau}B} & B^{**}\ar[d]^{T^{**}}\ar[r]^{\sbsb{\pi}B}
              & \clo B\ar@{-->}[d]^{\clo T} \ar[r] & 0 \\
      0\ar[r] & C\ar[r]_{\sbsb{\tau}C} & C^{**}\ar[r]_{\sbsb{\pi}C} & \clo C\ar[r] & 0
 }\end{equation}  \ns
\end{proof}

\begin{exam}\label{exam_bar_functor} Pasangan pemetaan $B \mapsto \clo B$, $T \mapsto \clo T$ dari
kategori $\cat{BAN_\infty}$ ke dirinya sendiri merupakan suatu
 \index{<unopu@$B \mapsto \clo B$, $T \mapsto \clo T$ (funktor)}%
 \index{funktor!dari suatu ruang Banach $B$ ke $\clo B$}%
funktor kovarian eksak.
\end{exam}

\begin{prop} Dalam kategori ruang Banach dan pemetaan linear kontinu, jika
barisan
 \[\begin{CD}
  0 \longrightarrow A @> S >> B @> T >> C \longrightarrow 0
  \end{CD}\]
eksak, maka diagram berikut komutatif dan memiliki baris-baris
serta kolom-kolom yang eksak.
 \[\bfig
    \iiixiii{'7777}[A`A^{**}`\overline A`B`B^{**}`\overline B`C`C^{**}`\overline C;%
``````S`{S^{**}}`{\overline S}`T`{T^{**}}`{\overline T}]
  \efig\]
\end{prop}

\begin{cor} Setiap subruang dari suatu ruang Banach refleksif bersifat refleksif.
\end{cor}

\begin{cor} Jika $M$ suatu subruang dari ruang Banach refleksif~$B$, maka $B/M$ refleksif.
\end{cor}

\begin{cor} Misalkan $M$ suatu subruang dari ruang Banach~$B$. Jika $M$ dan $B/M$ keduanya refleksif,
maka $B$ juga refleksif.
\end{cor}



















\begin{defn} Misalkan $T \in \ofml L(B,C)$, dengan $B$ dan $C$ ruang Banach. Definisikan
$\coker T$, yaitu
 \index{kokernel}%
 \index{kokernel@$\coker T$ (kokernel dari $T$)}%
\df{kokernel} dari~$T$, sebagai $C/\ran T$.
\end{defn}



Menentukan secara eksplisit dual dari suatu ruang Banach sering kali agak menantang. Berikut diberikan beberapa
contoh penting. Namun, pertama-tama perlu disampaikan satu peringatan. Notasi untuk objek-objek dalam ruang linear bernorma
yang terdiri atas barisan dapat menyesatkan. Ketika membahas kelengkapan ruang-ruang semacam itu, kita berhadapan dengan barisan dari
barisan bilangan kompleks. Demi memudahkan pembaca (dan mungkin juga diri Anda sendiri sebulan kemudian,
ketika Anda mencoba membaca sesuatu yang Anda tulis!), sebaiknya bedakan secara notasi antara barisan bilangan kompleks
dan barisan dari barisan. Ada banyak cara untuk melakukannya: subskrip ganda, subskrip dan
superskrip, notasi fungsional untuk barisan kompleks, \emph{dan sebagainya}.

Notasi bermasalah lain yang muncul dalam ruang barisan adalah $\sum_{k=1}^\infty a_k \vc e^k$, dengan $\vc e^k$
merupakan barisan yang bernilai $1$ pada koordinat ke-$k$ dan $0$ di semua koordinat lainnya. (Ingat bahwa dalam
Contoh~\ref{C063149}, kita menggunakan vektor-vektor ini sebagai basis ortonormal biasa (atau standar) untuk ruang Hilbert
$l_2$.) Tampaknya hampir tidak ada kebingungan yang timbul jika $\sum_{k=1}^\infty a_k \vc e^k$ dipandang
hanya sebagai singkatan untuk barisan bilangan kompleks $(a_k) = (a_1,a_2,a_3,\dots)$. Namun, jika seseorang hendak
memperlakukannya secara harfiah sebagai deret tak hingga, diperlukan kehati-hatian. Dalam ruang $\sbsb c0$ dan $l_1$, tidak ada
masalah. Deret itu memang merupakan limit, ketika $n$ membesar, dari $\sum_{k=1}^n a_k \vc e^k$. (Mengapa?) Namun, dalam ruang
$c$ dan $l_\infty$, misalnya, terjadi kegagalan yang sangat buruk. Untuk melihat apa yang dapat terjadi, tinjau barisan konstan
$(1,1,1,\dots)$ (yang merupakan anggota $c$ maupun $l_\infty$). Maka
$\lim_{n \sto\infty}\norm{(1,1,1,\dots) - \sum_{k=1}^n \vc e^k}$ gagal mendekati nol dalam kedua kasus tersebut.
Dalam $c$ maupun $l_\infty$, limitnya adalah~$1$.

\begin{exam}\label{exam_dual_C0} Ruang linear bernorma $\sbsb c0$ (lihat~\ref{C0231303})
 \index{c@$c_0$!sebagai ruang Banach}%
 \index{Banach!ruang!$c_0$ sebagai suatu}%
 \index{l@$l_1$!sebagai dual dari $c_0$}%
 \index{dual!dari $c_0$}%
merupakan ruang Banach dan dualnya adalah $l_1$ (lihat~\ref{C023191}).
\end{exam}

\begin{proof}[\emph{Petunjuk pembuktian}] Tinjau pemetaan
$T\colon {\sbsb c0}^* \sto l_1\colon f \mapsto \sum_{k=1}^\infty f(\vc e^k)\vc e^k$.
(Jangan lupa menunjukkan bahwa $T$ benar-benar memetakan ke~$l_1$.) Periksa bahwa $T$ merupakan isomorfisme isometrik.
Untuk membuktikan bahwa $T$ surjektif, mulailah dengan suatu vektor $b \in l_1$ dan tinjau fungsi
$f\colon \sbsb c0 \sto \C\colon a \mapsto \sum_{k=1}^\infty a_kb_k$.    \ns
\end{proof}

\begin{cor} Ruang $l_1$ merupakan suatu
 \index{l@$l_1$!sebagai ruang Banach}%
 \index{Banach!ruang!$l_1$ sebagai suatu}%
ruang Banach.
\end{cor}

Menarik bahwa ternyata ruang $\sbsb c0$ dan $c$ tidak isomorfik secara isometrik,
meskipun ruang-ruang dualnya isomorfik secara isometrik. Perhitungan dual dari $c$ secara teknis sedikit lebih rumit daripada
perhitungan~${\sbsb c0}^{\!*}$. Mungkin berguna untuk meninjau \emph{basis Schauder} bagi kedua ruang ini.

\begin{defn} Suatu barisan $(\vc e^k)$ dari vektor-vektor dalam suatu ruang Banach $B$ disebut
 \index{basis Schauder}%
 \index{basis!Schauder}%
\df{basis Schauder} bagi ruang tersebut jika untuk setiap $x \in B$ terdapat barisan skalar $(\alpha_k)$ yang unik
sedemikian sehingga $x = \sum_{k=1}^\infty \alpha_k \vc e^k$.
\end{defn}

Dua catatan mengenai definisi sebelumnya: Kata ``unik'' sangat penting; dan tidak benar bahwa
setiap ruang Banach memiliki basis Schauder.

\begin{exam} Suatu basis ortonormal dalam ruang Hilbert separabel merupakan basis Schauder bagi ruang tersebut. (Lihat
Proposisi~\ref{prop_unique_onbasis} dan~\ref{prop_Schauder_basis_Hsp}.)
\end{exam}

\begin{exam}\label{exam_dual_c0} Kita memilih
 \index{biasa!basis Schauder untuk~$\sbsb c0$}%
 \index{Schauder!basis!biasa (atau standar) untuk~$\sbsb c0$}%
 \index{basis!standar (atau biasa)}%
vektor-vektor yang disebut \df{standar} (atau
 \index{standar!basis Schauder untuk $\sbsb c0$}%
\df{biasa}) sebagai \df{vektor-vektor basis} untuk $\sbsb c0$ dengan cara yang sama seperti untuk~$l_2$ (lihat~\ref{C063149}).
Untuk setiap $n \in \N$, misalkan $\vc e^n$ adalah barisan dalam $l_2$ yang koordinat ke-$n$-nya bernilai $1$
dan semua koordinat lainnya bernilai~$0$. Maka $\{\vc e^n\colon n \in \N\}$ merupakan basis Schauder
bagi~$\sbsb c0$.
\end{exam}

Sepintas lalu, ruang $c$ dari semua barisan konvergen mungkin tampak jauh lebih besar daripada $\sbsb c0$,
yang hanya memuat barisan-barisan yang konvergen ke nol. Lagi pula, untuk setiap vektor
$x$ dalam $\sbsb c0$ terdapat tak terhitung banyaknya vektor berbeda dalam $c$ (cukup tambahkan sebarang bilangan kompleks
tak nol pada setiap suku barisan~$x$). Mungkinkah $c$, yang memuat tak terhitung banyaknya
salinan dari koleksi anggota~$\sbsb c0$ yang tak terhitung, memiliki basis Schauder (yang mesti terhitung)?
Setelah direnungkan, jawabannya adalah \emph{ya}. Kita hanya perlu menambahkan satu elemen
pada basis bagi~$\sbsb c0$ untuk menangani translasi.

\begin{exam} Misalkan $\vc e^1$, $\vc e^2$, \dots seperti dalam Contoh~\ref{exam_dual_c0} sebelumnya. Selain itu,
nyatakan dengan $\vc 1$ barisan konstan $(1,1,1,\dots)$ dalam~$c$. Maka $(\vc 1, \vc e^1, \vc e^2, \dots)$
merupakan basis Schauder bagi~$c$.
\end{exam}

\begin{exam}\label{exam_dual_c} Ruang linear bernorma $c$ (lihat~\ref{C0231302})
 \index{c@$c$!sebagai ruang Banach}%
 \index{Banach!ruang!$c$ sebagai suatu}%
 \index{l@$l_1$!sebagai dual dari $c$}%
 \index{dual!dari $c$}%
merupakan ruang Banach dan dualnya adalah $l_1$ (lihat~\ref{C023191}).
\end{exam}

\begin{proof}[\emph{Petunjuk pembuktian}]
\leavevmode\par\noindent
Pertama-tama, perhatikan bahwa untuk setiap $f \in c^*$, barisan
$\bigl(f(\vc e^k)\bigr)$ dapat dijumlahkan secara mutlak dan, sebagai akibatnya, $\sum_{k=1}^\infty a_k f(\vc e^k)$
konvergen untuk setiap barisan $a = (a_1,a_2,\dots)$ dalam~$c$. Perhatikan pula bahwa setiap $a \in c$ dapat dituliskan sebagai
$a = a_\infty \vc 1 + \sum_{k=1}^\infty (a_k -a_\infty)\vc e^k$, dengan $a_\infty := \lim_{k \sto \infty}a_k$
dan $\vc 1 := (1,1,1,\dots)$. Gunakan hal ini untuk membuktikan bahwa untuk setiap $f \in c^*$,
   \[ \abs{f(a)} \le \biggl[ \abs{\beta_f} + \sum_{k=1}^\infty \abs{f(\vc e^k)} \biggr] {\norm a}_\infty \]
dengan $\beta_f := f(\vc 1) - \sum_{k=1}^\infty f(\vc e^k)$. Untuk $t = (t_0,t_1,t_2,\dots) \in l_1$ dan
$a = (a_1,a_2,a_3,\dots) \in c$, definisikan
   \[ S_t(a) := t_0 a_\infty + \sum_{k=1}^\infty t_ka_k \]
dan periksa bahwa $S\colon t \mapsto S_t$ merupakan pemetaan linear kontinu dari $l_1$ ke~$c^*$.
Selanjutnya, untuk $f \in c^*$, definisikan
   \[ Tf := (\beta_f, f(\vc e^1), f(\vc e^2), \dots) \]
dan tunjukkan bahwa $T\colon f \mapsto Tf$ merupakan isometri linear dari $c^*$ ke~$l_1$. Sekarang, untuk melihat bahwa $c^*$ dan
$l_1$ isomorfik secara isometrik, cukup ditunjukkan bahwa $ST = \id{c^*}$ dan $TS = \id{l_1}$.  \ns
\end{proof}












\begin{defn} Suatu ruang topologis disebut
 \index{lokal!kompak}%
 \index{kompak!lokal}%
\df{kompak lokal} jika setiap titik dalam ruang tersebut memiliki lingkungan kompak.
\end{defn}

\begin{notn} Misalkan $X$ suatu ruang kompak lokal. Suatu fungsi bernilai skalar $f$ pada $X$ termasuk dalam $\fml C_0(X)$ jika
 \index{C@$\fml C_0(X)$!ruang fungsi kontinu yang lenyap di tak hingga}%
fungsi tersebut kontinu dan lenyap di tak hingga (yakni, jika untuk setiap bilangan $\epsilon > 0$ terdapat subhimpunan kompak
dari $X$ yang di luarnya berlaku $\abs{f(x)} < \epsilon$).
\end{notn}

\begin{exam}\label{exam_C0X} Jika $X$ suatu ruang Hausdorff kompak lokal, maka, terhadap operasi aljabar titik demi titik dan norma seragam,
 \index{Banach!ruang!$\fml C_0(X)$ sebagai suatu}%
 \index{C@$\fml C_0(X)$!sebagai ruang Banach}%
$\fml C_0(X)$ merupakan ruang Banach.
\end{exam}

Kita mengingat kembali salah satu teorema terpenting dari analisis real.

\begin{thm}[Representasi Riesz]\label{C057951} Misalkan $X$ suatu ruang Hausdorff kompak lokal.
 \index{Riesz!teorema representasi}%
Jika $\phi$ suatu fungsional linear terbatas pada $\fml C_0(X)$, maka terdapat ukuran Borel kompleks reguler
$\mu$ yang unik pada $X$ sedemikian sehingga
  \[ \phi(f) = \int_X f\,d\mu \]
untuk setiap $f \in \fml C_0(X)$ dan $\norm\mu = \norm\phi$.
\end{thm}

\begin{proof} Teorema ini dan akibat berikutnya disajikan dengan baik dalam Lampiran C pada buku analisis
fungsional karya Conway~\cite{Conway:1990}. Lihat juga~\cite{Folland:1984}, Teorema 7.17; \cite{HewittS:1965},
Teorema 20.47 dan 20.48; \cite{McDonaldW:1999}, Teorema 9.16; dan \cite{Rudin:1987}, Teorema 6.19.  \ns
\end{proof}

\begin{cor}\label{C057952} Jika $X$ suatu ruang Hausdorff kompak lokal, maka $\bigl(\fml C_0(X)\,\bigr)^*$
 \index{dual!dari $\fml C_0(X)$}%
 \index{C@$\fml C_0(X)$!dual dari}%
 \index{M@$\textbf{M}(X)$!sebagai dual dari $\fml C_0(X)$}%
isomorfik secara isometrik dengan~$\textbf{M}(X)$,
 \index{M@$\textbf{M}(X)$!ukuran Borel kompleks reguler pada~$X$}%
ruang Banach dari semua ukuran Borel kompleks reguler pada~$X$.
\end{cor}

\begin{cor} Jika $X$ suatu ruang Hausdorff kompak lokal, maka $\textbf{M}(X)$
 \index{M@$\textbf{M}(X)$!sebagai ruang Banach}%
 \index{Banach!ruang!$\textbf{M}(X)$ sebagai suatu}%
merupakan ruang Banach.
\end{cor}










Dalam~\ref{defn_nls_adjoint}, kita mendefinisikan adjoin dari suatu operator linear antara ruang-ruang Banach,
dan dalam~\ref{def000926}, kita mendefinisikan adjoin dari suatu operator ruang Hilbert.
Dalam kasus operator pada ruang Hilbert (yang juga merupakan ruang Banach),
wajar untuk menanyakan hubungan apa, jika ada, di antara kedua ``adjoin'' ini.
Keduanya tentu tidak dapat sama, karena yang pertama bekerja di antara ruang-ruang dual, sedangkan yang kedua
di antara ruang-ruang semula. Dalam proposisi berikut, kita menggunakan antiisomorfisme
$\psi$ yang didefinisikan dalam~\ref{000341} (lihat~\ref{0022057}) untuk menunjukkan bahwa kedua adjoin itu
``pada hakikatnya'' sama.

\begin{prop} Misalkan $T$ suatu operator pada ruang Hilbert $H$ dan $\psi$ antiisomorfisme yang
didefinisikan dalam~\ref{000341}. Jika kita (untuk sementara) menyatakan dual ruang Banach
dari $T$ dengan $T'\colon H^* \sto H^*$, maka $T' = \psi\, T^* \psi^{-1}$. Artinya,
diagram berikut komutatif.
   \[ \xy
          \Square/{->}`{<-}`{<-}`{->}/[H^*`H^*`H`H;T'`\psi`\psi`T^*]
      \endxy \]
\end{prop}

\begin{defn}\label{0022091} Misalkan $A$ suatu subhimpunan dari ruang Banach $B$. Maka
 \index{anihilator}%
\df{anihilator} dari $A$, yang dinotasikan
 \index{<perp@$A^\perp$ (anihilator dari suatu himpunan)}%
dengan~$A^\perp$, adalah $\{f \in B^*\colon f(a) = 0 \text{ untuk setiap $a \in A$}\}$.
\end{defn}

Konflik notasi antara \ref{0001511} dan \ref{0022091} dapat
menimbulkan kebingungan. Untuk melihat bahwa dalam ruang Hilbert, anihilator suatu subruang dan
komplemen ortogonal subruang itu ``pada hakikatnya'' merupakan hal yang sama, gunakan
antiisomorfisme isometrik $\psi$ antara $H$ dan $H^*$ yang dibahas dalam~\ref{0022057} untuk mengidentifikasi keduanya.

\begin{prop} Misalkan $M$ suatu subruang dari ruang Hilbert $H$. Jika kita (untuk sementara) menyatakan
anihilator dari $M$ dengan~$M^a$, maka $M^a = \psi^{\sto}\bigl(M^\perp\bigr)$.
\end{prop}
\section{Proyeksi dan Subruang Terkomplemen}
\begin{defn}
 Pemetaan linear terbatas dari suatu ruang Banach ke dirinya sendiri disebut
 \index{proyeksi!di ruang Banach}%
 \index{operator!proyeksi}%
\df{proyeksi} jika bersifat idempoten. Jika $E$ adalah pemetaan seperti itu, kita menyebutnya
proyeksi \df{pada $\ran E$ sepanjang $\ker E$}.
\end{defn}

Perhatikan bahwa ini sudah pasti \emph{bukan} hal yang sama dengan proyeksi di
ruang Hilbert. Proyeksi ruang Hilbert adalah \emph{proyeksi ortogonal}; yakni,
proyeksi itu swaadjoin sekaligus idempoten. Jangkauan suatu proyeksi ruang Hilbert
tegak lurus terhadap kernelnya. Karena istilah ``ortogonal'', ``swaadjoin'', dan
``tegak lurus'' tidak bermakna di ruang Banach, kita tidak akan mengharapkan kedua pengertian
``proyeksi'' itu identik. Namun, dalam kedua kasus tersebut proyeksi \emph{memang}
terbatas, linear, dan idempoten.

\begin{defn} Misalkan $M$ subruang linear tertutup dari suatu ruang Banach~$B$. Jika terdapat
subruang linear tertutup $N$ dari $B$ sedemikian sehingga $B = M \oplus N$, maka kita mengatakan bahwa $M$ adalah
 \index{subruang terkomplemen}%
 \index{subruang!terkomplemen}%
subruang \df{terkomplemen}, bahwa $N$ adalah
 \index{komplemen!dari subruang Banach}%
\df{komplemen (ruang Banach)}-nya, dan bahwa subruang $M$ dan $N$ bersifat
\df{komplementer}.
\end{defn}

\begin{prop}\label{C069814} Misalkan $M$ dan $N$ subruang komplementer dari suatu ruang Banach~$B$. Jelas bahwa setiap
elemen di $B$ dapat ditulis secara tunggal dalam bentuk $m + n$ untuk suatu $m \in M$ dan $n \in
N$. Definisikan pemetaan $E \colon B \sto B$ dengan $E(m + n) = m$ untuk $m \in M$ dan $n \in N$.
Maka $E$ adalah proyeksi pada $M$ sepanjang~$N$.
\end{prop}

\begin{proof}[\emph{Petunjuk pembuktian}] Gunakan~\ref{prop_cont_idem_op}.  \ns
\end{proof}

\begin{prop}\label{C069817} Jika $E$ adalah proyeksi pada suatu ruang Banach $B$, maka kernel dan jangkauannya merupakan
subruang-subruang komplementer dari~$B$.
\end{prop}

Berikut ini merupakan akibat langsung dari dua proposisi sebelumnya.

\begin{cor}\label{cor_compl_ran_proj} Suatu subruang dari ruang Banach terkomplemen jika dan hanya jika subruang itu
merupakan jangkauan suatu proyeksi.
\end{cor}

\begin{prop} Misalkan $B$ ruang Banach. Jika $E \colon B \sto B$ adalah proyeksi ke suatu subruang
$M$ sepanjang subruang $N$, maka $I - E$ adalah proyeksi ke $N$ sepanjang~$M$.
\end{prop}

\begin{prop}\label{C069824} Jika $M$ dan $N$ merupakan subruang komplementer dari suatu ruang Banach $B$, maka $B/M$ dan
$N$ isomorfik.
\end{prop}

\begin{prop}\label{prop_linv_Bsp} Misalkan $T \colon B \sto C$ pemetaan linear terbatas di antara ruang-ruang Banach.
 \index{kiri!invers!operator ruang Banach}%
Jika $T$ injektif dan $\ran T$ terkomplemen, maka $T$ mempunyai invers kiri.
\end{prop}

\begin{proof}[\emph{Petunjuk pembuktian}] Untuk petunjuk ini kita memperkenalkan notasi (yang agak rumit).
Jika $f\colon X \sto Y$ adalah fungsi di antara dua himpunan, kita definisikan
$f_\bullet \colon X \sto \ran f\colon x \mapsto f(x)$. Artinya, $f_\bullet$ tepat merupakan fungsi yang sama
dengan $f$, kecuali bahwa kodomainnya adalah jangkauan~$f$.

Karena $\ran T$ terkomplemen, terdapat proyeksi $E$ dari $C$ ke~$\ran T$
(menurut~\ref{cor_compl_ran_proj}). Tunjukkan bahwa $T_\bullet$ mempunyai invers kontinu.
Misalkan $S = {T_\bullet}^{-1}E_\bullet$.    \ns
\end{proof}

\begin{prop}\label{prop_rinv_Bsp} Misalkan $S \colon C \sto B$ pemetaan linear terbatas di antara ruang-ruang Banach.
 \index{kanan!invers!operator ruang Banach}%
Jika $S$ surjektif dan $\ker S$ terkomplemen, maka $S$ mempunyai invers kanan.
\end{prop}

\begin{proof}[\emph{Petunjuk pembuktian}] Karena $\ker S$ terkomplemen, terdapat $N \preccurlyeq C$ sedemikian sehingga
$C = N \oplus \ker S$. Tunjukkan bahwa pembatasan $S$ ke $N$ mempunyai invers kontinu dan komposisikan
invers ini dengan pemetaan inklusi dari $N$ ke~$C$.   \ns
\end{proof}

Hasil berikut merupakan konvers dari kedua proposisi~\ref{prop_linv_Bsp} dan~\ref{prop_rinv_Bsp}.

\begin{prop}\label{C069831} Misalkan $T \colon B \sto C$ dan $S\colon C \sto B$ pemetaan linear terbatas di antara ruang-ruang
Banach. Jika $ST = \id B$, maka
 \begin{enumerate}[label=\rm{(\alph*)}]
  \item[(a)] $S$ surjektif;
  \item[(b)] $T$ injektif;
  \item[(c)] $TS$ merupakan proyeksi pada $\ran T$ sepanjang~$\ker S$; dan
  \item[(d)] $C = \ran T \oplus \ker S$.
 \end{enumerate}
\end{prop}

\begin{prop} Jika dalam kategori $\cat{BAN_\infty}$ barisan
 \[ \begin{CD} 0 @>>> A @>j>> B @>f>> C @>>> 0
    \end{CD}\]
 dan
 \[ \begin{CD} 0 @>>> C @>g>> B @>p>> D @>>> 0
    \end{CD}\]
eksak dan jika $fg = \id{C}$, maka $pj$ adalah isomorfisme dari $A$ pada~$D$.
\end{prop}

\begin{proof}[\emph{Petunjuk pembuktian}] Gunakan \ref{C069831} untuk melihat bahwa $B = \ker p \oplus \ran j$. Tinjau pemetaan
$\bigl.p\bigr|_{\ran j}\,j_\bullet$ (notasi untuk $j_\bullet$ seperti dalam petunjuk untuk~\ref{prop_linv_Bsp}). \ns
\end{proof}

\begin{prop}\label{prop_merge_exactCD} Jika dalam kategori $\cat{BAN_\infty}$ diagram
 \begin{equation} \begin{CD} 0 @>>> A @>j>> B @>f>> C @>>> 0  \\
      @.   @VSVV   @VVUV   @.    @.    \\
 0 @>>> A' @>>j'> B' @>>f'> C' @>>> 0
    \end{CD}\end{equation}
dan
 \[ \begin{CD} 0 @>>> C @>g>> B @>p>> D @>>> 0   \\
      @.     @.   @VUVV   @VVTV     @.  \\
  0 @>>> C' @>>g'> B' @>>p'> D' @>>> 0
    \end{CD}\]
komutatif, jika baris-barisnya eksak, jika $fg = \id{C}$ dan jika $f'g' = \id{C'}$, maka diagram
 \[ \begin{CD}  A  @>pj>>  D   \\
       @VSVV   @VVTV  \\
     A'  @>>p'j'> D'
    \end{CD}\]
komutatif dan pemetaan $pj$ serta $p'j'$ merupakan isomorfisme.
\end{prop}

Dalam~\ref{exam_dual_functor} dan~\ref{exam_dualfunctor_exact} kita melihat bahwa dalam kategori ruang Banach
dan pemetaan linear kontinu, funktor dualitas $B \mapsto B^*$, $T \mapsto T^*$ merupakan funktor kontravarian eksak.
Dan dalam~\ref{exam_bar_functor} kita melihat bahwa dalam kategori yang sama pasangan pemetaan $B \mapsto \clo B$,
$T \mapsto \clo T$ merupakan funktor kovarian eksak. Dengan mengomposisikan kedua funktor ini diperoleh dua funktor
kontravarian eksak baru, yang akan kita lambangkan dengan $F$ dan $G$:
  \[ \xymatrix{B\ar[r]^F &{\clo B}^{\,*}}\,,\,\xymatrix{T\ar[r]^F &{\clo T}^{\,*}} \qquad \text{dan} \qquad
           \xymatrix{B\ar[r]^G &\clo{B^*}}\,,\,\xymatrix{T\ar[r]^G &\clo{T^*}}  \]
Wajar untuk menanyakan hubungan antara ${\clo B}^{\,*}$ dan $\clo{B^*}$. Proposisi berikut
memberi tahu kita bukan hanya bahwa kedua ruang ini selalu isomorfik, melainkan bahwa keduanya isomorfik secara alami. Artinya,
 \index{alami!ekuivalensi!antara ${\clo B}^{\,*}$ dan $\clo{B^*}$}%
funktor $F$ dan $G$ ekuivalen secara alami.

\begin{prop} Jika $T \in \ofml B(B,C)$, dengan $B$ dan $C$ ruang Banach, maka terdapat isomorfisme
$\sbsb{\phi}B \colon {\clo B}^{\,*} \sto \clo{B^*}$ dan $\sbsb{\phi}C \colon {\clo C}^{\,*} \sto \clo{C^*}$ sedemikian sehingga
   \[ {\clo T}^{\,*} = {\sbsb{\phi}B}^{\!\!\!-1}\,\, \clo{T^*}\,\,\sbsb{\phi}C\,. \]
\end{prop}

\begin{proof}[\emph{Petunjuk pembuktian}] Karena funktor dualitas eksak (lihat~\ref{exam_dualfunctor_exact}), kita dapat
mengambil dual dari diagram komutatif~\eqref{eqn_exactCD_Bbar} untuk memperoleh diagram komutatif kedua dengan baris-baris
eksak. Sekarang tuliskan diagram komutatif ketiga dengan mengganti setiap kemunculan $B$ dalam diagram~\eqref{eqn_exactCD_Bbar}
dengan~$B^*$. Periksa bahwa ${\sbsb{\tau}B}^{\!\!*}\,\sbsb{\tau}{B^*} = \id{B^*}$ sehingga Anda dapat menggunakan proposisi~\ref{prop_merge_exactCD}
pada dua diagram terakhir dan dengan demikian memperoleh diagram komutatif
     \[ \xy
          \Square/{<-}`{->}`{->}`{<-}/[{\clo B}^{\,*}`{\clo C}^{\,*}`\clo{B^*}`\clo{C^*};{\clo T}^{\,*}`\sbsb{\phi}B`\sbsb{\phi}C`\clo{T^*}]
      \endxy \]
dengan $\sbsb{\phi}B := \sbsb{\pi}{B^*}\,{\sbsb{\pi}B}^{\!\!*}$ suatu isomorfisme.   \ns
\end{proof}

\begin{cor} Suatu ruang Banach refleksif jika dan hanya jika dualnya refleksif.
\end{cor}

\begin{prop}\label{prop_fd_subsp_complm} Setiap subruang berdimensi hingga dari suatu ruang Banach terkomplemen.
\end{prop}

\begin{proof}[\emph{Petunjuk pembuktian}] Misalkan $F$ subruang berdimensi hingga dari suatu ruang Banach $B$ dan misalkan
$\{\vc e^1, \dots, \vc e^n\}$ basis Hamel untuk~$F$. Maka untuk setiap $x \in F$ terdapat skalar-skalar tunggal
$\alpha_1(x)$, \dots , $\alpha_n(x)$ sedemikian sehingga $x = \sum_{k=1}^n \alpha_k(x)\vc e^k$. Untuk setiap $1 \le k \le n$
verifikasi bahwa $\alpha_k$ merupakan fungsional linear terbatas, yang dapat diperluas ke seluruh~$B$. Misalkan $N$ adalah
irisan kernel-kernel dari perluasan ini dan tunjukkan bahwa $B = F \oplus N$.  \ns
\end{proof}

\begin{defn} Ingat bahwa dalam~\ref{000082} kita mendefinisikan \emph{kodimensi} suatu subruang dari ruang vektor
sebagai dimensi komplemen ruang vektornya (yang selalu ada menurut proposisi~\ref{0000824}). Kita menggunakan
istilah yang sama untuk ruang Banach:
 \index{kodimensi}%
\df{kodimensi} dari sebarang \emph{subruang vektor} suatu ruang Banach adalah dimensi \emph{komplemen ruang
vektornya}.
\end{defn}

Untuk korolar berikut (dari proposisi~\ref{prop_fd_subsp_complm}), ingatlah
konvensi~\ref{conv_subsp_Bsp} tentang penggunaan kata ``subruang'' dalam konteks ruang Banach.

\begin{cor} Dalam suatu ruang Banach, setiap subruang berkodimensi hingga terkomplemen.
\end{cor}





















Misalkan $M$ subruang tak nol dari suatu ruang Banach~$B$. Tinjau dua sifat berikut dari pasangan $B$ dan~~$M$.

\textbf{Sifat P:} \emph{Terdapat proyeksi bernorma $1$ dari $B$ ke $M$.}

\textbf{Sifat E:} \emph{Untuk setiap ruang Banach $C$, setiap pemetaan linear terbatas dari $M$ ke $C$ dapat diperluas menjadi pemetaan
linear terbatas dari $B$ ke $C$ tanpa memperbesar normanya.}

\begin{prop} Untuk subruang tak nol $M$ dari suatu ruang Banach $B$, sifat \textbf{P} dan \textbf{E} ekuivalen.
\end{prop}

\begin{proof}[\emph{Petunjuk pembuktian}] Untuk menunjukkan bahwa \textbf{E} mengakibatkan \textbf{P}, mulailah dengan memperluas pemetaan identitas
pada $M$ ke seluruh~$B$.   \ns
\end{proof}

















\section{Prinsip Keterbatasan Seragam}
\begin{defn} Misalkan $S$ suatu himpunan dan $V$ ruang linear bernorma. Suatu keluarga $\fml F$ dari
fungsi-fungsi dari $S$ ke $V$ disebut
 \index{titik demi titik!terbatas}%
 \index{terbatas!titik demi titik}%
\df{terbatas titik demi titik} jika untuk setiap $x \in S$ terdapat konstanta $M_x > 0$ sedemikian sehingga
$\norm{f(x)} \le M_x$ untuk setiap $f \in \fml F$.
\end{defn}

\begin{defn} Misalkan $V$ dan $W$ ruang linear bernorma. Suatu keluarga $\fml T$ dari pemetaan linear terbatas
dari $V$ ke $W$ disebut
 \index{seragam!terbatas}%
 \index{terbatas!seragam}%
\df{terbatas seragam} jika terdapat $M > 0$ sedemikian sehingga $\norm T \le M$ untuk setiap $T \in
\fml T$.
\end{defn}

\begin{exer}\label{C070114} Misalkan $V$ dan $W$ ruang linear bernorma. Jika suatu keluarga $\fml T$ dari
pemetaan linear terbatas dari $V$ ke $W$ terbatas seragam, maka keluarga itu terbatas titik demi titik.
\end{exer}

\emph{Prinsip keterbatasan seragam} adalah pernyataan bahwa konvers
dari~\ref{C070114} berlaku apabila $V$ lengkap. Untuk membuktikannya kita akan menggunakan
analog ``lokal'' dari pernyataan ini untuk fungsi-fungsi kontinu pada ruang metrik lengkap. Secara kasar, analog ini mengatakan bahwa jika
suatu keluarga fungsi bernilai real pada ruang metrik lengkap terbatas titik demi titik, maka keluarga itu terbatas
seragam pada suatu subhimpunan terbuka tak kosong dari ruang tersebut.

\begin{prop}\label{C070117} Misalkan $M$ ruang metrik lengkap yang tak kosong dan $\fml F \subseteq \fml C(M,\R)$.
Jika untuk setiap $x \in M$ terdapat konstanta $N_x > 0$ sedemikian sehingga $\abs{f(x)} \le N_x$ untuk
setiap $f \in \fml F$, maka terdapat himpunan terbuka tak kosong $U$ di $M$ dan bilangan $N > 0$ sedemikian
sehingga $\abs{f(u)} \le N$ untuk setiap $f \in \fml F$ dan setiap $u \in U$.
\end{prop}

\begin{proof}[\emph{Petunjuk pembuktian}] Untuk setiap bilangan asli $k$, misalkan $A_k =
\bigcap\{f^{\gets}\bigl(\,[-k,k]\,\bigr)\colon f \in \fml F\}$. Gunakan versi yang sama dari
\emph{teorema kategori Baire} yang kita gunakan dalam membuktikan \emph{teorema pemetaan terbuka} (\ref{C069414})
untuk menyimpulkan bahwa untuk sedikitnya satu $n \in \N$, himpunan $A_n$ mempunyai interior
tak kosong. Misalkan $U = \intr{A_n}$.   \ns
\end{proof}

\begin{thm}[Prinsip keterbatasan seragam]\label{thm_PUB} Misalkan $B$ ruang Banach dan $W$ ruang linear bernorma.
 \index{prinsip!keterbatasan seragam}%
 \index{keterbatasan!prinsip seragam}%
 \index{seragam!keterbatasan!prinsip}%
Jika suatu keluarga $\ofml T$ dari pemetaan linear terbatas dari $B$ ke $W$ terbatas titik demi
titik, maka keluarga itu terbatas seragam.
\end{thm}

\begin{proof}[\emph{Petunjuk pembuktian}] Untuk setiap $T \in \ofml T$, definisikan
  \[ \sbsb fT\colon B \sto \R\colon x \mapsto \norm{Tx}. \]
Gunakan proposisi sebelumnya untuk menunjukkan bahwa terdapat subhimpunan terbuka tak kosong $U$ dari $B$
dan konstanta $M > 0$ sedemikian sehingga $\sbsb fT(x) \le M$ untuk setiap $T \in \ofml T$ dan setiap
$x \in U$. Pilih titik $a \in B$ dan bilangan $r > 0$ sedemikian sehingga $\clo{B_r(a)} \subseteq
U$. Kemudian verifikasi bahwa
  \[ \norm{Ty} \le r^{-1}\bigl(\, \sbsb fT(a + ry) + \sbsb fT(a)\,\bigr) \le 2Mr^{-1} \]
untuk setiap $T \in \ofml T$ dan setiap $y \in B$ sedemikian sehingga $\norm y \le 1$. \ns
\end{proof}

\begin{thm}[Banach-Steinhaus] Misalkan $B$ ruang Banach, $W$ ruang linear bernorma,
 \index{teorema Banach-Steinhaus}%
dan $(T_n)$ barisan pemetaan linear terbatas dari $B$ ke~$W$. Jika limit titik demi titik
$\lim_{n \sto \infty} T_nx$ ada untuk setiap $x$ di~$B$, maka pemetaan $S\colon B \sto
W\colon x \mapsto \lim_{n \sto \infty} T_nx$ adalah pemetaan linear terbatas.
\end{thm}

\begin{proof}[\emph{Petunjuk pembuktian}] Untuk setiap $x \in B$, barisan $\bigl(\norm{T_nx}\bigr)_{n=1}^\infty$
konvergen, dan karena itu terbatas. Gunakan~\ref{thm_PUB}.   \ns
\end{proof}

\begin{defn} Suatu subhimpunan $A$ dari ruang linear bernorma $V$ disebut
 \index{terbatas!secara lemah!himpunan dalam ruang linear bernorma}%
 \index{lemah!keterbatasan!dalam ruang linear bernorma}%
 \index{w@$w$-terbatas}%
\df{terbatas secara lemah} ($w$-terbatas) jika $f^{\sto}(A)$ adalah subhimpunan terbatas dari medan skalar
$V$ untuk setiap $f \in V^*$.
\end{defn}

\begin{prop} Suatu subhimpunan $A$ dari ruang Hilbert $H$ terbatas secara lemah jika dan hanya jika untuk setiap $y \in H$
 \index{lemah!keterbatasan!dalam ruang Hilbert}%
himpunan $\{\langle a,y \rangle \colon a \in A\}$ terbatas.
\end{prop}

Proposisi sebelumnya kurang lebih sudah jelas. Proposisi berikutnya merupakan konsekuensi penting dari
\emph{prinsip keterbatasan seragam}.

\begin{prop}\label{C070131} Suatu subhimpunan dari ruang linear bernorma terbatas secara lemah jika dan hanya jika
subhimpunan itu terbatas.
\end{prop}

\begin{proof}[\emph{Petunjuk pembuktian}] Jika $A$ subhimpunan yang terbatas secara lemah dari ruang linear bernorma $V$
dan $f \in V^*$, apa yang dapat Anda katakan tentang $\sup\{\abs{a^{**}(f)}\colon a \in A\}$?  \ns
\end{proof}

\begin{exer} Sahabat baik Anda, Fred R. Dimm, kembali memerlukan bantuan. Kali ini ia mencemaskan
pernyataan (lihat~\ref{C070131}) bahwa suatu subhimpunan ruang linear bernorma terbatas jika
dan hanya jika subhimpunan itu terbatas secara lemah. Ia meninjau barisan $(a^n)$ dalam ruang Hilbert
$l^2$ yang didefinisikan oleh $a^n = \sqrt{n}\,\vc e^n$ untuk setiap $n \in \N$, dengan $\{\vc
e^n\colon n \in \N\}$ basis ortonormal biasa untuk~$l_2$ (lihat
contoh~\ref{C063149}). Ia melihat bahwa barisan itu jelas tidak terbatas (dengan topologi biasa
pada~$l^2$). Namun, baginya barisan itu tampak terbatas secara lemah. Jika tidak, menurut argumennya, maka
berdasarkan \emph{teorema Riesz-Fr\'echet} \ref{000342} akan terdapat barisan $x$ dalam
$l^2$ sedemikian sehingga himpunan $\{\abs{\langle a^n,x \rangle}\colon n \in\N\}$ tidak terbatas. Baginya ini tampak
tidak mungkin terjadi karena barisan semacam itu harus menurun lebih lambat daripada
barisan $\bigl(n^{-1/2}\bigr)$, yang sudah menurun terlalu lambat untuk menjadi anggota~$l^2$.
Tenangkan Fred dengan menemukan barisan yang sesuai. \emph{Petunjuk:} Gunakan barisan yang
memiliki banyak nol.
\end{exer}

\begin{defn} Suatu barisan $(x_n)$ dalam ruang linear bernorma $V$ disebut
 \index{lemah!barisan Cauchy}%
 \index{Cauchy!barisan!lemah}%
 \index{barisan!Cauchy lemah}%
\df{Cauchy secara lemah} jika barisan $(f(x_n))$ bersifat Cauchy untuk setiap $f$ dalam~$V^*$. Barisan
tersebut
 \index{lemah!konvergensi!barisan dalam ruang linear bernorma}%
 \index{konvergensi!lemah}%
 \index{barisan!konvergensi lemah dari suatu}%
\df{konvergen secara lemah} ke titik $a \in V$ jika $f(x_n) \sto f(a)$ untuk setiap $f \in V^*$
(yakni, jika barisan itu konvergen dalam topologi lemah yang diinduksi oleh anggota-anggota~$V^*$). Dalam hal ini
kita tulis
 \index{<arrowtow@$x_n \to^w a$ (konvergensi sekuensial lemah)}%
$x_n \to^w a$. Ruang $V$ disebut
 \index{lengkap!secara lemah}%
 \index{lemah!kelengkapan}%
\df{lengkap sekuensial secara lemah} jika setiap barisan Cauchy secara lemah di $V$ konvergen secara lemah.
\end{defn}

\begin{prop} Dalam ruang linear bernorma, setiap barisan Cauchy secara lemah terbatas.
\end{prop}

\begin{prop} Setiap ruang Banach refleksif lengkap sekuensial secara lemah.
\end{prop}

\begin{exam} Misalkan $f_n(t) = \left\{
                             \begin{array}{ll}
                               1 - nt, & \hbox{jika $0 \le t \le \frac1n$;} \\
                               0, & \hbox{jika $\frac1n < t \le 1$.}
                             \end{array}
                           \right.$ untuk setiap $n \in \N$. Barisan fungsi $(f_n)$ menunjukkan bahwa ruang Banach
$\fml C([0,1])$ tidak lengkap sekuensial secara lemah.
\end{exam}

\begin{cor} Ruang Banach $\fml C([0,1])$ tidak refleksif.
\end{cor}

\begin{prop} Misalkan $\fml E$ basis ortonormal untuk suatu ruang Hilbert~$H$ dan $(x_n)$ suatu
barisan di~$H$. Maka $x_n \to^w \vc 0$ jika dan hanya jika $(x_n)$ terbatas dan $\langle
x_n,e \rangle \sto 0$ untuk setiap $e \in \fml E$.
\end{prop}

\begin{defn} Misalkan $H$ ruang Hilbert. Suatu jaring $(T_\lambda)$ dalam $\ofml B(H)$
 \index{lemah!topologi operator!konvergensi dalam}%
 \index{operator!topologi!lemah}%
 \index{topologi!operator lemah}%
 \index{lemah!konvergensi!operator ruang Hilbert}%
\df{konvergen secara lemah} (atau \df{konvergen dalam topologi operator lemah}) ke $S \in \ofml B(H)$ jika
  \[ \langle T_\lambda x,y \rangle \sto \langle Sx,y \rangle \]
untuk setiap $x$, $y \in H$. Dalam hal ini kita tulis
$T_\lambda~\to^{\textbf{WOT}}~S$.
 \index{<arrowtowot@$T_\lambda~\to^{\textbf{WOT}}~S$}%

Suatu keluarga $\ofml T \subseteq \ofml B(H)$ disebut
 \index{lemah!keterbatasan!keluarga operator ruang Hilbert}%
 \index{terbatas!secara lemah}
 \index{terbatas!dalam topologi operator lemah}%
 \index{lemah!topologi operator!keterbatasan dalam}%
 \index{WOT@\textbf{WOT} (topologi operator lemah)}%
\df{terbatas secara lemah} (atau \df{terbatas dalam topologi operator lemah}) jika untuk setiap $x$, $y \in H$
terdapat konstanta positif $\alpha_{x,y}$ sedemikian sehingga
  \[\abs{\langle Tx,y \rangle} \le \alpha_{x,y}\]
untuk setiap $T \in \ofml T$.
\end{defn}

Sayangnya definisi sebelumnya dapat menjadi sumber kebingungan. Karena $\ofml B(H)$ adalah
ruang linear bernorma, ruang itu sudah mempunyai topologi ``lemah'' (lihat~\ref{C067434}). Lalu, bagaimana menangani dua
topologi berbeda pada $\ofml B(H)$ yang mempunyai nama sama? Sebagian penulis menyebut topologi yang baru saja didefinisikan sebagai
\emph{topologi operator lemah}, sehingga mencegah timbulnya kebingungan. Penulis lain berpendapat bahwa
topologi lemah dalam pengertian ruang linear bernorma begitu jarang berguna dalam teori operator pada ruang Hilbert
sehingga, kecuali diberi penjelas khusus, frasa \emph{topologi lemah} dalam konteks operator ruang Hilbert
seharusnya dipahami sebagai topologi yang baru saja didefinisikan. (Lihat, misalnya, komentar Halmos dalam dua
paragraf sebelum Soal 108 di~\cite{Halmos:1982}.) Demikian pula, kita sering melihat frasa \emph{suatu
himpunan terbatas yang terdiri atas operator-operator ruang Hilbert.} Kecuali diberi penjelas lain, yang dimaksud adalah bahwa himpunan
tersebut \emph{terbatas seragam}. Bagaimanapun juga, berhati-hatilah terhadap frasa-frasa ini ketika membaca.

Berikut ini konsekuensi penting dari \emph{prinsip keterbatasan seragam}~\ref{thm_PUB}.

\begin{prop} Setiap keluarga operator ruang Hilbert yang terbatas dalam topologi operator
lemah adalah terbatas seragam.
\end{prop}

\begin{proof}[\emph{Petunjuk pembuktian}] Misalkan $H$ ruang Hilbert dan $\ofml T$ keluarga operator pada~$H$
yang terbatas terhadap topologi operator lemah. Gunakan proposisi~\ref{C070131} untuk menunjukkan, bagi
$y \in H$ yang tetap, bahwa $\{T^*y\colon T \in \ofml T\}$ adalah subhimpunan terbatas dari~$H$. Gunakan hal ini (dan~\ref{C070131}
sekali lagi) untuk menunjukkan bahwa $\{Tx\colon \text{$T \in \ofml T$, $x \in H$, dan $\norm x \le 1$}\}$ terbatas di~$H$. \ns
\end{proof}

Seperti dinyatakan di atas, banyak penulis akan menyingkat pernyataan proposisi sebelumnya menjadi:
\emph{Setiap keluarga operator ruang Hilbert yang terbatas secara lemah adalah terbatas.}

\begin{defn}
Suatu jaring $(T_\lambda)$ dalam $\ofml B(H)$
 \index{kuat!topologi operator!konvergensi dalam}%
 \index{operator!topologi!kuat}%
 \index{topologi!operator kuat}%
 \index{kuat!konvergensi!operator ruang Hilbert}%
 \index{SOT@\textbf{SOT} (topologi operator kuat)}%
\df{konvergen secara kuat} (atau \df{konvergen dalam topologi operator kuat}) ke $S \in \ofml B(H)$ jika
  \[ T_\lambda x \sto Sx \]
untuk setiap $x \in H$. Dalam hal ini kita tulis $T_\lambda~\to^{\textbf{SOT}}~S$.
 \index{<arrowtosot@$T_\lambda~\to^{\textbf{SOT}}~S$}%

Konvergensi norma biasa dari operator-operator dalam $\ofml B(H)$ sering disebut
 \index{seragam!topologi operator!konvergensi dalam}%
 \index{operator!topologi!seragam}%
 \index{topologi!operator seragam}%
 \index{seragam!konvergensi!operator ruang Hilbert}%
\df{konvergensi seragam} (atau \df{konvergensi dalam topologi operator seragam}). Artinya, kita tulis
 \index{<arrowtouniform@$T_\lambda~\sto~S$}%
$T_\lambda \sto S$ jika $\norm{S - T_\lambda}\sto 0$.
\end{defn}

\begin{prop} Misalkan $H$ ruang Hilbert, $(T_n)$ barisan dalam $\ofml B(H)$,
dan $B \in \ofml B(H)$. Maka implikasi-implikasi berikut berlaku:
   \[ T_n \sto B \quad \implies \quad T_n~\to^{\textbf{SOT}}~B
                                \quad \implies \quad T_n~\to^{\textbf{WOT}}~B. \]
\end{prop}

\begin{prop} Misalkan $(S_n)$ dan $(T_n)$ barisan operator pada suatu ruang Hilbert~$H$. \\
Jika $S_n \to^{\textbf{WOT}} A$ dan $T_n \to^{\textbf{SOT}} B$, maka $S_nT_n \to^{\textbf{WOT}} AB$.
\end{prop}

\begin{proof}[\emph{Petunjuk pembuktian}] Tuliskan $S_nT_n - AB$ sebagai $S_nT_n - S_nB + S_nB - AB$.  \ns
\end{proof}

\begin{exam} Dalam proposisi sebelumnya, hipotesis $T_n \to^{\textbf{SOT}} B$ tidak dapat
diganti dengan $T_n~\to^{\textbf{WOT}}~B$.
\end{exam}

\begin{proof}[\emph{Petunjuk pembuktian}] Tinjau pangkat-pangkat operator geser unilateral.  \ns
\end{proof}






\endinput

 \chapter{OPERATOR KOMPAK}\label{chap_cpt_ops}

\section{Definisi dan Sifat-Sifat Dasar}

Kita akan melanjutkan pembahasan dari akhir Bab 5 dengan mempelajari kelas-kelas operator.
Konteks kita akan sedikit lebih umum; dalam hal operator kompak, kita akan
meninjau operator pada ruang Banach.  Pertama, mari kita ingat kembali beberapa definisi dan
fakta dasar dari kalkulus lanjut (atau mata kuliah pengantar analisis).

\begin{defn} Suatu subhimpunan $A$ dari ruang metrik $M$ disebut
 \index{total!keterbatasan}%
 \index{terbatas!total}%
\df{terbatas total} jika untuk setiap $\epsilon > 0$ terdapat suatu subhimpunan hingga $F$ dari $A$
sedemikian sehingga untuk setiap $a \in A$ terdapat titik $x \in F$ yang memenuhi $d(x,a) < \epsilon$.
Untuk ruang linear bernorma, definisi ini dapat dirumuskan kembali: Suatu subhimpunan $A$ dari ruang linear
bernorma $V$ disebut \df{terbatas total} jika untuk setiap bola terbuka $B$ di sekitar nol dalam $V$ terdapat
suatu subhimpunan hingga $F$ dari $A$ sedemikian sehingga $A \subseteq F + B$.  Hal ini memiliki parafrasa yang kurang lebih baku:
Suatu ruang bersifat terbatas total jika dapat diawasi oleh sejumlah hingga polisi dengan jarak pandang
sependek apa pun.  (Sebagian penulis menyebut sifat ini \emph{prakompak}.)
\end{defn}

\begin{prop} Suatu ruang metrik bersifat kompak jika dan hanya jika ruang itu lengkap dan terbatas total.
\end{prop}

\begin{prop} Suatu ruang metrik bersifat kompak jika dan hanya jika ruang itu lengkap dan terbatas total.
\end{prop}

\begin{defn} Suatu subhimpunan $A$ dari ruang topologis $X$ disebut
 \index{kompak!relatif}%
 \index{relatif kompak}%
\df{kompak relatif} jika tutupannya kompak.
\end{defn}

\begin{prop} Dalam ruang metrik, setiap himpunan kompak relatif adalah terbatas total.
\end{prop}

Dalam ruang metrik lengkap, kebalikannya berlaku.

\begin{prop} Dalam ruang metrik lengkap, setiap himpunan terbatas total adalah kompak relatif.
\end{prop}

\begin{prop} Dalam ruang metrik, setiap himpunan terbatas total adalah separabel.
\end{prop}

\begin{thm}[Teorema Arzel\`a-Ascoli]\label{AAThm} Misalkan $X$ ruang Hausdorff kompak. Suatu subhimpunan $\fml F$
dari $\fml C(X)$ terbatas total apabila subhimpunan itu terbatas (secara seragam) dan ekuikontinu.
\end{thm}

Berikutnya kita akan membandingkan topologi \emph{lemah} dan \emph{kuat} pada ruang linear bernorma.
Istilah
 \index{lemah!\emph{vs.} topologi kuat pada ruang linear bernorma}%
 \index{topologi!lemah!\emph{vs.} kuat pada ruang linear bernorma}%
 \index{kuat!\emph{vs.} topologi lemah pada ruang linear bernorma}%
 \index{konvensi!pada ruang linear bernorma \emph{topologi kuat} merujuk pada topologi norma}%
\emph{topologi lemah} akan merujuk pada topologi lemah biasa yang ditentukan oleh fungsional linear terbatas
pada ruang tersebut sebagaimana didefinisikan dalam~\ref{C067434}, sedangkan \emph{topologi kuat} akan merujuk pada topologi norma pada ruang tersebut.
Namun, perlu diingat bahwa kata ``kuat'' dalam konteks ruang linear bernorma, jika ditilik secara ketat,
sesungguhnya mubazir.  Kata itu digunakan semata-mata sebagai penegasan.  Misalnya, banyak penulis akan menyatakan proposisi berikut sebagai:
\emph{Setiap barisan yang konvergen secara lemah dalam ruang linear bernorma adalah terbatas.}

\begin{prop} Setiap barisan yang konvergen secara lemah dalam ruang linear bernorma adalah terbatas secara kuat.
\end{prop}

\begin{proof}[\emph{Petunjuk untuk bukti}] Gunakan \emph{prinsip keterbatasan seragam}~\ref{thm_PUB}. \ns
\end{proof}

\begin{defn} Kita mengatakan bahwa suatu pemetaan linear $T\colon V \sto W$ bersifat
 \index{lemah!kekontinuan}%
 \index{kontinu!secara lemah}%
 \index{linear!pemetaan!kekontinuan lemah suatu}%
\df{kontinu secara lemah} jika pemetaan itu membawa jaring yang konvergen secara lemah ke jaring yang konvergen secara lemah.
\end{defn}

\begin{prop} Setiap pemetaan linear yang kontinu (secara kuat) antara ruang-ruang linear bernorma bersifat kontinu secara lemah.
\end{prop}

\begin{defn} Suatu pemetaan linear $T\colon V \sto W$ antara ruang-ruang linear bernorma disebut
 \index{kompak!pemetaan linear}%
 \index{linear!pemetaan!kompak}%
 \index{operator!kompak}%
\df{kompak} jika pemetaan itu membawa himpunan terbatas ke himpunan kompak relatif. Secara ekuivalen, $T$ bersifat kompak jika pemetaan itu
memetakan bola satuan tertutup dalam $V$ ke suatu himpunan kompak relatif dalam~$W$.   Keluarga semua pemetaan linear kompak
dari $V$ ke $W$ dinotasikan
 \index{K@$\ofml K(V,W)$!keluarga pemetaan linear kompak}%
 \index{K@$\ofml K(V)$!keluarga operator kompak}%
dengan~$\ofml K(V,W)$; dan kita tulis $\ofml K(V)$ untuk~$\ofml K(V,V)$.
\end{defn}

\begin{prop} Setiap pemetaan linear kompak bersifat terbatas.
\end{prop}

\begin{prop} Suatu pemetaan linear $T\colon V \sto W$ antara ruang-ruang linear bernorma bersifat kompak jika dan hanya jika pemetaan itu membawa
setiap barisan terbatas dalam $V$ ke suatu barisan dalam $W$ yang memiliki subbarisan konvergen.
\end{prop}

\begin{exam} Setiap operator berperingkat hingga pada ruang Banach bersifat kompak.
\end{exam}

\begin{exam} Setiap fungsional linear terbatas pada ruang Banach $B$ merupakan pemetaan linear kompak dari $B$ ke~$\K$.
\end{exam}

\begin{exam} Misalkan $C$ ruang Banach $\fml C([0,1])$ dan $k \in \fml C([0,1] \times [0,1])$.  Untuk setiap $f \in C$
 \index{integral!operator}%
 \index{operator!integral}%
definisikan fungsi $Kf$ pada $[0,1]$ dengan
   \[ Kf(x) = \int_0^1 k(x,y)f(y)\,dy\,. \]
Maka operator integral $K$ merupakan operator kompak pada~$C$.
\end{exam}

\begin{proof}[\emph{Petunjuk untuk bukti}] Gunakan \emph{teorema Arzel\`a-Ascoli}~\ref{AAThm}.  \ns
\end{proof}

\begin{exam}\label{X_square_int_kernel} Misalkan $H$ ruang Hilbert $\lfs2{[0,1]}$ yang terdiri atas (kelas-kelas ekuivalensi) fungsi yang
terintegralkan kuadrat pada $[0,1]$ terhadap ukuran Lebesgue dan
 \index{integral!operator}%
 \index{operator!integral}%
$k$ suatu fungsi yang terintegralkan kuadrat
pada $[0,1] \times [0,1]$.  Untuk setiap $f \in H$ definisikan fungsi $Kf$ pada $[0,1]$ dengan
   \[ Kf(x) = \int_0^1 k(x,y)f(y)\,dy\,. \]
Maka operator integral $K$ merupakan operator kompak pada~$H$.
\end{exam}

\begin{prop} Suatu operator pada ruang Hilbert bersifat kompak jika dan hanya jika operator itu memetakan bola satuan tertutup dalam ruang tersebut
ke suatu himpunan kompak.
\end{prop}

\begin{cor}\label{cor_id_op_not_cpt} Operator identitas pada ruang Hilbert berdimensi tak hingga tidak kompak.
(Lihat contoh~\ref{X_l2ball_not_cpt}.)
\end{cor}

\begin{exam} Jika $B$ ruang Banach, keluarga $\ofml B(B)$ dari operator pada $B$ merupakan
 \index{Banach!aljabar!$\ofml B(B)$ sebagai suatu}%
 \index{B@$\ofml B(V)$!sebagai aljabar Banach beridentitas}%
aljabar Banach beridentitas dan
 \index{K@$\ofml K(V)$!sebagai ideal tertutup sejati}%
keluarga $\ofml K(B)$ dari operator kompak pada $B$ merupakan ideal tertutup dalam~$\ofml B(B)$.
Jika $B$ berdimensi tak hingga, ideal $\ofml K(B)$ adalah sejati.
\end{exam}

\begin{prop} Misalkan $T\colon B \sto C$ suatu pemetaan linear terbatas antara ruang-ruang Banach. Maka $T$ bersifat kompak
jika dan hanya jika $T^*$ bersifat kompak.
\end{prop}

\begin{defn}\label{0013} Suatu
 \index{C*@aljabar-$C^*$}%
 \index{aljabar!$C^*$-}%
\df{aljabar-$C^\ast$} adalah aljabar Banach $A$ dengan involusi yang memenuhi
   \[ \norm{a^*a} = \norm{a}^2 \]
untuk setiap $a \in A$.  Sifat norma ini biasanya disebut
 \index{C*@syarat-$C^*$}%
\emph{syarat-$C^*$}. Norma aljabar yang memenuhi syarat ini adalah
 \index{C*@norma-$C^*$}%
 \index{norma!$C^*$-}%
\df{norma-$C^*$}.  Suatu
 \index{C*@subaljabar-$C^*$}%
 \index{subaljabar!$C^*$-}%
\df{subaljabar-$C^*$} dari aljabar-$C^*$ $A$ adalah subaljabar-$*\,$ tertutup dari~$A$.

Kita menyatakan kategori aljabar-$C^*$ dan homomorfisme-$*\,$ dengan $\cat{CSA}$.
 \index{kategori!$\cat{CSA}$ sebagai suatu}%
 \index{CSA@$\cat{CSA}$!kategori}%
Sepintas mungkin tampak aneh bahwa dalam kategori ini, yang objek-objeknya memiliki sifat aljabar sekaligus topologis,
morfismenya hanya disyaratkan mempertahankan struktur aljabar. Namun, ternyata
setiap morfisme dalam $\cat{CSA}$ secara otomatis kontinu---bahkan, kontraktif (lihat
proposisi~\futurexref{12.3.16}{00152171})---dan
 \index{morfisme bijektif!dapat dibalik!dalam $\cat{CSA}$}%
setiap morfisme injektif merupakan isometri (lihat proposisi~\futurexref{12.3.17}{00152181}).  Jelas bahwa
(seperti pada definisi~\ref{00109}) setiap morfisme bijektif dalam kategori ini merupakan isomorfisme. Dengan demikian, setiap homomorfisme-$*\,$ bijektif
merupakan isomorfisme-$*\,$ isometrik.
\end{defn}

\begin{exam} Ruang vektor $\C$ dari bilangan kompleks, dengan perkalian bilangan kompleks biasa
 \index{C*@aljabar-$C^*$!$\C$ sebagai suatu}%
 \index{C@$\C$!sebagai aljabar-$C^*$}%
dan konjugasi kompleks $z \mapsto \conj z$ sebagai involusi, merupakan aljabar-$C^*$
komutatif beridentitas.
\end{exam}

\begin{exam} Jika $X$ suatu ruang Hausdorff kompak, aljabar $\fml C(X)$ dari fungsi kontinu
 \index{C*@aljabar-$C^*$!$\fml C(X)$ sebagai suatu}%
 \index{C@$\fml C(X)$!sebagai aljabar-$C^*$}%
bernilai kompleks pada $X$ merupakan aljabar-$C^*$ komutatif beridentitas apabila involusinya
diambil sebagai konjugasi kompleks.
\end{exam}

\begin{exam} Jika $X$ suatu ruang Hausdorff kompak lokal, aljabar Banach
 \index{C*@aljabar-$C^*$!$\fml C_0(X)$ sebagai suatu}%
 \index{C@$\fml C_0(X)$!sebagai aljabar-$C^*$}%
$\fml C_0(X) = \fml C_0(X,\C)$ dari fungsi kontinu bernilai kompleks pada $X$ yang lenyap
di tak hingga merupakan aljabar-$C^*$ komutatif (tidak mesti beridentitas) apabila involusinya
diambil sebagai konjugasi kompleks.
\end{exam}

\begin{exam} Jika $(X,\mu)$ suatu ruang ukur, aljabar $\lfs\infty(X,\mu)$ dari
 \index{C*@aljabar-$C^*$!$L_\infty(S)$ sebagai suatu}%
 \index{L@$\lfs\infty(S)$!sebagai aljabar-$C^*$}%
fungsi terukur bernilai kompleks yang terbatas secara esensial pada $X$ (sekali lagi dengan konjugasi
kompleks sebagai involusi) merupakan aljabar-$C^*$.  (Secara teknis, tentu saja, anggota
$\lfs\infty(X,\mu)$ adalah kelas-kelas ekuivalensi fungsi yang berbeda pada himpunan-himpunan berukuran nol.)
\end{exam}

\begin{exam}\label{001305fa} Aljabar $\ofml B(H)$ dari operator linear terbatas pada ruang Hilbert
 \index{C*@aljabar-$C^*$!$\ofml B(H)$ sebagai suatu}%
 \index{B@$\ofml B(V)$!sebagai aljabar-$C^*$}%
$H$ merupakan aljabar-$C^*$ beridentitas.  Keluarga $\ofml K(H)$ dari operator kompak merupakan ideal-$*\,$ tertutup
dalam aljabar tersebut.
\end{exam}

\begin{exam}\label{001306} Sebagai kasus khusus dari contoh sebelumnya, perhatikan bahwa himpunan
 \index{C*@aljabar-$C^*$!M@$\mathbf M_n$ sebagai suatu}%
 \index{Ma@$\mathbf M_n = \M n\C$!sebagai aljabar-$C^*$}%
$\mathbf M_n$ dari matriks bilangan kompleks berukuran $n \times n$ merupakan aljabar-$C^*$ beridentitas.
\end{exam}

\begin{prop}\label{prop_K_is_clo_FR} Dalam ruang Hilbert $H$, ideal $\ofml K(H)$ dari operator kompak merupakan tutupan
ideal $\ofml{FR}(H)$ dari operator berperingkat hingga.
\end{prop}

Selama kira-kira satu dasawarsa sebelum Perang Dunia Kedua, sejumlah matematikawan Polandia terkemuka berkumpul secara teratur di kedai
kopi di Lwow untuk membahas hasil dan mengajukan masalah.  Mula-mula mereka bertemu di \emph{Caf\'e Roma}, kemudian di
\emph{Scottish Caf\'e}.  Pada 1935 Stefan Banach menyediakan buku catatan berjilid tempat mereka dapat menuliskan masalah-masalah tersebut.
Dari sinilah \emph{Scottish Book}~\cite{Mauldin:1981} yang kini terkenal bermula.  Pada 6 November 1936, Stanislaw Mazur
menuliskan Masalah 153, yang dikenal sebagai \emph{masalah aproksimasi}; pada intinya masalah itu menanyakan apakah proposisi sebelumnya
berlaku untuk ruang Banach.  Dapatkah setiap operator kompak pada ruang Banach diaproksimasi dalam norma oleh operator berperingkat hingga?
Sebagai hadiah bagi pemecah masalah ini, Mazur menawarkan seekor angsa hidup (tergolong salah satu hadiah paling murah hati yang dijanjikan dalam
\emph{Scottish Book}).  Masalah tersebut tetap terbuka selama beberapa dasawarsa, hingga pada 1972 Per Enflo memberikan jawaban negatif,
setelah berhasil menemukan ruang Banach refleksif separabel yang di dalamnya aproksimasi semacam itu gagal.  Dalam sebuah upacara yang diadakan
tahun itu di Warsawa, Enflo menerima angsa hidupnya dari Mazur sendiri.  Hasil tersebut diterbitkan dalam Acta
Mathematica~\cite{Enflo:1973} pada~1973.

\begin{exam} Operator diagonal $\diag(1,\frac12,\frac13,\dots)$ merupakan operator kompak pada ruang Hilbert~$l_2$.
\end{exam}











\section{Isometri Parsial}
\begin{defn}  Suatu elemen $v$ dari aljabar-$C^*$ $A$ disebut
 \index{parsial!isometri}%
\df{isometri parsial} jika $v^*v$ merupakan proyeksi dalam~$A$.  (Karena $v^*v$ selalu
swaadjoin, cukup disyaratkan bahwa $v^*v$ idempoten.)  Elemen $v^*v$ adalah
\df{proyeksi}
 \index{awal!proyeksi}%
 \index{proyeksi!awal}%
 \index{parsial!isometri!proyeksi awal dari suatu}%
\df{awal} (atau
 \index{tumpuan!proyeksi}%
 \index{proyeksi!tumpuan}%
 \index{parsial!isometri!proyeksi tumpuan dari suatu}%
\df{tumpuan}) dari~$v$, sedangkan $vv^*$ adalah \df{proyeksi}
 \index{akhir!proyeksi}%
 \index{proyeksi!akhir}%
 \index{parsial!isometri!proyeksi akhir dari suatu}%
\df{akhir} (atau
 \index{jangkauan!proyeksi}%
 \index{proyeksi!jangkauan}%
 \index{parsial!isometri!proyeksi jangkauan dari suatu}%
\df{jangkauan}) dari~$v$.  (Sebagai konsekuensi langsung dari proposisi berikutnya,
jika $v$ suatu isometri parsial, maka $vv^*$ memang merupakan proyeksi.)
\end{defn}

\begin{prop}\label{pi0111} Jika $v$ suatu isometri parsial dalam aljabar-$C^*$, maka
 \begin{enumerate}[label=\rm{(\alph*)}]
  \item $vv^*v = v$\,; dan
  \item $v^*$ merupakan isometri parsial.
 \end{enumerate}
\end{prop}

\begin{proof}[\emph{Petunjuk untuk bukti}] Misalkan $z = v - vv^*v$ dan tinjau $z^*z$.  \ns
\end{proof}

\begin{prop}\label{pi0114} Misalkan $v$ suatu isometri parsial dalam aljabar-$C^*$~$A$.  Maka proyeksi awalnya
$p$ adalah proyeksi terkecil (terhadap urutan parsial~$\preceq$
pada $\fml P(A)$\,) sedemikian sehingga $vp = v$, dan proyeksi akhirnya $q$ adalah proyeksi
terkecil sedemikian sehingga $qv = v$.
\end{prop}

\begin{prop}\label{pi0117} Jika $V$ suatu isometri parsial pada ruang Hilbert $H$ (yakni, jika $V$ suatu isometri parsial
dalam aljabar-$C^*$~$\ofml B(H)$\,), maka proyeksi awal $V^*V$ adalah
proyeksi dari $H$ pada $(\ker V)^\perp$, dan proyeksi akhir $VV^*$ adalah proyeksi
dari $H$ pada~$\ran V$.
\end{prop}

Berdasarkan hasil sebelumnya, $(\ker V)^\perp$ disebut
 \index{awal!ruang}%
 \index{ruang!awal}%
 \index{ruang!tumpuan}%
 \index{tumpuan!ruang}%
\df{ruang} \df{awal} (atau \df{tumpuan}) dari~$V$, sedangkan $\ran V$ kadang-kadang disebut
 \index{akhir!ruang}%
 \index{ruang!akhir}%
\df{ruang akhir} dari~$V$.

\begin{prop}\label{pi0121} Suatu operator pada ruang Hilbert merupakan isometri parsial jika dan hanya jika operator tersebut merupakan
isometri pada komplemen ortogonal dari kernelnya.
\end{prop}





\begin{thm}[Dekomposisi Polar]\label{0022287} Jika $T$ suatu operator pada ruang
 \index{dekomposisi polar}%
 \index{dekomposisi!polar}%
Hilbert $H$, maka terdapat isometri parsial $V$ pada $H$ sedemikian sehingga
 \begin{enumerate}
    \item[(i)] ruang awal $V$ adalah $(\ker T)^\perp$,
    \item[(ii)] ruang akhir $V$ adalah $\clo{\ran T}$, dan
    \item[(iii)] $T = V\abs T$.
 \end{enumerate}
Dekomposisi $T$ ini unik dalam arti bahwa jika $T = V_0P$ dengan $V_0$ suatu
isometri parsial dan $P$ suatu operator positif sedemikian sehingga $\ker V_0 = \ker P$, maka $P =
\abs T$ dan $V_0 = V$.
\end{thm}

\begin{proof} Lihat \cite{Conway:1990}, VIII.3.11;\allowbreak \cite{Davidson:1996}, teorema I.8.1;
\allowbreak \cite{KadisonR:1983}, teorema 6.1.2;\allowbreak \cite{Murphy:1990}, teorema 2.3.4;
atau~\cite{Pedersen:1995}, teorema 3.2.17. \ns
\end{proof}

\begin{cor} Jika $T$ adalah suatu
  \index{dekomposisi polar!dari operator invertibel}%
 \index{operator!invertibel!dekomposisi polar dari}%
 \index{invertibel!operator!dekomposisi polar dari}%
operator invertibel pada suatu ruang Hilbert, maka isometri parsial dalam dekomposisi polar
(lihat~\ref{0022287}) bersifat uniter.
\end{cor}

\begin{prop} Jika $T$ adalah suatu
 \index{dekomposisi polar!dari operator normal}%
 \index{operator!normal!dekomposisi polar dari}%
 \index{normal!operator!dekomposisi polar dari}%
operator normal pada suatu ruang Hilbert $H$, maka terdapat operator uniter $U$ pada $H$
yang berkomutasi dengan $T$ dan memenuhi $T = U\abs T$.
\end{prop}

\begin{proof} Lihat \cite{Pedersen:1995}, proposisi 3.2.20.\ns \end{proof}

\begin{exer} Misalkan $(S, \fml A,\mu)$ suatu ruang ukur $\sigma$-hingga dan $\phi \in
L_\infty(\mu)$. Tentukan dekomposisi polar dari
 \index{operator perkalian!dekomposisi polar dari}%
 \index{operator!perkalian!dekomposisi polar dari}%
 \index{dekomposisi polar!dari operator perkalian}%
operator perkalian $M_\phi$ (lihat contoh~\ref{C063527}).
\end{exer}









\section{Operator Kelas Jejak}

Mari kita mulai dengan meninjau beberapa sifat standar \emph{jejak} suatu matriks $n \times n$.

\begin{defn}\label{tr0113} Misalkan $a = [a_{ij}] \in \mathbf M_n$. Definisikan $\tr a$, yaitu
 \index{jejak!dari suatu matriks}%
 \index{matriks!jejak dari suatu}%
 \index{tr@$\tr$ (jejak)}%
\df{jejak} dari $a$, dengan
  \[ \tr a = \sum_{i = 1}^n a_{ii}. \]
\end{defn}

\begin{prop}\label{tr0117} Fungsi $\tr\colon \mathbf M_n \sto \C$ bersifat linear.
\end{prop}

\begin{prop}\label{tr0121} Jika $a$, $b \in \mathbf M_n$, maka $\tr (ab) = \tr (ba)$.
\end{prop}

Jejak bersifat invarian terhadap keserupaan. Dua matriks $n \times n$, yaitu $a$ dan $b$, disebut
 \index{serupa!matriks}%
 \index{matriks!serupa}%
\df{serupa} jika terdapat matriks invertibel $s$ sedemikian sehingga $b = sas^{-1}$.

\begin{prop}\label{tr0124} Matriks-matriks dalam $\mathbf M_n$ yang serupa mempunyai jejak yang sama.
\end{prop}

\begin{defn}\label{tr0211} Misalkan $H$ suatu ruang Hilbert separabel, $(e_n)$ suatu basis ortonormal bagi~$H$, dan $T$
suatu operator positif pada~$H$. Definisikan
 \index{jejak}%
 \index{tr@$\tr$ (jejak)}%
\df{jejak} dari~$T$~dengan
    \[ \tr T:=\sum_{k=1}^\infty \langle Te_k,e_k \rangle. \]
(Jejak tidak harus berhingga.)
\end{defn}

\begin{prop} Jika $S$ dan $T$ merupakan operator positif pada suatu ruang Hilbert separabel dan $\alpha \ge 0$, maka
   \[  \tr(S + T) = \tr S + \tr T \]
dan
   \[ \tr(\alpha T) = \alpha \tr T. \]
\end{prop}

\begin{prop} Jika $T$ suatu operator pada ruang Hilbert separabel, maka $\tr(T^*T) = \tr(TT^*)$.
\end{prop}

Untuk menunjukkan bahwa definisi jejak operator positif yang diberikan di atas tidak bergantung pada basis ortonormal yang dipilih bagi
ruang Hilbert tersebut, kita memerlukan suatu hasil (yang diberikan berikut ini) yang bergantung pada fakta yang muncul kemudian dalam catatan ini: setiap operator positif pada ruang Hilbert
memiliki akar kuadrat (positif) (lihat contoh~\futurexref{11.5.7}{X_sqroot_op}).

\begin{prop}\label{tr0221} Pada keluarga operator positif di suatu ruang Hilbert separabel $H$, jejak bersifat invarian terhadap ekuivalensi uniter;
yakni, jika $T$ operator positif dan $U$ uniter, maka $\tr T = \tr(UTU^*)$.
\end{prop}

\begin{proof}[\emph{Petunjuk untuk bukti}] Dalam proposisi sebelumnya, ambil $S = UT^{\frac12}$.    \ns
\end{proof}

\begin{prop}\label{tr0224} Definisi \emph{jejak}~\ref{tr0211} dari suatu operator positif pada ruang Hilbert separabel tidak bergantung
pada pemilihan basis ortonormal bagi ruang tersebut.
\end{prop}

\begin{proof}[\emph{Petunjuk untuk bukti}] Misalkan $(e^k)$ dan $(f^k)$ basis-basis ortonormal bagi ruang tersebut dan $T$ suatu operator positif
pada ruang itu. Ambil $U$ sebagai operator uniter unik yang memenuhi $e^k = Uf^k$ untuk setiap $k \in \N$.  \ns
\end{proof}

\begin{defn} Misalkan $V$ suatu ruang vektor. Suatu subhimpunan $C$ dari $V$ disebut
 \index{kerucut}%
\df{kerucut} jika $\alpha C \subseteq C$ untuk setiap $\alpha \ge 0$. Suatu kerucut $C$ dalam $V$ disebut
 \index{kerucut!proper}%
 \index{proper!kerucut}%
\df{proper} jika $C \cap (-C) = \{\vc 0\}$.
\end{defn}

 \begin{prop} Suatu kerucut $C$ dalam ruang vektor bersifat konveks jika dan hanya jika $C + C \subseteq C$.
\end{prop}

\begin{exam} Pada suatu ruang Hilbert separabel $H$, keluarga semua operator positif dengan jejak hingga membentuk kerucut konveks proper dalam ruang
vektor~$\ofml B(H)$.
\end{exam}

\begin{defn} Dalam ruang Hilbert separabel $H$, subruang linear dari $\ofml B(H)$ yang direntang oleh operator-operator positif dengan jejak hingga
adalah keluarga $\ofml T(H)$ dari operator
 \index{jejak!kelas}%
 \index{operator!kelas jejak}%
 \index{T@$\ofml T(H)$ (operator kelas jejak pada~$H$)}%
\df{kelas jejak} pada~$H$.
\end{defn}

\begin{prop} Misalkan $H$ ruang Hilbert separabel dengan basis ortonormal~$(e_n)$. Jika $T$ operator kelas jejak pada $H$, maka
   \[ \tr T = \sum_{k=1}^\infty \langle Te_k,e_k \rangle \]
dan deret di ruas kanan konvergen mutlak.
\end{prop}
















\section{Operator Hilbert--Schmidt}

\begin{defn} Suatu operator $A$ pada ruang Hilbert separabel $H$ adalah operator
 \index{Hilbert!--Schmidt, operator}%
 \index{operator!Hilbert--Schmidt}%
\df{Hilbert--Schmidt} jika $A^*A$ termasuk kelas jejak; yakni, jika $\sum_{k=1}^\infty \norm{Ae_k}^2 < \infty$,
dengan $(e_k)$ suatu basis ortonormal untuk~$H$. Keluarga semua operator Hilbert--Schmidt pada $H$
 \index{HS@$\ofml{HS}(H)$ (operator Hilbert--Schmidt pada~$H$)}%
dinyatakan dengan~$\ofml{HS}(H)$.
\end{defn}

\begin{prop} Jika $H$ ruang Hilbert separabel, maka keluarga operator Hilbert--Schmidt pada $H$ merupakan ideal
dua sisi dalam~$\ofml B(H)$.
\end{prop}

\begin{exam} Keluarga operator Hilbert--Schmidt pada suatu ruang Hilbert $H$ dapat dijadikan ruang hasil kali dalam.
\end{exam}

\begin{proof}[\emph{Petunjuk untuk bukti}] Polarisasi. Jika $A$, $B \in \ofml{HS}(H)$,
tinjau $\sum_{k=0}^3 i^k \bigl(A + i^kB\bigr)^*(A + i^kB)$. \ns
\end{proof}

\begin{exam} Hasil kali dalam yang didefinisikan dalam contoh sebelumnya menjadikan $\ofml{HS}(H)$ suatu ruang Hilbert.
\end{exam}

\begin{prop} Pada ruang Hilbert separabel, setiap operator Hilbert--Schmidt bersifat kompak.
\end{prop}

\begin{proof}[\emph{Petunjuk untuk bukti}] Jika $(e_k)$ suatu basis ortonormal untuk ruang Hilbert tersebut, tinjau
operator $F_n := AP_n$, dengan $P_n$ proyeksi pada rentang linear $\{e_1, \dots, e_n\}$. \ns
\end{proof}

\begin{exam} Operator integral yang didefinisikan dalam contoh~\ref{X_square_int_kernel} merupakan operator Hilbert--Schmidt.
\end{exam}

\begin{exam}
    Operator Volterra yang didefinisikan dalam contoh~\ref{000319} merupakan operator Hilbert--Schmidt.
\end{exam}











\endinput

 \chapter{BEBERAPA TEORI SPEKTRAL}\label{spectrum}

Dalam bab ini, untuk pertama kalinya kita akan menganggap, kecuali dinyatakan lain, bahwa ruang
vektor dan aljabar berskalar kompleks, bukan real. Sebagian besar alasan untuk hal ini
 \index{konvensi!dalam bab~\ref{spectrum} skalar bersifat kompleks}%
adalah perlunya menggunakan \emph{teorema Liouville}~\ref{C073131} untuk membuktikan
proposisi~\ref{C073134}.

\section{Spektrum}

\begin{defn} Suatu elemen $a$ dari aljabar beridentitas $A$ disebut
 \index{invertibel!kiri!dalam aljabar}%
\df{invertibel kiri} jika terdapat elemen $a_l$ dalam $A$ (disebut \df{invers kiri}
dari~$a$) sedemikian sehingga $a_l a = \vc 1$, dan disebut
 \index{invertibel!kanan!dalam aljabar}%
\df{invertibel kanan} jika terdapat elemen $a_r$ dalam $A$ (disebut \df{invers
kanan} dari~$a$) sedemikian sehingga $aa_r = \vc 1$. Elemen tersebut disebut
 \index{invertibel!dalam aljabar}%
\df{invertibel} jika elemen itu sekaligus invertibel kiri dan invertibel kanan. Himpunan semua
elemen invertibel dari $A$ dinotasikan
 \index{invertibel@$\inv A$ (elemen invertibel dalam aljabar)}%
dengan~$\inv A$.
\end{defn}

Suatu elemen dari aljabar beridentitas memiliki paling banyak satu invers perkalian. Bahkan,
berlaku pernyataan yang lebih kuat berikut.

\begin{prop}\label{000521} Jika suatu elemen dari aljabar beridentitas memiliki invers kiri dan
invers kanan, maka kedua invers tersebut sama (sehingga elemen itu invertibel).
\end{prop}

Jika suatu elemen $a$ dari aljabar beridentitas invertibel, inversnya yang unik dinotasikan
dengan~$a^{-1}$.

\begin{prop} Jika $a$ elemen invertibel dari suatu aljabar beridentitas, maka $a^{-1}$ juga
 \index{invers!dari suatu invers}%
invertibel dan
   \[ \bigl(a^{-1}\bigr)^{-1} = a\,. \]
\end{prop}

\begin{prop} Jika $a$ dan $b$ elemen-elemen invertibel dari suatu aljabar beridentitas, maka hasil kali
 \index{invers!dari suatu hasil kali}%
$ab$ juga invertibel dan
  \[ (ab)^{-1} = b^{-1}a^{-1}\,. \]
\end{prop}

\begin{prop}\label{000524} Jika $a$ dan $b$ elemen-elemen invertibel dari suatu aljabar beridentitas,
 \index{invers!selisih antara}%
maka
  \[ a^{-1} - b^{-1} = a^{-1}(b - a)b^{-1}\,. \]
\end{prop}

\begin{prop} Misalkan $a$ dan $b$ elemen-elemen dari suatu aljabar beridentitas. Jika $ab$ dan
$ba$ keduanya invertibel, maka $a$ dan $b$ juga invertibel.
\end{prop}

\begin{proof}[\emph{Petunjuk untuk bukti}] Gunakan proposisi~\ref{000521}. \ns \end{proof}

\begin{cor} Misalkan $a$ elemen dari suatu aljabar-$*\,$ beridentitas. Maka $a$ invertibel jika dan hanya jika $a^*a$ dan $aa^*$
invertibel; dalam hal ini,
   \[ a^{-1} = \bigl(a^*a\bigr)^{-1}a^* = a^*\big(aa^*\bigr)^{-1}\,. \]
\end{cor}

\begin{prop}\label{000526} Misalkan $a$ dan $b$ elemen-elemen dari suatu aljabar beridentitas. Maka
$\vc 1 - ab$ invertibel jika dan hanya jika $\vc 1 - ba$ invertibel.
\end{prop}

\begin{proof} Jika $\vc 1 - ab$ invertibel, maka $\vc 1 + b(\vc 1 - ab)^{-1}a$ adalah
invers dari $\vc 1 - ba$.
\end{proof}

\begin{exam} Contoh paling sederhana dari aljabar kompleks ialah keluarga $\C$ dari bilangan
kompleks. Aljabar ini beridentitas sekaligus komutatif.
\end{exam}

\begin{exam}\label{000532} Jika $X$ suatu ruang topologis tak kosong, maka keluarga $\fml C(X)$
dari semua fungsi kontinu bernilai kompleks pada $X$ merupakan suatu
 \index{C@$\fml C(X)$!sebagai aljabar beridentitas}%
aljabar di bawah operasi penjumlahan, perkalian skalar, dan perkalian titik demi titik yang biasa.
Aljabar ini beridentitas sekaligus komutatif.
\end{exam}

\begin{exam} Jika $X$ suatu ruang Hausdorff kompak lokal yang tidak kompak, maka (di bawah
operasi titik demi titik) keluarga $\fml C_0(X)$ dari semua fungsi kontinu bernilai kompleks
pada $X$ yang lenyap di tak hingga merupakan aljabar komutatif.
Namun, keluarga ini \emph{bukan} suatu
 \index{C@$\fml C_0(X)$!sebagai aljabar tanpa identitas}%
aljabar beridentitas.
\end{exam}

\begin{exam}\label{000533} Keluarga
 \index{Ma@$\mathbf M_n = \M n\C$!matriks bilangan kompleks berukuran $n \times n$}%
$\mathbf M_n$ dari matriks bilangan kompleks berukuran $n \times n$ merupakan
 \index{Ma@$\mathbf M_n = \M n\C$!sebagai aljabar beridentitas}%
aljabar beridentitas di bawah operasi penjumlahan, perkalian skalar, dan perkalian matriks yang biasa.
Aljabar ini tidak komutatif jika $n > 1$.
\end{exam}

\begin{exam}\label{000533a} Misalkan $A$ suatu aljabar. Jadikan keluarga
 \index{Ma@$\M nA$!matriks elemen aljabar berukuran $n \times n$}%
$\M n{A}$ dari matriks elemen $A$ berukuran $n \times n$ sebagai suatu
 \index{Ma@$\M nA$!sebagai aljabar}%
aljabar dengan menggunakan aturan operasi matriks yang sama dengan aturan untuk~$\mathbf M_n$.
Dengan demikian, $\mathbf M_n$ hanyalah $\M n\C$. Aljabar $\M nA$ beridentitas jika dan hanya jika $A$ beridentitas.
\end{exam}

\begin{exam}\label{000534} Jika $V$ suatu ruang linear bernorma, maka
 \index{B@$\ofml B(V)$!sebagai aljabar beridentitas}%
$\ofml B(V)$ merupakan aljabar beridentitas. Jika $\dim V~>~1$, aljabar tersebut tidak komutatif.
\end{exam}

\begin{defn}\label{000535} Misalkan $a$ elemen dari aljabar beridentitas~$A$. Maka
 \index{spektrum}%
\df{spektrum} dari $a$, yang dinotasikan dengan
 \index{sigma@$\sigma_A(a)$ atau $\sigma(a)$ (spektrum dari~$a$)}%
$\sigma_A(a)$ atau cukup $\sigma(a)$, adalah himpunan semua bilangan kompleks $\lambda$ sedemikian sehingga
$a - \lambda\vc 1$ tidak invertibel.
\end{defn}

\begin{exam} Jika $z$ suatu elemen dari aljabar bilangan kompleks $\C$, maka
 \index{spektrum!dari bilangan kompleks}%
$\sigma(z) = \{z\}$.
\end{exam}

\begin{exam}\label{0005712} Misalkan $X$ suatu ruang Hausdorff kompak. Jika $f$ suatu elemen
 \index{spektrum!dari fungsi kontinu}%
dari aljabar $\fml C(X)$ yang terdiri atas fungsi-fungsi kontinu bernilai kompleks pada $X$, maka
spektrum $f$ adalah jangkauannya.
\end{exam}

\begin{exam} Misalkan $X$ suatu ruang topologis sembarang. Jika $f$ suatu elemen dari
 \index{spektrum!dari fungsi kontinu terbatas}%
aljabar $\fml C_b(X)$ yang terdiri atas fungsi-fungsi kontinu terbatas bernilai kompleks pada $X$, maka
spektrum $f$ adalah tutupan jangkauannya.
\end{exam}

\begin{exam} Misalkan $S$ suatu ruang ukur positif dan $f \in \lfs{\infty}{S}$ suatu
(kelas ekuivalensi) fungsi yang terbatas secara esensial pada~$S$. Maka
 \index{spektrum!dari fungsi terbatas secara esensial}%
spektrum $f$ adalah jangkauan esensialnya.
\end{exam}

\begin{exam} Keluarga $\mathbf M_3$ dari matriks bilangan kompleks berukuran $3 \times 3$ merupakan aljabar
beridentitas di bawah operasi matriks yang biasa. Spektrum matriks $\begin{bmatrix} 5
& -6 & -6 \\ -1 & 4 & 2 \\ 3 & -6 & -4 \end{bmatrix}$ adalah $\{1,2\}$.
\end{exam}

\begin{prop} Misalkan $a$ suatu elemen dari aljabar beridentitas sedemikian sehingga $a^2 = \vc 1$.
 \index{spektrum!dari elemen yang kuadratnya $1$}%
Maka salah satu dari hal berikut berlaku:
 \begin{enumerate}[label=\rm{(\alph*)}]
  \item $a = \vc 1$, dan dalam hal ini $\sigma(a) = \{1\}$; atau
  \item $a = -\vc 1$, dan dalam hal ini $\sigma(a) = \{-1\}$; atau
  \item $\sigma(a) = \{-1,1\}$.
 \end{enumerate}
\end{prop}

\begin{proof}[\emph{Petunjuk untuk bukti}] Pada (c), untuk membuktikan $\sigma(a) \subseteq \{-1,1\}$,
tinjau $\D\frac1{1 - \lambda^2}(a + \lambda \mathbf 1).$  \ns
\end{proof}

\begin{prop} Ingat bahwa suatu elemen $a$ dari aljabar disebut
 \index{idempoten}%
 \index{spektrum!dari elemen idempoten}%
\df{idempoten} jika $a^2 = a$. Misalkan $a$ suatu elemen idempoten dari aljabar beridentitas. Maka
salah satu dari hal berikut berlaku:
 \begin{enumerate}[label=\rm{(\alph*)}]
  \item $a = \vc 1$, dan dalam hal ini $\sigma(a) = \{1\}$; atau
  \item $a = \vc 0$, dan dalam hal ini $\sigma(a) = \{0\}$; atau
  \item $\sigma(a) = \{0,1\}$.
 \end{enumerate}
\end{prop}

\begin{proof}[\emph{Petunjuk untuk bukti}] Pada (c), untuk membuktikan $\sigma(a) \subseteq \{0,1\}$,
tinjau $\D \frac1{\lambda - \lambda^2}\bigl(a + (\lambda - 1)\vc 1\bigr).$ \ns
\end{proof}

\begin{prop} Misalkan $a$ suatu elemen invertibel dari aljabar beridentitas. Maka suatu bilangan kompleks
tak nol $\lambda$ termasuk dalam spektrum $a$ jika dan hanya jika $1/\lambda$ termasuk dalam
spektrum~$a^{-1}$.
\end{prop}

Proposisi berikut merupakan akibat langsung dari proposisi~\ref{000526}.

\begin{prop}\label{000575} Jika $a$ dan $b$ elemen-elemen dari suatu aljabar beridentitas, maka,
kecuali mungkin untuk $0$, spektrum $ab$ dan spektrum $ba$ sama.
\end{prop}

\begin{defn}\label{C035854} Suatu pemetaan $f \colon A \sto B$ antara aljabar-aljabar Banach disebut
 \index{Banach!aljabar!homomorfisme}%
 \index{homomorfisme!aljabar Banach}%
\df{homomorfisme (aljabar Banach)} jika pemetaan tersebut sekaligus merupakan homomorfisme
aljabar dan pemetaan kontinu antara $A$ dan $B$. Kita menotasikan kategori
 \index{BALG@$\cat{BALG}$!kategori}%
 \index{kategori!$\cat{BALG}$ sebagai}%
aljabar Banach dan homomorfisme aljabar kontinu dengan $\cat{BALG}$.
\end{defn}

\begin{prop} Misalkan $A$ dan $B$ aljabar-aljabar Banach. Maka
 \index{langsung!jumlah!aljabar Banach}%
 \index{Banach!aljabar!jumlah langsung}%
\df{jumlah langsung} dari $A$ dan $B$, yang dinotasikan dengan $A \oplus B$, didefinisikan sebagai himpunan $A
\times B$ yang dilengkapi dengan operasi titik demi titik:
   \begin{enumerate}[label=\rm{(\alph*)}]
     \item $(a,b) + (a',b') = (a + a',b + b')$;
     \item $(a,b)\,(a',b') = (aa',bb')$;
     \item $\alpha(a,b) = (\alpha a,\alpha b)$
   \end{enumerate}
untuk $a$,$a' \in A$, $b$,$b' \in B$, dan $\alpha \in \C$. Seperti dalam~\ref{C023414}, definisikan
$\norm{(a,b)} = \norm a + \norm b$. Hal ini menjadikan $A \oplus B$ suatu aljabar Banach.
\end{prop}

\begin{exer} Identifikasi hasil kali dan koproduk (jika ada) dalam kategori $\cat{BALG}$
yang objeknya aljabar Banach
 \index{hasil kali!dalam $\cat{BALG}$}%
 \index{Banach!aljabar!hasil kali}%
 \index{Banach!aljabar!koproduk}%
 \index{koproduk!dalam $\cat{BALG}$}%
 \index{BALG@$\cat{BALG}$!hasil kali dan koproduk dalam}%
dan morfismenya homomorfisme aljabar kontinu, serta dalam kategori yang objeknya aljabar Banach dan
morfismenya homomorfisme aljabar kontraktif. (Di sini,
 \index{kontraktif}%
\emph{kontraktif} berarti $\norm{Tx} \le \norm x$ untuk setiap~$x$.)
\end{exer}

\begin{prop} Dalam suatu aljabar Banach $A$, operasi penjumlahan dan perkalian (yang dipandang sebagai
pemetaan dari $A \oplus A$ ke $A$) bersifat kontinu, dan perkalian skalar (yang dipandang sebagai pemetaan dari
$\K \oplus A$ ke $A$) juga bersifat kontinu.
\end{prop}

\begin{prop}[Deret Neumann]\label{prop_Neumann_series} Misalkan $a$ suatu elemen dari aljabar Banach beridentitas.
 \index{deret Neumann}%
 \index{deret!Neumann}%
Jika $\norm a < 1$, maka
$\vc 1 - a \in \inv A$ dan $(\vc 1 - a)^{-1} = \sum_{k=0}^\infty a^k$.
\end{prop}

\begin{proof}[\emph{Petunjuk untuk bukti}] Tunjukkan bahwa barisan $(\vc 1, a, a^2, \dots)$
dapat dijumlahkan. (Dalam aljabar beridentitas kita mengartikan $a^0$ sebagai~$\vc 1$.)  \ns
\end{proof}

\begin{cor}\label{cor_Neumann_series} Jika $a$ suatu elemen dari aljabar Banach beridentitas dengan $\norm{\vc 1 - a} < 1$, maka $a$ invertibel dan
    \[ \norm{a^{-1}} \le \frac{1}{1 - \norm{\vc 1 - a}}\,.  \]
\end{cor}

\begin{cor}\label{cor2_Neumann_series} Misalkan $a$ suatu elemen invertibel dari aljabar Banach beridentitas $A$ dan $b \in A$,
serta $\norm{a - b} < \norm{a^{-1}}^{-1}$. Maka $b$ invertibel dalam~$A$.
\end{cor}

\begin{proof}[\emph{Petunjuk untuk bukti}] Gunakan korolari sebelumnya untuk menunjukkan bahwa $a^{-1}b$ invertibel.
\ns \end{proof}

\begin{cor} Jika $a$ termasuk dalam suatu aljabar Banach beridentitas dan $\norm a < 1$, maka
   \[ \norm{(\vc 1 - a)^{-1} - \vc 1} \le \frac{\norm a}{1 - \norm a}\,. \]
\end{cor}

\begin{prop}\label{000644fa} Jika $A$ aljabar Banach beridentitas, maka $\open{\inv A}A$.
\end{prop}

\begin{proof}[\emph{Petunjuk untuk bukti}] Misalkan $a \in \inv A$. Tunjukkan, untuk $h$ yang cukup kecil,
bahwa $\vc 1 - a^{-1}h$ invertibel.  \ns
\end{proof}

\begin{prop}\label{000646fa} Misalkan $A$ aljabar Banach beridentitas. Pemetaan $a \mapsto a^{-1}$ dari $\inv A$ ke
dirinya sendiri merupakan homeomorfisme.
\end{prop}

\begin{prop} Misalkan $A$ aljabar Banach beridentitas. Pemetaan $r\colon a \mapsto a^{-1}$ dari
$\inv A$ ke dirinya sendiri bersifat diferensiabel dan pada setiap elemen invertibel $a$, berlaku
$dr_a(h) = -a^{-1}ha^{-1}$ untuk setiap $h \in A$.
\end{prop}

\begin{proof}[\emph{Petunjuk untuk bukti}] Untuk pengantar singkat mengenai diferensiasi pada ruang
berdimensi tak hingga, lihat Bab 13 dari catatan saya~\cite{Erdman:2007} tentang Analisis Real.  \ns
\end{proof}

\begin{prop}\label{000661} Misalkan $a$ suatu elemen dari aljabar Banach beridentitas~$A$. Maka spektrum $a$
kompak dan $\abs\lambda \le \norm a$ untuk setiap $\lambda \in \sigma(a)$.
\end{prop}

\begin{proof}[\emph{Petunjuk untuk bukti}] Gunakan \emph{teorema Heine-Borel}. Untuk membuktikan bahwa
spektrum tertutup, perhatikan bahwa $(\sigma(a))^c = f^\gets(\inv A)$ dengan $f(\lambda) = a -
\lambda\mathbf 1$ untuk setiap bilangan kompleks~$\lambda$. Tunjukkan juga bahwa jika $\abs\lambda >
\norm a$, maka $\mathbf 1 - \lambda^{-1}a$ invertibel.  \ns
\end{proof}

\begin{defn} Misalkan $a$ suatu elemen dari aljabar Banach beridentitas.
 \index{resolven!pemetaan}%
\df{Pemetaan resolven} untuk $a$ didefinisikan oleh
   \[ R_a \colon \C \setminus \sigma(a) \sto A
                 \colon \lambda \mapsto (a  - \lambda\vc 1)^{-1}\,. \]
\end{defn}

\begin{defn} Misalkan $\open U\C$ dan $A$ suatu aljabar Banach beridentitas. Suatu fungsi $f \colon U \sto A$
 \index{analitik}%
bersifat \df{analitik} pada $U$ jika
   \[ f'(a) := \lim_{z \sto a} \frac{f(z) - f(a)}{z - a} \]
ada untuk setiap $a \in U$. Suatu fungsi bernilai kompleks yang analitik pada seluruh $\C$
disebut fungsi
 \index{entire}
\df{entire}.
\end{defn}

\begin{prop} Untuk $a$ suatu elemen dari aljabar Banach beridentitas $A$ dan $\phi$ suatu fungsional linear
terbatas pada~$A$, misalkan $f := \phi \circ R_a \colon \C \setminus \sigma(a) \sto \C$. Maka
 \begin{enumerate}[label=\rm{(\alph*)}]
  \item $f$ analitik pada domainnya; dan
  \item $f(\lambda) \sto 0$ ketika $\abs\lambda \sto \infty$.
 \end{enumerate}
\end{prop}

\begin{proof}[\emph{Petunjuk untuk bukti}] Untuk (i), perhatikan bahwa dalam aljabar beridentitas $a^{-1} - b^{-1}
= a^{-1}(b - a)b^{-1}$.   \ns
\end{proof}

\begin{thm}[Liouville]\label{C073131} Setiap fungsi entire yang terbatas pada $\C$ bersifat konstan.
 \index{teorema Liouville}%
\end{thm}

\begin{proof}[\emph{Komentar tentang bukti}] Pembuktian teorema ini memerlukan sedikit latar belakang
dalam teori fungsi analitik, yakni materi yang tidak dibahas dalam catatan ini. Teorema
tersebut tentu dibahas dalam setiap mata kuliah pengantar tentang variabel kompleks. Sebagai contoh, bukti yang
baik dapat ditemukan dalam \cite{Rudin:1987}, teorema 10.23.   \ns
\end{proof}

\begin{prop}\label{C073134} Spektrum setiap elemen dari aljabar Banach beridentitas tidak kosong.
\end{prop}

\begin{proof}[\emph{Petunjuk untuk bukti}] Gunakan argumen kontradiksi. Gunakan \emph{teorema Liouville}~\ref{C073131}
untuk menunjukkan bahwa $\phi \circ R_a$ konstan bagi setiap fungsional linear terbatas $\phi$ pada~$A$. Kemudian gunakan
(salah satu versi) \emph{teorema Hahn-Banach} untuk membuktikan bahwa $R_a$ konstan. Mengapa konstanta ini harus~$0$?   \ns
\end{proof}

\noindent\emph{Catatan.} Penting untuk diingat bahwa kita hanya bekerja dengan aljabar
\emph{kompleks}. Hasil ini \emph{salah} untuk aljabar Banach real. Contoh
tandingan sederhana adalah matriks $\begin{bmatrix}0 & 1 \\ -1 & 0 \end{bmatrix}$ yang dipandang sebagai
elemen aljabar Banach (real) $M_2$ yang terdiri atas semua matriks $2 \times 2$ dengan entri
bilangan real.

 \vskip 1 pt

\begin{thm}[Gelfand-Mazur]\label{C073141} Jika $A$ aljabar Banach beridentitas yang setiap
 \index{teorema Gelfand-Mazur}%
elemen taknolnya invertibel, maka $A$ isomorfik secara isometrik dengan~$\C$.
\end{thm}

\begin{proof}[\emph{Petunjuk untuk bukti}] Misalkan $B = \{ \lambda\vc 1\colon \lambda \in \C\}$. Gunakan
proposisi sebelumnya~\ref{C073134} untuk menunjukkan bahwa $B = A$.    \ns
\end{proof}

Pernyataan berikut merupakan konsekuensi langsung dari \emph{teorema Gelfand-Mazur}
\ref{C073141} dan proposisi~\ref{C037821}.

\begin{cor}\label{C073147} Jika $I$ ideal maksimal dalam aljabar Banach komutatif beridentitas $A$,
maka $A/I$ isomorfik secara isometrik dengan~$\C$.
\end{cor}

\begin{prop}\label{0006703fa} Misalkan $a$ suatu elemen dari aljabar beridentitas. Maka
   \[ \sigma(a^n) = [\sigma(a)]^n \]
untuk setiap $n \in \N$. (Notasi $[\sigma(a)]^n$ berarti $\{\lambda^n\colon \lambda \in
\sigma(a)\}$.)
\end{prop}

\begin{defn}\label{C073154} Misalkan $a$ suatu elemen dari aljabar beridentitas.
 \index{spektral!radius}%
 \index{radius!spektral}%
\df{radius spektral} dari~$a$, yang dinotasikan dengan $\rho(a)$, didefinisikan sebagai
$\sup\{\abs\lambda\colon \lambda \in \sigma(a)\}$.
\end{defn}

\begin{prop}\label{0006703} Jika $a$ suatu elemen dari aljabar Banach beridentitas, maka $\rho(a) \le
\norm a$ dan $\rho(a^n) = \bigl(\rho(a)\bigr)^n$.
\end{prop}

\begin{prop} Misalkan $\phi\colon A \sto B$ homomorfisme yang mempertahankan identitas di antara aljabar-aljabar beridentitas. Jika $a$ elemen invertibel dari $A$,
maka $\phi(a)$ invertibel dalam $B$ dan inversnya adalah~$\phi(a^{-1})$.
\end{prop}

\begin{cor} Jika $\phi\colon A \sto B$ homomorfisme yang mempertahankan identitas di antara aljabar-aljabar beridentitas, maka
   \[ \sigma(\phi(a)) \subseteq \sigma(a) \]
untuk setiap $a \in A$.
\end{cor}

\begin{cor} Jika $\phi\colon A \sto B$ homomorfisme yang mempertahankan identitas di antara aljabar-aljabar beridentitas, maka
   \[ \rho(\phi(a)) \le \rho(a) \]
untuk setiap $a \in A$.
\end{cor}

\begin{thm}[Rumus radius spektral]\label{C073157} Jika $a$ suatu elemen dari aljabar Banach
 \index{spektral!radius!rumus untuk}%
beridentitas, maka
   \[ \rho(a) = \inf\bigl\{\norm{a^n}^{1/n}\colon n \in \N\bigr\}
                                = \lim_{n \sto \infty} \norm{a^n}^{1/n}\,. \]
\end{thm}

\begin{exer} Dengan menggunakan rumus untuk operator Volterra $V$ dalam contoh~\ref{000319} sebagai
 \index{operator Volterra!dan persamaan integral}%
 \index{operator!Volterra!dan persamaan integral}%
acuan, pada ruang Banach~$\fml C([0,1])$ definisikan operator integral melalui rumus $Vf(x)=\int_0^x f(t)\,dt$.
 \begin{enumerate}
  \item[(a)] Hitung spektrum operator $V$. (\emph{Petunjuk.} Tunjukkan bahwa $\norm{V^n}
\le (n!)^{-1}$ untuk setiap $n \in \N$ dan gunakan \emph{rumus radius spektral}.)
  \item[(b)] Apa implikasi (a) terhadap kemungkinan memperoleh fungsi kontinu $f$ yang
memenuhi persamaan integral
   \[ f(x) - \mu\int_0^xf(t)\,dt = h(x), \]
dengan $\mu$ suatu skalar dan $h$ suatu fungsi kontinu yang diberikan pada~$[0,1]$?
  \item[(c)] Gunakan gagasan dalam (b) dan uraian deret Neumann untuk $\vc 1 - \mu V$ (lihat
proposisi~\ref{prop_Neumann_series}) untuk menghitung secara eksplisit (dan dalam bentuk fungsi elementer)
suatu solusi untuk persamaan integral
 \[f(x) - {\textstyle \frac12}\int_0^xf(t)\,dt = e^x.\]
 \end{enumerate}
\end{exer}

\begin{prop} Misalkan $A$ aljabar Banach beridentitas dan $B$ subaljabar tertutup yang memuat~$\vc 1_A$.
Maka
 \begin{enumerate}[label=\rm{(\alph*)}]
   \item $\inv B$ terbuka sekaligus tertutup dalam $B \cap \inv A$;
   \item $\sigma_A(b) \subseteq \sigma_B(b)$ untuk setiap $b \in B$; dan
   \item jika $b \in B$ dan $\sigma_A(b)$ tidak memiliki lubang (yakni, jika komplemennya dalam $\C$
terhubung), maka $\sigma_A(b) =    \sigma_B(b)$.
 \end{enumerate}
\end{prop}

\begin{proof}[\emph{Petunjuk untuk bukti}] Tinjau fungsi
$f_b\colon (\sigma_A(b))^c \sto B \cap \inv A\colon \lambda \mapsto b - \lambda \vc 1$.   \ns
\end{proof}















\section{Spektrum Operator Ruang Hilbert}
\begin{prop} Misalkan $H$ ruang Hilbert dan $T \in \ofml B(H)$ suatu operator swaadjoin. Pernyataan berikut
ekuivalen:
 \begin{enumerate}
    \item[(i)] $\sigma(T) \subseteq [0,\infty)$;
    \item[(ii)] terdapat operator $S$ pada $H$ sedemikian sehingga $T = S^*S$; dan
    \item[(iii)] $T$ positif.
 \end{enumerate}
\end{prop}

\begin{notn} Jika $A$ adalah himpunan bilangan kompleks, kita
 \index{<unopurstar@$A^*$ (himpunan konjugat kompleks dari $A \subseteq \C$)}%
mendefinisikan $A^* := \{\conj\lambda \colon \lambda \in A\}$. Hal ini dilakukan untuk menghindari kerancuan dengan tutupan himpunan~$A$.
\end{notn}

\begin{prop} Jika $T$ suatu operator ruang Banach, maka $\sigma(T^*) = \sigma(T)$, sedangkan jika operator tersebut merupakan operator ruang Hilbert,
maka $\sigma(T^*) = (\sigma(T))^*$.
\end{prop}

Dalam ruang berdimensi hingga, hanya ada satu cara bagi suatu bilangan kompleks $\lambda$ untuk menjadi anggota spektrum suatu operator~$T$.
Alasannya adalah bahwa hanya ada satu cara suatu operator, dalam hal ini $T - \lambda I$, pada ruang berdimensi hingga dapat gagal
menjadi invertibel (lihat paragraf sebelum definisi~\ref{defn_similar_vs_ops}). Di sisi lain, ada beberapa cara bagi $\lambda$
untuk menjadi anggota spektrum suatu operator pada ruang berdimensi tak hingga---karena ada beberapa cara suatu operator dapat gagal
menjadi invertibel. (Lihat, misalnya, proposisi~\ref{prop_nasc_op_invert}.) Secara historis, hal ini telah melahirkan sejumlah usulan
taksonomi spektrum suatu operator.

\begin{defn} Misalkan $T$ suatu operator pada ruang Hilbert (atau Banach)~$H$. Suatu bilangan kompleks $\lambda$ termasuk dalam
 \index{resolven!himpunan}%
\df{himpunan resolven} dari $T$ jika $T - \lambda I$ invertibel. Himpunan resolven adalah komplemen dari \emph{spektrum}
$T$. Himpunan bilangan kompleks $\lambda$ yang membuat $T - \lambda I$ tidak injektif disebut
 \index{spektrum titik}%
 \index{spektrum!titik}%
\df{spektrum titik} dari $T$ dan dinotasikan
 \index{sigma@$\sigma_p(T)$ (spektrum titik dari~$T$)}%
dengan $\sigma_p(T)$. Anggota spektrum titik dari $T$ adalah
 \index{nilai eigen}%
\df{nilai eigen} dari $T$. Himpunan bilangan kompleks $\lambda$ yang membuat $T - \lambda I$ tidak terbatas jauh dari nol disebut
 \index{spektrum titik aproksimatif}%
 \index{spektrum!titik aproksimatif}%
\df{spektrum titik aproksimatif} dari $T$ dan dinotasikan
 \index{sigma@$\sigma_{ap}(T)$ (spektrum titik aproksimatif dari~$T$)}%
dengan $\sigma_{ap}(T)$. Himpunan semua bilangan kompleks $\lambda$ sedemikian sehingga tutupan jangkauan $T - \lambda I$ merupakan subruang
sejati dari $H$ disebut
 \index{spektrum kompresi}%
 \index{spektrum!kompresi}%
\df{spektrum kompresi} dari $T$ dan dinotasikan
 \index{sigma@$\sigma_c(T)$ (spektrum kompresi dari~$T$)}%
dengan $\sigma_c(T)$. Adapun
 \index{spektrum residual}%
 \index{spektrum!residual}%
\df{spektrum residual} dari $T$, yang kita notasikan
 \index{sigma@$\sigma_r(T)$ (spektrum residual dari~$T$)}%
dengan $\sigma_r(T)$, adalah $\sigma_c(T) \setminus \sigma_p(T)$.
\end{defn}

\begin{prop} Jika $T$ suatu operator ruang Hilbert, maka $\sigma_p(T) \subseteq \sigma_{ap}(T) \subseteq \sigma(T)$.
\end{prop}

\begin{prop} Jika $T$ suatu operator ruang Hilbert, maka $\sigma(T) = \sigma_{ap}(T) \cup \sigma_c(T)$.
\end{prop}

\begin{prop} Jika $T$ suatu operator ruang Hilbert, maka $\sigma_p(T^*) = (\sigma_c(T))^*$.
\end{prop}

\begin{cor} Jika $T$ suatu operator ruang Hilbert, maka $\sigma(T) = \sigma_{ap}(T) \cup (\sigma_{ap}(T^*))^*$.
\end{cor}

\begin{thm}[Teorema Pemetaan Spektral]\label{SMThm} Jika $T$ suatu operator ruang Hilbert dan $p$ suatu polinomial, maka
   \[ \sigma(p(T)) = p(\sigma(T)). \]
\end{thm}

Teorema sebelumnya berlaku secara lebih umum untuk setiap fungsi analitik yang didefinisikan pada lingkungan spektrum~$T$.
(Untuk pembuktiannya, lihat~\cite{Conway:1990}, bab VII, teorema 4.10.) Hasil tersebut juga berlaku untuk bagian-bagian spektrum.

\begin{prop} Jika $T$ suatu operator ruang Hilbert dan $p$ suatu polinomial, maka
   \begin{enumerate}[label=\rm{(\alph*)}]
     \item $\sigma_p(p(T)) = p(\sigma_p(T))$,
     \item $\sigma_{ap}(p(T)) = p(\sigma_{ap}(T))$, dan
     \item $\sigma_c(p(T)) = p(\sigma_c(T))$,
   \end{enumerate}
\end{prop}

\begin{prop}
    Jika $T$ operator invertibel pada ruang Hilbert, maka $\sigma(T^{-1}) = (\sigma(T))^{-1}$.
\end{prop}

\begin{defn}\label{defn_similar_Hsp_ops} Seperti halnya pada ruang vektor, dua operator $R$ dan $T$ pada ruang Hilbert (atau Banach) $H$ dikatakan
 \index{serupa!operator}%
\df{serupa} jika terdapat operator invertibel $S$ pada $H$ sedemikian sehingga $R = STS^{-1}$.
\end{defn}

\begin{prop} Jika $S$ dan $T$ merupakan operator yang serupa pada ruang Hilbert, maka $\sigma(S) = \sigma(T)$, $\sigma_p(S) = \sigma_p(T)$,
$\sigma_{ap}(S) = \sigma_{ap}(T)$, dan $\sigma_c(S) = \sigma_c(T)$.
\end{prop}

\begin{prop} Jika $T$ merupakan operator normal pada ruang Hilbert, maka $\sigma_c(T) = \sigma_p(T)$, sehingga $\sigma_r(T)~=~\emptyset$.
\end{prop}

\begin{prop} Jika $T$ merupakan operator pada ruang Hilbert, maka $\sigma_{ap}(T)$ tertutup.
\end{prop}

\begin{exam} Misalkan $D = \diag(a_1,a_2,a_3,\dots)$ merupakan operator diagonal pada ruang Hilbert dan misalkan
 \index{spektrum!bagian-bagian!operator diagonal}%
 \index{operator!diagonal!bagian-bagian spektrum}%
 \index{diagonal!operator!bagian-bagian spektrum}%
$A = \{a_k\colon k\in\N\}$. Maka
   \begin{enumerate}[label=(\alph*)]
     \item $\sigma_p(D) = \sigma_c(D) = A$,
     \item $\sigma_{ap}(D) = \clo A$, dan
     \item $\sigma_r(D) = \emptyset$.
   \end{enumerate}
\end{exam}

\begin{exam} Misalkan $S$ merupakan operator pergeseran unilateral pada $l_2$
 \index{spektrum!bagian-bagian!pergeseran unilateral}%
 \index{operator!pergeseran unilateral!bagian-bagian spektrum}%
 \index{pergeseran unilateral!bagian-bagian spektrum}%
(lihat contoh~\ref{C063526}). Maka
    \begin{enumerate}[label=(\alph*)]
      \item $\sigma(S) = D$,
      \item $\sigma_p(S) = \emptyset$,
      \item $\sigma_{ap}(S) = C$,
      \item $\sigma_c(S) = B$,
      \item $\sigma(S^*) = D$,
      \item $\sigma_p(S^*) = B$,
      \item $\sigma_{ap}(S^*) = D$, dan
      \item $\sigma_c(S^*) = \emptyset$.
    \end{enumerate}
di mana $B$ adalah cakram satuan terbuka, $D$ adalah cakram satuan tertutup, dan $C$ adalah lingkaran satuan pada bidang kompleks.
\end{exam}

\begin{exam} Misalkan $(S,\sfml A, \mu)$ merupakan ruang ukuran positif $\sigma$-hingga, $\phi \in L_\infty(S,\mu)$, dan $M_\phi$ merupakan
 \index{spektrum!bagian-bagian!operator perkalian}%
 \index{operator!perkalian!bagian-bagian spektrum}%
 \index{operator perkalian!bagian-bagian spektrum}%
operator perkalian yang bersesuaian pada ruang Hilbert $\lfs2{S,\mu}$ (lihat contoh~\ref{C063527}). Maka
spektrum dan spektrum titik aproksimatif dari $M_\phi$ adalah jangkauan esensial~$\phi$, sedangkan spektrum titik
dan spektrum kompresi dari $M_\phi$ adalah $\{\lambda \in \C\colon \mu(\phi^{\gets}(\{\lambda\})) > 0\}$.
\end{exam}






















\endinput

 \chapter{RUANG VEKTOR TOPOLOGIS}

\section{Himpunan Seimbang dan Himpunan Penyerap}

\begin{defn} Suatu himpunan $S$ di dalam ruang vektor disebut
 \index{seimbang}%
\df{seimbang} (atau \df{melingkar}) jika $\Dsk S \subseteq S$.
\end{defn}

\begin{notn} Misalkan $\Dsk$ menyatakan cakram satuan tertutup
 \index{D@$\Dsk$ (cakram satuan tertutup di~$\C$)}%
pada bidang kompleks $\{z \in \C\colon \abs z \le 1\}$.
\end{notn}

\begin{prop} Misalkan $B$ suatu subhimpunan seimbang dari sebuah ruang vektor. Jika $\alpha$ dan $\beta$ adalah skalar sedemikian sehingga $\abs\alpha \le \abs\beta$,
maka $\alpha B \subseteq \beta B$.
\end{prop}

\begin{cor} Jika $B$ suatu subhimpunan seimbang dari sebuah ruang vektor dan $\abs\lambda =1$, maka $\lambda B = B$.
\end{cor}

\begin{defn} Jika $S$ suatu subhimpunan dari sebuah ruang vektor, maka himpunan $\Dsk S$ adalah
 \index{seimbang!selubung}%
 \index{selubung!seimbang}%
\df{selubung seimbang} dari~$S$.
\end{defn}

\begin{prop} Misalkan $S$ suatu subhimpunan dari ruang vektor~$V$.  Maka selubung seimbang dari $S$ adalah irisan semua
subhimpunan seimbang dari $V$ yang memuat~$S$; jadi, himpunan itu adalah subhimpunan seimbang terkecil dari $V$ yang memuat~$S$.
\end{prop}

\begin{defn} Suatu subhimpunan $A$ dari ruang vektor $V$
 \index{menyerap}%
\df{menyerap} suatu subhimpunan $B \subseteq V$ jika terdapat $M > 0$ sedemikian sehingga $B \subseteq tA$ setiap kali $\abs t \ge M$.
Himpunan $A$ bersifat
 \index{menyerap}%
\df{menyerap} (atau \df{radial}) jika himpunan itu menyerap setiap himpunan beranggota satu (secara ekuivalen, setiap himpunan hingga).
\end{defn}

Jadi, secara longgar, suatu himpunan bersifat menyerap jika setiap vektor dalam ruang tersebut termuat dalam suatu hasil dilatasi yang sesuai dari himpunan itu.

\begin{exam} Misalkan $B$ adalah gabungan sumbu-$x$ dan sumbu-$y$ dalam ruang vektor riil $\R^2$.  Maka $B$ seimbang, tetapi tidak konveks.
\end{exam}

\begin{exam} Dalam ruang vektor $\C$, sebuah elips dengan salah satu fokusnya di titik asal merupakan contoh himpunan penyerap yang tidak seimbang.
\end{exam}

\begin{exam} Dalam ruang vektor $\C^2$, himpunan $\Dsk \times \{0\}$ seimbang, tetapi tidak menyerap.
\end{exam}

\begin{prop} Suatu himpunan seimbang $B$ dalam ruang vektor $V$ bersifat menyerap jika dan hanya jika untuk setiap $x \in V$ terdapat skalar
$\alpha \ne 0$ sedemikian sehingga $\alpha x \in B$.
\end{prop}











\section{Filter}

Dalam mempelajari ruang vektor topologis, penggunaan bahasa \emph{filter} memudahkan kita mencirikan sifat-sifat topologis
ruang-ruang tersebut.  Bahasa ini merupakan sarana alternatif bagi jaring, tetapi sebagian besar ekuivalen dengannya, dan terutama berguna dalam situasi (nonmetrik) ketika barisan
tidak lagi banyak membantu.

Secara kasar, sebuah \emph{filter} adalah keluarga tak kosong dari himpunan-himpunan tak kosong yang tertutup terhadap irisan berhingga dan pengambilan superhimpunan.

\begin{defn} Suatu keluarga tak kosong $\sfml F$ dari subhimpunan-subhimpunan tak kosong dari suatu himpunan $S$ disebut
 \index{filter!pada suatu himpunan}%
\df{filter} pada $S$ jika
   \begin{enumerate}[label=(\alph*)]
     \item $A$, $B \in \sfml F$ mengakibatkan $A \cap B \in \sfml F$, dan
     \item jika $A \in \sfml F$ dan $A \subseteq B$, maka $B \in \sfml F$.
   \end{enumerate}
\end{defn}

\begin{exam} Salah satu contoh filter yang sederhana (dan, bagi keperluan kita, terpenting) adalah keluarga semua lingkungan suatu titik $a$ di sebuah
ruang topologis.  Keluarga ini adalah
 \index{lingkungan!filter pada suatu titik}%
 \index{filter!lingkungan}%
\df{filter lingkungan} pada~$a$. Filter ini akan
 \index{N@$\sfml N_a$ (filter lingkungan pada~$a$)}%
dilambangkan dengan $\sfml N_a$.
\end{exam}

\begin{defn} Suatu keluarga tak kosong $\sfml B$ yang terdiri atas subhimpunan-subhimpunan tak kosong dari himpunan $S$ disebut
 \index{basis filter}%
\df{basis filter} pada $S$ jika untuk setiap $A$, $B \in \sfml B$ terdapat $C \in \sfml B$ sedemikian sehingga $C \subseteq A \cap B$.

Kita mengatakan bahwa suatu subkeluarga $\sfml B$ dari filter $\sfml F$ adalah
 \index{basis filter!bagi suatu filter}%
 \index{basis!bagi suatu filter}%
\df{basis filter bagi} $\sfml F$ jika setiap anggota $\sfml F$ memuat suatu anggota~$\sfml B$.  (Penting untuk memeriksa bahwa
setiap subkeluarga seperti itu memang merupakan basis filter.)
\end{defn}

\begin{exam} Setiap filter adalah basis filter.
\end{exam}

\begin{exam} Misalkan $a$ sebuah titik dalam ruang topologis.  Setiap basis lingkungan pada $a$ adalah basis filter bagi filter $\sfml N_a$
yang terdiri atas semua lingkungan~$a$. Secara khusus, keluarga semua lingkungan terbuka dari $a$ adalah basis filter bagi~$\sfml N_a$.
\end{exam}

\begin{exam} Misalkan $\sfml B$ suatu basis filter pada himpunan $S$. Maka keluarga $\sfml F$ yang terdiri atas semua himpunan yang memuat anggota $\sfml B$
adalah filter pada $S$, dan $\sfml B$ adalah basis filter bagi~$\sfml F$.  Filter ini akan kita sebut filter yang
 \index{filter!yang dibangkitkan oleh suatu basis filter}%
\df{dibangkitkan oleh}~$\sfml B$.
\end{exam}

\begin{defn} Suatu basis filter (khususnya, suatu filter) $\sfml B$ pada ruang topologis $X$
 \index{konvergensi!suatu filter}%
 \index{filter!konvergensi suatu}%
 \index{konvergensi!suatu basis filter}%
 \index{basis filter!konvergensi suatu}%
 \index{limit!suatu basis filter}%
\df{konvergen} ke sebuah titik $a \in X$ jika setiap anggota $N$ dari $\sfml N_a$, yaitu filter lingkungan pada $a$, memuat suatu anggota~$\sfml B$.
Dalam hal ini, kita
 \index{<arrowto@$\sfml B \sto a$ (konvergensi suatu basis filter)}%
menuliskan $\sfml B \sto a$.
\end{defn}

\begin{prop} Suatu basis filter pada ruang topologis Hausdorff konvergen ke paling banyak satu titik.
\end{prop}

\begin{exam} Misalkan $(x_\lambda)$ sebuah jaring dalam himpunan~$S$ dan $\sfml F$ keluarga semua $A \subseteq S$ sedemikian sehingga jaring $(x_\lambda)$
pada akhirnya berada di~$A$.  Maka $\sfml F$ adalah filter pada~$S$.  Kita akan menyebutnya filter yang
 \index{filter!yang dibangkitkan oleh suatu jaring}%
\df{dibangkitkan oleh} jaring~$(x_\lambda)$.
\end{exam}

\begin{prop} Misalkan $(x_\lambda)$ suatu jaring dalam ruang topologis~$X$, misalkan $\sfml F$ filter yang dibangkitkan oleh~$(x_\lambda)$,
dan $a \in X$.  Maka $\sfml F \sto a$ jika dan hanya jika $x_\lambda \sto a$.
\end{prop}

\begin{exam} Misalkan $\sfml B$ suatu basis filter pada himpunan~$S$.  Misalkan $D$ himpunan semua pasangan terurut $(b,B)$ sedemikian sehingga $b \in B \in \sfml B$.
Untuk $(b_1,B_1)$, $(b_2,B_2) \in D$, definisikan $(b_1,B_1) \le (b_2,B_2)$ jika $B_2 \subseteq B_1$.  Ini membuat $D$ menjadi himpunan terarah.  Sekarang bentuklah
sebuah jaring dalam $S$ dengan mendefinisikan
   \[ x\colon D \sto S\colon (b,B) \mapsto b \]
yakni, tetapkan $x_\lambda = b$ untuk setiap $\lambda = (b,B) \in D$.  Maka $\bigl(x_\lambda\bigr)_{\lambda \in D}$ adalah sebuah jaring dalam~$S$.  Jaring ini adalah jaring yang
 \index{jaring!yang dibangkitkan oleh suatu basis filter}%
 \index{jaring!yang didasarkan pada suatu basis filter}%
\df{dibangkitkan oleh} (atau \df{didasarkan pada}) basis filter~$\sfml B$.
\end{exam}

\begin{prop} Misalkan $\sfml B$ suatu basis filter pada ruang topologis $X$, misalkan $(x_\lambda)$ adalah jaring yang dibangkitkan oleh $\sfml B$, dan misalkan $a \in X$.
Maka $x_\lambda \sto a$ jika dan hanya jika $\sfml B \sto a$.
\end{prop}













\section{Topologi Kompatibel}

\begin{defn} Misalkan $\sfml T$ suatu topologi pada ruang vektor~$X$.  Kita mengatakan bahwa $\sfml T$
 \index{kompatibel}%
\df{kompatibel} dengan struktur linear pada $X$ jika operasi penjumlahan $A\colon X \times X \sto X\colon (x,y) \mapsto x + y$
dan perkalian skalar $M\colon \K \times X \sto X\colon (\alpha,x) \mapsto \alpha x$ kontinu.  (Di sini, $X \times X$ dan
$\K \times X$ diperlengkapi dengan topologi hasil kali.)

Jika $X$ adalah ruang vektor yang dilengkapi dengan topologi $\sfml T$ yang kompatibel dengan struktur linearnya, kita mengatakan bahwa pasangan
$(X,\sfml T)$ adalah suatu
 \index{topologis!ruang vektor}%
 \index{vektor!ruang!topologis}%
 \index{ruang!vektor topologis}%
\df{ruang vektor topologis}.  (Sebagaimana tentu Anda duga, kita hampir selalu akan tetap menggunakan singkatan baku yang tidak logis seperti
``Misalkan $X$ suatu ruang vektor topologis \dots''.)  Kita lambangkan dengan $\cat{TVS}$ kategori ruang-ruang vektor topologis
 \index{kategori!$\cat{TVS}$ sebagai suatu}%
 \index{TVS@$\cat{TVS}$!kategori}%
dan pemetaan linear kontinu.
\end{defn}

\begin{exam} Misalkan $X$ ruang vektor tak nol.  Topologi diskret pada $X$ tidak kompatibel dengan struktur linear pada~$X$.
\end{exam}

\begin{prop} Misalkan $X$ ruang vektor topologis, $a \in X$, dan $\alpha \in \K$.
   \begin{enumerate}[label=\rm{(\alph*)}]
    \item Pemetaan $T_a\colon x \mapsto x + a$
 \index{translasi}%
(\df{translasi} oleh~$a$) dari $X$ ke dirinya sendiri merupakan suatu homeomorfisme.
    \item Jika $\alpha \ne 0$, pemetaan $x \mapsto \alpha x$ dari $X$ ke dirinya sendiri merupakan suatu homeomorfisme.
   \end{enumerate}
\end{prop}

\begin{cor} Suatu himpunan $U$ adalah lingkungan titik $x$ dalam ruang vektor topologis jika dan hanya jika $-x + U$ adalah lingkungan~$\vc 0$.
\end{cor}

\begin{cor} Misalkan $\sfml U$ suatu basis filter bagi filter lingkungan di titik asal dalam ruang vektor topologis~$X$.  Suatu subhimpunan $V$ dari $X$
terbuka jika dan hanya jika untuk setiap $x \in V$ terdapat $U \in \sfml U$ sedemikian sehingga $x + U \subseteq V$.
\end{cor}

Korolari sebelumnya menyatakan bahwa topologi suatu ruang vektor topologis sepenuhnya ditentukan oleh suatu basis filter bagi
filter lingkungan di titik asal ruang tersebut.  Basis filter semacam itu kita sebut
 \index{lokal!basis}%
 \index{basis!lokal}%
\df{basis lokal} (bagi topologi tersebut).

\begin{prop} Jika $U$ adalah lingkungan titik asal dalam ruang vektor topologis dan $0 \ne \lambda \in \K$, maka $\lambda U$ adalah
lingkungan titik asal.
\end{prop}

\begin{prop} Dalam ruang vektor topologis, setiap lingkungan $\vc 0$ memuat suatu lingkungan seimbang dari~$\vc 0$.
\end{prop}

\begin{prop} Dalam ruang vektor topologis, setiap lingkungan $\vc 0$ bersifat menyerap.
\end{prop}

\begin{prop} Jika $V$ adalah lingkungan $\vc 0$ dalam ruang vektor topologis, maka terdapat lingkungan $U$ dari $\vc 0$
sedemikian sehingga $U + U \subseteq V$.
\end{prop}

Proposisi berikutnya merupakan (hampir) kebalikan dari ketiga proposisi sebelumnya.

\begin{prop}\label{prop_fbase_induces_top} Misalkan $X$ ruang vektor dan $\sfml U$ suatu basis filter pada $X$ yang memenuhi syarat-syarat berikut:
   \begin{enumerate}[label=\rm{(\alph*)}]
     \item setiap anggota $\sfml U$ seimbang;
     \item setiap anggota $\sfml U$ bersifat menyerap; dan
     \item untuk setiap $V \in \sfml U$ terdapat $U \in \sfml U$ sedemikian sehingga $U + U \subseteq V$.
   \end{enumerate}
Maka terdapat topologi $\sfml T$ pada $X$ yang menjadikan $X$ suatu ruang vektor topologis dan $\sfml U$ suatu basis lokal bagi~$\sfml T$.
\end{prop}

\begin{exam} Dalam ruang linear bernorma, keluarga $\{C_r(\vc 0)\colon r > 0\}$ yang terdiri atas bola-bola tertutup berpusat di titik asal merupakan basis lokal bagi
topologi pada ruang tersebut.
\end{exam}

\begin{exam} Misalkan $\Omega$ suatu subhimpunan terbuka tak kosong dari $\R^n$.  Untuk setiap subhimpunan kompak $K \subseteq \Omega$ dan setiap $\epsilon > 0$, misalkan
$U_{K,\epsilon} = \{ f \in \fml C(\Omega)\colon \abs{f(x)} \le \epsilon \text{ jika } x \in K\}$.  Maka keluarga semua
himpunan $U_{K,\epsilon}$ tersebut merupakan basis lokal bagi suatu topologi pada~$\fml C(\Omega)$.  Topologi ini adalah
 \index{topologi!konvergensi seragam pada himpunan-himpunan kompak}%
 \index{seragam!konvergensi!pada himpunan-himpunan kompak}%
 \index{konvergensi!seragam!pada himpunan-himpunan kompak}%
 \index{himpunan-himpunan kompak!topologi konvergensi seragam pada}%
\df{topologi konvergensi seragam pada himpunan-himpunan kompak}.
\end{exam}

\begin{prop} Misalkan $A$ suatu subhimpunan dari ruang vektor topologis dan $0 \ne \alpha \in \K$.  Maka $\alpha \clo A = \clo{\alpha A}$.
\end{prop}

\begin{prop} Misalkan $S$ suatu himpunan dalam ruang vektor topologis.
   \begin{enumerate}[label=\rm{(\alph*)}]
    \item Jika $S$ seimbang, demikian pula tutupannya.
    \item Selubung seimbang dari $S$ mungkin tidak tertutup sekalipun $S$ sendiri tertutup.
   \end{enumerate}
\end{prop}

\begin{proof}[\emph{Petunjuk untuk bukti}] Untuk (b), misalkan $S$ adalah hiperbola di $\R^2$ yang persamaannya $xy = 1$.  \ns
\end{proof}

\begin{prop} Jika $M$ merupakan subruang vektor dari suatu ruang vektor topologis, maka~$\clo M$ juga merupakan subruang vektor.
\end{prop}

Jelas bahwa setiap himpunan penyerap dalam ruang vektor topologis harus memuat titik asal.  Namun, contoh berikut menunjukkan bahwa
himpunan tersebut belum tentu memuat suatu lingkungan titik asal.

\begin{exam}[Apel Schatz] Dalam ruang vektor real $\R^2$, himpunan
    \[ C_1\bigl((-1,0)\bigr) \cup C_1\bigl((1,0)\bigr) \cup \bigl(\{0\} \times [-1,1]\bigr) \]
bersifat menyerap, tetapi bukan merupakan lingkungan titik asal.
\end{exam}

Hasil berikut menyatakan bahwa dalam ruang vektor topologis, himpunan kompak dan himpunan tertutup dapat dipisahkan oleh himpunan terbuka jika keduanya saling lepas.

\begin{prop} Jika $K$ dan $C$ adalah dua subhimpunan tak kosong yang saling lepas dalam ruang vektor topologis $X$, dengan $K$ kompak dan $C$ tertutup, maka terdapat
lingkungan $V$ dari titik asal sedemikian sehingga $(K + V) \cap (C + V)~=~\emptyset$.
\end{prop}

\begin{cor} Jika $K$ adalah subhimpunan kompak dari ruang vektor topologis Hausdorff $X$ dan $C$ adalah subhimpunan tertutup dari~$X$, maka $K + C$
tertutup di~$X$.
\end{cor}

\begin{proof}[\emph{Petunjuk untuk bukti}] Misalkan $x \in (K + C)^c$.  Tinjau himpunan $x - K$ dan~$C$.  \ns
\end{proof}

\begin{cor} Setiap anggota suatu basis lokal $\sfml B$ bagi ruang vektor topologis memuat tutupan suatu anggota~$\sfml B$.
\end{cor}

\begin{prop} Misalkan $X$ ruang vektor topologis.  Maka pernyataan-pernyataan berikut ekuivalen.
  \begin{enumerate}[label=\rm{(\alph*)}]
    \item $\{\vc 0\}$ adalah himpunan tertutup di~$X$.
    \item $\{x\}$ adalah himpunan tertutup untuk setiap $x \in X$.
    \item Untuk setiap $x \ne 0$ di $X$, terdapat lingkungan titik asal yang tidak memuat $x$.
    \item $\bigcap \sfml U = \{0\}$, dengan $\sfml U$ suatu basis lokal bagi~$X$.
    \item $X$ adalah Hausdorff
  \end{enumerate}
\end{prop}

\begin{defn} Ruang topologis disebut
 \index{regular!ruang topologis}%
 \index{topologis!ruang!regular}%
\df{regular} jika titik-titik dan himpunan-himpunan tertutup dapat dipisahkan oleh himpunan terbuka; yaitu, jika $C$ adalah subhimpunan tertutup dari ruang tersebut dan $x$ adalah titik
dari ruang tersebut yang tidak berada dalam $C$, maka terdapat himpunan terbuka $U$ dan $V$ yang saling lepas sedemikian sehingga $C \subseteq U$ dan $x \in V$.
\end{defn}

\begin{prop} Setiap ruang vektor topologis Hausdorff bersifat regular.
\end{prop}

\begin{cor} Jika $X$ adalah ruang vektor topologis yang setiap titiknya merupakan himpunan tertutup, maka $X$ bersifat regular.
\end{cor}

\begin{prop} Misalkan $X$ ruang vektor topologis Hausdorff.  Jika $M$ adalah subruang dari $X$ dan $F$ adalah subruang vektor berdimensi hingga
dari $X$, maka $M + F$ tertutup di~$X$. (Khususnya, $F$ tertutup.)
\end{prop}

\begin{defn} Suatu subhimpunan $B$ dari ruang vektor topologis disebut
 \index{terbatas!himpunan dalam suatu TVS}%
\df{terbatas} jika untuk setiap lingkungan $U$ dari $\vc 0$ dalam ruang tersebut terdapat $\alpha > 0$ sedemikian sehingga $B \subseteq \alpha U$.
\end{defn}

\begin{prop} Suatu subhimpunan dari ruang vektor topologis bersifat terbatas jika dan hanya jika diserap oleh setiap lingkungan~$\vc 0$.
\end{prop}

\begin{prop} Suatu subhimpunan $B$ dari ruang vektor topologis bersifat terbatas jika dan hanya jika $\alpha_n x_n \sto \vc 0$ kapan pun $(x_n)$ adalah
barisan vektor dalam $B$ dan $(\alpha_n)$ adalah barisan skalar dalam $c_0$.
\end{prop}

\begin{prop} Jika suatu subhimpunan dari ruang vektor topologis bersifat terbatas, maka tutupannya juga terbatas.
\end{prop}

\begin{prop} Jika subhimpunan $A$ dan $B$ dari ruang vektor topologis bersifat terbatas, maka $A \cup B$ juga terbatas.
\end{prop}

\begin{prop} Jika subhimpunan $A$ dan $B$ dari ruang vektor topologis bersifat terbatas, maka $A + B$ juga terbatas.
\end{prop}

\begin{prop} Setiap subhimpunan kompak dari ruang vektor topologis adalah terbatas.
\end{prop}

Seperti halnya pada ruang linear bernorma, kita mengatakan bahwa pemetaan linear bersifat \emph{terbatas} jika memetakan himpunan terbatas ke himpunan terbatas.

\begin{defn} Suatu pemetaan linear $T\colon X \sto Y$ antara ruang-ruang vektor topologis disebut
 \index{terbatas!pemetaan linear}%
\df{terbatas} jika $T^\sto(B)$ merupakan subhimpunan terbatas dari $Y$ kapan pun $B$ merupakan subhimpunan terbatas dari~$X$.
\end{defn}

\begin{prop} Setiap pemetaan linear kontinu antara ruang-ruang vektor topologis bersifat terbatas.
\end{prop}

\begin{defn} Suatu filter $\sfml F$ pada subhimpunan $A$ dari ruang vektor topologis $X$ adalah
 \index{Cauchy!filter}%
 \index{filter!Cauchy}%
\df{filter Cauchy} jika untuk setiap lingkungan $U$ dari titik asal dalam $X$ terdapat $B \in \sfml F$ sedemikian sehingga
   \[ B - B \subseteq U\,. \]
\end{defn}

\begin{prop} Pada ruang vektor topologis, setiap filter yang konvergen merupakan filter Cauchy.
\end{prop}

\begin{defn} Suatu subhimpunan $A$ dari ruang vektor topologis $X$ disebut
 \index{lengkap!ruang vektor topologis}%
 \index{topologis!ruang vektor!lengkap}%
 \index{ruang!vektor topologis!lengkap}%
\df{lengkap} jika setiap filter Cauchy pada $A$ konvergen ke suatu titik di~$A$.
\end{defn}

\begin{prop} Setiap subhimpunan lengkap dari ruang vektor topologis Hausdorff adalah tertutup.
\end{prop}

\begin{prop} Jika $C$ merupakan subhimpunan tertutup dari suatu himpunan lengkap dalam ruang vektor topologis, maka $C$ lengkap.
\end{prop}



\begin{prop} Suatu pemetaan linear $T\colon X \sto Y$ antara ruang-ruang vektor topologis bersifat kontinu jika dan hanya jika pemetaan tersebut kontinu di nol.
\end{prop}



















\section{Hasil Bagi}

\begin{conv} Seperti halnya pada ruang Hilbert dan Banach, kata ``subruang'' dalam konteks ruang vektor topologis
 \index{subruang!dari ruang vektor topologis}%
 \index{konvensi!subruang ruang vektor topologis adalah tertutup}%
berarti \emph{subruang vektor tertutup}.
\end{conv}

\begin{defn} Misalkan $M$ suatu subruang vektor dari ruang vektor topologis~$X$ dan $\pi\colon X \sto X/M \colon x \mapsto [x]$
adalah pemetaan hasil bagi yang biasa. Pada ruang vektor hasil bagi $X/M$ kita definisikan suatu topologi, yang kita sebut
 \index{hasil bagi!topologi}%
 \index{hasil bagi!ruang vektor topologis}%
 \index{topologis!ruang vektor!hasil bagi}%
\df{topologi hasil bagi}, dengan menetapkan filter lingkungan di titik asal pada~$X/M$ sebagai citra oleh $\pi$ dari
filter lingkungan di titik asal pada~$X$. Artinya, kita menyatakan suatu subhimpunan $V$ dari $X/M$ sebagai lingkungan $\vc 0$ jika,
dan hanya jika, terdapat suatu lingkungan $U$ dari titik asal di $X$ sedemikian sehingga $V = \pi^{\sto}(U)$.
\end{defn}

\begin{prop} Misalkan $M$ suatu subruang vektor dari ruang vektor topologis~$X$. Jika $X/M$ diberi topologi hasil bagi, maka
pemetaan hasil bagi $\pi\colon X \sto X/M$ merupakan pemetaan terbuka sekaligus kontinu.
\end{prop}

\begin{exam} Misalkan $M$ adalah sumbu~$y$ dalam ruang vektor topologis $\R^2$ (dengan topologinya yang biasa). Identifikasikan $\R^2/M$ dengan
sumbu~$x$. Meskipun pemetaan hasil bagi memetakan himpunan terbuka ke himpunan terbuka, hiperbola $\{(x,y) \in \R^2\colon xy = 1\}$ menunjukkan
bahwa $\pi$ tidak harus memetakan himpunan tertutup ke himpunan tertutup.
\end{exam}

\begin{prop}\label{prop_quotient_top_strong} Misalkan $M$ suatu subruang vektor dari ruang vektor topologis~$X$. Topologi hasil bagi
pada $X/M$ adalah topologi terbesar yang membuat pemetaan hasil bagi kontinu.
\end{prop}

Proposisi berikut menjamin bahwa topologi hasil bagi, sebagaimana didefinisikan di atas, menjadikan ruang vektor hasil bagi suatu ruang vektor topologis.

\begin{prop} Misalkan $M$ suatu subruang vektor dari ruang vektor topologis~$X$. Topologi hasil bagi pada $X/M$ kompatibel dengan
operasi-operasi ruang vektor pada ruang tersebut.
\end{prop}

\begin{prop} Jika $M$ suatu subruang vektor dari ruang vektor topologis~$X$, maka ruang hasil bagi $X/M$ bersifat Hausdorff jika dan hanya
jika $M$ tertutup dalam~$X$.
\end{prop}

\begin{exer} Jelaskan pula bagaimana proposisi sebelumnya memungkinkan kita mengaitkan suatu ruang lain yang bersifat Hausdorff dengan
suatu ruang vektor topologis non-Hausdorff.
\end{exer}

\begin{prop} Misalkan $T\colon X \sto Y$ suatu pemetaan linear antara ruang-ruang vektor topologis dan $M$ suatu subruang dari $X$. Jika
 \index{hasil bagi!teorema!untuk $\cat{TVS}$}%
$M \subseteq \ker T$, maka terdapat fungsi unik $\wt T\colon X/M \sto Y$ sedemikian sehingga $T = \wt T \circ \pi$ (dengan $\pi$
adalah pemetaan hasil bagi). Fungsi $\wt T$ bersifat linear; fungsi ini kontinu jika dan hanya jika $T$ kontinu; fungsi ini terbuka jika dan hanya jika $T$~terbuka.
\end{prop}
















\section{Ruang Konveks Lokal dan Seminorma}

\begin{prop} Dalam suatu ruang vektor topologis, tutupan suatu himpunan konveks adalah konveks.
\end{prop}

\begin{prop} Dalam suatu ruang vektor topologis, setiap lingkungan konveks dari $\vc 0$ memuat suatu lingkungan seimbang dan konveks dari~$\vc 0$.
\end{prop}

\begin{prop} Suatu subhimpunan tertutup $C$ dari ruang vektor topologis bersifat konveks jika dan hanya jika $\frac12(x + y)$ termasuk dalam $C$ setiap kali $x$ dan $y$ juga demikian.
\end{prop}

\begin{prop} Jika suatu subhimpunan dari ruang vektor topologis bersifat konveks, maka tutupan dan interiornya juga konveks.
\end{prop}

\begin{defn} Suatu ruang vektor topologis disebut
 \index{lokal!kekonveksan}%
 \index{konveks!lokal}%
\df{konveks lokal} jika ruang tersebut memiliki basis lokal yang terdiri atas himpunan-himpunan konveks. Untuk ringkasnya, suatu ruang vektor topologis
yang konveks lokal cukup kita sebut sebagai
 \index{lokal!ruang konveks}%
\df{ruang konveks lokal}.
\end{defn}

\begin{prop} Jika $p$ suatu seminorma pada ruang vektor, maka $p(\vc 0) = 0$.
 \index{seminorma!sifat-sifat elementer}%
\end{prop}

\begin{prop} Jika $p$ suatu seminorma pada ruang vektor $V$, maka $\abs{p(x) - p(y)} \le p(x - y)$ untuk semua $x$, $y \in V$.
\end{prop}

\begin{cor} Jika $p$ suatu seminorma pada ruang vektor $V$, maka $p(x) \ge 0$ untuk semua $x \in V$.
\end{cor}

\begin{prop} Jika $p$ suatu seminorma pada ruang vektor $V$, maka $p^{\gets}(\{0\})$ adalah subruang vektor dari~$V$.
\end{prop}

\begin{defn} Misalkan $p$ suatu seminorma pada ruang vektor $V$ dan $\epsilon > 0$. Himpunan
 \index{B@$B_{p,\epsilon}$, $B_p$ (semibola terbuka yang ditentukan oleh~$p$)}%
$B_{p,\epsilon} := p^{\gets}\bigl([0,\epsilon)\bigr)$ kita sebut
 \index{terbuka!semibola}%
 \index{semibola}%
\df{semibola terbuka} berjari-jari $\epsilon$ di titik asal yang ditentukan oleh~$p$. Demikian pula, himpunan
\index{C@$C_{p,\epsilon}$, $C_p$ (semibola tertutup yang ditentukan oleh~$p$)}%
$C_{p,\epsilon} := p^{\gets}\bigl([0,\epsilon]\bigr)$ adalah
 \index{tertutup!semibola}%
\df{semibola tertutup} berjari-jari $\epsilon$ di titik asal yang ditentukan oleh~$p$. Untuk semibola terbuka dan tertutup berjari-jari satu yang ditentukan
oleh $p$, gunakan masing-masing notasi $B_p$ dan $C_p$.
\end{defn}

\begin{prop} Misalkan $\fml P$ suatu keluarga seminorma pada ruang vektor~$V$. Keluarga $\sfml U$ dari semua irisan hingga dari semibola-semibola terbuka
di titik asal $V$ menginduksi suatu topologi (yang tidak harus Hausdorff) pada $V$ yang kompatibel dengan struktur linear $V$ dan yang terhadapnya
$\sfml U$ merupakan suatu basis lokal.
\end{prop}

\begin{proof}[\emph{Petunjuk untuk bukti}] Proposisi~\ref{prop_fbase_induces_top}.  \ns
\end{proof}

\begin{prop} Misalkan $p$ seminorma pada ruang vektor topologis~$X$. Pernyataan berikut ekuivalen:
   \begin{enumerate}[label=\rm{(\alph*)}]
     \item $B_p$ terbuka dalam~$X$;
     \item $p$ kontinu di titik asal; dan
     \item $p$ kontinu.
   \end{enumerate}
\end{prop}

\begin{prop} Jika $p$ dan $q$ merupakan seminorma kontinu pada suatu ruang vektor topologis, maka $p + q$ dan~$\max\{p,q\}$ juga demikian.
\end{prop}

\begin{exam} Misalkan $N$ suatu subhimpunan yang bersifat menyerap, seimbang, dan konveks dari ruang vektor~$V$. Definisikan
   \[ \sbsb \mu N\colon V \sto \R \colon x \mapsto \inf\{\lambda > 0\colon x \in \lambda N\}. \]
Fungsi $\sbsb \mu N$ adalah seminorma pada~$V$. Fungsi ini disebut
 \index{Minkowski!fungsional}%
 \index{fungsional!Minkowski}%
 \index{mu@$\sbsb \mu N$ (fungsional Minkowski untuk~$N$)}%
\df{fungsional Minkowski} dari~$N$.
\end{exam}

\begin{prop} Jika $N$ suatu himpunan penyerap, seimbang, dan konveks dalam ruang vektor~$V$ dan $\sbsb \mu N$ adalah fungsional Minkowski dari $N$, maka
     \[ B_{\sbsb \mu{\!N}} \subseteq N \subseteq C_{\sbsb \mu{\!N}} \]
dan fungsional-fungsional Minkowski yang dibangkitkan oleh ketiga himpunan $B_{\sbsb \mu{\!N}}$, $N$, dan $C_{\sbsb \mu{\!N}}$ semuanya identik.
\end{prop}

\begin{prop} Jika $N$ suatu himpunan penyerap, seimbang, dan konveks dalam ruang vektor~$V$ dan $\sbsb \mu N$ adalah fungsional Minkowski dari $N$, maka
$\sbsb \mu N$ kontinu jika dan hanya jika $N$ merupakan lingkungan nol; dalam hal ini
  \[ B_{\sbsb \mu{\!N}} = \intr N \quad \text{dan} \quad  C_{\sbsb \mu{\!N}} = \clo N. \]
\end{prop}

\begin{cor} Misalkan $N$ suatu subhimpunan tertutup yang bersifat menyerap, seimbang, dan konveks dari ruang vektor topologis~$X$. Maka fungsional Minkowski
$\sbsb \mu N$ adalah seminorma unik pada $X$ sedemikian sehingga $N = C_{\sbsb \mu{\!N}}$.
\end{cor}

\begin{prop} Misalkan $p$ suatu seminorma pada ruang vektor $V$. Maka himpunan $B_p$ bersifat menyerap, seimbang, dan konveks. Selain itu, $p = \sbsb \mu{B_p}$.
\end{prop}

\begin{defn} Suatu keluarga seminorma $\fml P$ pada ruang vektor $V$ disebut
 \index{pemisah!keluarga seminorma}%
\df{pemisah} jika untuk setiap $x \in V$ dengan $x \ne \vc 0$ terdapat $p \in \fml P$ sedemikian sehingga $p(x) \ne 0$.
\end{defn}

\begin{prop} Dalam suatu ruang vektor topologis Hausdorff, misalkan $\sfml B$ suatu basis lokal yang terdiri atas himpunan terbuka, seimbang, dan konveks. Maka
$\{\sbsb \mu{\!N}\colon N \in \sfml B\}$ adalah suatu keluarga seminorma pemisah pada ruang tersebut; setiap anggotanya kontinu, dan
$N = B_{\sbsb \mu{\!N}}$ berlaku untuk setiap $N \in \sfml B$.
\end{prop}

Proposisi berikut memperlihatkan bagaimana banyak ruang konveks lokal yang paling penting muncul---dari suatu keluarga
seminorma pemisah.

\begin{prop}\label{prop_clbase_from_snorms} Misalkan $\fml P$ suatu keluarga seminorma pemisah pada ruang vektor~$V$. Untuk setiap
$p \in \fml P$ dan setiap $n \in \N$, misalkan $U_{p,n} = p^{\gets}\bigl([0,\frac1n)\bigr)$. Maka keluarga semua irisan hingga
dari himpunan-himpunan $U_{p,n}$ membentuk suatu basis lokal yang seimbang dan konveks bagi suatu topologi pada $V$ yang menjadikan $V$
suatu ruang konveks lokal dan setiap seminorma dalam $\fml P$ kontinu.
\end{prop}

\begin{exam}\label{X_LCS1} Misalkan $X$ suatu ruang Hausdorff yang kompak lokal. Untuk setiap subhimpunan kompak tak kosong $K$ dari $X$, definisikan
   \[ \sbsb pK\colon \fml C(X) \sto \R\colon f \mapsto \sup\{\abs{f(x)}\colon x \in K\}. \]
Maka keluarga $\fml P$ dari semua $\sbsb pK$, dengan $K$ suatu subhimpunan kompak tak kosong dari $X$, merupakan suatu keluarga seminorma
yang menjadikan $\fml C(X)$ suatu ruang konveks lokal.
\end{exam}

\begin{prop} Suatu pemetaan linear $T\colon V \sto W$ dari suatu ruang vektor topologis ke suatu ruang konveks lokal bersifat kontinu jika dan hanya jika
$p \circ T$ merupakan seminorma kontinu pada $V$ untuk setiap seminorma kontinu $p$ pada~$W$.
\end{prop}

\begin{prop} Misalkan $f$ suatu fungsional linear kontinu pada subruang $M$ dari ruang konveks lokal~$X$. Maka $f$ dapat diperluas menjadi suatu
fungsional linear kontinu pada seluruh~$X$.
\end{prop}



\begin{proof} Lihat~\cite{Conway:1990}, Proposisi IV.5.13.   \ns
\end{proof}



















\section{Ruang Fr\'echet}

\begin{defn} Ruang topologis $(X,\sfml T)$ disebut
 \index{dapat dimetrikkan}%
\df{dapat dimetrikkan} jika terdapat suatu metrik $d$ pada $X$ sedemikian sehingga topologi yang dibangkitkan oleh $d$ adalah~$\sfml T$.
\end{defn}

\begin{defn} Suatu metrik $d$ pada ruang vektor $V$ disebut
 \index{translasi!invarian}%
 \index{invarian!terhadap translasi}%
\df{invarian terhadap translasi} jika
   \[ d(x,y) = d(x + a, y + a) \]
untuk semua $x$, $y$, dan $a$ di~$V$.
\end{defn}

\begin{prop} Misalkan $(X,\sfml T)$ ruang vektor topologis Hausdorff yang memiliki basis lokal terhitung. Maka terdapat metrik
invarian terhadap translasi $d$ pada $X$ yang membangkitkan topologi $\sfml T$ dan yang bola-bola terbukanya di titik asal seimbang. Jika $X$ ruang konveks lokal,
maka $d$ dapat dipilih sedemikian sehingga, selain itu, bola-bola terbukanya konveks.
\end{prop}

\begin{proof} Lihat \cite{Rudin:1991}, Teorema 1.24.  \ns
\end{proof}

\begin{exam}\label{X_metric_from_seminorms} Misalkan $\{p_n\colon n \in \N\}$ suatu keluarga seminorma pemisah terhitung pada ruang
vektor~$X$, dan misalkan $\sfml T$ topologi pada $X$ yang dijamin oleh Proposisi~\ref{prop_clbase_from_snorms}. Dalam hal ini kita dapat secara eksplisit
mendefinisikan suatu metrik invarian terhadap translasi pada $X$ yang menginduksi topologi~$\sfml T$. Untuk setiap barisan $\bigl(\alpha_k\bigr) \in c_0$
dengan $\alpha_k>0$ untuk setiap $k\in\N$, suatu metrik dengan sifat-sifat yang diinginkan diberikan oleh
   \[ d(x,y) = \max_{k \in \N}\biggl\{\frac{\alpha_k p_k(x-y)}{1 + p_k(x-y)}\biggr\}. \]
\end{exam}

\begin{proof} Lihat \cite{Rudin:1991}, Catatan 1.38(c).  \ns
\end{proof}

\begin{defn} Sebuah
 \index{ruang Fr\'echet}%
\df{ruang Fr\'echet} adalah ruang konveks lokal yang lengkap dan dapat dimetrikkan.
\end{defn}

\begin{notn} Misalkan $A$ dan $B$ subhimpunan suatu ruang topologis. Kita menulis
 \index{<binrell@$A \prec B$ ($\clo A \subseteq \intr B$)}%
$A\prec B$ jika $\clo A \subseteq \intr B$.
\end{notn}

\begin{prop} Pernyataan berikut berlaku untuk subhimpunan $A$, $B$, dan $C$ dari suatu ruang topologis:
 \begin{enumerate}
  \item[(a)] jika $A \prec B$ dan $B \prec C$, maka $A \prec C$;
  \item[(b)] jika $A \prec C$ dan $B \prec C$, maka $A \cup B \prec C$; dan
  \item[(c)] jika $A \prec B$ dan $A \prec C$, maka $A \prec B \cap C$.
 \end{enumerate}
\end{prop}

Berikut adalah hasil standar dari topologi elementer.

\begin{prop}\label{prop_open_lim_cpt} Misalkan $\Omega$ subhimpunan terbuka tak kosong dari $\R^n$. Maka terdapat suatu barisan
$(K_n)$ yang terdiri atas subhimpunan kompak tak kosong dari $\Omega$ sedemikian sehingga $K_i \prec K_j$ setiap kali $i < j$ dan
\smash[b]{$\bigcup\limits_{i=1}^\infty K_i = \Omega$}.
\end{prop}

\begin{proof} Lihat \cite{Treves:1967}, Lema 10.1.  \ns
\end{proof}

\begin{exam} Misalkan $\Omega$ subhimpunan terbuka tak kosong dari suatu ruang Euklides $\R^n$, dan misalkan $\fml C(\Omega)$ keluarga semua
fungsi kontinu bernilai skalar pada~$\Omega$. Berdasarkan proposisi sebelumnya, kita dapat menuliskan $\Omega$ sebagai gabungan suatu koleksi
terhitung himpunan kompak tak kosong $K_1$, $K_2$, \dots sedemikian sehingga $K_i \prec K_j$ setiap kali $i < j$. Untuk setiap $n \in \N$ dan
$f \in \fml C(\Omega)$, misalkan $p_n(f) = \sup\{\abs{f(x)}\colon  x \in K_n\}$ seperti dalam Contoh~\ref{X_LCS1}. Menurut
Proposisi~\ref{prop_clbase_from_snorms}, keluarga seminorma ini menjadikan $\fml C(\Omega)$ ruang konveks lokal. Misalkan
$\sfml T$ topologi yang dihasilkan pada ruang ini. Berdasarkan Contoh~\ref{X_metric_from_seminorms}, fungsi
   \[ d(f,g) = \max_{k \in \N}\frac{2^{-k}p_k(f-g)}{1 + p_k(f-g)} \]
mendefinisikan suatu metrik pada $\fml C(\Omega)$ yang menginduksi topologi~$\sfml T$. Dengan metrik ini,
$\fml C(\Omega)$ adalah ruang Fr\'echet.
\end{exam}

\begin{notn}[Notasi multi-indeks]\label{mi_notn} Misalkan $\Omega$ subhimpunan terbuka dari suatu ruang Euklides~$\R^n$. Kita akan meninjau
fungsi-fungsi bernilai skalar yang dapat didiferensialkan
 \index{notasi multi-indeks}%
tak hingga pada $\Omega$; yakni fungsi pada $\Omega$ yang memiliki turunan untuk setiap orde di tiap titik~$\Omega$. Kita
akan menyebut fungsi-fungsi demikian sebagai
 \index{fungsi mulus}%
fungsi \df{mulus} pada~$\Omega$. Kelas semua fungsi mulus pada $\Omega$ dinotasikan dengan
 \index{C@$\fml C^\infty(\Omega)$!fungsi yang dapat didiferensialkan tak hingga}%
$\fml C^\infty(\Omega)$.

Kelas penting lain yang akan kita tinjau
 \index{C@$\fml C^\infty_c(\Omega)$!fungsi yang dapat didiferensialkan tak hingga dengan tumpuan kompak}%
adalah $\fml C_c^\infty(\Omega)$, yaitu keluarga semua fungsi bernilai skalar pada $\Omega$ yang dapat didiferensialkan tak hingga dan memiliki tumpuan kompak;
yakni fungsi yang bernilai nol di luar suatu subhimpunan kompak dari~$\Omega$. Fungsi-fungsi dalam keluarga ini sering disebut
 \index{fungsi uji}%
\df{fungsi uji} pada~$\Omega$. Notasi lain yang lebih singkat,
 \index{D@$\fml D(\Omega)$!ruang fungsi uji; fungsi yang dapat didiferensialkan tak hingga dengan tumpuan kompak}%
yang lazim digunakan untuk keluarga ini adalah~$\fml D(\Omega)$.

Kita akan menaruh perhatian pada operator diferensial pada ruang~$\fml D(\Omega)$. Misalkan $\bs\alpha = (\alpha_1, \dots, \alpha_n)$ suatu
 \index{multi-indeks}%
\df{multi-indeks}, yaitu suatu tupel-$n$ bilangan bulat dengan setiap $\alpha_k \ge 0$, dan misalkan
  \[ D^{\bs\alpha} = \biggl(\frac{\partial}{\partial x_1}\biggr)^{\alpha_1} \dots
                                                     \biggl(\frac{\partial}{\partial x_n}\biggr)^{\alpha_n}. \]
Yang dimaksud dengan
\index{D@$D^{\bs\alpha}$ (notasi multi-indeks untuk diferensiasi)}%
 \index{orde!multi-indeks}%
 \index{operator diferensial!yang bekerja pada fungsi mulus}%
\df{orde} operator diferensial $D^{\bs\alpha}$ adalah $\abs{\bs\alpha} := \sum_{k=1}^n \alpha_k$. (Ketika
$\abs{\bs\alpha} = 0$, kita memahami $D^{\bs\alpha}(f) = f$ untuk setiap fungsi~$f$.) Notasi alternatif untuk
 \index{<unopura@$f^{\bs\alpha}$ (notasi multi-indeks untuk diferensiasi)}%
$D^{\bs\alpha}f$ adalah~$f^{(\bs\alpha)}$.
\end{notn}

\begin{exam}\label{X_smooth_fcns_Frechet} Pada keluarga $\fml C^\infty(\Omega)$ dari fungsi-fungsi mulus pada suatu subhimpunan terbuka $\Omega$ dari
ruang Euklides~$\R^n$, definisikan keluarga seminorma $\{\sbsb{p}{K,j}\}$. Untuk setiap subhimpunan kompak $K$ dari $\Omega$ dan setiap
$j \in \N$, misalkan
   \[ \sbsb{p}{K,j}(f) = \sup\{\abs{D^{\bs\alpha}f(x)}\colon x \in K \text{ dan } \abs{\bs\alpha} \le j\}. \]
Keluarga seminorma ini menjadikan $\fml C^\infty(\Omega)$ ruang Fr\'echet. Topologi yang dihasilkan pada $\fml C^\infty(\Omega)$
kadang-kadang disebut
 \index{topologi!konvergensi seragam pada himpunan-himpunan kompak}%
 \index{seragam!konvergensi!pada himpunan-himpunan kompak}%
 \index{himpunan kompak!topologi konvergensi seragam pada}%
 \index{konvergensi!seragam!pada himpunan-himpunan kompak}%
\emph{topologi konvergensi seragam pada himpunan-himpunan kompak bagi fungsi-fungsi tersebut dan semua turunannya}.
\end{exam}

\begin{proof}[\emph{Petunjuk untuk bukti}] Untuk melihat bahwa ruang ini dapat dimetrikkan, gunakan Contoh~\ref{X_metric_from_seminorms} dan
Proposisi~\ref{prop_open_lim_cpt}.  \ns
\end{proof}

\begin{exam}\label{Schwartz_space} Misalkan $n \in \N$. Nyatakan dengan $\fml S = \fml S(\R^n)$
 \index{S@$\fml S = \fml S(\R^n)$ (ruang Schwartz)!sebagai ruang Fr\'echet}%
ruang semua fungsi mulus $f$ yang didefinisikan pada seluruh~$\R^n$ dan memenuhi
   \[ \lim_{\norm x \sto \infty}{\norm x}^k \abs{D^{\bs\alpha}f(x)} = 0 \]
untuk semua multi-indeks $\bs\alpha = (\alpha_1,\dots,\alpha_n) \in \N^n$ dan semua bilangan bulat $k \ge 0$. Fungsi-fungsi demikian sering disebut
sebagai
 \index{fungsi yang menurun cepat}%
 \index{fungsi!menurun cepat}%
 \index{turunan!fungsi dengan turunan yang menurun cepat}%
 \index{ruang!fungsi yang menurun cepat}%
\emph{fungsi dengan turunan yang menurun cepat}, atau lebih singkat lagi, \emph{fungsi yang menurun cepat}. Pada $\fml S$ kita gunakan
topologi yang didefinisikan oleh seminorma-seminorma
   \[ \sbsb p{m,k}(f) := \sup_{\abs{\bs\alpha} \le m}\bigl\{\sup_{x \in \R^n} \{(1 + {\norm x}^2)^k \abs{(D^{\bs\alpha}f)(x)}\}\bigr\} \]
untuk $k$, $m = 0,1,2,\dots$. Ruang yang dihasilkan adalah ruang Fr\'echet yang disebut
 \index{ruang Schwartz}%
 \index{ruang!Schwartz}%
\df{ruang Schwartz} dari~$\R^n$.
\end{exam}

\begin{proof}[\emph{Petunjuk untuk bukti}] Untuk bukti kelengkapan, lihat~\cite{Treves:1967}, Contoh IV, halaman 92--94. \ns
\end{proof}

\begin{prop} Suatu fungsi uji $\psi$ pada $\R$ merupakan turunan dari fungsi uji lain jika dan hanya jika $\int_{-\infty}^\infty\psi = 0$.
\end{prop}

\begin{prop} Misalkan $F$ ruang Fr\'echet yang topologinya diinduksi oleh suatu metrik invarian terhadap translasi. Maka setiap subhimpunan terbatas
dari $F$ memiliki diameter hingga.
\end{prop}

\begin{exam} Tinjau ruang Fr\'echet $\R$ dengan topologi yang diinduksi oleh metrik invarian terhadap translasi
\smash[b]{$d(x,y) = \dfrac{\abs{x - y}}{1 + \abs{x - y}}$}. Maka himpunan bilangan asli $\N$ memiliki diameter hingga, tetapi tidak terbatas.
\end{exam}




































\endinput

 \chapter{DISTRIBUSI}

\section{Limit Induktif}\label{sec_ind_limits}
\begin{defn}\label{defn_dir_sys} Misalkan $D$ suatu himpunan terarah terhadap pengurutan parsial~$\le$ dan misalkan $\{A_i\colon i \in D\}$ suatu keluarga
objek dalam suatu kategori~$\cat C$.  Anggap bahwa untuk setiap $i$, $j \in D$ dengan $i \le j$ terdapat suatu morfisme
$\phi_{j,i}\colon A_i \sto A_j$ yang memenuhi $\phi_{ki} = \phi_{kj}\phi_{ji}$ setiap kali $i \le j \le k$ dalam~$D$.  Anggap pula bahwa
$\phi_{ii}$ adalah pemetaan identitas pada $A_i$ untuk setiap $i \in D$. Maka keluarga berindeks $\vc A = \bigl(A_i\bigr)_{i \in D}$ dari
objek-objek bersama keluarga berindeks $\bs\phi = \bigl(\phi_{ji}\bigr)$ dari \emph{morfisme penghubung} disebut suatu
 \index{morfisme penghubung}%
 \index{morfisme!penghubung}%
 \index{terarah!sistem}%
 \index{sistem!terarah}%
\df{sistem terarah} dalam~$\cat C$.
\end{defn}

\begin{defn}\label{defn_ind_lim} Misalkan pasangan $(\vc A, \bs\phi)$ suatu sistem terarah (seperti di atas) dalam kategori~$\cat C$ yang
himpunan terarah dasarnya adalah~$D$.  Suatu
 \index{induktif!limit}%
 \index{limit!induktif}%
\df{limit induktif} (atau
 \index{langsung!limit}%
 \index{limit!langsung}%
\df{limit langsung}) dari sistem $(\vc A,\bs\phi)$ adalah pasangan $(L,\bs\mu)$ dengan $L$ suatu objek dalam $\cat C$
dan $\bs\mu = \bigl(\mu_i\bigr)_{i \in D}$ suatu keluarga morfisme $\mu_j\colon A_j \sto L$ dalam $\cat C$ yang memenuhi
 \begin{enumerate}[label=\rm{(\alph*)}]
  \item $\mu_i = \mu_j \phi_{ji}$ setiap kali $i \le j$ dalam $D$, dan
  \item jika $(M,\bs\lambda)$ suatu pasangan dengan $M$ suatu objek dalam $\cat C$ dan $\bs\lambda = \bigl(\lambda_i\bigr)_{i \in D}$
adalah suatu keluarga berindeks morfisme $\lambda_j\colon A_j \sto M$ dalam $\cat C$ yang memenuhi $\lambda_i = \lambda_j\phi_{ji}$ setiap kali $i \le j$ dalam~$D$,
maka terdapat morfisme tunggal $\psi\colon L \sto M$ sedemikian sehingga $\lambda_i = \psi\mu_i$ untuk setiap $i \in D$.
 \end{enumerate}
Dengan penyalahgunaan bahasa yang lazim, biasanya kita mengatakan bahwa $L$ adalah limit induktif dari sistem $\vc A = (A_i)$ dan
 \index{<arrowto@$\underrightarrow{\lim} A_i$ (limit induktif)}%
menulis $L = \underrightarrow{\lim} A_i$.

\vskip 15 pt

   \[\xymatrix@+15pt@C+25pt{&&&& L\ar@{-->}[dddd]^\psi \\
                   &&&&   \\
                  {\cdots}\ar[r] & A_i \ar[uurrr]^{\mu_i}\ar[ddrrr]^{\lambda_i}\ar[r]^(.6){\phi_{ji}} &
                       A_j\ar[uurr]^{\mu_j}\ar[ddrr]^{\lambda_j}\ar[r] & {\cdots} &  \\
                   &&&&   \\
                   &&&& M
    }\]
\end{defn}

Dalam suatu kategori (konkret) sebarang, limit induktif belum tentu ada.  Namun, jika ada, limit tersebut tunggal.

\begin{prop}\label{prop_ind_lim_unique} Limit induktif (jika ada dalam suatu kategori) bersifat tunggal (hingga isomorfisme).
\end{prop}

\begin{exam} Misalkan $S$ suatu objek tak kosong dalam kategori $\cat{SET}$. Tinjau keluarga $\sfml P(S)$ dari semua himpunan bagian $S$ yang diarahkan
oleh pengurutan parsial~$\subseteq$.  Setiap kali $A \subseteq B \subseteq S$, ambil morfisme penghubung $\sbsb\iota{BA}$ sebagai
pemetaan inklusi dari $A$ ke~$B$. Maka $\sfml P(S)$, bersama keluarga morfisme penghubung, membentuk suatu sistem terarah.  Limit
induktif dari sistem ini adalah~$S$.
\end{exam}

\begin{exam} Limit langsung ada dalam kategori $\cat{VEC}$ dari ruang vektor dan pemetaan linear.
\end{exam}

\begin{proof}[\emph{Petunjuk untuk bukti}] Misalkan $(\vc V, \bs\phi)$ suatu sistem terarah dalam $\cat{VEC}$ (seperti dalam definisi~\ref{defn_dir_sys})
dan $D$ adalah himpunan terarah dasarnya. Misalkan
$W=\bigoplus_{i\in D}V_i$ dan $\iota_i\colon V_i\sto W$ inklusi kanonik. Tetapkan
  \[ N=\operatorname{span}\{\iota_i(u)-\iota_j(\phi_{ji}(u)):
       i\le j,\quad u\in V_i\}. \]
Ambil $L=W/N$ dan $\mu_i=q\circ\iota_i$, dengan $q\colon W\sto W/N$ pemetaan hasil bagi, lalu verifikasikan sifat universalnya
(lihat contoh~\ref{C015544} dan~\ref{X_dir_sum_VEC}).  \ns
\end{proof}

\begin{exam} Barisan $\Z \to^{\vc 1} \Z \to^{\vc 2} \Z \to^{\vc 3} \Z \to^{\vc 4} \dots$
(dengan $\vc n \colon \Z \sto \Z$ memenuhi $\vc n(1) = n$) merupakan suatu barisan induktif grup-grup Abelian
yang limit induktifnya adalah himpunan $\Q$ dari bilangan rasional.
\end{exam}

\begin{exam} Barisan $\Z \to^{\vc 2} \Z \to^{\vc 2} \Z \to^{\vc 2} \Z \to^{\vc 2} \dots$
merupakan suatu barisan induktif grup-grup Abelian yang limit induktifnya adalah himpunan bilangan rasional
diadik.
\end{exam}









\section{Ruang-$LF$}\label{sec_LFsps}
Karena tujuan utama kita dalam bab ini adalah pengantar teori distribusi, mulai sekarang kita akan membatasi perhatian
pada jenis limit induktif yang sangat sederhana---\emph{limit induktif ketat dari suatu barisan induktif}. Secara khusus, kita akan
tertarik pada limit induktif ketat dari barisan ruang Fr\'echet yang meningkat.


Ingat bahwa topologi lemah pada suatu himpunan $S$ diinduksi oleh suatu keluarga fungsi $f_\alpha \colon S \sto
X_\alpha$ yang memetakan ke ruang-ruang topologis. Topologi ini adalah topologi terlemah pada $S$ yang membuat
semua fungsi tersebut kontinu (lihat latihan~\ref{C021421}). Gagasan dualnya adalah
 \index{topologi!kuat}%
 \index{kuat!topologi}%
\df{topologi kuat}, yakni topologi yang diinduksi oleh suatu keluarga fungsi yang memetakan \emph{dari}
ruang-ruang topologis $X_\alpha$ \emph{ke dalam}~$S$. Topologi ini adalah topologi terkuat yang membuat
semua fungsi tersebut kontinu. Salah satu topologi kuat yang telah kita jumpai adalah
\emph{topologi hasil bagi} (lihat proposisi~\ref{prop_quotient_top_strong}). Contoh penting lainnya,
yang akan kita perkenalkan berikutnya, adalah \emph{topologi limit induktif}.


\begin{defn} Suatu
 \index{induktif!barisan!ketat}%
 \index{barisan!induktif ketat}%
 \index{ketat!barisan induktif}%
\df{barisan induktif ketat} adalah barisan $(X_i)$ ruang konveks lokal sedemikian sehingga untuk setiap $i \in \N$
   \begin{enumerate}
     \item[(i)] $X_i$ adalah subruang tertutup dari $X_{i+1}$ dan
     \item[(ii)] topologi $\sfml T_i$ pada $X_i$ adalah pembatasan topologi $\sfml T_{i+1}$ pada $X_i$.
   \end{enumerate}
Ini, tentu saja, merupakan sistem terarah dengan morfisme penghubung $\phi_{ji}$ yang tidak lain adalah pemetaan inklusi dari $X_i$ ke $X_j$ untuk $i < j$.

Yang disebut
 \index{induktif!limit!ketat}%
 \index{limit!induktif ketat}%
 \index{ketat!limit induktif}%
\df{limit induktif ketat} dari barisan ruang semacam itu adalah gabungannya $L = \bigcup_{i=1}^\infty X_i$. Kita memberi $L$ topologi
konveks lokal terbesar yang membuat semua pemetaan inklusi $\psi_i\colon X_i \sto L$ kontinu. Inilah
 \index{induktif!limit!topologi}%
 \index{topologi!limit induktif}%
 \index{limit!induktif!topologi}%
\df{topologi limit induktif} pada~$L$.
\end{defn}

\begin{prop} Topologi limit induktif yang didefinisikan di atas ada.
\end{prop}

\begin{defn} Limit induktif ketat dari suatu barisan ruang Fr\'echet disebut
 \index{LF@ruang-$LF$}%
 \index{ruang!$LF$-}%
\df{ruang-$LF$}.
\end{defn}

\begin{exam} Ruang $\fml D(\Omega) = \fml C^\infty_c(\Omega)$ dari fungsi uji pada suatu himpunan terbuka $\Omega$ dalam ruang Euklides
adalah ruang-$LF$ (lihat notasi~\ref{mi_notn}).
\end{exam}

\begin{prop} Setiap ruang-$LF$ lengkap.
\end{prop}

\begin{proof} Lihat~\cite{Treves:1967}, Teorema 13.1.  \ns
\end{proof}

\begin{prop} Misalkan $T\colon L \sto Y$ suatu pemetaan linear dari ruang-$LF$ ke ruang konveks lokal. Andaikan bahwa $L = \underrightarrow{\lim}X_k$
dengan $(X_k)$ suatu barisan induktif ketat ruang Fr\'echet. Maka $T$ kontinu jika dan hanya jika pembatasannya $T\bigr|_{X_k}$
pada setiap $X_k$ kontinu.
\end{prop}

\begin{proof} Lihat~\cite{Treves:1967}, Proposisi~13.1, atau~\cite{Conway:1990}, Proposisi IV.5.7, atau~\cite{Grubb:2009}, Teorema B.18(c).  \ns
\end{proof}

\begin{cau} Fran\c{c}ois Treves, dalam bukunya yang ditulis dengan indah, \emph{Topological Vector Spaces, Distributions, and Kernels}, memperingatkan
para pembacanya tentang suatu kekeliruan yang sangat sulit dikenali: orang mudah sekali menganggap bahwa jika $M$ adalah subruang vektor tertutup dari ruang-$LF$
$L = \underrightarrow{\lim}X_k$, maka topologi yang diwarisi $M$ dari $L$ pasti sama dengan topologi limit induktif yang diperolehnya
dari barisan induktif ketat ruang Fr\'echet~$M \cap X_k$. Ternyata hal ini \emph{tidak} benar secara umum.
\end{cau}

\begin{prop} Misalkan $T\colon L \sto Y$ suatu pemetaan linear dari ruang-$LF$ ke ruang konveks lokal. Maka $T$ kontinu jika dan hanya jika
pemetaan tersebut kontinu sekuensial.
\end{prop}

















\section{Distribusi}
\begin{defn} Misalkan $\Omega$ suatu himpunan bagian terbuka tak kosong dari~$\R^n$. Suatu fungsi $f\colon \Omega \sto \R$ disebut
 \index{terintegralkan!secara lokal}%
 \index{secara lokal!terintegralkan}%
\df{terintegralkan secara lokal} jika $\int_K \abs{f}\,d\mu < \infty$ untuk setiap himpunan bagian kompak $K$ dari~$\Omega$. (Di sini $\mu$ adalah ukuran Lebesgue
berdimensi-$n$ yang biasa.) Keluarga semua fungsi tersebut pada $\Omega$ dinotasikan
 \index{L@$\locint\Omega$!fungsi yang terintegralkan secara lokal pada~$\Omega$}%
dengan~$\locint\Omega$.
\end{defn}

\begin{exam} Fungsi konstan $\vc 1$ pada garis real terintegralkan secara lokal, meskipun fungsi ini jelas tidak terintegralkan.
\end{exam}

Bahkan fungsi yang terintegralkan secara lokal tidak harus kontinu. Pada garis real, misalnya, fungsi karakteristik dari sebarang himpunan terukur Lebesgue
terintegralkan secara lokal. Bayangkan betapa memudahkannya jika kita dapat membicarakan secara bermakna \emph{turunan} dari sebarang fungsi demikian dan,
lebih baik lagi, diperbolehkan menurunkannya sesering yang kita inginkan. --- Selamat datang di dunia distribusi. Di sini
   \begin{enumerate}
       \item setiap fungsi yang terintegralkan secara lokal dapat dipandang sebagai suatu distribusi, dan
       \item setiap distribusi dapat diturunkan tak berhingga kali, dan lebih lanjut,
       \item semua aturan diferensiasi yang lazim dalam kalkulus dasar berlaku.
   \end{enumerate}

\begin{defn} Misalkan $\Omega$ suatu himpunan bagian terbuka tak kosong dari~$\R^n$. Suatu
 \index{distribusi}%
\df{distribusi} adalah fungsional linear kontinu pada ruang fungsi uji $\fml D(\Omega) = \fml C^\infty_c(\Omega)$ pada~$\Omega$.
Meskipun secara teknis fungsional demikian didefinisikan pada suatu ruang fungsi pada $\Omega$, kita akan mengikuti kebiasaan longgar yang lazim
dan menyebutnya sebagai \emph{distribusi pada~$\Omega$}. Himpunan semua distribusi pada $\Omega$, yaitu ruang dual dari $\fml D(\Omega)$,
 \index{D@$\fml D^*(\Omega)$ (ruang distribusi pada~$\Omega$)}%
akan dinotasikan dengan~$\fml D^*(\Omega)$.
\end{defn}

\begin{prop} Misalkan $\Omega$ suatu himpunan bagian terbuka dari~$\R^n$. Suatu fungsional linear $L$ pada ruang fungsi uji $\fml D(\Omega)$ pada $\Omega$ adalah
distribusi jika dan hanya jika $L(\phi_k) \sto 0$ setiap kali $(\phi_k)$ merupakan barisan fungsi uji sedemikian sehingga $\phi_k^{(\bs\alpha)} \sto 0$
secara seragam untuk setiap multi-indeks $\bs\alpha$ dan terdapat suatu himpunan kompak $K \subseteq \Omega$ yang memuat tumpuan setiap~$\phi_k$.
\end{prop}

\begin{proof} Lihat~\cite{Conway:1990}, proposisi IV.5.21 atau~\cite{Treves:1967}, proposisi 21.1(b).  \ns
\end{proof}

\begin{exam} Jika $\Omega$ suatu himpunan bagian terbuka tak kosong dari $\R^n$ dan $f$ suatu fungsi yang terintegralkan secara lokal pada $\Omega$, maka pemetaan
   \[ L_f\colon \fml D(\Omega) \sto \K \colon \phi \mapsto \int f\phi\,d\lambda \]
merupakan suatu distribusi. Distribusi demikian disebut
 \index{regular!distribusi}%
 \index{distribusi!regular}%
 \index{L@$L_f$ (distribusi regular yang bersesuaian dengan~$f$)}%
distribusi \df{regular}. Distribusi yang tidak regular disebut
 \index{singular!distribusi}%
 \index{distribusi!singular}%
distribusi \df{singular}.
\end{exam}

\emph{Catatan.} Misalkan $u$ suatu distribusi dan $\phi$ suatu fungsi uji pada $\Omega$. Notasi baku untuk $u(\phi)$ adalah
$\langle u,\phi \rangle$ dan $\langle \phi,u \rangle$. Kita juga menggunakan
 \index{<unoptop@$\tilde f$ (distribusi regular yang bersesuaian dengan~$f$)}%
notasi $\tilde f$ untuk $L_f$. Jadi, dalam contoh sebelumnya,

    \[ \langle \tilde f,\phi \rangle = \langle \phi,\tilde f \rangle
            = \tilde f(\phi) = \langle L_f,\phi \rangle
            = \langle \phi,L_f \rangle = L_f(\phi) \]
untuk $f \in \locint\Omega$.

\begin{exam} Jika $\mu$ suatu ukuran Borel regular pada himpunan bagian terbuka tak kosong $\Omega$ dari~$\R^n$ dan $f$ suatu fungsi yang terintegralkan secara lokal
terhadap ukuran tersebut pada $\Omega$, maka pemetaan
    \[ L_\mu\colon \fml D(\Omega) \sto \R\colon \phi \mapsto \int_\Omega f\phi \,d\mu \]
merupakan suatu distribusi pada $\Omega$.
\end{exam}

\begin{notn} Misalkan $a \in \R$. Definisikan $\mu_a$, yaitu
 \index{Dirac!ukuran}%
 \index{ukuran!Dirac}%
 \index{mua@$\mu_a$ (ukuran Dirac yang terkonsentrasi di~$a$)}%
\df{ukuran Dirac} yang terkonsentrasi di $a$, pada himpunan-himpunan Borel dari $\R$, dengan
   \[ \mu_a(B) = \begin{cases}  1, &\text{jika $a \in B$} \\
                                0, &\text{jika $a \notin B$}.
                 \end{cases} \]
Distribusi yang bersesuaian kita notasikan dengan $\delta_a$ dan kita sebut
 \index{Dirac!distribusi}%
 \index{distribusi!Dirac}%
 \index{delta@$\delta_a$, $\delta$ (distribusi Dirac)}%
\df{distribusi delta Dirac di $a$}. Jadi, $\delta_a = L_{\mu_a}$. Jika $a = 0$, kita menulis $\delta$ untuk $\delta_0$.
Secara historis, distribusi $\delta$ telah menyandang nama yang terkenal menyesatkan (meskipun secara teknis
tidak salah): \emph{fungsi delta Dirac}.
\end{notn}

\begin{exam} Jika $\delta_a$ adalah \emph{distribusi delta Dirac} di suatu titik $a \in \R$ dan $\phi$ suatu fungsi uji pada~$\R$,
maka $\delta_a(\phi) = \phi(a)$.
\end{exam}

\begin{notn} Misalkan $H$ menyatakan fungsi karakteristik interval $[0,\infty)$. Fungsi ini adalah
 \index{Heaviside!fungsi}%
 \index{fungsi!Heaviside}%
\df{fungsi Heaviside}. Distribusi yang bersesuaian, $\widetilde H = L_H$, adalah
 \index{Heaviside!distribusi}%
 \index{distribusi!Heaviside}%
 \index{H@$\wt H$ (distribusi Heaviside)}%
\df{distribusi Heaviside}.
\end{notn}

\begin{exam} Jika $\wt H$ adalah \emph{distribusi Heaviside} dan $\phi \in \fml D(\R)$, maka $\wt H(\phi) = \int_0^\infty \phi$.
\end{exam}

\begin{exam} Misalkan $f$ suatu fungsi yang dapat didiferensialkan pada $\R$ dan turunannya terintegralkan secara lokal. Maka
    \begin{equation}\label{deriv_reg_distr}
        L_{f\,'}(\phi) = -L_f(\phi\,')
    \end{equation}
untuk setiap $\phi \in \fml D(\R)$.
\end{exam}

\begin{proof}[\emph{Petunjuk untuk bukti}] Integrasi parsial. \ns
\end{proof}

Sekarang timbul sebuah pertanyaan penting: Bagaimana kita dapat mendefinisikan ``turunan'' suatu distribusi? Tentunya wajar jika kita menginginkan suatu
distribusi regular berperilaku hampir sama dengan fungsi yang melahirkannya. Artinya, kita ingin turunan $(L_f)'$ dari
suatu distribusi regular $L_f$ bersesuaian dengan distribusi regular yang timbul dari turunan fungsi~$f$. Dengan kata lain,
yang kita inginkan adalah $(L_f)' = L_{f\,'}$. Menurut contoh sebelumnya, ini menyarankan bahwa untuk fungsi $f$ yang dapat didiferensialkan
yang turunannya terintegralkan secara lokal, kita seharusnya mendefinisikan
   \begin{equation}\label{deriv_distr1}
        (L_f)'(\phi) = -L_f(\phi\,')
   \end{equation}
untuk setiap $\phi \in \fml D(\R)$. Pengamatan bahwa ruas kanan persamaan~\eqref{deriv_distr1} bermakna untuk
\emph{distribusi apa pun} memotivasi definisi berikut.

\begin{defn}\label{def_deriv_distr} Misalkan $u$ suatu distribusi pada himpunan bagian terbuka tak kosong $\Omega$ dari~$\R$. Kita mendefinisikan $u\,'$, yaitu
 \index{turunan!distribusi}%
 \index{distribusi!turunan}%
\df{turunan} dari $u$,
 \index{D@$Du$ (turunan suatu distribusi~$u$)}%
 \index{<unopur@$u'$ (turunan suatu distribusi~$u$)}%
dengan
    \[ u\,'(\phi) := -u(\phi\,') \]
untuk setiap fungsi uji $\phi$ pada~$\Omega$. Kita juga menggunakan notasi $Du$ untuk turunan dari~$u$.
\end{defn}

\begin{exam} \emph{Distribusi Dirac} $\delta$ adalah turunan dari \emph{distribusi Heaviside}~$\widetilde H$.
\end{exam}

\begin{exer} Misalkan $S$ suatu fungsi signum (atau tanda) pada $\R$; artinya, misalkan $S = 2H - 1$, dengan $H$ fungsi Heaviside.
Misalkan pula $g$ suatu fungsi pada $\R$ sedemikian sehingga $g\colon x \mapsto -x - 3$ untuk $x < 0$ dan $g\colon x \mapsto x + 5$ untuk $x \ge 0$.
   \begin{enumerate}[label=\rm{(\alph*)}]
     \item Carilah konstanta $\alpha$ dan $\beta$ sedemikian sehingga ${\tilde g}' = \alpha \widetilde S + \beta \delta$.
     \item Carilah suatu fungsi $a$ pada $\R$ sedemikian sehingga ${\tilde a}' = \wt S$.
   \end{enumerate}
\end{exer}

\begin{exer} Misalkan $f(x) = \left\{
                           \begin{array}{ll}
                             1 - \abs x, & \hbox{jika $\abs x \le 1$;} \\
                                 0,      & \hbox{selain itu.}
                           \end{array}
                         \right.$ Dengan menghitung kedua ruas persamaan~\eqref{deriv_reg_distr} secara terpisah, tunjukkan bahwa persamaan itu
benar secara klasik untuk setiap fungsi uji yang tumpuannya termuat dalam interval~$[-1,1]$.
\end{exer}

\begin{exer} Misalkan $f(x) = \sbsb\chi{(0,1)}$ untuk $x \in~\R$ dan $\phi \in \fml D\bigl((-1,1)\bigr)$. Bandingkan nilai $L_{f\,'}(\phi)$ dan
$\bigl(L_f\bigr)'(\phi)$.
\end{exer}

\begin{exer} Misalkan $f(x) = \left\{
                           \begin{array}{ll}
                             x, & \hbox{jika $0 \le x \le 1$;} \\
                             x + 2, & \hbox{jika $1 < x \le 2$;} \\
                             0, & \hbox{selain itu.}
                           \end{array}
                         \right.$ Bandingkan nilai $L_{f\,'}(\phi)$ dan $\bigl(L_f\bigr)'(\phi)$ untuk suatu fungsi uji~$\phi$.
\end{exer}

\begin{exer} Misalkan $h(x) = \left\{
                           \begin{array}{ll}
                             2x, & \hbox{jika $x < -1$;} \\
                             1, & \hbox{jika $-1 \le x < 1$;} \\
                             0, & \hbox{jika $x \ge 1$.}
                           \end{array}
                         \right.$ Carilah suatu fungsi $f$ yang terintegralkan secara lokal sedemikian sehingga ${\tilde f}\,' = \tilde h + 3\delta_1 - \delta_{-1}$.
\end{exer}

Definisi~\ref{def_deriv_distr} dapat diperluas ke turunan parsial dari distribusi tanpa kesulitan.

\begin{defn} Misalkan $u$ suatu distribusi pada suatu himpunan bagian terbuka tak kosong $\Omega$ dari $\R^n$ untuk suatu $n \ge 2$ dan $\bs\alpha$
suatu multi-indeks. Kita mendefinisikan
 \index{operator diferensial!yang bekerja pada distribusi}%
 \index{D@$D^{\bs\alpha}$ (notasi multi-indeks untuk diferensiasi)}%
\df{operator diferensial} $D^{\bs\alpha} = \bigl(\frac{\partial}{\partial x_1}\bigr)^{\alpha_1} \dots
\bigl(\frac{\partial}{\partial x_n}\bigr)^{\alpha_n}$ berorde $\abs{\bs\alpha}$ yang bekerja pada $u$ dengan
   \[ D^{\bs\alpha}u(\phi) := (-1)^{\abs{\bs\alpha}}u\bigl(\phi^{(\bs\alpha)}) \]
untuk setiap fungsi uji $\phi$ pada~$\Omega$.
 \index{<unopura@$u^{\bs\alpha}$ (notasi multi-indeks untuk diferensiasi)}%
Notasi alternatif untuk $D^{\bs\alpha}u$ adalah $u^{(\bs\alpha)}$.
\end{defn}

\begin{exer} Suatu \df{dipol} (dengan momen listrik 1 di titik asal pada $\R$) dapat dipandang sebagai ``limit'' ketika $\epsilon \sto 0^+$
 \index{dipol}%
dari suatu sistem $T_\epsilon$ yang terdiri atas dua muatan $-\frac1\epsilon$ dan $\frac1\epsilon$ yang masing-masing ditempatkan di $0$ dan $\epsilon$. Tunjukkan
bagaimana hal ini dapat mengarahkan kita untuk mendefinisikan dipol secara matematis sebagai $-\delta\,'$. \emph{Petunjuk.} Pandanglah $T_\epsilon$ sebagai
distribusi $\frac1\epsilon \delta_\epsilon - \frac1\epsilon \delta$.
\end{exer}
 
\begin{defn} Misalkan $X$ suatu ruang konveks lokal Hausdorff dan $X^*$ ruang dualnya, yaitu ruang semua
fungsional linear kontinu pada~$X$. Untuk setiap $f \in X^*$ definisikan
  \[ p_f\colon X \sto \K\colon x \mapsto \abs{f(x)}. \]
Jelas bahwa setiap $p_f$ merupakan seminorma pada~$X$. Seperti dalam kasus ruang linear bernorma, kita mendefinisikan
 \index{lemah!topologi!pada ruang konveks lokal}%
 \index{sigma@$\sigma(X,X^*)$ (topologi lemah pada~$X$)}%
\df{topologi lemah} pada $X$, yang dinotasikan dengan $\sigma(X,X^*)$, sebagai topologi lemah yang dibangkitkan oleh keluarga seminorma
$\{p_f\colon f \in X^*\}$ pada~$X$. Demikian pula, untuk setiap $x \in X$ definisikan
   \[ p_x\colon X^* \sto \K \colon f \mapsto \abs{f(x)}. \]
Sekali lagi, setiap $p_x$ merupakan seminorma pada~$X^*$, dan kita mendefinisikan
 \index{w*@$w^*$-topologi (topologi lemah-bintang)}%
 \index{sigma@$\sigma(X^*,X)$ (topologi-$w^*$ pada~$X^*$)}%
\df{topologi-$w^*$} pada $X^*$, yang dinotasikan dengan $\sigma(X^*,X)$, sebagai topologi lemah yang dibangkitkan oleh keluarga seminorma
$\{p_x\colon x \in X\}$ pada~$X^*$.
\end{defn}

\begin{prop} Jika ruang distribusi $\fml D^*(\Omega)$ pada subhimpunan terbuka $\Omega$ dari $\R^n$ diberi topologi-$w^*$,
maka suatu jaring $(u_\lambda)$ distribusi pada $\Omega$ konvergen ke distribusi $v$ pada $\Omega$ jika dan hanya jika
$\left\langle u_\lambda, \phi\right\rangle \sto \left\langle v, \phi\right\rangle$ untuk setiap fungsi uji $\phi$ pada $\Omega$.
\end{prop}

\begin{prop} Andaikan suatu jaring $(u_\lambda)$ distribusi konvergen ke distribusi~$v$. Maka ${u_\lambda}^{(\alpha)} \sto v^{(\alpha)}$
untuk setiap multi-indeks $\alpha$.
\end{prop}

\begin{proof} Lihat \cite{Rudin:1991}, Teorema 6.17; atau \cite{Yosida:1965}, Bab II, Seksi 3, Teorema; atau \cite{Grubb:2009}, Teorema 3.9;
atau~\cite{HirschL:1999}, Bab 8, Proposisi 2.4.
\end{proof}

\begin{prop} Jika $g$, $f_k \in \locint\Omega$ untuk setiap $k$ (dengan $\open{\Omega}{\R^n}$) dan $f_k \sto g$ secara seragam pada
setiap subhimpunan kompak dari $\Omega$, maka $\tilde f_k \sto \tilde g$.
\end{prop}

\begin{prop} Jika $f \in \fml C^\infty(\R^n)$, maka $\left(L_f\right)^{(\bs\alpha)} = L_{f^{(\bs\alpha)}}$ untuk setiap multi-indeks $\bs\alpha$.
\end{prop}

\begin{exer} Jika $\delta$ bukan suatu fungsi pada $\R$, lalu apa yang dimaksud orang ketika mereka menulis
   \[ \lim_{k \sto \infty} \frac{k}{\pi(1+k^2x^2)} = \delta(x)\quad? \]
Tunjukkan bahwa, jika ditafsirkan dengan tepat (sebagai suatu pernyataan tentang distribusi), ungkapan itu bahkan benar. \emph{Petunjuk.}
Jika $s_k(x) = k\,\pi^{-1}(1+k^2x^2)^{-1}$ dan $\phi$ suatu fungsi uji pada $\R$, maka
$\int_{-\infty}^\infty s_k\phi = \int_{-\infty}^\infty s_k\phi(0) + \int_{-\infty}^\infty s_k\eta$ dengan $\eta(x) = \phi(x) - \phi(0)$.
Tunjukkan bahwa $\int_{-\infty}^\infty s_k\eta \sto 0$ ketika $k \sto \infty$.
\end{exer}

\begin{exer} Misalkan $f_n(x) = \sin nx$ untuk $n \in \N$ dan $x \in \R$. Barisan $(f_n)$ tidak konvergen secara klasik;
tetapi barisan itu konvergen secara distribusional (ke $0$). Nyatakan pernyataan ini dengan tepat, lalu buktikan. \emph{Petunjuk.} Carilah antiturunannya.
\end{exer}

\begin{prop} Jika $\phi$ suatu fungsi mulus pada $\R$ dengan tumpuan kompak, maka
    \[ \lim_{n \sto \infty}\frac2\pi \int_{-\infty}^\infty \frac{n^3x}{\bigl(1 + n^2x^2\bigr)^2}\,\, \phi(x)\,dx = \phi'(0)\,.\]
\end{prop}

\begin{exer} Berikan makna pada ungkapan
   \[ 2\pi\sum_{n = -\infty}^\infty \delta(x-2\pi n) =  1 +  2\sum_{n = 1}^\infty \cos nx\,, \]
dan tunjukkan bahwa, jika ditafsirkan dengan tepat, ungkapan itu benar. \emph{Petunjuk.} Siapa pun yang cukup bejat untuk percaya bahwa $\delta(x)$ \emph{berarti}
sesuatu kemungkinan akan menafsirkan $\delta(x-a)$ sebagai hal yang sama dengan $\delta_a(x)$. Anda dapat memulai proses
merasionalisasikannya dengan menghitung deret Fourier fungsi $g$ yang didefinisikan oleh $g(x) = \frac1{4\pi}x^2 - \frac12 x$ pada $[0,2\pi]$
dan kemudian diperluas secara periodik ke $\R$. Anda mungkin perlu mencari teorema tentang konvergensi titik demi titik dan seragam dari deret Fourier
serta tentang pendiferensialan suku demi suku deret semacam itu.
\end{exer}

\begin{defn} Misalkan $\Omega$ suatu subhimpunan terbuka dari $\R^n$, $u \in \fml D^*(\Omega)$, dan $f \in \fml C^\infty(\Omega)$.
 \index{perkalian!distribusi dengan fungsi mulus}%
 \index{<binopconcat@$fu$ (hasil kali fungsi mulus dengan distribusi)}%
Definisikan
   \[ (fu)(\phi) := u(f\phi) \]
untuk semua fungsi uji $\phi$ pada~$\Omega$. Secara ekuivalen, kita dapat menulis $\langle fu,\phi \rangle = \langle u,f\phi \rangle$.
\end{defn}

\begin{prop} Fungsional $fu$ dalam definisi sebelumnya merupakan distribusi pada~$\Omega$.
\end{prop}

\begin{proof} Lihat~\cite{Rudin:1991}, halaman 159, Seksi 6.15.    \ns
\end{proof}






















\section{Konvolusi}

Seksi ini dan Seksi~\ref{section_Fourier_transform} memberikan pengantar yang \emph{sangat} singkat tentang konvolusi
dan transformasi Fourier distribusi. Untuk banyak pembuktian, serta uraian yang lebih terperinci, pembaca dirujuk
ke buku \emph{Functional Analysis} karya Walter Rudin yang elegan~\cite{Rudin:1991}.

Untuk menghindari munculnya terlalu banyak faktor berbentuk $\sqrt{2\pi}$ beserta kebalikannya dalam rumus-rumus, selama
sisa bab ini kita akan mengintegralkan terhadap
 \index{ukuran Lebesgue ternormalisasi}%
 \index{m@$m$ (ukuran Lebesgue ternormalisasi)}%
\df{ukuran Lebesgue ternormalisasi} $m$ pada~$\R$. Ukuran ini didefinisikan oleh
   \[ m(A) := \frac1{\sqrt{2\pi}}\lambda(A) \]
untuk setiap subhimpunan terukur Lebesgue $A$ dari $\R$ (dengan $\lambda$ ukuran Lebesgue biasa pada~$\R$).

Kita meninjau kembali beberapa fakta dasar dari analisis real mengenai konvolusi fungsi bernilai skalar. Ingat bahwa suatu fungsi bernilai kompleks atau
real diperluas pada $\R$ disebut
 \index{terintegralkan Lebesgue}%
 \index{terintegralkan}%
\emph{terintegralkan Lebesgue} jika $\int_\R\abs f\, dm < \infty$
(dengan $m$ ukuran Lebesgue ternormalisasi pada~$\R$).
 \index{L@$\lfs 1(\R)$!ruang fungsi terintegralkan Lebesgue}%
 \index{L@$\lfs 1(\R)$!sebagai ruang Banach}%
 \index{Banach!ruang!$\lfs 1(\R)$ sebagai}%
Dengan $\lfs 1(\R)$ kita nyatakan ruang Banach semua kelas ekuivalensi fungsi terintegralkan Lebesgue pada~$\R$; dua fungsi \emph{ekuivalen}
jika ukuran Lebesgue himpunan tempat keduanya berbeda adalah nol. Norma pada ruang Banach ini
   \index{<unopn@$\norm{\hphantom{x}}_1$ (norma-$1$)}%
diberikan oleh
   \[ {\norm f}_1 := \int_\R \abs f\,dm. \]

\begin{prop} Misalkan $f$ dan $g$ fungsi-fungsi terintegralkan Lebesgue pada~$\R$. Maka fungsi
   \[ y \mapsto f(x - y)g(y) \]
termasuk dalam $\lfs 1(\R)$ untuk hampir setiap~$x$. Definisikan suatu
 \index{<binopast@$f \ast g$ (konvolusi dua fungsi)}%
fungsi $f \ast g$ dengan
   \[ (f \ast g)(x) := \int_\R f(x - y)g(y)\,dm(y) \]
pada setiap titik tempat ruas kanan terdefinisi. Maka fungsi $f \ast g$ termasuk dalam $\lfs 1(\R)$.
\end{prop}

\begin{proof} Lihat~\cite{HewittS:1965}, Teorema 21.31.  \ns
\end{proof}

\begin{defn} Fungsi $f \ast g$ yang didefinisikan di atas adalah
 \index{konvolusi!fungsi-fungsi dalam $\lfs 1(\R)$}%
\df{konvolusi} dari $f$ dan~$g$.
\end{defn}

\begin{prop} Ruang Banach $\lfs 1(\R)$ dengan konvolusi merupakan aljabar Banach komutatif. Namun, aljabar ini \emph{tidak} beridentitas.
\end{prop}

\begin{proof} Lihat~\cite{HewittS:1965}, Teorema 21.34 dan 21.35.  \ns
\end{proof}

\begin{defn}\label{defn_Fourier_transform} Untuk setiap $f$ dalam $\lfs 1(\R)$, definisikan suatu fungsi $\hat f$ dengan
    \[ \hat f(x) = \int_{-\infty}^\infty f(t) e^{-itx}\,dm(t). \]
Fungsi $\hat f$ adalah
 \index{Fourier!transformasi}%
 \index{transformasi!Fourier}%
 \index{F@$\mathfrak F(f)$ (transformasi Fourier dari~$f$)}%
 \index{<unoptop@$\hat f$ (transformasi Fourier dari~$f$)}%
\df{transformasi Fourier} dari~$f$. Kadang-kadang kita menulis $\mathfrak Ff$ untuk~$\hat f$.
\end{defn}

\begin{thm}[Lemma Riemann--Lebesgue] Jika $f \in \lfs 1(\R)$,
 \index{lemma Riemann--Lebesgue}%
maka $\hat f \in \fml C_0(\R)$.
\end{thm}

\begin{proof} Lihat~\cite{HewittS:1965}, 21.39.   \ns
\end{proof}

\begin{prop} Transformasi Fourier
   \[ \mathfrak F\colon \lfs 1(\R) \sto \fml C_0(\R)\colon f \mapsto \hat f \]
merupakan homomorfisme aljabar Banach. Kedua aljabar itu tidak beridentitas.
\end{prop}

Untuk penerapan, segi paling berguna dari proposisi sebelumnya ialah bahwa transformasi Fourier
mengubah konvolusi dalam $\lfs 1(\R)$ menjadi
perkalian titik demi titik dalam~$\fml C_0(\R)$; yakni, $\wh{f\ast g} =\hat f\hat g$.

Berikutnya kita mendefinisikan apa yang dimaksud dengan mengambil konvolusi suatu distribusi dan suatu fungsi uji.

\begin{notn} Untuk suatu titik $x \in \R$ dan suatu fungsi bernilai skalar $\phi$ pada $\R$
 \index{phi@$\phi_x$}%
definisikan
   \[ \phi_x\colon \R \sto \K\colon y \mapsto \phi(x - y). \]
\end{notn}

\begin{prop} Jika $\phi$ dan $\psi$ merupakan fungsi-fungsi bernilai skalar pada $\R$ dan $x \in \R$, maka
   \[ (\phi + \psi)_x = \phi_x + \psi_x\,. \]
\end{prop}

\begin{prop}\label{prop_convol_distr1} Untuk $u \in \fml D^*(\R)$ dan $\phi \in \fml D(\R)$,
 \index{<binopast@$u \ast \phi$ (konvolusi distribusi dan fungsi uji)}%
tetapkan
   \[ (u \ast \phi)(x) := \langle u,\phi_x \rangle \]
untuk semua $x \in \R$. Maka $u \ast \phi \in \fml C^\infty(\R)$.
\end{prop}

\begin{proof} Lihat~\cite{Rudin:1991}, Teorema 6.30(b).  \ns
\end{proof}

\begin{defn} Ketika $u$ dan $\phi$ seperti dalam proposisi sebelumnya, fungsi $u \ast \phi$ adalah
 \index{konvolusi!distribusi dan fungsi uji}%
\df{konvolusi} dari $u$ dan~$\phi$.
\end{defn}

\begin{prop} Jika $u$ dan $v$ merupakan distribusi pada $\R$ dan $\phi$ suatu fungsi uji pada $\R$, maka
   \[ (u + v) \ast \phi = (u \ast \phi) + (v \ast \phi)\,. \]
\end{prop}

\begin{prop} Jika $u$ suatu distribusi pada $\R$ dan $\phi$ serta $\psi$ fungsi-fungsi uji pada $\R$, maka
   \[ u \ast (\phi + \psi) = (u \ast \phi) + (u \ast \psi)\,. \]
\end{prop}

\begin{exam} Jika $H$ fungsi Heaviside dan $\phi$ suatu fungsi uji pada~$\R$, maka
   \[ (\wt H \ast \phi)(x) = \int_{-\infty}^x \phi(t)\,dt. \]
\end{exam}

\begin{prop}\label{prop_convol_distr2} Jika $u \in \fml D^*(\R)$ dan $\phi \in \fml D(\R)$, maka
   \[ D^{\bs\alpha}(u \ast \phi) = (D^{\bs\alpha}u) \ast \phi = u \ast \bigl(D^{\bs\alpha}(\phi)\bigr). \]
\end{prop}

\begin{proof} Lihat~\cite{Rudin:1991}, Teorema 6.30(b).  \ns
\end{proof}

\begin{prop}\label{prop_convol_distr3} Jika $u \in \fml D^*(\R)$ dan $\phi$, $\psi \in \fml D(\R)$, maka
   \[ u \ast (\phi \ast \psi) = (u \ast \phi) \ast \psi. \]
\end{prop}

\begin{proof} Lihat~\cite{Rudin:1991}, Teorema 6.30(c).  \ns
\end{proof}

\begin{prop} Misalkan $u$, $v \in \fml D^*(\R)$. Jika
   \[ u \ast \phi = v \ast \phi \]
untuk semua $\phi \in \fml D(\R)$, maka $ u = v$.
\end{prop}

\begin{defn} Misalkan $\Omega$ suatu subhimpunan terbuka dari $\R^n$ dan $u \in \fml D^*(\R^n)$. Kita mengatakan bahwa ``$u = 0$ pada $\Omega$'' jika $u(\phi) = 0$
untuk setiap $\phi \in \fml D(\Omega)$. Dengan semangat yang sama, kita mengatakan bahwa ``dua distribusi $u$ dan $v$ sama pada $\Omega$'' jika $u - v = 0$
pada~$\Omega$. Suatu titik $x \in \R^n$ termasuk dalam
 \index{tumpuan!distribusi}%
 \index{supp@$\supp u$ (tumpuan distribusi~$u$)}%
\df{tumpuan} dari $u$ jika \emph{tidak ada} lingkungan terbuka bagi $x$ yang padanya $u = 0$. Nyatakan tumpuan $u$ dengan~$\supp u$.
\end{defn}

\begin{exam} Tumpuan distribusi delta Dirac $\delta$ adalah~$\{0\}$.
\end{exam}

\begin{exam} Tumpuan distribusi Heaviside $\wt H$ adalah~$[0,\infty)$.
\end{exam}

\begin{rem} Andaikan suatu distribusi $u$ pada $\R^n$ memiliki tumpuan kompak. Maka dapat ditunjukkan bahwa $u$ memiliki perluasan unik menjadi
fungsional linear kontinu pada~$\fml C^\infty(\R^n)$, yakni keluarga semua fungsi mulus pada~$\R^n$. Dalam kondisi ini, kesimpulan
Proposisi~\ref{prop_convol_distr1} tetap benar dan definisi berikut berlaku. Jika $u \in \fml D^*(\R^n)$ memiliki tumpuan kompak
dan $\phi \in \fml C^\infty(\R^n)$, maka kesimpulan Proposisi~\ref{prop_convol_distr2} tetap benar. Selain itu, jika $\psi$
merupakan fungsi uji pada $\R^n$, maka $u \ast \psi$ merupakan fungsi uji pada $\R^n$ dan kesimpulan~\ref{prop_convol_distr3} tetap berlaku.
Untuk pembahasan dan verifikasi menyeluruh atas pernyataan-pernyataan ini, lihat~\cite{Rudin:1991}, Teorema 6.24 dan~6.35.
\end{rem}

\begin{prop} Jika $u$ dan $v$ merupakan distribusi pada $\R^n$ dan paling sedikit salah satunya memiliki tumpuan kompak, maka terdapat distribusi tunggal
$u \ast v$ sedemikian sehingga
   \[ (u \ast v) \ast \phi = u \ast (v \ast \phi) \]
untuk semua fungsi uji $\phi$ pada~$\R^n$.
\end{prop}

\begin{proof} Lihat \cite{Rudin:1991}, Definisi 6.36\emph{ dan seterusnya}.   \ns
\end{proof}

\begin{defn} Dalam proposisi sebelumnya, distribusi $u \ast v$ adalah
 \index{konvolusi!dua distribusi}%
 \index{<binopast@$u \ast v$ (konvolusi dua distribusi)}%
\df{konvolusi} dari $u$ dan~$v$.
\end{defn}

\begin{exam} Jika $\delta$ distribusi delta Dirac, maka
   \[ \tilde 1 \ast \delta\,' = \tilde 0\,. \]
\end{exam}

\begin{exam} Jika $\delta$ distribusi delta Dirac dan $H$ fungsi Heaviside, maka
   \[ \delta\,' \ast \wt H = \delta\,. \]
\end{exam}

\begin{prop} Jika $u$ dan $v$ merupakan distribusi pada $\R^n$ dan paling sedikit salah satunya memiliki tumpuan kompak, maka $u \ast v = v \ast u$.
\end{prop}

\begin{proof} Lihat \cite{Rudin:1991}, Teorema 6.37(a).  \ns
\end{proof}

\begin{exer} Buktikan proposisi sebelumnya dengan asumsi bahwa $u$ dan $v$ \emph{keduanya} memiliki tumpuan kompak. \emph{Petunjuk.}
Gunakan $(u \ast v) \ast (\phi \ast \psi)$, dengan $\phi$ dan $\psi$ fungsi-fungsi uji, sambil mengingat bahwa konvolusi fungsi
bersifat komutatif.
\end{exer}

\begin{prop} Jika $u$ dan $v$ merupakan distribusi pada $\R^n$ dan paling sedikit salah satunya memiliki tumpuan kompak, maka
  \[ \supp(u \ast v) \subseteq \supp u + \supp v. \]
\end{prop}

\begin{proof} Lihat \cite{Rudin:1991}, Teorema 6.37(b).  \ns
\end{proof}

\begin{prop}\label{prop_assoc_convol_distr} Jika $u$, $v$, dan $w$ merupakan distribusi pada $\R^n$ dan paling sedikit dua di antaranya memiliki
tumpuan kompak, maka
    \[ u \ast (v \ast w) = (u \ast v) \ast w. \]
\end{prop}

\begin{exam} Ungkapan $\tilde 1 \ast \delta\,' \ast \wt H$ bersifat ambigu. Jadi, hipotesis mengenai tumpuan
distribusi-distribusi dalam proposisi sebelumnya tidak dapat dihilangkan.
\end{exam}

\begin{prop} Terhadap penjumlahan, perkalian skalar, dan konvolusi, keluarga semua distribusi pada $\R$ yang bertumpuan kompak
merupakan aljabar komutatif beridentitas.
\end{prop}













\section{Solusi Distribusional untuk Persamaan Diferensial Biasa}

\begin{notn} Misalkan $\Omega$ suatu subhimpunan terbuka tak kosong dari $\R$, $n \in \N$, dan $a_0$, $a_1$, \dots $a_n \in \fml C^\infty(\Omega)$.
Kita meninjau operator diferensial (biasa)
 \index{diferensial!operator}%
 \index{operator!diferensial}%
berikut:
   \[ L = \sum_{k=0}^n a_k(x)\frac{d^k}{dx^k} \]
beserta persamaan diferensial (biasa)
 \index{diferensial!persamaan}%
 \index{persamaan!diferensial}%
yang terkait dengannya,
   \begin{equation}\label{eqn_ODE1}
        Lu = v
   \end{equation}
dengan $u$ dan $v$ distribusi-distribusi pada~$\Omega$. Sehubungan dengan Persamaan~\eqref{eqn_ODE1}, kita juga dapat meninjau dua variasi:
persamaan
   \begin{equation}\label{eqn_ODE2}
        \langle Lu, \phi\rangle = \langle v, \phi\rangle
   \end{equation}
dengan $\phi$ suatu fungsi uji, dan persamaan
   \begin{equation}\label{eqn_ODE3}
        \langle u, L^*\phi\rangle = \langle v, \phi\rangle
   \end{equation}
dengan $\phi$ suatu fungsi uji dan $L^* := \displaystyle\sum_{k=0}^n (-1)^k \frac{d^k(a_k\phi)}{dx^k}$ merupakan
 \index{adjoin!formal}%
 \index{adjoin formal}%
\df{adjoin formal} dari~$L$.
\end{notn}

\begin{defn} Andaikan bahwa dalam persamaan-persamaan di atas distribusi $v$ diberikan. Jika terdapat distribusi regular $u$ yang
memenuhi~\eqref{eqn_ODE1}, kita mengatakan bahwa $u$ merupakan solusi
 \index{solusi klasik}%
 \index{solusi!klasik}%
\df{klasik} bagi persamaan diferensial itu. Jika $u$ merupakan distribusi regular yang bersesuaian dengan suatu fungsi yang
\emph{tidak} memenuhi~\eqref{eqn_ODE1} dalam pengertian klasik, tetapi \emph{memenuhi}~\eqref{eqn_ODE2} untuk setiap fungsi uji~$\phi$,
maka kita mengatakan bahwa $u$ merupakan solusi
 \index{lemah!solusi}%
 \index{solusi!lemah}%
\df{lemah} bagi persamaan diferensial itu. Selanjutnya, jika $u$ merupakan distribusi singular yang memenuhi~\eqref{eqn_ODE3} untuk setiap fungsi
uji~$\phi$, kita mengatakan bahwa $u$ merupakan solusi
 \index{solusi distribusional}%
 \index{solusi!distribusional}%
\df{distribusional} bagi persamaan diferensial itu. Keluarga solusi
 \index{solusi diperumum}%
 \index{solusi!diperumum}%
\df{diperumum} mencakup semua solusi klasik, lemah, dan distribusional.
\end{defn}

\begin{exam} Solusi diperumum bagi persamaan $\dfrac{du}{dx} = 0$ pada~$\R$ hanyalah distribusi konstan (yakni,
distribusi regular yang berasal dari fungsi konstan). Jadi, setiap solusi diperumum merupakan solusi klasik.
\end{exam}

\begin{proof}[\emph{Petunjuk untuk bukti}] Andaikan $u$ suatu solusi diperumum bagi persamaan itu. Pilih suatu fungsi uji $\sbsb\phi 0$ sedemikian
sehingga $\int_\R \sbsb\phi 0 = 1$. Misalkan $\phi$ suatu fungsi uji sebarang pada~$\R$. Definisikan
   \[ \psi(x) = \phi(x) - \sbsb\phi 0(x)\int_\R \phi \]
untuk semua $x \in \R$. Amati bahwa $\langle u,\psi \rangle = 0$, lalu hitung $\langle u,\phi \rangle$.  \ns
\end{proof}

\begin{exam} Distribusi Heaviside adalah solusi lemah persamaan diferensial $x\dfrac{du}{dx}~=~0$ pada~$\R$. Tentu saja,
distribusi-distribusi konstan merupakan solusi klasik.
\end{exam}

\begin{exam} Distribusi Dirac $\delta$ merupakan solusi distribusional bagi persamaan diferensial $x^2\dfrac{du}{dx} = 0$ pada~$\R$.
(Distribusi Heaviside merupakan solusi lemah; dan distribusi-distribusi konstan merupakan solusi klasik.)
\end{exam}

\begin{exam} Untuk $p = 0,1,2,\dots$, turunan ke-$p$, $\dfrac{d^p\delta}{dx^p}$, dari distribusi delta Dirac merupakan solusi
bagi persamaan diferensial $x^{p+2}\dfrac{du}{dx} = 0$ pada~$\R$.
\end{exam}

\begin{exer} Carilah suatu solusi distribusional bagi persamaan diferensial $x \dfrac{du}{dx} + u = 0$ pada~$\R$.
\end{exer}










\section{Transformasi Fourier}\label{section_Fourier_transform}

Dalam~\ref{defn_Fourier_transform}, kita mendefinisikan transformasi Fourier suatu fungsi $\lfs 1\R$. Cara ekuivalen untuk menyatakan
definisi ini adalah melalui konvolusi:
  \[ (\mathfrak F f)(x) = (f \ast \varepsilon^x)(0) \]
untuk $x \in \R$, dengan $\varepsilon^x$ fungsi $t \mapsto e^{itx}$.

\begin{defn} Kita ingin memperluas transformasi yang berguna ini ke distribusi. Namun, ternyata ruang $\fml D^*(\R)$ terlampau besar
untuk mendukung definisi yang wajar baginya. Dengan memulai dari himpunan ``fungsi uji'' yang lebih besar, yakni fungsi-fungsi dalam ruang Schwartz
$\fml S(\R)$ (yang memuat $\fml D(\R)$ sebagai subruang rapat), kita memperoleh ruang fungsional linear kontinu yang lebih kecil; pada ruang ini, untungnya, kita
dapat mendefinisikan suatu \emph{transformasi Fourier} secara alami.
\end{defn}

Sebelum mendefinisikan transformasi Fourier pada distribusi, mari kita ingat kembali dua fakta penting dari analisis real mengenai
transformasi Fourier yang bekerja pada fungsi mulus terintegralkan.

\begin{prop}[Rumus Inversi Fourier] Jika $f \in \fml S(\R)$,
 \index{Fourier!rumus inversi}%
maka
   \[ f(x) = \int_\R \hat f \varepsilon^x\,dm\,. \]
\end{prop}

\begin{prop}\label{prop_FT_isomor1} Transformasi Fourier
   \[ \mathfrak F\colon \fml S(\R) \sto \fml S(\R)\colon f \mapsto \hat f \]
merupakan isomorfisme ruang Fr\'echet (yakni, pemetaan linear kontinu bijektif dengan invers kontinu), dan untuk semua
$f \in \fml S(\R)$ serta $x \in \R$, kita mempunyai $\mathfrak F^2f(x) = f(-x)$, sehingga $\mathfrak F^4f = f$.
\end{prop}

\begin{proof} Pembuktian kedua hasil sebelumnya, bersama banyak informasi lain tentang transformasi yang menarik ini
dan beragam penerapannya pada persamaan diferensial biasa maupun parsial, dapat ditemukan di berbagai sumber.
Beberapa yang segera terlintas adalah~\cite{Cerda:2010}, Bab~7; \cite{Horvath:1966}, Bab~4, Seksi~11;
\cite{McDonaldW:1999}, Seksi~11.3; \cite{Rudin:1991}, Bab~7; \cite{Treves:1967}, Bab~25; dan \cite{Yosida:1965},
Bab~VI.   \ns
\end{proof}

\begin{defn} Fungsi inklusi $\iota\colon \fml D(\R) \sto \fml S(\R)$ bersifat kontinu. Maka pemetaan antara ruang-ruang dualnya
   \[ u\colon \fml S^*(\R) \sto \fml D^*(\R)\colon \tau \mapsto u_\tau = \tau \circ \iota \]
merupakan isomorfisme ruang vektor dari $\fml S^*(\R)$ ke suatu ruang distribusi pada~$\R$, yang kita sebut ruang
 \index{distribusi tempered}%
\df{distribusi tempered} (sebagian penulis memilih frasa
 \index{distribusi!temperate}%
\df{distribusi temperate}). Lazimnya, fungsional linear $\tau$ dan $u_\tau$ di atas diidentifikasi. Jadi, $\fml S^*(\R)$
sendiri dipandang sebagai ruang distribusi tempered.
\end{defn}

Jika Anda kesulitan memverifikasi salah satu pernyataan dalam definisi sebelumnya, lihat Teorema 7.10\emph{ dan seterusnya} dalam~\cite{Rudin:1991}.

\begin{exam} Setiap distribusi dengan tumpuan kompak merupakan distribusi tempered.
\end{exam}

\begin{defn} Definisikan
 \index{Fourier!transformasi}%
 \index{transformasi!Fourier}%
 \index{F@$\mathfrak F(u)$ (transformasi Fourier dari~$u$)}%
 \index{<unoptop@$\hat u$ (transformasi Fourier dari~$u$)}%
\df{transformasi Fourier} $\hat u$ dari suatu distribusi tempered $u \in \fml S^*(\R)$ dengan
   \[ \hat u(\phi) = u(\hat\phi) \]
untuk semua $\phi \in \fml S(\R)$. Notasi $\mathfrak Fu$ merupakan alternatif bagi~$\hat u$.
\end{defn}

\begin{prop} Jika $f \in \lfs1\R$, maka kita dapat memperlakukan $\tilde f$ sebagai distribusi tempered dan, dalam hal ini,
   \[ \hat{\tilde f} = \tilde{\hat f}\,. \]
\end{prop}

\begin{exam} Transformasi Fourier fungsi delta Dirac diberikan oleh
   \[ \hat\delta = \tilde{\vc 1}\,. \]
\end{exam}

\begin{exam} Transformasi Fourier distribusi regular $\tilde{\vc 1}$ diberikan oleh
   \[ \hat{\tilde{\vc 1}} = \delta\,. \]
\end{exam}

Proposisi~\ref{prop_FT_isomor1} memiliki generalisasi yang memuaskan bagi distribusi tempered.

\begin{prop} Transformasi Fourier $\mathfrak F\colon \fml S^*(\R) \sto \fml S^*(\R)$ merupakan pemetaan linear kontinu
bijektif yang inversnya juga kontinu. Lebih lanjut, $\mathfrak F^4$ merupakan pemetaan identitas pada~$\fml S^*(\R)$.
\end{prop}

Kita akan kembali sejenak ke pembahasan transformasi Fourier dalam bab berikutnya.
















\endinput

 \chapter{TEOREMA GELFAND-NAIMARK}\label{Gelfand_Naimark}

Seperti pada Bab~\ref{spectrum}, dalam bab ini---dan untuk sisa catatan ini---kita mengasumsikan bahwa medan
  \index{konvensi!dalam bab~\ref{Gelfand_Naimark} dan bab-bab berikutnya skalar bersifat kompleks}%
skalar kita adalah medan $\C$ bilangan kompleks.



\section{Ideal-Ideal Maksimal dalam $\fml C(X)$}
\begin{prop}\label{0006801} Jika $J$ adalah ideal sejati dalam suatu aljabar Banach beridentitas, maka
tutupannya juga merupakan ideal sejati.
\end{prop}

\begin{cor} Setiap ideal maksimal dalam aljabar Banach beridentitas bersifat tertutup.
\end{cor}

\begin{exam} Untuk setiap subhimpunan $C$ dari ruang topologis~$X$, himpunan
 \index{ideal!dalam $\fml C(X)$}%
 \index{J@$J_C$ (fungsi kontinu yang lenyap pada~$C$)}%
   \[ J_C := \bigl\{f \in \fml C(X)\colon f^{\sto}(C) = \{0\}\,\bigr\} \]
merupakan ideal dalam $\fml C(X)$. Lebih lanjut, $J_C \supseteq J_D$ apabila
$C~\subseteq~D~\subseteq~X$. (Selanjutnya kita akan menulis $J_x$ untuk ideal
$J_{\{x\}}$.)
\end{exam}

\begin{prop}\label{0006832} Misalkan $X$ ruang topologis kompak dan $I$ ideal sejati
dalam $\fml C(X)$. Maka terdapat $x \in X$ sedemikian sehingga $I \subseteq J_x$.
\end{prop}

\begin{prop} Misalkan $x$ dan $y$ titik-titik dalam ruang Hausdorff kompak. Jika $J_x \subseteq J_y$,
maka $x = y$.
\end{prop}

\begin{prop}\label{0006836} Misalkan $X$ ruang Hausdorff kompak. Subhimpunan $I$ dari $\fml C(X)$
merupakan ideal maksimal dalam $\fml C(X)$ jika dan hanya jika $I = J_x$ untuk suatu $x~\in~X$.
\end{prop}

\begin{cor} Jika $X$ ruang Hausdorff kompak, maka pemetaan $x \mapsto J_x$ dari $X$ ke
$\mx\,\fml C(X)$ adalah bijektif.
\end{cor}

Kekompakan merupakan unsur penting dalam proposisi~\ref{0006836}.

\begin{exam} Dalam aljabar Banach $\fml C_b\bigl(\,(0,1)\,\bigr)$ yang terdiri atas fungsi kontinu
terbatas pada interval $(0,1)$ terdapat ideal maksimal $I$ sedemikian sehingga untuk
\emph{tidak ada} titik $x \in (0,1)$ berlaku $I = J_x$. Misalkan $I$ ideal maksimal yang memuat
ideal $S$ dari semua fungsi $f$ dalam $\fml C_b\bigl(\,(0,1)\,\bigr)$ yang untuknya terdapat
lingkungan $U_f$ dari $0$ dalam $\R$ sedemikian sehingga $f(x) = 0$ untuk semua $x \in U_f \cap (0,1)$.
\end{exam}














\section{Ruang Karakter}
\begin{defn} Suatu
 \index{karakter}%
\df{karakter} (atau
 \index{fungsional!linear!multiplikatif}%
 \index{linear!fungsional!multiplikatif}%
\df{fungsional linear multiplikatif tak nol}) pada aljabar $A$ adalah homomorfisme tak nol
dari $A$ ke dalam~$\C$. Himpunan semua karakter pada $A$ dinotasikan
 \index{delta@$\Delta(A)$!ruang karakter dari~$A$}%
dengan~$\Delta A$.
\end{defn}

\begin{prop}\label{00068511} Misalkan $A$ aljabar beridentitas dan $\phi \in \Delta A$. Maka
 \begin{enumerate}[label=\rm{(\alph*)}]
  \item $\phi(\vc 1) = 1$;
  \item jika $a \in \inv A$, maka $\phi(a) \ne 0$;
  \item jika $a$
 \index{nilpoten}%
  \df{nilpoten} (yaitu, jika $a^n = 0$ untuk suatu $n \in \N$), maka $\phi(a) = 0$;
 \index{idempoten}%
  \item jika $a$ adalah
  \df{idempoten} (yaitu, jika $a^2 = a$), maka $\phi(a)$ adalah $0$ atau~$1$; dan
  \item $\phi(a) \in \sigma(a)$ untuk setiap $a \in A$.
 \end{enumerate}
\end{prop}
Kita catat sambil lalu bahwa bagian (e) dari proposisi sebelumnya tidak memberi kita cara mudah
untuk menunjukkan bahwa spektrum $\sigma(a)$ dari suatu elemen aljabar tidak kosong. Hal ini
bergantung pada pengetahuan bahwa $\Delta(A)$ tidak kosong.

\begin{exam} Pemetaan identitas adalah satu-satunya karakter pada aljabar~$\C$.
\end{exam}

\begin{exam} Misalkan $A$ adalah aljabar matriks berukuran $2 \times 2$, $a = \bigl[a_{ij}\bigr]$,
sedemikian sehingga $a_{12} = 0$. Aljabar ini memiliki tepat dua karakter, yaitu $\phi(a) = a_{11}$
dan $\psi(a) = a_{22}$.
\end{exam}

\begin{proof}[\emph{Petunjuk pembuktian}] Tuliskan $\begin{bmatrix} a & 0 \\ b & c \end{bmatrix}$ sebagai kombinasi linear
dari matriks-matriks $ u = \begin{bmatrix} 1 & 0 \\ 0 & 1 \end{bmatrix}$, $v = \begin{bmatrix} 0 & 0 \\ 1 & 0 \end{bmatrix}$,
dan \smash[t]{$w = \begin{bmatrix} 0 & 0 \\ 0 & 1 \end{bmatrix}$}. Gunakan proposisi~\ref{00068511}.   \ns
\end{proof}

\begin{exam} Aljabar semua matriks $2 \times 2$ dengan entri bilangan kompleks tidak memiliki karakter.
\end{exam}

\begin{prop} Misalkan $A$ aljabar beridentitas dan $\phi$ fungsional linear pada~$A$. Maka $\phi
\in \Delta A$ jika dan hanya jika $\ker\phi$ tertutup terhadap perkalian dan $\phi(\vc 1) =
1$.
\end{prop}

\begin{proof}[\emph{Petunjuk pembuktian}] Untuk arah sebaliknya, terapkan $\phi$ pada hasil kali
$a - \phi(a)\vc 1$ dan $b - \phi(b)\vc 1$ untuk $a$, $b \in A$. \ns
\end{proof}

\begin{prop}\label{00068522} Setiap fungsional linear multiplikatif pada aljabar Banach beridentitas
$A$ kontinu. Bahkan, jika $\phi \in \Delta(A)$, maka $\phi$ merupakan kontraksi dan
$\norm\phi = 1$.
\end{prop}

\begin{exam}\label{00068524} Misalkan $X$ ruang topologis dan $x \in X$. Kita mendefinisikan
 \index{evaluasi!fungsional}%
 \index{evaluasi@$\sbsb EXx$ atau $E_x$ (fungsional evaluasi di~$x$)}%
\df{fungsional evaluasi di $x$}, yang dinotasikan dengan $\sbsb EXx$, melalui
   \[ \sbsb EXx\colon \fml C(X) \sto \C\colon f \mapsto f(x)\,. \]
Fungsional ini merupakan karakter pada~$\fml C(X)$ dan kernelnya adalah $J_x$. Jika hanya ada
satu ruang topologis yang sedang dibahas, kita sederhanakan notasi $\sbsb EXx$ menjadi~$E_x$.
Jadi, khususnya, untuk $f \in \fml C(X)$ kita sering menulis $E_x(f)$ alih-alih
$\sbsb EXx(f)$ yang lebih panjang.
\end{exam}

\begin{defn} Dalam proposisi~\ref{00068522} kita menemukan bahwa setiap karakter pada aljabar Banach
beridentitas $A$ berada pada bola satuan dual~$A^*$. Karena itu kita dapat memberikan himpunan
himpunan $\Delta A$ karakter pada~$A$ dengan topologi-$w^*$ relatif yang diwarisinya
dari~$A^*$. Topologi ini disebut
 \index{Gelfand!topologi}%
 \index{topologi!Gelfand}%
\df{topologi Gelfand} pada $\Delta A$, dan ruang topologis yang dihasilkan disebut
 \index{ruang!karakter}%
 \index{karakter!ruang}%
 \index{ruang!struktur}%
 \index{struktur!ruang}%
\df{ruang karakter} (atau \emph{ruang pembawa}, atau \emph{ruang struktur}, atau \emph{ruang ideal maksimal}, atau \emph{spektrum}) dari~$A$.
\end{defn}

\begin{prop} Ruang karakter dari aljabar Banach beridentitas merupakan ruang Hausdorff kompak.
\end{prop}

\begin{exam}\label{000633} Misalkan $l_1(\Z)$ adalah keluarga semua barisan bilateral
   \[ (\dots,a_{-2},a_{-1},a_0,a_1,a_2,\dots) \]
yang
 \index{mutlak!terjumlahkan!barisan bilateral}%
\df{terjumlahkan mutlak}; yaitu, yang memenuhi $\sum_{k=-\infty}^{\infty} \abs{a_k} <
\infty$. Ini merupakan ruang Banach di bawah operasi penjumlahan dan perkalian skalar titik demi titik,
dengan norma
   \[ \norm a = \sum_{k=-\infty}^{\infty} \abs{a_k}\,. \]
Untuk $a$, $b \in l_1(\Z)$ definisikan $a \ast b$ sebagai barisan yang entri ke-$n$-nya
diberikan oleh
  \[ (a \ast b)_n = \sum_{k=-\infty}^{\infty} a_{n-k}\,b_k\,. \]
Operasi $\ast$ disebut
 \index{konvolusi!dalam $l_1(\Z)$}%
\df{konvolusi}. (Untuk melihat asal definisi ini, cobalah mengalikan deret pangkat
$\sum_{-\infty}^{\infty}a_kz^k$ dan $\sum_{-\infty}^{\infty}b_kz^k$.) Dengan operasi tambahan ini
$l_1(\Z)$ menjadi aljabar Banach komutatif beridentitas.

Ruang ideal maksimal dari aljabar Banach ini
 \index{ruang!ideal maksimal!dari $l_1(\Z)$}%
 \index{ruang karakter!dari $l_1(\Z)$}%
adalah (homeomorfik dengan) lingkaran satuan~$\T$.
\end{exam}

\begin{proof}[\emph{Petunjuk pembuktian}] Untuk setiap $z \in \T$ definisikan
   \[ \psi_z\colon l_1(\Z) \sto \C\colon a \mapsto \sum_{k=-\infty}^\infty a_kz^k\,. \]
Tunjukkan bahwa $\psi_z \in \Delta l_1(\Z)$. Kemudian tunjukkan bahwa pemetaan
   \[ \psi\colon \T \sto \Delta l_1(\Z)\colon z \mapsto \psi_z \]
adalah suatu homeomorfisme. \ns
\end{proof}

\begin{prop} Jika $\phi \in \Delta A$ dengan $A$ aljabar beridentitas, maka $\ker\phi$ merupakan
ideal maksimal dalam~$A$.
\end{prop}

\begin{proof}[\emph{Petunjuk pembuktian}] Untuk menunjukkan kemaksimalan, misalkan $I$ ideal dalam $A$
yang memuat $\ker\phi$ secara sejati. Pilih $z \in I \setminus \ker\phi$. Pertimbangkan elemen
$x - \bigl(\phi(x)/\phi(z)\bigr)\,z$, dengan $x$ elemen sebarang dari~$A$.
\ns
\end{proof}

\begin{prop} Suatu karakter pada aljabar beridentitas ditentukan sepenuhnya oleh kernelnya.
\end{prop}

\begin{proof}[\emph{Petunjuk pembuktian}] Misalkan $a$ elemen aljabar dan $\phi$ suatu karakter.
Untuk berapa banyak bilangan kompleks $\lambda$ dapat $a^2 - \lambda a$ menjadi anggota kernel~$\phi$? \ns
\end{proof}

\begin{cor} Jika $A$ aljabar beridentitas, maka pemetaan $\phi \mapsto \ker\phi$ dari
$\Delta A$ ke~$\mx A$ injektif.
\end{cor}

\begin{prop}\label{0007207} Misalkan $I$ ideal maksimal dalam aljabar Banach komutatif beridentitas~$A$.
Maka terdapat karakter pada $A$ yang kernelnya adalah~$I$.
\end{prop}

\begin{proof}[\emph{Petunjuk pembuktian}] Gunakan korolari~\ref{C073147}. Mengapa kita dapat memandang
pemetaan hasil bagi sebagai suatu karakter?  \ns
\end{proof}

\begin{cor}\label{0007207a} Jika $A$ aljabar Banach komutatif beridentitas, maka pemetaan
$\phi \mapsto \ker\phi$ merupakan bijeksi dari $\Delta A$ ke~$\mx A$.
\end{cor}

\begin{defn} Misalkan $A$ aljabar Banach komutatif beridentitas. Berdasarkan korolari sebelumnya,
kita dapat memberikan $\mx A$ suatu topologi yang membuatnya homeomorfik dengan ruang karakter~$\Delta A$.
Inilah
 \index{ruang!ideal maksimal}%
 \index{ruang!maksimal ideal}%
 \index{maksimal@$\mx A$!homeomorfik dengan $\Delta(A)$}%
\df{ruang ideal maksimal} dari~$A$. Karena $\Delta A$ dan $\mx A$ homeomorfik, lazim untuk
mengidentifikasi keduanya, sehingga $\Delta A$ sering disebut \emph{ruang ideal maksimal} dari~$A$.
\end{defn}

\begin{defn}\label{00073} Misalkan $X$ ruang Hausdorff kompak dan $x \in X$.
Ingat bahwa dalam contoh~\ref{00068524} kita mendefinisikan $\sbsb EXx$, yaitu \emph{fungsional
evaluasi} di~$x$, melalui
   \[ \sbsb EXx(f) := f(x) \]
untuk setiap $f \in \fml C(X)$. Pemetaan
   \[ \sbsb EX \colon X \sto \Delta\fml C(X)\colon x \mapsto \sbsb EXx \]
adalah
 \index{pemetaan evaluasi}%
 \index{evaluasi@$\sbsb EX$!pemetaan evaluasi}%
\df{pemetaan evaluasi} pada~$X$. Sebagaimana disebutkan sebelumnya, ketika hanya satu ruang topologis
yang sedang dipertimbangkan, biasanya kita menyingkat $\sbsb EX$ menjadi $E$ dan $\sbsb EXx$ menjadi~$E_x$.
\end{defn}

\begin{notn} Untuk menyatakan bahwa dua ruang topologis $X$ dan $Y$ homeomorfik, kita
 \index{<binrel@$X \approx Y$ ($X$ dan $Y$ homeomorfik)}%
menulis $X \approx Y$.
\end{notn}

\begin{prop}\label{000731} Misalkan $X$ ruang Hausdorff kompak. Maka
 \index{ruang!ideal maksimal!dari $\fml C(X)$}%
 \index{ruang karakter!dari $\fml C(X)$}%
pemetaan evaluasi pada~$X$
   \[ \sbsb EX \colon X \sto \Delta\fml C(X) \colon x \mapsto \sbsb EXx \]
merupakan homeomorfisme. Jadi kita memiliki
   \[ X \approx \Delta\fml C(X) \approx \mx\fml C(X)\,. \]
Lebih dari itu: bukan hanya setiap $\sbsb EX$ merupakan homeomorfisme antara ruang Hausdorff kompak,
melainkan $E$ sendiri merupakan suatu \emph{ekuivalensi natural} antara funktor---funktor identitas
dan funktor $\Delta\fml C$.
\end{prop}

Identifikasi antara ruang Hausdorff kompak $X$ dan ruang karakternya serta ruang ideal maksimalnya
begitu kuat sehingga banyak orang membicarakannya seolah-olah ketiganya benar-benar sama. Sangat lazim
untuk mendengar, misalnya, bahwa ``ideal maksimal dalam $\fml C(X)$ hanyalah titik-titik di~$X$''.
Walaupun secara harfiah tidak benar, ungkapan itu memang terdengar sedikit kurang mengintimidasi daripada
``ideal maksimal dari $\fml C(X)$ adalah kernel fungsional evaluasi pada titik-titik~$X$''.

\begin{prop}\label{000735} Misalkan $X$ dan $Y$ ruang Hausdorff kompak dan
 \index{funktor!$\fml C$ sebagai suatu}%
 \index{C@$\fml C$ (funktor)}%
$F \colon X \sto Y$ kontinu. Ingat bahwa dalam contoh~\ref{C029717} kita mendefinisikan $\fml C(F)$
pada $\fml C(Y)$ melalui
   \[ \fml C(F)(g) = g \circ F\]
untuk semua $g \in \fml C(Y)$. Maka
  \begin{enumerate}[label=\rm{(\alph*)}]
   \item $\fml C(F)$ memetakan $\fml C(Y)$ ke dalam $\fml C(X)$.
   \item Pemetaan $\fml C(F)$ merupakan homomorfisme aljabar Banach beridentitas yang kontraktif.
   \item $\fml C(F)$ injektif jika dan hanya jika $F$ surjektif.
   \item $\fml C(F)$ surjektif jika dan hanya jika $F$ injektif.
   \item Jika $X$ homeomorfik dengan $Y$, maka $\fml C(X)$ isomorfik secara isometrik
   dengan~$\fml C(Y)$.
  \end{enumerate}
\end{prop}

\begin{prop}\label{000736} Misalkan $A$ dan $B$ aljabar Banach komutatif beridentitas dan
 \index{funktor!$\Delta$ sebagai suatu}%
 \index{delta@$\Delta$ (funktor)}%
$T\colon A \sto B$ homomorfisme aljabar beridentitas. Definisikan $\Delta T$ pada $\Delta B$ melalui
   \[ \Delta T(\psi) = \psi \circ T \]
untuk semua $\psi \in \Delta B$. Maka
  \begin{enumerate}[label=\rm{(\alph*)}]
   \item $\Delta T$ memetakan $\Delta B$ ke dalam $\Delta A$.
   \item Pemetaan $\Delta T\colon \Delta B \sto \Delta A$ kontinu.
   \item Jika $T$ surjektif, maka $\Delta T$ injektif.
   \item Jika $T$ merupakan suatu isomorfisme (aljabar), maka $\Delta T$ merupakan homeomorfisme.
   \item Jika $A$ dan $B$ isomorfik (secara aljabar), maka $\Delta A$ dan $\Delta B$
   homeomorfik.
  \end{enumerate}
\end{prop}

\begin{cor} Misalkan $X$ dan $Y$ ruang Hausdorff kompak. Jika $\fml C(X)$ dan $\fml C(Y)$ isomorfik
secara aljabar, maka $X$ dan $Y$ homeomorfik.
\end{cor}

\begin{cor}\label{nasc_homeomor_CHSps} Dua ruang Hausdorff kompak homeomorfik jika dan hanya jika aljabar
fungsi kontinu bernilai kompleksnya isomorfik (secara aljabar).
\end{cor}

Korolari~\ref{nasc_homeomor_CHSps} menimbulkan sekumpulan masalah yang sangat menarik. Korolari itu menyarankan
bahwa untuk setiap sifat topologis ruang Hausdorff kompak $X$ seharusnya ada sifat aljabar yang bersesuaian
dari aljabar $\fml C(X)$, dan \emph{sebaliknya}. Namun, korolari tersebut tidak memberikan resep yang dapat
digunakan untuk menemukan korespondensi ini. Salah satu contoh yang sangat sederhana dari hasil semacam ini
diberikan di bawah sebagai proposisi~\ref{prop_connected_idempotent}. Untuk pembahasan yang indah mengenai
pertanyaan ini, lihat monograf Semadeni~\cite{Semadeni:1971}, khususnya tabel pada halaman~132--133.

\begin{cor} Misalkan $X$ dan $Y$ ruang Hausdorff kompak. Jika $\fml C(X)$ dan $\fml C(Y)$ isomorfik
secara aljabar, maka keduanya isomorfik secara isometrik.
\end{cor}

\begin{prop}\label{prop_connected_idempotent} Ruang topologis $X$ terhubung jika dan hanya jika aljabar $\fml C(X)$
tidak memuat idempoten nontrivial.
\end{prop}

(Idempoten \emph{trivial} suatu aljabar adalah $\vc 0$ dan~$\vc 1$.)













\section{Transformasi Gelfand}
\begin{defn} Misalkan $A$ aljabar Banach komutatif dan $a\in A$. Definisikan
   \[\sbsb{\Gamma}A a\colon \Delta A \sto \C \colon \phi \mapsto \phi(a) \]
untuk setiap $\phi \in \Delta(A)$. (Notasi alternatif: jika tidak menimbulkan kebingungan, kita
sering menulis $\Gamma a$ atau $\wh a$ untuk $\sbsb{\Gamma}A a$.) Pemetaan $\sbsb{\Gamma}A$ adalah
 \index{Gelfand!transformasi!pada aljabar Banach}%
 \index{Gamma@$\sbsb\Gamma A$, $\Gamma$ (transformasi Gelfand)}%
 \index{<unoptop@$\wh a$ (transformasi Gelfand dari~$a$)}%
\df{transformasi Gelfand pada~$A$}.
\end{defn}

Karena $\Delta A \subseteq A^*$, jelas bahwa $\sbsb{\Gamma}A a$ hanyalah restriksi $a^{**}$
pada ruang karakter~$A$. Lebih lanjut, topologi Gelfand pada $\Delta A$ adalah topologi-$w^*$ relatif,
yaitu topologi terlemah yang membuat $a^{**}$ kontinu pada $\Delta A$ untuk setiap $a \in A$,
sehingga $\sbsb{\Gamma}A a$ merupakan fungsi kontinu pada $\Delta A$. Jadi
$\sbsb{\Gamma}A \colon A \sto \fml C(\Delta A)$.

Sebagai ringkasan dan demi kemudahan, elemen $\sbsb{\Gamma}A a = \Gamma a = \wh a$ biasanya cukup
disebut
 \index{Gelfand!transformasi!dari suatu elemen}%
\emph{transformasi Gelfand dari~$a$}---karena frasa \emph{transformasi Gelfand pada $A$ yang dievaluasi
pada~$a$} terasa janggal.

\begin{defn} Kita mengatakan bahwa suatu keluarga fungsi $\fml F$ pada himpunan $S$
 \index{pemisahan!titik}%
\df{memisahkan titik-titik} pada $S$ (atau merupakan
 \index{pemisah!keluarga fungsi}%
\df{keluarga fungsi pemisah} pada $S$) jika untuk setiap pasangan elemen berbeda $x$
dan $y$ dari $S$ terdapat $f \in \fml F$ sedemikian sehingga $f(x) \ne f(y)$.
\end{defn}

\begin{prop} Misalkan $X$ ruang topologis kompak. Maka $\fml C(X)$ memisahkan titik-titik jika dan
hanya jika $X$ Hausdorff.
\end{prop}

\begin{prop}\label{0007402} Jika $A$ aljabar Banach komutatif, maka $\sbsb{\Gamma}A
\colon A \sto \fml C(\Delta A)$ merupakan homomorfisme aljabar yang kontraktif dan jangkauan
$\sbsb{\Gamma}A$ merupakan subaljabar pemisah dari $\fml C(\Delta A)$. Jika, sebagai tambahan,
$A$ beridentitas, maka $\Gamma$ bernorma satu, jangkauannya merupakan subaljabar beridentitas dari
$\fml C(\Delta A)$, dan pemetaan itu merupakan homomorfisme beridentitas.
\end{prop}

\begin{prop} Misalkan $a$ elemen aljabar Banach komutatif beridentitas~$A$. Maka $a$ invertibel dalam
$A$ jika dan hanya jika $\hat a$ invertibel dalam~$\fml C(\Delta A)$.
\end{prop}

\begin{prop}\label{00074043} Misalkan $A$ aljabar Banach komutatif beridentitas dan $a$ suatu
elemen dari~$A$. Maka $\ran\hat a = \sigma(a)$ dan $\norm{\hat a}_u~=~\rho(a)$.
\end{prop}

\begin{defn} Suatu elemen $a$ dari aljabar Banach disebut
 \index{kuasinilpoten}%
\df{kuasinilpoten} jika $\lim\limits_{n \sto \infty}\norm{a^n}^{1/n} = 0$.
\end{defn}

\begin{prop} Misalkan $a$ elemen aljabar Banach komutatif beridentitas~$A$. Maka pernyataan-pernyataan
berikut ekuivalen:
  \begin{enumerate}[label=\rm{(\alph*)}]
   \item $a$ kuasinilpoten;
   \item $\rho(a) = 0$;
   \item $\sigma(a) = \{0\}$;
   \item $\Gamma a = 0$;
   \item $\phi(a) = 0$ untuk setiap $\phi \in \Delta A$;
   \item $a \in \bigcap\mx A$.
  \end{enumerate}
\end{prop}

\begin{defn} Suatu aljabar Banach disebut
 \index{semisederhana}%
\df{semisederhana} jika tidak memiliki elemen kuasinilpoten tak nol.
\end{defn}

\begin{prop} Misalkan $A$ aljabar Banach komutatif beridentitas. Pernyataan-pernyataan berikut ekuivalen:
  \begin{enumerate}[label=\rm{(\alph*)}]
   \item $A$ semisederhana;
   \item jika $\rho(a) = 0$, maka $a = 0$;
   \item jika $\sigma(a) = \{0\}$, maka $a = 0$;
   \index{monomorfisme}%
   \item transformasi Gelfand $\sbsb{\Gamma}A$ merupakan monomorfisme (yaitu,
   homomorfisme injektif);
   \item jika $\phi(a) = 0$ untuk setiap $\phi \in \Delta A$, maka $a = 0$;
   \item $\bigcap\mx A = \{0\}$.
  \end{enumerate}
\end{prop}

\begin{prop}\label{00074083} Misalkan $A$ aljabar Banach komutatif beridentitas. Maka
pernyataan-pernyataan berikut ekuivalen:
  \begin{enumerate}[label=\rm{(\alph*)}]
    \item $\|a^2\| = \|a\|^2$ untuk semua $a \in A$;
    \item $\rho(a) = \|a\|$ untuk semua $a \in A$; dan
    \item transformasi Gelfand merupakan isometri; yaitu, $\norm{\wh a}_u = \|a\|$ untuk
    semua $a \in A$.
  \end{enumerate}
\end{prop}

\begin{exam} Transformasi Gelfand pada $l_1(\Z)$ bukan isometri.
\end{exam}

Ingat dari proposisi~\ref{000731} bahwa ketika $X$ ruang Hausdorff kompak, pemetaan evaluasi $E_X$
mengidentifikasi ruang $X$ dengan ruang ideal maksimal aljabar Banach~$\fml C(X)$. Jadi menurut proposisi~\ref{000735},
pemetaan $\fml CE_X$ mengidentifikasi aljabar $\fml C(\Delta(\fml C(X)))$ fungsi kontinu pada ruang ideal maksimal
ini dengan aljabar $\fml C(X)$ itu sendiri. Ternyata transformasi Gelfand pada aljabar $\fml C(X)$ hanyalah invers
dari pemetaan identifikasi ini.

\begin{exam} Misalkan $X$ ruang Hausdorff kompak. Maka transformasi Gelfand pada aljabar Banach
$\fml C(X)$ merupakan isomorfisme isometrik. Bahkan, pada $\fml C(X)$ transformasi Gelfand
\smash[b]{$\sbsb\Gamma X$} adalah $(\fml CE_X)^{-1}$.
\end{exam}

\begin{exam}\label{0007472} Ruang ideal maksimal aljabar Banach $L_1(\R)$
 \index{ruang!ideal maksimal!untuk $L_1(\R)$}%
adalah $\R$ sendiri dan transformasi Gelfand pada $L_1(\R)$ adalah transformasi Fourier.
\end{exam}

\begin{proof} Lihat~\cite{BaggettF:1979}, halaman 169--171.  \ns
\end{proof}

Fungsi $t \mapsto e^{it}$ merupakan bijeksi dari interval $[-\pi,\pi)$ ke
 \index{T@$\T$ (lingkaran satuan pada bidang kompleks)}%
lingkaran satuan~$\T$ pada bidang kompleks. Salah satu akibatnya adalah kita tidak perlu membedakan antara
  \begin{enumerate}[label=\rm{(\alph*)}]
     \item fungsi $2\pi$-periodik pada~$\R$,
     \item semua fungsi pada $[-\pi,\pi)$,
     \item fungsi $f$ pada $[-\pi,\pi]$ sedemikian sehingga $f(-\pi) = f(\pi)$, dan
     \item fungsi pada~$\T$.
  \end{enumerate}
Dalam kelanjutan, kita akan sering mengidentifikasi kelas-kelas fungsi ini tanpa penjelasan lebih lanjut.

Identifikasi lain yang berguna adalah antara lingkaran satuan $\T$ dalam $\C$ dan ruang ideal maksimal
aljabar $l_1(\Z)$. Homeomorfisme $\psi$ antara dua ruang Hausdorff kompak ini didefinisikan dalam contoh~\ref{000633}.
Dalam bekerja dengan transformasi Gelfand $\Gamma_a$ dari elemen $a \in l_1(\Z)$, sering secara teknis lebih mudah
memperlakukannya sebagai fungsi, sebut saja $G_a$, yang daerah asalnya adalah~$\T$, seperti ditunjukkan oleh diagram berikut.
  \[
    \xy
      \btriangle[\T`\Delta l_1(\Z)`\C;\psi`G_a`\Gamma_a]
    \endxy
  \]
Jadi untuk $a \in l_1(\Z)$ dan $z \in \T$ berlaku
  \[ G_a(z) = \Gamma_a(\psi_z) = \psi_z(a) = \sum_{k=-\infty}^\infty a_kz^k\,. \]

\begin{defn} Jika $f \in \fml L_1([-\pi,\pi))$, maka
  \index{Fourier!deret}%
\df{deret Fourier} untuk $f$ adalah deret
   \[ \sum_{n=-\infty}^\infty\wt f(n)\exp(int) \qquad -\pi \le t \le \pi \]
dengan barisan $\wt f$ didefinisikan oleh
   \[ \wt f(n) = {\textstyle\frac1{2\pi}}\int_{-\pi}^\pi f(t)\exp(-int)\,dt \]
untuk semua $n \in \Z$. Barisan dua arah tak hingga $\wt f$ adalah
  \index{Fourier!transformasi!pada $L_1([-\pi,\pi])$}%
\df{transformasi Fourier} dari $f$, dan bilangan $\wt f(n)$ adalah
  \index{Fourier!koefisien}%
\df{koefisien Fourier} ke-$n$ dari $f$. Jika $\wt f \in l_1(\Z)$, kita mengatakan bahwa $f$ memiliki
  \index{mutlak!konvergen!deret Fourier}%
  \index{Fourier!deret!konvergen mutlak}%
\df{deret Fourier yang konvergen mutlak}. Himpunan semua fungsi kontinu pada $\T$
dengan deret Fourier yang konvergen mutlak dinotasikan dengan
  \index{mutlak@$\fml{A\,C}(\T)$ (deret Fourier yang konvergen mutlak}%
$\fml{A\,C}(\T)$.
\end{defn}

\begin{prop} Jika $f$ fungsi kontinu $2\pi$-periodik pada $\R$ yang transformasi Fouriernya nol,
maka $f = 0$.
\end{prop}

\begin{cor} Transformasi Fourier pada $\fml C(\T)$ injektif.
\end{cor}

\begin{prop} Transformasi Fourier pada $\fml C(\T)$ merupakan invers kiri dari transformasi Gelfand
pada~$l_1(\Z)$.
\end{prop}

\begin{prop} Jangkauan transformasi Gelfand pada $l_1(\Z)$ adalah aljabar Banach komutatif beridentitas
$\fml{A\,C}(\T)$.
\end{prop}

\begin{prop} Terdapat fungsi kontinu yang deret Fouriernya divergen di~$0$.
\end{prop}

\begin{proof} Lihat, misalnya, \cite{HewittS:1965}, latihan 18.45.  \ns
\end{proof}

Apa yang dikatakan hasil sebelumnya mengenai transformasi Gelfand $\Gamma\colon l_1(\Z) \sto
\fml C(\T)$?

\vskip 3 pt

Misalkan suatu fungsi $f$ termasuk dalam $\fml{AC}(\T)$ dan tidak pernah nol. Maka $1/f$ tentu kontinu pada~$\T$,
tetapi apakah deret Fouriernya konvergen mutlak? Salah satu keberhasilan awal studi abstrak aljabar Banach
adalah pembuktian yang sangat sederhana atas pertanyaan ini, yang mula-mula diberikan oleh Norbert Wiener.
Bukti asli Wiener melalui perbandingan cukup sulit.

\begin{thm}[Teorema Wiener] Misalkan $f$ fungsi kontinu pada $\T$ yang tidak pernah
 \index{teorema Wiener}%
nol. Jika $f$ memiliki deret Fourier yang konvergen mutlak, maka kebalikannya~$1/f$ juga demikian.
\end{thm}

\begin{exam} Transformasi Laplace juga dapat dipandang sebagai kasus khusus transformasi Gelfand.
Untuk rinciannya, lihat~\cite{BaggettF:1979}, halaman 173--175.
\end{exam}













\section{Aljabar-$C^*$ Beridentitas}

\begin{prop} Jika $a$ elemen suatu aljabar-$*\,$, maka $\sigma(a^*) = \conj{\sigma(a)}$.
\end{prop}


\begin{prop}\label{001311} Dalam setiap aljabar-$C^*$, involusi merupakan isometri. Artinya,
 \index{involusi!merupakan isometri}%
$\norm{a^*} = \norm{a}$ untuk setiap elemen $a$ dalam aljabar.
\end{prop}

Dalam definisi~\ref{def_norm_alg} tentang \emph{aljabar bernorma}, kita memberikan syarat khusus bahwa
elemen identitas dari aljabar bernorma beridentitas harus bernorma satu. Dalam aljabar-$C^*$, syarat
ini bersifat berlebihan.

\begin{cor} Dalam aljabar-$C^*$ beridentitas, $\norm{\vc 1} = 1$.
\end{cor}

\begin{cor}\label{001311005fa} Setiap elemen uniter dalam aljabar-$C^*$ beridentitas bernorma satu.
\end{cor}

\begin{cor}\label{001311008} Jika $a$ elemen aljabar-$C^*$ $A$ sedemikian sehingga
$ab = \vc 0$ untuk setiap $b \in A$, maka $a = \vc 0$.
\end{cor}

\begin{prop}\label{00131101} Misalkan $a$ elemen normal dari aljabar-$C^*$.
Maka $\norm{a^2} = {\norm a}^2$. Jika aljabar tersebut beridentitas, maka $\rho(a)~=~\norm a$.
\end{prop}

\begin{cor} Jika $p$ proyeksi tak nol dalam aljabar-$C^*$, maka $\norm p = 1$.
\end{cor}

\begin{cor}\label{00131102} Misalkan $a \in A$, dengan $A$ suatu aljabar-$C^*$ komutatif. Maka
$\norm{a^2} = {\norm a}^2$ dan, jika sebagai tambahan $A$ beridentitas, maka $\rho(a)~=~\norm a$.
\end{cor}

\begin{cor}\label{00131103} Pada aljabar-$C^*$ komutatif beridentitas $A$, transformasi Gelfand
$\Gamma$ merupakan isometri; yaitu, $\norm{\Gamma_a}_u = \norm{\wh a}_u = \norm a$ untuk
setiap $a \in A$.
\end{cor}

\begin{cor}\label{00131104} Norma aljabar-$C^*$ beridentitas bersifat unik, dalam arti bahwa
untuk aljabar beridentitas $A$ dengan involusi, terdapat paling banyak satu norma yang menjadikan
$A$ sebuah aljabar-$C^*$.
\end{cor}

\begin{prop}\label{001312} Jika $a$ elemen swaadjoin dari aljabar-$C^*$ beridentitas,
 \index{spektrum!elemen swaadjoin}%
 \index{swaadjoin!spektrum elemen}%
maka $\sigma(a) \subseteq \R$.
\end{prop}

\begin{proof}[\emph{Petunjuk pembuktian}] Tuliskan $\lambda \in \sigma(a)$ dalam bentuk $\alpha + i\beta$, dengan $\alpha$, $\beta \in \R$.
Untuk setiap $n \in \N$ misalkan $b_n := a - \alpha \vc 1 + i n\beta\vc 1$. Tunjukkan bahwa untuk setiap $n$ skalar
$i(n+1)\beta$ termasuk dalam spektrum $b_n$ dan karena itu
   \[ {\abs{i(n+1)\beta}}^2 \le {\norm{a - \alpha\vc 1}}^2 + n^2\beta^2\,. \]   \ns
\end{proof}

\begin{cor}\label{001406} Misalkan $A$ aljabar-$C^*$ komutatif beridentitas. Jika $a \in A$ swaadjoin,
maka transformasi Gelfand $\wh a$ bernilai real.
\end{cor}

\begin{cor}\label{0014062} Setiap karakter pada aljabar-$C^*$ merupakan homomorfisme-$*\,$.
\end{cor}

\begin{cor}\label{0014062fa} Transformasi Gelfand pada aljabar-$C^*$ komutatif beridentitas $A$ merupakan
homomorfisme-$*\,$ beridentitas ke dalam~$\fml C(\Delta A)$.
\end{cor}

\begin{cor} Jangkauan transformasi Gelfand pada aljabar-$C^*$ komutatif beridentitas $A$ merupakan
subaljabar pemisah, swaadjoin, beridentitas dari~$\fml C(\Delta A)$.
\end{cor}

\begin{prop}\label{001313} Jika $u$ elemen uniter dari aljabar-$C^*$ beridentitas,
 \index{spektrum!elemen uniter}%
 \index{uniter!spektrum elemen}%
maka $\sigma(u) \subseteq \T$.
\end{prop}

\begin{prop} Transformasi Gelfand $\Gamma$ merupakan ekuivalensi natural antara funktor identitas
dan funktor $\fml C\Delta$ pada kategori aljabar-$C^*$ komutatif beridentitas dan homomorfisme-$*\,$
beridentitas.
\end{prop}













\section{Teorema Gelfand-Naimark}

Kita mengingat kembali suatu teorema baku dari analisis real.

\begin{thm}[Teorema Stone-Weierstrass]\label{00141} Misalkan $X$ ruang Hausdorff kompak. Setiap
 \index{teorema Stone-Weierstrass}%
subaljabar-$*\,$ pemisah beridentitas dari $\fml C(X)$ rapat.
\end{thm}

\begin{proof} Lihat \cite{Douglas:1972}, teorema 2.40. \ns \end{proof}

Versi pertama dari \emph{teorema Gelfand-Naimark} menyatakan bahwa setiap aljabar-$C^*$ komutatif
beridentitas adalah aljabar semua fungsi kontinu pada suatu ruang Hausdorff kompak. Tentu saja kata
\emph{adalah} dalam kalimat sebelumnya berarti \emph{isomorfik secara isometrik dan sebagai-$*$ dengan}.
Ruang Hausdorff kompak yang dimaksud adalah ruang karakter aljabar tersebut.

\begin{thm}[Teorema Gelfand-Naimark I]\label{00142} Misalkan $A$ aljabar-$C^*$ komutatif
 \index{Gelfand-Naimark!teorema I}%
beridentitas. Maka transformasi Gelfand $\Gamma_A\colon a \mapsto \wh a$ merupakan
isomorfisme-$*\,$ isometrik dari $A$ ke~$\fml C(\Delta A)$.
\end{thm}






\begin{defn} Misalkan $S$ subhimpunan tak kosong dari aljabar-$C^*$~$A$. Irisan
keluarga semua subaljabar-$C^*$ dari $A$ yang memuat $S$ disebut
 \index{C*@$C^*$-aljabar!yang dibangkitkan oleh suatu himpunan}%
\df{subaljabar-$C^*$ yang dibangkitkan oleh~$S$}. Kita menotasikannya dengan
 \index{C*@$C^*(S)$ ($C^*$-subalgebra yang dibangkitkan oleh~$S$)}%
$C^*(S)$. (Mudah dilihat bahwa irisan suatu keluarga subaljabar-$C^*$ memang merupakan
aljabar-$C^*$.) Dalam beberapa kasus kita sedikit menyingkat notasi: misalnya, jika $a \in A$ kita
menulis $C^*(a)$ untuk~$C^*(\{a\})$ dan $C^*(\vc 1,a)$ untuk $C^*(\{\vc 1,a\})$.
\end{defn}

\begin{prop} Misalkan $S$ subhimpunan tak kosong dari aljabar-$C^*$~$A$. Untuk setiap bilangan natural
$n$ definisikan himpunan $W_n$ sebagai himpunan semua elemen $a$ dari $A$ yang untuknya terdapat
$x_1$, $x_2$, \dots, $x_n$ dalam $S \cup S^*$ sedemikian sehingga $a = x_1x_2\cdots x_n$. Misalkan
$W = \bigcup_{n=1}^\infty W_n$. Maka
   \[ C^*(S) = \clo{\spn W}\,. \]
\end{prop}













\begin{thm}[Teorema Spektral Abstrak]\label{00144} Jika $a$ elemen normal dari aljabar-$C^*$
 \index{abstrak!teorema spektral}%
 \index{spektral!teorema!abstrak}%
beridentitas $A$, maka aljabar-$C^*$ $\fml C(\sigma(a))$ isomorfik secara isometrik dan sebagai-$*$
dengan $C^*(\vc 1,a)$.
\end{thm}

\begin{proof}[\emph{Petunjuk pembuktian}] Gunakan transformasi Gelfand dari $a$ untuk mengidentifikasi
ruang ideal maksimal $C^*(\vc 1,a)$ dengan spektrum~$a$. Terapkan funktor~$\fml C$.
Komposisikan pemetaan yang dihasilkan dengan $\Gamma^{-1}$, dengan $\Gamma$ transformasi Gelfand pada
aljabar-$C^*$~$C^*(\vc 1,a)$.   \ns
\end{proof}

Jika $\psi\colon \fml C(\sigma(a)) \sto C^*(\vc 1,a)$ adalah isomorfisme-$*\,$ pada teorema sebelumnya dan $f$
merupakan fungsi kontinu pada spektrum $a$, maka elemen $\psi(f)$ dalam aljabar $A$ biasanya dinotasikan
dengan~$f(a)$. Operasi yang mengaitkan fungsi kontinu $f$ dengan elemen $f(a)$ dalam $A$ ini disebut
 \index{kalkulus!fungsional}%
 \index{fungsional!kalkulus}%
\emph{kalkulus fungsional} yang terkait dengan~$a$.

\begin{exam} Misalkan pada teorema sebelumnya $\psi\colon \fml C(\sigma(a)) \sto
C^*(\vc 1,a)$ merealisasikan isomorfisme-$*\,$ isometrik. Maka citra di bawah $\psi$ dari fungsi konstan~$\vc 1$
pada spektrum $a$ adalah $\vc 1_A$, dan citra di bawah $\psi$ dari fungsi identitas $\lambda \mapsto \lambda$
pada spektrum $a$ adalah~$a$.
\end{exam}

\begin{exam}\label{X_sqroot_op} Misalkan $T$ operator normal pada ruang Hilbert $H$ yang spektrumnya termuat
dalam~$[0,\infty)$. Misalkan $\psi\colon \fml C(\sigma(T)) \sto C^*(I,T)$ merealisasikan isomorfisme-$*\,$
isometrik antara kedua aljabar-$C^*$ ini. Maka terdapat setidaknya satu operator $\sqrt T$ yang kuadratnya
adalah~$T$. Memang, setiap kali $f$ fungsi kontinu pada spektrum operator normal~$T$, kita dapat secara bermakna
membicarakan operator~$f(T)$.
\end{exam}

\begin{exam} Jika $N$ operator normal pada ruang Hilbert yang spektrumnya adalah $\{0,1\}$, maka
$N$ merupakan proyeksi ortogonal.
\end{exam}

\begin{exam} Jika $N$ operator normal pada ruang Hilbert yang spektrumnya termuat dalam lingkaran satuan~$\T$,
maka $N$ merupakan operator uniter.
\end{exam}

\begin{prop}[Teorema pemetaan spektral]\label{001446} Misalkan $a$ elemen swaadjoin
 \index{spektral!teorema pemetaan}%
dalam aljabar-$C^*$ beridentitas~$A$ dan $f$ fungsi kontinu bernilai kompleks pada spektrum
$a$. Maka
  \[ \sigma(f(a)) = f^{\sto}(\sigma(a))\,. \]
\end{prop}

Kalkulus fungsional memetakan fungsi kontinu ke fungsi kontinu dalam pengertian berikut.

\begin{prop}\label{001449} Misalkan $A$ aljabar-$C^*$ beridentitas, $K$ subhimpunan kompak tak kosong dari~$\R$,
dan $f\colon K \sto \C$ fungsi kontinu. Misalkan $\ofml H_K$ himpunan semua elemen swaadjoin dari $A$ yang spektrumnya
merupakan subhimpunan dari~$K$. Maka fungsi $f\colon \ofml H_K \sto A\colon h \mapsto f(h)$ kontinu.
\end{prop}

\begin{proof}[\emph{Petunjuk pembuktian}] Aproksimasi $f$ secara seragam dengan suatu polinom dan gunakan
\emph{teorema pemetaan spektral}~\ref{001446} serta argumen $\epsilon$ dibagi $3$.   \ns
\end{proof}













\endinput

 \chapter{BERTAHAN TANPA IDENTITAS}

\section{Unitalisasi Aljabar Banach}
\begin{defn}\label{00150005} Misalkan $A$ suatu aljabar. Definisikan
 \index{<binopunit@$A \bowtie \C$ (unitalisasi suatu aljabar)}%
$A \bowtie \C$ sebagai himpunan $A \times \C$ yang penjumlahan dan perkalian skalarnya
didefinisikan secara titik demi titik dan perkaliannya didefinisikan oleh
  \[ (a,\lambda)\cdot(b,\mu) = (ab + \mu a + \lambda b, \lambda\mu). \]
Jika aljabar $A$ dilengkapi dengan suatu involusi, definisikan involusi pada $A \bowtie \C$
secara titik demi titik; yaitu, $(a,\lambda)^* := (a^*, \conj\lambda)$. (Notasi $A \bowtie \C$
tidak baku.)
\end{defn}

\begin{prop}\label{0015001} Jika $A$ suatu aljabar, maka $A \bowtie \C$ adalah aljabar
beridentitas yang di dalamnya $A$ dibenamkan sebagai (yaitu, isomorfik dengan) suatu ideal
sedemikian sehingga $(A \bowtie \C)/A \cong \C$. Identitas $A \bowtie \C$ adalah $(0,1)$. Jika
$A$ suatu aljabar-$*\,$, maka $A \bowtie \C$ adalah aljabar-$*\,$ beridentitas yang di dalamnya
$A$ merupakan ideal-$*\,$ sedemikian sehingga $(A \bowtie \C)/A \cong \C$.
\end{prop}

\begin{defn} Aljabar $A \bowtie \C$ adalah
 \index{unitalisasi!suatu aljabar-$*\,$}%
\df{unitalisasi} dari aljabar (atau aljabar-$*\,$)~$A$.
\end{defn}

\begin{notn}\label{00150012} Dalam konstruksi sebelumnya, unsur-unsur aljabar
$A \bowtie \C$ secara teknis adalah pasangan terurut $(a,\lambda)$. Unsur-unsur itu biasanya
ditulis secara berbeda. Misalkan $\iota\colon A \sto A \bowtie \C \colon a \mapsto (a,0)$ dan
$\pi\colon A \bowtie \C \sto \C\colon (a,\lambda) \mapsto \lambda$. Karena $(0,1)$ adalah
identitas dalam $A \bowtie \C$, diperoleh
   \begin{align*} (a,\lambda) &= (a,0) + (\vc 0,\lambda) \\
                              &= \iota(a) + \lambda \vc 1_{A \bowtie \C}
   \end{align*}
Sudah lazim untuk memperlakukan $\iota$ sebagai pemetaan inklusi. Jadi, wajar untuk
menulis $(a,\lambda)$ sebagai $a + \lambda \vc 1_{A \bowtie \C}$ atau cukup sebagai
$a + \lambda \vc 1$. Tampaknya tidak timbul ambiguitas dengan menghilangkan acuan pada
identitas perkalian, sehingga notasi baku untuk pasangan $(a,\lambda)$ adalah $a + \lambda$.
\end{notn}

\begin{defn}\label{00150014} Misalkan $A$ suatu aljabar yang \emph{tak beridentitas}. Definisikan
 \index{spektrum!dalam aljabar tak beridentitas}%
\df{spektrum} suatu unsur $a \in A$ sebagai spektrum $a$ ketika dipandang sebagai unsur
dari~$A \bowtie \C$; yaitu, $\sigma_A(a) := \sigma_{A \bowtie \C}(a)$.
\end{defn}

\begin{defn}\label{00150015} Misalkan $A$ suatu aljabar bernorma (dengan atau tanpa involusi). Pada
unitalisasi $A \bowtie \C$ dari $A$, definisikan $\norm{(a,\lambda)} := \norm a + \abs\lambda$.
\end{defn}

\begin{prop}\label{001500151} Misalkan $A$ suatu aljabar bernorma. Pemetaan
 \index{unitalisasi!suatu aljabar bernorma}%
 \index{unitalisasi!suatu aljabar Banach-$*\,$}%
$(a,\lambda) \mapsto \norm a + \abs\lambda$ yang didefinisikan di atas adalah norma yang
menjadikan $A \bowtie \C$ suatu aljabar bernorma. Jika $A$ suatu aljabar Banach (berturut-turut,
aljabar Banach-$*\,$), maka $A \bowtie \C$ adalah aljabar Banach (berturut-turut, aljabar
Banach-$*\,$). Aljabar Banach (atau aljabar Banach-$*\,$) yang dihasilkan adalah
\df{unitalisasi} dari $A$ dan akan
 \index{<unitization@$A_e$ (unitalisasi suatu aljabar Banach)}%
dinotasikan dengan~$A_e$.
\end{prop}

Dengan definisi \emph{spektrum} yang diperluas ini, banyak fakta terdahulu untuk aljabar
Banach beridentitas tetap berlaku dalam tatanan yang lebih umum. Secara khusus, untuk
rujukan selanjutnya, kita nyatakan kembali butir-butir \ref{000661}, \ref{C073134},
\ref{0006703}, dan~\ref{C073157}.

\begin{prop}\label{001500157} Misalkan $a$ suatu unsur dari aljabar Banach~$A$. Maka spektrum
$a$ kompak dan $\abs\lambda \le \norm a$ untuk setiap $\lambda \in \sigma(a)$.
\end{prop}

\begin{prop}\label{00150016} Spektrum setiap unsur dari suatu aljabar Banach tidak kosong.
\end{prop}

\begin{prop}\label{00150017} Untuk setiap unsur $a$ dari suatu aljabar Banach, $\rho(a) \le
\norm a$ dan $\rho(a^n) = \bigr(\rho(a)\bigr)^n$.
\end{prop}

\begin{thm}[Rumus Radius Spektral]\label{00150018} Jika $a$ suatu unsur dari aljabar Banach,
 \index{spektral!radius!rumus untuk}%
maka
   \[ \rho(a) = \inf\bigl\{\norm{a^n}^{1/n}\colon n \in \N\bigr\}
              = \lim_{n \sto \infty} \norm{a^n}^{1/n}. \]
\end{thm}

\begin{defn} Misalkan $A$ suatu aljabar. Suatu ideal kiri $J$ dalam $A$ adalah
 \index{modular!ideal kiri}%
 \index{ideal!kiri modular}%
\df{ideal kiri modular} jika terdapat suatu unsur $u$ dalam $A$ sedemikian sehingga
$au - a \in J$ untuk setiap $a \in A$. Unsur $u$ demikian disebut
 \index{identitas!kanan!terhadap suatu ideal}%
\df{identitas kanan terhadap $J$}. Demikian pula, suatu ideal kanan $J$ dalam $A$ adalah
 \index{modular!ideal kanan}%
 \index{ideal!kanan modular}%
\df{ideal kanan modular} jika terdapat suatu unsur $v$ dalam $A$ sedemikian sehingga
$va - a \in J$ untuk setiap $a \in A$. Unsur $v$ demikian disebut
 \index{identitas!kiri!terhadap suatu ideal}%
\df{identitas kiri terhadap $J$}. Suatu ideal dua sisi $J$ adalah
 \index{modular!ideal}%
 \index{ideal!modular}%
\df{ideal modular} jika terdapat suatu unsur $e$ yang sekaligus merupakan identitas kiri
dan identitas kanan terhadap~$J$.
\end{defn}

\begin{prop} Suatu ideal $J$ dalam suatu aljabar bersifat modular jika dan hanya jika ideal
tersebut merupakan ideal kiri modular sekaligus ideal kanan modular.
\end{prop}

\begin{proof}[\emph{Petunjuk untuk pembuktian}] Tunjukkan bahwa jika $u$ suatu identitas kanan
terhadap $J$ dan $v$ suatu identitas kiri terhadap $J$, maka $vu$ sekaligus merupakan
identitas kanan dan identitas kiri terhadap~$J$. \ns
\end{proof}

\begin{prop} Suatu ideal $J$ dalam suatu aljabar $A$ bersifat modular jika dan hanya jika
aljabar hasil bagi $A/J$ beridentitas.
\end{prop}

\begin{exam} Misalkan $X$ suatu ruang Hausdorff kompak lokal. Untuk setiap $x \in X$, ideal
$J_x$ adalah ideal modular maksimal dalam aljabar-$C^*$ $\fml C_0(X)$ dari fungsi-fungsi
kontinu bernilai kompleks pada~$X$.
\end{exam}

\begin{proof} Menurut versi \emph{lema Urysohn} untuk ruang Hausdorff kompak lokal
(lihat, misalnya, \cite{Erdman:2007}, teorema 17.2.10), terdapat suatu fungsi $g \in
\fml C_0(X)$ sedemikian sehingga $g(x) = 1$. Jadi, $J_x$ bersifat modular karena $g$
merupakan suatu identitas terhadap~$J_x$. Karena $\fml C_0(X) = J_x \oplus\,\, \spn\{g\}$,
ideal $J_x$ mempunyai kodimensi $1$ dan oleh karena itu maksimal.
\end{proof}

\begin{prop} Jika $J$ suatu ideal modular sejati dalam aljabar Banach, maka penutupannya
juga demikian.
\end{prop}

\begin{cor} Setiap ideal modular maksimal dalam suatu aljabar Banach tertutup.
\end{cor}

\begin{prop} Setiap ideal modular sejati dalam suatu aljabar Banach termuat dalam suatu
ideal modular maksimal.
\end{prop}

\begin{prop}\label{00150144} Misalkan $A$ suatu aljabar Banach komutatif dan $\phi \in
\Delta A$. Maka $\ker\phi$ adalah ideal modular maksimal dalam~$A$ dan $A/\ker\phi$ adalah
suatu medan. Lebih lanjut, setiap ideal modular maksimal merupakan kernel dari tepat satu
karakter dalam~$\Delta A$.
\end{prop}

\begin{prop}\label{00150147} Setiap fungsional linear multiplikatif pada suatu aljabar
Banach komutatif $A$ bersifat kontinu. Bahkan, setiap karakter bersifat kontraktif.
\end{prop}

\begin{exam} Misalkan $A_e$ unitalisasi dari suatu aljabar Banach komutatif~$A$ (lihat
proposisi~\ref{001500151}).
 \index{phi@$\phi_\infty$ (suatu karakter pada unitalisasi suatu aljabar)}%
Definisikan
   \[ \phi_\infty\colon A_e \sto \C\colon (a,\lambda) \mapsto \lambda\,. \]
Maka $\phi_\infty$ adalah suatu karakter pada~$A_e$.
\end{exam}

\begin{prop}\label{00150154} Setiap karakter $\phi$ pada suatu aljabar Banach komutatif
$A$ mempunyai perpanjangan tunggal menjadi suatu
 \index{phi@$\phi_e$ (perpanjangan suatu karakter $\phi$ ke unitalisasi suatu aljabar)}%
karakter $\phi_e$ pada unitalisasi $A_e$ dari~$A$. Selain itu, pembatasan pada $A$ dari sebarang
karakter pada $A_e$, dengan pengecualian yang jelas yaitu $\phi_\infty$, adalah suatu
karakter pada~$A$.
\end{prop}

\begin{prop} Jika $A$ suatu aljabar Banach komutatif, maka
 \begin{enumerate}
    \item[(a)] $\Delta A$ adalah ruang Hausdorff kompak lokal,
    \item[(b)] $\Delta A_e = \Delta A \cup \{\phi_\infty\}$,
    \item[(c)] $\Delta A_e$ adalah kompaktifikasi satu titik dari~$\Delta A$, dan
    \item[(d)] pemetaan $\phi \mapsto \phi_e$ adalah homeomorfisme dari $\Delta A$ ke
$\Delta A_e \setminus \{\phi_\infty\}$.
 \end{enumerate}
Jika $A$ beridentitas (sehingga $\Delta A$ kompak), maka $\phi_\infty$ adalah titik
terisolasi dari~$\Delta A_e$.
\end{prop}

\begin{thm} Jika $A$ suatu aljabar Banach komutatif dengan ruang karakter yang tidak kosong,
maka transformasi Gelfand
  \[ \Gamma = \Gamma_A\colon A \sto \fml C_0(\Delta A)\colon a \mapsto \wh a = \Gamma_a \]
adalah suatu homomorfisme aljabar kontraktif dan $\rho(a) = \norm{\wh a}_u$. Lebih lanjut,
jika $A$ \emph{tidak} beridentitas, maka $\sigma(a) = \ran \wh a \cup \{0\}$ untuk setiap
$a \in A$.
\end{thm}

















\section{Barisan Eksak dan Ekstensi}
Prosedur yang cukup sederhana untuk unitalisasi aljabar Banach-$*\,$
(lihat~\ref{00150005} dan~\ref{00150015}) tidak dapat diterapkan begitu saja pada
aljabar-$C^*$. Norma yang didefinisikan dalam~\ref{00150015} tidak memenuhi syarat-$C^*$
(definisi~\ref{0013}). Ternyata unitalisasi aljabar-$C^*$ sedikit lebih rumit. Sebelum
menelaah perinciannya, kita melihat beberapa bahan persiapan mengenai \emph{barisan eksak}
dan \emph{ekstensi} aljabar-$C^*$.

Dalam dunia aljabar-$C^*$, barisan eksak didefinisikan pada dasarnya dengan cara yang sama
seperti barisan eksak aljabar Banach (lihat~\ref{defn_Bsp_exact_seq}).

\begin{defn} Suatu barisan aljabar-$C^*$ dan homomorfisme-$*\,$
   \[\xymatrix{{\cdots}\ar[r] & A_{n-1}\ar[r]^{\phi_n}
          & A_n\ar[r]^{\phi_{n+1}} & A_{n+1}\ar[r] & {\cdots}
   }\]
dikatakan
 \index{barisan!eksak}%
 \index{eksak!barisan}%
\df{eksak di} $A_n$ jika $\ran \phi_n = \ker \phi_{n+1}$. Suatu barisan disebut
\df{eksak} jika barisan tersebut eksak di setiap aljabar-$C^*$ penyusunnya. Suatu barisan
aljabar-$C^*$ dan homomorfisme-$*\,$ berbentuk
   \begin{equation}\label{001500202i2}\xymatrix{
      0 \ar[r] & A \ar[r]^\phi & E\ar[r]^\psi & B\ar[r] & 0
   }\end{equation}
 \index{barisan!eksak pendek}%
 \index{barisan eksak pendek}%
adalah suatu \df{barisan eksak pendek}. (Di sini $\vc 0$ menyatakan aljabar-$C^*$ trivial
berdimensi $0$, dan panah-panah tanpa label adalah homomorfisme-$*$ yang jelas.) Barisan
eksak pendek aljabar-$C^*$~\eqref{001500202i2} adalah
 \index{eksak!barisan!terbelah}%
 \index{barisan!eksak terbelah}%
 \index{eksak terbelah!barisan}%
\df{eksak terbelah} jika terdapat suatu homomorfisme-$*\,$ $\xi\colon B \sto E$ sedemikian
sehingga $\psi\circ\xi = \id B$.
\end{defn}

Definisi-definisi sebelumnya diberikan untuk kategori $\cat{CSA}$ dari aljabar-$C^*$ dan
homomorfisme-$*\,$. Tentu saja, tidak ada yang khusus mengenai kategori tertentu ini.
Barisan eksak bermakna dalam banyak situasi, misalnya dalam berbagai kategori ruang Banach,
aljabar Banach, ruang Hilbert, ruang vektor, grup Abel, modul, dan sebagainya.

Sering kali, dalam konteks aljabar-$C^*$, barisan eksak \eqref{001500202i2} disebut suatu
 \index{ekstensi}%
\df{ekstensi}. Sebagian penulis menyebutnya \emph{ekstensi $A$ oleh $B$} (misalnya,
\cite{Wegge-Olsen:1993} dan \cite{Davidson:1996}), sedangkan penulis lain mengatakan bahwa
barisan itu merupakan \emph{ekstensi $B$ oleh $A$} (\cite{Murphy:1990},
\cite{Fillmore:1996}, dan \cite{Blackadar:2006}). Dalam \cite{Fillmore:1996} dan
\cite{Blackadar:2006}, \emph{ekstensi} didefinisikan sebagai barisan~\eqref{001500202i2};
dalam \cite{Wegge-Olsen:1993}, ekstensi didefinisikan sebagai tripel terurut
$(\phi,E,\psi)$; dan dalam \cite{Murphy:1990} serta \cite{Davidson:1996}, ekstensi
didefinisikan sebagai aljabar-$C^*$ $E$ itu sendiri. Terlepas dari definisi formalnya,
umum dikatakan bahwa \emph{$E$ adalah ekstensi $A$ oleh $B$ (atau $B$ oleh $A$)}.

\begin{conv}\label{00151212} Dalam suatu aljabar-$C^*$, kata \emph{ideal} akan selalu
 \index{konvensi!ideal bersifat swaadjoin}%
 \index{konvensi!ideal bersifat tertutup}%
 \index{ideal!dalam suatu aljabar-$C^*$}%
berarti suatu \hbox{ideal-$*\,$} dua sisi tertutup (kecuali, tentu saja, jika yang
sebaliknya dinyatakan secara eksplisit). Sebentar lagi akan kita tunjukkan (dalam
proposisi~\ref{0019017}) bahwa syarat agar suatu ideal dalam aljabar-$C^*$ bersifat
swaadjoin sebenarnya berlebih. Suatu ideal aljabar dua sisi dari aljabar-$C^*$ yang tidak
harus tertutup akan disebut
 \index{ideal!aljabar}%
 \index{aljabar!ideal}%
\df{ideal aljabar}. Suatu ideal aljabar swaadjoin dari aljabar-$C^*$ akan disebut
 \index{ideal!aljabar-$\ast$}%
 \index{aljabar!ideal-$\ast$}%
\df{ideal-$\ast$ aljabar}.
\end{conv}

Bandingkan dua proposisi berikut dengan~\ref{00151231} dan~\ref{00151232} yang hampir identik.

\begin{prop}\label{00151231c} Tinjau diagram berikut dalam kategori aljabar-$C^*$ dan
homomorfisme-${}^*$
 \[\xymatrix{
      0\ar[r] & A\ar[d]^f\ar[r]^j & B\ar[d]^g\ar[r]^k
              & C\ar@{-->}[d]^h \ar[r] & 0 \\
      0\ar[r] & A'\ar[r]_{j'} & B'\ar[r]_{k'}
      & C'\ar[r] & 0
 }\]
Jika baris-barisnya eksak dan persegi kirinya komutatif, maka terdapat homomorfisme-${}^*$
$h\colon C \sto C'$ tunggal yang membuat persegi kanannya komutatif.
\end{prop}

\begin{prop}[Lema Lima]\label{00151232c}
 \index{lema lima}%
Andaikan dalam diagram aljabar-$C^*$ dan homomorfisme-${}^*$ berikut
 \[\xymatrix{
      0\ar[r] & A\ar[d]^f\ar[r]^j & B\ar[d]^g\ar[r]^k
              & C\ar[d]^h \ar[r] & 0 \\
      0\ar[r] & A'\ar[r]_{j'} & B'\ar[r]_{k'}
      & C'\ar[r] & 0
 }\]
baris-barisnya eksak dan persegi-perseginya komutatif. Buktikan hal-hal berikut.
 \begin{enumerate}
  \item Jika $g$ surjektif, maka $h$ juga surjektif.
  \item Jika $f$ surjektif dan $g$ injektif, maka $h$ injektif.
  \item Jika $f$ dan $h$ surjektif, maka $g$ juga surjektif.
  \item Jika $f$ dan $h$ injektif, maka $g$ juga injektif.
 \end{enumerate}
\end{prop}

\begin{defn} Misalkan $A$ dan $B$ aljabar-$C^*$. Kita mendefinisikan
 \index{jumlah langsung!aljabar-$C^*$}%
 \index{eksternal!jumlah langsung}%
 \index{C*@aljabar-$C^*$!jumlah langsung dari}%
 \index{<binopsum@$A \oplus B$ (jumlah langsung aljabar-$C^*$)}%
\df{jumlah langsung (eksternal)} dari $A$ dan $B$, yang dinotasikan dengan $A \oplus B$,
sebagai hasil kali Cartesius $A \times B$ dengan operasi-operasi aljabar yang didefinisikan
secara titik demi titik dan norma yang diberikan oleh
   \[ \norm{(a,b)} = \max\{\norm a,\norm b\} \]
untuk semua $a \in A$ dan $b \in B$. Notasi alternatif untuk unsur $(a,b)$ dalam
$A \oplus B$
 \index{<binopsum@$a \oplus b$ (suatu unsur dari $A \oplus B$)}%
adalah $a \oplus b$.
\end{defn}

\begin{exam} Misalkan $A$ dan $B$ aljabar-$C^*$. Maka jumlah langsung $A$ dan $B$ adalah
suatu aljabar-$C^*$ dan barisan berikut adalah barisan eksak pendek terbelah:
 \[\xymatrix{
    \vc 0 \ar[r] & A \ar[r]^-{\iota_1} & A \oplus B \ar[r]^-{\pi_2}
             & B \ar[r] & \vc 0
 }\]
\end{exam}
\noindent Pemetaan-pemetaan yang ditunjukkan di atas adalah pemetaan-pemetaan yang jelas:
   \[ \iota_1\colon A \sto A \oplus B\colon a \mapsto (a,\vc 0) \qquad
        \text{dan} \qquad  \pi_2 \colon A \oplus B \sto B: (a,b) \mapsto b\,. \]
 \index{ekstensi!jumlah langsung}%
 \index{jumlah langsung!ekstensi}
Ini adalah \df{ekstensi jumlah langsung}.

\begin{prop} Jika $A$ dan $B$ adalah aljabar-$C^*$ tak nol, maka
 \index{jumlah langsung!aljabar-$C^*$}%
 \index{produk!dalam $\cat{CSA}$}%
 \index{csa@$\cat{CSA}$!produk dalam}%
jumlah langsungnya $A \oplus B$ merupakan produk dalam kategori $\cat{CSA}$ dari
aljabar-$C^*$ dan homomorfisme-$*\,$ (lihat contoh~\ref{0007812}). Jumlah langsung itu
beridentitas jika dan hanya jika $A$ dan $B$ keduanya beridentitas.
\end{prop}

\begin{defn}\label{0015151} Misalkan $A$ dan $B$ aljabar-$C^*$ serta $E$ dan $E'$ merupakan
ekstensi $A$ oleh $B$. Ekstensi-ekstensi ini
 \index{ekuivalensi!kuat!ekstensi}%
 \index{ekstensi!ekuivalensi kuat}%
\df{ekuivalen kuat} jika terdapat suatu isomorfisme-${}^*$ $\theta\colon E \sto E'$
yang membuat diagram
 \[\xymatrix{
     0\ar[r] & A\ar[r]\ar@{=}[d] & E\ar[r]\ar[d]^\theta
             & B\ar[r]\ar@{=}[d] & 0 \\
     0\ar[r] & A\ar[r] & E'\ar[r] & B\ar[r] & 0
 }\]
komutatif.
\end{defn}

\begin{prop}\label{0015152} Dalam definisi sebelumnya, cukup disyaratkan bahwa $\theta$
merupakan suatu homomorfisme-${}^*$.
\end{prop}

\begin{prop}\label{0015155} Misalkan $A$ dan $B$ aljabar-$C^*$. Suatu ekstensi
 \[\xymatrix{
     \vc 0\ar[r] & A\ar[r]^\phi & E\ar[r]^\psi & B\ar[r] & \vc 0
 }\]
ekuivalen kuat dengan ekstensi jumlah langsung $A \oplus B$ jika dan hanya jika terdapat\newline
suatu \mbox{homomorfisme-${}^*$} $\nu\colon E \sto A$ sedemikian sehingga
$\nu\circ\phi = \id A$.
\end{prop}

\begin{prop}\label{0015157} Jika barisan-barisan aljabar-$C^*$
 \begin{equation}\label{0015157i}\xymatrix{
     0\ar[r] & A\ar[r]^\phi & E\ar[r]^\psi & B\ar[r] & 0
 }\end{equation}
dan
 \begin{equation}\label{0015157ii}\xymatrix{
     0\ar[r] & A\ar[r]^{\phi'} & E'\ar[r]^{\psi'} & B\ar[r] & 0
 }\end{equation}
ekuivalen kuat dan \eqref{0015157i} terbelah, maka \eqref{0015157ii} juga terbelah.
\end{prop}

















\section{Unitalisasi Aljabar-$C^*$}
Prosedur yang cukup sederhana untuk unitalisasi aljabar Banach-$*\,$
(lihat~\ref{00150005} dan~\ref{00150015}) tidak dapat diterapkan begitu saja pada
aljabar-$C^*$. Norma yang didefinisikan dalam~\ref{00150015} tidak memenuhi syarat-$C^*$
(definisi~\ref{0013}). Ternyata unitalisasi aljabar-$C^*$ sedikit lebih rumit.

\begin{prop} Kernel dari suatu homomorfisme-$*\,$ $\phi\colon A \sto B$ antara
aljabar-$C^*$ adalah suatu ideal-$*\,$ tertutup dalam~$A$ dan jangkauannya adalah suatu
subaljabar-$*\,$ dari~$B$.
\end{prop}

\begin{prop}\label{001315} Jika $a$ suatu unsur dalam aljabar-$C^*$, maka
  \begin{align*}
    \norm{a} &= \sup\{\norm{xa}\colon \norm{x} \le 1\} \\
             &= \sup\{\norm{ax}\colon \norm{x} \le 1\}.
  \end{align*}
\end{prop}

\begin{cor}\label{0013151} Jika $a$ suatu unsur dari aljabar-$C^*$ $A$, operator $L_a$,
yang disebut
 \index{kiri!operator perkalian}%
 \index{operator!perkalian kiri}%
\df{perkalian kiri oleh~$a$} dan didefinisikan oleh
   \[ L_a\colon A \sto A\colon x \mapsto ax \]
adalah suatu operator (linear terbatas) pada~$A$. Lebih lanjut, pemetaan
   \[ L\colon A \sto \ofml B(A)\colon a \mapsto L_a \]
sekaligus merupakan isometri dan homomorfisme aljabar.
\end{cor}

\begin{prop}\label{0015213} Misalkan $A$ suatu aljabar-$C^*$. Maka terdapat suatu
aljabar-$C^*$ beridentitas $\wt A$ yang di dalamnya $A$ dibenamkan sebagai suatu ideal
sedemikian sehingga barisan
 \begin{equation}\label{0015213i}
   \xymatrix{
     \vc 0\ar[r] & A\ar[r] & \wt A \ar[r] & \C \ar[r] & \vc 0
 }\end{equation}
eksak terbelah. Jika $A$ beridentitas, maka barisan \eqref{0015213i} ekuivalen kuat dengan
ekstensi jumlah langsung, sehingga $\wt A \cong A \oplus \C$. Jika $A$ \emph{tidak}
beridentitas, maka $\wt A$ tidak isomorfik dengan $A \oplus \C$.
\end{prop}

\begin{proof}[\emph{Petunjuk untuk pembuktian}] Pembuktian hasil ini agak rumit. Setiap orang
sebaiknya menelusuri semua perinciannya setidaknya sekali dalam hidupnya. Berikut ini
adalah garis besar suatu pembuktian.

Perhatikan bahwa kita berbicara tentang unitalisasi dari suatu aljabar-$C^*$ $A$
\emph{baik $A$ sudah memiliki unit (identitas perkalian) maupun belum}. Kita membagi
argumen menjadi dua kasus.

\vskip 8 pt

\noindent\textbf{Kasus 1: aljabar $A$ beridentitas.}
\begin{enumerate}
    \item Pada aljabar $A \bowtie \C$, definisikan
       \[ \norm{(a,\lambda)} := \max\{\norm{a + \lambda \vc 1_A}, \abs\lambda\} \]
dan misalkan $\wt A := A \bowtie \C$ bersama dengan norma ini.
    \item Buktikan bahwa pemetaan $(a,\lambda) \mapsto \norm{(a,\lambda)}$ adalah norma
pada~$A \bowtie \C$.
    \item Buktikan bahwa norma ini adalah norma aljabar.
    \item Tunjukkan bahwa norma ini, sesungguhnya, merupakan norma-$C^*$ pada
$A \bowtie \C$.
    \item Amati bahwa norma ini merupakan perpanjangan dari norma pada~$A$.
    \item Buktikan bahwa $A \bowtie \C$ adalah suatu aljabar-$C^*$ dengan memverifikasi
kelengkapan ruang metrik yang diinduksi oleh norma sebelumnya.
    \item Buktikan bahwa barisan
       \[ \vc 0 \to A \to^\iota \wt A \two/->`<-/^Q_\psi \C \to \vc 0  \]
eksak terbelah (dengan $\iota\colon a \mapsto (a,0)$, $Q\colon (a,\lambda) \mapsto
\lambda$, dan $\psi\colon \lambda \mapsto (0,\lambda)$).
    \item Buktikan bahwa $\cstariso{\wt A}{A \oplus \C}$.
\end{enumerate}

\vskip 8 pt

\noindent\textbf{Kasus 2: aljabar $A$ \emph{tidak} beridentitas.}
\begin{enumerate}
\setcounter{enumi}{8}
    \item Buktikan fakta sederhana berikut.
  \begin{lem} Misalkan $A$ suatu aljabar dan $B$ suatu aljabar bernorma. Jika
  $\phi\colon A \sto B$ adalah homomorfisme aljabar, maka fungsi
  $a \mapsto \norm a := \norm{\phi(a)}$ adalah seminorma pada~$A$. Fungsi tersebut adalah
  norma jika $\phi$ injektif.
  \end{lem}
    \item Ingat bahwa dalam~\ref{0013151} kita mendefinisikan operator $L_a$,
\emph{perkalian kiri oleh $a$}. Sekarang, misalkan
      \[ A^{\sharp} :=
        \{ L_a + \lambda I_A \in \ofml B(A)\colon a \in A \text{ dan } \lambda \in \C\} \]
dan tunjukkan bahwa $A^\sharp$ adalah suatu aljabar bernorma.
    \item Jadikan $A^\sharp$ suatu aljabar-$*\,$ dengan mendefinisikan
      \[ (L_a + \lambda I_A)^* := L_{a^*} + \conj\lambda I_A \]
untuk semua $a \in A$ dan $\lambda \in \C$.
    \item Definisikan
      \[ \phi\colon A \bowtie \C \sto A^\sharp\colon
                   (a,\lambda) \mapsto L_a + \lambda I_A \]
dan verifikasikan bahwa $\phi$ adalah suatu homomorfisme-$*\,$.
    \item Buktikan bahwa $\phi$ injektif.
    \item Gunakan (9) untuk melengkapi $A \bowtie \C$ dengan suatu norma yang
menjadikannya aljabar bernorma beridentitas. Misalkan $\wt A := A \bowtie \C$ dengan
    norma yang ditarik balik oleh $\phi$ dari $A^\sharp$.


    \item Verifikasikan fakta-fakta berikut.
      \begin{enumerate}
        \item Pemetaan $\phi\colon \wt A \sto A^\sharp$ adalah isomorfisme isometrik.
        \item $\ran L$ adalah subaljabar tertutup dari
$\ran \phi = A^\sharp \subseteq \ofml B(A)$.
        \item $I_A \notin \ran L$.
      \end{enumerate}
    \item Buktikan bahwa norma pada $\wt A$ memenuhi syarat-$C^*$.
    \item Buktikan bahwa $\wt A$ adalah suatu aljabar-$C^*$ beridentitas. (Untuk
menunjukkan bahwa $\wt A$ lengkap, kita hanya perlu menunjukkan bahwa $A^\sharp$ lengkap.
Untuk tujuan ini, misalkan $\bigl(\phi(a_n,\lambda_n)\bigr)_{n=1}^\infty$ suatu barisan
Cauchy dalam~$A^\sharp$. Untuk menunjukkan bahwa barisan ini konvergen, cukup ditunjukkan
bahwa barisan itu mempunyai suatu subbarisan konvergen. Dengan menunjukkan bahwa barisan
$(\lambda_n)$ terbatas, kita dapat mengekstrak darinya suatu subbarisan konvergen
$\bigl(\lambda_{n_k}\bigr)$. Buktikan bahwa $(L_{a_{n_k}})$ konvergen.)
    \item Buktikan bahwa barisan
        \begin{equation}\xymatrix{
     0\ar[r] & A\ar[r] & \wt A\ar[r] & \C\ar[r] & 0
        }\end{equation}
eksak terbelah.
    \item Aljabar-$C^*$ $\wt A$ \emph{tidak} isomorfik dengan $A \oplus \C$.
\end{enumerate}  \ns
\end{proof}

\begin{defn}\label{00152131} Aljabar-$C^*$ $\wt A$ yang dikonstruksi dalam proposisi
sebelumnya adalah
 \index{unitalisasi!suatu aljabar-$C^*$}%
\df{unitalisasi} dari~$A$.
\end{defn}

Perhatikan bahwa definisi \emph{spektrum} yang diperluas dalam~\ref{00150014} berlaku untuk
aljabar-$C^*$ karena penambahan identitas merupakan perkara aljabar murni dan sama
bagi aljabar-$C^*$ seperti bagi aljabar Banach umum. Dengan demikian, banyak fakta
terdahulu yang dinyatakan untuk aljabar-$C^*$ beridentitas tetap berlaku. Secara khusus,
untuk rujukan selanjutnya, kita nyatakan kembali butir-butir \ref{00131101},
\ref{00131102}, \ref{00131103}, \ref{00131104}, \ref{001312}, \ref{001406},
dan~\ref{0014062}.

\begin{prop}\label{00152141} Misalkan $a$ suatu unsur normal dari suatu aljabar-$C^*$. Maka
$\norm{a^2} = {\norm a}^2$ dan karena itu $\rho(a)~=~\norm a$.
\end{prop}

\begin{cor}\label{00152142} Jika $A$ suatu aljabar-$C^*$ komutatif, maka
$\norm{a^2} = {\norm a}^2$ dan $\rho(a)~=~\norm a$ untuk setiap $a \in A$.
\end{cor}

\begin{cor}\label{00152143} Pada suatu aljabar-$C^*$ komutatif $A$, transformasi Gelfand
$\Gamma$ adalah isometri; yaitu, $\norm{\Gamma_a}_u = \norm{\wh a}_u = \norm a$ untuk
setiap $a \in A$.
\end{cor}

\begin{cor}\label{00152144} Norma suatu aljabar-$C^*$ bersifat tunggal dalam arti bahwa, jika
diberikan suatu aljabar $A$ dengan involusi, paling banyak ada satu norma yang menjadikan
$A$ suatu aljabar-$C^*$.
\end{cor}

\begin{prop}\label{00152145} Jika $h$ suatu unsur swaadjoin dari suatu aljabar-$C^*$,
 \index{spektrum!unsur swaadjoin}%
 \index{swaadjoin!unsur!spektrum}%
maka $\sigma(h) \subseteq \R$.
\end{prop}

\begin{prop}\label{00152152} Jika $a$ suatu unsur swaadjoin dalam suatu aljabar-$C^*$,
maka transformasi Gelfandnya $\wh a$ bernilai real.
\end{prop}

\begin{prop}\label{00152153} Setiap karakter pada suatu aljabar-$C^*$ $A$ mempertahankan
involusi; dengan demikian, transformasi Gelfand $\sbsb{\Gamma}{\!\!A}$ adalah suatu
homomorfisme-$*$.
\end{prop}

Akibat langsung dari hasil-hasil sebelumnya adalah versi kedua dari
\emph{teorema Gelfand--Naimark}, yang menyatakan bahwa setiap aljabar-$C^*$ komutatif
(isomorfik-$*\,$ secara isometrik dengan) aljabar dari semua fungsi kontinu pada suatu ruang
Hausdorff kompak lokal yang lenyap di tak hingga. Seperti halnya pada versi pertama
teorema ini (lihat~\ref{00142}), ruang Hausdorff kompak lokal yang dimaksud adalah ruang
karakter dari aljabar tersebut.

\begin{thm}[Teorema Gelfand--Naimark II]\label{00152156} Misalkan $A$ suatu aljabar-$C^*$
komutatif.
 \index{Gelfand-Naimark!Teorema II}%
Maka transformasi Gelfand $\Gamma_A\colon a \mapsto \wh a$ adalah suatu
isomorfisme-$*\,$ isometrik dari $A$ pada $\fml C_0(\Delta A)$.
\end{thm}

Dari proposisi berikut, diperoleh bahwa proses unitalisasi bersifat funktorial.

\begin{prop} \label{00152161} Setiap homomorfisme-$*$ $\phi\colon A \sto B$ antara
aljabar-$C^*$ mempunyai suatu perpanjangan tunggal menjadi homomorfisme-$*$ beridentitas
$\wt \phi\colon \wt A \sto \wt B$ antara unitalisasi masing-masing.
\end{prop}

Berbeda dengan keadaan pada aljabar Banach umum, tidak ada perbedaan antara kategori
topologis dan kategori geometris aljabar-$C^*$. Salah satu aspek paling mengagumkan dari
teori aljabar-$C^*$ adalah bahwa homomorfisme-$*$ antara aljabar-aljabar demikian secara
otomatis kontinu---bahkan, kontraktif. Akibatnya, jika dua aljabar-$C^*$ isomorfik-$*$
secara aljabar, maka keduanya isomorfik secara isometrik.

\begin{prop}\label{00152171} Setiap homomorfisme-$*$ antara aljabar-$C^*$ bersifat kontraktif.
\end{prop}

\begin{prop}\label{00152181} Setiap homomorfisme-$*$ injektif antara aljabar-$C^*$ adalah
suatu isometri.
\end{prop}

\begin{prop}\label{00152191} Misalkan $X$ suatu ruang Hausdorff kompak lokal dan
$\wt X = X \cup \{\infty\}$ merupakan kompaktifikasi satu titiknya. Definisikan
  \[ \iota\colon \fml C_0(X) \sto \,\fml C(\wt X)\colon f \mapsto \wt f \]
dengan
 \[ \wt f(x) = \left\{%
  \begin{array}{ll}
     f(x), & \hbox{jika $x \in X$;} \\
     0, & \hbox{jika $x = \infty$.} \\
  \end{array}%
            \right. \]
Definisikan pula $E_\infty$ pada $\fml C(\wt X)$ dengan $E_\infty(g) = g(\infty)$. Maka
barisan
  \[ \vc 0 \to \fml C_0(X) \to^\iota \,\,\fml C(\wt X)
                                        \to^{E_\infty} \C \to \vc 0 \]
eksak.
\end{prop}

Dalam proposisi sebelumnya, kita menyebut $\wt X$ sebagai kompaktifikasi satu titik dari $X$
bahkan jika $X$ sejak awal kompak. Kebanyakan definisi \emph{kompaktifikasi} mensyaratkan
suatu ruang rapat dalam setiap kompaktifikasinya. (Lihat ulasan saya pada awal
bagian 17.3 dari~\cite{Erdman:2007}.) Sebelumnya kita telah menggunakan konvensi bahwa
unitalisasi dari suatu aljabar beridentitas memperoleh suatu identitas perkalian baru.
Demi konsistensi dengan pilihan ini, selanjutnya kita akan menganut konvensi bahwa
kompaktifikasi satu titik dari suatu ruang kompak memperoleh suatu titik tambahan
(terisolasi). (Lihat pula istilah yang kita perkenalkan
dalam~\futurexref{14.3.1}{0038614}.)

 \vskip 2 pt

Dari sudut pandang teorema Gelfand--Naimark~(\ref{00152156}), wawasan mendasar yang
disarankan oleh proposisi berikut adalah bahwa unitalisasi suatu aljabar-$C^*$ komutatif,
dalam arti tertentu, merupakan ``hal yang sama'' dengan kompaktifikasi satu titik dari
suatu ruang Hausdorff kompak lokal.

\begin{prop}\label{00152194} Jika $X$ suatu ruang Hausdorff kompak lokal, maka
aljabar-$C^*$ beridentitas $\bigl(\fml C_0(X)\bigr)^\sim$ dan $\fml C(\wt X)$
isomorfik-$*\,$ secara isometrik.
\end{prop}

\begin{proof} Definisikan
  \[ \theta\colon \bigl(\fml C_0(X)\bigr)^\sim \!\!\sto \fml C(\wt X)\colon
                         (f,\lambda) \mapsto \wt f + \lambda \vc 1_{\wt X} \]
(dengan $\vc 1_{\wt X}$ menyatakan fungsi konstan $1$ pada~$\wt X$). Kemudian tinjau diagram
  \[\xymatrix{
     0\ar[r] & \fml C_0(X)\ar[r]\ar@{=}[d] & \bigl(\fml C_0(X)\bigr)^\sim\ar[r]\ar[d]^\theta
             & \C\ar[r]\ar@{=}[d] & 0 \\
     0\ar[r] & \fml C_0(X)\ar[r]^\iota & \fml C(\wt X)\ar[r]^{E_\infty} & \C\ar[r] & 0
  }\]
Baris atas eksak menurut proposisi~\ref{0015213}, baris bawah eksak menurut
proposisi~\ref{00152191}, dan diagram tersebut jelas komutatif. Pemeriksaan bahwa $\theta$
adalah suatu homomorfisme-$*\,$ bersifat rutin. Oleh karena itu, $\theta$ adalah suatu
isomorfisme-$*\,$ isometrik menurut proposisi~\ref{0015152} dan
korolari~\ref{00152181}.
\end{proof}


















\section{Kuasi-Invers}
\begin{defn} Suatu elemen $b$ dari aljabar $A$ adalah \df{kuasi-invers kiri} bagi $a \in A$ jika
$ba = a + b$. Elemen tersebut adalah \df{kuasi-invers kanan} bagi $a$ jika $ab = a + b$. Jika $b$
sekaligus merupakan kuasi-invers kiri dan kuasi-invers kanan bagi $a$, maka $b$ adalah
 \index{kuasi-invers}%
\df{kuasi-invers} bagi~$a$. Jika $a$ memiliki kuasi-invers, kita lambangkan kuasi-invers itu dengan~$a'$.
\end{defn}

\begin{prop} Jika $b$ adalah kuasi-invers kiri bagi $a$ dalam aljabar $A$ dan $c$ adalah
kuasi-invers kanan bagi $a$ dalam $A$, maka $b = c$.
\end{prop}

\begin{prop} Misalkan $A$ aljabar beridentitas dan $a$, $b \in A$. Maka
 \begin{enumerate}
  \item[(a)] $a'$ ada jika dan hanya jika $(\vc 1 - a)^{-1}$ ada; dan
  \item[(b)] $b^{-1}$ ada jika dan hanya jika $(\vc 1 - b)'$ ada.
 \end{enumerate}
\end{prop}

\begin{proof}[\emph{Petunjuk bukti}] Untuk arah sebaliknya pada (a), tinjau
$a + c - ac$, dengan $c = \vc 1 - (\vc 1 - a)^{-1}$. \ns
\end{proof}

\begin{prop} Suatu elemen $a$ dari aljabar Banach $A$ bersifat kuasi-invertibel jika dan hanya jika
elemen itu bukan merupakan identitas terhadap ideal modular maksimal mana pun dalam~$A$.
\end{prop}

\begin{prop} Misalkan $A$ aljabar Banach dan $a \in A$. Jika $\rho(a) < 1$, maka $a'$ ada
dan $a' = -\sum_{k=1}^\infty a^k$.
\end{prop}

\begin{prop} Misalkan $A$ aljabar Banach dan $a \in A$. Jika $\norm{a} < 1$, maka $a'$ ada
dan
   \[ \frac{\norm{a}}{1 + \norm{a}} \le \norm{a'} \le \frac{\norm{a}}{1 - \norm{a}}. \]
\end{prop}

\begin{prop} Misalkan $A$ aljabar Banach beridentitas dan $a \in A$. Jika $\rho(\vc 1 - a) < 1$, maka
$a$ invertibel dalam~$A$.
\end{prop}

\begin{prop} Misalkan $A$ aljabar Banach dan $a$, $b \in A$. Jika $b'$ ada dan
$\norm{a} < (1 + \norm{b'})^{-1}$, maka $(a + b)'$ ada dan
  \[ \norm{(a + b)' - b'}
           \le \frac{\norm{a}\,(1 + \norm{b'})^2}{1 - \norm{a}\,(1 + \norm{b'})}\,.\]
\end{prop}
 \begin{proof}[\emph{Petunjuk bukti}] Tunjukkan terlebih dahulu bahwa $u = (a - b'a)'$ ada dan bahwa
 $u + b' - ub'$ adalah kuasi-invers kiri bagi~$a~+~b$. \ns
\end{proof}

Bandingkan hasil berikut dengan proposisi~\ref{000644fa} dan~\ref{000646fa}.
\begin{prop} Himpunan $Q_A$ yang terdiri atas
 \index{q@$Q_A$ (elemen kuasi-invertibel)}%
elemen-elemen kuasi-invertibel dari aljabar Banach $A$ adalah himpunan terbuka dalam $A$, dan pemetaan $a \mapsto
a'$ merupakan homeomorfisme dari $Q_A$ ke dirinya sendiri.
\end{prop}

\begin{notn} Jika $a$ dan $b$ adalah elemen dalam suatu aljabar, kita definisikan
 \index{<binop@$a \circ b$ (operasi dalam suatu aljabar)}%
$a \circ b := a + b - ab$.
\end{notn}

\begin{prop} Jika $A$ aljabar, maka $Q_A$ merupakan grup terhadap~$\circ$.
\end{prop}

\begin{prop} Jika $A$ aljabar Banach beridentitas dan $\inv A$ adalah himpunan elemen-elemen
invertibelnya, maka
   \[ \Psi\colon Q_A \sto \inv A\colon  a \mapsto \vc 1 - a \]
merupakan isomorfisme.
\end{prop}

\begin{defn} Jika $A$ aljabar dan $a \in A$, kita definisikan
 \index{spektrum-q}%
 \index{spektrum!q-}%
\df{spektrum-q} dari $a$ dalam $A$ dengan
  \[ \breve{\sigma}_A(a) := \{\lambda \in \C\colon  \lambda \ne 0
                      \text{ dan } \dfrac a\lambda \notin Q_A\}\cup\{0\}. \]
\end{defn}

Dalam aljabar beridentitas, definisi di atas ``hampir'' merupakan definisi yang biasa.
\begin{prop} Misalkan $A$ aljabar beridentitas dan $a \in A$. Maka untuk setiap $\lambda \ne 0$, berlaku
$\lambda \in \sigma(a)$ jika dan hanya jika $\lambda \in \breve{\sigma}(a)$.
\end{prop}

\begin{prop} Jika aljabar $A$ tak beridentitas dan $a \in A$, maka $\breve{\sigma}_A(a) =
\breve{\sigma}_{\wt A}(a)$, dengan $\wt A$ adalah unitalisasi dari~$A$.
\end{prop}























\section{Elemen Positif dalam Aljabar-$C^*$}
\begin{defn} Suatu elemen swaadjoin $a$ dari aljabar-$C^*$ $A$ disebut
 \index{positif!elemen dari suatu aljabar-$C^*$}%
\df{positif} jika $\sigma(a) \subseteq [0,\infty)$. Dalam hal ini kita tulis $a \ge \vc 0$.
Kita lambangkan himpunan semua elemen positif dari $A$
 \index{<plus@$A^+$ (kerucut positif)}%
dengan~$A^+$. Himpunan ini adalah
 \index{positif!kerucut}%
 \index{kerucut!positif}%
\df{kerucut positif} dari~$A$. Untuk setiap himpunan bagian $B$ dari $A$, misalkan $B^+ = B \cap A^+$. Kita akan
menggunakan kerucut positif untuk menginduksi suatu urutan parsial pada~$A$: kita tulis
 \index{<relation@$a \le b$ (relasi urutan dalam suatu aljabar-$C^*$)}%
$a \le b$ apabila $b - a \in A^+$.
\end{defn}

\begin{defn} Misalkan $\le$ suatu relasi pada himpunan tak kosong $S$. Jika relasi
 \index{<relation@$\le$ (notasi untuk suatu urutan parsial)}%
 $\le$ bersifat refleksif dan transitif, relasi tersebut adalah suatu
 \index{praurutan}%
 \index{urutan!pra-}%
\df{praurutan}. Jika $\le$ adalah praurutan dan juga antisimetris, relasi tersebut adalah suatu
 \index{parsial!urutan}%
 \index{urutan!parsial}%
\df{urutan parsial}.

Suatu urutan parsial $\le$ pada ruang vektor riil $V$
 \index{kompatibilitas!urutan dengan operasi}%
\df{kompatibel} dengan (atau
 \index{menghormati operasi}%
\df{menghormati}) operasi-operasi (penjumlahan dan perkalian skalar) pada~$V$ jika untuk semua $x$,
$y$, $z \in V$
 \begin{enumerate}
  \item[(a)] $x \le y$ mengakibatkan $x+z \le y+z$, dan
  \item[(b)] $x \le y$, $\alpha \ge 0$ mengakibatkan $\alpha x \le \alpha y$.
 \end{enumerate}
Suatu ruang vektor riil yang dilengkapi dengan urutan parsial yang kompatibel dengan operasi-operasi
ruang vektor disebut
 \index{terurut!ruang vektor}%
 \index{vektor!ruang!terurut}%
\df{ruang vektor terurut}.
\end{defn}

\begin{defn} Misalkan $V$ ruang vektor. Suatu himpunan bagian $C$ dari $V$ adalah
 \index{kerucut}%
\df{kerucut} dalam $V$ jika $\alpha C \subseteq C$ untuk setiap $\alpha \ge 0$. Suatu kerucut $C$ dalam $V$
disebut
 \index{proper!kerucut}%
 \index{kerucut!proper}%
\df{proper} jika $C \cap (-C) = \{\vc 0\}$.
\end{defn}

\begin{prop}\label{0018006} Suatu kerucut $C$ dalam ruang vektor bersifat konveks jika dan hanya jika $C + C
\subseteq C$.
\end{prop}

\begin{exam}\label{0018007} Jika $V$ ruang vektor terurut, maka himpunan
 \index{<dualp@$V^+$ (kerucut positif dalam suatu ruang vektor terurut)}%
   \[ V^+ := \{x \in V \colon x \ge \vc 0\} \]
adalah kerucut konveks proper dalam~$V$. Himpunan ini adalah
 \index{positif!kerucut!dalam suatu ruang vektor terurut}%
 \index{kerucut!positif}%
\df{kerucut positif} dari~$V$, dan anggota-anggotanya adalah
 \index{positif!elemen (dari suatu ruang vektor terurut)}%
\df{elemen positif} dari~$V$.
\end{exam}

\begin{prop}\label{0018008} Misalkan $V$ ruang vektor riil dan $C$ kerucut konveks proper
dalam~$V$. Definisikan $x \le y$ jika $y - x \in C$. Maka relasi $\le$ adalah urutan parsial
pada $V$ dan kompatibel dengan operasi-operasi ruang vektor pada~$V$. Relasi ini adalah
\df{urutan parsial
 \index{terinduksi!urutan parsial}%
 \index{parsial!urutan!terinduksi oleh kerucut konveks proper}%
yang diinduksi oleh} kerucut~$C$. Kerucut positif $V^+$ dari ruang vektor terurut yang dihasilkan
tepat sama dengan $C$.
\end{prop}

\begin{prop} Jika $a$ adalah elemen swaadjoin dari suatu aljabar-$C^*$ beridentitas dan $t \in \R$, maka
 \begin{enumerate}
   \item[(i)] $a \ge \vc 0$ apabila $\norm{a - t\vc 1} \le t$; dan
   \item[(ii)] $\norm{a - t\vc 1} \le t$ apabila $\norm a \le t$ dan $a \ge \vc 0$.
  \end{enumerate}
\end{prop}

\begin{exam} Kerucut positif dari suatu aljabar-$C^*$ $A$ adalah kerucut konveks proper tertutup
dalam ruang vektor riil~$\ofml H(A)$.
\end{exam}

\begin{prop}\label{0018015} Jika $a$ dan $b$ adalah elemen positif dari suatu aljabar-$C^*$
dan $ab = ba$, maka $ab$ positif.
\end{prop}

\begin{prop}\label{001804} Setiap elemen positif dari suatu aljabar-$C^*$ $A$ memiliki akar pangkat
$n$ positif yang tunggal ($n \in \N$). Artinya, jika $a \in A^+$, maka terdapat tepat satu
$b \in A^+$ sedemikian sehingga $b^n = a$.
\end{prop}

\begin{proof} \emph{Petunjuk.} Bagian eksistensi merupakan penerapan sederhana kalkulus fungsional
$C^*$ (yaitu, \emph{teorema spektral abstrak}~\ref{00144}). Elemen $b$
yang diberikan oleh kalkulus fungsional bersifat positif dalam aljabar $C^*(\vc 1,a)$. Jelaskan mengapa
elemen itu juga positif dalam~$A$. Argumen keunikan memerlukan perhatian yang saksama. \ns
\end{proof}

\begin{thm}[Dekomposisi Jordan]\label{001807} Jika $c$ adalah elemen swaadjoin dari suatu
 \index{dekomposisi Jordan}%
 \index{dekomposisi!Jordan}%
aljabar-$C^*$ $A$, maka terdapat elemen-elemen positif
 \index{<positive@$c^+$ (bagian positif dari $c$)}%
 \index{positif!bagian!dari suatu elemen swaadjoin}%
 \index{<positive@$c^-$ (bagian negatif dari $c$)}%
 \index{negatif!bagian!dari suatu elemen swaadjoin}%
$c^+$ dan $c^-$ dalam $A$ sedemikian sehingga $c = c^+ - c^-$ dan $c^+c^- = \vc 0$.
\end{thm}

\begin{lem} Jika $c$ adalah elemen dari suatu aljabar-$C^*$ sedemikian sehingga $-c^*c \ge \vc 0$, maka $c
= \vc 0$.
\end{lem}

\begin{prop} Jika $a$ adalah elemen dari suatu aljabar-$C^*$, maka $a^*a \ge \vc 0$.
\end{prop}

\begin{prop}\label{00180731} Jika $c$ adalah elemen dari suatu aljabar-$C^*$ $A$, maka pernyataan-pernyataan
berikut ekuivalen:
 \begin{enumerate}
    \item[(i)] $c \ge \vc 0$;
    \item[(ii)] terdapat $b \ge \vc 0$ sedemikian sehingga $c = b^2$; dan
    \item[(iii)] terdapat $a \in A$ sedemikian sehingga $c = a^*a$,
 \end{enumerate}
\end{prop}

\begin{exam}\label{00180741} Jika $T$ operator pada ruang Hilbert $H$, maka $T$ merupakan
elemen positif dari aljabar-$C^*$ $\ofml B(H)$ jika dan hanya jika $\langle Tx,x \rangle
\ge 0$ untuk setiap $x \in H$.
\end{exam}

\begin{proof}[\emph{Petunjuk bukti}] Mudah untuk menunjukkan bahwa jika $T \ge \vc 0$ dalam $\ofml B(H)$, maka
$\langle Tx,x \rangle \ge 0$ untuk setiap $x \in H$: gunakan proposisi~\ref{00180731} untuk
menuliskan $T$ sebagai $S^*S$ bagi suatu operator~$S$.

Sebaliknya, andaikan $\langle Tx,x \rangle \ge 0$ untuk setiap $x \in H$. Mudah
dilihat bahwa ini mengakibatkan $T$ swaadjoin. Gunakan \emph{teorema dekomposisi
Jordan}~\ref{001807} untuk menuliskan $T$ sebagai $T^+ - T^-$. Untuk sembarang $u \in H$, misalkan $x = T^-u$
dan periksa bahwa $0 \le \langle Tx,x \rangle = -\langle (T^-)^3u,u \rangle$. Sekarang $(T^-)^3$
adalah elemen positif dari $\ofml B(H)$. (Mengapa?) Simpulkan bahwa $(T^-)^3 = \vc 0$ dan
karena itu $T^- = \vc 0$. (Untuk rincian tambahan, lihat~\cite{DoranB:1986}, halaman~37.) \ns
\end{proof}

\begin{defn} Untuk sembarang elemen $a$ dari suatu aljabar-$C^*$, kita definisikan
 \index{mutlak!nilai!dalam suatu aljabar-$C^*$}%
 \index{<unop@$\abs a$ (nilai mutlak dalam suatu aljabar-$C^*$)}%
$\abs a$ sebagai $\sqrt{a^*a}$.
\end{defn}

\begin{prop} Jika $a$ adalah elemen swaadjoin dari suatu aljabar-$C^*$, maka
  \[ \abs a = a^+ + a^-\,. \]
\end{prop}

\begin{exam} Nilai mutlak dalam suatu aljabar-$C^*$ tidak harus
subaditif; yaitu, $\abs{a + b}$ tidak harus lebih kecil daripada $\abs{a}
+ \abs{b}$. Sebagai contoh, dalam $\M 2\C$, ambil $a = \begin{bmatrix} 1 & 1 \\
1 & 1 \end{bmatrix}$ dan $b = \begin{bmatrix} 0 & 0 \\ 0 & -2
\end{bmatrix}$. Maka $\abs{a} = a$, $\abs{b} = \begin{bmatrix} 0 &
0 \\ 0 & 2 \end{bmatrix}$, dan $\abs{a + b} = \begin{bmatrix} \sqrt{2} & 0 \\ 0 &
\sqrt{2}
\end{bmatrix}$. Jika $\abs{a + b} - \abs{a} - \abs{b}$ positif, maka, berdasarkan
contoh~\ref{00180741}, $\langle\,(\abs{a + b} - \abs{a} - \abs{b})\,x\,,\,x\,\rangle$
akan positif untuk setiap vektor $x \in \C^2$. Namun, hal ini tidak benar untuk $x = (1,0)$.
\end{exam}

\begin{prop}\label{0018201} Jika $\phi\colon A \sto B$ adalah homomorfisme-$*\,$ antara
aljabar-$C^*$, maka $\phi(a) \in B^+$ apabila $a \in A^+$. Jika $\phi$ adalah
isomorfisme-$*\,$, maka $\phi(a) \in B^+$ jika dan hanya jika $a \in A^+$.
\end{prop}

\begin{prop} Misalkan $a$ elemen swaadjoin dari suatu aljabar-$C^*$ $A$ dan $f$ fungsi kontinu
bernilai kompleks pada spektrum~$a$. Maka $f \ge \vc 0$ dalam $\fml
C(\sigma(a))$ jika dan hanya jika $f(a) \ge \vc 0$ dalam~$A$.
\end{prop}

\begin{prop}\label{0018203} Jika $a$ adalah elemen swaadjoin dari suatu aljabar-$C^*$ $A$, maka
$\norm a\,\vc 1_{\wt A} \pm a \ge \vc 0$ dalam~$\wt A$.
\end{prop}

\begin{prop}\label{0018204} Jika $a$ dan $b$ adalah elemen swaadjoin dari suatu
aljabar-$C^*$~$A$ dan $a \le b$, maka $x^*ax \le x^*bx$ untuk setiap $x \in A$.
\end{prop}

\begin{prop} Jika $a$ dan $b$ adalah elemen dari suatu aljabar-$C^*$ dengan $0 \le a \le b$, maka
$\norm a \le \norm b$.
\end{prop}

\begin{prop}\label{0018301} Misalkan $A$ aljabar-$C^*$ beridentitas dan $c \in A^+$. Maka
$c$ invertibel jika dan hanya jika $c \ge \epsilon\vc 1$ untuk suatu $\epsilon > 0$.
\end{prop}

\begin{prop} Misalkan $A$ aljabar-$C^*$ beridentitas dan $c \in A$. Jika $c \ge \vc 1$, maka $c$
invertibel dan $\vc 0 \le c^{-1} \le \vc 1$.
\end{prop}

\begin{prop} Jika $a$ adalah elemen positif yang invertibel dalam suatu aljabar-$C^*$ beridentitas, maka $a^{-1}$
positif.
\end{prop}

Selanjutnya kita tunjukkan bahwa notasi $a^{\ssst{-\frac12}}$ tidak ambigu.

\begin{prop} Misalkan $a \in A^+$, dengan $A$ aljabar-$C^*$ beridentitas. Jika $a$ invertibel, maka
$a^{\ssst{\frac12}}$ juga invertibel dan $\bigl(a^{\ssst{\frac12}}\bigr)^{-1} =
\bigl(a^{-1}\bigr)^{\ssst{\frac12}}$.
\end{prop}

\begin{cor}\label{00183313} Jika $a$ adalah elemen invertibel dalam suatu aljabar-$C^*$ beridentitas,
 \index{invers!dari suatu nilai mutlak}%
maka~$\abs a$ juga invertibel.
\end{cor}

\begin{prop}\label{0018401} Misalkan $a$ dan $b$ adalah elemen dari suatu aljabar-$C^*$. Jika
$\vc 0 \le a \le b$ dan $a$ invertibel, maka $b$ invertibel dan $b^{-1} \le
a^{-1}$.
\end{prop}

\begin{prop} Jika $a$ dan $b$ adalah elemen dari suatu aljabar-$C^*$ dan $\vc 0 \le a \le b$,
maka $\sqrt a \le \sqrt b$.
\end{prop}

\begin{exam} Misalkan $a$ dan $b$ elemen dari suatu aljabar-$C^*$ dengan $\vc 0 \le a \le b$.
\emph{Tidak} selalu berlaku bahwa $a^2 \le b^2$.
\end{exam}

\begin{proof}[\emph{Petunjuk bukti}] Dalam aljabar-$C^*$ $\mathbf M_2$, misalkan $a =
\begin{bmatrix} 1 & 0 \\ 0 & 0 \end{bmatrix}$ dan
$b = a + \tfrac12 \begin{bmatrix} 1 & 1 \\ 1 & 1 \end{bmatrix}$. \ns
\end{proof}
























\section{Identitas Aproksimatif}
\begin{defn} Suatu
 \index{identitas!aproksimatif}%
 \index{aproksimatif!identitas}%
\df{identitas aproksimatif} (atau
 \index{unit!aproksimatif}%
 \index{aproksimatif!unit}%
\df{unit aproksimatif}) dalam suatu aljabar-$C^*$ $A$ adalah jaring menaik
${(e_\lambda)}_{\lambda \in \Lambda}$ yang terdiri atas elemen-elemen positif dari $A$ sedemikian sehingga
$\norm{e_\lambda} \le 1$ dan $ae_\lambda \sto a$ (secara ekuivalen, $e_\lambda\, a \sto a$)
untuk setiap $a \in A$. Jika jaring tersebut ternyata merupakan barisan, kita memperoleh
 \index{sekuensial!identitas aproksimatif}%
 \index{aproksimatif!identitas!sekuensial}%
\df{identitas aproksimatif sekuensial}. Dalam literatur, berhati-hatilah terhadap definisi yang
beragam: banyak penulis menghilangkan syarat bahwa jaring tersebut menaik dan/atau bahwa jaring itu
terbatas.
\end{defn}

\begin{exam} Misalkan $A = \fml C_0(\R)$. Untuk setiap $n \in \N$, misalkan $U_n = (-n,n)$ dan
$e_n\colon \R \sto [0,1]$ suatu fungsi dalam $A$ yang dukungannya termuat dalam $U_{n+1}$
dan sedemikian sehingga $e_n(x) = 1$ untuk setiap $x \in U_n$. Maka $(e_n)$ adalah identitas
aproksimatif (sekuensial) bagi~$A$.
\end{exam}

\begin{prop} Jika $A$ adalah aljabar-$C^*$, maka himpunan
 \index{lambda@$\Lambda$ (anggota $A^+$ dalam bola satuan terbuka)}%
   \[ \Lambda := \{a \in A^+\colon \norm a < 1\} \]
merupakan himpunan terarah (terhadap urutan yang diwarisinya dari $\ofml H(A)$).
\end{prop}

\begin{proof}[\emph{Petunjuk bukti}] Periksa bahwa fungsi
  \[ f\colon [0,1) \sto \R^+\colon t \mapsto (1 - t)^{-1} - 1 = \frac t{1 - t} \]
menginduksi isomorfisme urutan $f\colon \Lambda \sto A^+$. (Suatu
 \index{urutan!isomorfisme}%
 \index{isomorfisme!urutan}%
\df{isomorfisme urutan} antara himpunan-himpunan terurut parsial adalah bijeksi yang mempertahankan urutan,
yang inversnya juga mempertahankan urutan.) Bukti yang saksama untuk hasil ini melibatkan pemeriksaan
terhadap cukup banyak rincian. \ns
\end{proof}

\begin{prop} Jika $A$ adalah aljabar-$C^*$, maka himpunan
  \[ \Lambda := \{a \in A^+\colon \norm a < 1\} \]
merupakan identitas aproksimatif bagi~$A$.
\end{prop}

\begin{cor}\label{00190161} Setiap aljabar-$C^*$ $A$ memiliki identitas aproksimatif.
 \index{aproksimatif!identitas!keberadaan}%
Jika $A$ separabel, maka $A$ memiliki identitas aproksimatif sekuensial.
\end{cor}

\begin{prop}\label{0019017} Setiap ideal aljabar (dua sisi) tertutup dalam suatu aljabar-$C^*$
bersifat swaadjoin.
\end{prop}

\begin{prop}\label{001902} Jika $J$ ideal dalam suatu aljabar-$C^*$ $A$, maka
 \index{hasil bagi!aljabar-$C^*$}%
$A/J$ adalah aljabar-$C^*$.
\end{prop}

\begin{proof} Lihat \cite{HigsonR:2000}, teorema 1.7.4, atau \cite{Davidson:1996}, halaman
13--14, atau~\cite{Fillmore:1996}, teorema 2.5.4. \ns
\end{proof}

\begin{exam} Jika $J$ ideal dalam suatu aljabar-$C^*$ $A$, maka barisan
  \[ \vc 0 \to J \to A \to^\pi A/J \to \vc 0 \]
adalah barisan eksak pendek.
\end{exam}

\begin{prop}\label{0019023} Jangkauan suatu homomorfisme-$*$ antara aljabar-$C^*$ bersifat tertutup
(dan karena itu merupakan aljabar-$C^*$).
\end{prop}

\begin{defn} Suatu subaljabar-$C^*$ $B$ dari aljabar-$C^*$ $A$ disebut
 \index{herediter}%
\df{herediter} jika $a \in B$ setiap kali $a \in A$, $b \in B$, dan $\vc 0 \le a \le b$.
\end{defn}

\begin{prop}\label{001915} Andaikan $x^*x \le a$ dalam suatu aljabar-$C^*$ $A$. Maka terdapat
$b \in A$ sedemikian sehingga $x = ba^{\frac14}$ dan $\norm b \le {\norm a}^\frac14$.
\end{prop}

\begin{proof} Lihat \cite{Davidson:1996}, halaman 13. \ns
\end{proof}

\begin{prop} Andaikan $J$ ideal dalam suatu aljabar-$C^*$ $A$, $j \in J^+$, dan
$a^*a \le j$. Maka $a \in J$. Jadi, ideal-ideal dalam aljabar-$C^*$ bersifat herediter.
\end{prop}

\begin{thm}\label{001922} Misalkan $A$ dan $B$ aljabar-$C^*$ dan $J$ ideal dalam~$A$.\newline
Jika $\phi$ \mbox{homomorfisme-$*\,$} dari $A$ ke $B$ dan $\ker\phi \supseteq J$, maka terdapat
tepat satu \mbox{homomorfisme-$*\,$} $\widetilde\phi\colon A/J \sto B$ yang membuat
diagram berikut komutatif.
 \[ \xymatrix@+30pt{A \ar[d]_*+{\pi} \ar[dr]^*+{\phi} & \\
            A/J \ar[r]_*+{\widetilde\phi} & B  } \]
Selain itu, $\widetilde\phi$ injektif jika dan hanya jika $\ker\phi = J$; dan
$\widetilde\phi$ surjektif jika dan hanya jika $\phi$ surjektif.
\end{thm}

\begin{cor}\label{001923} Jika $\vc 0 \to A \to^\phi E \to B \to \vc 0$ adalah barisan eksak
pendek dari aljabar-$C^*$, maka $E/\ran\phi$ dan $B$ isomorfik-$*\,$ secara isometrik.
\end{cor}

\begin{cor}\label{001924} Setiap aljabar-$C^*$ $A$ memiliki kodimensi satu dalam
unitalisasinya~$\wt A$; yaitu, $\dim \wt A/A = 1$.
\end{cor}

\begin{prop} Misalkan $A$ aljabar-$C^*$, $B$ subaljabar-$C^*$ dari $A$, dan $J$
ideal dalam~$A$. Maka
   \[  B/(B \cap J) \cong (B + J)/J\,. \]
\end{prop}

Misalkan $B$ subaljabar beridentitas dari suatu aljabar sembarang $A$. Jelas bahwa jika suatu
elemen $b \in B$ invertibel dalam $B$, maka elemen itu juga invertibel dalam~$A$. Ternyata
konversnya benar dalam aljabar-$C^*$: jika $b$ invertibel dalam $A$, maka inversnya
terletak dalam~$B$. Hal ini biasanya dinyatakan dengan mengatakan bahwa \emph{setiap subaljabar-$C^*$
beridentitas dari suatu aljabar-$C^*$ bersifat}
 \index{invers!tertutup}%
 \index{tertutup!invers}%
\df{tertutup terhadap invers}.

\begin{prop} Misalkan $B$ subaljabar-$C^*$ beridentitas dari suatu aljabar-$C^*$ $A$. Jika $b \in \inv(A)$,
maka $b^{-1} \in B$.
\end{prop}

\begin{cor}\label{0019283} Misalkan $B$ subaljabar-$C^*$ beridentitas dari suatu aljabar-$C^*$ $A$ dan
$b \in B$. Maka
   \[ \sigma_B(b) = \sigma_A(b)\,. \]
\end{cor}

\begin{cor}\label{0019286} Misalkan $\phi\colon A \sto B$ adalah monomorfisme-$C^*$ beridentitas dari
suatu aljabar-$C^*$ $A$ dan $a \in A$. Maka
   \[ \sigma(a) = \sigma(\phi(a))\,. \]
\end{cor}



\endinput

 \chapter{KONSTRUKSI GELFAND-NAIMARK-SEGAL}\label{gelfand_naimark_segal}

\section{Fungsional Linear Positif}
\begin{defn} Misalkan $A$ suatu aljabar dengan involusi. Untuk setiap fungsional linear $\tau$ pada
$A$ dan setiap $a \in A$ definisikan
 \index{<unopurstar@$\tau^\star$ (konjugat dari $\tau$)}%
$\tau^\star(a) = \conj{\tau(a^*)}$.  Kita mengatakan bahwa $\tau$ bersifat
 \index{Hermitian!fungsional linear}%
 \index{linear!fungsional!Hermitian}%
\df{Hermitian} jika $\tau^\star = \tau$.  Perhatikan bahwa suatu fungsional linear $\tau\colon A
\sto \C$ bersifat Hermitian jika dan hanya jika mempertahankan involusi; yaitu, $\tau^\star =
\tau$ jika dan hanya jika $\tau(a^*) = \conj{\tau(a)}$ untuk semua $a \in A$.
\end{defn}

\begin{cau} $\tau^\star$ yang didefinisikan di atas jangan dikacaukan dengan pemetaan adjoin
biasa $\tau^*\colon \C^* \sto A^*$. Konteks (atau penggunaan kaca pembesar) semestinya
memperjelas mana yang dimaksud.
\end{cau}

\begin{prop} Suatu fungsional linear $\tau$ pada aljabar-$C^*$ $A$ bersifat Hermitian jika dan hanya jika
$\tau(a)~\in~\R$ setiap kali $a$ swaadjoin.
\end{prop}

Seperti halnya selalu terjadi pada pemetaan di antara ruang-ruang vektor terurut, pemetaan
\emph{positif} ialah pemetaan yang membawa elemen positif ke elemen positif.

\begin{defn} Suatu fungsional linear $\tau$ pada aljabar-$C^*$ $A$ bersifat
 \index{positif!fungsional linear}%
 \index{linear!fungsional!positif}%
\df{positif} jika $\tau(a) \ge 0$ untuk setiap $a \in A$ yang memenuhi $a \ge \vc 0$.
\end{defn}


\begin{prop} Setiap fungsional linear positif pada suatu aljabar-$C^*$ bersifat Hermitian.
\end{prop}

\begin{prop} Keluarga fungsional linear positif merupakan suatu kerucut konveks proper dalam ruang
vektor real semua fungsional linear Hermitian pada suatu aljabar-$C^*$. Kerucut itu menginduksi suatu
urutan parsial pada ruang vektor tersebut: $\tau_1 \le \tau_2$ setiap kali $\tau_2 - \tau_1$
bersifat positif.
\end{prop}

\begin{defn} Suatu
 \index{keadaan}%
\df{keadaan} pada aljabar-$C^*$ $A$ adalah fungsional linear positif $\tau$ pada $A$ sedemikian sehingga
$\norm\tau = 1$. Jika $A$ beridentitas, syarat ini ekuivalen dengan $\tau(\vc 1_A) = 1$.
\end{defn}

\begin{exam} Misalkan $x$ suatu vektor dalam ruang Hilbert $H$. Definisikan
   \[ \omega_x:\ofml B(H) \sto \C\colon T \mapsto \langle Tx,x \rangle\,. \]
Maka $\omega_x$ merupakan fungsional linear positif pada $\ofml B(H)$. Jika $x$ suatu
vektor satuan, maka $\omega_x$ merupakan keadaan pada~$\ofml B(H)$. Suatu keadaan $\tau$ disebut
 \index{vektor!keadaan}%
 \index{keadaan!vektor}%
\df{keadaan vektor} jika $\tau = \omega_x$ untuk suatu vektor satuan~$x$.
\end{exam}

\begin{prop}[Ketaksamaan Schwarz]\label{0025105} Jika $\tau$ suatu fungsional linear positif pada
aljabar-$C^*$ $A$, maka
   \[ \abs{\tau(b^*a)}^2 \le \tau(a^*a)\tau(b^*b) \]
untuk semua $a$, $b \in A$.
 \index{ketaksamaan Schwarz}%
 \index{ketaksamaan!Schwarz}%
\end{prop}

\begin{prop} Suatu fungsional linear $\tau$ pada aljabar-$C^*$ beridentitas $A$ bersifat positif jika dan hanya jika
fungsional itu terbatas dan $\norm\tau = \tau(\vc 1_A)$.
\end{prop}

\begin{proof} Lihat \cite{KadisonR:1983}, halaman 256--257. \ns
\end{proof}



















\section{Representasi}
\begin{defn}\label{0028} Misalkan $A$ suatu aljabar-$C^*$. Suatu
 \index{representasi}%
\df{representasi} dari $A$ adalah pasangan $(\pi,H)$ dengan $H$ ruang Hilbert dan $\pi
\colon A \sto \ofml B(H)$ suatu homomorfisme-$*\,$. Biasanya cukup dikatakan bahwa $\pi$ merupakan
representasi dari~$A$. Apabila kita ingin menekankan peran ruang Hilbert tertentu itu,
kita mengatakan bahwa $\pi$ \emph{merupakan representasi dari $A$ pada~$H$.} Bergantung pada konteks, kita
dapat menulis $\pi_a$ atau $\pi(a)$ untuk operator ruang Hilbert yang merupakan citra
elemen aljabar $a$ di bawah~$\pi$. Suatu representasi $\pi$ dari $A$ pada $H$ disebut
 \index{representasi!tak terdegenerasi}%
 \index{tak terdegenerasi!representasi}%
\df{tak terdegenerasi} jika $\pi(A)H$ padat di~$H$.
\end{defn}

 \index{konvensi!representasi aljabar-$C*$ beridentitas bersifat beridentitas}%
\begin{conv} Kita menambahkan persyaratan berikut pada definisi sebelumnya: jika aljabar-$C^*$
$A$ beridentitas, maka suatu representasi dari $A$ harus berupa homomorfisme-$*\,$ \emph{beridentitas}.
\end{conv}

\begin{defn} Suatu representasi $\pi$ dari aljabar-$C^*$ $A$ pada ruang Hilbert $H$ disebut
 \index{representasi!setia}%
 \index{representasi setia}%
\df{setia} jika representasi itu injektif. Jika terdapat suatu vektor $x \in H$ sedemikian sehingga
$\pi^{\sto}(A)x = \{\pi_a(x)\colon a \in A\}$ padat di $H$, maka kita mengatakan bahwa
representasi $\pi$ bersifat
 \index{representasi!siklik}%
 \index{siklik!representasi}%
\df{siklik} dan bahwa $x$ merupakan
 \index{representasi!vektor siklik bagi suatu}%
 \index{siklik!vektor}%
\df{vektor siklik} bagi~$\pi$.
\end{defn}

\begin{exam}\label{002802} Misalkan $(S, \sfml A, \mu)$ suatu ruang ukur $\sigma$-hingga dan
$L_\infty = L_\infty(S, \sfml A, \mu)$ aljabar-$C^*$ dari fungsi-fungsi yang terbatas secara esensial dan
terukur terhadap $\mu$ pada~$S$. Seperti yang kita lihat dalam contoh~\ref{C063527}, untuk setiap $\phi \in
L_\infty$ operator perkalian $M_\phi$ yang bersesuaian merupakan operator pada
ruang Hilbert $L_2 = L_2(S, \sfml A, \mu)$. Pemetaan $M\colon L_\infty \sto \ofml
B(L_2)\colon \phi \mapsto M_\phi$ merupakan representasi setia dari aljabar-$C^*$
$L_\infty$ pada ruang Hilbert~$L_2$.
\end{exam}

\begin{exam} Misalkan $\fml C([0,1])$ aljabar-$C^*$ dari fungsi-fungsi kontinu
pada interval~$[0,1]$. Untuk setiap $\phi \in \fml C([0,1])$, operator perkalian
$M_\phi$ yang bersesuaian merupakan operator pada ruang Hilbert $L_2 = L_2([0,1])$
dari fungsi-fungsi pada~$[0,1]$ yang terintegralkan kuadrat terhadap ukuran Lebesgue.
Pemetaan $M\colon \fml C([0,1]) \sto \ofml B(L_2)\colon \phi \mapsto M_\phi$ merupakan
representasi setia dari aljabar-$C^*$ $\fml C([0,1])$ pada ruang Hilbert~$L_2$.
\end{exam}

\begin{exam} Andaikan $\pi$ suatu representasi dari aljabar-$C^*$ beridentitas $A$ pada
ruang Hilbert~$H$ dan $x$ suatu vektor satuan dalam~$H$. Jika $\omega _x$ merupakan
keadaan vektor yang bersesuaian pada $\ofml B(H)$, maka $\omega_x \circ \pi$ merupakan keadaan pada~$A$.
\end{exam}

\begin{exer} Misalkan $X$ suatu ruang Hausdorff kompak lokal. Carilah suatu representasi
isometrik (dan karena itu setia) $(\pi,H)$ dari aljabar-$C^*$ $C_0(X)$ pada suatu
ruang Hilbert~$H$.
\end{exer}

\begin{defn} Misalkan $\rho$ suatu keadaan pada aljabar-$C^*$ $A$. Maka
   \[ L_\rho := \{a \in A\colon \rho(a^*a) = 0\} \]
disebut
 \index{kiri!kernel}%
 \index{kernel!kiri}%
\df{kernel kiri} dari~$\rho$.
\end{defn}

Ingat bahwa sebagai bagian dari pembuktian \emph{ketaksamaan Schwarz}~\ref{0025105} bagi fungsional
linear positif, kita telah memverifikasi hasil berikut.

\begin{prop} Jika $\rho$ suatu keadaan pada aljabar-$C^*$ $A$, maka $\langle a,b \rangle_0 :=
\rho(b^*a)$ mendefinisikan suatu hasil kali dalam semu pada~$A$.
\end{prop}

\begin{cor} Jika $\rho$ suatu keadaan pada aljabar-$C^*$ $A$, maka kernel kirinya $L_\rho$ merupakan suatu
subruang vektor dari $A$ dan $\langle\,[a],[b]\,\rangle := \langle a,b \rangle_0$ mendefinisikan
suatu hasil kali dalam pada ruang vektor hasil bagi~$A/L_\rho$.
\end{cor}

\begin{prop} Misalkan $\rho$ suatu keadaan pada aljabar-$C^*$ $A$ dan $a \in L_\rho$. Maka $\rho(b^*a)
= 0$ untuk setiap $b \in A$.
\end{prop}

\begin{prop} Jika $\rho$ suatu keadaan pada aljabar-$C^*$ $A$, maka kernel kirinya $L_\rho$ merupakan suatu
ideal kiri tertutup dalam~$A$.
\end{prop}


























\section{Konstruksi GNS dan Teorema Gelfand-Naimark Ketiga}
Teorema berikut dikenal sebagai
 \index{Gelfand-Naimark-Segal!konstruksi}%
 \index{konstruksi!Gelfand-Naimark-Segal}%
\emph{konstruksi Gelfand-Naimark-Segal} (\emph{konstruksi GNS}).

\begin{thm}[Konstruksi GNS]\label{0029} Misalkan $\rho$ suatu keadaan pada aljabar-$C^*$ $A$. Maka
terdapat suatu representasi siklik $\pi_\rho$ dari $A$ pada ruang Hilbert $H_\rho$ dan suatu
vektor satuan siklik $x_\rho$ untuk $\pi_\rho$ sedemikian sehingga $\rho = \omega_{x_\rho} \circ
\pi_\rho$.
\end{thm}

\begin{notn} Dalam pembahasan berikut, $\pi_\rho$, $H_\rho$, dan $x_\rho$ adalah representasi
siklik, ruang Hilbert, dan vektor satuan siklik yang keberadaannya dijamin oleh konstruksi
GNS ketika diterapkan pada suatu keadaan $\rho$ yang diberikan pada suatu aljabar-$C^*$ $A$.
\end{notn}

\begin{prop} Misalkan $\rho$ suatu keadaan pada aljabar-$C^*$ $A$ dan $\pi$ suatu representasi siklik
dari $A$ pada ruang Hilbert $H$ sedemikian sehingga $\rho = \omega_x \circ \pi$ untuk suatu vektor
satuan siklik $x$ bagi~$\pi$. Maka terdapat suatu pemetaan uniter $U$ dari $H_\rho$ ke~$H$ sedemikian sehingga
$x = Ux_\rho$ dan $\pi(a) = U \pi_\rho(a) U^*$ untuk semua $a \in A$.
\end{prop}

\begin{defn} Misalkan $\{H_\lambda\colon \lambda \in \Lambda\}$ suatu keluarga ruang Hilbert.
Nyatakan dengan $\bigoplus\limits_{\lambda \in \Lambda} H_\lambda$ himpunan semua fungsi
$x\colon \Lambda \sto \bigcup\limits_{\lambda \in \Lambda} H_\lambda\colon \lambda
\mapsto x_\lambda$ sedemikian sehingga $x_\lambda \in H_\lambda$ untuk setiap $\lambda \in \Lambda$ dan
$\D\sum\limits_{\lambda \in \Lambda} \norm{x_\lambda}^2 < \infty$. Pada
$\bigoplus\limits_{\lambda \in \Lambda} H_\lambda$, definisikan penjumlahan dan perkalian
skalar secara titik demi titik; yaitu, $(x + y)_\lambda = x_\lambda + y_\lambda$ dan $(\alpha
x)_\lambda = \alpha x_\lambda$ untuk semua $\lambda \in \Lambda$, dan definisikan hasil kali dalam
dengan $\langle x,y \rangle = \sum_{\lambda \in \Lambda} \langle x_\lambda,y_\lambda\rangle$.
Operasi-operasi ini terdefinisi dengan baik dan menjadikan $\bigoplus\limits_{\lambda \in \Lambda}
H_\lambda$ suatu ruang Hilbert. Ruang ini adalah
 \index{jumlah langsung!ruang Hilbert}%
\df{jumlah langsung} dari ruang-ruang Hilbert~$H_\lambda$. Berbagai notasi untuk elemen-elemen jumlah
langsung ini digunakan dalam literatur: $x$, $(x_\lambda)$, $(x_\lambda)_{\lambda \in
\Lambda}$, dan $\oplus_\lambda x_\lambda$ lazim digunakan.

Sekarang andaikan bahwa $\{T_\lambda\colon \lambda \in \Lambda\}$ suatu keluarga operator ruang
Hilbert dengan $T_\lambda \in \ofml B(H_\lambda)$ untuk setiap $\lambda \in \Lambda$.
Andaikan pula bahwa $\sup\{\norm{T_\lambda}\colon \lambda \in \Lambda\} < \infty$. Maka
$T(x_\lambda)_{\lambda \in \Lambda} = (T_\lambda x_\lambda)_{\lambda \in \Lambda}$
mendefinisikan suatu operator pada ruang Hilbert $\bigoplus_\lambda H_\lambda$. Operator $T$
biasanya dinotasikan dengan $\bigoplus_\lambda T_\lambda$ dan disebut
 \index{jumlah langsung!operator ruang Hilbert}%
\df{jumlah langsung} dari operator-operator~$T_\lambda$.
\end{defn}

\begin{prop} Pernyataan-pernyataan dalam definisi sebelumnya bahwa
\smash[b]{$\bigoplus\limits_{\lambda \in \Lambda} H_\lambda$} merupakan ruang Hilbert dan
$\bigoplus_\lambda T_\lambda$ merupakan operator pada $\bigoplus\limits_{\lambda \in \Lambda}
H_\lambda$ adalah benar.
\end{prop}

\begin{exam} Misalkan $A$ suatu aljabar-$C^*$ dan $\{\pi_\lambda\colon \lambda \in \Lambda\}$ suatu
keluarga representasi dari $A$ pada ruang-ruang Hilbert $H_\lambda$ sedemikian sehingga $\pi_\lambda(a)
\in \ofml B(H_\lambda)$ untuk setiap $\lambda \in \Lambda$ dan setiap $a \in A$. Maka
  \[ \pi = \textstyle{ \bigoplus_\lambda \pi_\lambda\colon A \sto
       \ofml B\bigl(\bigoplus_\lambda H_\lambda\bigr)
       \colon a \mapsto \bigoplus_\lambda \pi_\lambda(a)} \]
merupakan representasi dari $A$ pada ruang Hilbert $\bigoplus_\lambda H_\lambda$. Representasi ini adalah
 \index{jumlah langsung!representasi}%
\df{jumlah langsung} dari representasi-representasi~$\pi_\lambda$.
\end{exam}

\begin{thm}\label{thm_exist_faith_rep} Setiap aljabar-$C^*$ memiliki suatu representasi setia.
\end{thm}

\begin{proof} Lihat \cite{Blackadar:2006}, Korolari II.6.4.10; \cite{Conway:1990}, halaman 253;
\cite{DoranB:1986}, Teorema 19.1; \cite{Fillmore:1996}, Teorema 5.4.1;
\cite{KadisonR:1983}, halaman 281; atau \cite{Murphy:1990}, Teorema 3.4.1. \ns
\end{proof}

Suatu perumusan ulang yang jelas dari teorema sebelumnya adalah versi ketiga dari
\emph{teorema Gelfand-Naimark}, yang menyatakan bahwa setiap aljabar-$C^*$ (pada hakikatnya)
merupakan suatu aljabar operator ruang Hilbert.

\begin{cor}[Teorema Gelfand-Naimark III]\label{002973} Setiap aljabar-$C^*$
 \index{Gelfand-Naimark!Teorema III}
isomorfik-$*\,$ secara isometrik dengan suatu subaljabar-$C^*$ dari $\ofml B(H)$ untuk suatu ruang
Hilbert~$H$.
\end{cor}


\endinput

 \chapter{ALJABAR PENGALI}\label{multiplier_algebras}

\section{Modul Hilbert}
\begin{notn}\label{0038014} Hasil kali dalam yang muncul sebelumnya dalam catatan ini
dan yang dijumpai dalam buku teks dan monograf standar tentang ruang Hilbert,
analisis fungsional, dan sebagainya, bersifat linear dalam variabel pertama dan linear konjugat dalam
variabel kedua. Kebanyakan ahli aljabar operator kontemporer memilih bekerja dengan objek
yang disebut modul Hilbert-$A$ \emph{kanan} ($A$ suatu aljabar-$C^*$). Untuk modul semacam itu,
ternyata lebih mudah menggunakan ``hasil kali dalam'' yang linear dalam variabel kedua
dan linear konjugat dalam variabel pertama. Walaupun perubahan konvensi ini mungkin menimbulkan
sedikit rasa jengkel, secara matematis pengaruhnya kecil. Tentu saja,
kita ingin ruang Hilbert menjadi contoh modul Hilbert-$\C$. Agar hal ini
mungkin, kita membekali ruang Hilbert, yang hasil kali dalamnya dinyatakan dengan
$\langle\hphantom{X},\hphantom{X}\rangle$, dengan ``hasil kali dalam'' baru yang didefinisikan
 \index{<bracketspointy@$\langle x\,\vert\,y \rangle$ (hasil kali dalam)}%
oleh $\langle x\,\vert\,y\rangle := \langle y,x \rangle$. ``Hasil kali dalam'' ini linear
dalam variabel kedua dan linear konjugat dalam variabel pertama. Saya akan berusaha konsisten dalam
menggunakan $\langle\hphantom{X},\hphantom{X}\rangle$ untuk hasil kali dalam standar dan
$\langle\hphantom{X}\,\vert\,\hphantom{X}\rangle$ untuk hasil kali dalam yang linear dalam
variabel keduanya.

Solusi lain (yang sangat umum) untuk persoalan ini adalah menetapkan bahwa hasil kali dalam
selalu linear dalam variabel keduanya dan ``mengoreksi'' buku teks, monograf, dan
makalah standar sesuai dengan ketetapan tersebut.
\end{notn}

\begin{conv} Berdasarkan uraian sebelumnya, selanjutnya kita akan menggunakan kata
``seskuilinear'' untuk berarti salah satu di antara
 \index{seskuilinear}%
 \index{konvensi!tentang \emph{seskuilinear}}%
\emph{linear dalam variabel pertama dan linear konjugat dalam variabel kedua} atau \emph{linear dalam
variabel kedua dan linear konjugat dalam variabel pertama}.
\end{conv}

\begin{defn}\label{0038017} Misalkan $A$ suatu aljabar-$C^*$ tak nol. Suatu ruang vektor $V$ adalah
\df{modul-$A$}
 \index{modul}%
 \index{amodule@modul-$A$}%
jika terdapat suatu pemetaan bilinear
  \[ B\colon V \times A \sto V\colon (x,a) \mapsto xa \]
sedemikian sehingga $x(ab) = (xa)b$ berlaku untuk semua $x \in V$ dan $a$, $b \in A$. Kita juga mensyaratkan
bahwa $x \vc 1_A = x$ untuk setiap $x \in V$ jika $A$ beridentitas. (Suatu fungsi dua variabel
 \index{bilinear}%
disebut \df{bilinear} jika fungsi itu linear dalam kedua variabelnya.)
\end{defn}

\begin{defn}\label{0038021} Untuk memperjelas sebagian materi berikutnya, akan
berguna untuk memiliki definisi formal berikut. Suatu
 \index{vektor!ruang}%
\df{ruang vektor (kompleks)} adalah tripel $(V,+,M)$ dengan $(V,+)$ suatu grup Abel dan
$M\colon \C \sto \Hom(V,+)$ suatu homomorfisme gelanggang beridentitas. (Seperti dapat diduga,
$\Hom(V,+)$ menyatakan gelanggang beridentitas dari endomorfisme grup $(V,+)$.) Jadi, suatu
 \index{aljabar}%
\df{aljabar} adalah kuadruplet terurut $(A,+,M,\,\cdot\,)$ dengan $(A,+,M)$ suatu ruang
vektor dan $\,\cdot\,\colon A \times A \sto A \colon (a,b) \mapsto a \cdot b = ab$ suatu
operasi biner pada $A$ yang memenuhi syarat-syarat dalam~\ref{defn_alg}.
\end{defn}

\begin{exer} Pastikan bahwa definisi \emph{ruang vektor} sebelumnya ekuivalen
dengan definisi yang biasa Anda gunakan.
\end{exer}

\begin{defn} Misalkan $(A,+,M,\,\cdot\,)$ suatu aljabar (kompleks). Maka $A^{\textrm{op}}$
 \index{<A@$A^{\textrm{op}}$ (aljabar lawan dari $A$)}%
adalah aljabar $(A,+,M,*)$ dengan $a*b = b\cdot a$ untuk semua $a$, $b \in A$. Aljabar ini adalah
 \index{aljabar lawan}%
 \index{aljabar!lawan}%
\df{aljabar lawan} dari~$A$. Aljabar itu hanyalah $A$ dengan urutan perkalian dibalik.
Jika $A$ dan $B$ adalah aljabar, maka setiap fungsi $f \colon A \sto B$ dengan cara yang jelas menginduksi
suatu fungsi dari $A$ ke $B^{\textrm{op}}$ (atau dari $A^{\textrm{op}}$ ke $B$,
atau dari $A^{\textrm{op}}$ ke $B^{\textrm{op}}$). Semua fungsi ini akan kita nyatakan
cukup dengan~$f$.
\end{defn}

\begin{defn}\label{0038027} Misalkan $A$ dan $B$ aljabar. Suatu fungsi
$\phi\colon A \sto B$ adalah suatu
 \index{antihomomorfisme}%
\df{antihomomorfisme} jika fungsi $\phi\colon A \sto B^{\textrm{op}}$ merupakan suatu homomorfisme.
Suatu antihomomorfisme bijektif adalah
 \index{antiisomorfisme}%
\df{antiisomorfisme}. Dalam contoh~\ref{0022057} kita menjumpai \emph{antiisomorfisme}
ruang Hilbert. Antiisomorfisme tersebut agak berbeda karena alih-alih membalik perkalian
(yang tidak dimiliki ruang Hilbert), antiisomorfisme itu mengonjugatkan skalar. Dalam kedua kasus,
antiisomorfisme bukanlah sesuatu yang sama sekali berbeda dari isomorfisme, melainkan justru
sesuatu yang sangat serupa.
\end{defn}

Gagasan modul-$A$ (dengan $A$ suatu aljabar) telah didefinisikan dalam~\ref{0038017}. Anda
mungkin lebih menyukai definisi alternatif berikut, yang lebih selaras dengan
definisi \emph{ruang vektor} dalam~\ref{0038021}.

\begin{defn}\label{0038031} Misalkan $A$ suatu aljabar. Kuadruplet terurut
 \index{modul}%
 \index{amodule@modul-$A$}%
$(V,+,M,\Phi)$ disebut \df{modul-$A$} jika $(V,+,M)$ suatu ruang vektor dan $\Phi\colon A \sto \ofml L(V)$ suatu
homomorfisme aljabar. Jika $A$ beridentitas, kita juga mensyaratkan bahwa $\Phi$ beridentitas.
\end{defn}

\begin{exer} Periksa bahwa definisi \emph{modul-$A$} dalam~\ref{0038017}
dan~\ref{0038031} ekuivalen.
\end{exer}

Sekarang kita nyatakan secara tepat apa artinya, ketika $A$ suatu aljabar-$C^*$, membekali modul-$A$ dengan
hasil kali dalam bernilai-$A$.

\begin{defn} Misalkan $A$ suatu aljabar-$C^*$. Suatu
 \index{hasil kali dalam semu!modul-$A$}%
 \index{modul!hasil kali dalam semu modul-$A$}%
 \index{amodule@modul-$A$!hasil kali dalam semu}%
\df{modul-$A$ hasil kali dalam semu} adalah suatu modul-$A$ $V$ beserta pemetaan
  \[ \beta\colon V \times V \sto A\colon (x,y) \mapsto \langle x\vert\, y\rangle \]
yang linear dalam variabel keduanya dan memenuhi
 \begin{enumerate}
    \item[(i)] $\langle x\vert\, ya\rangle = \langle x\vert\, y\rangle a$,
    \item[(ii)] $\langle x\vert\, y\rangle = \langle y\vert\, x\rangle^*$, dan
    \item[(iii)] $\langle x\vert\, x\rangle \ge \vc 0$
 \end{enumerate}
untuk semua $x$, $y \in V$ dan $a \in A$. Modul ini adalah suatu
 \index{hasil kali dalam!modul-$A$}%
 \index{modul!hasil kali dalam modul-$A$}%
 \index{amodule@modul-$A$!hasil kali dalam}%
\df{modul-$A$ hasil kali dalam} (atau suatu
 \index{pra-Hilbert!modul-$A$}%
 \index{modul!pra-Hilbert-$A$}%
 \index{amodule@modul-$A$!pra-Hilbert}%
\df{modul pra-Hilbert-$A$}) jika sebagai tambahan
 \begin{enumerate}
    \item[(iv)] $\langle x\vert\, x\rangle = \vc 0$ mengakibatkan $x = \vc 0$
 \end{enumerate}
untuk $x \in V$. Kita akan menyebut pemetaan $\beta$ sebagai
 \index{hasil kali dalam!bernilai-$A$}%
 \index{Aval@hasil kali dalam (semu) bernilai-$A$}%
\df{hasil kali dalam (semu) bernilai-$A$} pada~$V$.
\end{defn}

\begin{exam} Setiap ruang hasil kali dalam adalah modul-$\C$ hasil kali dalam.
\end{exam}

\begin{prop} Misalkan $A$ suatu aljabar-$C^*$ dan $V$ suatu modul-$A$ hasil kali dalam semu.
Hasil kali dalam semu $\langle\hphantom{x}\vert\,\hphantom{x}\rangle$ bersifat linear konjugat
dalam variabel pertamanya, baik secara harfiah maupun dalam arti bahwa $\langle va \vert\,
w\rangle = a^* \langle v \vert\, w\rangle$ untuk semua $v$, $w \in V$ dan $a \in A$.
\end{prop}

\begin{prop}[Ketaksamaan Schwarz---untuk modul-$A$ hasil kali dalam]\label{0038038} Misalkan
$V$ suatu modul-$A$ hasil kali dalam dengan $A$ suatu aljabar-$C^*$. Maka
  \[ \langle x\vert\, y \rangle^*\langle x\vert\, y \rangle
          \le \norm{\langle x\vert\, x \rangle}\,\langle y\vert\, y \rangle \]
untuk semua $x$, $y \in V$.
\end{prop}

\begin{proof}[\emph{Petunjuk pembuktian}] Tunjukkan bahwa, tanpa mengurangi keumuman,
kita dapat mengasumsikan bahwa
$\norm{\langle x\vert\, x\rangle} = 1$. Pertimbangkan elemen positif $\langle xa -
y\vert\, xa - y\rangle$ dengan $a = \langle x \vert\, y \rangle$. Gunakan
proposisi~\ref{0018203} dan~\ref{0018204}. \ns
\end{proof}

\begin{defn}\label{0038041} Untuk setiap elemen $v$ dari suatu modul-$A$ hasil kali dalam
(dengan $A$ suatu aljabar-$C^*$), definisikan
  \[ \norm v := \norm{\langle v \vert\, v \rangle}^{1/2}. \]
\end{defn}

\begin{prop}[Ketaksamaan Schwarz lainnya] Misalkan $A$ suatu aljabar-$C^*$ dan $V$ suatu
modul-$A$ hasil kali dalam. Maka untuk semua $v$, $w \in V$
  \[ \norm{\langle v \vert\, w\rangle} \le \norm v\,\norm w. \]
\end{prop}

\begin{cor} Jika $v$ dan $w$ merupakan elemen suatu modul-$A$ hasil kali dalam (dengan $A$ suatu
aljabar-$C^*$), maka $\norm{v + w} \le \norm v + \norm w$ dan pemetaan $x \mapsto \norm x$
merupakan norma pada~$V$.
\end{cor}

\begin{prop} Jika $A$ suatu aljabar-$C^*$ dan $V$ suatu modul-$A$ hasil kali dalam, maka
  \[ \norm{va} \le \norm v\, \norm a \]
untuk semua $v \in V$ dan $a \in A$.
\end{prop}

\begin{defn} Misalkan $A$ suatu aljabar-$C^*$ dan $V$ suatu modul-$A$ hasil kali dalam. Jika $V$
lengkap terhadap (metrik yang diinduksi oleh) norma yang didefinisikan dalam~\ref{0038041},
maka $V$ adalah suatu
 \index{amodule@modul-$A$!Hilbert}%
 \index{Hilbert!modul-$A$}%
\df{modul Hilbert-$A$}.
\end{defn}

\begin{exam}\label{0038111} Untuk $a$ dan $b$ dalam suatu aljabar-$C^*$ $A$,
 \index{Hilbert!modul-$A$!$A$ sebagai}%
 \index{amodule@modul-$A$!Hilbert!$A$ sebagai}%
definisikan
  \[ \langle a \vert\, b\rangle := a^*b\,. \]
Maka $A$ sendiri merupakan modul Hilbert-$A$. Setiap ideal kanan tertutup dalam $A$ juga merupakan
modul Hilbert-$A$.
\end{exam}

\begin{defn} Misalkan $V$ dan $W$ modul-modul Hilbert-$A$ dengan $A$ suatu aljabar-$C^*$. Suatu
pemetaan $T\colon V \sto W$ bersifat
 \index{alinear@linear-$A$}%
 \index{linear!$A$-}%
\df{linear-$A$} jika pemetaan itu linear dan jika $T(va) = T(v)a$ berlaku untuk semua $v \in V$ dan $a \in
A$. Pemetaan $T$ adalah suatu
 \index{amodule@modul-$A$!morfisme}%
 \index{modul!morfisme}%
 \index{morfisme!modul Hilbert-$A$}%
 \index{Hilbert!modul-$A$!morfisme}%
\df{morfisme modul Hilbert-$A$} jika pemetaan itu terbatas dan linear-$A$.
\end{defn}

Ingat dari~\ref{def000926} bahwa setiap operator ruang Hilbert memiliki adjoin.
Hal ini tidak berlaku untuk modul Hilbert-$A$.

\begin{defn} Misalkan $V$ dan $W$ modul-modul Hilbert-$A$ dengan $A$ suatu aljabar-$C^*$. Suatu
fungsi $T\colon V \sto W$ bersifat
 \index{dapat diadjoinkan}%
\df{dapat diadjoinkan} jika terdapat suatu fungsi $T^*\colon W \sto V$ yang memenuhi
  \[ \langle Tv \vert\, w \rangle = \langle v \vert\, T^*w \rangle \]
untuk semua $v \in V$ dan $w \in W$. Fungsi $T^*$, jika ada, adalah
 \index{adjoin}
\df{adjoin} dari~$T$. Nyatakan dengan
 \index{L@$\ofml L(V,W)$, $\ofml L(V)$ (pemetaan yang dapat diadjoinkan)}%
$\ofml L(V,W)$ keluarga pemetaan yang dapat diadjoinkan dari $V$ ke~$W$. Kita menyingkat $\ofml L(V,V)$
menjadi~$\ofml L(V)$.
\end{defn}

\begin{prop} Misalkan $V$ dan $W$ modul-modul Hilbert-$A$ dengan $A$ suatu aljabar-$C^*$. Jika suatu
fungsi $T\colon V \sto W$ dapat diadjoinkan, maka fungsi itu merupakan morfisme modul Hilbert-$A$.
Lebih lanjut, jika $T$ dapat diadjoinkan, maka adjoinnya juga dapat diadjoinkan dan $T^{**} = T$.
\end{prop}

\begin{exam} Misalkan $X$ interval satuan $[0,1]$ dan $Y = \{0\}$. Dengan topologi biasa,
$X$ merupakan ruang Hausdorff kompak dan $Y$ suatu subruang dari~$X$. Misalkan $A$ adalah
aljabar-$C^*$ $\fml C(X)$ dan $J_0$ adalah ideal $\{f \in A\colon f(0) = 0\}$ (lihat
proposisi~\ref{0006836}). Pandang $V = A$ dan $W = J_0$ sebagai modul-modul Hilbert-$A$ (lihat
contoh~\ref{0038111}). Maka pemetaan inklusi
$\iota\colon W \sto V$ merupakan morfisme modul Hilbert-$A$ yang tidak dapat diadjoinkan.
\end{exam}

\begin{prop} Misalkan $A$ suatu aljabar-$C^*$. Pasangan pemetaan $V \mapsto V$, $T \mapsto T^*$
merupakan funktor kontravarian dari kategori modul Hilbert-$A$ dan pemetaan yang dapat diadjoinkan
ke dirinya sendiri.
\end{prop}

\begin{prop}\label{0038244} Misalkan $A$ suatu aljabar-$C^*$ dan $V$ suatu modul Hilbert-$A$.
 \index{C*@$C^*$-aljabar!$\ofml L(V)$ sebagai suatu}%
 \index{L@$\ofml L(V)$!sebagai suatu aljabar-$C^*$}%
Maka $\ofml L(V)$ merupakan aljabar-$C^*$ beridentitas.
\end{prop}

\begin{notn} Misalkan $V$ dan $W$ modul-modul Hilbert-$A$ dengan $A$ suatu aljabar-$C^*$.
 \index{theta@$\Theta_{v,w}$}%
Untuk $v \in V$ dan $w \in W$, misalkan
  \[ \Theta_{v,w}\colon W \sto V\colon x \mapsto v\langle w \vert\, x\rangle\,. \]
(Bandingkan dengan~\ref{0022431fa}.)
\end{notn}

\begin{prop} Pemetaan $\Theta$ yang didefinisikan di atas bersifat seskuilinear.
\end{prop}

\begin{prop} Misalkan $V$ dan $W$ modul-modul Hilbert-$A$ dengan $A$ suatu aljabar-$C^*$. Untuk
setiap $v \in V$ dan $w \in W$, pemetaan $\Theta_{v,w}$ dapat diadjoinkan dan $(\Theta_{v,w})^*
= \Theta_{w,v}$.
\end{prop}

Proposisi berikut memperumum proposisi~\ref{prop_op_act_otimes1} dan~\ref{prop_op_act_otimes2}.

\begin{prop} Misalkan $U$, $V$, $W$, dan $Z$ modul-modul Hilbert-$A$ dengan $A$ suatu aljabar-$C^*$.
Jika $S \in \ofml L(Z,W)$ dan $T \in \ofml L(V,U)$, maka
  \[ T \Theta_{v,w} = \Theta_{Tv,w} \qquad \text{ dan } \qquad
                               \Theta_{v,w} S = \Theta_{v,S^*w} \]
untuk semua $v \in V$ dan $w \in W$.
  \[ Z \to^S W \to^{\Theta_{v,w}} V \to^T U \]
\end{prop}

\begin{prop} Misalkan $U$, $V$, dan $W$ modul-modul Hilbert-$A$ dengan $A$ suatu aljabar-$C^*$.
Andaikan $u \in U$; $v$, $v' \in V$; dan $w \in W$. Maka
  \[ \Theta_{u,v}\Theta_{v',w} = \Theta_{u\langle v \vert\,v' \rangle, w}
                               = \Theta_{u, w\langle v'\vert\,v \rangle}\,. \]
\end{prop}

\begin{notn}\label{0038331} Misalkan $A$ suatu aljabar-$C^*$ dan $V$ serta $W$
modul-modul Hilbert-$A$. Kita nyatakan dengan
 \index{K@$\ofml K(V,W)$, $\ofml K(V)$}%
$\ofml K(W,V)$ rentang linear tertutup dari $\{\Theta_{v,w}\colon v \in V \text{ dan } w \in
W\}$. Seperti biasa, kita menyingkat $\ofml K(V,V)$ menjadi~$\ofml K(V)$.
\end{notn}

\begin{prop}\label{0038334} Jika $A$ suatu aljabar-$C^*$ dan $V$
suatu modul Hilbert-$A$, maka $\ofml K(V)$ adalah ideal dalam aljabar-$C^*$~$\ofml L(V)$.
\end{prop}

Contoh berikut dimaksudkan sebagai pembenaran bagi praktik standar untuk mengidentifikasi suatu
aljabar-$C^*$ $A$ dengan $\ofml K(A)$.

\begin{exam}\label{0038337} Jika suatu aljabar-$C^*$ $A$ kita pandang sebagai modul-$A$ (lihat
contoh~\ref{0038111}), maka $\cstariso{\ofml K(A)}A$.
\end{exam}

\begin{proof}[\emph{Petunjuk pembuktian}] Seperti dalam akibat~\ref{0013151}, untuk setiap
$a \in A$ definisikan operator perkalian kiri $L_a\colon A \sto A\colon x \mapsto ax$. Tunjukkan
bahwa setiap operator semacam itu dapat diadjoinkan dan bahwa pemetaan $L\colon A \sto \ofml L(A) \colon
a \mapsto L_a$ merupakan isomorfisme-$*\,$ ke suatu subaljabar-$C^*$ dari~$\ofml L(A)$. Kemudian
verifikasikan bahwa $\ofml K(A)$ adalah tutupan citra oleh $L$ dari rentang linear hasil kali
elemen-elemen~$A$. \ns
\end{proof}

\begin{exam}\label{0038341} Misalkan $H$ suatu ruang Hilbert yang dipandang sebagai modul-$\C$. Maka
$\ofml K(H)$ (sebagaimana didefinisikan dalam~\ref{0038331}) adalah ideal operator kompak pada~$H$ (lihat
proposisi~\ref{prop_K_is_clo_FR}).
\end{exam}

\begin{proof} Lihat~\cite{RaeburnW:1998}, contoh 2.27. \ns \end{proof}

Contoh sebelumnya telah mendorong banyak peneliti, ketika menangani suatu modul Hilbert sebarang~$V$,
menyebut elemen $\ofml K(V)$ sebagai \emph{operator kompak}. Istilah ini
meragukan karena operator semacam itu belum tentu kompak. (Sebagai contoh, jika
suatu aljabar-$C^*$ beridentitas berdimensi tak hingga $A$ kita pandang sebagai modul-$A$, maka
$\Theta_{\vc 1,\vc 1} = I_A$, tetapi operator identitas pada $A$ tidak kompak---lihat
contoh~\ref{cor_id_op_not_cpt}.)

Fakta bahwa modul Hilbert-$C^*$ hanya digunakan secara terbatas dalam catatan ini hendaknya
tidak membuat Anda berpikir bahwa kajiannya sangat khusus dan/atau kurang penting.
Sebaliknya,
kajian tersebut sekarang merupakan bidang penelitian penting dan dinamis yang memiliki penerapan
pada bidang-bidang yang beragam seperti teori-$K$, aljabar-$C^*$ graf, grup kuantum, probabilitas
kuantum, berkas vektor, geometri nonkomutatif, topologi aljabar dan topologi geometri,
ruang dan aljabar operator, serta wavelet. Lihatlah
laman web Michael Frank~\cite{Frank:2010}, \emph{Hilbert C*-modules and related subjects---a guided
reference overview}, tempat ia mencantumkan 1531 referensi (pada pembaruan 11.09.10) berupa buku,
makalah, dan tesis tentang modul semacam itu dan mengelompokkannya menurut penerapan. Terdapat
grafik menarik (pada halaman 9) yang menggambarkan pertumbuhan bidang
matematika ini. Grafik tersebut mencakup materi dari upaya perintis pada tahun 1950-an dan awal 1960-an
(0--2 makalah per tahun) hingga saat catatan ini ditulis (sekitar 100 makalah per tahun).
































\section{Ideal Esensial}
\begin{exam} Jika $A$ dan $B$ adalah aljabar-$C^*$, maka $A$ (lebih tepatnya, $A \oplus \{\vc
0\}$) merupakan ideal dalam $A \oplus B$.
\end{exam}

\begin{conv} Sebagaimana ditunjukkan oleh contoh sebelumnya, lazimnya $A$ dipandang sebagai
 \index{konvensi!$A$ merupakan subhimpunan dari $A \oplus B$}%
subhimpunan dari $A \oplus B$. Selanjutnya kita akan melakukan ini tanpa menyebutkannya lagi.
\end{conv}

\begin{notn} Untuk suatu elemen $c$ dari aljabar $A$
 \index{Ic@$I_c$}%
tetapkan
  \[ I_c := \biggl\{a_0c + cb_0 + \sum_{k=1}^p a_kcb_k \colon
   \text{$p \in \N$, $a_0$, \dots, $a_p$, $b_0$, \dots, $b_p \in A$}\biggr\} \]
\end{notn}

\begin{prop} Jika $c$ merupakan elemen dari aljabar $A$, maka $I_c$ merupakan ideal
(aljabar) dalam~$A$.
\end{prop}

Perhatikan bahwa dalam proposisi sebelumnya tidak dinyatakan bahwa ideal aljabar $I_c$
harus sejati. Sangat mungkin $I_c = A$ (misalnya, ketika $c$ merupakan
elemen invertibel dari aljabar beridentitas).

\vskip 2 pt

\begin{defn} Misalkan $c$ suatu elemen dari aljabar-$C^*$ $A$.
 \index{Jc@$J_c$ (ideal utama (tertutup) yang memuat~$c$)}%
 \index{utama!ideal}%
 \index{ideal!utama}%
Definisikan $J_c$, yakni \df{ideal utama} yang memuat $c$, sebagai irisan dari
keluarga semua ideal (tertutup) dalam $A$ yang memuat~$c$. Jelas bahwa $J_c$ merupakan ideal terkecil
yang memuat~$c$.
\end{defn}

\begin{prop} Dalam suatu aljabar-$C^*$, tutupan suatu ideal aljabar merupakan ideal.
\end{prop}

\begin{exam} Tutupan suatu ideal aljabar sejati dalam aljabar-$C^*$ tidak harus merupakan
ideal sejati. Sebagai contoh, $l_c$, yaitu himpunan barisan bilangan kompleks yang
pada akhirnya nol, padat dalam aljabar-$C^*$ $l_0 = \fml C_0(\N)$. (Namun, ingat
proposisi~\ref{0006801}.)
\end{exam}

\begin{prop} Jika $c$ suatu elemen dari aljabar-$C^*$, maka $J_c = \clo{I_c}$.
\end{prop}

\begin{notn} Kita menggunakan suatu konvensi notasi standar. Jika $A$ dan $B$ merupakan subhimpunan tak kosong
dari suatu aljabar, maka $AB$ berarti rentang linear dari hasil kali elemen-elemen dalam $A$ dan
elemen-elemen dalam $B$; yaitu,
 \index{<AB@$AB$ (dalam aljabar, rentang hasil kali elemen-elemen)}%
 \index{konvensi!dalam suatu aljabar, $AB$ menyatakan rentang hasil kali}%
$AB = \spn\{ab \colon \text{$a \in A$ dan $b \in B$}\}$. (Perhatikan bahwa dalam
definisi~\ref{000551} tidak ada bedanya apakah $AJ$ diartikan sebagai himpunan
hasil kali elemen-elemen dalam $A$ dengan elemen-elemen dalam $J$ atau rentang dari himpunan tersebut.)
\end{notn}

\begin{prop} Jika $I$ dan $J$ merupakan ideal dalam suatu aljabar-$C^*$, maka $\clo{IJ} =
I \cap J$.
\end{prop}

\vskip 2 pt

Suatu aljabar-$C^*$ tak beridentitas $A$ dapat dibenamkan sebagai ideal dalam aljabar-$C^*$ beridentitas melalui
berbagai cara. Aljabar-$C^*$ beridentitas terkecil yang memuat $A$ adalah unitalisasinya~$\wt
A$ (lihat proposisi~\ref{0015213}). Tentu saja tidak ada aljabar-$C^*$ beridentitas terbesar tempat
$A$ dapat dibenamkan sebagai ideal, sebab jika $A$ dibenamkan sebagai ideal dalam suatu
aljabar-$C^*$ beridentitas $B$ dan $C$ adalah \emph{sebarang} aljabar-$C^*$ beridentitas, maka $A$ merupakan ideal dalam
aljabar-$C^*$ beridentitas $B \oplus C$ yang lebih besar lagi. Alasan unitalisasi yang lebih besar ini
tidak begitu menarik adalah karena irisan ideal $C$ dengan $A$ ialah~$\{\vc 0\}$.
Hal ini memotivasi definisi berikut.

\vskip 2 pt

\begin{defn} Suatu ideal $J$ dalam aljabar-$C^*$ $A$ disebut
 \index{esensial!ideal}%
 \index{ideal!esensial}%
\df{esensial} jika dan hanya jika $I \cap J \ne \vc 0$ untuk setiap ideal tak nol $I$ dalam~$A$.
\end{defn}

\begin{exam} Suatu aljabar-$C^*$ $A$ merupakan ideal esensial dalam unitalisasinya $\wt A$ jika dan
hanya jika $A$ \emph{tidak} beridentitas.
\end{exam}

\begin{exam} Jika $H$ suatu ruang Hilbert, maka ideal operator kompak $\ofml K(H)$
merupakan ideal esensial dalam aljabar-$C^*$~$\ofml B(H)$.
\end{exam}

\begin{defn} Misalkan $J$ suatu ideal dalam aljabar-$C^*$~$A$. Kita mendefinisikan
$J^\perp$, yakni
 \index{anihilator}%
 \index{<superscript@$J^\perp$ (anihilator dari $J$)}%
\df{anihilator} dari $J$, sebagai $\bigl\{a \in A\colon Ja = \{\vc 0\}\bigr\}$.
\end{defn}

\begin{prop} Jika $J$ suatu ideal dalam aljabar-$C^*$, maka $J^\perp$ juga demikian.
\end{prop}

\begin{prop}\label{0038465} Suatu ideal $J$ dalam aljabar-$C^*$ $A$ bersifat esensial jika dan hanya
jika $J^\perp = \{\vc 0\}$.
\end{prop}

\begin{prop}\label{0038468} Jika $J$ suatu ideal dalam aljabar-$C^*$, maka $\bigl(J \oplus
J^\perp\bigr)^\perp = \{\vc 0\}$.
\end{prop}

\begin{notn} Misalkan $f$ suatu fungsi bernilai kompleks (atau real) pada himpunan $S$. Maka
  \[ Z_f := \{ s \in S\colon f(s) = 0\}. \]
Ini adalah
 \index{himpunan nol}%
 \index{Zf@$Z_f$ (himpunan nol dari~$f$)}%
\df{himpunan nol} dari~$f$.
\end{notn}

 \vskip 2 pt

Misalkan $A$ suatu aljabar-$C^*$ komutatif tak beridentitas. Menurut
\emph{teorema Gelfand--Naimark kedua}~\ref{00152156}, terdapat suatu ruang Hausdorff kompak lokal tak kompak
$X$ sedemikian sehingga $A = \fml C_0(X)$. (Di sini, tentu saja, kita membolehkan
diri melakukan penyalahgunaan bahasa yang lazim: secara harfiah, persamaan yang ditunjukkan itu
seharusnya berupa isomorfisme-$*\,$ isometrik.) Sekarang misalkan $Y$ suatu ruang Hausdorff kompak
yang memuat $X$ sebagai subhimpunan terbuka dan misalkan $B = \fml C(Y)$. Maka $B$ suatu aljabar-$C^*$
komutatif beridentitas. Pandang $A$ sebagai ideal dalam $B$ melalui pemetaan
$\iota\colon A \sto B\colon f \mapsto \wt f$ dengan
  \[ \wt f(y) = \left\{%
                \begin{array}{ll}
                         f(y), & \hbox{jika $y \in X$;} \\
                         0, & \hbox{selainnya.} \\
                \end{array}%
                \right.\]
Perhatikan bahwa himpunan tertutup $X^c$ adalah $\bigcap\{Z_{\wt f}\,\colon f \in \fml C_0(X)\}$.

\vskip 2 pt

\begin{prop}\label{0038474} Gunakan notasi seperti pada paragraf sebelumnya. Maka
ideal $A$ bersifat esensial dalam $B$ jika dan hanya jika subhimpunan terbuka $X$ padat dalam~$Y$.
\end{prop}

Dengan demikian, sifat esensial suatu ideal dalam konteks unitalisasi aljabar-$C^*$
komutatif tak beridentitas berkorespondensi secara tepat dengan sifat suatu subruang terbuka
yang padat dalam konteks kompaktifikasi ruang Hausdorff kompak lokal
tak kompak.





















\section{Kompaktifikasi dan Unitalisasi}
Dalam definisi~\ref{00152131}, objek yang keberadaannya dibuktikan dalam
proposisi sebelumnya~\ref{0015213} kita sebut ``unitalisasi'' dari suatu aljabar-$C^*$. Penyebutan seolah-olah tunggal
tersebut jelas menyesatkan. Sebagaimana ruang topologis dapat memiliki banyak
kompaktifikasi berbeda, suatu aljabar-$C^*$ dapat memiliki banyak unitalisasi. Jika $A$ merupakan
aljabar-$C^*$ komutatif tak beridentitas, jelas dari akibat~\ref{001924} bahwa
unitalisasi $\wt A$ adalah unitalisasi terkecil yang mungkin dari~$A$. Demikian pula, dalam topologi,
jika $X$ suatu ruang Hausdorff kompak lokal tak kompak, maka kompaktifikasi
satu titiknya jelas merupakan kompaktifikasi terkecil yang mungkin dari~$X$. Ingat bahwa
dalam proposisi~\ref{00152194} kita membuktikan bahwa membangun unitalisasi terkecil
dari $A$ ``pada dasarnya'' sama dengan membangun kompaktifikasi terkecil
dari~$X$.

 \vskip 2 pt

Kadang-kadang lebih nyaman (seperti, misalnya, pada paragraf sebelumnya) untuk mengambil suatu
``unitalisasi'' dari aljabar yang sudah beridentitas dan kadang-kadang lebih nyaman untuk mengambil suatu
``kompaktifikasi'' dari ruang yang sudah kompak. Karena tampaknya tidak ada
istilah yang diterima secara universal, saya memperkenalkan bahasa berikut (jelas tidak standar, tetapi
saya harap berguna).

\begin{defn}\label{0038614} Misalkan $A$ dan $B$ aljabar-$C^*$ serta $X$ dan $Y$
ruang topologis Hausdorff. Kita mengatakan bahwa
 \begin{enumerate}
    \item $B$ merupakan suatu \df{unitalisasi} dari $A$ jika $B$ beridentitas dan $A$ ($*\,$-isomorfik
dengan) suatu subaljabar-$C^*$ dari~$B$;
 \index{unitalisasi}%
 \index{unitalisasi!esensial}%
 \index{esensial!unitalisasi}%
    \item $B$ merupakan suatu \df{unitalisasi esensial} dari $A$ jika $B$ beridentitas dan $A$
($*\,$-isomorfik dengan) suatu ideal esensial dari~$B$;
 \index{kompaktifikasi}%
 \index{kompaktifikasi!esensial}%
 \index{esensial!kompaktifikasi}%
    \item $Y$ merupakan suatu \df{kompaktifikasi} dari $X$ jika ruang itu kompak dan $X$ (homeomorfik
dengan) suatu subruang dari~$Y$; dan
    \item $Y$ merupakan suatu \df{kompaktifikasi esensial} dari $X$ jika ruang itu kompak dan $X$
(homeomorfik dengan) suatu subruang padat dari~$Y$.
 \end{enumerate}
\end{defn}

Kiranya beberapa kata perlu disampaikan mengenai bagian-bagian dalam tanda kurung pada definisi
sebelumnya. Jarang terjadi bahwa suatu ruang topologis $X$ secara harfiah merupakan
\emph{subhimpunan} dari kompaktifikasi tertentu atas $X$ atau bahwa suatu aljabar-$C^*$ $A$ merupakan
\emph{subhimpunan} dari unitalisasi tertentu atas~$A$. Walaupun tentu benar bahwa sering kali
lebih nyaman untuk memandang satu aljabar-$C^*$ sebagai subhimpunan dari aljabar lain ketika sebenarnya yang
pertama hanya $*\,$-isomorfik dengan subhimpunan dari yang kedua, ada pula keadaan ketika
perincian menjadi lebih jelas jika pembenaman yang sesungguhnya ditentukan. Jika rincian perbedaan
ini belum sepenuhnya akrab, dua definisi berikut dimaksudkan untuk membantu.

\begin{defn} Misalkan $A$ dan $B$ aljabar-$C^*$. Kita mengatakan bahwa $A$
\df{dibenamkan} dalam~$B$ jika terdapat suatu homomorfisme-$*\,$ injektif $\iota\colon A \sto
B$; yaitu, jika $A$ isomorfik-$*\,$ dengan suatu subaljabar-$C^*$ dari~$B$ (lihat
proposisi~\ref{0019023} dan~\ref{00152181}). Homomorfisme-$*\,$ injektif $\iota$
merupakan suatu
 \index{C*@pembenaman-$C^*$}%
 \index{pembenaman!aljabar-$C^*$}%
\df{pembenaman} dari $A$ ke dalam~$B$. Dalam situasi ini lazimnya $A$ dan
jangkauan~$\iota$ diperlakukan sebagai aljabar-$C^*$ yang identik. Pasangan $(B,\iota)$ merupakan suatu
 \index{unitalisasi}%
\df{unitalisasi} dari~$A$ jika $B$ merupakan aljabar-$C^*$ beridentitas dan $\iota\colon A \sto B$ merupakan suatu
pembenaman. Unitalisasi $(B,\iota)$ bersifat
  \index{unitalisasi!esensial}%
  \index{esensial!unitalisasi}%
\df{esensial} jika jangkauan $\iota$ merupakan ideal esensial dalam~$B$.
\end{defn}

\begin{defn} Misalkan $X$ dan $Y$ ruang topologis Hausdorff. Kita mengatakan bahwa $X$
\df{dibenamkan} dalam~$Y$ jika terdapat homeomorfisme $j$ dari $X$ ke suatu subruang dari~$Y$.
Homeomorfisme $j$ merupakan suatu
 \index{topologis!pembenaman}%
 \index{pembenaman!topologis}%
\df{pembenaman} dari $X$ ke dalam~$Y$. Seperti dalam aljabar-$C^*$, lazimnya
jangkauan $j$ diidentifikasi dengan ruang~$X$. Pasangan $(Y,j)$ merupakan suatu
 \index{kompaktifikasi}%
\df{kompaktifikasi} dari~$X$ jika $Y$ suatu ruang Hausdorff kompak dan $j\colon X \sto Y$
merupakan suatu pembenaman. Kompaktifikasi $(Y,j)$ bersifat
  \index{kompaktifikasi!esensial}%
  \index{esensial!kompaktifikasi}%
\df{esensial} jika jangkauan $j$ padat dalam~$Y$.
\end{defn}

Kita telah membahas unitalisasi terkecil dari suatu aljabar-$C^*$ dan kompaktifikasi terkecil
dari suatu ruang Hausdorff kompak lokal. Sekarang, bagaimana dengan aljabar beridentitas terbesar, atau bahkan
maksimal, yang memuat~$A$? Jelas bahwa objek seperti itu tidak ada, sebab jika $B$ suatu
aljabar beridentitas yang memuat~$A$, maka $B \oplus C$ juga demikian, dengan $C$ sebarang
aljabar-$C^*$ beridentitas. Demikian pula, tidak ada ruang kompak terbesar yang memuat~$X$: jika $Y$ suatu
ruang kompak yang memuat $X$, maka gabungan saling lepas topologis $Y \uplus K$ juga demikian,
dengan $K$ sebarang ruang kompak tak kosong. Namun, masuk akal untuk menanyakan apakah
terdapat suatu unitalisasi esensial maksimal dari aljabar-$C^*$ atau suatu kompaktifikasi esensial
maksimal dari ruang Hausdorff kompak lokal. Jawabannya adalah \emph{ya} dalam kedua
kasus. \emph{Kompaktifikasi Stone--\v{C}ech} $\beta(X)$ yang terkenal bersifat maksimal di antara
kompaktifikasi esensial dari ruang Hausdorff kompak lokal tak kompak~$X$. Perinciannya
dapat ditemukan dalam sebarang buku teks topologi yang baik. Salah satu uraian standar yang mudah dibaca
adalah~\cite{Willard:1968}, butir 19.3--19.12. Pendekatan yang lebih canggih menggunakan sebagian
analisis fungsional---lihat, misalnya, \cite{Conway:1990}, bab V, bagian~6. Ternyata
terdapat pula unitalisasi esensial maksimal dari suatu aljabar-$C^*$ tak beridentitas
$A$---objek ini disebut \emph{aljabar pengali} dari~$A$.

Kita mengatakan bahwa suatu unitalisasi esensial $M$ dari aljabar-$C^*$ $A$ bersifat \emph{maksimal} jika setiap
aljabar-$C^*$ yang memuat $A$ sebagai ideal esensial dapat dibenamkan dalam~$M$. Berikut pernyataan yang lebih
formal.

\begin{defn} Suatu unitalisasi esensial $(M,j)$ dari aljabar-$C^*$ $A$ disebut
 \index{maksimal!unitalisasi esensial}%
 \index{esensial!unitalisasi!maksimal}%
 \index{unitalisasi!esensial maksimal}%
\df{maksimal} jika untuk setiap pembenaman $\iota\colon A \sto B$ yang jangkauannya merupakan ideal esensial
dalam~$B$, terdapat homomorfisme-$*\,$ $\phi\colon B \sto M$ sedemikian sehingga $\phi
\circ \iota = j$.
\end{defn}

\begin{prop} Dalam definisi sebelumnya, homomorfisme-$*\,$ $\phi$, jika ada, harus
injektif.
\end{prop}

\begin{prop} Dalam definisi sebelumnya, homomorfisme-$*\,$ $\phi$, jika ada, harus
tunggal.
\end{prop}

Bandingkan definisi berikut dengan~\ref{0028}.

\begin{defn} Misalkan $A$ dan $B$ aljabar-$C^*$, dan $V$ modul Hilbert-$A$. Pemetaan
$\phi\colon B \sto \ofml L(V)$ yang merupakan homomorfisme-$*\,$ disebut
 \index{tak terdegenerasi!homomorfisme-$*\,$}%
\df{tak terdegenerasi} jika $\phi^\sto(B)\,V$ padat dalam~$V$.
\end{defn}

\begin{prop} Misalkan $A$, $B$, dan $J$ aljabar-$C^*$, $V$ suatu modul Hilbert-$B$, serta
$\iota\colon J \sto A$ suatu homomorfisme-$*\,$ injektif yang jangkauannya merupakan ideal dalam~$A$.
Jika $\phi\colon J \sto \ofml L(V)$ suatu homomorfisme-$*\,$ tak terdegenerasi, maka terdapat
perpanjangan tunggal dari $\phi$ menjadi homomorfisme-$*\,$ $\overline\phi\colon A \sto
\ofml L(V)$ yang memenuhi $\overline\phi \circ \iota = \phi$.
\end{prop}

\begin{prop}\label{0038621} Jika $A$ suatu aljabar-$C^*$ tak nol, maka $(\ofml L(A),L)$ merupakan
unitalisasi esensial maksimal dari~$A$. Unitalisasi ini tunggal dalam arti bahwa jika $(M,j)$ merupakan
unitalisasi esensial maksimal lain dari $A$, maka terdapat suatu isomorfisme-$*\,$ $\phi
\colon M \sto \ofml L(A)$ sedemikian sehingga $\phi \circ j = L$.
\end{prop}

\begin{defn}\label{0038641} Misalkan $A$ suatu aljabar-$C^*$. Kita mendefinisikan
 \index{aljabar pengali}%
 \index{aljabar!pengali}%
\df{aljabar pengali} dari $A$ sebagai keluarga $\ofml L(A)$ operator yang dapat diadjoinkan
pada~$A$. Mulai sekarang kita menyatakan keluarga ini
 \index{MA@$M(A)$ (aljabar pengali dari $A$)}%
dengan $M(A)$.
\end{defn}









\endinput

 \chapter{TEORI FREDHOLM}

\section{Alternatif Fredholm}

Pada tahun 1903, Erik Ivar Fredholm menerbitkan sebuah makalah rintisan tentang persamaan integral dalam jurnal
\emph{Acta Mathematica}. Di antara banyak hasil pentingnya terdapat teorema yang kini kita kenal sebagai
\emph{alternatif Fredholm}. Kita menyatakan suatu versi hasil ini dalam bahasa
yang tersedia bagi Fredholm pada masa itu.

\begin{prop}[Alternatif Fredholm I]\label{004011} Tetapkan $\lambda \in \C\setminus\{0\}$ dan misalkan $k$ suatu fungsi kontinu bernilai kompleks
pada persegi satuan $[0,1] \times [0,1]$. \textbf{Entah} persamaan-persamaan
 \index{alternatif Fredholm!versi I}%
tak homogen
 \begin{align}
   \lambda f(s) &- \int_0^1 k(s,t)f(t)\,dt = g(s)\label{004011i} \text{ dan}    \tag{1}\\
   \conj{\lambda} h(s) &- \int_0^1 \conj{k(t,s)}h(t)\,dt = j(s)\label{004011ii} \tag{2}
 \end{align}
memiliki solusi $f$ dan $h$ untuk setiap $g$ dan $j$ yang diberikan, masing-masing, dengan solusi-solusi tersebut
tunggal; dalam hal ini persamaan homogen yang bersesuaian
 \begin{align} \lambda f(s) &- \int_0^1 k(s,t)f(t)\,dt = 0\label{004011iii} \text{ dan} \tag{3}\\
        \conj{\lambda} h(s) &- \int_0^1 \conj{k(t,s)}h(t)\,dt = 0\label{004011iv} \tag{4}
 \end{align}
hanya memiliki solusi trivial; \textbf{---atau---}\\
persamaan homogen \eqref{004011iii} dan \eqref{004011iv} masing-masing memiliki
sejumlah hingga (dan tak nol) solusi bebas linear yang sama banyaknya, yaitu $f_1,\dots,f_n$ dan $h_1,\dots,h_n$,
secara berturut-turut; dalam hal ini persamaan tak homogen \eqref{004011i} dan
\eqref{004011ii} memiliki solusi jika dan hanya jika $g$ dan $j$ memenuhi
 \begin{align} \int_0^1 h_k(t)\conj{g(t)}\,dt &= 0\label{004011v} \text{ dan} \tag{5}\\
              \int_0^1 j(t) \conj{f_k(t)}\,dt &= 0\label{004011vi} \tag{6}
 \end{align}
untuk $k = 1, \dots, n$.
\end{prop}

Pada tahun 1906, David Hilbert telah menyadari bahwa pengintegralan itu sendiri sangat sedikit berkaitan dengan
kebenaran hasil ini. Ia menemukan bahwa yang penting adalah kekompakan
operator integral yang dihasilkan. (Pada awal abad ke-$20$
\emph{kekompakan} disebut \emph{kontinuitas lengkap}. Istilah \emph{ruang Hilbert}
telah digunakan sejak tahun 1911 sehubungan dengan ruang barisan~$l_2$. Baru pada
tahun 1929 John von Neumann memperkenalkan gagasan---dan aksioma yang
mendefinisikan---\emph{ruang Hilbert abstrak}.) Jadi, berikut suatu versi yang agak lebih mutakhir dari
\emph{alternatif Fredholm}.

\begin{prop}[Alternatif Fredholm II]\label{004031} Jika $K$ suatu operator kompak pada ruang
 \index{alternatif Fredholm!versi II}%
Hilbert, $\lambda \in \C\setminus\{0\}$, dan $T = \lambda I - K$, maka \textbf{entah} persamaan-persamaan
tak homogen
 \begin{align}
            Tf &= g\label{004031i} \text{ dan} \tag{1'}\\
            T^* h &= j\label{004031ii} \tag{2'}
 \end{align}
memiliki solusi $f$ dan $h$ untuk setiap $g$ dan $j$ yang diberikan, masing-masing, dengan solusi-solusi tersebut
tunggal; dalam hal ini persamaan homogen yang bersesuaian
 \begin{align}
            Tf &= 0\label{004031iii} \text{ dan} \tag{3'}\\
          T^*h &= 0\label{004031iv} \tag{4'}
 \end{align}
hanya memiliki solusi trivial; \textbf{---atau---}\\
persamaan homogen \eqref{004031iii} dan \eqref{004031iv} masing-masing memiliki
sejumlah hingga (dan tak nol) solusi bebas linear yang sama banyaknya, yaitu $f_1,\dots,f_n$ dan $h_1,\dots,h_n$,
secara berturut-turut; dalam hal ini persamaan tak homogen \eqref{004031i} dan
\eqref{004031ii} memiliki solusi jika dan hanya jika $g$ dan $j$ memenuhi
 \begin{align} h_k &\perp g\label{004031v} \text{ dan} \tag{5'}\\
                 j &\perp f_k\label{004031vi} \tag{6'}
 \end{align}
untuk $k = 1, \dots, n$.
\end{prop}

Perhatikan bahwa dengan menggunakan beberapa fakta elementer mengenai kernel dan jangkauan
operator serta keortogonalan dalam ruang Hilbert, kita dapat meringkas pernyataan
dari~\ref{004031} secara cukup banyak.

\begin{prop}[Alternatif Fredholm IIIa]\label{004033} Jika $T = \lambda I - K$ dengan
 \index{alternatif Fredholm!versi IIIa}%
$K$ suatu operator kompak pada ruang Hilbert dan $\lambda \in \C\setminus\{0\}$, maka
 \begin{enumerate}
  \item $T$ injektif jika dan hanya jika bersifat surjektif,
  \item $\ran T^* = (\ker T)^\perp$, dan
  \item $\dim\ker T = \dim \ker T^*$.
 \end{enumerate}
Selain itu, syarat (1) dan (2) berlaku untuk $T^*$ maupun~$T$.
\end{prop}


















\section{Alternatif Fredholm -- lanjutan}
\begin{defn}\label{004034} Suatu operator $T$ pada ruang Banach disebut
 \index{operator Riesz--Schauder}%
 \index{operator!Riesz--Schauder}%
\df{operator Riesz--Schauder} jika dapat ditulis dalam bentuk $T = S + K$ dengan $S$
invertibel dan $K$ kompak.
\end{defn}

Materi dalam bagian~\ref{sec_onbases} dan definisi sebelumnya memungkinkan kita
memperumum versi \emph{alternatif Fredholm} yang diberikan dalam~\ref{004033} ke ruang
Banach.

\begin{prop}[Alternatif Fredholm IIIb]\label{0040343} Jika $T$ suatu operator Riesz--Schauder
 \index{alternatif Fredholm!versi IIIb}%
pada ruang Banach, maka
 \begin{enumerate}
  \item $T$ injektif jika dan hanya jika bersifat surjektif,
  \item $\ran T^* = (\ker T)^\perp$, dan
  \item $\dim\ker T = \dim \ker T^* < \infty$.
 \end{enumerate}
Selain itu, syarat (1) dan (2) berlaku untuk $T^*$ maupun~$T$.
\end{prop}

\begin{prop} Jika $M$ subruang tertutup dari ruang Banach $B$, maka $M^\perp \cong (B/M)^*$.
\end{prop}

\begin{proof} Lihat \cite{Conway:1990}, halaman 89. \ns \end{proof}

\begin{defn} Jika $T\colon V \sto W$ suatu pemetaan linear antara ruang vektor, maka
 \index{kokernel}%
\df{kokernel} pemetaan tersebut didefinisikan oleh
  \[ \coker T = W/ \ran T\,. \]
Ingat bahwa
 \index{kodimensi}%
\df{kodimensi} dari suatu subruang $U$ dalam ruang vektor $V$ adalah dimensi dari~$V/U$. Jadi,
apabila $T$ suatu pemetaan linear, $\dim\coker T = \codim\ran T$.
\end{defn}

Dalam kategori $\cat{HIL}$ dari ruang Hilbert dan pemetaan linear kontinu, jangkauan suatu
morfisme tidak harus menjadi objek kategori tersebut. Secara khusus, jangkauan suatu operator
tidak harus tertutup.

\begin{exam}\label{004066} Operator ruang Hilbert
   \[ T\colon l_2 \sto l_2\colon (x_1,x_2,x_3,\dots) \mapsto
                                   (x_1, \tfrac12 x_2, \tfrac13 x_3, \dots) \]
bersifat injektif, swaadjoin, kompak, dan kontraktif,
 \index{operator!dengan jangkauan tak tertutup}%
tetapi jangkauannya, walaupun padat dalam $l_2$, tidak mencakup seluruh~$l_2$.
\end{exam}

Walaupun tidak langsung berkaitan dengan tujuan kita saat ini, inilah tempat yang tepat untuk mencatat
fakta bahwa jumlah dua subruang tertutup dari suatu ruang Hilbert tidak harus tertutup.

\begin{exam}\label{004071} Misalkan $T$ operator pada ruang Hilbert $l_2$ yang didefinisikan dalam
 \index{subruang!jumlah subruang tertutup tidak harus tertutup}%
contoh~\ref{004066}, $M = l_2 \oplus \{\vc 0\}$, dan $N$ graf dari~$T$. Maka $M$
dan $N$ keduanya merupakan subruang tertutup dari ruang Hilbert $l_2 \oplus l_2$, tetapi $M + N$
tidak tertutup.
\end{exam}

\begin{proof} Verifikasi contoh~\ref{004071} diperoleh dengan mudah dari hasil berikut
dan contoh~\ref{004066}. \ns
\end{proof}

\begin{prop} Misalkan $H$ suatu ruang Hilbert, $T \in \ofml B(H)$, $M = H \oplus \{\vc 0\}$, dan $N$
merupakan graf dari~$T$. Maka
 \begin{enumerate}
   \item[(a)] Himpunan $N$ merupakan subruang dari $H \oplus H$.
   \item[(b)] Operator $T$ injektif jika dan hanya jika $M \cap N = \{\,(\vc 0,\vc 0)\,\}$.
   \item[(c)] Jangkauan $T$ padat dalam $H$ jika dan hanya jika $M + N$ padat dalam $H \oplus H$.
   \item[(d)] Operator $T$ surjektif jika dan hanya jika $M + N = H \oplus H$.
 \end{enumerate}
\end{prop}

Kabar baik bagi teori operator Fredholm adalah bahwa operator dengan kokernel
berdimensi hingga secara otomatis mempunyai jangkauan tertutup.

\begin{prop}\label{00411} Jika suatu pemetaan linear terbatas $A\colon H \sto K$ antara ruang Hilbert
memiliki kokernel berdimensi hingga, maka $\ran A$ tertutup dalam~$K$.
\end{prop}

Dalam \emph{Alternatif Fredholm IIIb}~\ref{0040343}, kita mengamati bahwa syarat (1)
berlebih dan juga bahwa (2) berlaku untuk \emph{sebarang} operator ruang Banach dengan jangkauan tertutup.
Hal ini memungkinkan kita menyatakan kembali~\ref{0040343} secara lebih hemat.

\begin{prop}[Alternatif Fredholm IV]\label{00413}
 \index{alternatif Fredholm!versi IV}%
Jika $T$ suatu operator Riesz--Schauder pada ruang Banach, maka
 \begin{enumerate}
    \item $T$ mempunyai jangkauan tertutup dan
    \item  $\dim\ker T = \dim \ker T^* < \infty$.
 \end{enumerate}
\end{prop}














\section{Operator Fredholm}
\begin{defn} Misalkan $H$ adalah ruang Hilbert dan $\ofml K(H)$ merupakan ideal operator kompak di dalam
aljabar-$C^*$~$\ofml B(H)$. Maka aljabar hasil bagi (lihat proposisi~\ref{001902})
$\ofml Q(H) := \ofml B(H)/\ofml K(H)$ adalah
 \index{aljabar Calkin}%
 \index{aljabar!Calkin}%
 \index{Q@$\ofml Q(H)$ (aljabar Calkin)}%
\df{aljabar Calkin}. Seperti biasa, pemetaan hasil bagi dari $\ofml B(H)$ ke $\ofml
B(H)/\ofml K(H)$
 \index{pi@$\pi$ (pemetaan hasil bagi)}%
 \index{pi@$\pi(T)$ (elemen aljabar Calkin yang memuat $T$)}%
dilambangkan dengan $\pi$, sehingga jika $T \in \ofml B(H)$ maka $\pi(T) = [T] = T + \ofml K(H)$ adalah
elemen yang bersesuaian dengannya dalam aljabar Calkin. Suatu elemen $T \in \ofml B(H)$ adalah
 \index{Fredholm!operator}%
 \index{operator!Fredholm}%
\df{operator Fredholm} jika $\pi(T)$ invertibel di dalam~$\ofml Q(H)$. Keluarga
 \index{F@$\ofml F(H)$ (operator Fredholm pada $H$)}%
semua operator Fredholm pada $H$ dilambangkan dengan~$\ofml F(H)$.
\end{defn}

\begin{prop} Jika $H$ ruang Hilbert, maka $\ofml F(H)$ merupakan subhimpunan terbuka swaadjoin
dari~$\ofml B(H)$ dan tertutup terhadap
 \index{perturbasi}%
perturbasi kompak (yakni, jika $T$ Fredholm dan $K$ kompak, maka $T + K$ juga Fredholm).
\end{prop}

\begin{thm}[Teorema Atkinson]\label{004352}
 \index{teorema Atkinson}%
Suatu operator ruang Hilbert bersifat Fredholm jika dan hanya jika kernel dan kokernelnya
berdimensi hingga.
\end{thm}

\begin{proof} Lihat \cite{Arveson:2002}, teorema 3.3.2; \cite{Blackadar:2006}, teorema I.8.3.6;
\cite{Douglas:1972}, teorema 5.17; \cite{HigsonR:2000}, teorema 2.1.4;
atau~\cite{Wegge-Olsen:1993}, teorema 14.1.1. \ns
\end{proof}

Menariknya, dimensi kernel dan kokernel suatu operator Fredholm bukanlah besaran yang
sangat penting. Namun, selisih keduanya ternyata sangat penting.

\begin{defn} Jika $T$ adalah operator Fredholm pada suatu ruang Hilbert, maka
 \index{indeks!operator Fredholm}%
 \index{Fredholm!indeks (\seeonly{indeks})}%
\df{indeks Fredholm} operator tersebut (atau cukup \df{indeks}) didefinisikan oleh
   \[ \ind T := \dim\ker T - \dim\coker T\,. \]
\end{defn}

\begin{exam}\label{0044125} Setiap operator ruang Hilbert yang invertibel bersifat Fredholm
 \index{indeks!operator invertibel}%
 \index{operator!invertibel!indeks}%
 \index{invertibel!operator!indeks}%
dengan indeks nol.
\end{exam}

\begin{exam}\label{0044128} Indeks Fredholm dari
 \index{indeks!operator geser unilateral}%
 \index{operator geser unilateral!indeks}%
 \index{operator!geser unilateral!indeks}%
operator geser unilateral adalah~$-1$.
\end{exam}

\begin{exam} Untuk pemetaan antar-ruang, kita memakai konvensi standar: pemetaan linear disebut Fredholm jika
jangkauannya tertutup serta kernel dan kokernelnya berdimensi hingga. Dengan konvensi ini, setiap pemetaan linear $T\colon V \sto W$ antara ruang-ruang vektor berdimensi hingga bersifat Fredholm dan
$\ind T = \dim V - \dim W$. \end{exam}

\begin{exam}\label{004413} Jika $T$ adalah operator Fredholm pada suatu ruang Hilbert,
 \index{indeks!adjoin operator}%
 \index{operator!adjoin!indeks}%
 \index{adjoin!indeks}%
maka $\ind T^* = -\ind T$.
\end{exam}

\begin{exam} Indeks setiap operator Fredholm
 \index{indeks!operator Fredholm normal}%
 \index{operator!Fredholm normal!indeks}%
 \index{normal!operator Fredholm!indeks}%
normal adalah~$0$.
\end{exam}

\begin{lem}\label{004432} Misalkan $S\colon U \sto V$ dan $T\colon V \sto W$ adalah transformasi linear
antara ruang-ruang vektor. Maka (terdapat pemetaan linear sedemikian sehingga) barisan
berikut eksak.
 \[ \vc 0 \to \ker S \to \ker TS \to \ker T \to \coker S
                            \to \coker TS \to \coker T \to \vc 0\,. \]
\end{lem}

\begin{lem}\label{004433} Jika $V_0, V_1, \dots, V_n$ adalah ruang-ruang vektor berdimensi hingga
dan barisan
  \[ \vc 0 \to V_0 \to V_1 \to \dots \to V_n \to \vc 0 \]
eksak, maka
  \[ \sum_{k=0}^n (-1)^k \dim V_k = 0\,. \]
\end{lem}

\begin{prop} Misalkan $H$ ruang Hilbert berdimensi tak hingga. Maka himpunan $\ofml F(H)$ yang terdiri atas operator Fredholm
pada $H$ merupakan semigrup terhadap komposisi dan fungsi indeks $\ind$ adalah epimorfisme
dari $\ofml F(H)$ ke semigrup aditif bilangan bulat~$\Z$.
\end{prop}

\begin{proof} \emph{Petunjuk.} Gunakan \ref{004432}, \ref{004433}, \ref{0044125}, \ref{0044128},
dan~\ref{004413}. \ns
\end{proof}





















\section{Alternatif Fredholm -- Penutup}
Hasil berikut menyiratkan bahwa setiap operator Fredholm berindeks nol dapat ditulis sebagai
jumlah suatu operator invertibel dan operator kompak.

\begin{prop}\label{004712} Jika $T$ adalah operator Fredholm berindeks nol pada suatu ruang Hilbert,
maka terdapat isometri parsial berperingkat hingga $V$ sedemikian sehingga $T - V$ bersifat invertibel.
\end{prop}

\begin{proof} Lihat \cite{Pedersen:1995}, lemma 3.3.14 atau \cite{Wegge-Olsen:1993}, proposisi
14.1.3.  \ns
\end{proof}

\begin{lem} Jika $F$ adalah operator berperingkat hingga pada suatu ruang Hilbert, maka $I + F$ bersifat Fredholm dengan
indeks nol.
\end{lem}

\begin{proof} Lihat \cite{Pedersen:1995}, lemma 3.3.13 atau \cite{Wegge-Olsen:1993}, lemma 14.1.4. \ns
\end{proof}

\begin{notn} Misalkan $H$ adalah ruang Hilbert. Untuk setiap bilangan bulat $n$, kita melambangkan
 \index{F@$\ofml F_n(H)$, $\ofml F_n$ (operator Fredholm berindeks~$n$)}%
keluarga semua operator Fredholm berindeks $n$ pada $H$ dengan $\ofml F_n(H)$ atau cukup $\ofml
F_n$.
\end{notn}

\begin{prop}\label{004724} Untuk setiap ruang Hilbert berlaku $\ofml F_0 + \ofml K = \ofml F_0$.
\end{prop}

\begin{proof} Lihat \cite{Douglas:1972}, lemma 5.20 atau \cite{Wegge-Olsen:1993}, 14.1.5. \ns
\end{proof}

\begin{cor}[Alternatif Fredholm V]\label{004725}
 \index{alternatif Fredholm!versi V}%
Setiap operator Riesz--Schauder pada suatu ruang Hilbert bersifat Fredholm berindeks nol.
\end{cor}

\begin{proof} Ini segera mengikuti proposisi sebelumnya~\ref{004724} dan
contoh~\ref{0044125}.
\end{proof}

Kita sebenarnya telah membuktikan hasil yang lebih kuat: versi terakhir (yang cukup umum) dari
\emph{alternatif Fredholm}.

\begin{cor}[Alternatif Fredholm VI]\label{0047253} Suatu operator ruang Hilbert merupakan operator
 \index{alternatif Fredholm!versi VI}%
Riesz--Schauder jika dan hanya jika operator itu merupakan operator Fredholm berindeks nol.
\end{cor}

\begin{proof} Proposisi~\ref{004712} dan korolari~\ref{004725}.
\end{proof}

\begin{prop} Jika $T$ adalah operator Fredholm berindeks $n \in \Z$ pada ruang Hilbert $H$
dan $K$ operator kompak pada~$H$, maka $T + K$ juga merupakan operator Fredholm berindeks~$n$; yakni,
  \[ \ofml F_n + \ofml K = \ofml F_n\,. \]
\end{prop}

\begin{proof} Lihat \cite{HigsonR:2000}, proposisi 2.1.6; \cite{Pedersen:1995}, teorema
3.3.17; atau \cite{Wegge-Olsen:1993}, proposisi 14.1.6. \ns
\end{proof}

\begin{prop} Misalkan $H$ adalah ruang Hilbert. Maka $\ofml F_n(H)$ merupakan subhimpunan terbuka dari $\ofml B(H)$
untuk setiap bilangan bulat~$n$.
\end{prop}

\begin{proof} Lihat \cite{Wegge-Olsen:1993}, proposisi 14.1.8. \ns \end{proof}

\begin{defn}\label{004735} Suatu
 \index{lintasan}%
\df{lintasan} dalam ruang topologis $X$ adalah pemetaan kontinu dari interval $[0,1]$
ke~$X$. Dua titik $p$ dan $q$ dalam $X$ dikatakan
 \index{lintasan!titik-titik yang terhubung oleh}%
 \index{terhubung!oleh lintasan}%
\df{terhubung oleh lintasan} (atau
 \index{homotop!titik-titik dalam ruang topologis}%
\df{homotop dalam~$X$}) jika terdapat lintasan $f\colon [0,1] \sto X$ dalam $X$ sedemikian sehingga
$f(0) = p$ dan $f(1) = q$. Dalam hal ini kita
 \index{<binrelequivh@$p \sim_h q$ (homotopi titik-titik)}%
menulis $p \sim_h q$ dalam~$X$ (atau cukup $p \sim_h q$ ketika ruang $X$ jelas dari konteks).
\end{defn}


\begin{exam} Misalkan $a$ adalah elemen invertibel dalam aljabar-$C^*$ beridentitas $A$ dan $b$ adalah elemen $A$
sedemikian sehingga $\norm{a - b} < \norm{a^{-1}}^{-1}$. Maka $a \sim_h b$ dalam~$\inv(A)$.
\end{exam}

\begin{proof}[\emph{Petunjuk bukti}] Terapkan korolari~\ref{cor2_Neumann_series} pada titik-titik dalam ruas garis tertutup~$[a,b]$. \ns
\end{proof}

\begin{prop}\label{K002011} Relasi $\sim_h$ untuk ekuivalensi homotopi yang didefinisikan di atas merupakan suatu relasi
 \index{ekuivalensi!homotopi!titik-titik}%
ekuivalensi pada himpunan titik-titik suatu ruang topologis.
\end{prop}

\begin{defn} Jika $X$ adalah ruang topologis dan $\sim_h$ adalah relasi ekuivalensi homotopi,
maka kelas-kelas ekuivalensi yang dihasilkan adalah
 \index{lintasan!komponen}%
 \index{komponen!lintasan}%
\df{komponen lintasan} dari~$X$.
\end{defn}

Proposisi berikut mengidentifikasi komponen-komponen lintasan himpunan operator Fredholm sebagai
himpunan $\ofml F_n$ yang terdiri atas operator berindeks~$n$.

\begin{prop} Operator $S$ dan $T$ di dalam ruang operator Fredholm $\ofml F(H)$ pada
ruang Hilbert $H$ bersifat homotop dalam $\ofml F(H)$ jika dan hanya jika keduanya memiliki indeks yang sama.
\end{prop}

\begin{proof} Lihat \cite{Wegge-Olsen:1993}, korolari 14.1.9. \ns \end{proof}







\endinput

 \chapter{EKSTENSI}

\section{Operator Normal secara Esensial}
\begin{defn} Misalkan $T$ suatu operator pada ruang Hilbert~$H$.  
 \index{spektrum!esensial}%
 \index{esensial!spektrum}%
 \index{sigmae@$\sigma_e(T)$ (spektrum esensial dari $T$)}%
\df{spektrum esensial} dari $T$, yang dinotasikan dengan $\sigma_e(T)$, adalah spektrum citra
$T$ dalam aljabar Calkin $\ofml Q(H)$; yaitu,
   \[ \sigma_e(T) = \sigma_{\ofml Q(H)}(\pi(T))\,. \]
\end{defn}

\begin{prop} Jika $T$ suatu operator pada ruang Hilbert $H$, maka
  \[ \sigma_e(T) = \{\lambda \in \C\colon T - \lambda I \notin \ofml F(H)\}\,. \]
\end{prop}

\begin{prop} Spektrum esensial dari operator swaadjoin $T$ pada ruang Hilbert adalah gabungan
titik-titik akumulasi spektrum $T$ dengan nilai-nilai eigen $T$ yang memiliki
multiplisitas tak hingga.  Anggota-anggota $\sigma(T) \setminus \sigma_e(T)$ adalah nilai-nilai eigen
terisolasi dengan multiplisitas hingga.
\end{prop}

\begin{proof} Lihat \cite{HigsonR:2000}, proposisi 2.2.2. \ns \end{proof}

Berikut dua teorema dari analisis fungsional klasik.

\begin{thm}[Weyl]\label{005121} Jika $S$ dan $T$ merupakan operator pada suatu ruang Hilbert
 \index{teorema Weyl}%
yang selisihnya kompak, maka spektrum keduanya sama, kecuali mungkin pada nilai-nilai eigen.
\end{thm}

\begin{thm}[Weyl--von Neumann]\label{005124} Misalkan $T$ suatu operator swaadjoin pada
 \index{teorema Weyl--von Neumann}%
ruang Hilbert terpisahkan~$H$. Untuk setiap $\epsilon > 0$ terdapat operator $D$ yang dapat
didiagonalkan sedemikian sehingga $T - D$ kompak dan $\norm{T - D} < \epsilon$.
\end{thm}

\begin{proof} Lihat \cite{Conway:2000}, 38.1; \cite{Davidson:1996}, akibat II.4.2; atau
\cite{HigsonR:2000}, 2.2.5. \ns
\end{proof}

\begin{defn} Misalkan $H$ dan $K$ ruang-ruang Hilbert. Operator $S \in \ofml B(H)$ dan $T \in \ofml
B(K)$ disebut
 \index{esensial!ekuivalen uniter}%
 \index{uniter!ekuivalensi!esensial}%
 \index{ekuivalensi!uniter esensial}%
\df{ekuivalen uniter secara esensial} (atau
 \index{kompalen}%
\df{kompalen}) jika terdapat pemetaan uniter $U\colon H \sto K$ sedemikian sehingga $S - U^*TU$
merupakan operator kompak pada~$H$. (Kita memperluas definisi~\ref{000161} dan~\ref{0001612} dengan
cara yang jelas: $U \in \ofml B(H,K)$ disebut
 \index{uniter!pemetaan linear terbatas}%
 \index{terbatas!pemetaan linear!uniter}%
\df{uniter} jika $U^*U = I_H$ dan $UU^* = I_K$; dan $S$ serta $T$ disebut
 \index{uniter!ekuivalensi!operator}%
 \index{ekuivalensi!uniter}%
\df{ekuivalen secara uniter} jika terdapat pemetaan uniter $U\colon H \sto K$ sedemikian sehingga $S =
U^*TU$.)
\end{defn}

\begin{prop}\label{005134} Operator-operator swaadjoin $S$ dan $T$ pada ruang-ruang Hilbert terpisahkan berdimensi tak hingga
ekuivalen uniter secara esensial jika dan hanya jika keduanya memiliki spektrum esensial yang sama.
\end{prop}

\begin{proof} Lihat \cite{HigsonR:2000}, proposisi 2.2.4.  \ns \end{proof}

\begin{defn} Suatu operator ruang Hilbert $T$ disebut
 \index{esensial!normal}%
 \index{normal!secara esensial}
 \index{operator!normal secara esensial}%
\df{normal secara esensial} jika
 \index{komutator}%
 \index{<brackets@$[T,T^*]$ (komutator dari $T$)}%
\emph{komutator}-nya $[T,T^*] := TT^* - T^*T$ kompak. Dengan kata lain: $T$ normal secara
esensial jika citranya $\pi(T)$ merupakan elemen normal dari aljabar Calkin. Operator $T$
disebut
 \index{esensial!swaadjoin}%
 \index{swaadjoin!secara esensial}
 \index{operator!swaadjoin secara esensial}%
\df{swaadjoin secara esensial} jika $T - T^*$ kompak; yakni, jika $\pi(T)$
swaadjoin dalam aljabar Calkin.
\end{defn}

\begin{exam} Operator geser unilateral $S$ (lihat contoh~\ref{C063526})
 \index{geser unilateral!normalitas esensial}%
normal secara esensial (tetapi tidak normal).
\end{exam}





















\section{Operator Toeplitz}
\begin{defn} Misalkan
 \index{zeta@$\zeta$ (fungsi identitas pada lingkaran satuan)}%
$\zeta\colon \T \sto \T\colon z \mapsto z$ merupakan fungsi identitas pada lingkaran satuan.
Maka $\{\zeta^n\colon n \in \Z\}$ merupakan basis ortonormal bagi ruang Hilbert $L_2(\T)$ yang terdiri atas
(kelas-kelas ekuivalensi dari) fungsi-fungsi yang terintegralkan kuadrat pada $\T$ terhadap ukuran panjang
busur (yang dinormalisasi secara sesuai). Kita menyatakan dengan $H^2$ subruang dari $L_2(\T)$ yang merupakan
rentang linear tertutup dari $\{\zeta^n\colon n \ge 0\}$ dan dengan $P_+$ proyeksi
(ortogonal) pada $L_2(\T)$ yang jangkauan
(sekaligus kodomain)-nya adalah~$H^2$. Ruang $H^2$ merupakan
contoh suatu
 \index{ruang Hardy!$H^2$}%
 \index{ruang!Hardy}%
 \index{H@$H^2$}%
\emph{ruang Hardy}. Untuk setiap $\phi \in L_\infty(\T)$ kita mendefinisikan pemetaan $T_\phi$ dari
ruang Hilbert $H^2$ ke dirinya sendiri dengan $T_\phi = P_+M_\phi$ (dengan $M_\phi$ operator
perkalian yang didefinisikan dalam contoh~\ref{C063527}). Operator demikian disebut
 \index{Toeplitz!operator}%
 \index{operator!Toeplitz}%
 \index{T@$T_\phi$ (operator Toeplitz dengan simbol~$\phi$)}%
\df{operator Toeplitz} dan $\phi$ disebut
 \index{simbol!operator Toeplitz}%
 \index{Toeplitz!operator!simbol}%
\df{simbol}-nya. Jelas bahwa $T_\phi$ merupakan operator pada $H^2$ dan $\norm{T_\phi} \le \norm
{\phi}_\infty$.
\end{defn}

\begin{exam} Operator Toeplitz $T_\zeta$
 \index{geser unilateral!sebagai operator Toeplitz}%
 \index{Toeplitz!operator!geser unilateral sebagai}%
bertindak pada vektor-vektor basis $\zeta^n$ dari $H^2$ menurut $T_\zeta(\zeta^n) = \zeta^{n+1}$, sehingga operator ini
ekuivalen secara uniter dengan geser unilateral $S$.
\end{exam}

\begin{prop}\label{005216} Pemetaan $T\colon L_\infty(\T) \sto \ofml B(H^2)\colon \phi
\mapsto T_\phi$ bersifat positif, linear, dan melestarikan involusi.
\end{prop}

\begin{exam}\label{005217} Operator Toeplitz $T_\zeta$ dan $T_{\conj\zeta}$ menunjukkan bahwa
pemetaan $T$ dalam proposisi sebelumnya \emph{bukan} representasi dari $L_\infty(\T)$
pada~$\ofml B(H^2)$. (Bandingkan dengan contoh~\ref{002802}.)
\end{exam}

\begin{notn} Misalkan $H^\infty := H^2 \cap L_\infty(\T)$.
 \index{H@$H^\infty$}%
 \index{ruang Hardy!$H^\infty$}%
 \index{ruang!Hardy}%
Ini merupakan contoh lain dari suatu \emph{ruang Hardy}.
\end{notn}

\begin{prop} Suatu fungsi terbatas secara esensial $\phi$ pada lingkaran satuan termasuk dalam
$H^\infty$ jika dan hanya jika operator perkalian $M_\phi$ memetakan $H^2$ ke dalam~$H^2$.
\end{prop}

Meskipun (seperti telah kita lihat dalam contoh~\ref{005217}) pemetaan $T$ yang didefinisikan dalam
proposisi~\ref{005216} pada umumnya tidak multiplikatif, pemetaan ini multiplikatif untuk suatu
kelas fungsi yang luas.

\begin{prop} Jika $\phi \in L_\infty(\T)$ dan $\psi \in H^\infty$, maka $T_{\phi\psi} =
T_\phi T_\psi$ dan $T_{\conj\psi \phi} = T_{\conj\psi} T_\phi$.
\end{prop}

\begin{prop} Jika operator Toeplitz $T_\phi$ dengan simbol $\phi \in L_\infty(\T)$
invertibel, maka fungsi $\phi$ invertibel.
\end{prop}

\begin{proof} Lihat \cite{Douglas:1972}, proposisi 7.6. \ns \end{proof}

\begin{thm}[Inklusi Spektral Hartman--Wintner] Jika $\phi \in L_\infty(\T)$,
 \index{teorema inklusi spektral Hartman--Wintner}%
maka $\sigma(\phi) \subseteq \sigma(T_\phi)$ dan $\rho(T_\phi) = \norm{T_\phi} =
\norm{\phi}_\infty$.
\end{thm}

\begin{proof} Lihat \cite{Douglas:1972}, akibat 7.7 atau \cite{Murphy:1990}, teorema 3.5.7. \ns
\end{proof}

\begin{cor} Pemetaan $T\colon L_\infty(\T) \sto \ofml B(H^2)$ yang didefinisikan dalam
proposisi~\ref{005216} merupakan suatu isometri.
\end{cor}

\begin{prop} Suatu operator Toeplitz $T_\phi$ dengan simbol $\phi \in L_\infty(\T)$ kompak jika
dan hanya jika $\phi = \vc 0$.
\end{prop}

\begin{proof} Lihat \cite{Murphy:1990}, teorema 3.5.8.  \ns \end{proof}

\begin{prop} Jika $\phi \in \fml C(\T)$ dan $\psi \in L_\infty(\T)$, maka
semikomutator $T_\phi T_\psi - T_{\phi\psi}$ dan $T_\psi T_\phi - T_{\psi\phi}$ bersifat
kompak.
\end{prop}

\begin{proof} Lihat \cite{Arveson:2002}, proposisi 4.3.1; \cite{Davidson:1996}, akibat
V.1.4; atau \cite{Murphy:1990}, lema 3.5.9. \ns
\end{proof}

\begin{cor} Setiap operator Toeplitz dengan simbol kontinu bersifat normal secara esensial.
\end{cor}

\begin{defn} Andaikan $H$ ruang Hilbert terpisahkan dengan basis $\{e_0, e_1, e_2,
\dots\}$ dan $T$ suatu operator pada $H$ yang representasi matriks (tak hingga)-nya adalah
$[t_{ij}]$. Jika entri-entri dalam matriks ini hanya bergantung pada selisih indeks $i
- j$ (yakni, jika setiap diagonal yang sejajar dengan diagonal utama merupakan barisan konstan),
maka matriks tersebut disebut
 \index{Toeplitz!matriks}%
 \index{matriks!Toeplitz}%
\df{matriks Toeplitz}.
\end{defn}

\begin{prop} Misalkan $T$ suatu operator pada ruang Hilbert $H^2$. Representasi matriks
$T$ terhadap basis biasa $\{\zeta^n\colon n \ge 0\}$ merupakan matriks Toeplitz jika
dan hanya jika $S^*TS = T$ (dengan $S$ geser unilateral).
\end{prop}

\begin{proof} Lihat \cite{Arveson:2002}, proposisi 4.2.3.  \ns \end{proof}

\begin{prop} Jika $T_\phi$ suatu operator Toeplitz dengan simbol $\phi \in L_\infty(\T)$, maka
$S^*T_\phi S = T_\phi$. Sebaliknya, jika $R$ suatu operator pada ruang Hilbert $H^2$
sedemikian sehingga $S^*RS = R$, maka terdapat $\phi \in L_\infty(\T)$ tunggal sedemikian sehingga $R =
T_\phi$.
\end{prop}

\begin{proof} Lihat \cite{Arveson:2002}, teorema 4.2.4.  \ns \end{proof}

\begin{defn} 
 \index{Toeplitz!aljabar}%
 \index{aljabar!Toeplitz}%
 \index{T@$\ofml T$ (aljabar Toeplitz)}%
\df{aljabar Toeplitz} $\ofml T$ adalah subaljabar-$C^*$ dari $\ofml B(H^2)$ yang dibangkitkan oleh
operator geser unilateral; yaitu, $\ofml T = C^*(S)$.
\end{defn}

\begin{prop} Himpunan $\ofml K(H^2)$ dari operator-operator kompak pada $H^2$ merupakan ideal dalam aljabar
Toeplitz~$\ofml T$.
\end{prop}

\begin{prop} Aljabar Toeplitz terdiri atas semua perturbasi kompak dari operator Toeplitz
dengan simbol kontinu. Yaitu,
   \[ \ofml T = \{ T_\phi + K\colon \phi \in \fml C(\T) \text{ dan } K \in \ofml K(H^2)\} \,. \]
Lebih lanjut, jika $T_\phi + K = T_\psi + L$ dengan $\phi$, $\psi \in \fml C(\T)$ dan $K$, $L
\in \ofml K(H^2)$, maka $\phi = \psi$ dan $K = L$.
\end{prop}

\begin{proof} Lihat \cite{Arveson:2002}, teorema 4.3.2 atau \cite{Davidson:1996}, teorema V.1.5.
\ns \end{proof}

\begin{prop}\label{005251} Pemetaan $\pi \circ T\colon \fml C(\T) \sto \ofml Q(H^2)\colon \phi
\mapsto \pi(T_\phi)$ merupakan monomorfisme-$*\,$ beridentitas. Jadi pemetaan $\alpha\colon \fml C(\T)
\sto \ofml T/\ofml K(H^2)\colon \phi \mapsto \pi(T_\phi)$ menghasilkan isomorfisme
antara aljabar-$C^*$ $\fml C(\T)$ dan $\ofml T/\ofml K(H^2)$.
\end{prop}

\begin{proof} Lihat \cite{HigsonR:2000}, proposisi 2.3.3 atau \cite{Murphy:1990},
teorema 3.5.11.  \ns
\end{proof}

\begin{cor} Jika $\phi \in \fml C(\T)$, maka $\sigma_e(T_\phi) = \ran\phi$.
\end{cor}

\begin{prop}\label{005253} Barisan
   \[ \vc 0 \to \ofml K(H^2) \to \ofml T \to^\beta \fml C(\T) \to \vc 0 \]
bersifat eksak. Barisan ini tidak terbelah.
\end{prop}

\begin{proof} Pemetaan $\beta$ didefinisikan oleh $\beta(R) := \alpha^{-1}\bigl(\pi(R)\bigr)$ untuk
setiap $R \in \ofml T$ (dengan $\alpha$ isomorfisme yang didefinisikan dalam
proposisi sebelumnya~\ref{005251}). Lihat \cite{Arveson:2002}, catatan 4.3.3; \cite{Davidson:1996},
teorema V.1.5; atau \cite{HigsonR:2000}, halaman 35. \ns
\end{proof}

\begin{defn}\label{0052531} Barisan eksak pendek dalam proposisi sebelumnya~\ref{005253}
adalah
 \index{Toeplitz!ekstensi}%
 \index{ekstensi!Toeplitz}%
\df{ekstensi Toeplitz} dari $\fml C(\T)$ oleh~$\ofml K(H^2)$.
\end{defn}

\begin{rem} Suatu versi diagram ekstensi Toeplitz yang sering muncul kira-kira berbentuk
seperti
     \[ \vc 0 \to \ofml K(H^2) \to \ofml T \two/->`<-/^\beta_T \fml C(\T) \to \vc 0 \tag{1} \]
dengan $\beta$ seperti dalam~\ref{005253} dan $T$ adalah pemetaan $\phi \mapsto T_\phi$. Diagram ini
dapat disalahtafsirkan. Diagram tersebut mungkin memberi kesan kepada pembaca yang kurang waspada bahwa ini adalah ekstensi
terbelah, terutama karena memang benar bahwa $\beta \circ T = I_{\fml
C(\T)}$. Masalahnya, tentu saja, bahwa diagram ini bukan diagram dalam kategori
$\cat{CSA}$ dari aljabar-$C^*$ dan homomorfisme-$*\,$. Kita telah melihat dalam
contoh~\ref{005217} bahwa pemetaan $T$ bukan homomorfisme-$*\,$ karena tidak
selalu melestarikan perkalian. Elemen-elemen invertibel dalam $\fml C(\T)$ belum tentu terangkat menjadi
elemen-elemen invertibel dalam aljabar Toeplitz~$\ofml T$; jadi epimorfisme-$*\,$ $\beta$
tidak memiliki invers kanan dalam kategori~$\cat{CSA}$.

Sebagian penulis menangani masalah ini dengan mengatakan bahwa barisan (1) bersifat semiterbelah (lihat, misalnya,
Arveson~\cite{Arveson:2002}, halaman 112). Penulis lain meminjam istilah dari teori kategori,
yang menggunakan kata ``penampang'' untuk berarti ``memiliki invers kanan''.
Davidson~\cite{Davidson:1996}, misalnya, pada halaman 134 menyebut pemetaan $T$ sebagai
``penampang kontinu'', sedangkan Douglas~\cite{Douglas:1972}, pada halaman 179 menyebutnya
``penampang silang isometrik''. Meskipun memang benar bahwa $\beta$ memiliki $T$ sebagai invers
kanan dalam kategori $\cat{SET}$ dari himpunan dan pemetaan, bahkan juga dalam kategori ruang Banach
dan pemetaan linear terbatas, pemetaan tersebut tidak memiliki invers kanan dalam $\cat{CSA}$.
\end{rem}

\begin{prop} Jika $\phi$ suatu fungsi dalam $\fml C(\T)$, maka $T_\phi$ Fredholm jika dan hanya jika
fungsi itu tidak pernah bernilai nol.
\end{prop}

\begin{proof} Lihat \cite{Arveson:2002}, halaman 112, akibat 1; \cite{Davidson:1996}, teorema
V.1.6; \cite{Douglas:1972}, teorema 7.26; atau \cite{Murphy:1990}, akibat 3.5.12. \ns
\end{proof}

\begin{prop}\label{005271} Jika $\phi$ elemen invertibel dari $\fml C(\T)$, maka terdapat
bilangan bulat $n$ tunggal sedemikian sehingga $\phi = \zeta^n \exp \psi$ untuk suatu $\psi \in \fml
C(\T)$.
\end{prop}

\begin{proof} Lihat \cite{Arveson:2002}, proposisi 4.4.1 dan 4.4.2 atau \cite{Murphy:1990},
lema 3.5.14. \ns
\end{proof}

\begin{defn} Bilangan bulat yang keberadaannya dinyatakan dalam proposisi sebelumnya~\ref{005271}
adalah
 \index{bilangan lilit}%
 \index{w@$w(\phi)$ (bilangan lilit dari~$\phi$)}%
\df{bilangan lilit} dari fungsi invertibel $\phi \in \fml C(\T)$. Bilangan ini dinotasikan
dengan~$w(\phi)$.
\end{defn}

\begin{thm}[Teorema indeks Toeplitz]\label{005274}
 \index{Toeplitz!teorema indeks}%
Jika $T_\phi$ suatu operator Toeplitz dengan simbol kontinu yang tidak pernah bernilai nol, maka operator ini
merupakan operator Fredholm dan
   \[ \ind(T_\phi) = -w(\phi)\,. \]
\end{thm}

\begin{proof} Bukti elementer dapat ditemukan dalam \cite{Arveson:2002}, teorema 4.4.3, dan
\cite{Murphy:1990}. Bagi pembaca yang memiliki sedikit latar belakang tentang homotopi kurva,
proposisi~\ref{005271}, yang mengantar ke definisi \emph{bilangan lilit} di atas,
dapat dilewati. Merupakan fakta elementer dalam teori homotopi bahwa grup fundamental
$\pi_1(\C \setminus \{0\})$ dari bidang kompleks berlubang bersifat siklik tak hingga. Terdapat
isomorfisme $\tau$ dari $\pi_1(\C \setminus \{0\})$ ke $\Z$ yang mengaitkan bilangan bulat $+1$
dengan (kelas ekuivalensi yang memuat) fungsi $\zeta$. Hal ini memungkinkan kita untuk
mengaitkan setiap anggota invertibel $\phi$ dari $\fml C(\T)$ dengan bilangan bulat yang bersesuaian
di bawah $\tau$ dengan kelas ekuivalensinya. Bilangan bulat ini kita sebut
 \index{bilangan lilit}%
 \index{fundamental!grup}%
 \index{grup!fundamental}%
\emph{bilangan lilit} dari~$\phi$. Definisi demikian memungkinkan pemberian bukti yang lebih singkat
dan lebih anggun bagi \emph{teorema indeks Toeplitz}. Untuk pembahasan dengan pendekatan ini,
lihat \cite{Douglas:1972}, teorema 7.26, atau \cite{HigsonR:2000}, teorema 2.3.2. Untuk
latar belakang mengenai grup fundamental, lihat \cite{Massey:1967}, bab dua;
\cite{Wallace:1970}, lampiran A; atau \cite{Willard:1968}, bagian 32 dan 33. \ns
\end{proof}

\begin{thm}[Dekomposisi Wold]
 \index{dekomposisi Wold}%
 \index{dekomposisi!Wold}%
Setiap
 \index{sejati!isometri}%
 \index{isometri!sejati}%
isometri sejati (yakni, nonuniter) pada suatu ruang Hilbert merupakan jumlah langsung salinan-salinan
geser unilateral, atau merupakan jumlah langsung suatu operator uniter bersama salinan-salinan
geser tersebut.
\end{thm}

\begin{proof} Lihat \cite{Davidson:1996}, teorema V.2.1; \cite{Halmos:1982}, soal (dan
solusi) 149; atau \cite{Murphy:1990}, teorema 3.5.17. \ns
\end{proof}

\begin{thm}[Coburn]
 \index{teorema Coburn}%
Misalkan $v$ suatu isometri dalam aljabar-$C^*$ beridentitas~$A$.
Maka
terdapat
homomorfisme-$*\,$
beridentitas
$\tau$ tunggal dari aljabar Toeplitz $\ofml T$ ke $A$ sedemikian sehingga
$\tau(T_\zeta) = v$. Lebih lanjut, jika $vv^* \ne \vc 1$, maka $\tau$ merupakan suatu isometri.
\end{thm}

\begin{proof} Lihat \cite{Murphy:1990}, teorema 3.5.18, atau \cite{Davidson:1996},
teorema V.2.2. \ns
\end{proof}
%
%%
 
  
 











\section{Penjumlahan Ekstensi}

 \index{konvensi!ruang Hilbert separabel dan berdimensi tak hingga (mulai bagian Penjumlahan Ekstensi)}%
\begin{center}
 \fbox{\parbox{5.5in}{\textbf{Mulai sekarang semua ruang Hilbert akan separabel dan berdimensi
 tak hingga.}}}
\end{center}

\begin{prop} Suatu operator ruang Hilbert $T$ bersifat swaadjoin secara esensial jika dan hanya jika
operator itu merupakan perturbasi kompak dari suatu operator swaadjoin.
\end{prop}

\begin{prop}\label{005414} Dua operator ruang Hilbert $T_1$ dan $T_2$ yang swaadjoin secara esensial
ekuivalen uniter secara esensial jika dan hanya jika keduanya memiliki spektrum esensial yang
sama.
\end{prop}

Hasil yang analog tidak berlaku bagi operator normal secara esensial.

\begin{exam} Operator Toeplitz $T_\zeta$ dan $T_{\zeta^2}$ bersifat normal secara esensial dengan
spektrum esensial yang sama, tetapi keduanya tidak ekuivalen uniter secara esensial.
\end{exam}

\begin{defn} Sekarang kita membatasi perhatian pada suatu kelas khusus ekstensi. Jika $H$ merupakan
ruang Hilbert dan $A$ merupakan aljabar-$C^*$, kita mengatakan bahwa $(\ofml E,\phi)$ adalah suatu
 \index{ekstensi!$\ofml K(H)$ oleh $A$}%
\df{ekstensi $\ofml K = \ofml K(H)$ oleh $A$} jika $\ofml E$ merupakan subaljabar-$C^*$ beridentitas
dari $\ofml B(H)$ yang memuat $\ofml K$ dan $\phi$ merupakan homomorfisme-$*\,$ beridentitas sedemikian sehingga
barisan
 \begin{equation}\label{005431i}
    \vc 0 \to \ofml K \to^\iota \ofml E \to^\phi A \to \vc 0
 \end{equation}
(dengan $\iota$ pemetaan inklusi) eksak. Bahkan, kita hampir secara eksklusif akan
memperhatikan kasus $A = \fml C(X)$ untuk suatu ruang metrik kompak~$X$. Kita
mengatakan bahwa dua ekstensi $(\ofml E,\phi)$ dan $(\ofml E\,',\phi\,')$
 \index{ekuivalensi!ekstensi}%
 \index{ekstensi!ekuivalensi}%
\df{ekuivalen} jika terdapat suatu isomorfisme $\psi\colon \ofml E \sto \ofml E\,'$ yang
membuat diagram berikut komutatif.
 \begin{equation}\label{005431ii}
   \xymatrix{
     \vc 0\ar[r] & \ofml K\ar[r]^\iota\ar[d]_{\psi|_{\ofml K}} & \ofml E\ar[r]^\phi\ar[d]^\psi
                 & A\ar[r]\ar@{=}[d] & \vc 0 \\
     \vc 0\ar[r] & \ofml K\ar[r]^\iota & \ofml E\,'\ar[r]^{\phi\,'} & A\ar[r] & \vc 0
 }\end{equation}
Perhatikan bahwa definisi ini sedikit berbeda dari \emph{ekuivalensi kuat}
pada~\ref{0015151}. Keluarga semua kelas ekuivalensi ekstensi semacam ini kita lambangkan
 \index{ext@$\ext A$, $\ext X$!kelas ekuivalensi ekstensi}%
dengan~$\ext A$. Jika $A = \fml C(X)$, kita menulis $\ext X$, bukan~$\ext \fml C(X)$.
\end{defn}

\begin{defn} Jika $U\colon H_1 \sto H_2$ merupakan pemetaan uniter di antara ruang-ruang Hilbert, maka pemetaan
   \[ \ad_U\colon \ofml B(H_1) \sto \ofml B(H_2)\colon T \mapsto UTU^* \]
disebut
 \index{konjugasi}%
 \index{adU@$\ad_U$ (konjugasi oleh $U$)}%
\df{konjugasi} oleh~$U$.
\end{defn}

Jelas bahwa $\ad_U$ merupakan isomorfisme di antara aljabar-$C^*$ $\ofml B(H_1)$ dan
$\ofml B(H_2)$. Khususnya, jika $U$ merupakan operator uniter pada $H_1$, maka $\ad_U$ merupakan
automorfisme dari $\ofml K(H_1)$ maupun $\ofml B(H_1)$. Lebih lanjut, konjugasi merupakan
satu-satunya automorfisme aljabar-$C^*$~$\ofml K(H)$.

\begin{prop}\label{005436} Jika $H$ merupakan ruang Hilbert dan $\phi\colon \ofml K(H) \sto \ofml K(H)$
merupakan suatu automorfisme, maka $\phi = \ad_U$ untuk suatu operator uniter $U$ pada~$H$.
\end{prop}

\begin{proof} Lihat \cite{Davidson:1996}, lema V.6.1.  \ns \end{proof}

\begin{prop} Jika $H$ merupakan ruang Hilbert dan $A$ merupakan aljabar-$C^*$, maka ekstensi
$(\ofml E,\phi)$ dan $(\ofml E\,',\phi\,')$ dalam $\ext A$ ekuivalen jika dan hanya jika
terdapat suatu operator uniter $U$ dalam $\ofml B(H)$ sedemikian sehingga $\ofml E\,' = U \ofml E
U^*$ dan $\phi = \phi\,' \ad_U$.
\end{prop}

\begin{exam} Andaikan $T$ merupakan operator normal secara esensial pada suatu ruang Hilbert~$H$. Misalkan
 \index{E@$\ofml E_T$ (aljabar-$C^*$ yang dibangkitkan oleh $T$ dan $\ofml K(H)$)}%
$\ofml E_T$ merupakan aljabar-$C^*$ beridentitas yang dibangkitkan oleh $T$ dan~$\ofml K(H)$. Karena
$\pi(T)$ merupakan elemen normal aljabar Calkin, aljabar-$C^*$ beridentitas $\ofml
E_T/\ofml K(H)$ yang dibangkitkannya bersifat komutatif. Jadi \emph{teorema spektral
abstrak}~\ref{00144} memberi kita suatu isomorfisme aljabar-$C^*$ $\Psi\colon \fml C(\sigma_e(T))
\sto \ofml E_T/\ofml K(H)$. Misalkan $\sbsb{\phi}T = \Psi^{-1} \circ\pi\bigr|_{\ofml E_T}$.
Maka barisan
  \[ \vc 0 \to \ofml K(H) \to^\iota \ofml E_T \to^{\sbsb{\phi}T} \fml C(\sigma_e(T)) \to \vc 0 \]
eksak. Ini adalah
 \index{ekstensi!yang ditentukan oleh suatu operator normal secara esensial}%
\df{ekstensi $\ofml K(H)$ yang ditentukan oleh~$T$}.
\end{exam}

\begin{prop} Misalkan $T$ dan $T\,'$ operator normal secara esensial pada suatu ruang Hilbert~$H$. Kedua
operator itu ekuivalen uniter secara esensial jika dan hanya jika ekstensi yang
ditentukannya ekuivalen.
\end{prop}

\begin{prop} Jika $\ofml E$ merupakan aljabar-$C^*$ sedemikian sehingga $\ofml K(H) \subseteq \ofml E \subseteq
\ofml B(H)$ untuk suatu ruang Hilbert~$H$, $X$ merupakan subhimpunan kompak tak kosong dari $\C$, dan
$(\ofml E,\phi)$ merupakan ekstensi $\ofml K(H)$ oleh $\fml C(X)$, maka setiap elemen
dari $\ofml E$ bersifat normal secara esensial.
\end{prop}

\begin{defn} Jika $\phi_1\colon A_1 \sto B$ dan $\phi_2\colon A_2 \sto B$ merupakan
homomorfisme-$*\,$ beridentitas di antara aljabar-$C^*$, maka suatu
 \index{tarik balik}%
\df{tarik balik $A_1$ dan $A_2$ sepanjang $\phi_1$ dan $\phi_2$}, yang dilambangkan dengan
$(P,\pi_1,\pi_2)$, adalah aljabar-$C^*$ $P$ beserta sepasang
homomorfisme-$*\,$ beridentitas $\pi_1\colon P \sto A_1$ dan $\pi_2\colon P \sto A_2$ yang memenuhi
dua syarat berikut:
 \begin{enumerate}
    \item[(i)] diagram
        \[ \xy
          \Square[P`A_2`A_1`B;\pi_2`\pi_1`\phi_2`\phi_1]
      \endxy \]
bersifat komutatif; dan
    \item[(ii)] jika $\rho_1\colon Q \sto A_1$ dan $\rho_2\colon Q \sto A_2$ merupakan
homomorfisme-$*\,$ beridentitas dari aljabar-$C^*$ sedemikian sehingga diagram
        \[ \xy
          \Square[Q`A_2`A_1`B;\rho_2`\rho_1`\phi_2`\phi_1]
      \endxy \]
bersifat komutatif, maka terdapat suatu homomorfisme-$*\,$ beridentitas tunggal $\Psi\colon Q \sto P$
yang membuat diagram
       \[  \xy
         \pullback|brrb|<800,800>[P`A_2`A_1`B;\pi_2`\pi_1`\phi_2`\phi_1]%
         |amb|/{>}`{-->}`{>}/<500,500>[Q;\rho_2`\Psi`\rho_1]
       \endxy \]
komutatif.
 \end{enumerate}
\end{defn}

\begin{prop} Misalkan $H$ ruang Hilbert, $A$ aljabar-$C^*$ beridentitas, dan $\tau\colon
A \sto \ofml Q(H)$ monomorfisme-$*\,$ beridentitas. Maka, hingga isomorfisme,
terdapat tarik balik tunggal $(\ofml E,\pi_1,\pi_2)$ dari $A$ dan $\ofml B(H)$ sepanjang $\tau$ dan
$\pi$ sedemikian sehingga $(\ofml E,\pi_2)$ merupakan ekstensi $\ofml K(H)$ oleh $A$ yang membuat
diagram berikut komutatif.
 \begin{equation}\label{005464i}
   \xymatrix{
     \vc 0\ar[r] & \ofml K(H)\ar[r]^\iota\ar@{=}[d] & \ofml E\ar[r]^{\pi_2}\ar[d]^{\pi_1}
                 & A\ar[r]\ar[d]^\tau & \vc 0 \\
     \vc 0\ar[r] & \ofml K(H)\ar[r] & \ofml B(H)\ar[r]^{\pi} & \ofml Q(H)\ar[r] & \vc 0
 }\end{equation}
\end{prop}

\begin{proof}(Sketsa.) Misalkan $\ofml E = \{T \oplus a \in \ofml B(H) \oplus A\colon \tau(a) =
\pi(T)\}$, $\pi_1\colon \ofml E \sto \ofml B(H)\colon T \oplus a \mapsto T$, dan
$\pi_2\colon \ofml E \sto A\colon T \oplus a \mapsto a$. Ketunggalan (setiap
tarik balik) dibuktikan memakai ``omong kosong abstrak'' yang biasa.
\end{proof}

\begin{prop}\label{005466} Misalkan $H$ ruang Hilbert, $A$ aljabar-$C^*$ beridentitas, dan
$(\ofml E,\phi)$ suatu ekstensi $\ofml K(H)$ oleh~$A$. Maka terdapat suatu
monomorfisme-$*\,$ tunggal $\tau \colon A \sto \ofml Q(H)$ yang membuat
diagram~\eqref{005464i} komutatif.
\end{prop}

\begin{defn} Misalkan $H$ ruang Hilbert dan $A$ aljabar-$C^*$ beridentitas. Untuk $j=1,2$, misalkan
$\tau_j\colon A \sto \ofml Q(H)$ suatu monomorfisme-$*\,$ beridentitas. Kedua pemetaan ini
 \index{uniter!ekuivalensi!monomorfisme-$*\,$}%
 \index{ekuivalensi!uniter!monomorfisme-$*\,$}%
\df{ekuivalen secara uniter} jika terdapat suatu operator uniter $U$ pada $H$ sedemikian sehingga $\tau_2 =
\ad_U \tau_1$. Ekuivalensi uniter monomorfisme-$*\,$ beridentitas tentu merupakan suatu
relasi ekuivalensi. Kelas ekuivalensi yang memuat $\tau_1$ dilambangkan
dengan~$[\tau_1]$.
\end{defn}

\begin{prop} Misalkan $H$ ruang Hilbert dan $A$ aljabar-$C^*$ beridentitas. Dua ekstensi
$\ofml K(H)$ oleh $A$ ekuivalen jika dan hanya jika monomorfisme-$*\,$ beridentitas yang
bersesuaian dengannya (lihat~\ref{005466}) ekuivalen secara uniter.
\end{prop}

\begin{cor} Jika $H$ ruang Hilbert dan $A$ aljabar-$C^*$ beridentitas, terdapat suatu korespondensi
satu-satu di antara kelas-kelas ekuivalensi ekstensi $\ofml K(H)$ oleh $A$ dan
kelas-kelas ekuivalensi uniter monomorfisme-$*\,$ beridentitas dari $A$ ke aljabar
Calkin~$\ofml Q(H)$.
\end{cor}

\begin{conv} Berdasarkan akibat sebelumnya, kita akan memandang anggota $\ext A$ (atau
$\ext X$) sebagai
 \index{konvensi!E@$\ext A$ terdiri atas ekstensi atau morfisme}%
kelas ekuivalensi ekstensi atau kelas ekuivalensi uniter monomorfisme-$*\,$,
mana pun yang paling sesuai pada saat itu.
\end{conv}

\begin{prop}\label{0054702} Setiap ruang Hilbert $H$ (yang separabel dan berdimensi tak hingga)
isomorfik secara isometrik dengan $H \oplus H$. Jadi aljabar-$C^*$ $\ofml B(H)$ dan $\ofml
B(H \oplus H)$ isomorfik.
\end{prop}

\begin{defn}\label{0054715} Misalkan $\tau_1$, $\tau_2 \colon A \sto \ofml Q(H)$ merupakan
monomorfisme-$*\,$ beridentitas (dengan $A$ aljabar-$C^*$ dan $H$ ruang Hilbert). Kita mendefinisikan
suatu monomorfisme-$*\,$ beridentitas $\tau_1 \oplus \tau_2 \colon A \sto \ofml Q(H)$ melalui
  \[ (\tau_1 \oplus \tau_2)(a) := \rho\bigl(\tau_1(a) \oplus \tau_2(a)\bigr) \]
untuk setiap $a \in A$, dengan (seperti dalam Douglas\cite{Douglas:1980}) $\nu$ isomorfisme yang
ditetapkan dalam~\ref{0054702} dan $\rho$ pemetaan yang membuat diagram berikut
komutatif.
 \begin{equation}\label{0054715i}
  \xymatrix{
    \ofml B(H) \oplus \ofml B(H)\ar[r]\ar[dd]_{\pi \oplus \pi} & \ofml B(H \oplus H)\ar[r]^-\nu
                  & \ofml B(H)\ar[dd]^\pi   \\ & & & \\
    \ofml Q(H) \oplus \ofml Q(H)\ar[rr]_-\rho && \ofml Q(H)
 }\end{equation}
 Kemudian kita mendefinisikan operasi penjumlahan yang jelas pada $\ext A$:
   \[ [\tau_1] + [\tau_2] := [\tau_1 \oplus \tau_2] \]
untuk $[\tau_1]$, $[\tau_2] \in \ext A$.
\end{defn}

\begin{prop} Operasi penjumlahan (yang diberikan dalam~\ref{0054715}) pada $\ext A$ terdefinisi dengan baik
dan di bawah operasi ini $\ext A$ menjadi semigrup komutatif.
\end{prop}

\begin{defn}\label{0054811} Misalkan $A$ aljabar-$C^*$ beridentitas dan $\mathbf r\colon A \sto
\ofml B(H)$ representasi tak terdegenerasi dari $A$ pada suatu ruang Hilbert~$H$. Misalkan $P$
proyeksi ruang Hilbert dan $M$ jangkauan~$P$. Andaikan $P\mathbf r(a) -
\mathbf r(a)P \in \ofml K(H)$ untuk setiap $a~\in~A$. Lambangkan dengan $P_M$ proyeksi $P$
yang himpunan kodomainnya sama dengan jangkauannya; yaitu, $P_M\colon H \sto M \colon x \mapsto
Px$. Maka untuk setiap $a \in A$ kita mendefinisikan
 \index{operator!Toeplitz!abstrak}%
 \index{Toeplitz!operator!abstrak}%
 \index{abstrak!operator!Toeplitz}%
\df{operator Toeplitz abstrak} $T_a \in \ofml B(M)$
 \index{simbol!operator Toeplitz abstrak}%
 \index{T@$T_a$ (operator Toeplitz abstrak)}%
\df{bersimbol $a$ yang dikaitkan dengan} pasangan~$(\mathbf r,P)$ melalui $T_a = P_M\mathbf
r(a)\bigr|_M$.
\end{defn}

\begin{defn} Gunakan notasi seperti dalam definisi sebelumnya~\ref{0054811}. Kemudian kita mendefinisikan
 \index{ekstensi!Toeplitz!abstrak}%
 \index{Toeplitz!ekstensi!abstrak}%
 \index{abstrak!ekstensi!Toeplitz}%
\df{ekstensi Toeplitz abstrak $\tau_P$ yang dikaitkan dengan pasangan $(\mathbf r,P)$} melalui
   \[ \tau_P\colon A \sto \ofml Q(M)\colon a \mapsto \pi(T_a)\,. \]
\end{defn}

Perhatikan bahwa (monomorfisme-$*\,$ beridentitas yang dikaitkan dengan) ekstensi Toeplitz konkret
yang didefinisikan dalam~\ref{0052531} merupakan contoh ekstensi Toeplitz abstrak, dan
juga bahwa secara umum ekstensi Toeplitz abstrak tidak harus injektif.

\begin{defn} Misalkan $A$ aljabar-$C^*$ beridentitas dan $H$ ruang Hilbert. Dalam semangat
\cite{HigsonR:2000}, definisi 2.7.6, kita mengatakan bahwa suatu monomorfisme-$*\,$ beridentitas
$\tau\colon A \sto \ofml Q(H)$ bersifat
 \index{semiterbelah}%
\df{semiterbelah} jika terdapat suatu monomorfisme-$*\,$ beridentitas $\tau\,'\colon A \sto \ofml
Q(H)$ sedemikian sehingga $\tau \oplus \tau'$ terbelah.
\end{defn}

\begin{prop} Andaikan $A$ aljabar-$C^*$ beridentitas dan $H$ ruang Hilbert. Maka suatu
monomorfisme-$*\,$ beridentitas $\tau\colon A \sto \ofml Q(H)$ semiterbelah jika dan hanya jika
ekuivalen secara uniter dengan suatu ekstensi Toeplitz abstrak.
\end{prop}

\begin{proof} Lihat \cite{HigsonR:2000}, proposisi 2.7.10.  \ns \end{proof}

















\vskip .2 in


%\section{Hasil Kali Tensor Aljabar-$C^*$}


\vskip .2 in


















\section{Pemetaan Positif Lengkap}
Dalam contoh~\ref{001306} ditunjukkan bagaimana himpunan $\mathbf M_n$ semua matriks $n \times n$
atas bilangan kompleks dapat dijadikan aljabar-$C^*$. Sekarang kita memperumum contoh itu ke
aljabar $\M nA$ semua matriks $n \times n$ yang unsur-unsurnya berasal dari aljabar-$C^*$~$A$.

\begin{exam} Dalam contoh~\ref{000533a} dinyatakan bahwa jika $A$ merupakan aljabar-$C^*$, maka di bawah
operasi aljabar biasa, himpunan $\M nA$ semua matriks $n \times n$ yang unsur-unsurnya berasal dari
$A$ merupakan suatu aljabar. Aljabar ini dapat dijadikan aljabar-$*\,$ dengan mengambil transposisi
konjugat sebagai involusi. Artinya, definisikan
   \[ [ a_{ij}]^* := \bigl[ {a_{j\,i}}^*\bigr] \]
dengan $a_{ij} \in A$ untuk $1 \le i,j \le n$.

Misalkan $H$ ruang Hilbert. Untuk sementara, ruang itu tidak harus berdimensi tak hingga atau
separabel. Lambangkan dengan $H^n$ jumlah langsung rangkap-$n$-nya. Artinya,
   \[ H^n := \bigoplus_{k=1}^n H_k \]
dengan $H_k = H$ untuk $k = 1, \dots, n$. Misalkan $[T_{jk}] \in \mathbf M_n(\ofml B(H))$. Untuk
$\vc x = x_1 \oplus \dots \oplus x_n \in H^n$, definisikan
   \[ \vc T \colon H^n \sto H^n \colon \vc x \mapsto [T_{jk}] \vc x =
            \sum_{k=1}^n T_{1k}x_k \oplus \dots \oplus \sum_{k=1}^n T_{nk}x_k \,. \]
Maka $\vc T \in \ofml B(H^n)$. Lebih lanjut, pemetaan
   \[ \Psi \colon \mathbf M_n\bigl(\ofml B(H)\bigr) \sto \ofml B(H^n) \colon
             [T_{jk}] \mapsto \vc T \]
merupakan isomorfisme aljabar-$*\,$. Gunakan $\Psi$ untuk memindahkan norma operator pada $\ofml
B(H^n)$ ke~$\mathbf M_n\bigl(\ofml B(H)\bigr)$. Artinya, definisikan
   \[ \bignorm{\,[T_{jk}]\,} := \norm{\vc T}\,. \]
Ini menjadikan $\mathbf M_n\bigl(\ofml B(H)\bigr)$ aljabar-$C^*$ yang isomorfik dengan~$\ofml
B(H^n)$.

Sekarang andaikan $A$ aljabar-$C^*$ sebarang. Gunakan versi III Teorema
Gelfand--Naimark pada~\ref{002973}. Dengan teorema ini, kita dapat mengidentifikasi $A$ dengan
subaljabar-$C^*$ dari $\ofml B(H)$ untuk suatu ruang Hilbert $H$ dan, akibatnya,
 \index{M@$\M nA$!sebagai aljabar-$C^*$}%
 \index{C*@aljabar-$C^*$!M@$\M nA$ sebagai suatu}%
$\M nA$ dengan suatu subaljabar-$C^*$ dari $\mathbf M_n\bigl(\ofml B(H)\bigr)$. Dengan
identifikasi ini, $\M nA$ menjadi aljabar-$C^*$. (Perhatikan bahwa menurut akibat~\ref{00131104},
norma pada aljabar-$C^*$ bersifat tunggal; jadi jelas bahwa norma pada $\mathbf M_n(A)$
tidak bergantung pada cara khusus $A$ direpresentasikan sebagai subaljabar-$C^*$ dari
operator pada suatu ruang Hilbert.)
\end{exam}

\begin{exam} Misalkan $k \in \N$ dan $A = \mathbf M_k$. Maka untuk setiap $n \in \N$,
   \[ \mathbf M_n(A) = \mathbf M_n(\mathbf M_k) \cong \mathbf M_{nk}\,. \]
Isomorfismenya jelas: cukup hapus kurung siku bagian dalam. Sebagai contoh,
isomorfisme dari $\mathbf M_2(\mathbf M_2)$ ke $\mathbf M_4$ diberikan oleh
  \[ \begin{bmatrix} {} & {} \\
          \begin{bmatrix}
                          a_{11} & a_{12} \\
                          a_{21} & a_{22}
          \end{bmatrix}       &
          \begin{bmatrix} b_{11} & b_{12} \\
                          b_{21} & b_{22}
          \end{bmatrix}       \\*[25pt]
          \begin{bmatrix} c_{11} & c_{12} \\
                          c_{21} & c_{22}
          \end{bmatrix}       &
          \begin{bmatrix} d_{11} & d_{12} \\
                          d_{21} & d_{22}
          \end{bmatrix}       \\
                    {}  & {}
     \end{bmatrix}
    \quad \mapsto \quad
          \begin{bmatrix} a_{11} & a_{12} & b_{11} & b_{12} \\
                          a_{21} & a_{22} & b_{21} & b_{22} \\
                          c_{11} & c_{12} & d_{11} & d_{12} \\
                          c_{21} & c_{22} & d_{21} & d_{22}
          \end{bmatrix}
  \]
\end{exam}

\begin{defn} Suatu pemetaan $\phi\colon A \sto B$ di antara aljabar-$C^*$ disebut \df{positif} jika
pemetaan itu membawa elemen positif ke elemen positif; yakni, jika $\phi(a) \in B^+$ setiap kali
$a \in A^+$.
\end{defn}

\begin{exam} Setiap homomorfisme-$*\,$ di antara aljabar-$C^*$ bersifat positif. (Lihat
proposisi~\ref{0018201}.)
\end{exam}

\begin{prop} Setiap pemetaan linear positif di antara aljabar-$C^*$ bersifat terbatas.
\end{prop}

\begin{notn} Misalkan $A$ dan $B$ aljabar-$C^*$ serta $n \in \N$. Suatu pemetaan linear $\phi\colon A \sto
B$ menginduksi pemetaan linear $\phi^{(\,n)} := \phi \otimes \id{\mathbf M_n}\colon A \otimes
\mathbf M_n \sto B \otimes \mathbf M_n$. Seperti telah kita lihat, $A \otimes \mathbf M_n$ dapat
diidentifikasi dengan $\mathbf M_n(A)$. Jadi kita dapat memandang $\phi^{(\,n)}$ sebagai pemetaan dari
$\mathbf M_n(A)$ ke $\mathbf M_n(B)$ yang didefinisikan oleh $\phi^{(\,n)}\big(\,[a_{jk}]\,\bigr) =
\bigl[ \phi(a_{jk})\bigr]$.
\end{notn}

Mudah dilihat bahwa jika $\phi$ mempertahankan perkalian, maka setiap $\phi^{(\,n)}$ juga demikian,
dan jika $\phi$ mempertahankan involusi, maka setiap $\phi^{(\,n)}$ juga demikian.
Akan tetapi, kepositifan tidak selalu dipertahankan, sebagaimana ditunjukkan contoh~\ref{005824} berikut.

\begin{defn} Dalam $\mathbf M_n$,
 \index{standar!unit matriks}%
 \index{matriks!unit}%
 \index{unit!matriks}%
\df{unit matriks standar} adalah matriks-matriks $e^{jk}$ ($1 \le j,k \le n$) yang entri pada
baris ke-$j$ dan kolom ke-$k$ bernilai $1$, sedangkan semua entri lainnya
bernilai~$0$.
\end{defn}

\begin{exam}\label{005824} Misalkan $A = \mathbf M_2$. Maka $\phi\colon A \sto A\colon a \mapsto a^t$
(dengan $a^t$ transposisi~$a$) merupakan pemetaan positif. Pemetaan $\phi^{(2)}\colon
\mathbf M_2(A) \sto \mathbf M_2(A)$ tidak positif. Untuk melihatnya, misalkan $e^{11}$, $e^{12}$,
$e^{21}$, dan $e^{22}$ unit matriks standar untuk~$A$. Maka
  \[ \phi^{(2)}\left(\begin{bmatrix} e^{11} & e^{12} \\ e^{21} & e^{22} \end{bmatrix}\right)
    = \begin{bmatrix} \phi(e^{11}) & \phi(e^{12}) \\ \phi(e^{21}) & \phi(e^{22}) \end{bmatrix}
    = \begin{bmatrix} e^{11} & e^{21} \\ e^{12} & e^{22} \end{bmatrix}
    = \begin{bmatrix} 1 & 0 & 0 & 0 \\
                      0 & 0 & 1 & 0 \\
                      0 & 1 & 0 & 0 \\
                      0 & 0 & 0 & 1
      \end{bmatrix} \]
yang tidak positif.
\end{exam}

\begin{defn} Suatu pemetaan linear $\phi\colon A \sto B$ di antara aljabar-$C^*$ beridentitas disebut
 \index{npositif@$n$-positif}%
 \index{positif!$n$-}%
\df{$n$-positif} (untuk $n \in \N$) jika $\phi^{(\,n)}$ positif. Pemetaan itu disebut
 \index{positif!lengkap}%
 \index{lengkap!positif}%
\df{positif lengkap} jika bersifat $n$-positif untuk setiap $n \in \N$.
\end{defn}

\begin{exam} Setiap homomorfisme-$*\,$ di antara aljabar-$C^*$ bersifat positif lengkap.
\end{exam}

\begin{prop}\label{00582811} Misalkan $A$ aljabar-$C^*$ beridentitas dan $a \in A$. Maka $\norm a \le 1$ jika dan hanya jika
$\begin{bmatrix} \vc 1 & a \\ a^* & \vc 1 \end{bmatrix}$ positif dalam~$\mathbf M_2(A)$.
\end{prop}

\begin{prop} Misalkan $A$ aljabar-$C^*$ beridentitas dan $a$, $b \in A$. Maka $a^*a \le b$ jika
dan hanya jika $\begin{bmatrix} \vc 1 & a \\ a^* & b \end{bmatrix}$ positif dalam~$\mathbf
M_2(A)$.
\end{prop}

\begin{prop} Setiap pemetaan $2$-positif beridentitas di antara aljabar-$C^*$ beridentitas bersifat kontraktif.
\end{prop}

\begin{prop}[Ketaksamaan Kadison]\label{0058283} Jika $\phi\colon A \sto B$ merupakan pemetaan
$2$-positif beridentitas
 \index{ketaksamaan Kadison}%
 \index{ketaksamaan!Kadison}%
di antara aljabar-$C^*$ beridentitas, maka
   \[ \bigl(\phi(a)\bigr)^*\phi(a) \le \phi(a^*a) \]
untuk setiap $a \in A$.
\end{prop}

\begin{defn} Mudah dilihat bahwa jika $\phi\colon A \sto B$ merupakan pemetaan linear terbatas
di antara aljabar-$C^*$, maka $\phi^{(\,n)}$ terbatas untuk setiap~$n \in \N$. Jika
$\norm{\phi}_{\mathrm{cb}} := \sup\{\norm{\phi^{(\,n)}} \colon n \in \N\} < \infty$, maka
$\phi$ disebut
 \index{lengkap!terbatas}%
 \index{terbatas!lengkap}%
 \index{<unopn@$\norm{\hphantom{\phi}}_{\mathrm{cb}}$ (norma pemetaan terbatas lengkap)}%
\df{terbatas lengkap}.
\end{defn}

\begin{prop}\label{00582921} Jika $\phi\colon A \sto B$ merupakan pemetaan linear beridentitas dan positif lengkap
di antara aljabar-$C^*$ beridentitas, maka $\phi$ terbatas lengkap dan
   \[ \norm{\phi(\vc 1)} = \norm\phi = \norm\phi_{\mathrm{cb}}\,. \]
\end{prop}

\begin{proof} Lihat \cite{Paulsen:2002}. \ns \end{proof}

\begin{thm}[Teorema dilatasi Stinespring]\label{005831} Misalkan $A$ aljabar-$C^*$ beridentitas,
 \index{Stinespring!teorema dilatasi}%
 \index{dilatasi!teorema Stinespring}%
$H$ ruang Hilbert, dan $\phi\colon A \sto \ofml B(H)$ pemetaan linear beridentitas.
Maka $\phi$ positif lengkap jika dan hanya jika terdapat suatu ruang Hilbert
$H_0$, suatu isometri $V\colon H \sto H_0$, dan suatu representasi tak terdegenerasi $\vc
r\colon A \sto \ofml B(H_0)$ sedemikian sehingga $\phi(a) = V^*\vc r(a)V$ untuk setiap $a \in A$.
\end{thm}

\begin{proof} Lihat \cite{HigsonR:2000}, teorema 3.1.3, atau \cite{Paulsen:2002}, teorema 4.1. \ns
\end{proof}

\begin{notn} Misalkan $A$ aljabar-$C^*$, $f\colon A \sto \ofml B(H_0)$, dan
$g\colon A \sto \ofml B(H)$, dengan $H_0$ dan $H$ ruang Hilbert. Kita menulis
 \index{<relation@$g \lesssim f$ (relasi di antara pemetaan bernilai operator)}%
$g \lesssim f$ jika terdapat suatu isometri $V\colon H \sto H_0$ sedemikian sehingga $g(a) -
V^*f(a)V \in \ofml K(H)$ untuk setiap $a \in A$. Relasi $\lesssim$ merupakan praurutan,
tetapi bukan urutan parsial; yakni, relasi itu refleksif dan transitif, tetapi tidak
antisimetris.
\end{notn}

\begin{thm}[Voiculescu]\label{0058335} Misalkan $A$ aljabar-$C^*$ beridentitas yang separabel,
 \index{teorema Voiculescu}%
$\vc r\colon A \sto \ofml B(H_0)$ representasi tak terdegenerasi, dan
$\phi\colon A \sto \ofml B(H)$ pemetaan linear beridentitas dan positif lengkap (dengan $H$ dan $H_0$ ruang Hilbert). Jika $\phi(a) =
\vc 0$ setiap kali $\vc r(a) \in \ofml K(H_0)$, maka $\phi \lesssim \vc r$.
\end{thm}

\begin{proof} Buktinya panjang dan rumit, tetapi ``elementer''. Lihat \cite{HigsonR:2000},
bagian 3.5--3.6. \ns
\end{proof}

\begin{prop}\label{0058338} Misalkan $A$ aljabar-$C^*$ beridentitas yang separabel, $H$ ruang Hilbert,
dan $\tau_1$, $\tau_2 \colon A \sto \ofml Q(H)$ monomorfisme-$*\,$ beridentitas. Jika
$\tau_2$ terbelah, maka $\tau_1 + \tau_2$ ekuivalen secara uniter dengan~$\tau_1$.
\end{prop}

\begin{proof} Lihat \cite{HigsonR:2000}, teorema 3.4.7. \ns \end{proof}

\begin{cor}\label{0058339} Andaikan $A$ aljabar-$C^*$ beridentitas yang separabel, $H$ ruang
Hilbert, dan $\tau \colon A \sto \ofml Q(H)$ monomorfisme-$*\,$ beridentitas yang terbelah.
Maka $[\tau]$ merupakan identitas aditif dalam~$\ext A$.
\end{cor}

Perhatikan bahwa suatu monomorfisme-$*\,$ beridentitas $\tau\colon A \sto \ofml Q(H)$ dari aljabar-$C^*$
beridentitas yang separabel ke aljabar Calkin dari suatu ruang Hilbert $H$ bersifat semiterbelah jika dan
hanya jika $[\tau]$ mempunyai invers aditif dalam $\ext A$. Proposisi berikut mengatakan bahwa
hal ini terjadi jika dan hanya jika $\tau$ mempunyai suatu pengangkatan positif lengkap ke~$\ofml B(H)$.

\begin{prop}\label{0058411} Misalkan $A$ aljabar-$C^*$ beridentitas yang separabel dan $H$ ruang Hilbert.
Suatu monomorfisme-$*\,$ beridentitas $\tau\colon A \sto \ofml Q(H)$ semiterbelah jika dan hanya jika
terdapat suatu pemetaan linear beridentitas dan positif lengkap $\wt\tau\colon A \sto \ofml B(H)$ sedemikian sehingga
diagram
  \[
    \xy
      \dtriangle/<-`->`->/[\ofml B(H)`A`\ofml Q(H);\wt\tau`\pi`\tau]
    \endxy
   \]
komutatif.
\end{prop}

\begin{proof} Lihat \cite{HigsonR:2000}, teorema 3.1.5.  \ns \end{proof}

\begin{defn} Suatu aljabar-$C^*$ $A$ disebut
 \index{nuklir}%
 \index{C*@aljabar-$C^*$!nuklir}%
\df{nuklir} jika untuk setiap aljabar-$C^*$ $B$ terdapat suatu norma-$C^*$ tunggal pada hasil kali
tensor aljabaris $A \odot B$.
\end{defn}

\begin{exam}\label{0058513} Setiap aljabar-$C^*$ berdimensi hingga bersifat nuklir.
 \index{nuklir!aljabar berdimensi hingga bersifat}%
\end{exam}

\begin{proof} Lihat \cite{Murphy:1990}, teorema 6.3.9.  \ns \end{proof}

\begin{exam}\label{0058514} Aljabar-$C^*$ $\mathbf M_n$
 \index{M@$\mathbf M_n$!sebagai aljabar-$C^*$ nuklir}%
 \index{nuklir!$\mathbf M_n$ bersifat}%
bersifat nuklir.
\end{exam}
\begin{proof} Lihat \cite{Blackadar:2006}, II.9.4.2.  \ns \end{proof}

\begin{exam}\label{0058515} Jika $H$ ruang Hilbert,
 \index{K@$\ofml K(H)$!sebagai aljabar-$C^*$ nuklir}%
 \index{nuklir!$\ofml K(H)$ bersifat}%
maka aljabar-$C^*$ $\ofml K(H)$ dari operator kompak bersifat nuklir.
\end{exam}
\begin{proof} Lihat \cite{Murphy:1990}, contoh 6.3.2.  \ns \end{proof}

\begin{exam}\label{0058517} Jika $X$ ruang Hausdorff kompak, maka aljabar-$C^*$
 \index{kontinu@$\fml C(X)$!sebagai aljabar-$C^*$ nuklir}%
 \index{nuklir!$\fml C(X)$ bersifat}%
$\fml C(X)$ bersifat nuklir.
\end{exam}

\begin{proof} Lihat \cite{Paulsen:2002}, proposisi 12.9.  \ns \end{proof}

\begin{exam}\label{0058519} Setiap aljabar-$C^*$ komutatif bersifat nuklir.
 \index{nuklir!aljabar-$C^*$ komutatif bersifat}%
\end{exam}

\begin{proof} Lihat \cite{Fillmore:1996}, teorema 7.4.1, atau \cite{Wegge-Olsen:1993}, teorema
T.6.20. \ns
\end{proof}

Kesimpulan hasil berikut dikenal sebagai
 \index{lengkap!positif!sifat pengangkatan}%
\emph{sifat pengangkatan positif lengkap}.

\begin{prop}\label{0058541} Misalkan $A$ aljabar-$C^*$ beridentitas, separabel, dan nuklir, serta $J$ suatu
ideal separabel dalam aljabar-$C^*$ beridentitas~$B$. Maka untuk setiap pemetaan positif lengkap
$\phi\colon A \sto B/J$, terdapat suatu pemetaan positif lengkap $\wt\phi\colon A \sto B$
yang membuat diagram
  \[
    \xy
      \dtriangle/<-`->`->/[B`A`B/J;\wt\phi``\phi]
    \endxy
   \]
komutatif.
\end{prop}

\begin{proof} Lihat \cite{HigsonR:2000}, teorema 3.3.6.  \ns \end{proof}

\begin{cor}\label{0058543} Jika $A$ merupakan aljabar-$C^*$ beridentitas, separabel, dan nuklir,
 \index{ext@$\ext A$, $\ext X$!merupakan grup Abelian}%
maka $\ext A$ merupakan grup Abelian.
\end{cor}







\endinput

 \chapter{FUNKTOR $\mathbf{\emph{K}_0}$}

Sekarang kita meninjau secara singkat apa yang disebut teori-$K$ untuk aljabar-$C^*$. Barangkali pengantar terbaik untuk pokok bahasan ini adalah~\cite{RordamLL:2000},
yang beberapa bagian elementernya diikuti dengan cukup saksama dalam catatan ini. Pengantar lain yang sangat mudah dibaca adalah~\cite{Wegge-Olsen:1993}.

Diberikan suatu aljabar-$C^*$ $A$, kita terutama akan memperhatikan dua grup yang dikenal sebagai $K_0(A)$ dan~$K_1(A)$. Bab ini membahas
grup yang pertama, yakni grup proyeksi di $A$ dengan operasi penjumlahan. Tentu saja, jika kita ingin menjumlahkan proyeksi sembarang, kita langsung
menghadapi kesulitan serius. Sebelumnya telah kita tunjukkan (dalam proposisi~\ref{prop_sum_projs}) bahwa agar proyeksi dapat dijumlahkan, proyeksi-proyeksi itu harus komutatif
(dan hasil kalinya nol)! Penyelesaian dilema ini mencerminkan cara berpikir yang lazim dalam teori-$K$ secara umum. Jika dua proyeksi tidak
komutatif, singkirkan penghalang yang mencegah keduanya menjadi komutatif. Jika tidak cukup ruang bagi keduanya untuk saling melewati, berilah keduanya lebih banyak
ruang. Jangan bersikeras memandang keduanya sebagai makhluk yang mencoba hidup dalam dunia matriks $1 \times 1$ berentri di~$A$ yang sesak tanpa harapan.
Biarkan keduanya, misalnya, menjelajahi dunia matriks $2 \times 2$ yang jauh lebih lapang. % SOURCE-CORRECTION: CH17-C001
Untuk mewujudkan gagasan ini secara teknis, kita mencari relasi ekuivalensi $\sim$ yang mengidentifikasi proyeksi $p$ di $A$ dengan matriks $\begin{bmatrix} p & \vc 0 \\ \vc 0 & \vc 0 \end{bmatrix}$
dan $\begin{bmatrix} \vc 0 & \vc 0 \\ \vc 0 & p \end{bmatrix}$. Sebagai motivasi heuristik saja---bukan sebagai konsekuensi suatu kongruensi perkalian yang telah dibuktikan---hambatan itu tampak hilang: jika $p$ dan $q$ proyeksi di~$A$, maka % SOURCE-CORRECTION: CH17-C002
    \[ pq \sim \begin{bmatrix} p & \vc 0 \\ \vc 0 & \vc 0 \end{bmatrix} \begin{bmatrix} \vc 0 & \vc 0 \\ \vc 0 & q \end{bmatrix}  \\
                   = \begin{bmatrix} \vc 0 & \vc 0 \\ \vc 0 & \vc 0 \end{bmatrix}           \\
                   =  \begin{bmatrix} \vc 0 & \vc 0 \\ \vc 0 &  q \end{bmatrix}\begin{bmatrix} p & \vc 0 \\ \vc 0 & \vc 0 \end{bmatrix}  \\
                   \sim qp\,. \]
Dalam gambaran heuristik ini, $p$ dan $q$ komutatif modulo relasi ekuivalensi $\sim$, dan kita dapat mendefinisikan penjumlahan dengan sesuatu seperti
$p \oplus q = \begin{bmatrix} p & \vc 0 \\ \vc 0 & q \end{bmatrix}$.

Cukup jelas bahwa perangkat sederhana di atas tidak akan langsung menghasilkan sesuatu yang menyerupai grup Abelian, tetapi ini merupakan suatu permulaan. Jika Anda menduga
bahwa pada akhirnya konstruksi $K_0(A)$ mungkin agak rumit, dugaan Anda sepenuhnya benar.




\section{Relasi Ekuivalensi pada Proyeksi}

Kita mendefinisikan beberapa relasi ekuivalensi, yang semuanya sesuai untuk proyeksi dalam aljabar-$C^*$ beridentitas.

\begin{defn} Dalam suatu aljabar beridentitas, elemen $a$ dan $b$ disebut
 \index{serupa!elemen suatu aljabar}%
\df{serupa} jika terdapat elemen invertibel $s$ sedemikian sehingga $b = sas^{-1}$.
 \index{<binrelequivs@$a \sim_s b$ (keserupaan elemen suatu aljabar)}%
Dalam hal ini kita menulis $a \sim_s b$.
\end{defn}

\begin{prop} Relasi keserupaan $\sim_s$ yang didefinisikan di atas merupakan relasi ekuivalensi pada elemen-elemen suatu aljabar beridentitas.
\end{prop}

\begin{defn} Dalam suatu aljabar-$C^*$, elemen $a$ dan $b$ disebut
 \index{uniter!ekuivalensi}%
\df{ekuivalen secara uniter} jika terdapat suatu elemen uniter $u \in \wt A$ sedemikian sehingga $b = uau^*$. % SOURCE-CORRECTION: CH17-C003
 \index{<binrelequivu@$a \sim_u b$ (ekuivalensi uniter)}%
Dalam hal ini kita menulis $a \sim_u b$.
\end{defn}

\begin{prop} Relasi ekuivalensi uniter $\sim_u$ yang didefinisikan di atas merupakan relasi ekuivalensi pada elemen-elemen suatu aljabar-$C^*$.
\end{prop}

\begin{prop}\label{0060224fa} Elemen $a$ dan $b$ dari suatu aljabar-$C^*$ beridentitas $A$ ekuivalen secara uniter jika dan hanya jika
terdapat elemen $u \in \ofml U(A)$ sedemikian sehingga $b = uau^*$.
\end{prop}

\begin{proof}[\emph{Petunjuk untuk bukti}] Andaikan terdapat $v \in \ofml U(\wt A)$ sedemikian sehingga $b = vav^*$. Berdasarkan proposisi~\ref{0015213},
terdapat $u \in A$ dan $\alpha \in \C$ sedemikian sehingga $v = u + \alpha \vc j$ (dengan $\vc j = \vc 1_{\wt A} - \vc 1_A$). Tunjukkan
bahwa $\abs\alpha = 1$, $u$ uniter, dan $b = uau^*$.

Untuk arah sebaliknya, andaikan terdapat $u \in \ofml U(A)$ sedemikian sehingga $b = uau^*$. Misalkan $v = u + \vc j$. \ns
\end{proof}

\begin{prop}[Dekomposisi polar]\label{0060141} Untuk setiap elemen invertibel $s$ dari suatu
 \index{dekomposisi polar!elemen invertibel aljabar-$C^*$}%
 \index{dekomposisi!polar}%
aljabar-$C^*$ beridentitas, terdapat suatu elemen uniter $\omega(s)$ dalam aljabar tersebut sedemikian sehingga % SOURCE-CORRECTION: CH17-C004
   \[   s = \omega(s) \abs s\,. \]
\end{prop}

\begin{proof}[\emph{Petunjuk untuk bukti}] Gunakan Akibat~\ref{00183313}.   \ns
\end{proof}

\begin{prop}\label{K002021} Jika $A$ suatu aljabar-$C^*$ beridentitas, maka fungsi $\omega\colon \inv A \sto \ofml U(A)$
yang didefinisikan dalam proposisi sebelumnya kontinu.
\end{prop}

\begin{proof}[\emph{Petunjuk untuk bukti}] Gunakan proposisi~\ref{000646fa} dan~\ref{001449}.   \ns
\end{proof}

\begin{prop}\label{K002024} Misalkan $s = u\abs s$ dekomposisi polar dari elemen invertibel $s$ dalam aljabar-$C^*$ beridentitas~$A$.
Maka $u \sim_h s$ dalam $\inv A$.
\end{prop}

\begin{proof}[\emph{Petunjuk untuk bukti}] Untuk $0 \le t \le 1$, misalkan $c_t = u(t\abs s + (1 - t)\vc 1)$. Simpulkan dari
proposisi~\ref{0018301} bahwa terdapat $\epsilon \in (0,1]$ sedemikian sehingga $\abs s \ge \epsilon \vc 1$. Gunakan proposisi yang sama
untuk menunjukkan bahwa $t\abs s + (1 - t)\vc 1$ invertibel untuk setiap $t \in [0,1]$.   \ns
\end{proof}

\begin{prop}\label{K002027} Misalkan $u$ dan $v$ elemen-elemen uniter dalam aljabar-$C^*$ beridentitas~$A$. Jika $u \sim_h v$ dalam $\inv A$,
maka $u \sim_h v$ dalam~$\ofml U(A)$.
\end{prop}

\begin{prop}\label{K002031} Misalkan $u_1$, $u_2$, $u_3$, dan $u_4$ elemen-elemen uniter dalam aljabar-$C^*$ beridentitas~$A$.
Jika $u_1 \sim_h u_2$ dalam $\ofml U(A)$ dan $u_3 \sim_h u_4$ dalam $\ofml U(A)$, maka $u_1u_3 \sim_h u_2u_4$ dalam~$\ofml U(A)$.
\end{prop}

\begin{exam}\label{0060116} Misalkan $h$ suatu elemen swaadjoin dari aljabar-$C^*$ beridentitas~$A$. Maka $\exp(ih)$ uniter dan
homotop dengan~$\vc 1$ dalam~$\ofml U(A)$.
\end{exam}

\begin{proof}[\emph{Petunjuk untuk bukti}] Pertimbangkan lintasan $c\colon [0,1] \sto \ofml U(A)\colon t \mapsto \exp(ith)$. Gunakan % SOURCE-CORRECTION: CH17-C005
proposisi~\ref{001449}.  \ns
\end{proof}

\begin{notn} Jika $p$ dan $q$ proyeksi dalam suatu aljabar-$C^*$ $A$, kita menulis $p \sim q$
 \index{ekuivalensi Murray--von Neumann}%
 \index{ekuivalensi!Murray--von Neumann}%
 \index{<binrelequiv@$p \sim q$ (ekuivalensi Murray--von Neumann)}%
($p$ \df{ekuivalen Murray--von Neumann} dengan~$q$) jika terdapat elemen $v \in A$
sedemikian sehingga $v^*v = p$ dan $vv^* = q$. Perhatikan bahwa $v$ demikian secara otomatis merupakan isometri
parsial. Kita akan menyebutnya
 \index{realisasi ekuivalensi Murray--von Neumann}%
 \index{ekuivalensi!Murray--von Neumann!realisasi}%
 \index{Murray--von Neumann!ekuivalensi!realisasi}%
\emph{isometri parsial yang merealisasikan ekuivalensi}.
\end{notn}

\begin{prop}\label{0060216} Relasi ekuivalensi Murray--von Neumann $\sim$ merupakan
relasi ekuivalensi pada keluarga $\fml P(A)$ dari proyeksi-proyeksi dalam suatu aljabar-$C^*$~$A$.
\end{prop}

\begin{prop}\label{0060144} Misalkan $a$ dan $b$ elemen-elemen swaadjoin dari suatu aljabar-$C^*$ beridentitas. Jika $a \sim_s b$, maka $a \sim_u b$.
Bahkan, jika $b = sas^{-1}$ dan $s = u\abs s$ adalah dekomposisi polar dari $s$, maka $b = uau^*$.
\end{prop}

\begin{proof}[\emph{Petunjuk untuk bukti}] Andaikan terdapat elemen invertibel $s$ sedemikian sehingga $b = sas^{-1}$. Misalkan $s = u\abs s$
dekomposisi polar dari~$s$. Tunjukkan bahwa $a$ komutatif dengan ${\abs s}^2$, sehingga juga komutatif dengan setiap elemen dalam aljabar $C^*(\vc 1,{\abs s}^2)$.
Secara khusus, $a$ komutatif dengan~${\abs s}^{-1}$. Dari sini diperoleh $uau^* = b$.   \ns
\end{proof}

\begin{prop}\label{0060221} Jika $p$ dan $q$ proyeksi dalam suatu aljabar-$C^*$, maka
   \[ p \sim_h q \implies p \sim_u q \implies p \sim q\,. \]
\end{prop}

\begin{proof}[\emph{Petunjuk untuk bukti}] Dalam petunjuk ini $\vc 1 = \vc 1_{\wt A}$. Untuk implikasi pertama, tunjukkan bahwa tanpa mengurangi keumuman
kita dapat mengandaikan $\norm{p - q} < \frac12$. Misalkan $s = pq + (\vc 1 - p)(\vc 1 - q)$. Buktikan bahwa $p \sim_s q$. Untuk
itu, tuliskan $s - \vc 1$ sebagai $p(q - p) + (\vc 1 - p)((\vc 1 - q) - (\vc 1 - p))$ dan gunakan akibat~\ref{cor_Neumann_series}.
Kemudian gunakan proposisi~\ref{0060144}.

Untuk membuktikan implikasi kedua, perhatikan bahwa jika $upu^* = q$ untuk suatu elemen uniter dalam~$\wt A$, maka $up$ merupakan isometri parsial
dalam~$A$.  \ns
\end{proof}

Secara umum, kebalikan dari implikasi kedua, $p \sim q \implies p \sim_u q$, dalam proposisi sebelumnya tidak berlaku (lihat % SOURCE-CORRECTION: CH17-C024
contoh~\ref{0060233} di bawah). Namun, untuk proyeksi $p$ dan $q$ dalam suatu aljabar-$C^*$ beridentitas, jika berlaku baik $p \sim q$ maupun
$\vc 1 - p \sim \vc 1 - q$, maka kita dapat menyimpulkan $p \sim_u q$.

\begin{prop}\label{0060224} Misalkan $p$ dan $q$ proyeksi dalam suatu aljabar-$C^*$ beridentitas~$A$. Maka $p \sim_u q$ jika dan hanya jika
$p \sim q$ dan $\vc 1 - p \sim \vc 1 - q$.
\end{prop}

\begin{proof}[\emph{Petunjuk untuk bukti}] Untuk arah sebaliknya, andaikan isometri parsial $v$ merealisasikan ekuivalensi $p \sim q$
dan $w$ merealisasikan $\vc 1 - p \sim \vc 1 - q$. Pertimbangkan elemen $u = v + w + \vc j$ dalam~$\wt A$.   \ns
\end{proof}

\begin{defn} Suatu elemen $s$ dari aljabar-$C^*$ beridentitas disebut
 \index{isometri!dalam aljabar-$C^*$ beridentitas}%
\df{isometri} jika $s^*s = \vc 1$.
\end{defn}

\begin{exer} Jelaskan mengapa terminologi dalam definisi sebelumnya masuk akal.
\end{exer}

Tidak sulit untuk melihat bahwa kebalikan dari implikasi kedua dalam proposisi~\ref{0060221} pada umumnya gagal.

\begin{exam}\label{0060233} Jika $p$ dan $q$ proyeksi dalam suatu aljabar-$C^*$, maka
   \[ p \sim q \,\,\,\,\not\!\!\!\!\implies p \sim_u q\,. \]
Sebagai contoh, jika $s$ merupakan isometri tak uniter (seperti geser unilateral), maka $s^*s \sim ss^*$, tetapi $s^*s \nsim_u ss^*$.
\end{exam}

\begin{proof}[\emph{Petunjuk untuk bukti}] Buktikan, dan ingatlah, bahwa tidak ada proyeksi tak nol yang dapat ekuivalen Murray--von Neumann
dengan proyeksi nol.   \ns
\end{proof}

\begin{rem} Kebalikan dari implikasi pertama, $p \sim_u q \implies p \sim_h q$, dalam proposisi~\ref{0060221} juga tidak berlaku % SOURCE-CORRECTION: CH17-C025
secara umum untuk proyeksi dalam suatu aljabar-$C^*$. Namun, contoh yang menggambarkan gejala ini tidak mudah ditemukan. Untuk melihat
apa yang terlibat, lihat~\cite{RordamLL:2000}, contoh 2.2.9 dan 11.3.4.
\end{rem}

Alangkah baiknya jika implikasi-implikasi dalam proposisi~\ref{0060221} dapat dibalik. Namun, seperti telah kita lihat, hal itu tidak terjadi.
Salah satu cara menghadapi fakta yang membandel, sebagaimana dicatat pada awal bab ini, adalah memberi objek-objek matematika
yang sedang kita kaji ruang yang lebih luas untuk bergerak. Beralihlah ke matriks. Proposisi~\ref{0060236} dan~\ref{0060244} merupakan contoh
cara kerja teknik ini. Dalam arti tertentu, kita dapat membuat implikasi-implikasi dalam~\ref{0060221} berbalik arah.

\begin{notn} Jika $a_1$, $a_2$, \dots $a_n$ elemen-elemen dari suatu aljabar-$C^*$ $A$, maka
 \index{diagonal@$\diag(a_1,\dots,a_n)$ (matriks diagonal dalam $\M nA$)}%
$\diag(a_1,a_2,\dots,a_n)$ adalah matriks diagonal dalam $\M nA$ yang diagonal utamanya terdiri atas
elemen-elemen $a_1$, \dots, $a_n$. Notasi ini juga kita gunakan untuk matriks blok. Sebagai
contoh, jika $a$ matriks $m \times m$ dan $b$ matriks $n \times n$, maka
$\diag(a,b)$ adalah matriks $(m+n) \times (m+n)$ $\begin{bmatrix} a & \vc 0 \\ \vc 0 & b
\end{bmatrix}$. (Tentu saja, $\vc 0$ di sudut kanan atas matriks ini adalah matriks $m \times n$ yang semua
entrinya nol, sedangkan $\vc 0$ di sudut kiri bawah adalah matriks $n \times m$ yang setiap entrinya nol.)
\end{notn}

\begin{prop}\label{0060236} Misalkan $p$ dan $q$ proyeksi dalam suatu aljabar-$C^*$~$A$. Maka
   \[ p \sim q \implies \diag(p,\vc 0) \sim_u \diag(q,\vc 0) \text{ dalam $\M 2A$.} \]
\end{prop}

\begin{proof}[\emph{Petunjuk untuk bukti}] Misalkan $v$ isometri parsial dalam $A$ yang merealisasikan ekuivalensi Murray--von Neumann $p \sim q$.
Pertimbangkan matriks $u = \begin{bmatrix} v  &  \vc 1_{\wt A} - q  \\  \vc 1_{\wt A} - p  &  v^* \end{bmatrix}$.  \ns
\end{proof}

Ingat dari akibat~\ref{001311005fa} bahwa spektrum setiap elemen uniter dari suatu aljabar-$C^*$ beridentitas terletak pada lingkaran satuan.
Hebatnya, jika spektrumnya bukan seluruh lingkaran, elemen itu homotop dengan~$\vc 1$ dalam~$\ofml U(A)$.

\begin{prop}\label{0060121} Misalkan $u$ suatu elemen uniter dalam aljabar-$C^*$ beridentitas. Jika $\sigma(u) \ne \T$, maka $u \sim_h \vc 1$ dalam~$\ofml U(A)$.
\end{prop}

\begin{proof}[\emph{Petunjuk untuk bukti}] Pilih $\theta \in \R$ sedemikian sehingga $\exp(i\theta) \notin \sigma(u)$. Maka terdapat fungsi kontinu yang tunggal
   \[ \phi\colon \sigma(u) \sto (\theta, \theta + 2\pi)\colon \exp(it) \mapsto t\,. \]
Misalkan $h = \phi(u)$ dan gunakan contoh~\ref{0060116}.   \ns
\end{proof}

\begin{exam}\label{0060129a} Jika $A$ suatu aljabar-$C^*$ beridentitas, maka $\begin{bmatrix} \vc 0  &  \vc 1  \\  \vc 1  &  \vc 0  \end{bmatrix} \sim_h
\begin{bmatrix} \vc 1  &  \vc 0  \\ \vc 0  & \vc 1 \end{bmatrix}$ dalam~$\ofml U(\M 2A)$.
\end{exam}

Matriks $J = \begin{bmatrix} \vc 0  &  \vc 1  \\ \vc 1 & \vc 0 \end{bmatrix}$ sangat berguna jika dipakai bersama
proposisi~\ref{K002031} untuk membentuk homotopi antara matriks-matriks dalam $\ofml U(A)$. Dalam memeriksa beberapa contoh
berikutnya, matriks ini sangat membantu. Perhatikan bahwa mengalikan matriks $2 \times 2$ dari kanan dengan $J$ menukar kolom-kolomnya; mengalikan
dari kiri dengan $J$ menukar baris-barisnya; dan mengalikan dari kiri serta kanan dengan $J$ menukar elemen-elemen pada kedua diagonal.

\begin{exam}\label{0060129b} Jika $u$ dan $v$ elemen-elemen uniter dalam suatu aljabar-$C^*$ beridentitas~$A$, maka
   \[ \diag(u,v) \sim_h \diag(v,u) \]
dalam~$\ofml U(\M 2A)$.
\end{exam}

\begin{exam}\label{0060129c} Jika $u$ dan $v$ elemen-elemen uniter dalam suatu aljabar-$C^*$ beridentitas $A$, maka
  \[ \diag(u,v) \sim_h \diag(uv, \vc 1) \]
dalam $\ofml U(\M 2A)$.
\end{exam}

\begin{exam}\label{0060129d} Jika $u$ dan $v$ elemen-elemen uniter dalam suatu aljabar-$C^*$ beridentitas~$A$, maka
   \[ \diag(uv,\vc 1) \sim_h \diag(vu, \vc 1) \]
dalam~$\ofml U(\M 2A)$.
\end{exam}

\begin{prop}\label{0060244} Misalkan $p$ dan $q$ proyeksi dalam suatu aljabar-$C^*$~$A$. Maka
   \[ p \sim_u q \implies \diag(p,\vc 0) \sim_h \diag(q,\vc 0) \text{ dalam $\ofml P(\M 2A)$.} \]
\end{prop}

\begin{exam}\label{0060247fa} Dua proyeksi dalam $\mathbf M_n$ ekuivalen Murray--von Neumann jika dan hanya jika keduanya memiliki teras yang sama,
yang pada gilirannya ekuivalen dengan jangkauan keduanya berdimensi sama.
\end{exam}

\begin{proof}[\emph{Petunjuk untuk bukti}] Gunakan proposisi~\ref{tr0121}--\ref{tr0124} dan~\ref{pi0117}--\ref{pi0121}.   \ns
\end{proof}

Dalam contoh~\ref{000533a} kita memperkenalkan aljabar matriks $n \times n$ $\M nA$, dengan $A$ suatu aljabar. Jika $A$ suatu
aljabar-$*\,$ dan $n \in \N$, kita dapat memperlengkapi $\M nA$ dengan involusi secara alami. Involusi itu didefinisikan, seperti yang dapat diduga,
sebagai analog dari ``transposisi konjugat'': jika $\vc a \in \M nA$, maka
   \[ \vc a^* = \bigl[a_{ij}\bigr]^* := \bigl[{a_{ji}}^*\bigr]. \]

Jika $A$ suatu aljabar-$C^*$, kita dapat memperkenalkan norma pada $\M nA$ yang membuatnya juga menjadi aljabar-$C^*$.
 \index{Ma@$\M nA$!sebagai aljabar-$C^*$}%
 \index{C*@aljabar-$C^*$!$\M nA$ sebagai suatu}%
Pilih representasi setia $\phi\colon A \sto \ofml B(H)$ dari $A$, dengan $H$ suatu ruang Hilbert. Untuk setiap $n \in \N$
definisikan
   \[ \phi_n\colon \M nA \sto \ofml B(H^n)\colon \bigl[a_{ij}\bigr] \mapsto \bigl[ \phi(a_{ij})\bigr] \]
dengan $H^n$ jumlah langsung $n$ rangkap dari $H$ dengan dirinya sendiri. Kemudian definisikan norma suatu elemen $\vc a \in \M nA$ sebagai
norma operator $\phi_n(\vc a) \in \ofml B(H^n)$. Yakni, $\norm{\vc a} := \norm{\phi_n(\vc a)}$.

Norma ini merupakan norma-$C^*$ pada $\M nA$. Norma tersebut tidak bergantung pada representasi khusus yang kita pilih berkat
ketunggalan norma-$C^*$ (lihat akibat~\ref{00152144}).

\begin{prop} Misalkan $A$ suatu aljabar-$C^*$ dan $\vc a \in \M nA$. Maka
   \[ \max\{\norm{a_{ij}}\colon 1 \le i,j \le n\} \le \norm{\vc a} \le \sum_{i,j = 1}^n \norm{a_{ij}}. \]
\end{prop}

\begin{proof}[\emph{Petunjuk untuk bukti}] Karena setiap representasi setia merupakan isometri (lihat proposisi~\ref{00152181}),
Anda dapat menyederhanakan persoalan dengan memandang representasi $\phi\colon A \sto \ofml B(H)$ yang dibahas di atas sebagai
pemetaan inklusi.

Untuk ketaksamaan pertama, bagi $v \in H$ definisikan vektor $\vc E_j(v) \in H^n$ sebagai $n$-tupel
yang semua entrinya nol kecuali entri ke-$j$, yang bernilai~$v$. Kemudian periksa bahwa
   \[ \norm{a_{ij}v} \le \norm{\vc a\vc E_j(v)} \le \norm{\vc a} \]
setiap kali $\vc a \in \M nA$, $v \in H$, dan $\norm v \le 1$.

Untuk ketaksamaan kedua, tunjukkan bahwa
   \[ {\norm{\vc a\vc v}}^2 \le \sum_{i=1}^n\biggl(\sum_{j=1}^n \norm{a_{ij}}\biggr)^2 \le \biggl(\sum_{i=1}^n\sum_{j=1}^n \norm{a_{ij}}\biggr)^2 \]
setiap kali $\vc a \in \M nA$, $\vc v \in H^n$, dan $\norm{\vc v} \le 1$.    \ns
\end{proof}


























\section{Semigrup Proyeksi}

\begin{notn} Jika $A$ suatu aljabar-$C^*$ dan $n \in \N$, kita menetapkan
 \index{P@$\fml P_n(A)$, $\fml P_\infty(A)$ (proyeksi dalam aljabar matriks)}%
  \begin{align*}
       \fml P_n(A) &= \fml P(\M nA) \quad\text{ dan} \\
       \fml P_\infty(A) &= \bigcup_{n=1}^\infty \fml P_n(A).
  \end{align*}
\end{notn}

Sekarang kita memperluas ekuivalensi Murray--von Neumann ke matriks-matriks yang ukurannya berbeda.

\begin{notn} Jika $A$ suatu aljabar-$C^*$, misalkan
 \index{Ma@$\M {n,m}A$ (matriks dengan entri dari~$A$)}%
$\M {n,m}A$ menyatakan himpunan matriks $n \times m$ berentri dalam~$A$.
\end{notn}

Selanjutnya kita memperluas ekuivalensi Murray--von Neumann ke matriks proyeksi yang ukurannya tidak harus sama.

\begin{defn}\label{0060316fa} Jika $A$ suatu aljabar-$C^*$, $p \in \fml P_m(A)$, dan $q \in \fml P_n(A)$, kita menetapkan
 \index{ekuivalensi Murray--von Neumann}%
 \index{ekuivalensi!Murray--von Neumann}%
 \index{<binrelequiv@$p \sim q$ (ekuivalensi Murray--von Neumann)}%
$p \sim q$ jika terdapat $v \in \M {n,m}A$ sedemikian sehingga
  \[ v^*v = p \qquad\text{ dan } \qquad  vv^* = q. \]
Kita akan memperluas penggunaan nama \emph{ekuivalensi Murray--von Neumann} untuk relasi baru ini pada~$\fml P_\infty(A)$.
\end{defn}

\begin{prop}\label{0060317} Relasi $\sim$ yang didefinisikan di atas merupakan relasi
ekuivalensi pada~$\fml P_\infty(A)$.
\end{prop}

\begin{defn} Untuk setiap aljabar-$C^*$ $A$, kita mendefinisikan operasi biner $\oplus$ pada $\fml
P_\infty(A)$ dengan
   \[ p \oplus q = \diag(p,q)\,. \]
Jadi, jika $p \in \fml P_m(A)$ dan $q \in \fml P_n(A)$, maka $p \oplus q \in \fml
P_{m+n}(A)$.
\end{defn}

Dalam proposisi berikut, $\vc 0_n$ adalah identitas aditif dalam~$\fml P_n(A)$.

\begin{prop}\label{0060324} Misalkan $A$ suatu aljabar-$C^*$ dan $p \in \fml P_\infty(A)$.
 \index{<sum@$p \oplus q$ (operasi biner dalam $\fml P_\infty(A)$)}%
Maka $p \sim p \oplus \vc 0_n$ untuk setiap $n \in \N$.
\end{prop}

\begin{prop}\label{0060327} Misalkan $A$ suatu aljabar-$C^*$ dan $p$, $p'$, $q$, $q' \in
\fml P_\infty(A)$. Jika $p \sim p'$ dan $q \sim q'$, maka $p \oplus q \sim p' \oplus q'$.
\end{prop}

\begin{prop}\label{0060331} Misalkan $A$ suatu aljabar-$C^*$ dan $p$, $q \in \fml P_\infty(A)$.
Maka $p \oplus q \sim q \oplus p$.
\end{prop}

\begin{prop}\label{0060334} Misalkan $A$ suatu aljabar-$C^*$ dan $p$, $q \in \fml P_n(A)$ untuk
suatu $n$. Jika $p \perp q$, maka $p + q$ suatu proyeksi dalam $\M nA$ dan $p + q \sim p \oplus q$.
\end{prop}

\begin{prop}\label{0060337} Jika $A$ suatu aljabar-$C^*$, maka operasi $\oplus$ pada $\fml P_\infty(A)$ bersifat
asosiatif secara ketat dan komutatif hanya hingga ekuivalensi Murray--von Neumann. % SOURCE-CORRECTION: CH17-C006
\end{prop}

\begin{notn}\label{0060411} Jika $A$ suatu aljabar-$C^*$ dan $p \in \fml P_\infty(A)$, misalkan
 \index{<bracsqr@$\sbsb{[p]}{\fml D}$ (kelas ekuivalensi $\sim$ dari $p$)}%
$\sbsb{[p\,]}{\fml D}$ adalah kelas ekuivalensi yang memuat $p$ dan ditentukan oleh
relasi ekuivalensi~$\sim$. Tetapkan pula
 \index{D@$\fml D(A)$ (himpunan kelas ekuivalensi proyeksi)}%
$\fml D(A) := \{\,\sbsb{[p\,]}{\fml D}\colon p \in \fml P_\infty(A)\,\}$.
\end{notn}

\begin{defn}\label{0060414} Misalkan $A$ suatu aljabar-$C^*$. Definisikan operasi biner $+$ pada
$\fml D(A)$ dengan
 \index{<binop@$\sbsb{[p\,]}{\fml D} + \sbsb{[q\,]}{\fml D}$ (penjumlahan dalam $\fml D(A)$)}%
   \[ \sbsb{[p\,]}{\fml D} + \sbsb{[q\,]}{\fml D} := \sbsb{[p \oplus q\,]}{\fml D} \]
dengan $p$, $q \in \fml P_\infty(A)$.
\end{defn}

\begin{prop}\label{0060417} Operasi $+$ yang didefinisikan dalam~\ref{0060414} terdefinisi dengan baik
dan menjadikan $\fml D(A)$ suatu semigrup komutatif.
\end{prop}

\begin{exam}\label{0060421} Semigrup $\fml D(\C)$ isomorfik dengan semigrup aditif bilangan bulat tak negatif; yaitu, % SOURCE-CORRECTION: CH17-C007
    \[ \fml D(\C) \cong \Z^+ = \{0,1,2,\dots\}\,.  \]
\end{exam}

\begin{proof}[\emph{Petunjuk untuk bukti}] Gunakan contoh~\ref{0060247fa}. \ns
\end{proof}

\begin{exam}\label{0060424} Jika $H$ suatu ruang Hilbert terpisahkan berdimensi tak hingga, maka % SOURCE-CORRECTION: CH17-C008
   \[ \fml D(\ofml B(H)) \cong \Z^+ \cup \{\infty\}. \]
(Gunakan penjumlahan biasa pada $\Z^+$ dan tetapkan $n + \infty = \infty + n = \infty$ setiap kali $n
\in \Z^+ \cup \{\infty\}$.)
\end{exam}

\begin{exam}\label{0060427} Untuk aljabar-$C^*$ $\C \oplus \C$, kita mempunyai
   \[ \fml D(\C \oplus \C) \cong \Z^+ \oplus \Z^+\,. \]
\end{exam}


























\section{Konstruksi Grothendieck}
Dalam konstruksi medan bilangan real, kita dapat menggunakan teknik yang sama untuk beralih dari (semigrup aditif) bilangan asli ke
(grup aditif semua) bilangan bulat seperti yang kita gunakan untuk beralih dari (semigrup multiplikatif) bilangan bulat tak nol ke (grup multiplikatif
semua) bilangan rasional tak nol. Teknik ini disebut \emph{konstruksi Grothendieck}.

\begin{defn}\label{0062111} Misalkan $(S,+)$ suatu semigrup komutatif. Definisikan relasi
 \index{<binrelequiv@$(a,b) \sim (c,d)$ (ekuivalensi Grothendieck)}%
$\sim$ pada $S \times S$ dengan
   \[ (a,b) \sim (c,d) \text{ jika terdapat $k \in S$ sedemikian sehingga } a + d + k = b + c + k. \]
\end{defn}

\begin{prop}\label{0062114} Relasi $\sim$ yang didefinisikan di atas merupakan relasi ekuivalensi.
\end{prop}

\begin{notn} Untuk relasi ekuivalensi $\sim$ yang didefinisikan dalam~\ref{0062111}, kelas ekuivalensi
yang memuat pasangan $(a,b)$ akan dinotasikan dengan % SOURCE-CORRECTION: CH17-C009
 \index{<bracpnt@$\langle a,b \rangle$ (pasangan dalam konstruksi Grothendieck)}%
$\langle a,b \rangle$, bukan dengan~$[(a,b)]$.
\end{notn}

\begin{defn} Misalkan $(S,+)$ suatu semigrup komutatif. Pada $G(S)$, definisikan operasi biner
 \index{penjumlahan!pada grup Grothendieck}%
(yang juga dinotasikan dengan $+$) sebagai berikut:
  \[ \langle a,b \rangle + \langle c,d \rangle := \langle a + c,b + d\rangle\,. \]
\end{defn}

\begin{prop}\label{0062119} Operasi $+$ yang didefinisikan di atas terdefinisi dengan baik, dan di bawah
operasi ini $G(S)$ menjadi grup Abelian. % SOURCE-CORRECTION: CH17-C010
\end{prop}

Grup Abelian $(G(S),+)$ disebut
 \index{G@$G(S)$ (grup Grothendieck dari $S$)}%
 \index{Grothendieck!grup}%
 \index{grup!Grothendieck}%
\df{grup Grothendieck} dari~$S$.

\begin{prop}\label{0062124} Untuk suatu semigrup $S$ dan sembarang $a \in S$, definisikan
pemetaan
 \index{gamma@$\sbsb{\gamma}S$ (pemetaan Grothendieck)}%
   \[ \sbsb{\gamma}S \colon S \sto G(S) \colon s \mapsto \langle s + a, a \rangle\,. \]
Pemetaan $\sbsb{\gamma}S$, yang disebut
 \index{Grothendieck!pemetaan}%
\df{pemetaan Grothendieck}, terdefinisi dengan baik dan merupakan homomorfisme semigrup.
\end{prop}

Pernyataan bahwa pemetaan Grothendieck terdefinisi dengan baik berarti bahwa definisinya tidak bergantung
pada pilihan~$a$. Kita sering hanya menulis $\gamma$ untuk~$\sbsb{\gamma}S$.

\begin{exam}\label{0062126} Baik $\N = \{1,2,3,\dots\}$ maupun $\Z^+ = \{0,1,2,\dots\}$ merupakan
semigrup komutatif di bawah penjumlahan. Keduanya menghasilkan grup Grothendieck yang sama,
   \[G(\N) = G(\Z^+) = \Z\,. \]
\end{exam}

Tidak ada yang menjamin bahwa grup Grothendieck dari suatu semigrup sembarang akan sangat menarik.

\begin{exam}\label{0062127} Misalkan $S$ semigrup aditif komutatif $\Z^+ \cup \{\infty\}$. Maka $G(S) = \{\vc 0\}$.
\end{exam}

\begin{exam}\label{0062128} Misalkan $\Z_0$ semigrup multiplikatif (komutatif) bilangan bulat tak nol.
Maka $G(\Z_0) = \Q_0$, yaitu grup multiplikatif Abelian dari bilangan rasional tak
nol.
\end{exam}

\begin{prop}\label{0062141} Jika $S$ suatu semigrup komutatif, maka
   \[ G(S) = \{\gamma(s) - \gamma(t)\colon s,t \in S\}\,. \]
\end{prop}

\begin{prop}\label{0062142} Jika $r$, $s \in S$, dengan $S$ suatu semigrup komutatif, maka $\gamma(r) = \gamma(s)$ jika dan hanya jika
terdapat $t \in S$ sedemikian sehingga $r + t = s + t$.
\end{prop}

\begin{defn} Suatu semigrup komutatif $S$ memiliki
 \index{sifat pembatalan}%
\df{sifat pembatalan} jika setiap kali $r$, $s$, $t \in S$ memenuhi $r + t = s + t$, berlaku $r = s$.
\end{defn}

\begin{cor}\label{0062144} Misalkan $S$ suatu semigrup komutatif. Pemetaan Grothendieck $\sbsb\gamma S\colon S \sto G(S)$ injektif
jika dan hanya jika $S$ memiliki sifat pembatalan.
\end{cor}

Proposisi berikut menyatakan sifat universal grup Grothendieck.

\begin{prop}\label{0062151} Misalkan $S$ suatu semigrup (aditif) komutatif dan $G(S)$ grup Grothendieck-nya.
 \index{universal!sifat!pemetaan Grothendieck}%
 \index{Grothendieck!pemetaan!sifat universal}%
Jika $H$ suatu grup Abelian dan $\phi\colon S \sto H$ suatu pemetaan aditif, maka terdapat
homomorfisme grup tunggal $\psi\colon G(S) \sto H$ sedemikian sehingga diagram berikut
komutatif.
 \begin{equation}\label{0062151i}
   \xy
     \qtriangle/{>}`{>}`{-->}/[S`\abs{G(S)}`\abs{H};\gamma`\phi`\abs{\psi}]
     \morphism(1100,500)|r|/{-->}/<0,-500>[G(S)`H;\psi]
   \endxy
 \end{equation}
\end{prop}

Dalam diagram sebelumnya, $\abs{G(S)}$ dan $\abs{H}$ hanyalah $G(S)$ dan $H$ yang dipandang sebagai
semigrup, sedangkan $\abs{\psi}$ adalah homomorfisme semigrup yang bersesuaian. Dengan kata lain,
funktor pelupa $\abs{\hphantom{G}}$ ``melupakan'' hanya identitas dan invers, tetapi tidak
operasi penjumlahan. Jadi, segitiga di sebelah kiri merupakan diagram komutatif dalam
kategori semigrup dan homomorfisme semigrup.

\begin{prop}\label{0062154} Misalkan $\phi\colon S \sto T$ homomorfisme semigrup
komutatif. Maka pemetaan $\sbsb{\gamma}T \circ \phi\colon S \sto G(T)$ aditif. Berdasarkan
proposisi~\ref{0062151}, terdapat homomorfisme grup tunggal $G(\phi)\colon G(S)
\sto G(T)$ sedemikian sehingga diagram berikut komutatif.
   \[ \xy
          \Square[S`T`G(S)`G(T);\phi`\sbsb{\gamma}S`\sbsb{\gamma}T`G(\phi)]
      \endxy \]
\end{prop}

\begin{prop}\label{0062155} Pasangan pemetaan $S \mapsto G(S)$, yang memetakan semigrup komutatif ke
 \index{funktor!Grothendieck}%
 \index{Grothendieck!konstruksi!sifat funktorial}%
grup Grothendieck yang bersesuaian, dan $\phi \mapsto G(\phi)$, yang memetakan homomorfisme semigrup ke
homomorfisme-homomorfisme grup (sebagaimana didefinisikan dalam~\ref{0062154}), merupakan funktor kovarian dari kategori % SOURCE-CORRECTION: CH17-C011
semigrup komutatif dan homomorfisme semigrup ke kategori grup Abelian
dan homomorfisme grup.
\end{prop}

Salah satu keuntungan kecil dari penyertaan funktor pelupa yang agak bertele-tele dalam
diagram~\eqref{0062151i} adalah bahwa hal itu memungkinkan kita memandang pemetaan Grothendieck
$\gamma\colon S \mapsto \sbsb{\gamma}S$ sebagai transformasi alami antarfungtor.

\begin{cor}[Kealamian Pemetaan Grothendieck]\label{0062157} Misalkan $\abs{\hphantom{G}}$ funktor pelupa
pada grup Abelian yang ``melupakan'' identitas dan invers, tetapi tidak
operasi grup, seperti dalam~\ref{0062151}. Maka $\abs{G(\hphantom{S})}$ merupakan funktor kovarian dari
kategori semigrup komutatif dan homomorfisme semigrup ke dirinya sendiri.
 \index{Grothendieck!pemetaan!kealamian}%
 \index{alami!transformasi!pemetaan Grothendieck sebagai suatu}%
Selanjutnya, pemetaan Grothendieck $\gamma\colon S \mapsto \sbsb{\gamma}S$ merupakan transformasi
alami dari funktor identitas ke funktor~$\abs{G(\hphantom{S})}$.
  \[ \xy
          \Square[S`T`\abs{G(S)}`\abs{G(T)};\phi`\sbsb{\gamma}S`\sbsb{\gamma}T`\abs{G(\phi)}]
      \endxy \]
\end{cor}





\section{Grup $\mathbf{\emph{K}_0}$ untuk Aljabar-$C^*$ Beridentitas}

\begin{defn} Misalkan $A$ aljabar-$C^*$ beridentitas. Misalkan
 \index{K0@$K_0(A)$!untuk $A$ beridentitas}%
$K_0(A) := G(\fml D(A))$, grup Grothendieck dari semigrup $\fml D(A)$ yang didefinisikan dalam~\ref{0060411} dan~\ref{0060414},
dan definisikan
 \index{<bracsqr@$[p\,]$ (citra $[p\,]_{\fml D}$ di bawah pemetaan Grothendieck)}%
   \[ [\hphantom{P}]\colon \fml P_\infty(A) \sto K_0(A)
           \colon p \mapsto \sbsb{\gamma}{\fml D(A)}\bigl([p\,]_{\fml D}\bigr)\,. \]
\end{defn}

Proposisi berikut merupakan konsekuensi langsung dari definisi ini.

\begin{prop}\label{0062188fa} Misalkan $A$ aljabar-$C^*$ beridentitas dan $p$, $q \in \fml P_\infty(A)$. Jika $p$ dan $q$
ekuivalen Murray--von Neumann, maka $[p\,] = [q\,]$ dalam~$K_0(A)$.
\end{prop}

\begin{defn} Misalkan $A$ aljabar-$C^*$ beridentitas dan $p$, $q \in \fml P_\infty(A)$. Kita mengatakan bahwa $p$
 \index{ekuivalensi!stabil}%
 \index{ekuivalensi stabil}%
\df{ekuivalen secara stabil} dengan $q$ dan menulis
 \index{<binrelequivst@$p \sim_{st} q$ (ekuivalensi stabil)}%
$p \sim_{st} q$ jika terdapat proyeksi $r \in \fml P_\infty(A)$ sedemikian sehingga $p \oplus r \sim q \oplus r$.
\end{defn}

\begin{prop} Ekuivalensi stabil $\sim_{st}$ merupakan relasi ekuivalensi pada $\fml P_\infty(A)$.
\end{prop}

\begin{prop} Misalkan $A$ aljabar-$C^*$ beridentitas dan $p$, $q \in \fml P_\infty(A)$. Maka $p
\sim_{st} q$ jika dan hanya jika $p \oplus \vc 1_n \sim q \oplus \vc 1_n$ untuk suatu $n \in \N$.
\end{prop}

Di sini, tentu saja, $\vc 1_n$ adalah identitas perkalian dalam $\M nA$.

\begin{prop}[Gambaran Standar $K_0(A)$ ketika $A$ beridentitas]\label{0062224} Jika $A$ suatu
aljabar-$C^*$ beridentitas, maka
 \index{standar!gambaran $K_0(A)$!ketika $A$ beridentitas}%
 \index{gambaran (\seeonly{gambaran standar}))}%
   \begin{align*}
         K_0(A) &= \{ [p\,] - [q\,]\colon p,q \in \fml P_\infty(A)\}  \\
                &= \{ [p\,] - [q\,]\colon  p,q \in \fml P_n(A) \text{ untuk suatu $n \in \N$}\}\,.
   \end{align*}
\end{prop}

\begin{prop} Misalkan $A$ aljabar-$C^*$ beridentitas dan $p$, $q \in \fml P_\infty(A)$. Maka $[p
\oplus q\,] = [p\,] + [q\,]$.
\end{prop}

\begin{prop} Misalkan $A$ aljabar-$C^*$ beridentitas dan $p$, $q \in \fml P_n(A)$. Jika $p \sim_h q$
dalam $\fml P_n(A)$, maka $[p\,] = [q\,]$.
\end{prop}

\begin{prop} Misalkan $A$ aljabar-$C^*$ beridentitas dan $p$, $q \in \fml P_n(A)$. Jika $p \perp q$
dalam $\fml P_n(A)$, maka $p + q \in \fml P_n(A)$ dan $[p + q\,] = [p\,] + [q\,]$.
\end{prop}

\begin{prop} Misalkan $A$ aljabar-$C^*$ beridentitas dan $p$, $q \in \fml P_\infty(A)$. Maka
$[p\,] = [q\,]$ jika dan hanya jika $p \sim_{st} q$.
\end{prop}

 Proposisi berikut menetapkan suatu sifat universal dari $K_0(A)$ ketika $A$ beridentitas. Di dalamnya, funktor pelupa $\abs{\hphantom{G}}$
 adalah funktor yang dijelaskan dalam proposisi~\ref{0062151}.

\begin{prop}\label{0062244} Misalkan $A$ aljabar-$C^*$ beridentitas, $G$ grup Abelian, dan $\nu\colon \fml P_\infty(A) \sto \abs G$
homomorfisme semigrup yang memenuhi
 \index{universal!sifat!dari $K_0$ (kasus beridentitas)}%
 \index{K0@$K_0$!sifat universal dari}%
 \begin{enumerate}[label=\rm{(\alph*)}]
    \item $\nu(\vc 0_A) = \vc 0_G$ dan
    \item jika $p \sim_h q$ dalam $\fml P_n(A)$ untuk suatu $n$, maka $\nu(p) = \nu(q)$.
 \end{enumerate}
Maka terdapat homomorfisme grup tunggal $\wt\nu\colon K_0(A) \sto G$ sedemikian sehingga diagram berikut komutatif.
   \[\xy
      \qtriangle/{>}`{>}`{-->}/<1000,700>[\fml P_\infty(A)`\abs{K_0(A)}`\abs G;{[\hphantom{P}]}`\nu`\abs{\wt\nu}]
      \morphism(1600,700)|r|/{-->}/<0,-700>[K_0(A)`G;\wt\nu]
     \endxy\]
\end{prop}

\begin{proof}[\emph{Petunjuk untuk bukti}] Misalkan $\tau\colon \fml D(A) \sto \abs G\colon \sbsb{[p\,]}{\fml D} \mapsto \nu(p)$. % SOURCE-CORRECTION: CH17-C026
Periksa bahwa $\tau$ terdefinisi dengan baik dan merupakan homomorfisme semigrup. Kemudian gunakan proposisi~\ref{0062151}.  \ns
\end{proof}

\begin{defn} Homomorfisme-$*\,$ $\phi\colon A \sto B$ di antara aljabar-$C^*$ diperluas, untuk
setiap $n \in \N$, menjadi homomorfisme-$*\,$ $\phi\colon \M nA \sto \M nB$ dan juga (karena homomorfisme-$*\,$ membawa proyeksi
ke proyeksi) memiliki pembatasan berupa pemetaan $\phi$ dari $\fml P_\infty(A)$ ke $\fml P_\infty(B)$. % SOURCE-CORRECTION: CH17-C012
Untuk homomorfisme-$*\,$ $\phi$ semacam itu, definisikan
   \[ \nu\colon \fml P_\infty(A) \sto K_0(B)\colon p \mapsto [\phi(p)]\,. \]
Maka $\nu$ merupakan homomorfisme semigrup yang memenuhi syarat (a) dan (b) dalam proposisi~\ref{0062244}, sehingga
terdapat homomorfisme grup tunggal
 \index{K0@$K_0(\phi)$!kasus beridentitas}%
$K_0(\phi)\colon K_0(A) \sto K_0(B)$ sedemikian sehingga $K_0(\phi)\bigl([p\,]\bigr) = \nu(p)$ untuk setiap $p \in \fml P_\infty(A)$.
\end{defn}

\begin{prop} Pasangan pemetaan $A \mapsto K_0(A)$, $\phi \mapsto K_0(\phi)$ merupakan funktor kovarian
dari kategori aljabar-$C^*$ beridentitas dan homomorfisme-$*\,$ ke kategori grup Abelian dan
homomorfisme grup. Lebih lanjut, untuk semua homomorfisme-$*\,$ $\phi\colon A \sto B$ di antara aljabar-$C^*$ beridentitas, diagram
berikut komutatif.
  \[\xy
         \Square[\fml P_\infty(A)`\fml
         P_\infty(B)`K_0(A)`K_0(B);\phi`{[\hphantom{P}]}`{[\hphantom{P}]}`K_0(\phi)]
  \endxy\]
\end{prop}

\begin{notn} Untuk aljabar-$C^*$ $A$ dan $B$, misalkan
 \index{hom@$\Hom(A,B)$ (himpunan homomorfisme-$*\,$ di antara aljabar-$C^*$)}%
$\Hom(A,B)$ adalah keluarga semua homomorfisme-$*\,$ dari $A$ ke~$B$.
\end{notn}

\begin{defn} Misalkan $A$ dan $B$ aljabar-$C^*$ serta $a \in A$. Untuk $\phi$, $\psi \in \Hom(A,B)$,
misalkan
   \[ d_a(\phi,\psi) = \norm{\phi(a) - \psi(a)}\,. \]
Maka $d_a$ merupakan pseudometrik pada $\Hom (A,B)$. Topologi yang dibangkitkan oleh keluarga $\{d_a\colon a \in A\}$ adalah
 \index{topologi norma-titik}%
 \index{topologi!norma-titik}%
\df{topologi norma-titik} pada~$\Hom(A,B)$. (Untuk setiap $a \in A$, misalkan $\sfml B_a$ adalah keluarga bola terbuka dalam $\Hom(A,B)$ yang dibangkitkan oleh
pseudometrik~$d_a$. Keluarga $\bigcup\{\sfml B_a\colon a \in A\}$ merupakan subbasis bagi topologi norma-titik.)
\end{defn}

\begin{defn} Misalkan $A$ dan $B$ aljabar-$C^*$. Kita mengatakan bahwa homomorfisme-$*\,$ $\phi$, $\psi\colon A \sto B$
 \index{homotop!homomorfisme-$*\,$}%
 \index{<star@homomorfisme-$*\,$!homotopi dari}%
 \index{<binrelequivh@$\phi\sim_h\psi$ (homotopi homomorfisme-$*\,$)}%
\df{homotop}, dan menulis $\phi \sim_h \psi$, jika terdapat suatu fungsi
   \[ c\colon [0,1] \times A \sto B\colon (t,a) \mapsto c_t(a) \]
sedemikian sehingga
  \begin{enumerate}[label=\rm{(\alph*)}]
     \item untuk setiap $t \in [0,1]$, pemetaan $c_t\colon A \sto B \colon a \mapsto c_t(a)$ merupakan homomorfisme-$*\,$,
     \item untuk setiap $a \in A$, pemetaan $c(a)\colon [0,1] \sto B\colon t \mapsto c_t(a)$ merupakan lintasan (kontinu) dalam~$B$,
     \item $c_0 = \phi$, dan
     \item $c_1 = \psi$.
  \end{enumerate}
\end{defn}

\begin{prop} Dua homomorfisme-$*\,$ $\phi_0$, $\phi_1\colon A \sto B$ di antara aljabar-$C^*$
homotop jika dan hanya jika terdapat lintasan kontinu norma-titik dari $\phi_0$ ke $\phi_1$ dalam~$\Hom(A,B)$.
\end{prop}

\begin{defn} Kita mengatakan bahwa aljabar-$C^*$ $A$ dan $B$
 \index{homotop!ekuivalensi!homomorfisme-$*\,$}%
 \index{ekuivalensi!homotopi!homomorfisme-$*\,$}%
\df{ekuivalen secara homotopi} jika terdapat homomorfisme-$*\,$ $\phi\colon A \sto B$ dan $\psi\colon B \sto A$ sedemikian sehingga $\psi
\circ \phi \sim_h \id A$ dan $\phi \circ \psi \sim_h \id B$. Suatu aljabar-$C^*$ disebut
 \index{kontraktibel!aljabar-$C^*$}%
\df{kontraktibel} jika ekuivalen secara homotopi dengan~$\{\vc 0\}$.
\end{defn}

\begin{prop}\label{0062327} Misalkan $A$ dan $B$ aljabar-$C^*$ beridentitas. Jika $\phi \sim_h \psi$
dalam $\Hom(A,B)$, maka $K_0(\phi) = K_0(\psi)$.
\end{prop}

%\begin{proof} Lihat \cite{RordamLL:2000}, proposisi 3.2.6.  \ns \end{proof}

\begin{prop}\label{0062331} Jika aljabar-$C^*$ beridentitas $A$ dan $B$
 \index{homotopi!invariansi $K_0$}%
 \index{K0@$K_0$!invariansi homotopi dari}%
ekuivalen secara homotopi, maka $K_0(A) \cong K_0(B)$.
\end{prop}

\begin{prop}\label{0062334} Jika $A$ suatu aljabar-$C^*$ beridentitas, maka barisan eksak terbelah
   \begin{equation}\label{0062334i}
      \vc 0 \to A \to^\iota \wt A \two/->`<-/^Q_\psi \C \to \vc 0
   \end{equation}
\emph{(lihat \ref{0015213}, Kasus 1, butir (7)\,)} menginduksi barisan eksak terbelah lainnya
   \begin{equation}\label{0062334ii}
      \vc 0 \to K_0(A) \to^{K_0(\iota)} K_0(\wt A) \two/->`<-/^{K_0(Q)}_{K_0(\psi)} K_0(\C) \to \vc 0 \,.
   \end{equation}
\end{prop}

\begin{proof}[\emph{Petunjuk untuk bukti}] Definisikan dua fungsi tambahan
   \[ \mu\colon \wt A \sto A\colon a + \lambda \vc j \mapsto a \]
dan
   \[ \psi' \colon \C \sto \wt A \colon \lambda \mapsto \lambda \vc j\,. \]
Kemudian periksa bahwa
   \begin{enumerate}[label=\rm{(\alph*)}]
     \item $\mu \circ \iota = \id A$,
     \item $\iota \circ \mu + \psi' \circ Q = \id{\wt A}$,
     \item $Q \circ \iota = \vc 0$, dan
     \item $Q \circ \psi' = \id{\C}$. % SOURCE-CORRECTION: CH17-C013
   \end{enumerate}   \ns
\end{proof}

\begin{exam} Untuk setiap $n \in \N$,
 \index{K0@$K_0(A)$!adalah $\Z$ ketika $A = \mathbf M_n$}%
$K_0(\mathbf M_n) \cong \Z$.
\end{exam}

\begin{exam} Jika $H$ ruang Hilbert terpisahkan berdimensi tak hingga,
 \index{K0@$K_0(A)$!adalah $\vc 0$ ketika $A = \ofml B(H)$}%
maka $K_0(\ofml B(H)) \cong \vc 0$.
\end{exam}

\begin{defn} Ingat bahwa suatu ruang topologis $X$ disebut
 \index{kontraktibel!ruang topologis}%
\df{kontraktibel} jika terdapat titik $a$ dalam ruang tersebut dan fungsi kontinu $f\colon [0,1] \times X \sto X$ sedemikian sehingga $f(1,x) = x$
dan $f(0,x) = a$ untuk setiap $x \in X$.
\end{defn}

\begin{exam} Jika $X$ ruang Hausdorff kompak yang kontraktibel, maka
 \index{K0@$K_0(A)$!adalah $\Z$ ketika $A = \fml C(X)$, $X$ kontraktibel}%
$K_0(\fml C(X)) \cong \Z$.
\end{exam}


















\section{$\mathbf{\emph{K}_0}(A)$---Kasus Tak Beridentitas}
\begin{defn} Misalkan $A$ aljabar-$C^*$ tak beridentitas. Ingat bahwa barisan eksak terbelah
 ~\eqref{0062334i} bagi unitalisasi $A$ menginduksi barisan eksak terbelah
~\eqref{0062334ii} di antara grup-grup $K_0$ yang bersesuaian.
 \index{K0@$K_0(A)$!untuk $A$ tak beridentitas}%
Definisikan
   \[ \pi:=Q,\qquad K_0(A) = \ker(K_0(\pi))\,. \] % SOURCE-CORRECTION: CH17-C014
\end{defn}

\begin{prop}\label{0062413} Untuk aljabar-$C^*$ tak beridentitas $A$, pemetaan $[\hphantom{P}] \colon \fml
P_\infty(A) \sto K_0(\wt A)$ dapat dipandang sebagai pemetaan dari $\fml P_\infty(A)$ ke dalam~$K_0(A)$.
\end{prop}

\begin{prop}\label{0062416} Untuk aljabar-$C^*$ beridentitas maupun tak beridentitas, barisan
  \[ \vc 0 \to K_0(A) \to K_0(\wt A) \to K_0(\C) \to \vc 0 \]
bersifat eksak.
\end{prop}

\begin{prop}\label{0062418} Untuk aljabar-$C^*$ beridentitas maupun tak beridentitas, grup $K_0(A)$ (isomorfik
dengan) $\ker(K_0(\pi))$.
\end{prop}

\begin{prop}\label{0062424} Jika $\phi\colon A \sto B$ merupakan homomorfisme-$*\,$ di antara aljabar-$C^*$, maka
terdapat homomorfisme grup tunggal % SOURCE-CORRECTION: CH17-C015
 \index{K0@$K_0(\phi)$!kasus tak beridentitas}%
$K_0(\phi)$ yang membuat diagram berikut komutatif.
   \[\xy
     \hSquares/>`>`.>`>`=`>`>/[K_0(A)`K_0(\wt A)`K_0(\C)`K_0(B)`K_0(\wt
     B)`K_0(\C);`K_0(\pi_A)`K_0(\phi)`K_0(\wt\phi)```K_0(\pi_B)]
   \endxy\]
\end{prop}

\begin{prop}\label{0062511} Pasangan pemetaan $A \mapsto K_0(A)$, $\phi \mapsto K_0(\phi)$ merupakan
 \index{funktor!$K_0$}%
 \index{K0@$K_0$!sebagai funktor kovarian}%
funktor kovarian dari kategori $\cat{CSA}$ aljabar-$C^*$ dan homomorfisme-$*\,$ ke
kategori grup Abelian dan homomorfisme grup.
\end{prop}

Dalam proposisi~\ref{0062327} dan~\ref{0062331}, kita menyatakan invariansi homotopi dari
funktor $K_0$ untuk aljabar-$C^*$ beridentitas. Sekarang kita memperluas hasil tersebut ke sebarang
aljabar-$C^*$.

\begin{prop}\label{0062521} Misalkan $A$ dan $B$ aljabar-$C^*$. Jika $\phi \sim_h \psi$
dalam $\Hom(A,B)$, maka $K_0(\phi) = K_0(\psi)$.
\end{prop}

%\begin{proof} Lihat \cite{RordamLL:2000}, proposisi 4.1.4.  \ns \end{proof}

\begin{prop}\label{0062524} Jika aljabar-$C^*$ $A$ dan $B$ ekuivalen secara homotopi,
maka $K_0(A) \cong K_0(B)$.
\end{prop}

%\begin{proof} Lihat \cite{RordamLL:2000}, proposisi 4.1.4.  \ns \end{proof}

\begin{defn} Misalkan $\pi:=Q$ dan $\lambda:=\psi$ adalah homomorfisme-$*\,$ yang berturut-turut berperan sebagai pemetaan hasil bagi dan penampang kanonik dalam barisan eksak
terbelah~\eqref{0062334i} bagi unitalisasi suatu aljabar-$C^*$~$A$. Definisikan % SOURCE-CORRECTION: CH17-C016
 \index{skalar!pemetaan}%
\df{pemetaan skalar} $s\colon \wt A \sto \wt A$ untuk $\wt A$ dengan $s :=  \lambda \circ \pi$. Setiap
anggota $\wt A$ dapat ditulis dalam bentuk $a + \alpha \vc 1_{\wt A}$ untuk suatu $a \in A$ dan
$\alpha \in \C$. Perhatikan bahwa $s(a + \alpha \vc 1_{\wt A}) = \alpha \vc 1_{\wt A}$ dan bahwa $x
- s(x) \in A$ untuk setiap $x \in \wt A$. Untuk setiap bilangan asli $n$, pemetaan skalar menginduksi
pemetaan yang bersesuaian $s = s_n\colon \M n{\wt A} \sto \M n{\wt A}$. Suatu elemen $x \in \M n{\wt
A}$ merupakan
 \index{skalar!elemen}%
\df{elemen skalar} dari $\M n{\wt A}$ jika $s(x) = x$.
\end{defn}

\begin{prop}[Gambaran Standar $K_0(A)$ untuk sebarang $A$]\label{0062544} Jika $A$ suatu
aljabar-$C^*$,
 \index{standar!gambaran $K_0(A)$!untuk sebarang $A$}%
 \index{gambaran (\seeonly{gambaran standar}))}%
maka
     \[ K_0(A) = \{\,[p\,] - [s(p)]\colon p \in \fml P_\infty(\wt A)\}\,. \] % SOURCE-CORRECTION: CH17-C017
\end{prop}























\section{Sifat Keeksakan dan Stabilitas Funktor $K_0$}
\begin{defn} Funktor kovarian $F$ dari kategori $\cat A$ ke kategori $\cat B$ disebut
 \index{eksak terbelah!funktor}%
 \index{funktor!eksak terbelah}%
\df{eksak terbelah} jika membawa barisan eksak terbelah ke barisan eksak terbelah. Funktor tersebut disebut
 \index{funktor eksak separuh}%
 \index{funktor!eksak separuh}%
\df{eksak separuh} jika setiap kali barisan
   \[ \vc 0 \to A_1 \to^j A_2 \to^k A_3 \to \vc 0 \]
eksak dalam $\cat A$, maka
   \[ F(A_1) \to^{F(j)} F(A_2) \to^{F(k)} F(A_3) \]
eksak dalam~$\cat B$.
\end{defn}

\begin{prop}\label{0062664} Funktor $K_0$ eksak separuh.
\end{prop}

%\begin{proof} Lihat~\cite{RordamLL:2000}, proposisi 4.3.2.  \ns \end{proof}

\begin{prop}\label{0062667} Funktor $K_0$ eksak terbelah.
\end{prop}

%\begin{proof} Lihat~\cite{RordamLL:2000}, proposisi 4.3.3.  \ns \end{proof}

\begin{prop}\label{0062671} Funktor $K_0$ mempertahankan jumlah langsung. Artinya, jika $A$ dan $B$
aljabar-$C^*$, maka $K_0(A \oplus B) = K_0(A) \oplus K_0(B)$.
\end{prop}

\begin{exam}\label{0062674} Jika $A$ aljabar-$C^*$, maka $K_0(\wt A) = K_0(A) \oplus \Z$.
\end{exam}

Meskipun eksak terbelah sekaligus eksak separuh, funktor $K_0$ tidak eksak. Masing-masing dari dua
contoh berikut cukup untuk menunjukkan hal ini.

\begin{exam}\label{0062677} Barisan
   \[ \vc 0 \to \fml C_0\bigl((0,1)\bigr) \to^\iota \fml C\bigl([0,1]\bigr)
                               \to^\psi \C \oplus \C \to \vc 0 \]
dengan $\psi(f) = \bigl(f(0),f(1)\bigr)$ jelas eksak; tetapi $K_0(\psi)$ tidak surjektif.
\end{exam}

\begin{exam}\label{0062679} Jika $H$ ruang Hilbert berdimensi tak hingga, barisan eksak % SOURCE-CORRECTION: CH17-C018
  \[ \vc 0 \to \ofml K(H) \to^\iota \ofml B(H) \to^\pi \ofml Q(H) \to \vc 0 \]
yang terkait dengan aljabar Calkin $\ofml Q(H)$ bersifat eksak, tetapi $K_0(\iota)$ tidak injektif.
(Contoh ini memerlukan fakta yang belum kita turunkan: $K_0(\ofml K(H)) \cong \Z$; untuk itu, lihat~\cite{RordamLL:2000}, Akibat~6.4.2.)
\end{exam}

Berikutnya adalah sifat stabilitas penting dari funktor~$K_0$.
\begin{prop}\label{0062681} Jika $A$ aljabar-$C^*$, maka $K_0(A) \cong K_0(\M nA)$.
\end{prop}

%\begin{proof} Lihat~\cite{RordamLL:2000}, proposisi 4.3.8.   \ns \end{proof}





















\section{Limit Induktif}
Dalam bagian~\ref{sec_ind_limits}, kita mendefinisikan \emph{limit induktif} untuk sebarang kategori, dan dalam
bagian~\ref{sec_LFsps}, kita membatasi perhatian pada limit induktif ketat dari ruang vektor topologis
konveks lokal, yang berguna dalam konteks teori distribusi. Dalam bagian
ini, kita menelaah limit induktif dari barisan aljabar-$C^*$.

\begin{defn} Dalam sebarang kategori, suatu
 \index{induktif!barisan}%
 \index{barisan!induktif}%
\df{barisan induktif} adalah pasangan $(A,\phi)$, dengan $A = (A_j)$ barisan objek dan
$\phi = (\phi_j)$ barisan morfisme sedemikian sehingga $\phi_j\colon A_j \sto A_{j+1}$ untuk
setiap~$j$. Suatu
 \index{induktif!limit}%
 \index{limit!induktif}%
\df{limit induktif} (atau
 \index{langsung!limit}%
 \index{limit!langsung}%
\df{limit langsung}) dari barisan $(A,\phi)$ adalah pasangan $(L,\mu)$, dengan $L$ suatu objek dan
$\mu = (\mu_n)$ barisan morfisme $\mu_j\colon A_j \sto L$ yang memenuhi
 \begin{enumerate}
  \item[(i)] $\mu_j = \mu_{j+1} \circ \phi_j$ untuk setiap $j \in \N$, dan
  \item[(ii)] jika $(M,\lambda)$ adalah pasangan dengan $M$ suatu objek dan $\lambda = (\lambda_j)$
suatu barisan morfisme $\lambda_j\colon A_j \sto M$ yang memenuhi $\lambda_j = \lambda_{j+1}
\circ \phi_j$, maka terdapat morfisme tunggal $\psi\colon L \sto M$ sedemikian sehingga $\lambda_j
= \psi \circ \mu_j$ untuk setiap $j \in \N$.
 \end{enumerate}
Dengan sedikit penyalahgunaan bahasa, biasanya kita mengatakan bahwa $L$ adalah limit induktif dari barisan $(A_j)$ dan
 \index{limit@$\underrightarrow{\lim} A_j$ (limit induktif)}%
menulis $L = \underrightarrow{\lim} A_j$.

\vskip 15 pt

   \[\xymatrix@+15pt@C+25pt{&&&& L\ar@{-->}[dddd]^\psi \\
                   &&&&   \\
                  {\cdots}\ar[r] & A_j \ar[uurrr]^{\mu_j}\ar[ddrrr]^{\lambda_j}\ar[r]^(.6){\phi_j} &
                       A_{j+1}\ar[uurr]^{\mu_{j+1}}\ar[ddrr]^{\lambda_{j+1}}\ar[r] & {\cdots} &  \\
                   &&&&   \\
                   &&&& M
    }\]
\end{defn}

\begin{prop} Limit induktif dari barisan induktif (jika ada dalam suatu kategori) tunggal (hingga isomorfisme).
\end{prop}

\begin{prop} Setiap barisan induktif $(A,\phi)$ dari aljabar-$C^*$
 \index{cstaralg@aljabar-$C^*$!limit induktif}%
memiliki limit induktif. (Demikian pula setiap barisan induktif dari grup Abelian.)
\end{prop}

\begin{proof} Lihat \cite{Blackadar:2006}, II.8.2.1; \cite{RordamLL:2000}, proposisi 6.2.4; atau
\cite{Wegge-Olsen:1993}, lampiran~L.  \ns
\end{proof}

\begin{prop} Jika $(L,\mu)$ merupakan limit induktif dari barisan induktif $(A,\phi)$ dari
aljabar-$C^*$, maka $L = \clo{\bigcup {\mu_n}^\sto(A_n)}$ dan $\norm{\mu_m(a)} = \lim_{n \sto
\infty}\norm{\phi_{n,m}(a)} = \inf_{n \ge m}\norm{\phi_{n,m}(a)}$ untuk semua $m \in \N$ dan $a
\in A$.
\end{prop}

\begin{proof} Lihat \cite{RordamLL:2000}, proposisi 6.2.4.   \ns \end{proof}

\begin{defn} Suatu
 \index{aljabar-$C^*$ berdimensi hingga secara aproksimatif}%
 \index{aljabar-AF}%
 \index{hingga!berdimensi!secara aproksimatif}%
 \index{cstaralg@aljabar-$C^*$!aljabar-AF}%
\df{aljabar-$C^*$ berdimensi hingga secara aproksimatif} (singkatnya aljabar-AF) adalah limit induktif
dari suatu barisan aljabar-$C^*$ berdimensi hingga.
\end{defn}

\begin{exam} Jika $(A_n)$ merupakan barisan meningkat yang terdiri atas subaljabar-$C^*$ dari suatu aljabar-$C^*$ $D$
(dengan $\iota_n\colon A_n \sto A_{n+1}$ pemetaan inklusi untuk setiap~$n$), maka $(A,\iota)$
merupakan barisan induktif aljabar-$C^*$ yang limit induktifnya adalah $(B,j)$, dengan $B =
\clo{\bigcup_{n=1}^\infty A_n}$ dan $j_n\colon A_n \sto B$ pemetaan inklusi untuk setiap~$n$. % SOURCE-CORRECTION: CH17-C019
\end{exam}

\begin{exam} Barisan $\mathbf M_1 = \C \to^{\phi_1} \mathbf M_2 \to^{\phi_2} \mathbf M_3
\to^{\phi_3} \dots$ (dengan $\phi_n(a) = \diag(a,0)$ untuk setiap $n \in \N$ dan $a \in \mathbf
M_n$) merupakan barisan induktif aljabar-$C^*$ yang limit induktifnya adalah aljabar-$C^*$
$\ofml K(H)$ yang terdiri atas operator kompak pada ruang Hilbert terpisahkan berdimensi tak hingga~$H$. % SOURCE-CORRECTION: CH17-C020
\end{exam}

\begin{proof} Lihat \cite{RordamLL:2000}, bagian 6.4.  \ns \end{proof}

\begin{exam} Barisan $\Z \to^{\vc 1} \Z \to^{\vc 2} \Z \to^{\vc 3} \Z \to^{\vc 4} \dots$
(dengan $\vc n \colon \Z \sto \Z$ memenuhi $\vc n(1) = n$) merupakan barisan induktif grup
Abelian yang limit induktifnya adalah himpunan $\Q$ dari bilangan rasional.
\end{exam}

\begin{exam} Barisan $\Z \to^{\vc 2} \Z \to^{\vc 2} \Z \to^{\vc 2} \Z \to^{\vc 2} \dots$
merupakan barisan induktif grup Abelian yang limit induktifnya adalah himpunan bilangan rasional
diadik.
\end{exam}

Hasil berikut disebut
 \index{kekontinuan!dari $K_0$}%
 \index{K0@$K_0$!kekontinuan dari}%
\emph{sifat kekontinuan $K_0$}.

\begin{prop} Jika $(A,\phi)$ barisan induktif aljabar-$C^*$, maka
   \[ K_0\bigl(\,\underrightarrow{\lim}\, A_n\bigr) = \underrightarrow{\lim}\,K_0(A_n)\,. \]
\end{prop}

\begin{proof} Lihat \cite{RordamLL:2000}, teorema 6.3.2.   \ns \end{proof}

\begin{exam} Jika $H$ ruang Hilbert,
 \index{K0@$K_0(A)$!adalah $\Z$ ketika $A = \ofml K(H)$}%
maka $K_0(\ofml K(H)) \cong \Z$.
\end{exam}





















\section{Diagram Bratteli}
\begin{prop} Homomorfisme-$*\,$ tak nol dari $M_k$ ke $M_n$ ada hanya jika $n \ge k$; dalam
hal ini, semuanya tepat merupakan pemetaan berbentuk
   \[ \phi\colon a \mapsto u\,\diag(a, a, \dots, a, \vc 0)\, u^* \] % SOURCE-CORRECTION: CH17-C021
dengan $u$ matriks uniter. Di sini terdapat $m$ salinan $a$, dan $\vc 0$ adalah matriks nol berukuran $r \times
r$, dengan $n = mk + r$. Bilangan $m$ adalah
 \index{multiplisitas!dari homomorfisme-$*\,$}%
\df{multiplisitas} dari~$\phi$.
\end{prop}

\begin{proof} Lihat \cite{Power:1992}, akibat 1.3.  \ns \end{proof}

\begin{exam} Contoh homomorfisme-$*\,$ dari $M_2$ ke $M_7$ adalah
 \[  A =
  \begin{bmatrix} a & b \\ c & d  \end{bmatrix} \mapsto
    \begin{bmatrix}
      a & b & 0 & 0 & 0 & 0 & 0 \\
      c & d & 0 & 0 & 0 & 0 & 0 \\
      0 & 0 & a & b & 0 & 0 & 0 \\
      0 & 0 & c & d & 0 & 0 & 0 \\
      0 & 0 & 0 & 0 & 0 & 0 & 0 \\
      0 & 0 & 0 & 0 & 0 & 0 & 0 \\
      0 & 0 & 0 & 0 & 0 & 0 & 0
    \end{bmatrix} = \diag(A,A, \vc 0)\]
 dengan $m = 2$ dan $r = 3$.
\end{exam}

\enlargethispage{\baselineskip}

\begin{prop}\label{0068117} Setiap aljabar-$C^*$ berdimensi hingga $A$ isomorfik dengan
jumlah langsung aljabar matriks
   \[ A \simeq M_{k_1} \oplus \dots \oplus M_{k_r}.\]
Misalkan $A$ dan $B$ aljabar-$C^*$ berdimensi hingga, sehingga
   \[ A \simeq M_{k_1} \oplus \dots \oplus M_{k_r} \qquad \text{dan}
                          \qquad B \simeq M_{n_1} \oplus \dots \oplus M_{n_s} \]
dan misalkan $\phi\colon A \sto B$ homomorfisme-$*\,$ beridentitas. Maka $\phi$
ditentukan (hingga ekuivalensi uniter dalam~$B$) oleh matriks $s \times r$ $\vc m = \bigl[
m_{ij}\bigr]$ yang berentri bilangan bulat tak negatif sedemikian sehingga % SOURCE-CORRECTION: CH17-C022
   \begin{equation}\label{0068117i}
    \vc m\, \vc k = \vc n.
   \end{equation}
\emph{Di sini $\vc k = (k_1, \dots,k_r)$, $\vc n = (n_1, \dots, n_s)$, dan bilangan $m_{ij}$ adalah
multiplisitas pemetaan}
 \[\xymatrix@1{M_{k_j} \ar[r] & A \ar[r]^\phi & B \ar[r] & M_{n_i}
 }\]
\end{prop}

\begin{proof} Lihat \cite{Davidson:1996}, teorema III.1.1.  \ns \end{proof}

\begin{defn} Misalkan $\phi$ seperti dalam proposisi~\ref{0068117} sebelumnya. Suatu
 \index{diagram Bratteli}%
 \index{diagram!Bratteli}%
\df{diagram Bratteli} untuk $\phi$ terdiri atas dua baris (atau kolom) simpul yang diberi label
$k_j$ dan $n_i$, bersama $m_{ij}$ sisi yang menghubungkan $k_j$ ke $n_i$ (untuk $1 \le j \le
r$ dan $1 \le i \le s$).
\end{defn}

\begin{exam} Misalkan $\phi\colon \C \oplus \C \sto M_4 \oplus M_3$ diberikan oleh
 \[ (\lambda, \mu) \mapsto \left(
      \begin{bmatrix}
        \lambda & 0 & 0 & 0 \\
        0 & \lambda & 0 & 0 \\
        0 & 0 & \lambda & 0 \\
        0 & 0 & 0 & \mu
      \end{bmatrix}\, , \,
      \begin{bmatrix}
        \lambda & 0 & 0 \\
        0 & \lambda & 0 \\
        0 & 0 & \mu
      \end{bmatrix}\right)\, .\]
Maka $m_{11} = 3$, $m_{12} = 1$, $m_{21} = 2$, dan $m_{22} = 1$. Perhatikan bahwa
   \[\vc m \,\vc k = \begin{bmatrix} 3 & 1 \\ 2 & 1 \end{bmatrix}\,
             \begin{bmatrix} 1 \\ 1 \end{bmatrix} =
             \begin{bmatrix} 4 \\ 3 \end{bmatrix} = \vc n.\]
Suatu diagram Bratteli untuk $\phi$ adalah
    \[\xymatrix@=90pt{1 \ar@<.7ex>[r]\ar[r]\ar@<-.7ex>[r]\ar@<.5ex>[dr]\ar@<-.5ex>[dr] & 4 \\
                1 \ar[ur]\ar[r] & 3
     }\]
\end{exam}

\begin{exam} Misalkan $\phi\colon \C \oplus M_2 \sto M_3 \oplus M_5 \oplus M_2$
diberikan oleh
 \[ (\lambda, b) \mapsto \left(
    \begin{bmatrix}
        \lambda & 0  \\
            0 &  b
      \end{bmatrix}\, , \,
      \begin{bmatrix}
        \lambda & 0 & 0  \\
              0 & b & 0  \\
              0 & 0 &  b
      \end{bmatrix} \,,\, b \, \right).
    \]
Maka $m_{11} = 1$, $m_{12} = 1$, $m_{21} = 1$, $m_{22} = 2$, $m_{31} = 0$, dan $m_{32} = 1$.
Perhatikan bahwa
 \[\vc m \,\vc k = \begin{bmatrix} 1 & 1 \\ 1 & 2 \\ 0 & 1 \end{bmatrix}\,
             \begin{bmatrix} 1 \\ 2 \end{bmatrix} =
             \begin{bmatrix} 3 \\ 5 \\2 \end{bmatrix} = \vc n.\]
Suatu diagram Bratteli untuk $\phi$ adalah
 \[\xymatrix@=20pt@C+30pt{&& 3                    \\
                   1 \ar[urr]\ar[drr] &&   \\
                   && 5                    \\
         2\ar[uuurr]\ar@<+.5ex>[urr]\ar@<-.5ex>[urr]\ar[drr] && \\
                   && 2
   }\]
\end{exam}

\begin{exam} Misalkan $\phi\colon M_3 \oplus M_2 \oplus M_2 \sto M_9 \oplus M_7$
diberikan oleh

\vskip 10 pt

 \[ (a,b,c) \mapsto \left(
    \begin{bmatrix}
            a & \vc 0 & \vc 0 \\
        \vc 0 &     a & \vc 0 \\
        \vc 0 & \vc 0 & \vc 0
      \end{bmatrix}\, , \,
      \begin{bmatrix}
            a & \vc 0 & \vc 0  \\
        \vc 0 &     b & \vc 0  \\
        \vc 0 & \vc 0 &     c
      \end{bmatrix}\, \right)\, .\]
Maka $m_{11} = 2$, $m_{12} = 0$, $m_{13} = 0$, $m_{21} = 1$, $m_{22} = 1$, dan $m_{23} = 1$.
Perhatikan bahwa kali ini terjadi sesuatu yang ganjil: kita mempunyai $\vc m\,\vc k \ne \vc n$:
 \[\vc m \,\vc k = \begin{bmatrix} 2 & 0 & 0 \\ 1 & 1 & 1 \end{bmatrix}\,
             \begin{bmatrix} 3 \\ 2 \\ 2 \end{bmatrix} =
             \begin{bmatrix} 6 \\ 7 \end{bmatrix} \le
             \begin{bmatrix} 9 \\ 7 \end{bmatrix} = \vc n.\]
Apa masalahnya di sini? Nah, hasil yang dinyatakan sebelumnya berlaku untuk homomorfisme-$*\,$
\emph{beridentitas}, sedangkan $\phi$ ini \emph{tidak} beridentitas. Secara umum, inilah hasil terbaik
yang dapat kita harapkan: $\vc m\,\vc k \le \vc n$.

\noindent Berikut diagram Bratteli yang dihasilkan untuk $\phi$:
   \[\xymatrix@=+20pt@C+20pt@R-7pt{
        3\ar@<+.5ex>[drr]\ar@<-.5ex>[drr]\ar[dddrr] && \\
                            && 9                   \\
                           2\ar[drr] &&            \\
                            && 7                   \\
                           2\ar[urr] &&
   }\]
\end{exam}

\begin{exam} Jika $K$ himpunan Cantor, maka $\fml C(K)$ merupakan aljabar-AF.
 \index{AF@aljabar-AF!$\fml C(K)$ merupakan aljabar-AF ($K$ himpunan Cantor)}%
Untuk melihatnya, tuliskan $K$ sebagai irisan dari suatu keluarga menurun yang terdiri atas himpunan bagian tertutup $K_j$, yang masing-masing
terdiri atas $2^j$ subinterval tertutup yang saling lepas dari $[0,1]$. Untuk setiap $j \ge 0$, misalkan
$A_j$ subaljabar yang terdiri atas fungsi-fungsi dalam $\fml C(K)$ yang konstan pada setiap
interval penyusun~$K_j$. Jadi $A_j \simeq \C^{2^j}$ untuk setiap~$j \ge 0$. Pembenaman
 \[\phi_j \colon A_j \sto A_{j+1}\colon
      (a_1, a_2, \dots, a_{2^j}) \mapsto
        (a_1, a_1, a_2, a_2,\dots, a_{2^j}, a_{2^j})\]
memecah setiap proyeksi minimal menjadi jumlah dua proyeksi minimal. Untuk $\phi_0$, matriks
$\vc m_0$ dari ``multiplisitas parsial'' yang bersesuaian adalah $\begin{bmatrix} 1 \\ 1
\end{bmatrix}$; matriks $\vc m_1$ yang bersesuaian dengan $\phi_1$ adalah
$\begin{bmatrix} 1 & 0\\ 1 & 0 \\ 0 & 1 \\ 0 & 1 \end{bmatrix}$; dan seterusnya. Jadi diagram Bratteli
untuk limit induktif $\fml C(K) = \underrightarrow{\lim}A_j$ adalah:

\vskip 20 pt

 \[\xymatrix@=+7pt@C+20pt@R-7pt{
     &&& 1 & \\
     && 1\ar[ur]\ar[dr] && \\
     &&& 1 & \\
     & 1\ar[uur]\ar[ddr] &&& \\
     &&& 1 & \\
     && 1\ar[ur]\ar[dr] && \\
     &&& 1 & \\
     1\ar[uuuur]\ar[ddddr] &&&& \cdots \\
     &&& 1 & \\
     && 1\ar[ur]\ar[dr] && \\
     &&& 1 & \\
     & 1\ar[uur]\ar[ddr] &&& \\
     &&&1 & \\
     && 1\ar[ur]\ar[dr] && \\
     &&& 1 &
   }\]
\end{exam}

\begin{exam} Berikut suatu contoh dari apa yang disebut
 \index{aljabar CAR}%
 \index{cstaralg@aljabar-$C^*$!aljabar CAR}%
aljabar CAR (CAR = Relasi Antikomutasi Kanonik). Untuk $j \ge 0$, misalkan $A_j = % SOURCE-CORRECTION: CH17-C023
M_{\spsp{2}j}$ dan
 \[\phi_j\colon A_j \to A_{j+1}\colon \vc a \mapsto
      \begin{bmatrix} \vc a & 0 \\ 0 & \vc a \end{bmatrix}.\]
``Matriks'' multiplisitas $\vc m$ untuk setiap $\phi_j$ hanyalah matriks $1 \times 1$~$[2]$.
Kita melihat bahwa untuk setiap $j$
 \[\vc m\, \vc k(j) = [2]\,[2^j] = [2^{j+1}] = \vc n(j).\]
Ini menghasilkan diagram Bratteli berikut

\vskip 15pt

 \[\xymatrix{1 \ar@<.5ex>[r]\ar@<-.5ex>[r] & 2 \ar@<.5ex>[r]\ar@<-.5ex>[r] & 4
        \ar@<.5ex>[r]\ar@<-.5ex>[r] & 8 \ar@<.5ex>[r]\ar@<-.5ex>[r] & 16
        \ar@<.5ex>[r]\ar@<-.5ex>[r] & 32
        \ar@<.5ex>[r]\ar@<-.5ex>[r]& {\cdots}
   }\]
\vskip 15pt
\noindent untuk limit induktif $C = \underrightarrow{\lim}A_j$.

\textbf{Namun}, jangkauan $\phi_j$ termuat dalam subaljabar $B_{j+1} \simeq M_{2^j}
\oplus M_{2^j}$ dari~$A_{j+1}$. (Ambil $B_0 = A_0 = \C$.) Jadi, untuk $j \in \N$, kita dapat memandang
$\phi_j$ sebagai pemetaan dari $B_j$ ke~$B_{j+1}$:
 \[ \phi_j\colon(\vc b, \vc c) \mapsto
    \left(\begin{bmatrix} \vc b & 0 \\ 0 & \vc c \end{bmatrix},
          \begin{bmatrix} \vc b & 0 \\ 0 & \vc c \end{bmatrix}
              \right)\,.\]
Sekarang matriks multiplisitas $\vc m$ untuk setiap $\phi_j$ adalah $\begin{bmatrix} 1 & 1 \\ 1 & 1
\end{bmatrix}$, dan kita melihat bahwa untuk setiap $j$
 \[\vc m\, \vc k(j) =
   \begin{bmatrix} 1 & 1 \\ 1 & 1 \end{bmatrix}\,
        \begin{bmatrix} 2^{j-1} \\ 2^{j-1} \end{bmatrix} =
   \begin{bmatrix} 2^j \\ 2^j \end{bmatrix} = \vc n(j).\]
Karena $C = \underrightarrow{\lim}A_j = \underrightarrow{\lim}B_j$, kita memperoleh diagram Bratteli kedua
(yang cukup berbeda) untuk aljabar-AF~$C$.
   \[\xymatrix@-10pt{&& 1 \ar[rr]\ar[rrdd] && 2 \ar[rr]\ar[rrdd] && 4 \ar[rr]\ar[rrdd]
                                && 8 \ar[rr]\ar[rrdd] && 16 \ar[rr]\ar[rrdd]&& \cdots\\
              1 \ar[urr]\ar[drr]&&&&&&&&&&  \\
             && 1 \ar[rr]\ar[rruu] && 2 \ar[rr]\ar[rruu] && 4 \ar[rr]\ar[rruu]
                                && 8 \ar[rr]\ar[rruu] && 16 \ar[rr]\ar[rruu]&& \cdots \\
   }\]
\end{exam}

\begin{exam} Ini adalah
 \index{aljabar Fibonacci}%
 \index{cstaralg@aljabar-$C^*$!aljabar Fibonacci}%
\df{aljabar Fibonacci}. Untuk $j \in \N$, definisikan barisan $\bigl(p_j\bigr)$ dan
$\bigl(q_j\bigr)$ dengan relasi rekursi yang dikenal:
 \begin{align*}
       p_1 = &q_1 = 1, \\
       p_{j+1} = &p_j + q_j, \text{ dan } \\
       q_{j+1} &= p_j.
 \end{align*}
Untuk semua $j \in \N$, misalkan $A_j = M_{p_j} \oplus M_{q_j}$ dan

\vskip 10 pt

 \[\phi_j\colon A_j \sto A_{j+1}\colon (a,b) \mapsto
     \left(\begin{bmatrix} a & 0 \\ 0 & b
          \end{bmatrix},  a \right). \]
Matriks multiplisitas $\vc m$ untuk setiap $\phi_j$ adalah $\begin{bmatrix} 1 & 1 \\ 1 & 0
\end{bmatrix}$, dan untuk setiap $j$
 \[\vc m\, \vc k(j) =
   \begin{bmatrix} 1 & 1 \\ 1 & 0 \end{bmatrix}\,
        \begin{bmatrix} p_j \\ q_j \end{bmatrix} =
   \begin{bmatrix} p_j + q_j \\ p_j \end{bmatrix} =
        \begin{bmatrix} p_{j+1} \\ q_{j+1} \end{bmatrix} =
                          \vc n(j).\]
Diagram Bratteli yang dihasilkan untuk $F = \underrightarrow{\lim}A_j$ adalah
   \[\xymatrix@-10pt{&& 1 \ar[rr]\ar[rrdd] && 2 \ar[rr]\ar[rrdd] &&
         3 \ar[rr]\ar[rrdd] && 5 \ar[rr]\ar[rrdd] && 8 \ar[rr]\ar[rrdd]&& \cdots\\
              1 \ar[urr]\ar[drr]&&&&&&&&&&  \\
             && 1 \ar[rruu] && 1 \ar[rruu] && 2 \ar[rruu] &&
              3 \ar[rruu] && 5 \ar[rruu]&& \cdots \\
   }\]
\end{exam}













\endinput



%\enlargethispage{\baselineskip}






\part{PENDAMPING PENGUASAAN DAN JEMBATAN SPEKTRAL}

% O008 compact-spectral/SVD bridge -- separately authored component.
% Model provenance: OpenAI Codex gpt-5.6-sol, Ultra, atas arahan pengguna.
% Component license: CC BY-SA 4.0.
% Not authored or endorsed by John M. Erdman or Portland State University.

\chapter{JEMBATAN SPEKTRAL-KOMPAK DAN NILAI SINGULAR}
\label{o008bridge:compact-spectral-svd}

\begin{center}
\small
Komponen asli terpisah, CC BY-SA 4.0. Ditulis dengan bantuan OpenAI Codex
gpt-5.6-sol, Ultra, atas arahan pengguna. Komponen ini bukan tulisan John M.
Erdman dan tidak menyiratkan dukungan beliau atau Portland State University.
\end{center}

Bab~\ref{chap_cpt_ops} telah mengembangkan operator kompak dan dekomposisi
polar, sedangkan Bab~\ref{spectrum} telah mengembangkan spektrum. Bab tentang
teori Fredholm kemudian membuktikan alternatif Fredholm bagi operator
Riesz--Schauder, sifat jangkauan tertutup, serta teori indeks. Tujuan jembatan
ini bukan mengulang teori Fredholm tersebut, melainkan menghubungkan hasilnya
ke geometri eigenvektor operator kompak dan selanjutnya ke dekomposisi nilai
singular.

Semua ruang dalam bab ini merupakan ruang kompleks. Simbol $H_1$ dan $H_2$
menyatakan ruang Hilbert, sedangkan $B$ menyatakan ruang Banach. Jika
$u\in H_2$ dan $v\in H_1$, kita memakai notasi peringkat-satu
\[
  u\otimes v\colon H_1\longrightarrow H_2,
  \qquad
  (u\otimes v)x=\langle x,v\rangle u.
\]

\section{Konsekuensi spektral Riesz--Schauder}

% O008-BRIDGE-ID: O008-BRIDGE-CS-DEF-001
\begin{defn}\label{o008bridge:spectral-multiplicity}
Jika $T\in\ofml B(B)$ dan $\lambda\in\C$, maka
\df{ruang eigen} untuk~$\lambda$ adalah $\ker(T-\lambda I)$. Dimensinya
disebut \df{multiplikitas geometrik} nilai eigen~$\lambda$.
\end{defn}

% O008-BRIDGE-ID: O008-BRIDGE-CS-THM-001
\begin{thm}[Konsekuensi spektral Riesz--Schauder]
\label{o008bridge:riesz-schauder-spectrum}
Misalkan $B$ ruang Banach kompleks berdimensi tak hingga dan
$K\in\ofml K(B)$. Maka:
\begin{enumerate}[label=\rm{(\alph*)}]
  \item setiap $\lambda\in\sigma(K)\setminus\{0\}$ merupakan nilai eigen;
  \item setiap ruang eigen untuk $\lambda\neq0$ berdimensi hingga;
  \item untuk setiap $r>0$, hanya ada berhingga banyak nilai eigen berbeda
        yang memenuhi $|\lambda|\geq r$;
  \item $\sigma(K)\setminus\{0\}$ paling banyak terhitung, setiap titiknya
        terisolasi dalam spektrum, dan satu-satunya titik akumulasi yang
        mungkin adalah~$0$;
  \item $0\in\sigma(K)$.
\end{enumerate}
\end{thm}

\begin{proof}
Ambil $\lambda\neq0$. Operator
$\lambda I-K$ adalah operator Riesz--Schauder dalam pengertian
Bab~15. Alternatif Fredholm~\ref{0040343} menyatakan bahwa operator itu
injektif jika dan hanya jika surjektif. Jika injektif, ia bijektif dan teorema
pemetaan terbuka memberi invers terbatas. Jadi, bila
$\lambda\in\sigma(K)$, operator $\lambda I-K$ tidak injektif; dengan kata
lain, $\lambda$ adalah nilai eigen. Alternatif Fredholm yang sama juga memberi
\[
 \dim\ker(\lambda I-K)<\infty,
\]
yang membuktikan (a) dan~(b) tanpa mengulang argumen Fredholm.

Untuk (c), andaikan terdapat nilai eigen berbeda
$\lambda_1,\lambda_2,\ldots$ dengan $|\lambda_n|\geq r>0$, dan pilih
eigenvektor $x_n\neq0$. Eigenvektor untuk nilai eigen berbeda bebas linear.
Tuliskan $M_n=\spn\{x_1,\ldots,x_n\}$. Menurut lema Riesz, kita dapat memilih
$y_n\in M_n$ dengan $\|y_n\|=1$ dan
\[
  \dist(y_n,M_{n-1})>\frac12.
\]
Karena $(K-\lambda_n I)M_n\subseteq M_{n-1}$, berlaku
$y_n-\lambda_n^{-1}Ky_n\in M_{n-1}$. Untuk $m<n$ juga
$\lambda_m^{-1}Ky_m\in M_m\subseteq M_{n-1}$. Akibatnya
\[
 \left\|\lambda_n^{-1}Ky_n-\lambda_m^{-1}Ky_m\right\|>\frac12.
\]
Namun barisan $(\lambda_n^{-1}y_n)$ terbatas oleh~$1/r$, sehingga
kekompakan~$K$ mengharuskan $(\lambda_n^{-1}Ky_n)$ mempunyai subbarisan
konvergen. Ini bertentangan dengan pemisahan di atas. Jadi (c) benar.

Terakhir,
\[
 \sigma(K)\setminus\{0\}
 =\bigcup_{m=1}^{\infty}
   \{\lambda\in\sigma(K):|\lambda|\geq1/m\}
\]
adalah gabungan terhitung himpunan-himpunan berhingga. Pernyataan isolasi dan
akumulasi mengikuti langsung dari (c), sehingga (d) terbukti. Jika
$0\notin\sigma(K)$, maka $K$ invertibel dan
$I=K^{-1}K$ kompak. Hal ini mustahil pada ruang Banach berdimensi tak hingga:
lema Riesz menghasilkan barisan vektor satuan yang tidak mempunyai
subbarisan konvergen. Jadi operator identitas tidak kompak, dan (e) terbukti.
\end{proof}

% O008-BRIDGE-ID: O008-BRIDGE-CS-REM-001
\begin{rem}\label{o008bridge:no-fredholm-duplication}
Teorema~\ref{o008bridge:riesz-schauder-spectrum} memakai alternatif Fredholm
yang sudah terdapat dalam Bab~15 hanya pada dua titik: kesetaraan
injektif--surjektif dan keberhinggaan dimensi kernel. Jembatan ini tidak
mengulang pembentukan aljabar Calkin, teorema Atkinson, indeks Fredholm, atau
klasifikasi komponen lintasan operator Fredholm.
\end{rem}

% O008-BRIDGE-ID: O008-BRIDGE-CS-EXAM-001
\begin{exam}\label{o008bridge:diagonal-compact-example}
Pada $\ell^2$, tetapkan
\[
 K(x_1,x_2,\ldots)=
   \left(x_1,\frac{x_2}{2},\frac{x_3}{3},\ldots\right).
\]
Operator ini kompak dan swaadjoin. Nilai eigennya adalah $1/n$, masing-masing
bermultiplikitas satu, dan
$\sigma(K)=\{0\}\cup\{1/n:n\in\N\}$. Contoh ini memperlihatkan bahwa nol
memang dapat menjadi titik akumulasi dan bahwa jangkauan operator kompak tidak
harus tertutup.
\end{exam}

\section{Teorema spektral kompak swaadjoin}

% O008-BRIDGE-ID: O008-BRIDGE-CS-LEM-001
\begin{lem}\label{o008bridge:max-eigenvalue}
Jika $T\in\ofml K(H)$ swaadjoin dan $T\neq0$, maka $T$ mempunyai nilai eigen
real $\lambda$ dengan $|\lambda|=\|T\|$.
\end{lem}

\begin{proof}
Untuk operator swaadjoin, hasil tentang jangkauan numerik di Bab~5 memberi
\[
 \|T\|=\sup_{\|x\|=1}|\langle Tx,x\rangle|.
\]
Pilih vektor satuan $x_n$ sehingga nilai mutlak di ruas kanan menuju
$\|T\|$. Karena bentuk kuadratik swaadjoin bernilai real, setelah mengambil
subbarisan terdapat $\alpha\in\{\|T\|,-\|T\|\}$ dengan
$\langle Tx_n,x_n\rangle\to\alpha$. Untuk $\alpha=\|T\|$, misalnya,
\begin{align*}
 \|(T-\alpha I)x_n\|^2
 &=\|Tx_n\|^2-2\alpha\langle Tx_n,x_n\rangle+\alpha^2\\
 &\leq2\alpha\bigl(\alpha-\langle Tx_n,x_n\rangle\bigr)
 \longrightarrow0.
\end{align*}
Kasus $\alpha=-\|T\|$ sama dengan memakai $T-\alpha I=T+\|T\|I$.
Kekompakan~$T$ memberi subbarisan $(Tx_n)$ yang konvergen. Karena
$(T-\alpha I)x_n\to0$ dan $\alpha\neq0$, subbarisan $(x_n)$ sendiri
konvergen ke suatu vektor satuan~$x$. Dengan mengambil limit diperoleh
$Tx=\alpha x$. Jadi $\lambda=\alpha$ adalah nilai eigen real dan
$|\lambda|=\|T\|$.
\end{proof}

% O008-BRIDGE-ID: O008-BRIDGE-CS-LEM-002
\begin{lem}\label{o008bridge:orthogonal-eigenspaces}
Jika $T$ swaadjoin, maka ruang eigen untuk dua nilai eigen berbeda saling
ortogonal. Selain itu, komplemen ortogonal setiap ruang eigen invarian
terhadap~$T$.
\end{lem}

\begin{proof}
Misalkan $Tx=\lambda x$ dan $Ty=\mu y$. Nilai $\lambda$ dan $\mu$ real, dan
\[
 \lambda\langle x,y\rangle
 =\langle Tx,y\rangle
 =\langle x,Ty\rangle
 =\mu\langle x,y\rangle.
\]
Jika $\lambda\neq\mu$, maka $\langle x,y\rangle=0$. Jika $z$ tegak lurus
ruang eigen $E_\lambda$, maka untuk setiap $x\in E_\lambda$,
\[
 \langle Tz,x\rangle=\langle z,Tx\rangle
 =\lambda\langle z,x\rangle=0,
\]
sehingga $Tz\in E_\lambda^\perp$.
\end{proof}

% O008-BRIDGE-ID: O008-BRIDGE-CS-THM-002
\begin{thm}[Teorema spektral untuk operator kompak swaadjoin]
\label{o008bridge:compact-selfadjoint-spectral-theorem}
Misalkan $T\in\ofml K(H)$ swaadjoin. Terdapat himpunan ortonormal
$(e_n)$ yang terdiri atas eigenvektor untuk nilai eigen tak nol
$(\lambda_n)$ sedemikian sehingga
\begin{enumerate}[label=\rm{(\alph*)}]
  \item setiap nilai eigen tak nol diulang menurut multiplikitasnya yang
        hingga;
  \item bila daftarnya tak hingga, $|\lambda_n|\to0$;
  \item $H=\ker T\oplus\clo{\spn\{e_n:n\geq1\}}$;
  \item untuk setiap $x\in H$,
        \[
          Tx=\sum_n\lambda_n\langle x,e_n\rangle e_n;
        \]
  \item jika daftar tak hingga, tetapkan
        $T_N=\sum_{n=1}^N\lambda_n(e_n\otimes e_n)$; jika daftar berhingga
        memiliki panjang~$m$, tetapkan
        $T_N=\sum_{n=1}^{\min\{N,m\}}\lambda_n(e_n\otimes e_n)$.
        Dengan urutan $|\lambda_1|\geq|\lambda_2|\geq\cdots$, berlaku
        $\|T-T_N\|=\sup_{n>N}|\lambda_n|\to0$, dengan supremum himpunan
        indeks kosong didefinisikan sebagai~$0$.
\end{enumerate}
Dengan menambahkan sembarang basis ortonormal bagi $\ker T$, kita memperoleh
basis ortonormal seluruh~$H$ yang terdiri atas eigenvektor~$T$.
\end{thm}

\begin{proof}
Jika $T=0$, ambil daftar kosong dan hasilnya langsung benar. Jika
$\dim H<\infty$, teorema spektral kompleks berdimensi hingga dari
Bab~1, Teorema~\ref{00017}, memberikan basis ortonormal eigenvektor; setelah
eigenvektor bernilai eigen nol dipindahkan ke $\ker T$, semua butir di atas
mengikuti, termasuk konvensi ekor kosong pada~(e). Karena itu, selanjutnya
andaikan $\dim H=\infty$ dan $T\neq0$.
Teorema~\ref{o008bridge:riesz-schauder-spectrum} mengatakan bahwa himpunan
nilai eigennya yang tak nol paling banyak terhitung dan setiap ruang eigennya
berdimensi hingga. Pilih satu basis ortonormal dari setiap ruang eigen tak nol.
Menurut Lemma~\ref{o008bridge:orthogonal-eigenspaces}, gabungan semua basis
tersebut ortonormal, sehingga dapat disusun sebagai barisan $(e_n)$ dengan
nilai eigen terkait~$(\lambda_n)$. Teorema Riesz--Schauder yang sama memberi
$|\lambda_n|\to0$ bila barisan itu tak hingga.

Misalkan $M=\clo{\spn\{e_n:n\geq1\}}$. Komplemen $M^\perp$ invarian. Jika
$T|_{M^\perp}$ tidak nol, Lemma~\ref{o008bridge:max-eigenvalue} akan memberi
eigenvektor tak nol lain yang tegak lurus semua~$e_n$, bertentangan dengan
pemilihan seluruh ruang eigen tak nol. Jadi $T|_{M^\perp}=0$ dan
$M^\perp\subseteq\ker T$. Sebaliknya, setiap vektor dalam $\ker T$ tegak
lurus semua ruang eigen dengan nilai eigen tak nol, sehingga
$\ker T=M^\perp$. Ini membuktikan (c).

Uraikan $x=x_0+\sum_n\langle x,e_n\rangle e_n$ dengan
$x_0\in\ker T$. Kontinuitas $T$ memberi deret pada (d). Untuk ekor operator,
ortonormalitas memberi
\[
 \|(T-T_N)x\|^2
 =\sum_{n>N}|\lambda_n|^2|\langle x,e_n\rangle|^2
 \leq\left(\sup_{n>N}|\lambda_n|^2\right)\|x\|^2.
\]
Jadi $\|T-T_N\|\leq\sup_{n>N}|\lambda_n|$. Arah sebaliknya diperoleh dengan
menerapkan $T-T_N$ pada setiap $e_n$ dengan $n>N$ dan mengambil supremum.
Maka (e) benar dan $T$ merupakan limit norma operator-operator spektral
berperingkat hingga tersebut.
\end{proof}

% O008-BRIDGE-ID: O008-BRIDGE-CS-COR-001
\begin{cor}\label{o008bridge:positive-compact-spectral}
Jika $T$ juga positif dan $T\neq0$, maka semua $\lambda_n>0$ dan
\[
 \|T\|=\lambda_1,
 \qquad
 Tx=\sum_n\lambda_n\langle x,e_n\rangle e_n.
\]
\end{cor}

\begin{proof}
Untuk eigenvektor satuan~$e_n$,
$\lambda_n=\langle Te_n,e_n\rangle\geq0$. Karena daftar hanya memuat nilai
eigen tak nol, semuanya positif. Pernyataan norma mengikuti pengurutan dalam
teorema.
\end{proof}

\section{Nilai singular dan SVD kompak}

% O008-BRIDGE-ID: O008-BRIDGE-CS-DEF-002
\begin{defn}\label{o008bridge:absolute-value-singular-values}
Untuk $T\in\ofml K(H_1,H_2)$, tetapkan
\[
 |T|=(T^*T)^{1/2}.
\]
Nilai eigen positif $s_n$ dari $|T|$, diulang menurut multiplikitas dan
diurutkan tak-menaik, disebut \df{nilai singular}~$T$. Secara ekuivalen,
$s_n^2$ adalah nilai eigen positif dari~$T^*T$.
\end{defn}

Operator $T^*T$ kompak, positif, dan swaadjoin. Karena itu
Teorema~\ref{o008bridge:compact-selfadjoint-spectral-theorem} berlaku pada
$T^*T$. Akar positifnya~$|T|$ juga kompak: melalui kalkulus fungsional,
fungsi akar pada spektrum dapat didekati seragam oleh polinom yang bernilai
nol di~$0$, sehingga $|T|$ merupakan limit norma operator-operator kompak.
Jadi teorema tersebut berlaku pula pada~$|T|$. Kernel ketiga operator terkait oleh
\[
 \ker|T|=\ker(T^*T)=\ker T,
\]
sebab $\langle T^*Tx,x\rangle=\|Tx\|^2$.

Dekomposisi nilai singular untuk operator kompak memperoleh bentuk berikut.

% O008-BRIDGE-ID: O008-BRIDGE-CS-THM-003
\begin{thm}[SVD untuk operator kompak]
\label{o008bridge:compact-svd}
Misalkan $T\in\ofml K(H_1,H_2)$. Terdapat daftar hingga (mungkin kosong) atau
barisan nilai singular $s_1\geq s_2\geq\cdots>0$ dan keluarga ortonormal
$(v_n)$ dalam~$H_1$ serta $(u_n)$ dalam~$H_2$ sedemikian sehingga
\begin{align*}
  T v_n&=s_nu_n, & T^*u_n&=s_nv_n,\\
  Tx&=\sum_n s_n\langle x,v_n\rangle u_n,
      &&x\in H_1.
\end{align*}
Lebih tepat,
\begin{enumerate}[label=\rm{(\alph*)}]
  \item $(v_n)$ merupakan basis ortonormal
        $\clo{\ran T^*}=(\ker T)^\perp$;
  \item $(u_n)$ merupakan basis ortonormal
        $\clo{\ran T}=(\ker T^*)^\perp$;
  \item dalam norma operator,
        \[
          T=\sum_n s_n(u_n\otimes v_n);
        \]
  \item untuk daftar tak hingga tetapkan
        $T_N=\sum_{n=1}^Ns_n(u_n\otimes v_n)$, sedangkan untuk daftar
        berhingga sepanjang~$m$ tetapkan
        $T_N=\sum_{n=1}^{\min\{N,m\}}s_n(u_n\otimes v_n)$; maka
        \[
          \|T-T_N\|=s_{N+1}\longrightarrow0,
        \]
        dengan konvensi $s_{N+1}=0$ setelah daftar berhingga berakhir,
        termasuk untuk daftar kosong ketika $T=0$.
\end{enumerate}
\end{thm}

\begin{proof}
Terapkan teorema spektral kompak swaadjoin pada operator positif~$T^*T$.
Pilih eigenvektor ortonormal $v_n$ dan nilai eigen
$s_n^2>0$ sehingga
\[
 T^*Tv_n=s_n^2v_n.
\]
Tetapkan $u_n=s_n^{-1}Tv_n$. Maka
\[
 \langle u_n,u_m\rangle
 =\frac{1}{s_ns_m}\langle T^*Tv_n,v_m\rangle
 =\delta_{nm},
\]
jadi $(u_n)$ ortonormal, dan
$T^*u_n=s_n^{-1}T^*Tv_n=s_nv_n$.

Teorema spektral memberi bahwa $(v_n)$ merentang
$(\ker(T^*T))^\perp=(\ker T)^\perp$. Karena $T$ nol pada $\ker T$, ekspansi
Fourier setiap $x\in H_1$ menghasilkan
\[
 Tx=\sum_n\langle x,v_n\rangle Tv_n
   =\sum_n s_n\langle x,v_n\rangle u_n.
\]
Jadi $\clo{\ran T}$ termuat dalam rentang tertutup $(u_n)$, sedangkan setiap
$u_n=s_n^{-1}Tv_n$ termasuk dalam $\ran T$; ini membuktikan (a) dan~(b).

Untuk setiap vektor satuan~$x$,
\[
 \|(T-T_N)x\|^2
 =\sum_{n>N}s_n^2|\langle x,v_n\rangle|^2
 \leq s_{N+1}^2.
\]
Kesamaan dicapai pada $x=v_{N+1}$ bila suku itu ada. Maka
$\|T-T_N\|=s_{N+1}\to0$, yang sekaligus membuktikan (c) dan~(d).
\end{proof}

% O008-BRIDGE-ID: O008-BRIDGE-CS-COR-002
\begin{cor}[Pendekatan peringkat hingga terbaik]
\label{o008bridge:best-rank-n}
Jika $R\colon H_1\to H_2$ berperingkat paling banyak~$N$, maka
\[
  \|T-R\|\geq s_{N+1}.
\]
Karena $\|T-T_N\|=s_{N+1}$, pemotongan SVD~$T_N$ merupakan pendekatan
peringkat paling banyak~$N$ yang terbaik dalam norma operator.
\end{cor}

\begin{proof}
Jika daftar nilai singular mempunyai paling banyak~$N$ suku, maka
$s_{N+1}=0$ dan ketaksamaannya langsung. Andaikan $s_{N+1}$ ada.
Subruang $E=\spn\{v_1,\ldots,v_{N+1}\}$ berdimensi~$N+1$, sedangkan
$R|_E$ berperingkat paling banyak~$N$. Jadi terdapat vektor satuan
$x\in E\cap\ker R$. Tuliskan $x=\sum_{j=1}^{N+1}c_jv_j$. Karena
$s_j\geq s_{N+1}$,
\[
 \|(T-R)x\|^2=\|Tx\|^2
 =\sum_{j=1}^{N+1}s_j^2|c_j|^2
 \geq s_{N+1}^2\sum_{j=1}^{N+1}|c_j|^2
 =s_{N+1}^2.
\]
Dengan mengambil supremum atas vektor satuan diperoleh batas bawah yang
diminta; teorema sebelumnya menunjukkan bahwa~$T_N$ mencapainya.
\end{proof}

% O008-BRIDGE-ID: O008-BRIDGE-CS-PROP-001
\begin{prop}[Hubungan dengan dekomposisi polar]
\label{o008bridge:svd-polar}
Dalam dekomposisi polar $T=U|T|$, isometri parsial~$U$ memenuhi
$Uv_n=u_n$ untuk semua indeks singular. Secara khusus,
\[
 |T|=\sum_n s_n(v_n\otimes v_n),
 \qquad
 U|T|=\sum_n s_n(u_n\otimes v_n)=T.
\]
\end{prop}

\begin{proof}
Karena $|T|v_n=s_nv_n$ dan $Tv_n=s_nu_n$, persamaan
$T=U|T|$ memberi $s_nUv_n=s_nu_n$. Nilai $s_n$ positif, sehingga
$Uv_n=u_n$. Kedua ekspansi kemudian mengikuti dari teorema spektral bagi
$|T|$ dan konvergensi norma dalam
Teorema~\ref{o008bridge:compact-svd}.
\end{proof}

% O008-BRIDGE-ID: O008-BRIDGE-CS-EXAM-002
\begin{exam}[Matriks sebagai kasus berhingga]
Jika $H_1=\C^n$ dan $H_2=\C^m$, semua operator otomatis kompak. Setelah
$(v_n)$ dan $(u_n)$ diperluas menjadi basis ortonormal pada domain dan
kodomain, matriks operator mengambil bentuk
\[
  T=U\Sigma V^*,
\]
dengan $U$ dan $V$ uniter serta $\Sigma$ diagonal-persegi-panjang yang entri
tak nolnya adalah $s_1\geq s_2\geq\cdots$. Jadi teorema di atas tepat
memperluas SVD matriks ke operator kompak antar-ruang Hilbert.
\end{exam}

\section{Peta penggunaan}

Urutan konseptualnya bergerak dari operator kompak (Bab~7), melalui spektrum
(Bab~8), menuju Teorema~\ref{o008bridge:riesz-schauder-spectrum}, teorema
spektral kompak swaadjoin, lalu SVD. Prasyarat pembuktiannya lebih rinci:
\begin{enumerate}[label=\rm{(\alph*)}]
  \item pembuktian Riesz--Schauder memakai alternatif Fredholm yang telah
        dibuktikan pada Bab~15;
  \item teorema spektral kompak swaadjoin memakai hasil tersebut bersama
        geometri ruang Hilbert dari Bab~4;
  \item konstruksi $|T|=(T^*T)^{1/2}$ dan SVD juga memakai akar positif dan
        kalkulus fungsional kontinu dari Bab~11.
\end{enumerate}
Peta ini ditempatkan setelah seluruh buku agar tidak menimbulkan ketergantungan
melingkar dan tidak menggandakan pembuktian Fredholm dari Bab~15.


% O001 selected central reader-work results -- separately authored proofs.
% Model provenance: OpenAI Codex gpt-5.6-sol, Ultra, atas arahan pengguna.
% Component license: CC BY-SA 4.0.
% Not authored or endorsed by John M. Erdman or Portland State University.

\chapter{DUKUNGAN PENGUASAAN UNTUK HASIL KERJA-PEMBACA TERPILIH}

\begin{center}
\small
Sepuluh pembuktian asli terpisah, CC BY-SA 4.0. Ditulis dengan bantuan
OpenAI Codex gpt-5.6-sol, Ultra, atas arahan pengguna. Pernyataan dan petunjuk
berasal dari edisi terjemahan karya John M. Erdman; pembuktian lengkap ini
bukan tulisan beliau dan tidak menyiratkan dukungan beliau atau Portland State
University.
\end{center}

% O001-READER-WORK-SOLUTION-ID: O001-FAOA-2015-CH01-RW-001-SOLUTION
% SOURCE-RESULT-ID: FAOA-2015-CH01-NODE-0070
% SOURCE-HINT-ID: FAOA-2015-CH01-NODE-0071
% RESULT-TARGET-SHA256: 42ef1dbc69da73f7b5a0e6f64304b7e60174116ffe724d7f402b1f4be2b51632
% HINT-TARGET-SHA256: 9922dc89fda69dbcd1243bb117b5773d60030ad17056bb1be739d9ae4634ba4f
\begin{o001readerwork}
  {O001-FAOA-2015-CH01-RW-001-SOLUTION}
  {FAOA-2015-CH01-NODE-0070}
  {FAOA-2015-CH01-NODE-0071}
  {42ef1dbc69da73f7b5a0e6f64304b7e60174116ffe724d7f402b1f4be2b51632}
  {9922dc89fda69dbcd1243bb117b5773d60030ad17056bb1be739d9ae4634ba4f}
\begin{o001result}
\begin{thm}[Ketaksamaan Schwarz]\label{00015002} Misalkan $s$ suatu hasil kali dalam semu yang didefinisikan pada ruang vektor~$V$.
 \index{ketaksamaan Schwarz}%
 \index{ketaksamaan!Schwarz}%
Maka ketaksamaan
   \[ \abs{s(x,y)} \le \norm x \,\norm y. \]
berlaku untuk semua vektor $x$ dan~$y$ dalam~$V$.
\end{thm}
\end{o001result}
\begin{o001sourcehint}
\begin{proof}[\emph{Petunjuk pembuktian}] Tetapkan $x$, $y \in V$. Untuk setiap $\alpha \in \K$, kita mengetahui bahwa
  \begin{equation}\label{Schwarz_ineq}
      0 \le s(x - \alpha y, x - \alpha y) .
  \end{equation}
Uraikan ruas kanan~\eqref{Schwarz_ineq} menjadi empat suku. Jika $s(y,x)=0$, kesimpulannya langsung. Jika tidak,
tuliskan $s(y,x)$ dalam bentuk polar: $s(y,x) = r e^{i\theta}$, dengan $r > 0$ dan $\theta \in \R$.
Kemudian, dalam ketaksamaan yang diperoleh, tinjau $\alpha$ yang berbentuk $e^{-i\theta}t$ dengan
$t \in \R$. Perhatikan bahwa kini ruas kanan~\eqref{Schwarz_ineq} merupakan polinom kuadrat dalam~$t$. Apa yang dapat Anda katakan tentang
diskriminannya? \ns
\end{proof}
\end{o001sourcehint}
\begin{o001answer}
Diskriminan polinom kuadrat yang selalu tak negatif harus tak positif; hasilnya
$|s(x,y)|^2\leq s(x,x)s(y,y)=\|x\|^2\|y\|^2$.
\end{o001answer}
\begin{o001proof}
Tuliskan $s(y,x)=re^{i\theta}$ dengan $r\geq0$. Karena $s$ Hermitian,
$s(x,y)=re^{-i\theta}$. Untuk $t\in\R$, ambil
$\alpha=e^{-i\theta}t$. Seskuilinearitas memberi
\begin{align*}
0&\leq s(x-\alpha y,x-\alpha y)\\
 &=s(x,x)-\overline\alpha s(x,y)-\alpha s(y,x)
   +|\alpha|^2s(y,y)\\
 &=\|x\|^2-2rt+t^2\|y\|^2.
\end{align*}
Jika $\|y\|>0$, polinom kuadrat ini tak negatif untuk setiap~$t$, sehingga
diskriminannya tak positif:
$4r^2-4\|x\|^2\|y\|^2\leq0$. Jika $\|y\|=0$, polinom linear
$\|x\|^2-2rt$ tak mungkin tak negatif untuk semua~$t$ kecuali $r=0$;
ketaksamaan yang sama tetap berlaku. Karena $r=|s(x,y)|$, mengambil akar
kuadrat membuktikan ketaksamaan Schwarz, termasuk kasus seminorma nol.
\end{o001proof}
\end{o001readerwork}

% O001-READER-WORK-SOLUTION-ID: O001-FAOA-2015-CH03-RW-001-SOLUTION
% SOURCE-RESULT-ID: FAOA-2015-CH03-NODE-0157
% SOURCE-HINT-ID: FAOA-2015-CH03-NODE-0158
% RESULT-TARGET-SHA256: 8237a739b34dd84596960ed828967e88de8c127cf9a24022dfd6cf4bda2be6cb
% HINT-TARGET-SHA256: 8a9f7610446c0e7aed01ed8eded7c0b6745888c5263661e4479f4f03d4730404
\begin{o001readerwork}
  {O001-FAOA-2015-CH03-RW-001-SOLUTION}
  {FAOA-2015-CH03-NODE-0157}
  {FAOA-2015-CH03-NODE-0158}
  {8237a739b34dd84596960ed828967e88de8c127cf9a24022dfd6cf4bda2be6cb}
  {8a9f7610446c0e7aed01ed8eded7c0b6745888c5263661e4479f4f03d4730404}
\begin{o001result}
\begin{thm}[Teorema Hahn--Banach I]\label{HBTI} Misalkan $M$ suatu subruang dari ruang vektor real $V$ dan $p$
 \index{teorema Hahn--Banach}%
suatu fungsional sublinear pada~$V$. Jika $f$ suatu fungsional linear pada $M$ sedemikian sehingga $f \le p$ pada $M$, maka
$f$ memiliki perpanjangan $g$ ke seluruh $V$ sedemikian sehingga $g \le p$ pada~$V$.
\end{thm}
\end{o001result}
\begin{o001sourcehint}
\begin{proof}[\emph{Petunjuk pembuktian}] Dalam teorema ini, notasi $h \le p$ berarti bahwa $\dom h \subseteq \dom p = V$
dan bahwa $h(x) \le p(x)$ untuk semua $x \in \dom h$. Dalam hal ini kita mengatakan bahwa $h$
 \index{didominasi}%
\df{didominasi} oleh~$p$. Kita menulis
 \index{<binrelle@$f \preccurlyeq g$ ($g$ merupakan perpanjangan~$f$)}%
$f \preccurlyeq g$ jika $g$ merupakan perpanjangan~$f$ (yakni,
jika $f(x) = g(x)$ untuk semua $x \in \dom f$).

Misalkan $\fml S$ keluarga semua fungsional linear bernilai real pada subruang-subruang $V$ yang merupakan perpanjangan~$f$ dan
didominasi oleh~$p$. Keluarga $\fml S$ terurut parsial oleh $\preccurlyeq$. Gunakan \emph{lema Zorn} untuk menghasilkan suatu
elemen maksimal $g$ dari~$\fml S$. Pembuktian selesai jika Anda dapat menunjukkan bahwa $\dom{g} = V$.

Untuk itu, andaikan demi kontradiksi bahwa terdapat vektor $c \in V$ yang tidak termasuk dalam~$D = \dom{g}$. Buktikan terlebih dahulu bahwa
   \begin{equation}\label{eq_HBTI}
      g(y) - p(y - c) \le - g(x) + p(x + c)
   \end{equation}
untuk semua $x$, $y \in D$. Simpulkan bahwa terdapat bilangan $\gamma$ yang terletak di antara supremum semua bilangan
di ruas kiri~\eqref{eq_HBTI} dan infimum semua bilangan di ruas kanan.

Misalkan $E$ rentang dari $D \cup \{c\}$. Definisikan $h$ pada $E$ dengan $h(x + \alpha c) = g(x) + \alpha\gamma$ untuk $x \in D$ dan
$\alpha \in \R$. Sekarang tunjukkan bahwa $h \in \fml S$. (Ambil $\alpha > 0$ dan tinjau secara terpisah $h(x + \alpha c)$ dan $h(x - \alpha c)$.) \ns
\end{proof}
\end{o001sourcehint}
\begin{o001answer}
Keluarga semua perpanjangan yang didominasi mempunyai elemen maksimal menurut
lema Zorn; elemen maksimal itu harus berdomain seluruh~$V$ karena setiap
vektor di luar domain memungkinkan perpanjangan satu dimensi yang masih
didominasi oleh~$p$.
\end{o001answer}
\begin{o001proof}
Misalkan $\fml S$ terdiri atas semua pasangan $(D,h)$ dengan $M\subseteq D$
subruang $V$, $h\colon D\to\R$ linear, $f\preccurlyeq h$, dan $h\leq p$.
Urutkan dengan perpanjangan. Himpunan ini tidak kosong karena memuat~$(M,f)$.
Untuk setiap rantai, gabungan domainnya adalah subruang dan fungsi gabungan
terdefinisi dengan baik: dua anggota rantai dapat dibandingkan, sehingga
nilainya cocok pada irisan. Fungsi gabungan linear, memperpanjang~$f$, dan
tetap didominasi oleh~$p$. Jadi setiap rantai mempunyai batas atas. Lema Zorn
memberi elemen maksimal $(D,g)$.

Andaikan $D\neq V$ dan pilih $c\notin D$. Untuk $x,y\in D$,
\[
 g(x)+g(y)=g(x+y)\leq p(x+y)
 \leq p(x+c)+p(y-c).
\]
Maka
\[
 g(y)-p(y-c)\leq-g(x)+p(x+c).
\]
Definisikan
\[
 A=\sup_{y\in D}\{g(y)-p(y-c)\},\qquad
 B=\inf_{x\in D}\{-g(x)+p(x+c)\}.
\]
Ketaksamaan tadi menunjukkan $A\leq B$; kedua batas bernilai real karena
masing-masing keluarga tidak kosong dan saling membatasi. Pilih
$\gamma\in[A,B]$.

Karena $c\notin D$, setiap elemen $E=D+\R c$ mempunyai penulisan tunggal
$x+\alpha c$. Tetapkan
$h(x+\alpha c)=g(x)+\alpha\gamma$. Fungsi~$h$ linear dan memperpanjang~$g$.
Tinggal memeriksa dominasi. Jika $\alpha>0$, dari $\gamma\leq B$ dengan
$x/\alpha\in D$ diperoleh
\[
 g(x/\alpha)+\gamma\leq p(x/\alpha+c),
\]
sehingga, oleh homogenitas positif~$p$,
$h(x+\alpha c)\leq p(x+\alpha c)$. Jika $\alpha=-\beta<0$, dari
$\gamma\geq A$ dengan $x/\beta\in D$ diperoleh
\[
 g(x/\beta)-\gamma\leq p(x/\beta-c),
\]
dan setelah dikalikan~$\beta$ didapat
$h(x-\beta c)\leq p(x-\beta c)$. Kasus $\alpha=0$ adalah dominasi~$g$.
Jadi $(E,h)\in\fml S$ memperluas $(D,g)$ secara ketat, bertentangan dengan
maksimalitas. Karena itu $D=V$, dan~$g$ adalah perpanjangan yang diminta.
\end{o001proof}
\end{o001readerwork}

% O001-READER-WORK-SOLUTION-ID: O001-FAOA-2015-CH04-RW-001-SOLUTION
% SOURCE-RESULT-ID: FAOA-2015-CH04-NODE-0064
% SOURCE-HINT-ID: FAOA-2015-CH04-NODE-0065
% RESULT-TARGET-SHA256: 2676db6dfe542cf8bd2de4096013eb24533e755933d9953dd70b09a91e91f3d2
% HINT-TARGET-SHA256: bb8f98e29368268dd55b5c76f34a2060e5a0a38b303eccb5d877395800b7fbfa
\begin{o001readerwork}
  {O001-FAOA-2015-CH04-RW-001-SOLUTION}
  {FAOA-2015-CH04-NODE-0064}
  {FAOA-2015-CH04-NODE-0065}
  {2676db6dfe542cf8bd2de4096013eb24533e755933d9953dd70b09a91e91f3d2}
  {bb8f98e29368268dd55b5c76f34a2060e5a0a38b303eccb5d877395800b7fbfa}
\begin{o001result}
\begin{thm}[Teorema Riesz--Fr\'echet]\label{000342} Jika $f \in H^*$ dengan $H$ ruang Hilbert,
 \index{teorema Riesz--Fr\'echet}%
maka terdapat vektor unik $a$ dalam $H$ sehingga $f = \psi_a$. Selain itu,
$\norm a = \norm{\psi_a}$.
\end{thm}
\end{o001result}
\begin{o001sourcehint}
\begin{proof}[\emph{Petunjuk pembuktian}] Untuk kasus $f \ne 0$, pilih vektor satuan $z$ dalam $(\ker f)^\perp$.
Perhatikan bahwa untuk setiap $x \in H$, vektor $x - \frac{f(x)}{f(z)}\,z$ termasuk dalam $\ker f$. \ns
\end{proof}
\end{o001sourcehint}
\begin{o001answer}
Jika $f\neq0$, pilih vektor satuan $z\in(\ker f)^\perp$ dan ambil
$a=\overline{f(z)}z$; maka $f(x)=\langle x,a\rangle$ untuk semua~$x$, dan
$\|a\|=\|f\|$.
\end{o001answer}
\begin{o001proof}
Jika $f=0$, ambil $a=0$. Andaikan $f\neq0$. Kernel~$f$ tertutup dan proper,
sehingga $(\ker f)^\perp$ memuat suatu vektor satuan~$z$. Nilai $f(z)$ tidak
nol, sebab jika nol maka $z\in\ker f\cap(\ker f)^\perp=\{0\}$.

Untuk setiap $x\in H$, vektor
$w=x-f(x)f(z)^{-1}z$ termasuk dalam $\ker f$. Karena $z\perp\ker f$,
\[
 0=\langle w,z\rangle
  =\langle x,z\rangle-\frac{f(x)}{f(z)}.
\]
Jadi $f(x)=f(z)\langle x,z\rangle$. Dengan konvensi bahwa hasil kali dalam
linear dalam variabel pertama, jika $a=\overline{f(z)}z$ maka
\[
 \langle x,a\rangle=f(z)\langle x,z\rangle=f(x).
\]
Ketaksamaan Schwarz memberi $\|f\|\leq\|a\|$. Jika $a\neq0$, penerapan pada
$x=a/\|a\|$ memberi $|f(x)|=\|a\|$, sehingga $\|f\|=\|a\|$; kasus
$a=0$ sudah selesai. Jika $\psi_a=\psi_b$, maka
$\langle x,a-b\rangle=0$ untuk setiap~$x$. Mengambil $x=a-b$ memberi
$a=b$, jadi vektor representasi unik.
\end{o001proof}
\end{o001readerwork}

% O001-READER-WORK-SOLUTION-ID: O001-FAOA-2015-CH05-RW-001-SOLUTION
% SOURCE-RESULT-ID: FAOA-2015-CH05-NODE-0023
% SOURCE-HINT-ID: FAOA-2015-CH05-NODE-0024
% RESULT-TARGET-SHA256: 93218a7be4fa5cc031f93ac7990bf692f03cbf101a28e4539817f9f4f3fe6fb7
% HINT-TARGET-SHA256: c05412a2b16f484469ac8e68a65ea3695e647c245369758bd8795660de5b29e9
\begin{o001readerwork}
  {O001-FAOA-2015-CH05-RW-001-SOLUTION}
  {FAOA-2015-CH05-NODE-0023}
  {FAOA-2015-CH05-NODE-0024}
  {93218a7be4fa5cc031f93ac7990bf692f03cbf101a28e4539817f9f4f3fe6fb7}
  {c05412a2b16f484469ac8e68a65ea3695e647c245369758bd8795660de5b29e9}
\begin{o001result}
\begin{prop} Misalkan $\phi \colon H \times K \sto \C$ fungsional seskuilinear terbatas pada
hasil kali dua ruang Hilbert tak nol. Maka terdapat pemetaan linear terbatas yang unik $T \in \ofml B(H,K)$ dan
$S \in \ofml B(K,H)$ sedemikian sehingga
    \[ \phi(x,y) = \langle Tx,y \rangle = \langle x,Sy \rangle \]
untuk semua $x \in H$ dan $y \in K$. Selain itu, $\norm T = \norm S = \norm\phi$.
\end{prop}
\end{o001result}
\begin{o001sourcehint}
\begin{proof}[\emph{Petunjuk pembuktian}] Tunjukkan bahwa untuk setiap $x \in H$, pemetaan
$y \mapsto \conj{\phi(x,y)}$ adalah fungsional linear terbatas pada~$K$. \ns
\end{proof}
\end{o001sourcehint}
\begin{o001answer}
Teorema Riesz--Fréchet yang diterapkan pada tiap irisan variabel menghasilkan
operator unik $T$ dan~$S$; definisi norma seskuilinear memberi
$\|T\|=\|S\|=\|\phi\|$.
\end{o001answer}
\begin{o001proof}
Untuk $x\in H$ tetap, fungsi
$F_x(y)=\overline{\phi(x,y)}$ linear pada~$K$ dan memenuhi
$|F_x(y)|\leq\|\phi\|\|x\|\|y\|$. Teorema Riesz--Fréchet memberi vektor
tunggal $Tx\in K$ dengan $F_x(y)=\langle y,Tx\rangle$. Setelah dikonjugatkan,
\[
 \phi(x,y)=\langle Tx,y\rangle.
\]
Linearitas $\phi$ dalam~$x$ dan ketunggalan vektor representasi menunjukkan
$T(\alpha x+\beta z)=\alpha Tx+\beta Tz$. Selain itu,
$\|Tx\|=\|F_x\|\leq\|\phi\|\|x\|$, jadi $T$ terbatas dan
$\|T\|\leq\|\phi\|$. Sebaliknya,
$|\phi(x,y)|=|\langle Tx,y\rangle|\leq\|T\|\|x\|\|y\|$, sehingga
$\|\phi\|\leq\|T\|$. Jadi $\|T\|=\|\phi\|$.

Untuk $y\in K$ tetap, fungsi $x\mapsto\phi(x,y)$ adalah fungsional linear
terbatas pada~$H$. Riesz--Fréchet memberi vektor tunggal $Sy\in H$ sehingga
$\phi(x,y)=\langle x,Sy\rangle$. Karena $\phi$ linear konjugat dalam~$y$ dan
hasil kali dalam juga linear konjugat dalam variabel kedua, ketunggalan
menunjukkan bahwa $S$ linear. Argumen norma yang sama memberi
$\|S\|=\|\phi\|$. Jika operator lain mempunyai salah satu identitas tersebut,
ketunggalan dalam teorema Riesz--Fréchet pada setiap irisan memaksanya sama
dengan $T$ atau~$S$. Dengan demikian kedua operator unik.
\end{o001proof}
\end{o001readerwork}

% O001-READER-WORK-SOLUTION-ID: O001-FAOA-2015-CH06-RW-001-SOLUTION
% SOURCE-RESULT-ID: FAOA-2015-CH06-NODE-0024
% SOURCE-HINT-ID: FAOA-2015-CH06-NODE-0025
% RESULT-TARGET-SHA256: ad2506613fe9a301018275484eb844333cf2b5efa5a268c3a54ff1a2c9882b64
% HINT-TARGET-SHA256: 888706ad7f71a8381be45b7bcf76b7f3e4d5156dc7bcdb7566f2680c443fd808
\begin{o001readerwork}
  {O001-FAOA-2015-CH06-RW-001-SOLUTION}
  {FAOA-2015-CH06-NODE-0024}
  {FAOA-2015-CH06-NODE-0025}
  {ad2506613fe9a301018275484eb844333cf2b5efa5a268c3a54ff1a2c9882b64}
  {888706ad7f71a8381be45b7bcf76b7f3e4d5156dc7bcdb7566f2680c443fd808}
\begin{o001result}
\begin{thm}[Teorema Alaoglu]\label{C067441} Bola satuan tertutup dalam dual suatu ruang linear bernorma
 \index{teorema Alaoglu}%
kompak dalam topologi-$w^*$.
\end{thm}
\end{o001result}
\begin{o001sourcehint}
\begin{proof}[\emph{Petunjuk pembuktian}] Misalkan $V$ suatu ruang linear bernorma, dan nyatakan dengan $[V^*]_1$ bola satuan tertutup dari
ruang dualnya.  Untuk setiap $x \in V$, definisikan $D_x = \{\lambda \in \K \colon \abs\lambda \le \norm x\}$.  Hasil kali
\smash[b]{$\mathbf P = \prod\limits_{x \in V} D_x$} kompak berdasarkan \emph{teorema Tychonoff}.  Tinjau fungsi
   \[ \Phi\colon [V^*]_1 \sto \mathbf P \colon f \mapsto \bigl(f(x)\bigr)_{x \in V} \]
dengan $[V^*]_1$ dilengkapi topologi-$w^*$ dan $\mathbf P$ dilengkapi topologi hasil kali (lihat~\ref{C021437}).  Di sini kita menulis
$\bigl(f(x)\bigr)_{x \in V}$ untuk keluarga bilangan terindeks (atau barisan tergeneralisasi), bukan notasi yang lebih lazim
$\bigl(f_x\bigr)_{x \in V}$.  (Untuk pembahasan lebih lanjut tentang keluarga terindeks, lihat~\cite{Erdman:2007}, Bagian~7.2.)  Mudah dilihat bahwa
$\Phi$ injektif.  Untuk menunjukkan bahwa pemetaan itu kontinu, misalkan $\bigl(f_\lambda\bigr)_{\lambda \in \Lambda}$ suatu jaring dalam $[V^*]_1$
dan $g$ suatu fungsi dalam $[V^*]_1$ sedemikian sehingga $f_\lambda \to^{w^*} g$.  Perhatikan bahwa hal ini terjadi jika dan hanya jika
$\pi_x(f_\lambda) \sto \pi_x(g)$ untuk setiap $x \in V$ (dengan $\pi_x$ proyeksi koordinat pada~$\mathbf P$---lihat~\ref{C015531}).
Terakhir, misalkan $\bigl(f_\lambda\bigr)_{\lambda \in \Lambda}$ suatu jaring dalam $[V^*]_1$ sedemikian sehingga $\bigl(\Phi(f_\lambda)\bigr)_{\lambda \in \Lambda}$
konvergen dalam~$\mathbf P$.  Tunjukkan bahwa jaring $\bigl(f_\lambda(x)\bigr)_{\lambda \in \Lambda}$ konvergen dalam $\K$ untuk setiap $x \in V$.
Misalkan $g(x) = \lim\limits_\lambda f_\lambda(x)$ untuk setiap $x \in V$, lalu buktikan bahwa $g$ linear, $\norm g \le 1$,
$f_\lambda \to^{w^*} g$ dalam $[V^*]_1$, dan $\Phi(f_\lambda) \sto \Phi(g)$ dalam~$\mathbf P$.  \ns
\end{proof}
\end{o001sourcehint}
\begin{o001answer}
Bola dual tertutup tertanam secara homeomorfik sebagai subhimpunan tertutup
dari hasil kali kompak $\prod_{x\in V}\{\lambda:|\lambda|\leq\|x\|\}$;
karena itu ia kompak dalam topologi-$w^*$.
\end{o001answer}
\begin{o001proof}
Untuk setiap $x\in V$, cakram tertutup
$D_x=\{\lambda\in\K:|\lambda|\leq\|x\|\}$ kompak. Teorema Tychonoff memberi
kekompakan ruang hasil kali $\mathbf P=\prod_{x\in V}D_x$. Jika
$[V^*]_1$ adalah bola satuan dual, definisikan
\[
 \Phi(f)=(f(x))_{x\in V}.
\]
Ketaksamaan $|f(x)|\leq\|x\|$ memastikan bahwa $\Phi(f)\in\mathbf P$.
Pemetaan ini injektif. Berdasarkan definisi topologi-$w^*$, suatu jaring
$f_\lambda$ konvergen ke~$f$ jika dan hanya jika setiap koordinat
$f_\lambda(x)$ konvergen ke~$f(x)$. Jadi $\Phi$ adalah homeomorfisme dari
$[V^*]_1$ ke citranya dengan topologi subruang.

Kita tunjukkan bahwa citra itu tertutup. Misalkan jaring
$\Phi(f_\lambda)$ konvergen di~$\mathbf P$ ke $z=(z_x)_{x\in V}$ dan tetapkan
$g(x)=z_x$. Limit koordinat mempertahankan penjumlahan dan perkalian skalar:
\[
 g(\alpha x+\beta y)
 =\lim_\lambda f_\lambda(\alpha x+\beta y)
 =\alpha g(x)+\beta g(y).
\]
Selain itu $|g(x)|\leq\|x\|$ karena $z_x\in D_x$. Maka $g$ linear terbatas
dengan $\|g\|\leq1$, sehingga $g\in[V^*]_1$, dan $z=\Phi(g)$. Dengan
kriteria jaring untuk ketertutupan, $\Phi([V^*]_1)$ tertutup dalam~$\mathbf P$.
Citra itu kompak sebagai subhimpunan tertutup ruang kompak; homeomorfisme
$\Phi$ kemudian membuktikan kekompakan-$w^*$ bola dual.
\end{o001proof}
\end{o001readerwork}

% O001-READER-WORK-SOLUTION-ID: O001-FAOA-2015-CH06-RW-002-SOLUTION
% SOURCE-RESULT-ID: FAOA-2015-CH06-NODE-0030
% SOURCE-HINT-ID: FAOA-2015-CH06-NODE-0031
% RESULT-TARGET-SHA256: 7325fa61fae025e9f01d4aa48c36f6c41f02c4d0f6f073f1d2ac750a9bd1c70d
% HINT-TARGET-SHA256: 9b01db1c39da23fc93882972564dcf6c8ba7d482ad352c2c6efb50f6fe58fe89
\begin{o001readerwork}
  {O001-FAOA-2015-CH06-RW-002-SOLUTION}
  {FAOA-2015-CH06-NODE-0030}
  {FAOA-2015-CH06-NODE-0031}
  {7325fa61fae025e9f01d4aa48c36f6c41f02c4d0f6f073f1d2ac750a9bd1c70d}
  {9b01db1c39da23fc93882972564dcf6c8ba7d482ad352c2c6efb50f6fe58fe89}
\begin{o001result}
\begin{thm}[Teorema pemetaan terbuka]\label{C069414} Setiap surjeksi linear terbatas antara
 \index{terbuka!pemetaan!teorema}%
 \index{morfisme bijektif!dapat diinverskan!dalam $\cat{BAN_\infty}$}%
ruang-ruang Banach merupakan pemetaan terbuka.
\end{thm}
\end{o001result}
\begin{o001sourcehint}
\begin{proof}[\emph{Petunjuk pembuktian}] Misalkan $T\colon B \sto C$ suatu surjeksi linear terbatas antara ruang-ruang Banach.
Mulailah dengan membuktikan klaim berikut: yang perlu kita tunjukkan hanyalah bahwa $T\bigl((B)_1\bigr)$ memuat suatu bola terbuka di sekitar
titik asal dalam~$C$.  Untuk sementara, misalkan kita mengetahui bahwa bola semacam itu ada.  Namai bola itu~$W$.  Diberikan suatu subhimpunan terbuka $U$ dari $B$, kita
ingin menunjukkan bahwa citranya oleh $T$ terbuka dalam~$C$.  Untuk itu, ambil sebarang titik $y$ dalam $T(U)$ dan pilih $x$
dalam $U$ sedemikian sehingga $y = Tx$.  Karena $U - x$ merupakan lingkungan titik asal dalam $B$, kita dapat mencari bilangan $r > 0$ sedemikian sehingga
$r\,(B)_1 \subseteq U - x$.  Buktikan bahwa $T(U)$ terbuka dengan menunjukkan bahwa $y + rW \subseteq T(U)$.

Untuk membuktikan klaim di atas, misalkan $V = T\bigl((B)_1\bigr)$ dan perhatikan bahwa karena $T$ surjektif, $C = T(B) =
\bigcup_{k=1}^\infty k\clo V$.  Gunakan versi \emph{teorema kategori Baire} yang menyatakan bahwa jika suatu ruang metrik lengkap takkosong
merupakan gabungan suatu barisan himpunan tertutup, maka sekurang-kurangnya salah satu himpunan tersebut memiliki interior takkosong.
(Lihat, misalnya, Teorema 29.3.21 dalam~\cite{Erdman:2007}.)  Pilih $k \in \N$ sedemikian sehingga $\bigl(k\clo V\bigr)^o \neq \emptyset$.
Dengan demikian, terdapat $y \in \clo V$ dan $r > 0$ sedemikian sehingga $y + (C)_r \subseteq k\clo V$.  Dari kekonveksan
$k\clo V$, tidak sulit melihat bahwa $(C)_r \subseteq k\clo V$.  (Tuliskan sebarang titik $z$ dalam $(C)_r$ sebagai
$\frac12(y + z) + \frac12(-y + z)$.) Sekarang diperoleh $(C)_{r/k} \subseteq \clo V$.  Gunakan Proposisi~\ref{prop_for_OMT}.
\ns \end{proof}
\end{o001sourcehint}
\begin{o001answer}
Teorema kategori Baire memberi bola di sekitar nol dalam tutupan citra bola
satuan; argumen deret konvergen menghapus tanda tutupan. Linearitas kemudian
memindahkan dan menskalakan bola itu di sekitar setiap titik citra.
\end{o001answer}
\begin{o001proof}
Misalkan $V=T((B)_1)$. Karena $T$ surjektif,
$C=\bigcup_{k\geq1}k\clo V$. Teorema kategori Baire memberi suatu~$k$ sehingga
$k\clo V$ mempunyai interior tak kosong. Maka terdapat $y\in C$ dan $r>0$
dengan $y+(C)_r\subseteq k\clo V$. Himpunan $k\clo V$ konveks dan simetris;
karena itu juga $-y+(C)_r\subseteq k\clo V$. Untuk $z\in(C)_r$, kedua vektor
$y+z$ dan $-y+z$ berada di $k\clo V$, sehingga rata-ratanya~$z$ juga berada
di sana. Jadi $(C)_{r/k}\subseteq\clo V$.

Kita buktikan langkah teknis bahwa bila $(C)_s\subseteq\clo{T((B)_1)}$, maka
$(C)_s\subseteq T((B)_1)$. Dengan penskalaan cukup ambil $s=1$. Untuk
$z\in(C)_1$, pilih $0<\delta<1$ sehingga $\|z\|<1-\delta$, dan tetapkan
$q=z/(1-\delta)$. Secara induktif, dari hipotesis tutupan pilih
$u_n\in\delta^{n-1}(B)_1$ agar
\[
 \left\|q-T\sum_{j=1}^{n}u_j\right\|<\delta^n.
\]
Pemilihan dimulai dengan mendekati~$q$ oleh $T((B)_1)$; pada langkah
berikutnya, bagi residu dengan $\delta^n$, dekati di bola satuan, lalu
skalakan kembali. Deret $\sum u_n$ konvergen mutlak dalam ruang Banach~$B$.
Jika $x=\sum u_n$, kontinuitas~$T$ dan residu yang menuju nol memberi
$Tx=q$, sedangkan $\|x\|<\sum_{n\geq1}\delta^{n-1}=1/(1-\delta)$.
Jadi $z=T((1-\delta)x)$ dengan $\|(1-\delta)x\|<1$. Langkah teknis terbukti,
dan karena itu $V$ memuat suatu bola terbuka~$W$ di sekitar nol.

Sekarang ambil $U\subseteq B$ terbuka dan $y=Tx\in T(U)$. Ada $a>0$ dengan
$x+a(B)_1\subseteq U$. Maka
\[
 y+aW\subseteq Tx+T(a(B)_1)\subseteq T(U).
\]
Jadi setiap titik $T(U)$ mempunyai lingkungan yang termuat dalam $T(U)$;
dengan demikian $T(U)$ terbuka.
\end{o001proof}
\end{o001readerwork}

% O001-READER-WORK-SOLUTION-ID: O001-FAOA-2015-CH06-RW-003-SOLUTION
% SOURCE-RESULT-ID: FAOA-2015-CH06-NODE-0045
% SOURCE-HINT-ID: FAOA-2015-CH06-NODE-0046
% RESULT-TARGET-SHA256: 6b71222dff40bb8011df9a11d4042434fb0957662a12b18711a6ecdd53291517
% HINT-TARGET-SHA256: 23c62195283597b3dfd4933a536d993e287129d5f29981be51eecf10015f9000
\begin{o001readerwork}
  {O001-FAOA-2015-CH06-RW-003-SOLUTION}
  {FAOA-2015-CH06-NODE-0045}
  {FAOA-2015-CH06-NODE-0046}
  {6b71222dff40bb8011df9a11d4042434fb0957662a12b18711a6ecdd53291517}
  {23c62195283597b3dfd4933a536d993e287129d5f29981be51eecf10015f9000}
\begin{o001result}
\begin{thm}[Teorema graf tertutup]\label{thm_CGT} Misalkan $T\colon B \sto C$ suatu pemetaan linear antara ruang-ruang Banach.
 \index{tertutup!teorema graf}%
Jika graf $T$ tertutup, maka $T$ kontinu.
\end{thm}
\end{o001result}
\begin{o001sourcehint}
\begin{proof}[\emph{Petunjuk pembuktian}] Misalkan $G = \{(x,Tx)\colon x \in B\}$ graf~$T$. Terapkan
(Korolar~\ref{C069417} dari) \emph{teorema pemetaan terbuka} pada pemetaan
  \[ \pi\colon G \sto B\colon (x,Tx) \mapsto x\,. \]
\ns \end{proof}
\end{o001sourcehint}
\begin{o001answer}
Graf tertutup adalah ruang Banach; proyeksi graf ke domain merupakan bijeksi
linear terbatas, sehingga inversnya terbatas oleh teorema pemetaan terbuka.
Komposisi invers itu dengan proyeksi kodomain adalah~$T$.
\end{o001answer}
\begin{o001proof}
Lengkapi $B\times C$ dengan norma produk, misalnya
$\|(x,y)\|=\|x\|+\|y\|$. Ruang ini Banach. Karena graf
$G=\{(x,Tx):x\in B\}$ tertutup, $G$ adalah subruang Banach. Pemetaan
\[
 \pi\colon G\to B,
 \qquad
 \pi(x,Tx)=x,
\]
linear, bijektif, dan terbatas karena $\|x\|\leq\|(x,Tx)\|$. Teorema pemetaan
terbuka, dalam bentuk teorema invers terbatas, menyatakan bahwa
$\pi^{-1}\colon B\to G$ terbatas. Proyeksi koordinat kedua
$q\colon G\to C$, $q(x,y)=y$, juga terbatas. Namun
\[
 T=q\circ\pi^{-1}.
\]
Jadi $T$ adalah komposisi dua pemetaan linear terbatas dan karena itu kontinu.
\end{o001proof}
\end{o001readerwork}

% O001-READER-WORK-SOLUTION-ID: O001-FAOA-2015-CH06-RW-004-SOLUTION
% SOURCE-RESULT-ID: FAOA-2015-CH06-NODE-0136
% SOURCE-HINT-ID: FAOA-2015-CH06-NODE-0137
% RESULT-TARGET-SHA256: 0f6ce55313015e56ae0678821dcb6e5cad7d434ff09b62120fa0b3e95ab3354c
% HINT-TARGET-SHA256: d7ed276423d393da1ea3f084b5faf2f4d3ebba7da0a3902ad41dee24d81e6983
\begin{o001readerwork}
  {O001-FAOA-2015-CH06-RW-004-SOLUTION}
  {FAOA-2015-CH06-NODE-0136}
  {FAOA-2015-CH06-NODE-0137}
  {0f6ce55313015e56ae0678821dcb6e5cad7d434ff09b62120fa0b3e95ab3354c}
  {d7ed276423d393da1ea3f084b5faf2f4d3ebba7da0a3902ad41dee24d81e6983}
\begin{o001result}
\begin{thm}[Prinsip keterbatasan seragam]\label{thm_PUB} Misalkan $B$ ruang Banach dan $W$ ruang linear bernorma.
 \index{prinsip!keterbatasan seragam}%
 \index{keterbatasan!prinsip seragam}%
 \index{seragam!keterbatasan!prinsip}%
Jika suatu keluarga $\ofml T$ dari pemetaan linear terbatas dari $B$ ke $W$ terbatas titik demi
titik, maka keluarga itu terbatas seragam.
\end{thm}
\end{o001result}
\begin{o001sourcehint}
\begin{proof}[\emph{Petunjuk pembuktian}] Untuk setiap $T \in \ofml T$, definisikan
  \[ \sbsb fT\colon B \sto \R\colon x \mapsto \norm{Tx}. \]
Gunakan proposisi sebelumnya untuk menunjukkan bahwa terdapat subhimpunan terbuka tak kosong $U$ dari $B$
dan konstanta $M > 0$ sedemikian sehingga $\sbsb fT(x) \le M$ untuk setiap $T \in \ofml T$ dan setiap
$x \in U$. Pilih titik $a \in B$ dan bilangan $r > 0$ sedemikian sehingga $\clo{B_r(a)} \subseteq
U$. Kemudian verifikasi bahwa
  \[ \norm{Ty} \le r^{-1}\bigl(\, \sbsb fT(a + ry) + \sbsb fT(a)\,\bigr) \le 2Mr^{-1} \]
untuk setiap $T \in \ofml T$ dan setiap $y \in B$ sedemikian sehingga $\norm y \le 1$. \ns
\end{proof}
\end{o001sourcehint}
\begin{o001answer}
Teorema kategori Baire memberi satu bola tempat seluruh keluarga terbatas oleh
satu konstanta; linearitas memindahkan estimasi itu ke bola satuan dan memberi
$\sup_{T\in\ofml T}\|T\|<\infty$.
\end{o001answer}
\begin{o001proof}
Untuk $n\in\N$, tetapkan
\[
 A_n=\{x\in B:\|Tx\|\leq n\text{ untuk setiap }T\in\ofml T\}.
\]
Setiap $A_n$ tertutup sebagai irisan prabayangan bola tertutup oleh pemetaan
kontinu. Keterbatasan titik demi titik menyatakan
$B=\bigcup_{n\geq1}A_n$. Karena $B$ Banach, teorema kategori Baire memberi
$N$ sehingga $A_N$ mempunyai interior tak kosong. Pilih $a\in B$ dan $r>0$
sedemikian sehingga bola tertutup $\clo{B_r(a)}$ termuat dalam~$A_N$.

Jika $\|y\|\leq1$ dan $T\in\ofml T$, maka $a$ serta $a+ry$ berada dalam
$A_N$. Linearitas dan ketaksamaan segitiga memberi
\[
 r\|Ty\|=\|T(a+ry)-Ta\|
 \leq\|T(a+ry)\|+\|Ta\|\leq2N.
\]
Jadi $\|Ty\|\leq2N/r$ seragam untuk semua $T$ dan semua vektor dalam bola
satuan. Dengan mengambil supremum atas~$y$, diperoleh
$\|T\|\leq2N/r$ untuk setiap $T\in\ofml T$. Dengan demikian keluarga
terbatas seragam.
\end{o001proof}
\end{o001readerwork}

% O001-READER-WORK-SOLUTION-ID: O001-FAOA-2015-CH07-RW-001-SOLUTION
% SOURCE-RESULT-ID: FAOA-2015-CH07-NODE-0018
% SOURCE-HINT-ID: FAOA-2015-CH07-NODE-0019
% RESULT-TARGET-SHA256: 457ea425abb80f594c4cc4e579674722ce2be7ae22e3d7c523a1049ad7d323f9
% HINT-TARGET-SHA256: 6f4e9644b3dc23c58bb2b84c5d77764cc8828e283d03c093b98951b736cb1d1a
\begin{o001readerwork}
  {O001-FAOA-2015-CH07-RW-001-SOLUTION}
  {FAOA-2015-CH07-NODE-0018}
  {FAOA-2015-CH07-NODE-0019}
  {457ea425abb80f594c4cc4e579674722ce2be7ae22e3d7c523a1049ad7d323f9}
  {6f4e9644b3dc23c58bb2b84c5d77764cc8828e283d03c093b98951b736cb1d1a}
\begin{o001result}
\begin{exam} Misalkan $C$ ruang Banach $\fml C([0,1])$ dan $k \in \fml C([0,1] \times [0,1])$.  Untuk setiap $f \in C$
 \index{integral!operator}%
 \index{operator!integral}%
definisikan fungsi $Kf$ pada $[0,1]$ dengan
   \[ Kf(x) = \int_0^1 k(x,y)f(y)\,dy\,. \]
Maka operator integral $K$ merupakan operator kompak pada~$C$.
\end{exam}
\end{o001result}
\begin{o001sourcehint}
\begin{proof}[\emph{Petunjuk untuk bukti}] Gunakan \emph{teorema Arzel\`a-Ascoli}~\ref{AAThm}.  \ns
\end{proof}
\end{o001sourcehint}
\begin{o001answer}
Citra bola satuan oleh~$K$ terbatas seragam dan ekuikontinu; teorema
Arzelà--Ascoli membuat tutupannya kompak, sehingga~$K$ operator kompak.
\end{o001answer}
\begin{o001proof}
Linearitas~$K$ langsung dari linearitas integral. Jika
$M=\|k\|_\infty$, maka untuk $\|f\|_\infty\leq1$,
\[
 |Kf(x)|\leq\int_0^1|k(x,y)|\,dy\leq M.
\]
Jadi citra bola satuan terbatas seragam. Karena $k$ kontinu pada persegi
kompak, ia kontinu seragam. Diberikan $\epsilon>0$, ada $\delta>0$ sehingga
$|x-x'|<\delta$ mengakibatkan
\[
 \sup_{0\leq y\leq1}|k(x,y)-k(x',y)|<\epsilon.
\]
Untuk semua $f$ dalam bola satuan,
\[
 |Kf(x)-Kf(x')|
 \leq\int_0^1|k(x,y)-k(x',y)|\,|f(y)|\,dy
 <\epsilon.
\]
Dengan demikian keluarga $K((C)_1)$ ekuikontinu. Teorema Arzelà--Ascoli
menyatakan bahwa tutupannya kompak dalam $\fml C([0,1])$ dengan norma
supremum. Jadi $K$ membawa bola satuan ke himpunan relatif kompak dan,
menurut definisi, $K$ merupakan operator kompak.
\end{o001proof}
\end{o001readerwork}

% O001-READER-WORK-SOLUTION-ID: O001-FAOA-2015-CH08-RW-001-SOLUTION
% SOURCE-RESULT-ID: FAOA-2015-CH08-NODE-0052
% SOURCE-HINT-ID: FAOA-2015-CH08-NODE-0053
% RESULT-TARGET-SHA256: c357c420d40d3fc1aebb795d7a32686947ef6b99c15f0ba276cb122d3b67df7c
% HINT-TARGET-SHA256: 896fca1413df66e29c5b7f5161d8f4d7de60c2a6c90fc3c3a5323625bed07609
\begin{o001readerwork}
  {O001-FAOA-2015-CH08-RW-001-SOLUTION}
  {FAOA-2015-CH08-NODE-0052}
  {FAOA-2015-CH08-NODE-0053}
  {c357c420d40d3fc1aebb795d7a32686947ef6b99c15f0ba276cb122d3b67df7c}
  {896fca1413df66e29c5b7f5161d8f4d7de60c2a6c90fc3c3a5323625bed07609}
\begin{o001result}
\begin{prop}\label{C073134} Spektrum setiap elemen dari aljabar Banach beridentitas tidak kosong.
\end{prop}
\end{o001result}
\begin{o001sourcehint}
\begin{proof}[\emph{Petunjuk untuk bukti}] Gunakan argumen kontradiksi. Gunakan \emph{teorema Liouville}~\ref{C073131}
untuk menunjukkan bahwa $\phi \circ R_a$ konstan bagi setiap fungsional linear terbatas $\phi$ pada~$A$. Kemudian gunakan
(salah satu versi) \emph{teorema Hahn-Banach} untuk membuktikan bahwa $R_a$ konstan. Mengapa konstanta ini harus~$0$?   \ns
\end{proof}
\end{o001sourcehint}
\begin{o001answer}
Jika spektrum kosong, resolven bernilai-aljabar terdefinisi di seluruh~$\C$.
Setiap komposisinya dengan fungsional dual terbatas dan entire, sehingga nol
oleh Liouville dan perilaku di tak hingga; Hahn--Banach lalu memaksa resolven
sendiri nol, bertentangan dengan sifat invers.
\end{o001answer}
\begin{o001proof}
Andaikan $\sigma(a)=\emptyset$. Maka fungsi resolven
\[
 R_a(\lambda)=(a-\lambda\vc1)^{-1}
\]
terdefinisi dan analitik pada seluruh~$\C$. Untuk
$|\lambda|>\|a\|$, deret Neumann memberi
\[
 R_a(\lambda)
 =-\frac1\lambda\left(\vc1-\frac a\lambda\right)^{-1},
 \qquad
 \|R_a(\lambda)\|\leq\frac1{|\lambda|-\|a\|}.
\]
Jadi resolven menuju nol ketika $|\lambda|\to\infty$. Pada setiap cakram
kompak ia terbatas karena kontinu.

Ambil $\phi\in A^*$. Fungsi skalar $\phi\circ R_a$ entire dan terbatas di
seluruh bidang: di luar suatu cakram gunakan estimasi tadi, sedangkan di dalam
cakram gunakan kekompakan. Teorema Liouville membuatnya konstan, dan limit di
tak hingga menunjukkan bahwa konstanta itu nol. Jadi
$\phi(R_a(\lambda))=0$ untuk setiap $\phi\in A^*$ dan setiap~$\lambda$.
Teorema Hahn--Banach memisahkan titik-titik aljabar Banach, sehingga
$R_a(\lambda)=0$. Tetapi
$(a-\lambda\vc1)R_a(\lambda)=\vc1$, yang mustahil bila resolven nol.
Kontradiksi ini membuktikan bahwa $\sigma(a)$ tidak kosong.
\end{o001proof}
\end{o001readerwork}


\chapter{SOLUSI UNTUK SELURUH LATIHAN SUMBER}

Pernyataan latihan berikut tetap merupakan materi sumber John M. Erdman.
Solusi lengkap merupakan komponen asli terpisah berlisensi CC BY-SA 4.0,
ditulis dengan bantuan OpenAI Codex gpt-5.6-sol, Ultra atas arahan pengguna.
Solusi bukan tulisan Erdman dan tidak menyiratkan dukungan beliau atau
Portland State University.

% Bab: FAOA-2015-CH01
% Provenans: OpenAI Codex gpt-5.6-sol, Ultra, atas arahan pengguna
% Lisensi komponen solusi asli: CC BY-SA 4.0
% Non-endorsement: solusi ini bukan karya John M. Erdman dan tidak menyiratkan dukungannya.
% Jumlah solusi: 6
\markboth{SOLUSI O001 UNTUK BAB 1: ALJABAR LINEAR DAN TEOREMA SPEKTRAL}{SOLUSI O001 UNTUK BAB 1: ALJABAR LINEAR DAN TEOREMA SPEKTRAL}
\section*{Solusi O001 untuk Bab 1: Aljabar Linear dan Teorema Spektral}
\addcontentsline{toc}{section}{Solusi O001 untuk Bab 1: Aljabar Linear dan Teorema Spektral}
\noindent\textit{Catatan komponen.}
Keenam solusi berikut ditulis terpisah dari teks sumber, dilisensikan di bawah
CC BY-SA 4.0, dan berprovenans
\emph{OpenAI Codex gpt-5.6-sol, Ultra, atas arahan pengguna}.
Solusi-solusi ini bukan tulisan John M. Erdman dan tidak menyiratkan dukungannya.

% O001-SOLUTION-ID: O001-FAOA-2015-CH01-EX-001-SOLUTION
% SOURCE-EXERCISE-ID: FAOA-2015-CH01-NODE-0003
% STATEMENT-TARGET-SHA256: aebec53b01a9f71876a10649836fb04c2425385c4bf14115a19b2d06bb946b9e
\begin{o001solution}
  {O001-FAOA-2015-CH01-EX-001-SOLUTION}
  {FAOA-2015-CH01-NODE-0003}
  {aebec53b01a9f71876a10649836fb04c2425385c4bf14115a19b2d06bb946b9e}
\begin{o001statement}
Definisi \emph{ruang vektor} yang ditemukan dalam banyak buku dasar kurang lebih
sebagai berikut: suatu \emph{ruang vektor atas} $\K$ adalah himpunan $V$ beserta operasi
penjumlahan dan perkalian skalar yang memenuhi aksioma-aksioma berikut:
 \begin{enumerate}
  \item[(1)] jika $x$, $y \in V$, maka $x + y \in V$;
  \item[(2)] $(x + y) + z = x + (y + z)$ untuk setiap
 \index{asosiatif}%
$x$, $y$, $z \in V$ (asosiativitas);
  \item[(3)] terdapat $\vc 0 \in V$ sehingga $x + \vc 0 = x$ untuk setiap
 \index{identitas!aditif}%
$x \in V$ (keberadaan identitas aditif);
  \item[(4)] untuk setiap $x\in V$ terdapat $-x \in V$ sehingga
 \index{aditif!invers}%
 \index{invers!aditif}%
$x + (-x) = \vc 0$ (keberadaan invers aditif);
  \item[(5)] $x + y = y + x$ untuk setiap $x$, $y \in V$
 \index{komutatif}%
(komutativitas);
  \item[(6)] jika $\alpha \in \K$ dan $x \in V$, maka $\alpha x \in V$;
  \item[(7)] $\alpha (x + y) = \alpha x + \alpha y$ untuk setiap $\alpha \in \K$
dan setiap $x$, $y \in V$;
  \item[(8)] $(\alpha + \beta)x = \alpha x + \beta x$ untuk setiap $\alpha$,
$\beta \in \K$ dan setiap $x \in V$;
  \item[(9)] $(\alpha\beta)x = \alpha(\beta x)$ untuk setiap $\alpha$,
$\beta \in \K$ dan setiap $x \in V$; dan
  \item[(10)] $1\,x = x$ untuk setiap $x \in V$.
 \end{enumerate}
Buktikan bahwa definisi ini ekuivalen dengan definisi yang diberikan di atas dalam~\ref{00001}.
\end{o001statement}
\begin{o001answer}
Aksioma (1)--(5) tepat menyatakan bahwa $(V,+)$ adalah grup Abelian, sedangkan
aksioma (6)--(10) tepat menyatakan bahwa pemetaan
\[
 M\colon\K\longrightarrow\Hom(V),\qquad M(\alpha)(x)=\alpha x,
\]
merupakan homomorfisme gelanggang beridentitas. Jadi kedua definisi menentukan
struktur yang sama.
\end{o001answer}
\begin{o001proof}
Andaikan lebih dahulu aksioma (1)--(10) dipenuhi. Aksioma (1) memastikan bahwa
penjumlahan merupakan operasi pada $V$; aksioma (2)--(5), bersama keberadaan
operasi itu, memberi asosiativitas, identitas, invers, dan komutativitas.
Karena itu $(V,+)$ adalah grup Abelian.

Untuk $\alpha\in\K$, definisikan $M(\alpha)\colon V\to V$ dengan
$M(\alpha)(x)=\alpha x$. Aksioma (6) memastikan kodomainnya benar, dan
aksioma (7) memberi
\[
 M(\alpha)(x+y)=M(\alpha)x+M(\alpha)y.
\]
Jadi $M(\alpha)\in\Hom(V)$. Penjumlahan dalam $\Hom(V)$ didefinisikan titik
demi titik dan perkaliannya ialah komposisi. Oleh aksioma (8)--(10), untuk
setiap $\alpha,\beta\in\K$ berlaku
\[
 M(\alpha+\beta)=M(\alpha)+M(\beta),\qquad
 M(\alpha\beta)=M(\alpha)M(\beta),\qquad
 M(1)=I_V.
\]
Maka $M$ adalah homomorfisme gelanggang beridentitas seperti dalam
Definisi~\ref{00001}.

Sebaliknya, andaikan $(V,+)$ grup Abelian dan
$M\colon\K\to\Hom(V)$ homomorfisme gelanggang beridentitas. Sifat grup
Abelian memberi (1)--(5). Tetapkan $\alpha\in\K$ dan tulis
$\alpha x:=M(\alpha)(x)$. Karena $M(\alpha)$ adalah pemetaan $V\to V$, kita
memperoleh (6); karena $M(\alpha)$ homomorfisme grup, kita memperoleh (7).
Kesamaan
\[
 M(\alpha+\beta)=M(\alpha)+M(\beta),\quad
 M(\alpha\beta)=M(\alpha)\circ M(\beta),\quad
 M(1)=I_V
\]
masing-masing memberi (8), (9), dan (10) setelah diterapkan pada sebarang
$x\in V$. Kedua konstruksi tersebut saling membalik, sehingga definisinya
ekuivalen.
\end{o001proof}
\end{o001solution}

% O001-SOLUTION-ID: O001-FAOA-2015-CH01-EX-002-SOLUTION
% SOURCE-EXERCISE-ID: FAOA-2015-CH01-NODE-0058
% STATEMENT-TARGET-SHA256: 47a91f1c72eaedd58f83411463d0bd5444618c6107a05854fdc7924c856e0dbd
\begin{o001solution}
  {O001-FAOA-2015-CH01-EX-002-SOLUTION}
  {FAOA-2015-CH01-NODE-0058}
  {47a91f1c72eaedd58f83411463d0bd5444618c6107a05854fdc7924c856e0dbd}
\begin{o001statement}
Misalkan $T$ adalah operator pada $\R^3$ yang representasi matriksnya
    \[
      \begin{bmatrix}
               2  &  0  &  0 \\
              -1  &  3  &  2 \\
               1  & -1  &  0
      \end{bmatrix}.
    \]
  \begin{enumerate}
    \item Tentukan polinomial karakteristik dan polinomial minimal dari~$T$.
    \item Apa yang dapat disimpulkan dari bentuk polinomial minimal tersebut?
    \item Tentukan matriks yang mendiagonalkan~$T$. Apa bentuk diagonal $T$ yang dihasilkan oleh matriks ini?
    \item Tentukan (representasi matriks dari)~$\sqrt T$.
  \end{enumerate}
\end{o001statement}
\begin{o001answer}
Dengan $A=[T]$,
\[
 c_T(\lambda)=(\lambda-1)(\lambda-2)^2,\qquad
 m_T(\lambda)=(\lambda-1)(\lambda-2).
\]
Jadi $T$ dapat didiagonalkan. Salah satu matriks pendiagonal ialah
\[
 S=\begin{bmatrix}0&1&2\\-1&1&0\\1&0&1\end{bmatrix},
 \qquad S^{-1}AS=\diag(1,2,2).
\]
Salah satu akar kuadrat real yang diperoleh dengan memilih akar positif pada
kedua nilai eigen adalah
\[
 \sqrt A=(2-\sqrt2)I+(\sqrt2-1)A
 =\begin{bmatrix}
 \sqrt2&0&0\\
 1-\sqrt2&2\sqrt2-1&2\sqrt2-2\\
 \sqrt2-1&1-\sqrt2&2-\sqrt2
 \end{bmatrix}.
\]
\end{o001answer}
\begin{o001proof}
Ekspansi sepanjang baris pertama memberi
\[
\det(\lambda I-A)
 =(\lambda-2)
   \det\begin{bmatrix}\lambda-3&-2\\1&\lambda\end{bmatrix}
 =(\lambda-2)(\lambda^2-3\lambda+2)
 =(\lambda-1)(\lambda-2)^2.
\]
Penyelesaian sistem eigen menghasilkan
\[
\ker(A-I)=\spn\{(0,-1,1)\}
\]
dan
\[
\ker(A-2I)=\spn\{(1,1,0),(2,0,1)\}.
\]
Ketiga vektor yang ditampilkan bebas linear; memang matriks $S$ yang
menjadikan ketiganya sebagai kolom memiliki determinan $-1$. Jadi terdapat
basis eigen dan
$S^{-1}AS=\diag(1,2,2)$. Karena kedua nilai eigen benar-benar muncul,
polinomial minimal harus habis dibagi $(\lambda-1)(\lambda-2)$; karena $A$
dapat didiagonalkan, tidak diperlukan faktor berulang. Ini membuktikan rumus
untuk $m_T$ dan kesimpulan pada bagian kedua.

Misalkan $E_1$ dan $E_2$ proyeksi spektral untuk nilai eigen $1$ dan $2$.
Dari $I=E_1+E_2$ dan $A=E_1+2E_2$ diperoleh
$E_1=2I-A$ serta $E_2=A-I$. Karena $E_jE_k=0$ untuk $j\ne k$ dan
$E_j^2=E_j$, matriks
\[
 R:=E_1+\sqrt2\,E_2=(2-\sqrt2)I+(\sqrt2-1)A
\]
memenuhi
\[
 R^2=E_1+2E_2=A.
\]
Substitusi matriks $A$ memberi matriks yang dinyatakan dalam jawaban.
Perhitungan $R^2=A$ juga merupakan pemeriksaan langsung atas semua tanda.
\end{o001proof}
\end{o001solution}

% O001-SOLUTION-ID: O001-FAOA-2015-CH01-EX-003-SOLUTION
% SOURCE-EXERCISE-ID: FAOA-2015-CH01-NODE-0063
% STATEMENT-TARGET-SHA256: 36a90568d574839d10490f6cccf49001911a58d38d7d248542829133d7a33ae9
\begin{o001solution}
  {O001-FAOA-2015-CH01-EX-003-SOLUTION}
  {FAOA-2015-CH01-NODE-0063}
  {36a90568d574839d10490f6cccf49001911a58d38d7d248542829133d7a33ae9}
\begin{o001statement}
Misalkan $T$ adalah operator pada $\R^3$ yang representasi matriksnya
 \[ \begin{bmatrix}
        3  &  1  & -1 \\
        2  &  2  & -1 \\
        2  &  2  &  0
  \end{bmatrix}. \]

 \begin{enumerate}
  \item Tentukan polinomial karakteristik dan polinomial minimal dari~$T$.
  \item Tentukan ruang-ruang eigen dari~$T$.
  \item Tentukan bagian yang dapat didiagonalkan $D$ dan bagian nilpoten $N$ dari~$T$.
  \item Tentukan matriks yang mendiagonalkan~$D$. Apa bentuk diagonal
dari $D$ yang dihasilkan oleh matriks ini?
  \item Tunjukkan bahwa $D$ berkomutasi dengan~$N$.
 \end{enumerate}
\end{o001statement}
\begin{o001answer}
Dengan $A=[T]$,
\[
 c_T(\lambda)=(\lambda-1)(\lambda-2)^2,\qquad
 m_T(\lambda)=(\lambda-1)(\lambda-2)^2.
\]
Ruang eigennya ialah
\[
 \ker(A-I)=\spn\{(1,0,2)\},\qquad
 \ker(A-2I)=\spn\{(1,1,2)\}.
\]
Dekomposisi Jordan--Chevalley $A=D+N$ diberikan oleh
\[
D=\begin{bmatrix}1&1&0\\0&2&0\\-2&2&2\end{bmatrix},
\qquad
N=\begin{bmatrix}2&0&-1\\2&0&-1\\4&0&-2\end{bmatrix}.
\]
Matriks
\[
S=\begin{bmatrix}1&1&0\\0&1&0\\2&0&1\end{bmatrix}
\quad\text{memenuhi}\quad
S^{-1}DS=\diag(1,2,2),
\]
dan $DN=ND=\begin{bmatrix}4&0&-2\\4&0&-2\\8&0&-4\end{bmatrix}$.
\end{o001answer}
\begin{o001proof}
Perhitungan determinan memberi
\[
\det(\lambda I-A)=(\lambda-1)(\lambda-2)^2.
\]
Reduksi baris pada $A-I$ dan $A-2I$ masing-masing memberi ruang eigen yang
ditampilkan dalam jawaban. Khususnya, nilai eigen $2$ bermultiplisitas
aljabar dua tetapi ruang eigennya berdimensi satu. Oleh karena itu faktor
$(\lambda-2)$ dalam polinomial minimal harus berpangkat dua. Kedua nilai
eigen muncul, sehingga
$m_T(\lambda)=(\lambda-1)(\lambda-2)^2$.

Definisikan proyeksi ke ruang eigen untuk nilai $1$ sepanjang ruang eigen
tergeneralisasi untuk nilai $2$ dengan
\[
 E_1=(A-2I)^2
 =\begin{bmatrix}1&-1&0\\0&0&0\\2&-2&0\end{bmatrix}.
\]
Identitas polinomial minimal memberi $(A-I)E_1=0$, dan perkalian langsung
memberi $E_1^2=E_1$. Maka bagian yang dapat didiagonalkan adalah
\[
 D=1E_1+2(I-E_1)=2I-E_1
 =\begin{bmatrix}1&1&0\\0&2&0\\-2&2&2\end{bmatrix},
\]
dan
\[
 N=A-D
 =\begin{bmatrix}2&0&-1\\2&0&-1\\4&0&-2\end{bmatrix}.
\]
Perkalian matriks memberikan $N^2=0$, jadi $N$ nilpoten.

Kolom-kolom $S$ adalah $(1,0,2)$, $(1,1,0)$, dan $(0,0,1)$.
Kolom pertama merupakan vektor eigen $D$ untuk nilai $1$, sedangkan dua
kolom terakhir membentuk basis ruang eigen $D$ untuk nilai $2$.
Karena $\det S=1$, matriks itu invertibel dan menghasilkan bentuk diagonal
yang dinyatakan di atas. Terakhir, perkalian pada kedua urutan memberi
\[
 DN=ND=
 \begin{bmatrix}4&0&-2\\4&0&-2\\8&0&-4\end{bmatrix},
\]
yang memeriksa komutativitas secara langsung.
\end{o001proof}
\end{o001solution}

% O001-SOLUTION-ID: O001-FAOA-2015-CH01-EX-004-SOLUTION
% SOURCE-EXERCISE-ID: FAOA-2015-CH01-NODE-0076
% STATEMENT-TARGET-SHA256: f08e0ac70bab35e80596474dab5772dd1ab96f6eb0a60ed50265915217a9256b
\begin{o001solution}
  {O001-FAOA-2015-CH01-EX-004-SOLUTION}
  {FAOA-2015-CH01-NODE-0076}
  {f08e0ac70bab35e80596474dab5772dd1ab96f6eb0a60ed50265915217a9256b}
\begin{o001statement}
Misalkan $0 < a < b$.
 \begin{enumerate}
  \item Jika $f$ dan $g$ merupakan fungsi kontinu bernilai real pada
interval $[a,b]$, maka
   \[\left(\int_a^b f(x)g(x)\,dx\right)^2
             \le \int_a^b\bigl(f(x)\bigr)^2\,dx \cdot
                       \int_a^b\bigl(g(x)\bigr)^2\,dx.\]

\vskip 3 pt

  \item Gunakan bagian (a) untuk mencari bilangan $M$ dan $N$ (yang bergantung pada $a$
dan~$b$) sedemikian sehingga
   \[M(b-a) \le \ln\left(\frac ba\right) \le N(b-a).\]

\vskip 3 pt

  \item Gunakan bagian (b) untuk menunjukkan bahwa $2.5 \le e \le 3$.
 \end{enumerate}
\end{o001statement}
\begin{o001answer}
Bagian (a) adalah ketaksamaan Schwarz untuk hasil kali dalam integral. Pilihan
yang memenuhi bagian (b) ialah
\[
 M=\frac{2}{a+b},\qquad N=\frac1{\sqrt{ab}},
\]
sehingga
\[
\frac{2(b-a)}{a+b}\le\ln\frac ba
\le\frac{b-a}{\sqrt{ab}}.
\]
Dengan $(a,b)=(1,3)$ dan $(a,b)=(2,5)$ diperoleh
$\frac52<e\le3$, yang khususnya memberi $2.5\le e\le3$.
\end{o001answer}
\begin{o001proof}
Pada ruang fungsi kontinu bernilai real di $[a,b]$, tetapkan
\[
\langle f,g\rangle=\int_a^b f(x)g(x)\,dx.
\]
Ketaksamaan Schwarz memberi
$|\langle f,g\rangle|^2\le\langle f,f\rangle\langle g,g\rangle$.
Karena kuadrat bilangan real sama dengan kuadrat nilai mutlaknya, inilah
ketaksamaan pada bagian (a).

Mengikuti petunjuk pembuktian sumber, ambil
$f(x)=\sqrt{x}$ dan $g(x)=1/\sqrt{x}$. Bagian (a) memberi
\[
(b-a)^2
\le\left(\int_a^b x\,dx\right)\left(\int_a^b\frac{dx}{x}\right)
=\frac{(b-a)(a+b)}2\ln\frac ba.
\]
Karena $b-a>0$, pembagian menghasilkan
\[
\ln\frac ba\ge\frac{2(b-a)}{a+b}.
\]
Selanjutnya, dengan $f(x)=1/x$ dan $g(x)=1$, bagian (a) memberi
\[
\left(\ln\frac ba\right)^2
\le\left(\int_a^b\frac{dx}{x^2}\right)(b-a)
=\left(\frac1a-\frac1b\right)(b-a)
=\frac{(b-a)^2}{ab}.
\]
Semua suku yang diakarkan nonnegatif, sehingga
\[
\ln\frac ba\le\frac{b-a}{\sqrt{ab}}.
\]
Ini membuktikan pilihan $M$ dan $N$.

Untuk $(a,b)=(1,3)$, batas bawah memberi $\ln3\ge1$, sehingga, oleh
monotonisitas fungsi eksponensial, $3\ge e$. Untuk $(a,b)=(2,5)$, batas atas
memberi
\[
\ln\frac52\le\frac3{\sqrt{10}}<1
\]
karena $9<10$. Maka $\frac52<e$. Kedua hasil tersebut membuktikan batas yang
diminta.
\end{o001proof}
\end{o001solution}

% O001-SOLUTION-ID: O001-FAOA-2015-CH01-EX-005-SOLUTION
% SOURCE-EXERCISE-ID: FAOA-2015-CH01-NODE-0079
% STATEMENT-TARGET-SHA256: caeb07172c1e9193b4b6ba208c8dc34e8fea8474e2369b3a7187f64c418bd73b
\begin{o001solution}
  {O001-FAOA-2015-CH01-EX-005-SOLUTION}
  {FAOA-2015-CH01-NODE-0079}
  {caeb07172c1e9193b4b6ba208c8dc34e8fea8474e2369b3a7187f64c418bd73b}
\begin{o001statement}
Tunjukkan mengapa langsung jelas dari definisi bahwa $\norm{\vc 0} = 0$ dan bahwa norma
(dan seminorma) hanya dapat mengambil nilai nonnegatif.
\end{o001statement}
\begin{o001answer}
Homogenitas memberi $\norm{\vc0}=0$, dan ketaksamaan segitiga bersama
$\norm{-x}=\norm{x}$ memberi $0\le2\norm{x}$. Jadi setiap nilai norma maupun
seminorma nonnegatif.
\end{o001answer}
\begin{o001proof}
Untuk sebarang $x\in V$, homogenitas absolut menghasilkan
\[
\norm{\vc0}=\norm{0x}=|0|\,\norm{x}=0.
\]
Selain itu, $\norm{-x}=|-1|\,\norm{x}=\norm{x}$. Dengan menerapkan ketaksamaan
segitiga pada $x+(-x)=\vc0$, diperoleh
\[
0=\norm{\vc0}=\norm{x+(-x)}
\le\norm{x}+\norm{-x}=2\norm{x}.
\]
Jadi $\norm{x}\ge0$. Argumen ini hanya memakai homogenitas dan ketaksamaan
segitiga, sehingga berlaku tanpa perubahan untuk seminorma.
\end{o001proof}
\end{o001solution}

% O001-SOLUTION-ID: O001-FAOA-2015-CH01-EX-006-SOLUTION
% SOURCE-EXERCISE-ID: FAOA-2015-CH01-NODE-0125
% STATEMENT-TARGET-SHA256: ccf740d0df0eda4816be273a62026efa2cf1f2b433b990a40551ed07a9cc0f61
\begin{o001solution}
  {O001-FAOA-2015-CH01-EX-006-SOLUTION}
  {FAOA-2015-CH01-NODE-0125}
  {ccf740d0df0eda4816be273a62026efa2cf1f2b433b990a40551ed07a9cc0f61}
\begin{o001statement}
Misalkan $N = \dfrac13\begin{bmatrix}
                                4+2i & 1-i  & 1-i  \\
                                1-i  & 4+2i & 1-i  \\
                                1-i  & 1-i  & 4+2i \end{bmatrix}$.

\vskip 3 pt
 \begin{enumerate}
  \item Tunjukkan bahwa matriks $N$ bersifat normal.

\vskip 3 pt

  \item Dengan demikian, menurut \emph{teorema spektral}, $N$
dapat dituliskan sebagai kombinasi linear proyeksi-proyeksi ortogonal.
Tunjukkan secara eksplisit cara melakukannya (dan kerjakan perhitungannya).
 \end{enumerate}
\end{o001statement}
\begin{o001answer}
Definisikan
\[
P=\frac13\begin{bmatrix}1&1&1\\1&1&1\\1&1&1\end{bmatrix},
\qquad
Q=I-P=\frac13\begin{bmatrix}2&-1&-1\\-1&2&-1\\-1&-1&2\end{bmatrix}.
\]
Maka $P$ dan $Q$ adalah proyeksi ortogonal yang saling ortogonal,
$P+Q=I$, dan
\[
N=2P+(1+i)Q.
\]
Jadi $N$ normal. Nilai eigennya adalah $2$ pada
$\spn\{(1,1,1)\}$ dan $1+i$ pada
$\{(z_1,z_2,z_3):z_1+z_2+z_3=0\}$.
\end{o001answer}
\begin{o001proof}
Matriks $P$ memenuhi $P^*=P$ dan $P^2=P$, sehingga merupakan proyeksi
ortogonal ke $\spn\{(1,1,1)\}$. Karena $Q=I-P$, berlaku pula
$Q^*=Q$, $Q^2=Q$, serta
\[
PQ=QP=0,\qquad P+Q=I.
\]
Perhitungan entri pada
\[
2P+(1+i)Q
\]
memberi entri diagonal
$\frac23+\frac{2(1+i)}3=\frac{4+2i}{3}$ dan entri nondiagonal
$\frac23-\frac{1+i}{3}=\frac{1-i}{3}$. Jadi kombinasi itu persis matriks
$N$ dalam soal.

Karena $P$ dan $Q$ swaadjoin,
\[
N^*=2P+(1-i)Q.
\]
Dengan relasi proyeksi di atas,
\[
NN^*=4P+(1+i)(1-i)Q=4P+2Q=N^*N.
\]
Ini membuktikan kenormalan tanpa mengandaikannya. Persamaan
$N=2P+(1+i)Q$ sekaligus merupakan dekomposisi spektral yang diminta:
$P$ dan $Q$ adalah proyeksi-proyeksi ortogonal ke ruang eigen yang
bersesuaian.
\end{o001proof}
\end{o001solution}

% Bab: FAOA-2015-CH03
% Provenans: OpenAI Codex gpt-5.6-sol, Ultra, atas arahan pengguna
% Lisensi komponen solusi asli: CC BY-SA 4.0
% Non-endorsement: solusi ini bukan karya John M. Erdman dan tidak menyiratkan dukungannya.
% Jumlah solusi: 7
\markboth{SOLUSI O001 UNTUK BAB 3: RUANG LINEAR BERNORMA}{SOLUSI O001 UNTUK BAB 3: RUANG LINEAR BERNORMA}
\section*{Solusi O001 untuk Bab 3: Ruang Linear Bernorma}
\addcontentsline{toc}{section}{Solusi O001 untuk Bab 3: Ruang Linear Bernorma}
\noindent\textit{Catatan komponen.}
Ketujuh solusi berikut ditulis terpisah dari teks sumber, dilisensikan di bawah
CC BY-SA 4.0, dan berprovenans
\emph{OpenAI Codex gpt-5.6-sol, Ultra, atas arahan pengguna}.
Solusi-solusi ini bukan tulisan John M. Erdman dan tidak menyiratkan dukungannya.

% O001-SOLUTION-ID: O001-FAOA-2015-CH03-EX-001-SOLUTION
% SOURCE-EXERCISE-ID: FAOA-2015-CH03-NODE-0041
% STATEMENT-TARGET-SHA256: 1f693891073e0afb70f5764271328dd22974b4c3b33055ec494a61b7713f62e2
\begin{o001solution}
  {O001-FAOA-2015-CH03-EX-001-SOLUTION}
  {FAOA-2015-CH03-NODE-0041}
  {1f693891073e0afb70f5764271328dd22974b4c3b33055ec494a61b7713f62e2}
\begin{o001statement}
\label{C023191} Misalkan
 \index{l@$l_1$!sebagai ruang linear bernorma}%
 \index{l@$l_1$!barisan yang dapat dijumlahkan mutlak}%
 \index{bernorma!ruang linear!$l_1$ sebagai}%
$l_1$ merupakan himpunan semua barisan $x = (x_n)$ dari bilangan kompleks sedemikian
sehingga deret $\sum x_k$ konvergen mutlak. Jadikan $l_1$ ruang vektor dengan definisi
penjumlahan dan perkalian skalar titik demi titik yang biasa. Untuk setiap $x \in l_1$,
definisikan
 \index{<unopn@$\norm{\hphantom{x}}_1$ (norma-$1$)}%
 \index{<unopn@$\norm{\hphantom{x}}_\infty$ (norma seragam)}%
   \begin{align*}
      \norm{x}_1 &:= \sum_{k=1}^\infty \abs{x_k} \\
   \intertext{dan}
      \norm{x}_\infty &:= \sup\{\abs{x_k}\colon k \in \N\}.
   \end{align*}
(Masing-masing adalah
 \index{satu@norma-$1$!dari suatu barisan}%
 \index{norma!$1$- (dari suatu barisan)}%
norma-$1$ dan
 \index{seragam!norma!dari suatu barisan}%
\emph{norma seragam}.) Tunjukkan bahwa $\norm{\hphantom{x}}_1$ dan
$\norm{\hphantom{x}}_\infty$ keduanya merupakan norma pada~$l_1$. Kemudian buktikan atau
sangkal:
 \begin{enumerate}
  \item[(a)] Jika suatu barisan $(x^i)$ dari vektor-vektor dalam ruang linear bernorma
$\bigl(l_1, \norm{\hphantom{x}}_1\bigr)$ konvergen, maka barisan tersebut juga konvergen
dalam $\bigl(l_1, \norm{\hphantom{x}}_\infty\bigr)$.
  \item[(b)] Jika suatu barisan $(x^i)$ dari vektor-vektor dalam ruang linear bernorma
$\bigl(l_1, \norm{\hphantom{x}}_\infty\bigr)$ konvergen, maka barisan tersebut juga
konvergen dalam $\bigl(l_1, \norm{\hphantom{x}}_1\bigr)$.
 \end{enumerate}
\end{o001statement}
\begin{o001answer}
Kedua fungsi tersebut adalah norma pada $l_1$, dan
$\norm{x}_\infty\le\norm{x}_1$ untuk setiap $x\in l_1$. Karena itu (a) benar.
Pernyataan (b) salah: barisan
\[
x^i=(\underbrace{1/i,\ldots,1/i}_{i\ \text{suku}},0,0,\ldots)
\]
konvergen ke nol dalam norma seragam, tetapi
$\norm{x^i}_1=1$ untuk setiap $i$.
\end{o001answer}
\begin{o001proof}
Untuk $x\in l_1$, setiap $|x_k|$ tidak melebihi
$\sum_{j=1}^\infty|x_j|$, sehingga supremumnya berhingga. Jadi kedua rumus
memang bernilai real berhingga pada seluruh $l_1$.

Untuk norma-$1$, ketaknegatifan dan homogenitas mengikuti dari sifat nilai
mutlak. Jika $\norm{x}_1=0$, setiap suku $|x_k|$ bernilai nol, sehingga
$x=0$. Ketaksamaan segitiga titik demi titik dan penjumlahan deret
nonnegatif memberi
\[
\norm{x+y}_1
=\sum_{k=1}^\infty|x_k+y_k|
\le\sum_{k=1}^\infty(|x_k|+|y_k|)
=\norm{x}_1+\norm{y}_1.
\]
Maka $\norm{\hphantom{x}}_1$ adalah norma.

Untuk norma seragam, ketaknegatifan dan homogenitas juga mengikuti dari nilai
mutlak. Jika $\norm{x}_\infty=0$, maka $|x_k|=0$ untuk setiap $k$, sehingga
$x=0$. Selain itu,
\[
|x_k+y_k|\le|x_k|+|y_k|
\le\norm{x}_\infty+\norm{y}_\infty
\]
untuk setiap $k$; pengambilan supremum membuktikan ketaksamaan segitiga.
Jadi $\norm{\hphantom{x}}_\infty$ juga norma.

Karena
\[
\norm{z}_\infty=\sup_k|z_k|
\le\sum_{k=1}^\infty|z_k|=\norm{z}_1,
\]
konvergensi $x^i\to x$ dalam norma-$1$ mengakibatkan
$\norm{x^i-x}_\infty\le\norm{x^i-x}_1\to0$. Ini membuktikan (a).

Untuk vektor $x^i$ dalam jawaban, setiap vektor bertumpuan hingga, jadi
$x^i\in l_1$. Akan tetapi,
\[
\norm{x^i}_\infty=\frac1i\longrightarrow0,
\qquad
\norm{x^i}_1=i\frac1i=1.
\]
Dengan demikian $x^i\to0$ dalam norma seragam tetapi tidak dalam norma-$1$,
yang menyangkal (b).
\end{o001proof}
\end{o001solution}

% O001-SOLUTION-ID: O001-FAOA-2015-CH03-EX-002-SOLUTION
% SOURCE-EXERCISE-ID: FAOA-2015-CH03-NODE-0058
% STATEMENT-TARGET-SHA256: c063bcd29e442aa5670f59b396568ce8af222f0543fab7295967f54568271718
\begin{o001solution}
  {O001-FAOA-2015-CH03-EX-002-SOLUTION}
  {FAOA-2015-CH03-NODE-0058}
  {c063bcd29e442aa5670f59b396568ce8af222f0543fab7295967f54568271718}
\begin{o001statement}
Periksalah pernyataan-pernyataan yang belum dibuktikan dalam paragraf
sebelumnya. Secara khusus, buktikan bahwa
 \begin{enumerate}
    \item[(a)] $\cat{NLS_{\boldsymbol\infty}}$ merupakan suatu kategori.
    \item[(b)] Suatu isomorfisme dalam $\cat{NLS_{\boldsymbol\infty}}$ merupakan
isomorfisme ruang vektor sekaligus homeomorfisme.
    \item[(c)] Suatu pemetaan linear antara ruang-ruang linear bernorma merupakan isometri
jika dan hanya jika pemetaan tersebut mempertahankan norma.
    \item[(d)] $\cat{NLS_{\vc 1}}$ merupakan suatu kategori.
    \item[(e)] Suatu isomorfisme dalam $\cat{NLS_{\vc 1}}$ merupakan isometri sekaligus
isomorfisme ruang vektor.
    \item[(f)] Kategori $\cat{NLS_{\vc 1}}$ merupakan subkategori dari
$\cat{NLS_{\boldsymbol\infty}}$ dalam arti bahwa setiap morfisme kategori pertama
termasuk dalam kategori kedua.
 \end{enumerate}
\end{o001statement}
\begin{o001answer}
Pemetaan linear kontinu tertutup terhadap identitas dan komposisi, demikian
pula kontraksi linear. Isomorfisme dalam
$\cat{NLS_{\boldsymbol\infty}}$ adalah isomorfisme linear yang pemetaan dan
inversnya kontinu. Isomorfisme dalam $\cat{NLS_{\vc1}}$ adalah isometri
linear bijektif. Untuk pemetaan linear, sifat isometri ekuivalen dengan
pelestarian norma, dan setiap kontraksi linear merupakan pemetaan linear
kontinu.
\end{o001answer}
\begin{o001proof}
\begin{enumerate}
\item[(a)] Untuk setiap ruang linear bernorma $V$, identitas $I_V$ linear dan
kontinu. Jika $S\colon U\to V$ dan $T\colon V\to W$ linear kontinu, maka
$TS$ linear kontinu; dalam bentuk norma operator,
$\norm{TS}\le\norm{T}\norm{S}$. Asosiativitas komposisi dan hukum identitas
diwarisi dari komposisi fungsi. Semua aksioma kategori karena itu dipenuhi.

\item[(b)] Jika $T\colon V\to W$ isomorfisme dalam
$\cat{NLS_{\boldsymbol\infty}}$, terdapat morfisme
$S\colon W\to V$ dengan $ST=I_V$ dan $TS=I_W$. Maka $T$ bijektif,
$S=T^{-1}$, dan keduanya linear, jadi $T$ isomorfisme ruang vektor.
Karena $T$ dan $T^{-1}$ kontinu, $T$ juga homeomorfisme.

\item[(c)] Jika $T$ linear dan isometri, maka $T0=0$ dan
\[
\norm{Tx}=\norm{Tx-T0}=\norm{x-0}=\norm{x}.
\]
Sebaliknya, jika $T$ mempertahankan norma, linearitas memberi
\[
\norm{Tx-Ty}=\norm{T(x-y)}=\norm{x-y}
\]
untuk setiap $x,y$, sehingga $T$ isometri.

\item[(d)] Identitas memenuhi $\norm{I_Vx}=\norm{x}$, jadi merupakan
kontraksi. Jika $S$ dan $T$ kontraksi, maka
\[
\norm{TSx}\le\norm{Sx}\le\norm{x},
\]
sehingga komposisinya kontraksi. Asosiativitas dan hukum identitas kembali
diwarisi dari fungsi. Maka $\cat{NLS_{\vc1}}$ adalah kategori.

\item[(e)] Jika $T$ isomorfisme dalam $\cat{NLS_{\vc1}}$, baik $T$ maupun
$T^{-1}$ merupakan kontraksi. Jadi
\[
\norm{Tx}\le\norm{x}
=\norm{T^{-1}Tx}\le\norm{Tx}.
\]
Dengan demikian $\norm{Tx}=\norm{x}$ untuk setiap $x$. Bagian (c)
menunjukkan bahwa $T$ isometri; persamaan invers menunjukkan bahwa $T$
sekaligus isomorfisme ruang vektor.

\item[(f)] Jika $T$ kontraksi linear, maka
$\norm{Tx}\le\norm{x}$, jadi $T$ terbatas dengan $\norm T\le1$ dan karena
itu kontinu. Objek, identitas, dan komposisinya sama dengan yang diwarisi
dari $\cat{NLS_{\boldsymbol\infty}}$. Jadi setiap morfisme kategori
geometris merupakan morfisme kategori topologis, tepat dalam arti
subkategori yang dinyatakan dalam soal.
\end{enumerate}
\end{o001proof}
\end{o001solution}

% O001-SOLUTION-ID: O001-FAOA-2015-CH03-EX-003-SOLUTION
% SOURCE-EXERCISE-ID: FAOA-2015-CH03-NODE-0073
% STATEMENT-TARGET-SHA256: 730caa86f3c8bec4a40678e8832cce25c3d05b021a5fe0770bcff3a73f4a3ec2
\begin{o001solution}
  {O001-FAOA-2015-CH03-EX-003-SOLUTION}
  {FAOA-2015-CH03-NODE-0073}
  {730caa86f3c8bec4a40678e8832cce25c3d05b021a5fe0770bcff3a73f4a3ec2}
\begin{o001statement}
Verifikasilah pernyataan-pernyataan dalam Definisi~\ref{quo002def}. Secara khusus,
tunjukkan bahwa $\sim$ merupakan relasi ekuivalensi, bahwa penjumlahan dan perkalian skalar pada
himpunan kelas-kelas ekuivalensi terdefinisi dengan baik, bahwa terhadap operasi-operasi ini $V/M$ merupakan
ruang vektor, bahwa fungsi yang di sini disebut ``pemetaan hasil bagi'' memang merupakan pemetaan hasil bagi
dalam pengertian~\ref{quo001def}, dan bahwa pemetaan hasil bagi ini bersifat linear.
\end{o001statement}
\begin{o001answer}
Relasi $x\sim y\iff y-x\in M$ adalah relasi ekuivalensi. Rumus
$[x]+[y]=[x+y]$ dan $\alpha[x]=[\alpha x]$ tidak bergantung pada wakil,
menjadikan $V/M$ ruang vektor, dan pemetaan
$\pi(x)=[x]$ adalah pemetaan linear surjektif. Lebih lanjut, setiap fungsi
$h\colon V/M\to W$ yang membuat $h\pi$ linear harus linear; jadi $\pi$
adalah pemetaan hasil bagi dalam pengertian~\ref{quo001def}.
\end{o001answer}
\begin{o001proof}
Karena $M$ subruang, $0\in M$, tertutup terhadap invers aditif, dan tertutup
terhadap penjumlahan. Maka $x-x=0\in M$, sehingga $x\sim x$. Jika
$x\sim y$, maka $y-x\in M$ dan $x-y=-(y-x)\in M$, sehingga $y\sim x$.
Jika $x\sim y$ dan $y\sim z$, maka
$z-x=(z-y)+(y-x)\in M$, sehingga $x\sim z$. Jadi $\sim$ relasi
ekuivalensi.

Andaikan $[x]=[x']$ dan $[y]=[y']$. Berarti $x'-x,y'-y\in M$, sehingga
\[
(x'+y')-(x+y)=(x'-x)+(y'-y)\in M.
\]
Jadi $[x'+y']=[x+y]$. Untuk $\alpha\in\K$,
\[
\alpha x'-\alpha x=\alpha(x'-x)\in M,
\]
sehingga $[\alpha x']=[\alpha x]$. Kedua operasi dengan demikian
terdefinisi dengan baik.

Setiap aksioma ruang vektor sekarang dapat diperiksa pada wakilnya. Sebagai
contoh,
\[
([x]+[y])+[z]=[(x+y)+z]=[x+(y+z)]=[x]+([y]+[z]),
\]
$[0]$ adalah identitas aditif, $[-x]$ invers dari $[x]$, dan
\[
\alpha([x]+[y])=[\alpha(x+y)]
=[\alpha x]+[\alpha y].
\]
Komutativitas, aksioma skalar lainnya, dan identitas skalar mengikuti dengan
perhitungan wakil yang sama. Maka $V/M$ adalah ruang vektor.

Pemetaan $\pi\colon V\to V/M$, $\pi(x)=[x]$, surjektif karena setiap kelas
memiliki wakil. Selain itu,
\[
\pi(x+y)=[x+y]=[x]+[y],\qquad
\pi(\alpha x)=[\alpha x]=\alpha[x],
\]
jadi $\pi$ linear.

Terakhir, ambil sebarang ruang vektor $W$ dan fungsi
$h\colon V/M\to W$ sedemikian sehingga $h\circ\pi$ linear. Untuk semua
$[x],[y]\in V/M$ dan $\alpha\in\K$,
\[
\begin{aligned}
h([x]+[y])
 &=h(\pi(x+y))
  =(h\pi)(x+y)
  =(h\pi)(x)+(h\pi)(y)
  =h([x])+h([y]),\\
h(\alpha[x])
 &=h(\pi(\alpha x))
  =(h\pi)(\alpha x)
  =\alpha(h\pi)(x)
  =\alpha h([x]).
\end{aligned}
\]
Jadi $h$ linear. Ini persis sifat universal pada
Definisi~\ref{quo001def}, sehingga $\pi$ memang pemetaan hasil bagi.
\end{o001proof}
\end{o001solution}

% O001-SOLUTION-ID: O001-FAOA-2015-CH03-EX-004-SOLUTION
% SOURCE-EXERCISE-ID: FAOA-2015-CH03-NODE-0094
% STATEMENT-TARGET-SHA256: c2de90c5ffcfd0c7af44f29f75839337bc75b7c09430a8da07d70509a7db2ad1
\begin{o001solution}
  {O001-FAOA-2015-CH03-EX-004-SOLUTION}
  {FAOA-2015-CH03-NODE-0094}
  {c2de90c5ffcfd0c7af44f29f75839337bc75b7c09430a8da07d70509a7db2ad1}
\begin{o001statement}
Berikan bukti yang \emph{sangat} singkat (tanpa $\epsilon$, $\delta$, ataupun himpunan terbuka) bahwa jika
$(x_n)$ dan $(y_n)$ merupakan barisan-barisan dalam suatu ruang linear bernorma yang masing-masing konvergen ke $a$ dan $b$,
maka $x_n + y_n \sto a + b$.
\end{o001statement}
\begin{o001answer}
\[
\norm{(x_n+y_n)-(a+b)}
\le\norm{x_n-a}+\norm{y_n-b}\longrightarrow0.
\]
Jadi $x_n+y_n\to a+b$.
\end{o001answer}
\begin{o001proof}
Dengan mengelompokkan suku dan memakai ketaksamaan segitiga,
\[
\norm{(x_n+y_n)-(a+b)}
=\norm{(x_n-a)+(y_n-b)}
\le\norm{x_n-a}+\norm{y_n-b}.
\]
Kedua suku di ruas kanan menuju nol berdasarkan kedua hipotesis, sehingga
ruas kiri juga menuju nol. Ini adalah definisi konvergensi dalam norma dan
tidak memakai $\epsilon$, $\delta$, ataupun himpunan terbuka.
\end{o001proof}
\end{o001solution}

% O001-SOLUTION-ID: O001-FAOA-2015-CH03-EX-005-SOLUTION
% SOURCE-EXERCISE-ID: FAOA-2015-CH03-NODE-0118
% STATEMENT-TARGET-SHA256: a54a95af268dbf3cafc70d6ee756695be4fe8a45d3f1d175b15012cfa280f0d5
\begin{o001solution}
  {O001-FAOA-2015-CH03-EX-005-SOLUTION}
  {FAOA-2015-CH03-NODE-0118}
  {a54a95af268dbf3cafc70d6ee756695be4fe8a45d3f1d175b15012cfa280f0d5}
\begin{o001statement}
Apa saja koproduk dalam kategori $\cat{NLS_{\boldsymbol\infty}}$?
\end{o001statement}
\begin{o001answer}
Koproduk suatu keluarga ruang linear bernorma ada tepat ketika hanya
berhingga banyak anggotanya yang bukan ruang nol. Untuk
$V_1,\ldots,V_n$, satu modelnya ialah jumlah langsung aljabar
\[
\bigoplus_{k=1}^nV_k
\quad\text{dengan}\quad
\norm{(x_1,\ldots,x_n)}=\sum_{k=1}^n\norm{x_k},
\]
beserta injeksi koordinat. Koproduk keluarga kosong adalah ruang nol.
Jika terdapat tak hingga banyak objek bukan nol, koproduknya tidak ada dalam
kategori dengan semua pemetaan linear terbatas sebagai morfisme.
\end{o001answer}
\begin{o001proof}
Untuk keluarga hingga $V_1,\ldots,V_n$, injeksi koordinat
$\iota_k\colon V_k\to\bigoplus_{j=1}^nV_j$ merupakan pemetaan linear
terbatas. Jika $f_k\colon V_k\to W$ adalah pemetaan linear terbatas,
definisikan
\[
F(x_1,\ldots,x_n)=\sum_{k=1}^nf_k(x_k).
\]
Pemetaan ini satu-satunya pemetaan linear dengan $F\iota_k=f_k$, dan
\[
\norm{F(x_1,\ldots,x_n)}
\le\sum_{k=1}^n\norm{f_k}\norm{x_k}
\le\left(\max_{1\le k\le n}\norm{f_k}\right)
   \sum_{k=1}^n\norm{x_k}.
\]
Jadi $F$ terbatas dan sifat universal koproduk dipenuhi. Ruang nol tidak
mengubah konstruksi ini, sehingga keluarga dengan hanya berhingga banyak
anggota bukan nol mempunyai koproduk yang sama.

Sekarang andaikan keluarga $(V_\lambda)_{\lambda\in\Lambda}$ mempunyai tak
hingga banyak anggota bukan nol dan mempunyai koproduk hipotetis
$(C,(\iota_\lambda))$. Untuk setiap indeks $\lambda$ dengan
$V_\lambda\ne0$, gunakan keluarga morfisme ke $V_\lambda$ yang pada
$V_\lambda$ adalah identitas dan pada semua komponen lain adalah nol.
Sifat universal menghasilkan $p_\lambda\colon C\to V_\lambda$ dengan
$p_\lambda\iota_\lambda=I_{V_\lambda}$; khususnya
$\iota_\lambda\ne0$.

Pilih indeks-indeks berbeda $\lambda_1,\lambda_2,\ldots$ dengan
$V_{\lambda_n}\ne0$. Untuk setiap $n$, tetapkan
$f_{\lambda_n}=n\,\iota_{\lambda_n}\colon V_{\lambda_n}\to C$, dan tetapkan
$f_\lambda=0$ pada indeks lain. Masing-masing $f_\lambda$ adalah morfisme
terbatas. Sifat universal akan menghasilkan pemetaan linear terbatas
$F\colon C\to C$ dengan
\[
F\iota_{\lambda_n}=n\,\iota_{\lambda_n}.
\]
Ambil $x_n$ sehingga $\iota_{\lambda_n}x_n\ne0$. Maka
\[
\norm F\ge
\frac{\norm{F\iota_{\lambda_n}x_n}}
     {\norm{\iota_{\lambda_n}x_n}}
=n
\]
untuk setiap $n$, yang mustahil bagi operator terbatas. Jadi koproduk tak
ada bila tak hingga banyak komponen bukan nol.
\end{o001proof}
\end{o001solution}

% O001-SOLUTION-ID: O001-FAOA-2015-CH03-EX-006-SOLUTION
% SOURCE-EXERCISE-ID: FAOA-2015-CH03-NODE-0119
% STATEMENT-TARGET-SHA256: c18a12a34660fdaf55f44c9058084c0551cb4ad4560ec59d554f7a37201caf37
\begin{o001solution}
  {O001-FAOA-2015-CH03-EX-006-SOLUTION}
  {FAOA-2015-CH03-NODE-0119}
  {c18a12a34660fdaf55f44c9058084c0551cb4ad4560ec59d554f7a37201caf37}
\begin{o001statement}
Tentukan hasil kali dan koproduk dua ruang $V$ dan $W$ dalam kategori
$\cat{NLS_{\vc 1}}$ yang terdiri atas ruang-ruang linear bernorma dan kontraksi-kontraksi linear.
\end{o001statement}
\begin{o001answer}
Hasil kali adalah $V\times W$ dengan norma maksimum
\[
\norm{(v,w)}_{\mathrm{prod}}=\max\{\norm v,\norm w\}
\]
dan proyeksi koordinat. Koproduk adalah $V\oplus W$ dengan norma jumlah
\[
\norm{(v,w)}_{\mathrm{koprod}}=\norm v+\norm w
\]
dan injeksi koordinat.
\end{o001answer}
\begin{o001proof}
Pada norma maksimum, kedua proyeksi
$\pi_V(v,w)=v$ dan $\pi_W(v,w)=w$ adalah kontraksi. Jika
$f\colon X\to V$ dan $g\colon X\to W$ kontraksi, satu-satunya pemetaan
linear $H\colon X\to V\times W$ yang memenuhi
$\pi_VH=f$ dan $\pi_WH=g$ ialah $H(x)=(f(x),g(x))$. Pemetaan ini kontraksi
karena
\[
\norm{H(x)}_{\mathrm{prod}}
=\max\{\norm{f(x)},\norm{g(x)}\}\le\norm{x}.
\]
Jadi objek tersebut adalah hasil kali.

Pada norma jumlah, injeksi
$\iota_V(v)=(v,0)$ dan $\iota_W(w)=(0,w)$ adalah kontraksi, bahkan isometri.
Jika $f\colon V\to X$ dan $g\colon W\to X$ kontraksi, satu-satunya pemetaan
linear $K\colon V\oplus W\to X$ dengan
$K\iota_V=f$ dan $K\iota_W=g$ ialah
$K(v,w)=f(v)+g(w)$. Ketaksamaan
\[
\norm{K(v,w)}
\le\norm{f(v)}+\norm{g(w)}
\le\norm v+\norm w
=\norm{(v,w)}_{\mathrm{koprod}}
\]
menunjukkan bahwa $K$ kontraksi. Jadi objek kedua adalah koproduk.
\end{o001proof}
\end{o001solution}

% O001-SOLUTION-ID: O001-FAOA-2015-CH03-EX-007-SOLUTION
% SOURCE-EXERCISE-ID: FAOA-2015-CH03-NODE-0163
% STATEMENT-TARGET-SHA256: 9ceb84511f36b87a9eb590f586023123c303e4bc85ba4a4938930a5c6aca9926
\begin{o001solution}
  {O001-FAOA-2015-CH03-EX-007-SOLUTION}
  {FAOA-2015-CH03-NODE-0163}
  {9ceb84511f36b87a9eb590f586023123c303e4bc85ba4a4938930a5c6aca9926}
\begin{o001statement}
Kawan Anda, Fred R. Dimm, mengira bahwa ia telah menemukan pembuktian langsung yang sangat sederhana untuk
versi teorema Hahn--Banach yang diberikan dalam Teorema~\ref{HBTIII}: Suatu fungsional
linear kontinu $f$ yang didefinisikan pada subruang (linear) $M$ dari ruang linear bernorma~$V$ dapat
diperpanjang menjadi fungsional linear kontinu $g$ pada seluruh $V$ tanpa memperbesar
normanya. Menurutnya, kita hanya perlu mengambil $N$ sebagai sebarang komplemen ruang vektor
dari~$M$ dan mengambil $B$ sebagai suatu basis (Hamel) untuk $N$. Definisikan $g(e) = 0$ untuk setiap vektor $e
\in B$ dan $g = f$ pada~$M$. Kemudian perpanjang $g$ secara linear ke seluruh~$V$. Tentunya, kata Fred, $g$ linear berdasarkan konstruksinya dan normanya tidak bertambah karena kita
telah menjadikannya nol di tempat-tempat yang sebelumnya belum didefinisikan. Tunjukkan kepada Fred letak kekeliruannya.
\end{o001statement}
\begin{o001answer}
Perpanjangan Fred adalah $g=f\circ P$, dengan $P$ proyeksi aljabar ke $M$
sepanjang komplemen $N$. Proyeksi itu tidak harus kontraktif, bahkan tidak
harus terbatas. Jadi menaruh nilai nol pada $N$ tidak mengendalikan norma
$g$. Sudah pada $\R^2$ dapat terjadi $\norm g>\norm f$.
\end{o001answer}
\begin{o001proof}
Ambil $V=\R^2$ dengan norma Euclid,
\[
M=\spn\{(1,0)\},\qquad
f(t,0)=t,\qquad
N=\spn\{(1,1)\}.
\]
Fungsional $f$ mempunyai norma $1$ pada $M$. Pilih basis Hamel
$B=\{(1,1)\}$ untuk $N$ dan lakukan persis konstruksi Fred: tetapkan
$g(1,1)=0$ dan $g=f$ pada $M$. Karena
\[
(x,y)=(x-y,0)+y(1,1),
\]
perpanjangan linearnya adalah
\[
g(x,y)=x-y.
\]
Norma fungsional ini pada ruang Euclid adalah
\[
\norm g=\sqrt{1^2+(-1)^2}=\sqrt2>1=\norm f.
\]
Sebagai pemeriksaan langsung, pada
$u=(1,-1)/\sqrt2$ berlaku $\norm u=1$ tetapi $|g(u)|=\sqrt2$.

Kesalahannya ialah bahwa vektor $m+n$ dapat mempunyai norma jauh lebih kecil
daripada norma komponen $m$. Menjadikan $g$ nol pada $n$ tidak menghilangkan
kemungkinan pembatalan geometris itu. Secara umum, jika
$P\colon M\oplus N\to M$ adalah proyeksi aljabar, konstruksi Fred memberi
$g=f\circ P$; suatu komplemen aljabar sebarang tidak menjamin
$\norm P\le1$, dan dalam ruang berdimensi tak hingga $P$ bahkan dapat tidak
terbatas. Karena itu konstruksi tersebut tidak menjamin kesamaan norma maupun
kekontinuan.
\end{o001proof}
\end{o001solution}

% Lapisan solusi O001 -- Bab 4 (Ruang Hilbert)
% 10 solusi asli terpisah; bukan tulisan John M. Erdman.
% Provenans model: OpenAI Codex gpt-5.6-sol, Ultra, atas arahan pengguna.
% Lisensi komponen solusi: CC BY-SA 4.0.
% Tidak menyiratkan dukungan John M. Erdman atau Portland State University.

\markboth{SOLUSI O001 UNTUK BAB 4: RUANG HILBERT}{SOLUSI O001 UNTUK BAB 4: RUANG HILBERT}
\section*{Solusi O001 untuk Bab 4: Ruang Hilbert}
\addcontentsline{toc}{section}{Solusi O001 untuk Bab 4: Ruang Hilbert}
\begin{center}
\small
Komponen solusi asli terpisah, CC BY-SA 4.0. Ditulis dengan bantuan
OpenAI Codex gpt-5.6-sol, Ultra, atas arahan pengguna. Pernyataan latihan
berasal dari edisi terjemahan karya John M. Erdman; solusi ini bukan tulisan
beliau dan tidak menyiratkan dukungan beliau atau Portland State University.
\end{center}

% O001-SOLUTION-ID: O001-FAOA-2015-CH04-EX-001-SOLUTION
% SOURCE-EXERCISE-ID: FAOA-2015-CH04-NODE-0023
% STATEMENT-TARGET-SHA256: a4f5325f6a535b95e804661fb3f2a11e0b9b8d5a613e45903338613c2f6e4c70
\begin{o001solution}
  {O001-FAOA-2015-CH04-EX-001-SOLUTION}
  {FAOA-2015-CH04-NODE-0023}
  {a4f5325f6a535b95e804661fb3f2a11e0b9b8d5a613e45903338613c2f6e4c70}
\begin{o001statement}
 Apakah jumlah langsung yang didefinisikan di atas (dengan morfisme yang sesuai) merupakan hasil kali dalam kategori $\cat{HSp}$
yang objeknya ruang Hilbert dan morfismenya pemetaan linear terbatas?  Apakah jumlah langsung itu merupakan koproduk?
\end{o001statement}
\begin{o001answer}
Jumlah langsung ortogonal suatu keluarga $(H_i)_{i\in I}$ merupakan hasil kali
sekaligus koproduk dalam $\cat{HSp}$ jika dan hanya jika hanya berhingga banyak
$H_i$ yang tak nol.
\end{o001answer}
\begin{o001proof}
Jika hanya $H_1,\ldots,H_m$ yang tak nol, proyeksi koordinat
$p_i\colon\bigoplus_iH_i\to H_i$
bersifat terbatas. Jika $f_i\colon K\to H_i$ terbatas, maka
\[
 f(x)=(f_1x,\ldots,f_mx)
\]
adalah satu-satunya pemetaan yang memenuhi $p_i f=f_i$, dan
$\norm{f}^2\leq\sum_i\norm{f_i}^2$. Jadi jumlah langsung berhingga adalah
hasil kali. Dengan inklusi koordinat $j_i$, untuk pemetaan terbatas
$g_i\colon H_i\to K$, rumus
$g(x_1,\ldots,x_m)=\sum_i g_i x_i$ memberikan satu-satunya morfisme dengan
$g j_i=g_i$. Jadi objek yang sama juga koproduk.

Sebaliknya, andaikan tak hingga banyak $H_i$ yang tak nol. Pilih indeks
berbeda $i_1,i_2,\ldots$ dan vektor satuan $u_n\in H_{i_n}$. Untuk sifat
hasil kali, ambil $K=\K$, definisikan
$f_{i_n}(\alpha)=\alpha u_n$, dan ambil $f_i=0$ pada indeks lain. Setiap
$f_i$ terbatas. Namun pemetaan yang dipaksakan oleh semua proyeksi harus
mengirim $1$ ke keluarga yang koordinat $i_n$-nya sama dengan~$u_n$ untuk
semua~$n$; keluarga itu tidak dapat dijumlahkuadratkan. Jadi morfisme
pemfaktor tidak ada.

Untuk sifat koproduk, definisikan
$g_{i_n}(x)=\langle x,u_n\rangle$ dan $g_i=0$ pada indeks lain. Setiap
$g_i$ terbatas. Jika morfisme pemfaktor $g$ ada, maka untuk
\[
 x_N=\frac1{\sqrt N}\sum_{n=1}^{N}j_{i_n}u_n
\]
berlaku $\|x_N\|=1$, tetapi $g(x_N)=\sqrt N$. Ini bertentangan dengan
keterbatasan~$g$. Maka kedua sifat universal gagal tepat ketika terdapat tak
hingga banyak summan tak nol.
\end{o001proof}
\end{o001solution}

% O001-SOLUTION-ID: O001-FAOA-2015-CH04-EX-002-SOLUTION
% SOURCE-EXERCISE-ID: FAOA-2015-CH04-NODE-0053
% STATEMENT-TARGET-SHA256: 22c06694e8a8ce55ab43731d60ce3c9656a0b7af895ba5cbb039796129395f13
\begin{o001solution}
  {O001-FAOA-2015-CH04-EX-002-SOLUTION}
  {FAOA-2015-CH04-NODE-0053}
  {22c06694e8a8ce55ab43731d60ce3c9656a0b7af895ba5cbb039796129395f13}
\begin{o001statement}
Suatu fungsi $f\colon \R \sto \C$ berbentuk
   \[ f(t) = a_0 + \sum_{k=1}^n (a_k \cos kt + b_k \sin kt) \]
dengan $a_0, \dots, a_n, b_1, \dots, b_n$ bilangan kompleks disebut
 \index{polinom trigonometri}%
 \index{polinom!trigonometri}%
\df{polinom trigonometri}.
 \begin{enumerate}
  \item Jelaskan mengapa fungsi bernilai kompleks yang $2\pi$-periodik pada garis real $\R$
sering diidentifikasi dengan fungsi bernilai kompleks pada lingkaran satuan~$\T$.
  \item Beberapa penulis menyebut \emph{polinom trigonometri} sebagai fungsi berbentuk
    \[ f(t) = \sum_{k=-n}^n c_k e^{ikt} \]
 dengan $c_{-n}, \dots, c_{-1}, c_0, c_1, \dots, c_n$ bilangan kompleks. Jelaskan dengan cermat mengapa bentuk ini persis
sama dengan definisi sebelumnya.
  \item Jelaskan alasan penggunaan istilah \emph{polinom trigonometri}.
  \item Buktikan bahwa setiap fungsi kontinu yang $2\pi$-periodik pada $\R$ merupakan limit seragam
dari suatu barisan polinom trigonometri.
 \end{enumerate}
\end{o001statement}
\begin{o001answer}
(1) Identifikasi itu melalui $t\mapsto e^{it}$ dan ruang hasil bagi
$\R/(2\pi\Z)\cong\T$. (2) Kedua bentuk terkait oleh
\[
 c_0=a_0,\qquad c_k=\frac{a_k-i b_k}{2},\qquad
 c_{-k}=\frac{a_k+i b_k}{2}\quad(k\geq1).
\]
(3) Bentuk eksponensial merupakan polinom Laurent pada $z=e^{it}$.
(4) Jumlah Fej\'er dari deret Fourier $f$ adalah polinom trigonometri dan
konvergen seragam ke~$f$.
\end{o001answer}
\begin{o001proof}
Pemetaan $q\colon\R\to\T$, $q(t)=e^{it}$, bernilai sama tepat pada kelas
kongruensi modulo $2\pi$. Karena itu setiap fungsi $2\pi$-periodik mempunyai
faktorisasi tunggal melalui $\R/(2\pi\Z)$, yang secara alami diidentifikasi
dengan~$\T$.

Identitas Euler memberikan
\[
 \cos kt=\frac{e^{ikt}+e^{-ikt}}2,
 \qquad
 \sin kt=\frac{e^{ikt}-e^{-ikt}}{2i}.
\]
Substitusi memberi rumus koefisien dalam jawaban. Sebaliknya,
$a_k=c_k+c_{-k}$ dan $b_k=i(c_k-c_{-k})$, sehingga kedua keluarga fungsi
persis sama. Dengan $z=e^{it}$, bentuk itu adalah pembatasan
$\sum_{k=-n}^n c_kz^k$ pada $\abs z=1$; inilah alasan istilah
``polinom trigonometri''.

Tinggal dibuktikan pernyataan aproksimasi. Untuk $f$ kontinu dan
$2\pi$-periodik, tetapkan
\[
 \widehat f(k)=\frac1{2\pi}\int_{-\pi}^{\pi}f(s)e^{-iks}\,ds
\]
dan jumlah Fej\'er
\[
 \sigma_N f(t)=\sum_{k=-N}^{N}\left(1-\frac{\abs k}{N+1}\right)
                 \widehat f(k)e^{ikt}.
\]
Ini polinom trigonometri. Perkalian hingga dan perubahan variabel langsung
memberi
\[
 \sigma_N f(t)=\frac1{2\pi}\int_{-\pi}^{\pi}F_N(u)f(t-u)\,du,
 \qquad
 F_N(u)=\frac1{N+1}
 \left(\frac{\sin((N+1)u/2)}{\sin(u/2)}\right)^2,
\]
dengan nilai kontinu $F_N(0)=N+1$. Kernel ini tak negatif dan
$\frac1{2\pi}\int_{-\pi}^{\pi}F_N=1$. Maka
\[
 \abs{\sigma_N f(t)-f(t)}
 \leq\frac1{2\pi}\int_{-\pi}^{\pi}
 F_N(u)\abs{f(t-u)-f(t)}\,du.
\]
Untuk $\epsilon>0$, kekontinuan seragam $f$ memberi $\delta\in(0,\pi)$
sehingga faktor selisih kurang dari $\epsilon/2$ apabila $\abs u<\delta$.
Pada $\delta\leq\abs u\leq\pi$ berlaku
\[
 F_N(u)\leq\frac1{(N+1)\sin^2(\delta/2)}.
\]
Bagian integral tersebut jadi paling besar
$2\norm f_\infty/((N+1)\sin^2(\delta/2))$, terlepas dari~$t$, dan kurang
dari $\epsilon/2$ untuk $N$ cukup besar. Jadi
$\norm{\sigma_Nf-f}_\infty\to0$.
\end{o001proof}
\end{o001solution}

% O001-SOLUTION-ID: O001-FAOA-2015-CH04-EX-003-SOLUTION
% SOURCE-EXERCISE-ID: FAOA-2015-CH04-NODE-0078
% STATEMENT-TARGET-SHA256: 4f8e6fe2fc3fbbf3c5cd40c97d7a3c2295778e82b32add1bf6cf200560d0a1b4
\begin{o001solution}
  {O001-FAOA-2015-CH04-EX-003-SOLUTION}
  {FAOA-2015-CH04-NODE-0078}
  {4f8e6fe2fc3fbbf3c5cd40c97d7a3c2295778e82b32add1bf6cf200560d0a1b4}
\begin{o001statement}
Pada suatu sore, kawan Anda, Fred R. Dimm, datang membawa masalah. ``Begini,'' katanya,
``Pada $\lfs 2 = \lfs 2(\R,\R)$, fungsional evaluasi $E_0$, yang mengevaluasi setiap anggota $\lfs 2$
di~$0$, jelas merupakan fungsional linear terbatas. Jadi, menurut Teorema Representasi Riesz,
seharusnya ada fungsi $g$ dalam $\lfs 2$ sedemikian sehingga $f(0) = \langle f,g \rangle = \int_{-\infty}^\infty fg$
untuk semua $f$ dalam~$\lfs 2$. Namun, itu adalah sifat `fungsi~$\delta$', yang katanya
tidak ada. Di mana letak kesalahan saya?'' Berilah Freddy nasihat (yang membantu).
\end{o001statement}
\begin{o001answer}
$E_0$ bukan fungsional pada $\lfs2(\R)$: anggota $\lfs2$ adalah kelas
ekuivalensi hampir-di-mana-mana, sehingga nilai di satu titik tidak
terdefinisi. Bahkan pada fungsi kontinu terintegralkan-kuadrat, evaluasi di
nol tidak terbatas terhadap norma $\lfs2$.
\end{o001answer}
\begin{o001proof}
Jika dua fungsi hanya berbeda nilainya di nol, keduanya menentukan anggota
$\lfs2(\R)$ yang sama, tetapi ``evaluasi di nol'' akan memberi nilai berbeda.
Jadi $E_0$ tidak terdefinisi pada ruang itu.

Masalah tersebut tidak hilang dengan memilih fungsi kontinu sebagai wakil.
Untuk $n\geq1$, ambil
\[
 f_n(t)=\max\{1-n\abs t,0\}.
\]
Fungsi ini kontinu, $f_n(0)=1$, dan perhitungan langsung memberi
\[
 \norm{f_n}_2^2
 =2\int_0^{1/n}(1-nt)^2\,dt=\frac{2}{3n}\longrightarrow0.
\]
Jika evaluasi terbatas terhadap norma $\lfs2$, seharusnya
$1=\abs{E_0f_n}\leq\norm{E_0}\norm{f_n}_2\to0$, suatu kontradiksi.
Karena hipotesis Teorema Representasi Riesz gagal, tidak ada vektor
$g\in\lfs2$ yang merepresentasikan evaluasi titik tersebut.
\end{o001proof}
\end{o001solution}

% O001-SOLUTION-ID: O001-FAOA-2015-CH04-EX-004-SOLUTION
% SOURCE-EXERCISE-ID: FAOA-2015-CH04-NODE-0079
% STATEMENT-TARGET-SHA256: 7a531d47180bfa90ec00b09db0665bc412ad4270a85442b46f63ce4efa10ab53
\begin{o001solution}
  {O001-FAOA-2015-CH04-EX-004-SOLUTION}
  {FAOA-2015-CH04-NODE-0079}
  {7a531d47180bfa90ec00b09db0665bc412ad4270a85442b46f63ce4efa10ab53}
\begin{o001statement}
Masih pada sore yang sama, Freddy kembali. Setelah mempertimbangkan
nasihat yang Anda berikan (dalam soal sebelumnya), ia merevisi
pertanyaannya dengan menjadikan domain $E_0$ sebagai himpunan fungsi bernilai real yang terbatas dan
kontinu pada~$\R$. Dengan sabar Anda
menjelaskan kepada Freddy bahwa sekarang setidaknya ada dua kekeliruan dalam caranya
menggunakan `fungsi~$\delta$'. Apakah kedua kekeliruan tersebut?
\end{o001statement}
\begin{o001answer}
Dengan norma supremum, $C_b(\R)$ memang membuat $E_0$ terbatas, tetapi
$C_b(\R)$ bukan ruang Hilbert, sehingga teorema representasi Hilbert--Riesz
tidak berlaku. Jika ia justru bermaksud memakai hasil kali dalam
$\int fg$, hasil kali itu bahkan tidak terdefinisi pada semua $C_b(\R)$;
setelah domain dibatasi agar masuk $L_2$, evaluasi tetap tidak terbatas.
\end{o001answer}
\begin{o001proof}
Pada $C_b(\R)$ dengan norma supremum berlaku
$\abs{f(0)}\leq\norm f_\infty$, jadi evaluasi memang kontinu. Namun norma
supremum tidak berasal dari hasil kali dalam: misalnya fungsi konstan
$f=1$ dan $g$ dengan nilai dalam $[0,1]$, $g(0)=1$, dan tumpuan kompak
dapat dipilih sehingga
\[
 \norm{f+g}_\infty^2+\norm{f-g}_\infty^2=4+1\neq4
 =2\norm f_\infty^2+2\norm g_\infty^2.
\]
Jadi hukum jajargenjang gagal dan Teorema Representasi Riesz untuk ruang
Hilbert tidak dapat diterapkan.

Di sisi lain, rumus $\langle f,g\rangle=\int_\R fg$ tidak menentukan hasil
kali dalam pada seluruh $C_b(\R)$: fungsi konstan $1$ saja tidak berada
dalam $L_2(\R)$. Jika domain diganti dengan $C_b(\R)\cap L_2(\R)$ dan norma
$L_2$, barisan tenda $f_n$ dari solusi sebelumnya masih memenuhi
$f_n(0)=1$ dan $\norm{f_n}_2\to0$. Maka evaluasi tidak terbatas. Evaluasi
titik secara benar direpresentasikan oleh ukuran titik Dirac dalam teori
representasi untuk fungsional pada ruang fungsi kontinu, bukan oleh suatu
fungsi $L_2$ dalam argumen Freddy.
\end{o001proof}
\end{o001solution}

% O001-SOLUTION-ID: O001-FAOA-2015-CH04-EX-005-SOLUTION
% SOURCE-EXERCISE-ID: FAOA-2015-CH04-NODE-0080
% STATEMENT-TARGET-SHA256: 40661e351b7cad48baf7c9d8e4b91c629468788e4e93a6722bd384dc6146a546
\begin{o001solution}
  {O001-FAOA-2015-CH04-EX-005-SOLUTION}
  {FAOA-2015-CH04-NODE-0080}
  {40661e351b7cad48baf7c9d8e4b91c629468788e4e93a6722bd384dc6146a546}
\begin{o001statement}
Mimpi buruk itu berlanjut. Sekarang pukul 23.00 dan Freddy kembali lagi.
Kali ini (setelah, untungnya, menyerah soal fungsi~$\delta$) ia ingin
menerapkan teorema representasi pada fungsional ``evaluasi
pada nol'' pada ruang~$l_2(\Z)$. Dapatkah ia melakukannya? Jelaskan.
\end{o001statement}
\begin{o001answer}
Ya. Evaluasi koordinat nol adalah fungsional terbatas bernorma satu pada
$l_2(\Z)$ dan direpresentasikan oleh vektor basis standar $e^0$:
$E_0(x)=x_0=\langle x,e^0\rangle$.
\end{o001answer}
\begin{o001proof}
Berbeda dari $L_2(\R)$, anggota $l_2(\Z)$ benar-benar merupakan barisan
dengan koordinat yang terdefinisi. Untuk setiap $x=(x_k)_{k\in\Z}$,
\[
 \abs{x_0}^2\leq\sum_{k\in\Z}\abs{x_k}^2=\norm x_2^2,
\]
sehingga $E_0$ linear dan $\norm{E_0}\leq1$. Pada $e^0$ berlaku
$\norm{e^0}=1$ dan $E_0(e^0)=1$, jadi normanya tepat satu. Dengan konvensi
hasil kali dalam edisi ini,
\[
 \langle x,e^0\rangle=\sum_{k\in\Z}x_k\conj{e^0_k}=x_0.
\]
Jadi semua hipotesis teorema representasi terpenuhi dan representernya
adalah~$e^0$.
\end{o001proof}
\end{o001solution}

% O001-SOLUTION-ID: O001-FAOA-2015-CH04-EX-006-SOLUTION
% SOURCE-EXERCISE-ID: FAOA-2015-CH04-NODE-0091
% STATEMENT-TARGET-SHA256: 986b9bb676932d7ef63f7ab92f49a8fa51ab6068b73b4daedb6c0c01a7c2ec38
\begin{o001solution}
  {O001-FAOA-2015-CH04-EX-006-SOLUTION}
  {FAOA-2015-CH04-NODE-0091}
  {986b9bb676932d7ef63f7ab92f49a8fa51ab6068b73b4daedb6c0c01a7c2ec38}
\begin{o001statement}
Jelaskan mengapa topologi pada ruang Hilbert yang dibangkitkan oleh konvergensi lemah jaring sesungguhnya merupakan
topologi lemah dalam pengertian topologis yang lazim.
\end{o001statement}
\begin{o001answer}
Topologi itu tepat sama dengan topologi awal
$\sigma(H,H^*)$: topologi terkecil yang membuat setiap fungsional linear
terbatas kontinu.
\end{o001answer}
\begin{o001proof}
Untuk setiap $y\in H$, definisikan $\phi_y(x)=\langle x,y\rangle$. Teorema
Representasi Riesz menyatakan bahwa keluarga $\{\phi_y:y\in H\}$ adalah
seluruh dual kontinu~$H^*$. Topologi awal yang ditentukan keluarga ini
mempunyai basis lingkungan di $a$ berupa
\[
 U(a;y_1,\ldots,y_m;\epsilon)
 =\{x:\abs{\langle x-a,y_j\rangle}<\epsilon,
             \ 1\leq j\leq m\}.
\]
Menurut sifat umum topologi awal, suatu jaring $(x_\lambda)$ konvergen ke
$a$ dalam topologi tersebut jika dan hanya jika
$\phi_y(x_\lambda)\to\phi_y(a)$ untuk setiap $y\in H$. Kondisi ini persis
$\langle x_\lambda,y\rangle\to\langle a,y\rangle$ untuk setiap~$y$, yaitu
definisi konvergensi lemah dalam soal. Karena konvergensi jaring menentukan
topologi secara tunggal, topologi yang dibangkitkan oleh konvergensi itu
adalah $\sigma(H,H^*)$, topologi lemah dalam arti topologis.
\end{o001proof}
\end{o001solution}

% O001-SOLUTION-ID: O001-FAOA-2015-CH04-EX-007-SOLUTION
% SOURCE-EXERCISE-ID: FAOA-2015-CH04-NODE-0092
% STATEMENT-TARGET-SHA256: 848170d90169079965a64c432bbc9a6b1bfaf0ab895ba7558340a4a60524926c
\begin{o001solution}
  {O001-FAOA-2015-CH04-EX-007-SOLUTION}
  {FAOA-2015-CH04-NODE-0092}
  {848170d90169079965a64c432bbc9a6b1bfaf0ab895ba7558340a4a60524926c}
\begin{o001statement}
Misalkan $H$ ruang Hilbert dan $(x_{{}_{\sst\lambda}})$ jaring dalam $H$.
 \begin{enumerate}
  \item Jika $x_{{}_{\sst\lambda}} \to^s  a$ dalam $H$, maka $x_{{}_{\sst\lambda}} \to^w a$ dalam $H$.
  \item Topologi kuat (norma) pada $H$ lebih kuat (lebih besar) daripada topologi lemah.
  \item Tunjukkan dengan contoh bahwa kebalikan dari (a) tidak berlaku.
  \item Apakah norma pada $H$ kontinu secara lemah?
  \item Jika $x_{{}_{\sst\lambda}} \to^w  a$ dalam $H$ dan jika $\norm{x_{{}_{\sst\lambda}}} \to \norm{a}$, maka
$x_{{}_{\sst\lambda}}\to^s a$ dalam~$H$.
  \item Jika pemetaan linear $T\colon H \sto H$ kontinu (secara kuat), maka pemetaan tersebut kontinu secara lemah.
 \end{enumerate}
\end{o001statement}
\begin{o001answer}
(1), (2), (5), dan (6) benar. Untuk (3), basis standar $(e^n)$ di $l_2$
konvergen lemah ke nol tetapi tidak dalam norma. Untuk (4), norma kontinu
secara lemah tepat ketika $H$ berdimensi hingga.
\end{o001answer}
\begin{o001proof}
(1) Dari ketaksamaan Cauchy--Schwarz,
\[
 \abs{\langle x_\lambda-a,y\rangle}
 \leq\norm{x_\lambda-a}\norm y\longrightarrow0
\]
untuk setiap $y\in H$. (2) Jadi setiap jaring yang konvergen dalam topologi
norma juga konvergen lemah; ekuivalen dengan mengatakan bahwa topologi norma
lebih kuat.

(3) Dalam $l_2$, untuk setiap $y=(y_n)\in l_2$ berlaku
$\langle e^n,y\rangle=\conj{y_n}\to0$, sebab koordinat setiap barisan
$l_2$ menuju nol. Namun $\norm{e^n}=1$, sehingga $e^n$ tidak konvergen dalam
norma ke nol.

(4) Contoh yang sama menunjukkan bahwa, jika $H$ berdimensi tak hingga,
norma tidak kontinu lemah di nol: pilih suatu barisan ortonormal $(e^n)$;
maka $e^n\to^w0$ tetapi $\norm{e^n}=1\not\to0$. Jika $H$ berdimensi hingga,
ambil basis ortonormal $e^1,\ldots,e^m$. Konvergensi lemah berarti konvergensi
setiap koordinat, dan
$\norm x^2=\sum_{j=1}^m\abs{\langle x,e^j\rangle}^2$; jadi topologi lemah
sama dengan topologi norma dan fungsi norma kontinu.

(5) Karena hasil kali dalam kontinu pada setiap koordinat lemah yang tetap,
\[
 \norm{x_\lambda-a}^2
 =\norm{x_\lambda}^2+\norm a^2
  -2\Re\langle x_\lambda,a\rangle\longrightarrow0.
\]
(6) Untuk $y\in H$, pemetaan
$x\mapsto\langle Tx,y\rangle$ adalah fungsional linear terbatas, sebab nilai
mutlaknya paling besar $\norm T\norm y\norm x$. Oleh Riesz, ada $z_y\in H$
sehingga $\langle Tx,y\rangle=\langle x,z_y\rangle$. Jika
$x_\lambda\to^w a$, maka
$\langle Tx_\lambda,y\rangle\to\langle Ta,y\rangle$ untuk setiap~$y$;
jadi $Tx_\lambda\to^wTa$. Kriteria jaring untuk kontinuitas menyelesaikan
pembuktian.
\end{o001proof}
\end{o001solution}

% O001-SOLUTION-ID: O001-FAOA-2015-CH04-EX-008-SOLUTION
% SOURCE-EXERCISE-ID: FAOA-2015-CH04-NODE-0093
% STATEMENT-TARGET-SHA256: 4c8d88aed6251f2083c6029080a78a9b8a175cd792a4d2e114382a46518d4f4f
\begin{o001solution}
  {O001-FAOA-2015-CH04-EX-008-SOLUTION}
  {FAOA-2015-CH04-NODE-0093}
  {4c8d88aed6251f2083c6029080a78a9b8a175cd792a4d2e114382a46518d4f4f}
\begin{o001statement}
Misalkan $H$ ruang Hilbert berdimensi tak hingga. Tinjau subhimpunan-subhimpunan berikut dari~$H$:
 \begin{enumerate}[label=\rm{(\alph*)}]
  \item bola satuan terbuka;
  \item bola satuan tertutup;
  \item sfer satuan;
  \item subruang linear tertutup.
 \end{enumerate}
Tentukan untuk setiap himpunan tersebut apakah masing-masing bersifat: tertutup secara kuat, tertutup secara lemah,
terbuka secara kuat, terbuka secara lemah, kompak secara kuat, kompak secara lemah. (Salah satu bagian
terlalu sulit untuk saat ini: \emph{Teorema Alaoglu}~\futurexref{6.2.9}{C067441} akan menunjukkan bahwa
bola satuan tertutup kompak secara lemah.)
\end{o001statement}
\begin{o001answer}
Dengan urutan sifat ``tertutup kuat, tertutup lemah, terbuka kuat, terbuka
lemah, kompak kuat, kompak lemah'', jawabannya adalah
\[
\begin{array}{c|cccccc}
 &\mathrm{TK}&\mathrm{TL}&\mathrm{BK}&\mathrm{BL}&\mathrm{KK}&\mathrm{KL}\\ \hline
\{\norm x<1\}&\text{tidak}&\text{tidak}&\text{ya}&\text{tidak}&\text{tidak}&\text{tidak}\\
\{\norm x\leq1\}&\text{ya}&\text{ya}&\text{tidak}&\text{tidak}&\text{tidak}&\text{ya}\\
\{\norm x=1\}&\text{ya}&\text{tidak}&\text{tidak}&\text{tidak}&\text{tidak}&\text{tidak}.
\end{array}
\]
Untuk subruang linear tertutup $M$: ia tertutup kuat dan lemah; ia terbuka
kuat maupun lemah tepat jika $M=H$; dan ia kompak kuat maupun lemah tepat
jika $M=\{0\}$.
\end{o001answer}
\begin{o001proof}
Kontinuitas norma memberi keterbukaan bola terbuka serta ketertutupan bola
tertutup dan sfer dalam topologi kuat. Ketiganya tidak mempunyai sifat kuat
lain yang tercantum: bola terbuka tidak tertutup; dua himpunan lainnya tidak
terbuka; dan tidak satu pun kompak kuat. Untuk klaim terakhir, suatu barisan
ortonormal terletak di bola tertutup dan sfer serta mempunyai jarak
$\sqrt2$ antarsuku. Kelipatan tetap, misalnya $e^n/2$, memberi argumen yang
sama di dalam bola terbuka.

Bola tertutup bersifat tertutup lemah karena
\[
 \{x:\norm x\leq1\}
 =\bigcap_{\norm y\leq1}\{x:\abs{\langle x,y\rangle}\leq1\},
\]
suatu irisan himpunan tertutup lemah. Ia juga kompak lemah. Berikut bukti
lokal klaim terakhir. Petakan bola itu ke
$\prod_{y\in H}\{z\in\C:\abs z\leq\norm y\}$ melalui
$x\mapsto(\langle x,y\rangle)_y$. Produk cakram kompak ini kompak menurut
teorema kekompakan hasil kali. Setiap titik tutupan koordinat menentukan
fungsional konjugat-linear $\Phi(y)$ yang memenuhi
$\abs{\Phi(y)}\leq\norm y$ dan mempertahankan penjumlahan serta perkalian
skalar; sifat-sifat itu adalah persamaan tertutup dalam produk. Teorema
Representasi Riesz memberi $x\in H$, $\norm x\leq1$, dengan
$\Phi(y)=\langle x,y\rangle$. Jadi citra bola tertutup sendiri tertutup di
dalam produk kompak dan karenanya kompak; topologi yang diwarisi tepat
topologi lemah.

Setiap lingkungan lemah dasar hanya membatasi berhingga banyak hasil kali
dalam. Karena $H$ berdimensi tak hingga, terdapat vektor satuan $v$ yang
ortogonal terhadap semua vektor penguji itu. Maka lingkungan tersebut
memuat $x+tv$ untuk setiap skalar~$t$. Akibatnya bola terbuka, bola tertutup,
dan sfer tidak terbuka lemah. Barisan ortonormal memenuhi $e^n\to^w0$;
karena itu bola terbuka tidak tertutup lemah (kelipatan
$(1-1/n)e^1$ juga memberi limit batas), dan sfer tidak tertutup lemah.
Dalam ruang Hausdorff, himpunan kompak harus tertutup; jadi bola terbuka dan
sfer tidak kompak lemah.

Terakhir, untuk subruang tertutup $M$, teorema proyeksi memberi
$M=M^{\perp\perp}$. Dengan demikian
\[
 M=\bigcap_{y\in M^\perp}\{x:\langle x,y\rangle=0\},
\]
sehingga $M$ tertutup lemah. Subruang yang mempunyai interior dalam suatu
topologi ruang vektor harus sama dengan seluruh ruang; jadi $M$ terbuka
dalam salah satu topologi tepat jika $M=H$. Jika $M$ memuat $u\neq0$, fungsi
lemah-kontinu $x\mapsto\langle x,u\rangle$ tidak terbatas pada
$\{tu:t\in\K\}\subset M$, sedangkan citra kontinu himpunan kompak harus
kompak dan terbatas. Jadi $M$ tidak kompak lemah, dan terlebih lagi tidak
kompak kuat. Subruang nol adalah kompak dalam kedua topologi.
\end{o001proof}
\end{o001solution}

% O001-SOLUTION-ID: O001-FAOA-2015-CH04-EX-009-SOLUTION
% SOURCE-EXERCISE-ID: FAOA-2015-CH04-NODE-0094
% STATEMENT-TARGET-SHA256: 6ac5283840818281ac10974dbec5da9b992478b82ef530203afa9f84a1cfcb77
\begin{o001solution}
  {O001-FAOA-2015-CH04-EX-009-SOLUTION}
  {FAOA-2015-CH04-NODE-0094}
  {6ac5283840818281ac10974dbec5da9b992478b82ef530203afa9f84a1cfcb77}
\begin{o001statement}
Misalkan $\{e^n\colon n \in \N\}$ basis biasa untuk $l_2$. Untuk
setiap $n \in \N$, misalkan $a_n = \sqrt n\, e^n$. Manakah di antara pernyataan berikut
yang benar?
 \begin{enumerate}
  \item $0$ adalah titik akumulasi lemah dari himpunan
$\{a_n\colon n \in \N\}$.
  \item $0$ adalah titik gugus lemah dari barisan~$(a_n)$.
  \item $0$ adalah limit lemah dari suatu subbarisan dari~$(a_n)$.
 \end{enumerate}
\end{o001statement}
\begin{o001answer}
Pernyataan (1) dan (2) benar, sedangkan (3) salah.
\end{o001answer}
\begin{o001proof}
Ambil lingkungan lemah dasar nol yang ditentukan oleh
$y^1,\ldots,y^m\in l_2$ dan $\epsilon>0$. Tetapkan
$s_n=\sum_{j=1}^m\abs{y^j_n}^2$. Karena $\sum_ns_n<\infty$, untuk setiap
$N$ ada $n\geq N$ dengan $ns_n<\epsilon^2$: jika tidak, deret ekor
$\sum_ns_n$ mendominasi kelipatan deret harmonik. Untuk indeks itu,
\[
 \abs{\langle a_n,y^j\rangle}
 =\sqrt n\,\abs{y^j_n}\leq\sqrt{ns_n}<\epsilon
 \qquad(1\leq j\leq m).
\]
Jadi setiap lingkungan lemah nol memuat $a_n$ untuk indeks yang arbitrer
besar. Karena $0$ bukan salah satu $a_n$, nol adalah titik akumulasi lemah
himpunannya, dan karena kunjungan terjadi setelah setiap indeks, nol juga
titik gugus lemah barisan.

Sekarang ambil sebarang subbarisan $(a_{n_k})$. Jika ia konvergen lemah ke
$x$, maka untuk setiap koordinat tetap $r$, akhirnya $n_k\neq r$, sehingga
$\langle a_{n_k},e^r\rangle=0$ dan $\langle x,e^r\rangle=0$. Kelengkapan
basis memaksa $x=0$. Pilih lagi subsubbarisan dengan
$n_{k_j}\geq2^j$, dan definisikan $y\in l_2$ dengan
$y_{n_{k_j}}=n_{k_j}^{-1/2}$ serta koordinat lain nol. Memang
$\sum_j\abs{y_{n_{k_j}}}^2\leq\sum_j2^{-j}<\infty$, tetapi
\[
 \langle a_{n_{k_j}},y\rangle=1
\]
untuk setiap~$j$. Subsubbarisan itu tidak konvergen lemah ke nol, sehingga
subbarisan semula pun tidak mungkin konvergen lemah. Maka (3) salah.
\end{o001proof}
\end{o001solution}

% O001-SOLUTION-ID: O001-FAOA-2015-CH04-EX-010-SOLUTION
% SOURCE-EXERCISE-ID: FAOA-2015-CH04-NODE-0097
% STATEMENT-TARGET-SHA256: 4d967e20b7456eb8ff5a320ebaa0e09a669d541a90efd55b495fa12f0ecb5237
\begin{o001solution}
  {O001-FAOA-2015-CH04-EX-010-SOLUTION}
  {FAOA-2015-CH04-NODE-0097}
  {4d967e20b7456eb8ff5a320ebaa0e09a669d541a90efd55b495fa12f0ecb5237}
\begin{o001statement}
Dalam definisi sebelumnya, rujukan kepada funktor pelupa sering
dihilangkan dan diagram yang menyertainya sering digambar sebagai berikut:
 \[\xy
     \qtriangle/{>}`{>}`{-->}/[S`F`A;\iota`f`\wt f_\iota]
   \endxy\]
Diagram ini tentu terlihat jauh lebih sederhana. Menurut Anda, mengapa saya memilih versi yang
lebih rumit?
\end{o001statement}
\begin{o001answer}
Versi lengkap menjaga tipe setiap objek dan panah: $S,\abs F,\abs A$ serta
$\iota,f,\abs{\widetilde f_\iota}$ berada di kategori himpunan, sedangkan
$F,A$ dan $\widetilde f_\iota$ berada di kategori konkret. Diagram ringkas
mencampurkan dua kategori dan menyembunyikan peran funktor pelupa.
\end{o001answer}
\begin{o001proof}
Dalam kategori konkret $\cat C$, $F$ dan $A$ adalah objek kategori, bukan
secara harfiah himpunan. Pemetaan $\iota\colon S\to\abs F$ dan
$f\colon S\to\abs A$ hanyalah fungsi antarhimpunan dan pada umumnya bukan
morfisme $\cat C$. Sebaliknya, pemetaan universal
$\widetilde f_\iota\colon F\to A$ adalah morfisme dalam $\cat C$, dan
persamaan yang benar pada himpunan dasar ialah
\[
 \abs{\widetilde f_\iota}\circ\iota=f.
\]
Diagram sederhana menulis $S,F,A$ seolah-olah ketiganya objek satu kategori
dan menulis ketiga panah seolah-olah bertipe sama. Diagram yang lebih rumit
memperlihatkan baik segitiga pada himpunan dasar maupun morfisme kategori
yang dipetakan ke sana oleh funktor pelupa. Karena itu komutativitas dan
sifat universalnya bertipe benar dan tidak bergantung pada penyalahgunaan
notasi yang hanya aman setelah pembaca memahami identifikasi tersebut.
\end{o001proof}
\end{o001solution}

% Lapisan solusi O001 -- Bab 5 (Operator pada Ruang Hilbert)
% 4 solusi asli terpisah; bukan tulisan John M. Erdman.
% Provenans model: OpenAI Codex gpt-5.6-sol, Ultra, atas arahan pengguna.
% Lisensi komponen solusi: CC BY-SA 4.0.
% Tidak menyiratkan dukungan John M. Erdman atau Portland State University.

\markboth{SOLUSI O001 UNTUK BAB 5: OPERATOR PADA RUANG HILBERT}{SOLUSI O001 UNTUK BAB 5: OPERATOR PADA RUANG HILBERT}
\section*{Solusi O001 untuk Bab 5: Operator pada Ruang Hilbert}
\addcontentsline{toc}{section}{Solusi O001 untuk Bab 5: Operator pada Ruang Hilbert}
\begin{center}
\small
Komponen solusi asli terpisah, CC BY-SA 4.0. Ditulis dengan bantuan
OpenAI Codex gpt-5.6-sol, Ultra, atas arahan pengguna. Pernyataan latihan
berasal dari edisi terjemahan karya John M. Erdman; solusi ini bukan tulisan
beliau dan tidak menyiratkan dukungan beliau atau Portland State University.
\end{center}

% O001-SOLUTION-ID: O001-FAOA-2015-CH05-EX-001-SOLUTION
% SOURCE-EXERCISE-ID: FAOA-2015-CH05-NODE-0034
% STATEMENT-TARGET-SHA256: ee6308dea448da22e9841970db96f86a106225688b6a5aef98ca4b39f1d20a35
\begin{o001solution}
  {O001-FAOA-2015-CH05-EX-001-SOLUTION}
  {FAOA-2015-CH05-NODE-0034}
  {ee6308dea448da22e9841970db96f86a106225688b6a5aef98ca4b39f1d20a35}
\begin{o001statement}
Misalkan $\N_3 = \{1,2,3\}$ dan $\mu$ ukuran pencacahan pada $\N_3$.
 \begin{enumerate}
  \item[(a)] Identifikasilah ruang Hilbert~$L_2(\N_3,\mu)$.
  \item[(b)] Identifikasilah operator perkalian pada~$L_2(\N_3,\mu)$.
  \item[(c)] Misalkan $T$ operator linear pada $\C^3$ yang representasi
matriksnya adalah
   \[ \begin{bmatrix}
          2  &  0  &  1 \\
          0  &  2  & -1 \\
          1  & -1  &  1
      \end{bmatrix}.\]
Gunakan operator $T$ untuk mengilustrasikan rumusan teorema spektral sebelumnya.
 \end{enumerate}
\end{o001statement}
\begin{o001answer}
(a) $L_2(\N_3,\mu)=\C^3$ dengan
$\langle x,y\rangle=\sum_{j=1}^3x_j\conj{y_j}$. (b) Operator perkalian
adalah $M_\phi=\diag(\phi(1),\phi(2),\phi(3))$. (c) Matriks $T$ mempunyai
nilai eigen $0,2,3$ dengan basis ortonormal
\[
 u_0=\frac{(-1,1,2)}{\sqrt6},\qquad
 u_2=\frac{(1,1,0)}{\sqrt2},\qquad
 u_3=\frac{(1,-1,1)}{\sqrt3}.
\]
Jadi $T$ ekuivalen secara uniter dengan operator perkalian oleh fungsi
$\phi(1)=0$, $\phi(2)=2$, $\phi(3)=3$ pada~$L_2(\N_3,\mu)$.
\end{o001answer}
\begin{o001proof}
Karena ukuran setiap singleton sama dengan satu, kelas fungsi pada $\N_3$
ditentukan oleh tiga nilainya dan
\[
 \norm f_2^2=\sum_{j=1}^3\abs{f(j)}^2.
\]
Ini memberikan identifikasi isometrik dalam (a). Perkalian titik demi titik
oleh $\phi$ kemudian mempunyai matriks diagonal yang dinyatakan dalam (b),
dan setiap operator perkalian diperoleh dengan cara itu.

Untuk matriks $A$ pada (c), ekspansi determinan memberi
\[
 \det(\lambda I-A)=\lambda(\lambda-2)(\lambda-3).
\]
Substitusi langsung menunjukkan
\[
 Au_0=0,\qquad Au_2=2u_2,\qquad Au_3=3u_3.
\]
Ketiga vektor yang ditampilkan bernorma satu dan hasil kali dalam setiap
pasangannya nol, sehingga membentuk basis ortonormal. Definisikan
$U\colon L_2(\N_3,\mu)\to\C^3$ dengan
$U\delta_1=u_0$, $U\delta_2=u_2$, dan $U\delta_3=u_3$. Maka $U$ uniter dan,
untuk $j=1,2,3$,
\[
 TU\delta_j=\phi(j)U\delta_j=UM_\phi\delta_j.
\]
Karena vektor-vektor itu suatu basis, $T=UM_\phi U^{-1}$. Ini tepat
rumusan teorema spektral: operator normal (di sini bahkan swaadjoin)
ekuivalen secara uniter dengan operator perkalian.
\end{o001proof}
\end{o001solution}

% O001-SOLUTION-ID: O001-FAOA-2015-CH05-EX-002-SOLUTION
% SOURCE-EXERCISE-ID: FAOA-2015-CH05-NODE-0038
% STATEMENT-TARGET-SHA256: a4478e29765c8fb04487e45bd6b6fd8fdd22d706e22b5fab1b6b7f7294f1952c
\begin{o001solution}
  {O001-FAOA-2015-CH05-EX-002-SOLUTION}
  {FAOA-2015-CH05-NODE-0038}
  {a4478e29765c8fb04487e45bd6b6fd8fdd22d706e22b5fab1b6b7f7294f1952c}
\begin{o001statement}
Misalkan $V$ adalah \emph{operator Volterra} yang didefinisikan di atas.
 \begin{enumerate}
  \item[(a)] Tunjukkan bahwa $V$ injektif.
  \item[(b)] Hitunglah $V^*$.
  \item[(c)] Tentukan $V + V^*$ beserta jangkauannya.
 \end{enumerate}
\end{o001statement}
\begin{o001answer}
$V$ injektif,
\[
 (V^*g)(t)=\int_t^1g(x)\,dx,
 \qquad
 ((V+V^*)f)(x)=\int_0^1f(t)\,dt,
\]
dan $\ran(V+V^*)$ adalah subruang satu-dimensi semua fungsi konstan pada
$[0,1]$.
\end{o001answer}
\begin{o001proof}
Untuk $f\in L_2([0,1])$, fungsi
$F(x)=\int_0^x f(t)\,dt$ mempunyai wakil kontinu mutlak dan $F'=f$ hampir
di mana-mana. Jika $Vf=0$ sebagai anggota $L_2$, maka $F=0$ hampir di
mana-mana. Kekontinuan $F$ memaksa $F=0$ di setiap titik, sehingga
$f=F'=0$ hampir di mana-mana. Jadi $V$ injektif.

Teorema Fubini dapat diterapkan karena, dengan Cauchy--Schwarz, integral
mutlak pada segitiga $0\leq t\leq x\leq1$ berhingga. Untuk $f,g\in L_2$,
\begin{align*}
 \langle Vf,g\rangle
 &=\int_0^1\left(\int_0^x f(t)\,dt\right)g(x)\,dx\\
 &=\int_0^1 f(t)\left(\int_t^1g(x)\,dx\right)dt.
\end{align*}
Karena ruangnya real, ini sama dengan
$\langle f,V^*g\rangle$ untuk rumus $V^*$ dalam jawaban. Akhirnya,
\[
 ((V+V^*)f)(x)
 =\int_0^xf(t)\,dt+\int_x^1f(t)\,dt
 =\int_0^1f(t)\,dt.
\]
Jadi setiap citra konstan. Sebaliknya, fungsi konstan $f=c$ menghasilkan
fungsi konstan bernilai $c$, sehingga jangkauannya tepat semua fungsi
konstan.
\end{o001proof}
\end{o001solution}

% O001-SOLUTION-ID: O001-FAOA-2015-CH05-EX-003-SOLUTION
% SOURCE-EXERCISE-ID: FAOA-2015-CH05-NODE-0051
% STATEMENT-TARGET-SHA256: a360c061f372f76160bf31105405bc47db3322f73be4c99931a6f4e167c6eb0d
\begin{o001solution}
  {O001-FAOA-2015-CH05-EX-003-SOLUTION}
  {FAOA-2015-CH05-NODE-0051}
  {a360c061f372f76160bf31105405bc47db3322f73be4c99931a6f4e167c6eb0d}
\begin{o001statement}
Misalkan $\fml E$ adalah basis ortonormal bagi ruang Hilbert~$H$. Bahas transformasi linear $T$ pada $H$
sedemikian sehingga $T(e) = 0$ untuk setiap $e \in \fml E$. Berikan contoh tak trivial.
\end{o001statement}
\begin{o001answer}
Jika $T$ terbatas, maka $T=0$. Namun pada ruang Hilbert berdimensi tak
hingga terdapat transformasi linear tak kontinu dan tak nol yang menghilang
pada setiap vektor suatu basis ortonormal.
\end{o001answer}
\begin{o001proof}
Rentang linear $\spn\fml E$ padat dalam~$H$. Jika $T$ terbatas, untuk setiap
$x\in H$ pilih jaring (atau barisan jika basisnya terhitung) $x_\alpha$ dari
$\spn\fml E$ yang konvergen dalam norma ke~$x$. Hipotesis memberi
$Tx_\alpha=0$, sehingga kekontinuan menghasilkan $Tx=0$.

Untuk contoh tak trivial, ambil $H=l_2$, $\fml E=\{e^n:n\in\N\}$, dan
\[
 u=(1,1/2,1/3,\ldots)\in l_2.
\]
Vektor $u$ tidak berada dalam rentang linear hingga $\spn\fml E$, sehingga
$\fml E\cup\{u\}$ bebas linear. Perluas himpunan ini menjadi basis Hamel
$\fml B$ untuk $l_2$. Definisikan $T$ pada basis Hamel dengan
\[
 T(u)=e^1,\qquad T(b)=0\quad(b\in\fml B\setminus\{u\}),
\]
lalu perluas secara linear. Maka $T(e^n)=0$ untuk setiap~$n$, tetapi
$T(u)=e^1\neq0$. Transformasi ini harus tak terbatas, sebab seandainya
terbatas, argumen kepadatan pada paragraf pertama akan memaksanya sama
dengan nol. Dalam dimensi hingga, basis ortonormal sekaligus basis Hamel,
sehingga contoh tak trivial tidak mungkin ada.
\end{o001proof}
\end{o001solution}

% O001-SOLUTION-ID: O001-FAOA-2015-CH05-EX-004-SOLUTION
% SOURCE-EXERCISE-ID: FAOA-2015-CH05-NODE-0075
% STATEMENT-TARGET-SHA256: 84852f969329993492bb16e9519447052b491750458f20a6faa9aa1a87069596
\begin{o001solution}
  {O001-FAOA-2015-CH05-EX-004-SOLUTION}
  {FAOA-2015-CH05-NODE-0075}
  {84852f969329993492bb16e9519447052b491750458f20a6faa9aa1a87069596}
\begin{o001statement}
Misalkan $H$ ruang Hilbert kompleks dan $T$ operator positif pada~$H$.  Definisikan
   \[ \langle x,y \rangle_1 := \langle Tx,y \rangle \qquad \text{ dan } \qquad \norm{x}_1 \equiv \sqrt{\langle x,x \rangle_1} \]
untuk semua $x \in H$.
 \begin{enumerate}
  \item[(a)] Buktikan bahwa $\norm{\hphantom O}_1$ merupakan norma jika dan hanya jika $\ker T = \{0\}$.
  \item[(b)] Misalkan $\ker T = \{0\}$. Tunjukkan bahwa norma $\norm{\hphantom O}$ dan
$\norm{\hphantom O}_1$ menginduksi topologi yang sama pada~$H$ jika dan hanya jika $T$ invertibel dalam
$\ofml B(H)$. \emph{Petunjuk.}  Ketika terdapat dua topologi pada~$H$, sering kali berguna untuk
mempertimbangkan operator identitas pada~$H$.
  \item[(c)] Misalkan $T$ invertibel dalam $\ofml B(H)$.  Untuk $S \in \ofml B(H)$, carilah suatu rumus bagi
adjoin dari $S$ terhadap hasil kali dalam $\langle \hphantom x , \hphantom y \rangle_1$.
 \end{enumerate}
\end{o001statement}
\begin{o001answer}
(a) $\norm{\hphantom O}_1$ adalah norma tepat ketika $\ker T=\{0\}$.
(b) Kedua norma menginduksi topologi yang sama tepat ketika $T$ invertibel.
(c) Jika $S^{*_1}$ menyatakan adjoin terhadap hasil kali dalam baru, maka
\[
 S^{*_1}=T^{-1}S^*T.
\]
\end{o001answer}
\begin{o001proof}
Bentuk $[x,y]=\langle Tx,y\rangle$ bersifat seskuilinear dan Hermitian
karena $T=T^*$, serta $[x,x]\geq0$. Ketaksamaan Cauchy--Schwarz untuk bentuk
positif memberi
\[
 \abs{[x,y]}^2\leq[x,x][y,y].
\]
Jika $[x,x]=0$, ruas kanan menunjukkan $\langle Tx,y\rangle=0$ untuk semua
$y$, jadi $Tx=0$. Sebaliknya, $Tx=0$ langsung memberi $[x,x]=0$. Maka
bentuk itu definit positif, dan akar kuadratnya norma, tepat ketika
$\ker T=\{0\}$. Ini membuktikan (a), termasuk ketaksamaan segitiga yang
merupakan konsekuensi Cauchy--Schwarz bagi hasil kali dalam~$[\ ,\ ]$.

Selalu berlaku
\[
 \norm{x}_1^2=\langle Tx,x\rangle\leq\norm T\,\norm x^2,
\]
jadi identitas $I\colon(H,\norm{\hphantom O})\to
(H,\norm{\hphantom O}_1)$ kontinu. Kedua topologi sama tepat ketika
identitas sebaliknya kontinu, yaitu ketika ada $c>0$ sehingga
\[
 \norm x\leq c\norm{x}_1\quad(x\in H). \tag{1}
\]
Jika (1) berlaku dan $\delta=c^{-2}$, maka
\[
 \delta\norm x^2\leq\langle Tx,x\rangle
 \leq\norm{Tx}\norm x,
\]
sehingga $\norm{Tx}\geq\delta\norm x$. Jadi $T$ injektif dan jangkauannya
tertutup. Selain itu
$(\ran T)^\perp=\ker T^*=\ker T=\{0\}$, sehingga jangkauannya padat. Ia
karena itu sama dengan~$H$, dan invers $T^{-1}$ terbatas dengan norma paling
besar $\delta^{-1}$.

Sebaliknya, andaikan $T$ invertibel. Inversnya positif: jika $y=Tx$, maka
\[
 \langle T^{-1}y,y\rangle=\langle x,Tx\rangle
 =\langle Tx,x\rangle\geq0.
\]
Gunakan Cauchy--Schwarz untuk bentuk positif $[u,v]=\langle Tu,v\rangle$:
\begin{align*}
 \norm x^2
 &=\langle T(T^{-1}x),x\rangle\\
 &\leq
 \langle T(T^{-1}x),T^{-1}x\rangle^{1/2}
 \langle Tx,x\rangle^{1/2}\\
 &=\langle x,T^{-1}x\rangle^{1/2}\norm{x}_1
 \leq\norm{T^{-1}}^{1/2}\norm x\,\norm{x}_1.
\end{align*}
Untuk $x\neq0$, pembagian dengan $\norm x$ memberi
$\norm x\leq\norm{T^{-1}}^{1/2}\norm{x}_1$, yaitu (1). Ini menyelesaikan
(b).

Untuk (c), tetapkan $A=T^{-1}S^*T$, yang terbatas karena ketiga faktornya
terbatas. Untuk semua $x,y\in H$,
\begin{align*}
 \langle Sx,y\rangle_1
 &=\langle TSx,y\rangle
  =\langle Sx,Ty\rangle
  =\langle x,S^*Ty\rangle\\
 &=\langle Tx,T^{-1}S^*Ty\rangle
  =\langle x,Ay\rangle_1.
\end{align*}
Jadi, menurut definisi adjoin dalam hasil kali dalam baru,
$S^{*_1}=A=T^{-1}S^*T$. Urutan faktor ini juga memeriksa bahwa rumus sesuai
dengan konvensi edisi bahwa hasil kali dalam linear pada variabel pertama.
\end{o001proof}
\end{o001solution}

% Lapisan solusi O001 -- Bab 6 (Ruang Banach)
% 6 solusi asli terpisah; bukan tulisan John M. Erdman.
% Provenans model: OpenAI Codex gpt-5.6-sol, Ultra, atas arahan pengguna.
% Lisensi komponen solusi: CC BY-SA 4.0.
% Tidak menyiratkan dukungan John M. Erdman atau Portland State University.

\markboth{SOLUSI O001 UNTUK BAB 6: RUANG BANACH}{SOLUSI O001 UNTUK BAB 6: RUANG BANACH}
\section*{Solusi O001 untuk Bab 6: Ruang Banach}
\addcontentsline{toc}{section}{Solusi O001 untuk Bab 6: Ruang Banach}
\begin{center}
\small
Komponen solusi asli terpisah, CC BY-SA 4.0. Ditulis dengan bantuan
OpenAI Codex gpt-5.6-sol, Ultra, atas arahan pengguna. Pernyataan latihan
berasal dari edisi terjemahan karya John M. Erdman; solusi ini bukan tulisan
beliau dan tidak menyiratkan dukungan beliau atau Portland State University.
\end{center}

% O001-SOLUTION-ID: O001-FAOA-2015-CH06-EX-001-SOLUTION
% SOURCE-EXERCISE-ID: FAOA-2015-CH06-NODE-0013
% STATEMENT-TARGET-SHA256: bd4185ea9f3d1d399ef89f65145ab56c16f94038a48e3a3c894f2fca587e84ca
\begin{o001solution}
  {O001-FAOA-2015-CH06-EX-001-SOLUTION}
  {FAOA-2015-CH06-NODE-0013}
  {bd4185ea9f3d1d399ef89f65145ab56c16f94038a48e3a3c894f2fca587e84ca}
\begin{o001statement}
Misalkan $V$ dan $W$ ruang-ruang linear bernorma, $U$ subhimpunan konveks takkosong dari $V$,
dan $f \colon U \sto W$.  Tanpa menggunakan bentuk apa pun dari \emph{teorema nilai rata-rata}, tunjukkan
bahwa jika $df = \vc 0$ pada $U$, maka $f$ konstan pada~$U$.  (Untuk definisi $df$, yaitu
\emph{diferensial} dari~$f$, lihat Bagian 13.1--2 dalam catatan saya~\cite{Erdman:2007}.)
\end{o001statement}
\begin{o001answer}
Untuk setiap $x,y\in U$ berlaku $f(y)=f(x)$; jadi $f$ konstan pada seluruh~$U$.
\end{o001answer}
\begin{o001proof}
Ambil $x,y\in U$ dan tulis $v=y-x$. Jika $v=0$, klaim segera berlaku.
Jika $v\ne0$, kekonveksan $U$ menjamin bahwa
\[
  x+tv\in U\qquad(0\le t\le1).
\]
Definisikan $g(t)=f(x+tv)$. Dari definisi diferensial dan $df_{x+tv}=0$
diperoleh, langsung tanpa teorema nilai rata-rata,
\[
 \lim_{h\to0}\frac{\norm{g(t+h)-g(t)}}{|h|}=0
\]
untuk $t\in(0,1)$, dengan limit satu sisi pada kedua titik ujung.

Tetapkan $\varepsilon>0$. Misalkan $E$ adalah himpunan semua $s\in[0,1]$
sedemikian sehingga
\[
  \norm{g(t)-g(0)}\le\varepsilon t
  \quad\hbox{untuk setiap }0\le t\le s.
\]
Limit satu sisi di $0$ menunjukkan bahwa $E$ memuat suatu interval
$[0,\delta]$. Misalkan $c=\sup E$. Kontinuitas $g$ menunjukkan bahwa
ketaksamaan tersebut tetap berlaku untuk semua $t\le c$. Jika $c<1$,
limit di $c$ memberi $\eta>0$ sehingga, untuk $0\le t\le1$ dan
$|t-c|<\eta$,
\[
  \norm{g(t)-g(c)}\le\varepsilon|t-c|.
\]
Karena itu, untuk $c\le t<\min(1,c+\eta)$,
\[
 \norm{g(t)-g(0)}
 \le \norm{g(t)-g(c)}+\norm{g(c)-g(0)}
 \le \varepsilon(t-c)+\varepsilon c=\varepsilon t.
\]
Hal ini memperluas interval yang termasuk dalam $E$ melewati $c$, suatu
kontradiksi. Maka $c=1$ dan
$\norm{f(y)-f(x)}=\norm{g(1)-g(0)}\le\varepsilon$. Karena
$\varepsilon$ sebarang, $f(y)=f(x)$. Kedua titik dipilih sebarang, sehingga
$f$ konstan pada~$U$.
\end{o001proof}
\end{o001solution}

% O001-SOLUTION-ID: O001-FAOA-2015-CH06-EX-002-SOLUTION
% SOURCE-EXERCISE-ID: FAOA-2015-CH06-NODE-0038
% STATEMENT-TARGET-SHA256: 47dec53a61225c4a15f0ee33785dcea99ea9c145ba44a8f0f3dd82d79505b2dd
\begin{o001solution}
  {O001-FAOA-2015-CH06-EX-002-SOLUTION}
  {FAOA-2015-CH06-NODE-0038}
  {47dec53a61225c4a15f0ee33785dcea99ea9c145ba44a8f0f3dd82d79505b2dd}
\begin{o001statement}
Apakah terdapat suatu barisan $(a_n)$ dari bilangan kompleks yang memenuhi
syarat berikut: suatu barisan $(x_n)$ dalam $\C$ dapat dijumlahkan secara mutlak jika dan hanya jika
$(a_nx_n)$ terbatas? \emph{Petunjuk.} Tinjau $ T\colon l_\infty \sto l_1\colon (x_n)
\mapsto (x_n/a_n)$.
\end{o001statement}
\begin{o001answer}
Tidak terdapat barisan $(a_n)$ dengan sifat tersebut.
\end{o001answer}
\begin{o001proof}
Andaikan barisan semacam itu ada dan tulis
$Z=\{n\in\N:a_n=0\}$. Himpunan $Z$ harus berhingga: jika $Z$ tak
berhingga, barisan indikator $Z$ tidak termasuk $l_1$, sedangkan
$(a_nx_n)$ identik nol dan karena itu terbatas. Misalkan
$I=\N\setminus Z$; himpunan ini tak berhingga.

Pada ruang Banach $l_1(I)$, definisikan
\[
 A:l_1(I)\longrightarrow l_\infty(I),\qquad
 (Ax)_n=a_nx_n.
\]
Hipotesis memastikan bahwa $A$ terdefinisi pada seluruh $l_1(I)$. Grafnya
tertutup: jika $x^{(k)}\to x$ dalam $l_1(I)$ dan
$Ax^{(k)}\to y$ dalam $l_\infty(I)$, maka konvergensi koordinat memberi
$y_n=a_nx_n$ untuk setiap $n\in I$. Berikut pembuktian lokal bahwa ini
membuat $A$ terbatas, tanpa memakai teorema graf tertutup yang baru muncul
sesudah latihan ini. Graf $G(A)$ adalah ruang Banach karena tertutup dalam
$l_1(I)\times l_\infty(I)$. Proyeksi
$\pi_1:G(A)\to l_1(I)$ merupakan bijeksi linear terbatas. Korolar teorema
pemetaan terbuka memberi bahwa $\pi_1^{-1}$ terbatas. Karena proyeksi kedua
$\pi_2:G(A)\to l_\infty(I)$ juga terbatas,
$A=\pi_2\pi_1^{-1}$ terbatas. Dengan menguji vektor satuan koordinat
diperoleh
\[
  |a_n|\le \norm A=:C\qquad(n\in I).
\]

Sebaliknya, definisikan
\[
 T:l_\infty(I)\longrightarrow l_1(I),\qquad
 (Ty)_n=\frac{y_n}{a_n}.
\]
Untuk $y\in l_\infty(I)$, perluas $Ty$ dengan nol pada $Z$. Hasil kali
$a_n(Ty)_n$ terbatas, sehingga arah ``jika'' dalam hipotesis menjamin
$Ty\in l_1(I)$. Argumen konvergensi koordinat yang sama menunjukkan bahwa
graf $T$ tertutup. Dengan menerapkan lagi argumen ruang-graf dan teorema
pemetaan terbuka dari paragraf sebelumnya, $T$ terbatas. Khusus untuk barisan konstan
$y_n=1$ diperoleh
\[
  \sum_{n\in I}\frac1{|a_n|}<\infty.
\]
Namun $|a_n|\le C$ pada himpunan tak berhingga $I$, sehingga setiap suku
$1/|a_n|\ge1/C$ (dan $C>0$). Deret tersebut harus divergen, bertentangan
dengan kesimpulan sebelumnya. Jadi barisan yang diminta tidak ada.
\end{o001proof}
\end{o001solution}

% O001-SOLUTION-ID: O001-FAOA-2015-CH06-EX-003-SOLUTION
% SOURCE-EXERCISE-ID: FAOA-2015-CH06-NODE-0041
% STATEMENT-TARGET-SHA256: 4f5917f5a92603b55d9972e86c260a0307895ba7080637f5d07e796026fe043b
\begin{o001solution}
  {O001-FAOA-2015-CH06-EX-003-SOLUTION}
  {FAOA-2015-CH06-NODE-0041}
  {4f5917f5a92603b55d9972e86c260a0307895ba7080637f5d07e796026fe043b}
\begin{o001statement}
Tinjau suatu persamaan diferensial berbentuk
   \[ y'' + p_1\,y' + p_0\,y = q  \tag{$\ast$} \]
dengan $p_0$, $p_1$, dan $q$ fungsi-fungsi bernilai real kontinu (yang tetap) pada interval
$[a,b]$.  Nyatakan secara tepat, lalu buktikan, pernyataan bahwa \emph{solusi-solusi $(\ast)$
bergantung secara kontinu pada nilai awal}.  Anda boleh menggunakan tanpa bukti teorema standar berikut:
untuk setiap titik $c$ dalam $[a,b]$ dan setiap pasangan bilangan real $a_0$ dan $a_1$,
terdapat solusi tunggal untuk $(\ast)$ sedemikian sehingga $y(c) = a_0$ dan $y'(c) = a_1$.
\emph{Petunjuk.} Untuk $c \in [a,b]$ yang tetap, tinjau pemetaan
  \[ S\colon \fml C^2([a,b]) \sto \fml C([a,b]) \times \R^2 \colon
                                        f \mapsto \bigl(Tf,f(c),f'(c)\bigr) \]
dengan $\fml C^2([a,b])$ ruang Banach dari Contoh~\ref{C069431} dan $T$ pemetaan
linear terbatas dari Contoh~\ref{C069432}.
\end{o001statement}
\begin{o001answer}
Untuk setiap $c\in[a,b]$ terdapat konstanta $K_c<\infty$ sehingga, jika
$y_{u,v}$ dan $y_{\widetilde u,\widetilde v}$ adalah solusi dengan ruas kanan
$q$ yang sama dan data awal masing-masing $(u,v)$ dan
$(\widetilde u,\widetilde v)$ di $c$, maka
\[
 \norm{y_{u,v}-y_{\widetilde u,\widetilde v}}_{\fml C^2}
 \le K_c\bigl(|u-\widetilde u|+|v-\widetilde v|\bigr).
\]
Jadi ketergantungan pada nilai awal bahkan bersifat Lipschitz dalam norma
$\fml C^2$.
\end{o001answer}
\begin{o001proof}
Lengkapi $\fml C([a,b])\times\R^2$ dengan norma
\[
  \norm{(r,u,v)}=\norm r_u+|u|+|v|.
\]
Ruang ini Banach. Untuk $c$ yang tetap, pemetaan
\[
 S:\fml C^2([a,b])\longrightarrow\fml C([a,b])\times\R^2,
 \qquad S f=(Tf,f(c),f'(c))
\]
linear dan terbatas: keterbatasan $T$ diberikan oleh Contoh~\ref{C069432},
sedangkan $|f(c)|\le\norm f_u$ dan $|f'(c)|\le\norm{f'}_u$.

Teorema keberadaan dan ketunggalan yang diizinkan dalam soal menunjukkan
bahwa $S$ bijektif. Memang, untuk setiap $(r,u,v)$ terdapat tepat satu
$f\in\fml C^2([a,b])$ dengan $Tf=r$, $f(c)=u$, dan $f'(c)=v$.
Karena domain dan kodomainnya Banach, teorema pemetaan terbuka menyatakan
bahwa $S^{-1}$ terbatas. Tuliskan $K_c=\norm{S^{-1}}$.

Sekarang $S y_{u,v}=(q,u,v)$ dan
$S y_{\widetilde u,\widetilde v}=(q,\widetilde u,\widetilde v)$. Linearitas
memberi
\[
 y_{u,v}-y_{\widetilde u,\widetilde v}
 =S^{-1}(0,u-\widetilde u,v-\widetilde v),
\]
dan pengambilan norma menghasilkan tepat estimasi dalam jawaban. Estimasi
itu sekaligus mengendalikan fungsi, turunan pertama, dan turunan kedua
secara seragam; khususnya, konvergensi data awal mengakibatkan konvergensi
solusi dalam $\fml C^2([a,b])$.
\end{o001proof}
\end{o001solution}

% O001-SOLUTION-ID: O001-FAOA-2015-CH06-EX-004-SOLUTION
% SOURCE-EXERCISE-ID: FAOA-2015-CH06-NODE-0044
% STATEMENT-TARGET-SHA256: 402b55b02ea5879d8c8ab502089ac5bfa33972450b373ea5be17ecef04361dfa
\begin{o001solution}
  {O001-FAOA-2015-CH06-EX-004-SOLUTION}
  {FAOA-2015-CH06-NODE-0044}
  {402b55b02ea5879d8c8ab502089ac5bfa33972450b373ea5be17ecef04361dfa}
\begin{o001statement}
Jika $T\colon B \sto C$ suatu pemetaan linear kontinu antara ruang-ruang Banach, maka grafnya tertutup.
\end{o001statement}
\begin{o001answer}
Graf $G(T)=\{(x,Tx):x\in B\}$ tertutup dalam $B\times C$.
\end{o001answer}
\begin{o001proof}
Ambil barisan $(x_n,Tx_n)$ dalam $G(T)$ yang konvergen di $B\times C$ ke
suatu $(x,y)$. Maka $x_n\to x$ dalam $B$ dan $Tx_n\to y$ dalam $C$.
Kontinuitas $T$ juga memberi $Tx_n\to Tx$. Karena limit di ruang bernorma
$C$ tunggal, $y=Tx$. Jadi $(x,y)=(x,Tx)\in G(T)$. Produk ruang bernorma
bersifat metrik, sehingga tertutup secara barisan ekuivalen dengan tertutup;
akibatnya $G(T)$ tertutup.
\end{o001proof}
\end{o001solution}

% O001-SOLUTION-ID: O001-FAOA-2015-CH06-EX-005-SOLUTION
% SOURCE-EXERCISE-ID: FAOA-2015-CH06-NODE-0133
% STATEMENT-TARGET-SHA256: 43e22044a2ac89ed49aa32918c9786ed99670261fafce6eb918ea3414704ab25
\begin{o001solution}
  {O001-FAOA-2015-CH06-EX-005-SOLUTION}
  {FAOA-2015-CH06-NODE-0133}
  {43e22044a2ac89ed49aa32918c9786ed99670261fafce6eb918ea3414704ab25}
\begin{o001statement}
\label{C070114} Misalkan $V$ dan $W$ ruang linear bernorma. Jika suatu keluarga $\fml T$ dari
pemetaan linear terbatas dari $V$ ke $W$ terbatas seragam, maka keluarga itu terbatas titik demi titik.
\end{o001statement}
\begin{o001answer}
Jika $M>0$ memenuhi $\norm T\le M$ untuk semua $T\in\fml T$, maka untuk
setiap $x\in V$ berlaku $\norm{Tx}\le M\norm x$ bagi semua
$T\in\fml T$; misalnya $M_x=M\norm x+1$ adalah batas titik demi titik.
\end{o001answer}
\begin{o001proof}
Ambil $x\in V$. Dari definisi norma operator, untuk setiap
$T\in\fml T$ diperoleh
\[
  \norm{Tx}\le\norm T\,\norm x\le M\norm x<M\norm x+1=M_x.
\]
Konstanta $M_x>0$ hanya bergantung pada $x$, bukan pada $T$. Ini persis
definisi bahwa $\fml T$ terbatas titik demi titik. Tidak diperlukan
kelengkapan $V$ atau~$W$ untuk arah implikasi ini.
\end{o001proof}
\end{o001solution}

% O001-SOLUTION-ID: O001-FAOA-2015-CH06-EX-006-SOLUTION
% SOURCE-EXERCISE-ID: FAOA-2015-CH06-NODE-0144
% STATEMENT-TARGET-SHA256: 6cf97638fdea6d6113149f32f2a310d17d66f855f8b7b863003b30d6e15c694a
\begin{o001solution}
  {O001-FAOA-2015-CH06-EX-006-SOLUTION}
  {FAOA-2015-CH06-NODE-0144}
  {6cf97638fdea6d6113149f32f2a310d17d66f855f8b7b863003b30d6e15c694a}
\begin{o001statement}
Sahabat baik Anda, Fred R. Dimm, kembali memerlukan bantuan. Kali ini ia mencemaskan
pernyataan (lihat~\ref{C070131}) bahwa suatu subhimpunan ruang linear bernorma terbatas jika
dan hanya jika subhimpunan itu terbatas secara lemah. Ia meninjau barisan $(a^n)$ dalam ruang Hilbert
$l^2$ yang didefinisikan oleh $a^n = \sqrt{n}\,\vc e^n$ untuk setiap $n \in \N$, dengan $\{\vc
e^n\colon n \in \N\}$ basis ortonormal biasa untuk~$l_2$ (lihat
contoh~\ref{C063149}). Ia melihat bahwa barisan itu jelas tidak terbatas (dengan topologi biasa
pada~$l^2$). Namun, baginya barisan itu tampak terbatas secara lemah. Jika tidak, menurut argumennya, maka
berdasarkan \emph{teorema Riesz-Fr\'echet} \ref{000342} akan terdapat barisan $x$ dalam
$l^2$ sedemikian sehingga himpunan $\{\abs{\langle a^n,x \rangle}\colon n \in\N\}$ tidak terbatas. Baginya ini tampak
tidak mungkin terjadi karena barisan semacam itu harus menurun lebih lambat daripada
barisan $\bigl(n^{-1/2}\bigr)$, yang sudah menurun terlalu lambat untuk menjadi anggota~$l^2$.
Tenangkan Fred dengan menemukan barisan yang sesuai. \emph{Petunjuk:} Gunakan barisan yang
memiliki banyak nol.
\end{o001statement}
\begin{o001answer}
Ambil
\[
 x_n=\begin{cases}
       k/2^k,&n=4^k\text{ untuk suatu }k\in\N,\\
       0,&\text{selain itu}.
     \end{cases}
\]
Maka $x\in l^2$, tetapi
$\abs{\langle a^{4^k},x\rangle}=k$, sehingga himpunan nilai hasil kali
dalam tersebut tidak terbatas.
\end{o001answer}
\begin{o001proof}
Karena koordinat taknol hanya muncul pada $n=4^k$, diperoleh
\[
 \norm x_2^2=\sum_{k=1}^{\infty}\frac{k^2}{4^k}<\infty.
\]
Jadi $x\in l^2$. Untuk $n=4^k$, vektor
$a^n=\sqrt n\,\vc e^n=2^k\vc e^{4^k}$. Terlepas dari konvensi mengenai
peubah mana yang linear dalam hasil kali dalam kompleks, nilai mutlaknya
adalah
\[
 \abs{\langle a^{4^k},x\rangle}
 =2^k\abs{x_{4^k}}=2^k\frac{k}{2^k}=k.
\]
Nilai ini menuju tak hingga. Dengan demikian barisan jarang tersebut
memberikan fungsional Riesz--Fr\'echet yang menyaksikan bahwa
$\{a^n:n\in\N\}$ tidak terbatas secara lemah, sesuai dengan fakta bahwa
$\norm{a^n}=\sqrt n$ tidak terbatas.
\end{o001proof}
\end{o001solution}

% Lapisan solusi O001 -- Bab 7 (Operator Kompak)
% 1 solusi asli terpisah; bukan tulisan John M. Erdman.
% Provenans model: OpenAI Codex gpt-5.6-sol, Ultra, atas arahan pengguna.
% Lisensi komponen solusi: CC BY-SA 4.0.
% Tidak menyiratkan dukungan John M. Erdman atau Portland State University.

\markboth{SOLUSI O001 UNTUK BAB 7: OPERATOR KOMPAK}{SOLUSI O001 UNTUK BAB 7: OPERATOR KOMPAK}
\section*{Solusi O001 untuk Bab 7: Operator Kompak}
\addcontentsline{toc}{section}{Solusi O001 untuk Bab 7: Operator Kompak}
\begin{center}
\small
Komponen solusi asli terpisah, CC BY-SA 4.0. Ditulis dengan bantuan
OpenAI Codex gpt-5.6-sol, Ultra, atas arahan pengguna. Pernyataan latihan
berasal dari edisi terjemahan karya John M. Erdman; solusi ini bukan tulisan
beliau dan tidak menyiratkan dukungan beliau atau Portland State University.
\end{center}

% O001-SOLUTION-ID: O001-FAOA-2015-CH07-EX-001-SOLUTION
% SOURCE-EXERCISE-ID: FAOA-2015-CH07-NODE-0045
% STATEMENT-TARGET-SHA256: d34a802dca82cf381a646022b8b8cec4c40d52cfbb09ae67d5a13d39e1d48396
\begin{o001solution}
  {O001-FAOA-2015-CH07-EX-001-SOLUTION}
  {FAOA-2015-CH07-NODE-0045}
  {d34a802dca82cf381a646022b8b8cec4c40d52cfbb09ae67d5a13d39e1d48396}
\begin{o001statement}
Misalkan $(S, \fml A,\mu)$ suatu ruang ukur $\sigma$-hingga dan $\phi \in
L_\infty(\mu)$. Tentukan dekomposisi polar dari
 \index{operator perkalian!dekomposisi polar dari}%
 \index{operator!perkalian!dekomposisi polar dari}%
 \index{dekomposisi polar!dari operator perkalian}%
operator perkalian $M_\phi$ (lihat contoh~\ref{C063527}).
\end{o001statement}
\begin{o001answer}
Pada $L_2(S,\mu)$, misalkan $E=\{s:|\phi(s)|>0\}$ hingga himpunan nol dan
definisikan
\[
 u(s)=\begin{cases}\phi(s)/|\phi(s)|,&s\in E,\\0,&s\notin E.\end{cases}
\]
Maka
\[
  |M_\phi|=M_{|\phi|},\qquad M_\phi=M_uM_{|\phi|},
\]
dan $M_u$ adalah isometri parsial dengan proyeksi awal sekaligus proyeksi
akhir $M_{\mathbf1_E}$.
\end{o001answer}
\begin{o001proof}
Dari Contoh~\ref{C063527}, ${M_\phi}^*=M_{\overline\phi}$. Oleh karena itu
\[
 {M_\phi}^*M_\phi=M_{|\phi|^2}
 \quad\hbox{dan}\quad
 |M_\phi|=({M_\phi}^*M_\phi)^{1/2}=M_{|\phi|},
\]
karena $M_{|\phi|}$ positif dan kuadratnya $M_{|\phi|^2}$.

Fungsi $u$ terukur, terbatas secara esensial, dan
$|u|^2=\mathbf1_E$ hampir di mana-mana. Maka
\[
  M_u^*M_u=M_{|u|^2}=M_{\mathbf1_E},
\]
sebuah proyeksi ortogonal; jadi $M_u$ isometri parsial. Perhitungan titik
demi titik memberi
\[
 M_uM_{|\phi|}=M_{u|\phi|}=M_\phi.
\]
Kernel $M_\phi$ terdiri atas kelas fungsi yang nol hampir di mana-mana pada
$E$, sehingga komplemen ortogonalnya adalah $L_2(E)$, yakni jangkauan
$M_{\mathbf1_E}$. Ini ruang awal $M_u$.

Untuk memeriksa ruang akhir yang disyaratkan dekomposisi polar, ambil
$g\in L_2(E)$ dan tetapkan
$g_n=g\mathbf1_{\{|\phi|\ge1/n\}}$. Maka $g_n\to g$ dalam $L_2$ melalui
konvergensi terdominasi, sedangkan
\[
 g_n=M_\phi\left(\frac{g_n}{\phi}\right)
 \quad\hbox{dan}\quad
 \frac{g_n}{\phi}\in L_2.
\]
Jadi $\overline{\ran M_\phi}=L_2(E)$. Karena $M_uM_u^*=M_{\mathbf1_E}$,
ruang akhir $M_u$ juga $L_2(E)$. Semua syarat dekomposisi polar, termasuk
kasus $\phi=0$ hampir di mana-mana, telah terpenuhi.
\end{o001proof}
\end{o001solution}

% Lapisan solusi O001 -- Bab 8 (Spektrum)
% 2 solusi asli terpisah; bukan tulisan John M. Erdman.
% Provenans model: OpenAI Codex gpt-5.6-sol, Ultra, atas arahan pengguna.
% Lisensi komponen solusi: CC BY-SA 4.0.
% Tidak menyiratkan dukungan John M. Erdman atau Portland State University.

\markboth{SOLUSI O001 UNTUK BAB 8: SPEKTRUM}{SOLUSI O001 UNTUK BAB 8: SPEKTRUM}
\section*{Solusi O001 untuk Bab 8: Spektrum}
\addcontentsline{toc}{section}{Solusi O001 untuk Bab 8: Spektrum}
\begin{center}
\small
Komponen solusi asli terpisah, CC BY-SA 4.0. Ditulis dengan bantuan
OpenAI Codex gpt-5.6-sol, Ultra, atas arahan pengguna. Pernyataan latihan
berasal dari edisi terjemahan karya John M. Erdman; solusi ini bukan tulisan
beliau dan tidak menyiratkan dukungan beliau atau Portland State University.
\end{center}

% O001-SOLUTION-ID: O001-FAOA-2015-CH08-EX-001-SOLUTION
% SOURCE-EXERCISE-ID: FAOA-2015-CH08-NODE-0031
% STATEMENT-TARGET-SHA256: e16770f5d0da46b3543a0b001816b53299c752a00a3a6db2d47672ff3d0309cf
\begin{o001solution}
  {O001-FAOA-2015-CH08-EX-001-SOLUTION}
  {FAOA-2015-CH08-NODE-0031}
  {e16770f5d0da46b3543a0b001816b53299c752a00a3a6db2d47672ff3d0309cf}
\begin{o001statement}
Identifikasi hasil kali dan koproduk (jika ada) dalam kategori $\cat{BALG}$
yang objeknya aljabar Banach
 \index{hasil kali!dalam $\cat{BALG}$}%
 \index{Banach!aljabar!hasil kali}%
 \index{Banach!aljabar!koproduk}%
 \index{koproduk!dalam $\cat{BALG}$}%
 \index{BALG@$\cat{BALG}$!hasil kali dan koproduk dalam}%
dan morfismenya homomorfisme aljabar kontinu, serta dalam kategori yang objeknya aljabar Banach dan
morfismenya homomorfisme aljabar kontraktif. (Di sini,
 \index{kontraktif}%
\emph{kontraktif} berarti $\norm{Tx} \le \norm x$ untuk setiap~$x$.)
\end{o001statement}
\begin{o001answer}
Dalam $\cat{BALG}$ dengan homomorfisme kontinu, hasil kali $A$ dan $B$
adalah $A\times B$ dengan operasi titik demi titik dan, misalnya, norma
$\norm{(a,b)}=\norm a+\norm b$; kategori ini tidak mempunyai koproduk
biner pada umumnya. Dalam kategori homomorfisme kontraktif, hasil kali
adalah $A\times B$ dengan norma maksimum, sedangkan koproduknya adalah
produk bebas Banach kontraktif $A*B$, yaitu penyelesaian jumlah-$l_1$
dari kata-kata tensor berselang-seling dalam $A$ dan~$B$.
\end{o001answer}
\begin{o001proof}
Untuk homomorfisme kontinu, proyeksi koordinat dari $A\times B$ kontinu.
Jika $f:C\to A$ dan $g:C\to B$ kontinu, satu-satunya pemetaan yang membuat
diagram hasil kali komutatif ialah
\[
  h:C\longrightarrow A\times B,\qquad h(c)=(f(c),g(c)).
\]
Pemetaan ini homomorfisme dan
$\norm{h(c)}\le(\norm f+\norm g)\norm c$, sehingga kontinu. Ini membuktikan
sifat universal hasil kali. Norma maksimum juga memberi objek hasil kali
yang isomorfik dalam kategori ini.

Untuk melihat bahwa koproduk tidak selalu ada dalam kategori morfisme
kontinu, ambil aljabar Banach satu dimensi $E=\C$ dengan perkalian nol.
Andaikan $(P,i,j)$ adalah koproduk dua salinan $E$, dan tulis
$p=i(1)$ serta $q=j(1)$. Untuk setiap $\lambda>0$, pemetaan
\[
 f_\lambda(z)=\lambda z e_{12},\qquad
 g_\lambda(z)=\lambda z e_{21}
\]
adalah homomorfisme kontinu $E\to M_2(\C)$ karena
$e_{12}^2=e_{21}^2=0$. Sifat universal akan memberi homomorfisme kontinu
$H_\lambda:P\to M_2(\C)$ dengan nilai-nilai tersebut. Maka
\[
 H_\lambda((pq)^n)=\lambda^{2n}e_{11}.
\]
Dengan norma operator pada $M_2(\C)$ dan submultiplikativitas norma di
$P$, kontinuitas $H_\lambda$ akan memaksa
\[
 \lambda^{2n}\le\norm{H_\lambda}
       (\norm p\,\norm q)^n\qquad(n\ge1).
\]
Pemetaan universal ke aljabar perkalian-nol menunjukkan $p,q\ne0$.
Memilih $\lambda^2>\norm p\,\norm q$ membuat ketaksamaan terakhir mustahil
saat $n\to\infty$. Jadi kategori ini tidak mempunyai semua koproduk biner.

Dalam kategori homomorfisme kontraktif, beri $A\times B$ norma
\[
  \norm{(a,b)}_{\max}=\max\{\norm a,\norm b\}.
\]
Proyeksinya kontraktif, dan pasangan $h=(f,g)$ kontraktif apabila $f$ dan
$g$ kontraktif. Jadi inilah hasil kali pada kategori tersebut.

Untuk koproduk kontraktif, bagi setiap kata takkosong berselang-seling
$\epsilon=(\epsilon_1,\ldots,\epsilon_n)$ dengan huruf $A,B$, bentuk hasil
kali tensor projektif
\[
 X_{\epsilon_1}\mathbin{\widehat{\otimes}_{\pi}}\cdots
 \mathbin{\widehat{\otimes}_{\pi}}X_{\epsilon_n},
 \qquad X_A=A,\quad X_B=B,
\]
lalu ambil jumlah langsung-$l_1$ atas semua kata tersebut. Perkalian
dilakukan dengan menggabungkan kata; jika dua huruf pada sambungan sama,
kedua faktor batas dikalikan di dalam $A$ atau $B$. Submultiplikativitas
norma aljabar dan sifat norma projektif membuat operasi ini kontraktif dan
memanjang ke penyelesaiannya, yang kita notasikan $A*B$.

Penyertaan ke suku satu-huruf bersifat isometrik. Jika
$f:A\to C$ dan $g:B\to C$ kontraktif, pada tensor elementer definisikan
\[
 x_1\otimes\cdots\otimes x_n\longmapsto
 F_{\epsilon_1}(x_1)\cdots F_{\epsilon_n}(x_n),
 \qquad F_A=f,\quad F_B=g.
\]
Norma hasilnya paling besar daripada hasil kali norma faktor, sehingga
pemetaan pada setiap suku dan pada jumlah-$l_1$ kontraktif. Pemetaan ini
adalah homomorfisme dan memanjang secara unik ke $A*B$; keunikannya berlaku
karena kata-kata dalam kedua citra penyertaan rapat. Jadi $A*B$ memenuhi
sifat universal koproduk kontraktif.
\end{o001proof}
\end{o001solution}

% O001-SOLUTION-ID: O001-FAOA-2015-CH08-EX-002-SOLUTION
% SOURCE-EXERCISE-ID: FAOA-2015-CH08-NODE-0064
% STATEMENT-TARGET-SHA256: 8542ce320f7fb680060f8d11f114b381d061752b828a1f4a5e5f5080a8287e9c
\begin{o001solution}
  {O001-FAOA-2015-CH08-EX-002-SOLUTION}
  {FAOA-2015-CH08-NODE-0064}
  {8542ce320f7fb680060f8d11f114b381d061752b828a1f4a5e5f5080a8287e9c}
\begin{o001statement}
Dengan menggunakan rumus untuk operator Volterra $V$ dalam contoh~\ref{000319} sebagai
 \index{operator Volterra!dan persamaan integral}%
 \index{operator!Volterra!dan persamaan integral}%
acuan, pada ruang Banach~$\fml C([0,1])$ definisikan operator integral melalui rumus $Vf(x)=\int_0^x f(t)\,dt$.
 \begin{enumerate}
  \item[(a)] Hitung spektrum operator $V$. (\emph{Petunjuk.} Tunjukkan bahwa $\norm{V^n}
\le (n!)^{-1}$ untuk setiap $n \in \N$ dan gunakan \emph{rumus radius spektral}.)
  \item[(b)] Apa implikasi (a) terhadap kemungkinan memperoleh fungsi kontinu $f$ yang
memenuhi persamaan integral
   \[ f(x) - \mu\int_0^xf(t)\,dt = h(x), \]
dengan $\mu$ suatu skalar dan $h$ suatu fungsi kontinu yang diberikan pada~$[0,1]$?
  \item[(c)] Gunakan gagasan dalam (b) dan uraian deret Neumann untuk $\vc 1 - \mu V$ (lihat
proposisi~\ref{prop_Neumann_series}) untuk menghitung secara eksplisit (dan dalam bentuk fungsi elementer)
suatu solusi untuk persamaan integral
 \[f(x) - {\textstyle \frac12}\int_0^xf(t)\,dt = e^x.\]
 \end{enumerate}
\end{o001statement}
\begin{o001answer}
\begin{enumerate}
 \item[(a)] $\sigma(V)=\{0\}$.
 \item[(b)] Untuk setiap skalar $\mu$ dan setiap
 $h\in\fml C([0,1])$, terdapat tepat satu solusi, yaitu
 \[
   f=(\vc1-\mu V)^{-1}h=\sum_{n=0}^{\infty}\mu^nV^nh.
 \]
 \item[(c)] Solusi eksplisitnya adalah
 $f(x)=2e^x-e^{x/2}$.
\end{enumerate}
\end{o001answer}
\begin{o001proof}
Induksi dan teorema Fubini untuk integral kontinu memberi, untuk $n\ge1$,
\[
 (V^nf)(x)=\frac1{(n-1)!}\int_0^x(x-t)^{n-1}f(t)\,dt.
\]
Karena $0\le x\le1$,
\[
 |(V^nf)(x)|\le\frac{x^n}{n!}\norm f_u
 \le\frac1{n!}\norm f_u,
\]
sehingga $\norm{V^n}\le1/n!$. Rumus radius spektral menghasilkan
\[
 \rho(V)\le\lim_{n\to\infty}(1/n!)^{1/n}=0.
\]
Spektrum operator pada ruang Banach kompleks takkosong, jadi
$\sigma(V)\subseteq\{0\}$ memaksa $\sigma(V)=\{0\}$. Secara langsung,
$0$ memang berada dalam spektrum karena $V$ tidak surjektif: setiap fungsi
dalam jangkauannya bernilai nol di~$0$.

Untuk setiap $\mu$, $\sigma(\mu V)=\{0\}$ (termasuk $\mu=0$), sehingga
$\vc1-\mu V$ invertibel. Lebih konkret, estimasi di atas memberi
\[
 \sum_{n=0}^{\infty}\norm{\mu^nV^n}
 \le\sum_{n=0}^{\infty}\frac{|\mu|^n}{n!}=e^{|\mu|}<\infty.
\]
Karena itu deret operator $R=\sum_{n\ge0}\mu^nV^n$ konvergen dalam norma.
Mengalikan jumlah parsial dengan $\vc1-\mu V$ menghasilkan
$\vc1-\mu^{N+1}V^{N+1}$, yang menuju $\vc1$ dalam norma; urutan perkalian
yang lain sama. Maka $R=(\vc1-\mu V)^{-1}$ dan untuk setiap $h$ terdapat
tepat satu solusi $f=Rh$.

Untuk $h(x)=e^x$ dan $n\ge1$, induksi dari integrasi menunjukkan
\[
 (V^ne^\cdot)(x)=e^x-\sum_{k=0}^{n-1}\frac{x^k}{k!}.
\]
Dengan $\mu=1/2$, penataan ulang sah karena konvergensi mutlak dan seragam,
dan
\begin{align*}
 f(x)
 &=e^x+\sum_{n=1}^{\infty}2^{-n}
       \left(e^x-\sum_{k=0}^{n-1}\frac{x^k}{k!}\right)\\
 &=2e^x-\sum_{k=0}^{\infty}\frac{x^k}{k!}
       \sum_{n=k+1}^{\infty}2^{-n}
 =2e^x-\sum_{k=0}^{\infty}\frac{(x/2)^k}{k!}
 =2e^x-e^{x/2}.
\end{align*}
Sebagai pemeriksaan langsung,
\[
 \frac12\int_0^x(2e^t-e^{t/2})\,dt=e^x-e^{x/2},
\]
sehingga $f(x)-\frac12\int_0^xf(t)\,dt=e^x$.
\end{o001proof}
\end{o001solution}

% Lapisan solusi O001 -- Bab 9 (Ruang Vektor Topologis)
% 1 solusi asli terpisah; bukan tulisan John M. Erdman.
% Provenans model: OpenAI Codex gpt-5.6-sol, Ultra, atas arahan pengguna.
% Lisensi komponen solusi: CC BY-SA 4.0.
% Tidak menyiratkan dukungan John M. Erdman atau Portland State University.

\markboth{SOLUSI O001 UNTUK BAB 9: RUANG VEKTOR TOPOLOGIS}{SOLUSI O001 UNTUK BAB 9: RUANG VEKTOR TOPOLOGIS}
\section*{Solusi O001 untuk Bab 9: Ruang Vektor Topologis}
\addcontentsline{toc}{section}{Solusi O001 untuk Bab 9: Ruang Vektor Topologis}
\begin{center}
\small
Komponen solusi asli terpisah, CC BY-SA 4.0. Ditulis dengan bantuan
OpenAI Codex gpt-5.6-sol, Ultra, atas arahan pengguna. Pernyataan latihan
berasal dari edisi terjemahan karya John M. Erdman; solusi ini bukan tulisan
beliau dan tidak menyiratkan dukungan beliau atau Portland State University.
\end{center}

% O001-SOLUTION-ID: O001-FAOA-2015-CH09-EX-001-SOLUTION
% SOURCE-EXERCISE-ID: FAOA-2015-CH09-NODE-0072
% STATEMENT-TARGET-SHA256: 69f9a89b1b499eb3fee7d7d92605c324b0dabd8f66918754428f27cd5482fdfe
\begin{o001solution}
  {O001-FAOA-2015-CH09-EX-001-SOLUTION}
  {FAOA-2015-CH09-NODE-0072}
  {69f9a89b1b499eb3fee7d7d92605c324b0dabd8f66918754428f27cd5482fdfe}
\begin{o001statement}
Jelaskan pula bagaimana proposisi sebelumnya memungkinkan kita mengaitkan suatu ruang lain yang bersifat Hausdorff dengan
suatu ruang vektor topologis non-Hausdorff.
\end{o001statement}
\begin{o001answer}
Untuk ruang vektor topologis $X$, ruang yang terkait secara kanonik adalah
Hausdorffisasi
\[
  X_{\mathrm H}:=X/\overline{\{\vc0\}}.
\]
Ruang ini Hausdorff, dan setiap pemetaan linear kontinu dari $X$ ke ruang
vektor topologis Hausdorff berfaktor secara unik melalui pemetaan hasil bagi
$X\to X_{\mathrm H}$.
\end{o001answer}
\begin{o001proof}
Misalkan $N=\overline{\{\vc0\}}$. Karena penjumlahan, invers aditif, dan
perkalian skalar kontinu, tutupan subruang vektor tetap merupakan subruang
vektor; jadi $N$ adalah subruang tertutup dari~$X$. Proposisi sebelumnya
langsung memberi bahwa ruang hasil bagi $X/N$ dengan topologi hasil bagi
bersifat Hausdorff. Hal ini tetap berlaku ketika $X$ sendiri non-Hausdorff;
dalam kasus tersebut $N\ne\{\vc0\}$.

Konstruksi ini kanonik. Jika $Y$ ruang vektor topologis Hausdorff dan
$T:X\to Y$ linear kontinu, maka $\{\vc0\}$ tertutup dalam $Y$, sehingga
$\ker T=T^{-1}(\{\vc0\})$ tertutup dalam $X$. Karena $\vc0\in\ker T$,
\[
  N=\overline{\{\vc0\}}\subseteq\ker T.
\]
Maka rumus $\widetilde T([x])=T(x)$ terdefinisi dengan baik dan memberi
pemetaan linear unik $\widetilde T:X/N\to Y$ dengan
$T=\widetilde T\circ\pi$. Kontinuitas $\widetilde T$ mengikuti sifat
topologi hasil bagi: untuk setiap himpunan terbuka $O\subseteq Y$,
$\pi^{-1}(\widetilde T^{-1}(O))=T^{-1}(O)$ terbuka. Jadi $X/N$ bukan hanya
suatu ruang Hausdorff yang terkait, melainkan hasil bagi Hausdorff universal
dari~$X$.
\end{o001proof}
\end{o001solution}

% Lapisan solusi O001 -- Bab 10 (Distribusi)
% 11 solusi asli terpisah; bukan tulisan John M. Erdman.
% Provenans model: OpenAI Codex gpt-5.6-sol, Ultra, atas arahan pengguna.
% Lisensi komponen solusi: CC BY-SA 4.0.
% Tidak menyiratkan dukungan John M. Erdman atau Portland State University.

\markboth{SOLUSI O001 UNTUK BAB 10: DISTRIBUSI}{SOLUSI O001 UNTUK BAB 10: DISTRIBUSI}
\section*{Solusi O001 untuk Bab 10: Distribusi}
\addcontentsline{toc}{section}{Solusi O001 untuk Bab 10: Distribusi}
\begin{center}
\small
Komponen solusi asli terpisah, CC BY-SA 4.0. Ditulis dengan bantuan
OpenAI Codex gpt-5.6-sol, Ultra, atas arahan pengguna. Pernyataan latihan
berasal dari edisi terjemahan karya John M. Erdman; solusi ini bukan tulisan
beliau dan tidak menyiratkan dukungan beliau atau Portland State University.
\end{center}

% O001-SOLUTION-ID: O001-FAOA-2015-CH10-EX-001-SOLUTION
% SOURCE-EXERCISE-ID: FAOA-2015-CH10-NODE-0034
% STATEMENT-TARGET-SHA256: 993392326124fdc4d83c44ae53cdb333a2ca2c9911da027d164621951aafc62a
\begin{o001solution}
  {O001-FAOA-2015-CH10-EX-001-SOLUTION}
  {FAOA-2015-CH10-NODE-0034}
  {993392326124fdc4d83c44ae53cdb333a2ca2c9911da027d164621951aafc62a}
\begin{o001statement}
Misalkan $S$ suatu fungsi signum (atau tanda) pada $\R$; artinya, misalkan $S = 2H - 1$, dengan $H$ fungsi Heaviside.
Misalkan pula $g$ suatu fungsi pada $\R$ sedemikian sehingga $g\colon x \mapsto -x - 3$ untuk $x < 0$ dan $g\colon x \mapsto x + 5$ untuk $x \ge 0$.
   \begin{enumerate}[label=\rm{(\alph*)}]
     \item Carilah konstanta $\alpha$ dan $\beta$ sedemikian sehingga ${\tilde g}' = \alpha \widetilde S + \beta \delta$.
     \item Carilah suatu fungsi $a$ pada $\R$ sedemikian sehingga ${\tilde a}' = \wt S$.
   \end{enumerate}
\end{o001statement}
\begin{o001answer}
$\alpha=1$, $\beta=8$, dan salah satu pilihan ialah $a(x)=\lvert x\rvert$.
Lebih umum, $a(x)=\lvert x\rvert+C$ untuk sembarang konstanta~$C$.
\end{o001answer}
\begin{o001proof}
Pada masing-masing setengah garis, turunan klasik $g$ sama dengan $-1$ untuk
$x<0$ dan $1$ untuk $x>0$, yaitu sama hampir di mana-mana dengan~$S$.
Sementara itu, loncatan $g$ di titik asal adalah
\[
  g(0+)-g(0-)=5-(-3)=8.
\]
Rumus turunan distribusional bagi fungsi mulus-sepotong karena itu memberi
\[
  {\tilde g}'=\widetilde S+8\delta,
\]
sehingga $\alpha=1$ dan $\beta=8$.

Untuk bagian kedua, ambil $a(x)=\lvert x\rvert$. Fungsi ini kontinu di nol,
jadi tidak menghasilkan suku delta, dan turunannya sama dengan $-1$ di
$(-\infty,0)$ serta $1$ di $(0,\infty)$. Maka
${\tilde a}'=\widetilde S$. Penambahan konstanta tidak mengubah turunan
distribusional.
\end{o001proof}
\end{o001solution}

% O001-SOLUTION-ID: O001-FAOA-2015-CH10-EX-002-SOLUTION
% SOURCE-EXERCISE-ID: FAOA-2015-CH10-NODE-0035
% STATEMENT-TARGET-SHA256: f48002bdb810711a169089b4f615782e5af6798ecf6f38a7be8141ef70a103e5
\begin{o001solution}
  {O001-FAOA-2015-CH10-EX-002-SOLUTION}
  {FAOA-2015-CH10-NODE-0035}
  {f48002bdb810711a169089b4f615782e5af6798ecf6f38a7be8141ef70a103e5}
\begin{o001statement}
Misalkan $f(x) = \left\{
                           \begin{array}{ll}
                             1 - \abs x, & \hbox{jika $\abs x \le 1$;} \\
                                 0,      & \hbox{selain itu.}
                           \end{array}
                         \right.$ Dengan menghitung kedua ruas persamaan~\eqref{deriv_reg_distr} secara terpisah, tunjukkan bahwa persamaan itu
benar secara klasik untuk setiap fungsi uji yang tumpuannya termuat dalam interval~$[-1,1]$.
\end{o001statement}
\begin{o001answer}
Kedua ruas bernilai
\[
  \int_{-1}^{0}\phi(x)\,dx-\int_{0}^{1}\phi(x)\,dx .
\]
\end{o001answer}
\begin{o001proof}
Turunan klasik $f$ ada hampir di mana-mana dan diberikan oleh
\[
 f'(x)=
 \begin{cases}
  1,&-1<x<0,\\
 -1,&0<x<1,\\
  0,&\text{selain itu}.
 \end{cases}
\]
Oleh sebab itu,
\[
 L_{f'}(\phi)=\int_{-1}^{0}\phi(x)\,dx-
               \int_{0}^{1}\phi(x)\,dx .
\]
Di sisi lain, dengan memecah integral di nol dan mengintegralkan per bagian,
\begin{align*}
 (L_f)'(\phi)
 &=-\int_{-1}^{0}(1+x)\phi'(x)\,dx
   -\int_{0}^{1}(1-x)\phi'(x)\,dx\\
 &=\left[-\phi(0)+\int_{-1}^{0}\phi(x)\,dx\right]
   +\left[\phi(0)-\int_{0}^{1}\phi(x)\,dx\right]\\
 &=\int_{-1}^{0}\phi(x)\,dx-\int_{0}^{1}\phi(x)\,dx.
\end{align*}
Suku batas di nol saling menghapus. Dengan demikian
$(L_f)'=L_{f'}$, sebagaimana dinyatakan oleh~\eqref{deriv_reg_distr} untuk
fungsi kontinu dan mulus-sepotong ini.
\end{o001proof}
\end{o001solution}

% O001-SOLUTION-ID: O001-FAOA-2015-CH10-EX-003-SOLUTION
% SOURCE-EXERCISE-ID: FAOA-2015-CH10-NODE-0036
% STATEMENT-TARGET-SHA256: b4a47ce3926ad911b776a8e31eb2efc713e5e925d51c317a06458160b0c7cbbc
\begin{o001solution}
  {O001-FAOA-2015-CH10-EX-003-SOLUTION}
  {FAOA-2015-CH10-NODE-0036}
  {b4a47ce3926ad911b776a8e31eb2efc713e5e925d51c317a06458160b0c7cbbc}
\begin{o001statement}
Misalkan $f(x) = \sbsb\chi{(0,1)}$ untuk $x \in~\R$ dan $\phi \in \fml D\bigl((-1,1)\bigr)$. Bandingkan nilai $L_{f\,'}(\phi)$ dan
$\bigl(L_f\bigr)'(\phi)$.
\end{o001statement}
\begin{o001answer}
$L_{f'}(\phi)=0$, sedangkan $(L_f)'(\phi)=\phi(0)$; jadi pada
$(-1,1)$ berlaku $(L_f)'=\delta_0$ dan keduanya tidak sama pada umumnya.
\end{o001answer}
\begin{o001proof}
Turunan klasik $f'$ sama dengan nol hampir di mana-mana, sehingga distribusi
reguler yang ditentukannya memenuhi $L_{f'}(\phi)=0$. Sebaliknya,
\[
 (L_f)'(\phi)=-L_f(\phi')=-\int_0^1\phi'(x)\,dx
 =\phi(0)-\phi(1)=\phi(0).
\]
Karena $\phi$ bertumpuan kompak di dalam $(-1,1)$, nilainya di~$1$ adalah
nol. Jadi turunan distribusional menangkap loncatan di~$0$, yang diabaikan
oleh turunan klasik hampir di mana-mana.
\end{o001proof}
\end{o001solution}

% O001-SOLUTION-ID: O001-FAOA-2015-CH10-EX-004-SOLUTION
% SOURCE-EXERCISE-ID: FAOA-2015-CH10-NODE-0037
% STATEMENT-TARGET-SHA256: 76fa760413de68845660f5155e23b736072da6a018d66a8e53d649fa3f39a8eb
\begin{o001solution}
  {O001-FAOA-2015-CH10-EX-004-SOLUTION}
  {FAOA-2015-CH10-NODE-0037}
  {76fa760413de68845660f5155e23b736072da6a018d66a8e53d649fa3f39a8eb}
\begin{o001statement}
Misalkan $f(x) = \left\{
                           \begin{array}{ll}
                             x, & \hbox{jika $0 \le x \le 1$;} \\
                             x + 2, & \hbox{jika $1 < x \le 2$;} \\
                             0, & \hbox{selain itu.}
                           \end{array}
                         \right.$ Bandingkan nilai $L_{f\,'}(\phi)$ dan $\bigl(L_f\bigr)'(\phi)$ untuk suatu fungsi uji~$\phi$.
\end{o001statement}
\begin{o001answer}
\[
 L_{f'}(\phi)=\int_0^2\phi(x)\,dx,
 \qquad
 (L_f)'(\phi)=\int_0^2\phi(x)\,dx+2\phi(1)-4\phi(2).
\]
\end{o001answer}
\begin{o001proof}
Di luar himpunan berukuran nol $\{0,1,2\}$, turunan klasik $f$ sama dengan
$1$ pada $(0,2)$ dan nol di luar interval tersebut. Maka
$L_{f'}(\phi)=\int_0^2\phi$. Fungsi $f$ tidak meloncat di~$0$, tetapi
loncatannya di~$1$ dan~$2$ berturut-turut adalah
\[
 f(1+)-f(1-)=3-1=2,
 \qquad
 f(2+)-f(2-)=0-4=-4.
\]
Turunan distribusional fungsi mulus-sepotong adalah distribusi reguler dari
turunan hampir-di-mana-mana ditambah besar setiap loncatan kali delta pada
titik loncatan. Karena itu
\[
 (L_f)'=L_{\sbsb\chi{(0,2)}}+2\delta_1-4\delta_2,
\]
yang memberikan rumus jawaban setelah diterapkan pada~$\phi$.
\end{o001proof}
\end{o001solution}

% O001-SOLUTION-ID: O001-FAOA-2015-CH10-EX-005-SOLUTION
% SOURCE-EXERCISE-ID: FAOA-2015-CH10-NODE-0038
% STATEMENT-TARGET-SHA256: f52569f7102a885e40cd395f3d064ee179df2aacdc336e54e891140188da963f
\begin{o001solution}
  {O001-FAOA-2015-CH10-EX-005-SOLUTION}
  {FAOA-2015-CH10-NODE-0038}
  {f52569f7102a885e40cd395f3d064ee179df2aacdc336e54e891140188da963f}
\begin{o001statement}
Misalkan $h(x) = \left\{
                           \begin{array}{ll}
                             2x, & \hbox{jika $x < -1$;} \\
                             1, & \hbox{jika $-1 \le x < 1$;} \\
                             0, & \hbox{jika $x \ge 1$.}
                           \end{array}
                         \right.$ Carilah suatu fungsi $f$ yang terintegralkan secara lokal sedemikian sehingga ${\tilde f}\,' = \tilde h + 3\delta_1 - \delta_{-1}$.
\end{o001statement}
\begin{o001answer}
Salah satu pilihan adalah
\[
 f(x)=
 \begin{cases}
 x^2,&x<-1,\\
 x+1,&-1\leq x<1,\\
 5,&x\geq1.
 \end{cases}
\]
Setiap solusi lain sama hampir di mana-mana dengan $f+C$ untuk suatu konstanta
global~$C$; pada tingkat distribusi, kebebasannya tepat satu konstanta global.
\end{o001answer}
\begin{o001proof}
Pada tiga interval terbuka yang bersangkutan, turunan klasik fungsi yang
diberikan berturut-turut adalah $2x$, $1$, dan~$0$, jadi bagian regulernya
adalah $\tilde h$. Di~$-1$, limit kanan dan kiri adalah
\[
 f(-1+)=0,
 \qquad
 f(-1-)=1,
\]
sehingga loncatannya $-1$. Di~$1$, limit kanan dan kiri adalah $5$ dan~$2$,
sehingga loncatannya~$3$. Oleh rumus turunan fungsi mulus-sepotong,
\[
 {\tilde f}\,'=\tilde h-\delta_{-1}+3\delta_1.
\]
Jika antiturunan pada ketiga interval diberi konstanta terpisah, syarat kedua
loncatan memaksa konstanta tersebut berbeda tepat seperti pada rumus di atas,
hingga penambahan satu konstanta yang sama pada semua bagian. Perubahan nilai
fungsi pada himpunan berukuran nol tidak mengubah distribusi reguler, sehingga
pernyataan keunikan untuk fungsi harus dipahami hampir di mana-mana.
\end{o001proof}
\end{o001solution}

% O001-SOLUTION-ID: O001-FAOA-2015-CH10-EX-006-SOLUTION
% SOURCE-EXERCISE-ID: FAOA-2015-CH10-NODE-0040
% STATEMENT-TARGET-SHA256: c64f5a0efcf4bce85f5462849dd1c213bdb4a6fde85b288a8891a0e99b5db78c
\begin{o001solution}
  {O001-FAOA-2015-CH10-EX-006-SOLUTION}
  {FAOA-2015-CH10-NODE-0040}
  {c64f5a0efcf4bce85f5462849dd1c213bdb4a6fde85b288a8891a0e99b5db78c}
\begin{o001statement}
Suatu \df{dipol} (dengan momen listrik 1 di titik asal pada $\R$) dapat dipandang sebagai ``limit'' ketika $\epsilon \sto 0^+$
 \index{dipol}%
dari suatu sistem $T_\epsilon$ yang terdiri atas dua muatan $-\frac1\epsilon$ dan $\frac1\epsilon$ yang masing-masing ditempatkan di $0$ dan $\epsilon$. Tunjukkan
bagaimana hal ini dapat mengarahkan kita untuk mendefinisikan dipol secara matematis sebagai $-\delta\,'$. \emph{Petunjuk.} Pandanglah $T_\epsilon$ sebagai
distribusi $\frac1\epsilon \delta_\epsilon - \frac1\epsilon \delta$.
\end{o001statement}
\begin{o001answer}
$T_\epsilon\to-\delta'$ dalam $\fml D^*(\R)$ ketika
$\epsilon\to0^+$.
\end{o001answer}
\begin{o001proof}
Untuk setiap $\phi\in\fml D(\R)$,
\[
 T_\epsilon(\phi)
 =\frac{\phi(\epsilon)-\phi(0)}{\epsilon}
 \longrightarrow \phi'(0).
\]
Menurut definisi turunan distribusi,
$\delta'(\phi)=-\delta(\phi')=-\phi'(0)$, sehingga
$(-\delta')(\phi)=\phi'(0)$. Jadi limit dua muatan tersebut, dalam arti
konvergensi titik demi titik pada setiap fungsi uji, tepat sama dengan
$-\delta'$.
\end{o001proof}
\end{o001solution}

% O001-SOLUTION-ID: O001-FAOA-2015-CH10-EX-007-SOLUTION
% SOURCE-EXERCISE-ID: FAOA-2015-CH10-NODE-0047
% STATEMENT-TARGET-SHA256: fd8775418583e8d49dd07a07086c23720b69fc2cd2a0946237eb620641240fcf
\begin{o001solution}
  {O001-FAOA-2015-CH10-EX-007-SOLUTION}
  {FAOA-2015-CH10-NODE-0047}
  {fd8775418583e8d49dd07a07086c23720b69fc2cd2a0946237eb620641240fcf}
\begin{o001statement}
Jika $\delta$ bukan suatu fungsi pada $\R$, lalu apa yang dimaksud orang ketika mereka menulis
   \[ \lim_{k \sto \infty} \frac{k}{\pi(1+k^2x^2)} = \delta(x)\quad? \]
Tunjukkan bahwa, jika ditafsirkan dengan tepat (sebagai suatu pernyataan tentang distribusi), ungkapan itu bahkan benar. \emph{Petunjuk.}
Jika $s_k(x) = k\,\pi^{-1}(1+k^2x^2)^{-1}$ dan $\phi$ suatu fungsi uji pada $\R$, maka
$\int_{-\infty}^\infty s_k\phi = \int_{-\infty}^\infty s_k\phi(0) + \int_{-\infty}^\infty s_k\eta$ dengan $\eta(x) = \phi(x) - \phi(0)$.
Tunjukkan bahwa $\int_{-\infty}^\infty s_k\eta \sto 0$ ketika $k \sto \infty$.
\end{o001statement}
\begin{o001answer}
Maknanya ialah $L_{s_k}\to\delta$ dalam $\fml D^*(\R)$, yaitu
\[
 \int_{\R}s_k(x)\phi(x)\,dx\longrightarrow\phi(0)
 \qquad\text{untuk setiap }\phi\in\fml D(\R).
\]
\end{o001answer}
\begin{o001proof}
Dengan substitusi $y=kx$ diperoleh
\[
 \int_{\R}s_k(x)\phi(x)\,dx
 =\frac1\pi\int_{\R}\frac{\phi(y/k)}{1+y^2}\,dy.
\]
Untuk setiap $y$, $\phi(y/k)\to\phi(0)$, sedangkan integrannya secara mutlak
dibatasi oleh
$\|\phi\|_\infty\big/\bigl(\pi(1+y^2)\bigr)$, sebuah fungsi yang
terintegralkan. Teorema konvergensi terdominasi memberi
\[
 \lim_{k\to\infty}\int_{\R}s_k\phi
 =\frac{\phi(0)}\pi\int_{\R}\frac{dy}{1+y^2}
 =\phi(0)=\delta(\phi).
\]
Pernyataan tersebut dengan demikian merupakan pernyataan tentang limit
distribusi, bukan limit titik demi titik fungsi.
\end{o001proof}
\end{o001solution}

% O001-SOLUTION-ID: O001-FAOA-2015-CH10-EX-008-SOLUTION
% SOURCE-EXERCISE-ID: FAOA-2015-CH10-NODE-0048
% STATEMENT-TARGET-SHA256: debe9140cb2f7d0021122f265e6c1e892dbe064fb8578e7594839762ed156c72
\begin{o001solution}
  {O001-FAOA-2015-CH10-EX-008-SOLUTION}
  {FAOA-2015-CH10-NODE-0048}
  {debe9140cb2f7d0021122f265e6c1e892dbe064fb8578e7594839762ed156c72}
\begin{o001statement}
Misalkan $f_n(x) = \sin nx$ untuk $n \in \N$ dan $x \in \R$. Barisan $(f_n)$ tidak konvergen secara klasik;
tetapi barisan itu konvergen secara distribusional (ke $0$). Nyatakan pernyataan ini dengan tepat, lalu buktikan. \emph{Petunjuk.} Carilah antiturunannya.
\end{o001statement}
\begin{o001answer}
Untuk setiap $\phi\in\fml D(\R)$,
$\int_{\R}\sin(nx)\phi(x)\,dx\to0$; yakni $L_{f_n}\to0$ dalam
$\fml D^*(\R)$.
\end{o001answer}
\begin{o001proof}
Karena $-\cos(nx)/n$ adalah antiturunan $\sin(nx)$ dan $\phi$ bertumpuan
kompak, integrasi per bagian menghasilkan
\[
 \int_{\R}\sin(nx)\phi(x)\,dx
 =\frac1n\int_{\R}\cos(nx)\phi'(x)\,dx.
\]
Oleh karena itu,
\[
 \left|\int_{\R}\sin(nx)\phi(x)\,dx\right|
 \leq\frac{\|\phi'\|_{L^1}}{n}\longrightarrow0.
\]
Inilah tepat definisi konvergensi distribusional ke distribusi nol.
\end{o001proof}
\end{o001solution}

% O001-SOLUTION-ID: O001-FAOA-2015-CH10-EX-009-SOLUTION
% SOURCE-EXERCISE-ID: FAOA-2015-CH10-NODE-0050
% STATEMENT-TARGET-SHA256: ec002d6bdad3dd72b45e60fa31af1c8ef543b8b53ad21cee638924f22b6cfd48
\begin{o001solution}
  {O001-FAOA-2015-CH10-EX-009-SOLUTION}
  {FAOA-2015-CH10-NODE-0050}
  {ec002d6bdad3dd72b45e60fa31af1c8ef543b8b53ad21cee638924f22b6cfd48}
\begin{o001statement}
Berikan makna pada ungkapan
   \[ 2\pi\sum_{n = -\infty}^\infty \delta(x-2\pi n) =  1 +  2\sum_{n = 1}^\infty \cos nx\,, \]
dan tunjukkan bahwa, jika ditafsirkan dengan tepat, ungkapan itu benar. \emph{Petunjuk.} Siapa pun yang cukup bejat untuk percaya bahwa $\delta(x)$ \emph{berarti}
sesuatu kemungkinan akan menafsirkan $\delta(x-a)$ sebagai hal yang sama dengan $\delta_a(x)$. Anda dapat memulai proses
merasionalisasikannya dengan menghitung deret Fourier fungsi $g$ yang didefinisikan oleh $g(x) = \frac1{4\pi}x^2 - \frac12 x$ pada $[0,2\pi]$
dan kemudian diperluas secara periodik ke $\R$. Anda mungkin perlu mencari teorema tentang konvergensi titik demi titik dan seragam dari deret Fourier
serta tentang pendiferensialan suku demi suku deret semacam itu.
\end{o001statement}
\begin{o001answer}
Dalam $\fml D^*(\R)$ berlaku
\[
 2\pi\sum_{k\in\Z}\delta_{2\pi k}
 =\lim_{N\to\infty}\left(1+2\sum_{n=1}^{N}\cos(nx)\right).
\]
Kedua ruas yang diterapkan pada $\phi$ bernilai
$2\pi\sum_{k\in\Z}\phi(2\pi k)$.
\end{o001answer}
\begin{o001proof}
Ambil $\phi\in\fml D(\R)$ dan periodisasikan fungsi itu dengan
\[
 P\phi(t)=\sum_{k\in\Z}\phi(t+2\pi k).
\]
Pada setiap titik hanya berhingga banyak suku yang tidak nol secara lokal,
sehingga $P\phi$ mulus dan $2\pi$-periodik. Koefisien Fourier-nya adalah
\[
 c_n=\frac1{2\pi}\int_0^{2\pi}P\phi(t)e^{-int}\,dt
 =\frac1{2\pi}\int_{\R}\phi(t)e^{-int}\,dt.
\]
Integrasi per bagian berulang menunjukkan bahwa $(c_n)$ menurun lebih cepat
daripada setiap pangkat $|n|^{-1}$. Jadi deret Fourier tersebut konvergen
mutlak dan seragam. Dengan mengevaluasinya di~$0$, kita memperoleh rumus
penjumlahan Poisson
\[
 2\pi\sum_{k\in\Z}\phi(2\pi k)
 =\sum_{n\in\Z}\int_{\R}\phi(t)e^{-int}\,dt.
\]
Penggantian $n$ oleh $-n$ tidak mengubah jumlah atas semua bilangan bulat,
sehingga ruas kanan juga sama dengan
\[
 \lim_{N\to\infty}\int_{\R}
 \left(\sum_{n=-N}^{N}e^{int}\right)\phi(t)\,dt
 =\lim_{N\to\infty}\int_{\R}
 \left(1+2\sum_{n=1}^{N}\cos(nt)\right)\phi(t)\,dt.
\]
Sementara itu, ruas kiri adalah penerapan distribusi
$2\pi\sum_{k\in\Z}\delta_{2\pi k}$ pada~$\phi$; jumlah ini berhingga untuk
setiap fungsi uji. Identitas berlaku untuk setiap~$\phi$, sehingga berlaku
dalam arti distribusi.
\end{o001proof}
\end{o001solution}

% O001-SOLUTION-ID: O001-FAOA-2015-CH10-EX-010-SOLUTION
% SOURCE-EXERCISE-ID: FAOA-2015-CH10-NODE-0086
% STATEMENT-TARGET-SHA256: 0f28bc54f7a5f5bffbd03ec67bc3d3d57cd2e3851430c8c294375d7fe333d618
\begin{o001solution}
  {O001-FAOA-2015-CH10-EX-010-SOLUTION}
  {FAOA-2015-CH10-NODE-0086}
  {0f28bc54f7a5f5bffbd03ec67bc3d3d57cd2e3851430c8c294375d7fe333d618}
\begin{o001statement}
Buktikan proposisi sebelumnya dengan asumsi bahwa $u$ dan $v$ \emph{keduanya} memiliki tumpuan kompak. \emph{Petunjuk.}
Gunakan $(u \ast v) \ast (\phi \ast \psi)$, dengan $\phi$ dan $\psi$ fungsi-fungsi uji, sambil mengingat bahwa konvolusi fungsi
bersifat komutatif.
\end{o001statement}
\begin{o001answer}
Jika $u$ dan $v$ bertumpuan kompak, maka $u\ast v=v\ast u$.
\end{o001answer}
\begin{o001proof}
Karena $u$ dan $v$ bertumpuan kompak, $u\ast\psi$ dan $v\ast\phi$
merupakan fungsi uji untuk setiap $\phi,\psi\in\fml D(\R^n)$. Dengan definisi
konvolusi dua distribusi, sifat asosiatif konvolusi distribusi--fungsi-uji,
dan komutativitas konvolusi fungsi, diperoleh
\begin{align*}
 (u\ast v)\ast(\phi\ast\psi)
 &=u\ast\bigl(v\ast(\phi\ast\psi)\bigr)\\
 &=u\ast\bigl((v\ast\phi)\ast\psi\bigr)\\
 &=u\ast\bigl(\psi\ast(v\ast\phi)\bigr)\\
 &=(u\ast\psi)\ast(v\ast\phi).
\end{align*}
Di sisi lain, karena $\phi\ast\psi=\psi\ast\phi$, perhitungan yang sama
dengan $u,v$ dan $\phi,\psi$ dipertukarkan memberi
\[
 (v\ast u)\ast(\phi\ast\psi)
 =(v\ast\phi)\ast(u\ast\psi).
\]
Dua ruas terakhir sama karena keduanya merupakan konvolusi fungsi uji.

Sekarang tetapkan $w=u\ast v-v\ast u$. Ambil sembarang
$\eta\in\fml D(\R^n)$ dan suatu aproksimasi identitas
$(\rho_\epsilon)$ dalam $\fml D(\R^n)$. Hasil di atas dengan
$\phi=\eta$ dan $\psi=\rho_\epsilon$ menunjukkan
\[
 w\ast(\eta\ast\rho_\epsilon)=0.
\]
Karena $\eta\ast\rho_\epsilon\to\eta$ dalam topologi fungsi uji dan distribusi
kontinu pada topologi itu, kita memperoleh $w\ast\eta=0$. Hal ini berlaku
untuk setiap~$\eta$. Proposisi ketunggalan tepat sebelum definisi tumpuan
kemudian memberi $w=0$, yaitu $u\ast v=v\ast u$.
\end{o001proof}
\end{o001solution}

% O001-SOLUTION-ID: O001-FAOA-2015-CH10-EX-011-SOLUTION
% SOURCE-EXERCISE-ID: FAOA-2015-CH10-NODE-0099
% STATEMENT-TARGET-SHA256: d7c771909a0b07a97813b77756de75190669a00de4df0212ffbcdc10f4c996f3
\begin{o001solution}
  {O001-FAOA-2015-CH10-EX-011-SOLUTION}
  {FAOA-2015-CH10-NODE-0099}
  {d7c771909a0b07a97813b77756de75190669a00de4df0212ffbcdc10f4c996f3}
\begin{o001statement}
Carilah suatu solusi distribusional bagi persamaan diferensial $x \dfrac{du}{dx} + u = 0$ pada~$\R$.
\end{o001statement}
\begin{o001answer}
Distribusi delta Dirac $u=\delta$ merupakan suatu solusi.
\end{o001answer}
\begin{o001proof}
Untuk setiap fungsi uji~$\phi$,
\[
 (x\delta')(\phi)=\delta'(x\phi)
 =-(x\phi)'(0)=-\phi(0)=-\delta(\phi).
\]
Jadi $x\delta'=-\delta$, dan akibatnya
\[
 x\frac{d\delta}{dx}+\delta=x\delta'+\delta=0
\]
dalam $\fml D^*(\R)$.
\end{o001proof}
\end{o001solution}

% Lapisan solusi O001 -- Bab 13 (Konstruksi GNS)
% 1 solusi asli terpisah; bukan tulisan John M. Erdman.
% Provenans model: OpenAI Codex gpt-5.6-sol, Ultra, atas arahan pengguna.
% Lisensi komponen solusi: CC BY-SA 4.0.
% Tidak menyiratkan dukungan John M. Erdman atau Portland State University.

\markboth{SOLUSI O001 UNTUK BAB 13: KONSTRUKSI GNS}{SOLUSI O001 UNTUK BAB 13: KONSTRUKSI GNS}
\section*{Solusi O001 untuk Bab 13: Konstruksi GNS}
\addcontentsline{toc}{section}{Solusi O001 untuk Bab 13: Konstruksi GNS}
\begin{center}
\small
Komponen solusi asli terpisah, CC BY-SA 4.0. Ditulis dengan bantuan
OpenAI Codex gpt-5.6-sol, Ultra, atas arahan pengguna. Pernyataan latihan
berasal dari edisi terjemahan karya John M. Erdman; solusi ini bukan tulisan
beliau dan tidak menyiratkan dukungan beliau atau Portland State University.
\end{center}

% O001-SOLUTION-ID: O001-FAOA-2015-CH13-EX-001-SOLUTION
% SOURCE-EXERCISE-ID: FAOA-2015-CH13-NODE-0018
% STATEMENT-TARGET-SHA256: 84c22fa77eaccb119f0537a7807cc91265e0dc1e0c2aae2137eadc8497924228
\begin{o001solution}
  {O001-FAOA-2015-CH13-EX-001-SOLUTION}
  {FAOA-2015-CH13-NODE-0018}
  {84c22fa77eaccb119f0537a7807cc91265e0dc1e0c2aae2137eadc8497924228}
\begin{o001statement}
Misalkan $X$ suatu ruang Hausdorff kompak lokal. Carilah suatu representasi
isometrik (dan karena itu setia) $(\pi,H)$ dari aljabar-$C^*$ $C_0(X)$ pada suatu
ruang Hilbert~$H$.
\end{o001statement}
\begin{o001answer}
Ambil $H=\ell^2(X)$ dan definisikan
$(\pi(f)\xi)(x)=f(x)\xi(x)$. Maka $\pi$ adalah representasi isometrik dan
setia dari $C_0(X)$ pada~$H$.
\end{o001answer}
\begin{o001proof}
Untuk setiap $f\in C_0(X)$ dan $\xi\in\ell^2(X)$, perkalian titik demi titik
memberi
\[
 \|\pi(f)\xi\|_2^2
 =\sum_{x\in X}|f(x)|^2|\xi(x)|^2
 \leq\|f\|_\infty^2\|\xi\|_2^2.
\]
Jadi $\pi(f)$ operator terbatas dan $\|\pi(f)\|\leq\|f\|_\infty$. Definisi
titik demi titik langsung memberi
\[
 \pi(fg)=\pi(f)\pi(g),
 \qquad
 \pi(\overline f)=\pi(f)^*,
\]
serta linearitas~$\pi$.

Untuk setiap $x\in X$, misalkan $e_x$ adalah vektor basis satuan yang bernilai
$1$ di~$x$ dan nol di tempat lain. Maka
\[
 \|\pi(f)e_x\|_2=|f(x)|.
\]
Akibatnya $\|\pi(f)\|\geq |f(x)|$ untuk setiap~$x$, dan dengan mengambil
supremum diperoleh $\|\pi(f)\|\geq\|f\|_\infty$. Bersama estimasi sebelumnya,
$\|\pi(f)\|=\|f\|_\infty$. Jadi $\pi$ isometrik; khususnya
$\pi(f)=0$ memaksa $f=0$, sehingga representasi itu setia. Konstruksi ini
berlaku tanpa perlu memilih ukuran pada~$X$, bahkan ketika $X$ tidak dapat
dihitung.
\end{o001proof}
\end{o001solution}

% Lapisan solusi O001 -- Bab 14 (Aljabar Pengali)
% 2 solusi asli terpisah; bukan tulisan John M. Erdman.
% Provenans model: OpenAI Codex gpt-5.6-sol, Ultra, atas arahan pengguna.
% Lisensi komponen solusi: CC BY-SA 4.0.
% Tidak menyiratkan dukungan John M. Erdman atau Portland State University.

\clearpage
\markboth{SOLUSI O001 UNTUK BAB 14: ALJABAR PENGALI}{SOLUSI O001 UNTUK BAB 14: ALJABAR PENGALI}
% AMSBook mengambil kepala halaman genap dari topmark halaman sebelumnya.
% Gaya satu halaman ini memastikan halaman pembuka Bab 14 memakai kepala
% Bab 14 tanpa memindahkan tanda tersebut ke halaman terakhir Bab 13.
\makeatletter
\def\ps@oonechapterxiv{%
  \ps@headings
  \def\@evenhead{%
    \setTrue{runhead}%
    \normalfont\scriptsize
    \rlap{\thepage}\hfil
    \def\thanks{\protect\thanks@warning}%
    SOLUSI O001 UNTUK BAB 14: ALJABAR PENGALI\hfil}%
  \def\@oddhead{%
    \setTrue{runhead}%
    \normalfont\scriptsize\hfil
    \def\thanks{\protect\thanks@warning}%
    SOLUSI O001 UNTUK BAB 14: ALJABAR PENGALI\hfil\llap{\thepage}}%
}
\makeatother
\thispagestyle{oonechapterxiv}
\section*{Solusi O001 untuk Bab 14: Aljabar Pengali}
\addcontentsline{toc}{section}{Solusi O001 untuk Bab 14: Aljabar Pengali}
\begin{center}
\small
Komponen solusi asli terpisah, CC BY-SA 4.0. Ditulis dengan bantuan
OpenAI Codex gpt-5.6-sol, Ultra, atas arahan pengguna. Pernyataan latihan
berasal dari edisi terjemahan karya John M. Erdman; solusi ini bukan tulisan
beliau dan tidak menyiratkan dukungan beliau atau Portland State University.
\end{center}

% O001-SOLUTION-ID: O001-FAOA-2015-CH14-EX-001-SOLUTION
% SOURCE-EXERCISE-ID: FAOA-2015-CH14-NODE-0005
% STATEMENT-TARGET-SHA256: cd4d713e827bc5f3452e91f669d72db723d317f124b43c853eda067d3b95539e
\begin{o001solution}
  {O001-FAOA-2015-CH14-EX-001-SOLUTION}
  {FAOA-2015-CH14-NODE-0005}
  {cd4d713e827bc5f3452e91f669d72db723d317f124b43c853eda067d3b95539e}
\begin{o001statement}
Pastikan bahwa definisi \emph{ruang vektor} sebelumnya ekuivalen
dengan definisi yang biasa Anda gunakan.
\end{o001statement}
\begin{o001answer}
Pemetaan $M$ dan perkalian skalar biasa saling menentukan melalui
$\alpha x=M(\alpha)(x)$; aksioma homomorfisme gelanggang beridentitas bagi
$M$ tepat ekuivalen dengan aksioma perkalian skalar.
\end{o001answer}
\begin{o001proof}
Andaikan mula-mula $(V,+,M)$ memenuhi definisi formal dan tetapkan
$\alpha x=M(\alpha)(x)$. Karena setiap $M(\alpha)$ merupakan endomorfisme
grup Abel~$(V,+)$,
\[
 \alpha(x+y)=\alpha x+\alpha y.
\]
Karena $M$ homomorfisme gelanggang,
\begin{align*}
 (\alpha+\beta)x
 &=M(\alpha+\beta)x
  =\bigl(M(\alpha)+M(\beta)\bigr)x
  =\alpha x+\beta x,\\
 (\alpha\beta)x
 &=M(\alpha\beta)x
  =M(\alpha)M(\beta)x
  =\alpha(\beta x),\\
 1x&=M(1)x=x.
\end{align*}
Grup Abel memberi seluruh aksioma penjumlahan vektor, dan ketiga identitas di
atas beserta aditivitas $M(\alpha)$ memberi seluruh aksioma perkalian skalar
kompleks yang biasa.

Sebaliknya, mulai dari ruang vektor kompleks dalam pengertian biasa dan
definisikan $M(\alpha)(x)=\alpha x$. Distributivitas terhadap penjumlahan
vektor menunjukkan $M(\alpha)\in\Hom(V,+)$. Distributivitas terhadap
penjumlahan skalar memberi $M(\alpha+\beta)=M(\alpha)+M(\beta)$,
asosiativitas perkalian skalar memberi
$M(\alpha\beta)=M(\alpha)M(\beta)$, dan aksioma satuan memberi
$M(1)=I_V$. Jadi $M$ homomorfisme gelanggang beridentitas. Kedua konstruksi
jelas saling invers.
\end{o001proof}
\end{o001solution}

% O001-SOLUTION-ID: O001-FAOA-2015-CH14-EX-002-SOLUTION
% SOURCE-EXERCISE-ID: FAOA-2015-CH14-NODE-0009
% STATEMENT-TARGET-SHA256: 167b8865f270fa360a46726a6ce2259e3433b08b036fc661bf498ef6242b93dc
\begin{o001solution}
  {O001-FAOA-2015-CH14-EX-002-SOLUTION}
  {FAOA-2015-CH14-NODE-0009}
  {167b8865f270fa360a46726a6ce2259e3433b08b036fc661bf498ef6242b93dc}
\begin{o001statement}
Periksa bahwa definisi \emph{modul-$A$} dalam~\ref{0038017}
dan~\ref{0038031} ekuivalen.
\end{o001statement}
\begin{o001answer}
Dengan perkalian operator standar, kedua definisi tidak ekuivalen persis
seperti tercetak: definisi~\ref{0038031} harus memakai antihomomorfisme
$\Phi\colon A\to\ofml L(V)$, atau secara ekuivalen homomorfisme
$A^{\mathrm{op}}\to\ofml L(V)$. Setelah koreksi itu, korespondensinya ialah
$\Phi(a)x=xa$.
\end{o001answer}
\begin{o001proof}
Mulailah dengan modul kanan dalam~\ref{0038017} dan definisikan operator
$\Phi(a)$ melalui
\[
 \Phi(a)x=xa.
\]
Bilinearitas aksi menunjukkan bahwa $\Phi(a)$ linear dan bahwa
$a\mapsto\Phi(a)$ linear. Akan tetapi, aksioma modul kanan memberi
\[
 \Phi(ab)x=x(ab)=(xa)b=\Phi(b)\Phi(a)x,
\]
jadi
\[
 \Phi(ab)=\Phi(b)\Phi(a).
\]
Dengan kata lain, $\Phi$ adalah antihomomorfisme
$A\to\ofml L(V)$, atau homomorfisme
$A^{\mathrm{op}}\to\ofml L(V)$. Jika $A$ beridentitas, syarat
$x1_A=x$ tepat sama dengan $\Phi(1_A)=I_V$.

Sebaliknya, andaikan $\Phi\colon A\to\ofml L(V)$ linear, beridentitas bila
diperlukan, dan memenuhi $\Phi(ab)=\Phi(b)\Phi(a)$. Tetapkan
$xa=\Phi(a)x$. Linearitas $\Phi$ dan setiap operator $\Phi(a)$ memberi
bilinearitas, sedangkan
\[
 (xa)b=\Phi(b)\Phi(a)x=\Phi(ab)x=x(ab).
\]
Syarat satuan juga langsung cocok. Kedua konstruksi ini saling invers, sehingga
memberi ekuivalensi yang benar.

Bahwa kata ``homomorfisme'' pada~\ref{0038031} tidak dapat dipertahankan untuk
aljabar tak komutatif dapat dilihat dari contoh $A=M_2(\C)$ dan $V=A$ dengan
aksi kanan biasa. Jika $R_a(x)=xa$, maka
$R_aR_b=R_{ba}$, bukan $R_{ab}$ pada umumnya. Misalnya, untuk
$a=E_{12}$ dan $b=E_{21}$, diperoleh $ab=E_{11}$ dan $ba=E_{22}$.
Jadi pernyataan tercetak memerlukan koreksi arah perkalian di atas; solusi ini
tidak menyamarkan ketidaksesuaian tersebut.
\end{o001proof}
\end{o001solution}

% Lapisan solusi O001 -- Bab 17 (Funktor K0)
% 1 solusi asli terpisah; bukan tulisan John M. Erdman.
% Provenans model: OpenAI Codex gpt-5.6-sol, Ultra, atas arahan pengguna.
% Lisensi komponen solusi: CC BY-SA 4.0.
% Tidak menyiratkan dukungan John M. Erdman atau Portland State University.

\markboth{SOLUSI O001 UNTUK BAB 17: FUNKTOR $K_0$}{SOLUSI O001 UNTUK BAB 17: FUNKTOR $K_0$}
\section*{Solusi O001 untuk Bab 17: Funktor $K_0$}
\addcontentsline{toc}{section}{Solusi O001 untuk Bab 17: Funktor $K_0$}
\begin{center}
\small
Komponen solusi asli terpisah, CC BY-SA 4.0. Ditulis dengan bantuan
OpenAI Codex gpt-5.6-sol, Ultra, atas arahan pengguna. Pernyataan latihan
berasal dari edisi terjemahan karya John M. Erdman; solusi ini bukan tulisan
beliau dan tidak menyiratkan dukungan beliau atau Portland State University.
\end{center}

% O001-SOLUTION-ID: O001-FAOA-2015-CH17-EX-001-SOLUTION
% SOURCE-EXERCISE-ID: FAOA-2015-CH17-NODE-0026
% STATEMENT-TARGET-SHA256: 76065558fdb1aff0f098a8215cec17050381cb2ca560da7994991909e0f3f8cd
\begin{o001solution}
  {O001-FAOA-2015-CH17-EX-001-SOLUTION}
  {FAOA-2015-CH17-NODE-0026}
  {76065558fdb1aff0f098a8215cec17050381cb2ca560da7994991909e0f3f8cd}
\begin{o001statement}
Jelaskan mengapa terminologi dalam definisi sebelumnya masuk akal.
\end{o001statement}
\begin{o001answer}
Jika $s^*s=1$, maka perkalian kiri oleh~$s$ mempertahankan norma; dalam setiap
representasi-$*$ setia, operator yang mewakili~$s$ juga merupakan isometri
dalam arti operator Hilbert.
\end{o001answer}
\begin{o001proof}
Untuk setiap $x$ dalam aljabar-$C^*$ beridentitas~$A$, identitas-$C^*$ memberi
\[
 \|sx\|^2
 =\|(sx)^*(sx)\|
 =\|x^*s^*sx\|
 =\|x^*x\|
 =\|x\|^2.
\]
Jadi operator perkalian kiri $L_s\colon x\mapsto sx$ mempertahankan norma.
Secara setara, jika $\pi\colon A\to\ofml L(H)$ suatu representasi-$*$ setia,
maka untuk setiap $\xi\in H$,
\[
 \|\pi(s)\xi\|^2
 =\langle\pi(s)^*\pi(s)\xi,\xi\rangle
 =\langle\pi(s^*s)\xi,\xi\rangle
 =\|\xi\|^2.
\]
Maka $\pi(s)$ adalah operator isometrik, yang menjelaskan terminologinya.
Syarat $s^*s=1$ tidak memaksa $ss^*=1$, sehingga isometri ini tidak harus
surjektif; geser unilateral adalah contoh standar.
\end{o001proof}
\end{o001solution}



\backmatter


 \clearpage
 \phantomsection
        %allows TOC links to find bibliography correctly
 \bibliographystyle {amsplain}
 \bibliography{functional_analysis_op_algs_bib}
     %\addcontentsline{toc}{chapter}{Bibliography}


 \cleardoublepage
 \phantomsection
         %allows TOC links to find index correctly
 \printindex
 \addcontentsline{toc}{chapter}{Index}




\end{document}
